% Μαθηματικά Α' Γυμνασίου -- Τεύχος Α
% Μεταγλώττιση: xelatex teyxos_a.tex (δύο φορές, για το περιεχόμενο)
\documentclass[11pt,twoside,openright]{book}
\usepackage{glossa}
\hypersetup{pdftitle={Μαθηματικά Α' Γυμνασίου},pdfauthor={Ιωάννης Ψωμιάδης}}

% ---------- περιεχόμενα ----------
\usepackage{tocloft}
\renewcommand{\cftpartfont}{\bfseries\color{gblue}}
\renewcommand{\cftchapfont}{\normalfont}
\renewcommand{\cftchappagefont}{\normalfont}
\setlength{\cftbeforepartskip}{1.4em}
\setlength{\cftbeforechapskip}{0.25em}
\renewcommand{\cftdot}{.}
\renewcommand{\cftchapleader}{\cftdotfill{\cftdotsep}}
\renewcommand{\cftpartpresnum}{}
\renewcommand{\contentsname}{Περιεχόμενα}
\renewcommand{\cftpartdotsep}{\cftnodots}
\setcounter{tocdepth}{1}
\setlength{\cftsecindent}{1.5em}\setlength{\cftsecnumwidth}{2.8em}\renewcommand{\cftchapfont}{\bfseries}\renewcommand{\cftchappagefont}{\bfseries}\setlength{\cftbeforesecskip}{0.1em}
\let\oldpart\part
\newcommand{\epanalipsi}[1]{\clearpage\markboth{Επανάληψη κεφαλαίου #1}{Επανάληψη κεφαλαίου #1}%
  \addcontentsline{toc}{section}{\protect\numberline{}Επανάληψη του κεφαλαίου}%
  \setcounter{exer}{0}\renewcommand{\exprefix}{Ε#1}\vspace*{14mm}{\color{gblue}\bfseries\large Επανάληψη}\par\vspace{3mm}%
  {\raggedright\chaptertitlefont Επανάληψη κεφαλαίου #1\par}\vspace{6mm}}

\begin{document}
\frontmatter
\pagestyle{empty}
% ---------- εξώφυλλο ----------
\thispagestyle{empty}
\begin{center}
\vspace*{26mm}
{\color{ggray}\small Η Γλώσσα των Μαθηματικών}\par\vspace{16mm}
{\bfseries\fontsize{24}{30}\selectfont Μαθηματικά Α’ Γυμνασίου\par}\vspace{9mm}
{\small Μέρος Α΄: Αριθμητική – Άλγεβρα \ · \ Μέρος Β΄: Γεωμετρία\par}\vspace{9mm}
{\color{gblue}\large Βοήθημα με νόημα: ασκήσεις που θέλεις να λύσεις\par}\vspace{40mm}
Ιωάννης Ψωμιάδης\par\vspace{2mm}
{\color{ggray}\scriptsize Σύμφωνα με το Σχολικό Βιβλίο της Α’ Γυμνασίου\par}\vspace{6mm}
{\small Οκτώβριος 2026\par}
\end{center}
\cleardoublepage
\pagestyle{fancy}
\fancyhead[RO]{\color{ggray}\small Πριν ξεκινήσεις}
\fancyhead[LE]{\color{ggray}\small Η Γλώσσα των Μαθηματικών}
\pagenumbering{roman}

\chapter*{}\vspace{-22mm}
{\color{gblue}\bfseries\Large Πριν ξεκινήσεις}\par\vspace{4mm}
{\color{gblue}\bfseries\large Ένα γράμμα σε σένα}\par

Μπήκες στο γυμνάσιο. Άλλο κτίριο, άλλα πρόσωπα, ένα πρόγραμμα με διαφορετικό μάθημα σε κάθε ώρα και έναν καθηγητή για καθένα. Αν έχεις διαβάσει τα βιβλία της Ε’ και της ΣΤ’, θα αναγνωρίσεις τα παιδιά, τον τρόπο και τη φιλοσοφία. Αν όχι, δεν πειράζει: όλα όσα χρειάζεσαι ξαναλέγονται εδώ, όπου χρειάζεται.

Τα περισσότερα βιβλία μαθηματικών σου λένε πώς λύνεται μια άσκηση. Αυτό το βιβλίο θέλει να σου πει κάτι άλλο: γιατί αξίζει να τη λύσεις.

Ξέρεις ήδη να μιλάς ελληνικά. Ίσως λίγο αγγλικά. Με αυτές τις γλώσσες περιγράφεις τον κόσμο: τι βλέπεις, τι θέλεις, τι σε ενοχλεί, τι σε κάνει να γελάς. Τα μαθηματικά είναι κι αυτά μια γλώσσα. Περιγράφουν τον ίδιο κόσμο, απλώς με άλλες λέξεις. Όταν γράφεις $5>3$, λες την ίδια πρόταση με το «το πέντε είναι μεγαλύτερο από το τρία». Το σύμβολο $>$ είναι μια λέξη. Το $=$ είναι μια λέξη. Ακόμα και ο αριθμός 1.000.000 είναι μια λέξη, που στα ελληνικά τη λέμε «ένα εκατομμύριο».

Στο δημοτικό έμαθες λέξεις. Στην ΣΤ’ το λεξιλόγιό σου μεγάλωσε. Φέτος μαθαίνεις γραμματική: τους κανόνες που λένε ποιες προτάσεις επιτρέπονται και ποιες όχι. Γιατί το $-3$ υπάρχει, αφού δεν μπορείς να έχεις «μείον τρία μήλα»; Γιατί $(-1)\cdot(-1)=1$; Γιατί ισχύει το «κάθε τρίγωνο έχει άθροισμα γωνιών 180°» και πώς το ξέρουμε χωρίς να τα μετρήσουμε όλα; Φέτος για πρώτη φορά δεν θα σε ρωτήσω μόνο «πόσο κάνει;». Θα σε ρωτήσω και «γιατί;».

Σε αυτό το βιβλίο θα μάθεις τη γλώσσα των μαθηματικών όπως μαθαίνει κανείς κάθε γλώσσα: λέξη λέξη, κανόνα κανόνα, με παραδείγματα από τη ζωή σου και με πολλή εξάσκηση. Δεν θα σου ζητήσω να αποστηθίσεις κανόνες που δεν καταλαβαίνεις. Κάθε κανόνας θα έχει έναν λόγο ύπαρξης, και θα τον δεις πριν δεις τον κανόνα.

Ένα ακόμη. Οι ασκήσεις εδώ δεν σου ζητούν να λύσεις κάτι επειδή έτσι πρέπει. Σου ζητούν να βοηθήσεις κάποιον, να κερδίσεις ένα στοίχημα, να αποδείξεις ότι έχεις δίκιο, να λύσεις μια διαφωνία. Τα μαθηματικά είναι το εργαλείο. Ο λόγος είναι πάντα δικός σου.

{\color{gblue}\bfseries\large Η παρέα}\par

Το Σεπτέμβρη η παρέα της ΣΤ’ βρέθηκε στο γυμνάσιο, στο ίδιο τμήμα, την Α3. Μαζί της βρέθηκαν και δύο παιδιά από το άλλο δημοτικό της περιοχής, και ένα κορίτσι που δεν ήξερε κανέναν. Οι δύο γειτονιές έγιναν μία τάξη. Η παρέα αυτή θα σε συνοδεύει σε όλο το γυμνάσιο, και τα προβλήματά της θα μεγαλώνουν μαζί της.

\begin{itemize}[leftmargin=1.4em,itemsep=0.2ex]
\item Η \textbf{Ζωή} ρωτάει πάντα «γιατί;» και δεν σταματάει στην πρώτη απάντηση.
\item Ο \textbf{Ηλίας} μετράει, δοκιμάζει, φτιάχνει πίνακες. Όταν κάνει λάθος, το βρίσκει μόνος του, αρκεί να τον αφήσεις.
\item Η \textbf{Μυρτώ} έχει το πρακτικό μυαλό της κουζίνας της γιαγιάς: αν μια απάντηση δεν βγάζει νόημα στην πραγματικότητα, δεν τη δέχεται.
\item Ο \textbf{Στέλιος} μαθαίνει έναν κανόνα και τον εφαρμόζει με ενθουσιασμό παντού, ακόμα και εκεί που δεν ταιριάζει.
\item Η \textbf{Φανή} κρατάει σημειώσεις για όλους και ρωτάει όσα οι άλλοι δεν τολμούν.
\item Ο \textbf{Βασίλης}, ο πρώην πρόεδρος του μαθητικού συμβουλίου, έχει πάντα άποψη και την υποστηρίζει με σιγουριά, κάποτε και με περισσότερη απ' όση χρειάζεται.
\item Ο \textbf{Νίκος}, από το άλλο δημοτικό, παίζει μπάσκετ και βιάζεται. Ό,τι γίνεται γρήγορα, το κάνει πρώτος. Ό,τι γίνεται σωστά, το κάνει δεύτερος.
\item Η \textbf{Αγγελική}, από το άλλο δημοτικό, ζωγραφίζει και παίζει μουσική. Βλέπει σχήματα και μοτίβα εκεί που οι άλλοι βλέπουν μόνο αριθμούς.
\item Η \textbf{Πολυξένη} είναι η καινούργια. Ήρθε το καλοκαίρι από τη Θεσσαλονίκη. Κρατάει ένα τετράδιο με κωδικούς και μοτίβα, μιλάει λίγο και ρωτάει σπάνια. Όταν ρωτάει, η ερώτησή της είναι αυτή που χρειαζόταν να ακουστεί.
\item Η \textbf{γιαγιά Ρήνη}, ο \textbf{παππούς} και ο \textbf{θείος Λευτέρης} είναι πάντα εκεί, με τον κήπο, το χωριό και τα πλακάκια τους.
\end{itemize}

Και υπάρχει η \textbf{κυρία Αλεξάνδρα}, η καθηγήτρια των μαθηματικών τους. Δεν δίνει ποτέ έτοιμη απάντηση. Δίνει ερώτηση.

\clearpage
{\color{gblue}\bfseries\large Πώς είναι φτιαγμένο το βιβλίο}\par

Όλες οι παράγραφοι έχουν την ίδια σειρά. Μόλις τη μάθεις, θα ξέρεις πάντα πού βρίσκεσαι.

\begin{enumerate}[leftmargin=2em,itemsep=0.5ex]
\item \textbf{Η ανάγκη.} Μια μικρή ιστορία με ένα πρόβλημα που δεν λύνεται χωρίς το εργαλείο του κεφαλαίου. Πρώτα νιώθεις ότι κάτι λείπει. Μετά σου το δίνω.
\item \textbf{Η λέξη.} Η νέα έννοια στα ελληνικά: τι σημαίνει, πού τη συναντάς. Μετά ο ακριβής ορισμός της.
\item \textbf{Η μετάφραση.} Ένας πίνακας με δύο στήλες. Αριστερά μια πρόταση στα ελληνικά, δεξιά η ίδια πρόταση στα μαθηματικά.
\item \textbf{Η θεωρία.} Οι κανόνες, με όλη την ακρίβεια που χρειάζονται. Είναι γραμμένοι σύμφωνα με το Σχολικό Βιβλίο σου, ώστε να βρίσκεις αμέσως την αντίστοιχη παράγραφο στην τάξη.
\item \textbf{Λυμένα παραδείγματα.} Κάθε λύση γράφεται ως σκέψη, σε τρία βήματα: Τι ξέρω. Τι θέλω. Πώς τα συνδέω. Στο τέλος ελέγχουμε αν η απάντηση βγάζει νόημα.
\item \textbf{Ασκήσεις σε τρεις ζώνες.}
  \begin{itemize}[leftmargin=1.2em,itemsep=0.2ex]
  \item \emph{Το κατάλαβα:} εφαρμόζεις τον κανόνα κατευθείαν.
  \item \emph{Το χρησιμοποιώ:} λύνεις ένα πρόβλημα από τη ζωή.
  \item \emph{Το σκέφτομαι:} μια ερώτηση που σε βάζει να σκεφτείς πιο μακριά, συχνά προς το επόμενο κεφάλαιο.
  \end{itemize}
\item \textbf{Λύσεις.} Όλες στο τέλος του βιβλίου, γραμμένες με τον ίδιο τρόπο σκέψης. Κοίτα τες αφού προσπαθήσεις, όχι πριν.
\end{enumerate}

Το βιβλίο χωρίζεται σε δύο μέρη, όπως και το Σχολικό σου: \textbf{Μέρος Α΄, Αριθμητική – Άλγεβρα}, και \textbf{Μέρος Β΄, Γεωμετρία}. Κάθε μέρος έχει κεφάλαια και κάθε κεφάλαιο \emph{παραγράφους} με την ίδια αρίθμηση που έχουν στο Σχολικό: η παράγραφος 1.4 του βιβλίου αυτού αντιστοιχεί στην Α.1.4 του Σχολικού. Οι ασκήσεις αριθμούνται με την παράγραφο: η άσκηση 1.4.7 είναι η έβδομη της παραγράφου 1.4.

Στο τέλος κάθε κεφαλαίου θα βρεις μια \textbf{Επανάληψη} με έναν «χάρτη»: ένα πίνακα που σου λέει σε ποιες παραγράφους να γυρίσεις, αν κάποια άσκηση σε δυσκόλεψε.

{\color{gblue}\bfseries\large Για τον γονέα}\par

Το βιβλίο αυτό γράφτηκε ώστε να το διαβάζετε μαζί με το παιδί σας. Δεν χρειάζεται να θυμάστε τα μαθηματικά του γυμνασίου.

Ξεκινήστε από την «Ανάγκη» και τη «Λέξη». Είναι γραμμένες σε απλά ελληνικά. Αν τις καταλάβετε εσείς, μπορείτε να τις εξηγήσετε. Η «Μετάφραση» δείχνει πώς η ίδια ιδέα γράφεται με σύμβολα, και εκεί συνήθως κολλάνε τα παιδιά: όχι στην ιδέα, αλλά στη μετάφρασή της.

Τρεις συμβουλές από δέκα χρόνια διδασκαλίας:

\begin{itemize}[leftmargin=1.4em,itemsep=0.4ex]
\item Ζητήστε από το παιδί να σας εξηγήσει μια λύση με δικά του λόγια. Αν μπορεί να την πει στα ελληνικά, την έχει καταλάβει.
\item Ένα λάθος είναι πληροφορία, όχι αποτυχία. Ρωτήστε «πώς το σκέφτηκες;» πριν πείτε «είναι λάθος».
\item Οι ασκήσεις της ζώνης «Το σκέφτομαι» δεν είναι υποχρεωτικές. Είναι για τις μέρες που υπάρχει όρεξη.
\end{itemize}

{\color{gblue}\bfseries\large Ποιος το έγραψε}\par

Σπούδασα μαθηματικά και δίδαξα ως φροντιστής για περίπου δέκα χρόνια. Τα επόμενα χρόνια δούλεψα στη βιομηχανία, κάνοντας μη καταστροφικούς ελέγχους: ελέγχω σωλήνες, δοχεία, δεξαμενές και συγκολλήσεις χωρίς να τα κόψουμε, για να ξέρουμε αν είναι ασφαλή.

Σε αυτά τα χρόνια έμαθα κάτι που δεν μου το είχε πει κανένα σχολικό βιβλίο. Τα μαθηματικά δεν είναι μάθημα. Είναι ο τρόπος που οι άνθρωποι συμφωνούν για το τι συμβαίνει. Όταν δύο μηχανικοί διαφωνούν, δεν κερδίζει αυτός που φωνάζει πιο δυνατά. Κερδίζει αυτός που φέρνει τη μέτρηση.

Αυτό το βιβλίο είναι η προσπάθειά μου να σου δώσω νωρίς αυτό το εργαλείο. Και, αν τα καταφέρω, λίγη από την αγάπη που έχω για αυτή τη γλώσσα.

\hfill Ιωάννης Ψωμιάδης

\cleardoublepage
\fancyhead[RO]{\color{ggray}\small Περιεχόμενα}
{\color{gblue}\bfseries\Large Περιεχόμενα}\par\vspace{3mm}
\makeatletter
\@starttoc{toc}
\makeatother


\mainmatter
\pagestyle{fancy}
\fancyhead[RO]{\color{ggray}\small\leftmark}
\meros{Α}{ΑΡΙΘΜΗΤΙΚΗ – ΑΛΓΕΒΡΑ}{}
\kefalaio{Οι φυσικοί αριθμοί}{Από το λεξιλόγιο στη γραμματική}{%
\begin{tabular}{@{}l@{\hspace{1em}}p{0.82\linewidth}@{}}
1.1 & Φυσικοί αριθμοί, διάταξη, στρογγυλοποίηση\\[2pt]
1.2 & Πρόσθεση, αφαίρεση και πολλαπλασιασμός φυσικών αριθμών\\[2pt]
1.3 & Δυνάμεις φυσικών αριθμών\\[2pt]
1.4 & Ευκλείδεια διαίρεση, διαιρετότητα\\[2pt]
1.5 & Χαρακτήρες διαιρετότητας, ΜΚΔ, ΕΚΠ, ανάλυση σε γινόμενο πρώτων παραγόντων\\[2pt]
 & Επανάληψη του κεφαλαίου\\
\end{tabular}}
\paragrafos{Φυσικοί αριθμοί, διάταξη, στρογγυλοποίηση: η σειρά που δεν τελειώνει}{Α.1.1, «Φυσικοί αριθμοί – Διάταξη φυσικών – Στρογγυλοποίηση»}

\begin{anagke}
Πρώτη μέρα στο γυμνάσιο. Στην αίθουσα της Α3 η Ζωή, ο Ηλίας, η Μυρτώ, ο Στέλιος, η Φανή και ο Βασίλης, ο πρώην πρόεδρος του μαθητικού συμβουλίου της ΣΤ’, βρίσκουν στα θρανία δύο γνωστά πρόσωπα από το άλλο δημοτικό της περιοχής: τον Νίκο και την Αγγελική, που τους έχουν δει σε αγώνες μπάσκετ και σε γιορτές. Στο τελευταίο θρανίο κάθεται ένα κορίτσι που δεν το ξέρει κανείς. Την λένε Πολυξένη, ήρθε το καλοκαίρι από τη Θεσσαλονίκη και δεν έχει πει ακόμα ούτε λέξη.

Στην είσοδο ένα χαρτί λέει: «Το σχολείο μας έχει 387 μαθητές». Ο Νίκος το φωτογραφίζει και στέλνει στη μαμά του: «Το σχολείο έχει 400 παιδιά!». Η μαμά του απαντάει με ένα χαμογελαστό «στρογγύλεψες».

Η κυρία Αλεξάνδρα, η καθηγήτρια των μαθηματικών, δεν ανοίγει βιβλίο. Γράφει στον πίνακα δύο αριθμούς: τους μαθητές του σχολείου, 387, και το υψόμετρο του Ολύμπου, 2.917 μέτρα. «Έχετε να τους πείτε σε έναν φίλο σας. Ποιον θα πείτε ακριβώς και ποιον περίπου;»

«Όλους ακριβώς», λέει ο Στέλιος. «Τι νόημα έχει το “περίπου”;» Η Ζωή σηκώνει το χέρι. «Αν πεις στον φίλο σου ότι περπάτησες 8.412 βήματα, θα θυμάται μόνο το 8 και θα μπερδεύεται με τα υπόλοιπα. Το “περίπου 8.000” το κρατάς». Ο Βασίλης χαμογελάει: «Εγώ λέω πάντα “γύρω στα τετρακόσια”. Με το μάτι». Τότε ακούγεται για πρώτη φορά η φωνή της Πολυξένης, χαμηλά: «Και πώς ξέρουμε ότι όλοι θα βγάλουμε το ίδιο “περίπου”; Αν ο Νίκος κάνει το 387 τετρακόσια, και ο Βασίλης “γύρω στα τριακόσια ογδόντα”;»

Όλοι γυρίζουν να την κοιτάξουν. Η κυρία Αλεξάνδρα χαμογελάει: η ερώτηση είναι ακριβώς αυτή που ήθελε να ακούσει. Πότε επιτρέπεται το «περίπου», και πώς γίνεται με κανόνα, ώστε όλοι να κάνουμε το ίδιο;
\end{anagke}

\begin{lexi}{φυσικός, στρογγυλοποίηση}{φυσικός, στρογγυλοποίηση}
Οι \lexiword{φυσικοί} αριθμοί λέγονται έτσι επειδή είναι οι πρώτοι που βρίσκει κανείς στη φύση, όταν μετράει πράγματα: μήλα, παιδιά, βήματα. Κανείς δεν μέτρησε ποτέ «ενάμισι παιδί». Το μηδέν μπήκε στην παρέα αργότερα. Στο σχολείο θα το βρεις μέσα στους φυσικούς.

Η \lexiword{στρογγυλοποίηση} είναι το να κάνεις έναν αριθμό «στρογγυλό», δηλαδή να τον αλλάξεις με έναν κοντινό του που τελειώνει σε μηδενικά: 387 γίνεται 400. Στρογγυλός σημαίνει τακτοποιημένος, χωρίς αιχμές. Δεν είναι μαγικό: είναι \emph{συμφωνία}. Όλοι κάνουμε τα ίδια βήματα και βγάζουμε το ίδιο αποτέλεσμα.

Δύο ακόμα λέξεις που θα ξανασυναντήσεις: \lexiword{άρτιος} είναι ο αριθμός που «ταιριάζει» τέλεια σε ζευγάρια, και \lexiword{περιττός} αυτός που «περισσεύει» ένας. Τους λέμε και ζυγούς και μονούς.
\end{lexi}

\begin{metafrasi}
Ο αριθμός που ακολουθεί το 9. & $9+1=10$\\
Το πέντε είναι μεγαλύτερο από το τρία. & $5>3$\\
Ο αριθμός $x$ βρίσκεται ανάμεσα στο 3 και στο 8. & $3<x<8$\\
Τρία εκατομμύρια, σαράντα χιλιάδες και επτά. & 3.040.007\\
Το 3.486.917 στρογγυλοποιημένο στις χιλιάδες. & $3.486.917\approx 3.487.000$\\
Ο αριθμός είναι άρτιος. & Διαιρείται με το 2\\
\end{metafrasi}

\tmima{Η θεωρία}

\begin{orismos}{φυσικοί αριθμοί}
Οι αριθμοί $0,\,1,\,2,\,3,\,4,\,\dots,\,99,\,100,\,\dots$ λέγονται \textbf{φυσικοί αριθμοί}. Κάθε φυσικός έχει έναν επόμενο και, εκτός από το 0, έναν προηγούμενο.
\end{orismos}

Οι φυσικοί χωρίζονται σε δύο κατηγορίες: \textbf{άρτιους} (ή ζυγούς), που διαιρούνται με το 2, και \textbf{περιττούς} (ή μονούς), που δεν διαιρούνται. Ο άρτιος τελειώνει πάντα σε 0, 2, 4, 6 ή 8.

\tmimasmall{Δέκα ψηφία, άπειροι αριθμοί}

Με τα δέκα ψηφία $0,\,1,\,2,\,\dots,\,9$ γράφουμε όσους αριθμούς θέλουμε. Αυτό γίνεται επειδή η αξία ενός ψηφίου εξαρτάται και από τη \textbf{θέση} του, δηλαδή από τη \textbf{δεκαδική τάξη}: μονάδες, δεκάδες, εκατοντάδες, χιλιάδες, δεκάδες χιλιάδων, εκατοντάδες χιλιάδων, εκατομμύρια. Το 0 κρατάει μια θέση άδεια: στο 6.040.205 τα τρία μηδενικά λένε ότι δεν υπάρχουν εκατοντάδες χιλιάδων, χιλιάδες και δεκάδες.

\tmimasmall{Διάταξη: ποιος είναι μεγαλύτερος}

Γράφουμε $=$ για «ίσος με», $<$ για «μικρότερος από» και $>$ για «μεγαλύτερος από». Δύο φυσικούς μπορούμε πάντα να τους συγκρίνουμε.

\begin{kanonas}{πώς συγκρίνω δύο φυσικούς}
\begin{enumerate}[leftmargin=1.6em,itemsep=0.2ex,topsep=2pt]
\item Μετράω τα ψηφία. Ο αριθμός με περισσότερα ψηφία είναι μεγαλύτερος (αν δεν ξεκινάει με 0).
\item Αν έχουν ίδιο πλήθος ψηφίων, τα κοιτάζω από αριστερά προς τα δεξιά. Στο πρώτο ψηφίο που διαφέρουν, μεγαλύτερος είναι αυτός με το μεγαλύτερο ψηφίο.
\end{enumerate}
\end{kanonas}

Έτσι $73.409>73.094$, γιατί τα δύο πρώτα ψηφία ταιριάζουν και στο τρίτο το 4 είναι μεγαλύτερο από το 0. Και $812<10.250$, γιατί το 812 έχει λιγότερα ψηφία. Τους φυσικούς μπορούμε να τους βάλουμε σε σειρά: \textbf{αύξουσα}, από τον μικρότερο στον μεγαλύτερο, ή \textbf{φθίνουσα}, αντίστροφα. Δεν υπάρχει μεγαλύτερος φυσικός αριθμός, γιατί κάθε φυσικός έχει επόμενο. Ο μικρότερος είναι το 0.

\tmimasmall{Η αριθμητική ευθεία}

Η διάταξη μας επιτρέπει να βάλουμε τους φυσικούς πάνω σε μια ευθεία. Διαλέγουμε ένα σημείο Ο για το 0 (την αρχή) και ένα σημείο Α δεξιά του για το 1. Το τμήμα ΟΑ είναι η \textbf{μονάδα μέτρησης}. Με το ίδιο μήκος βρίσκουμε τα σημεία των 2, 3, 4, \dots

\begin{center}
\begin{tikzpicture}[x=1.15cm,y=1cm]
\draw[-{Stealth}] (-0.3,0)--(9.6,0);
\foreach \i in {0,...,8}{\draw (\i,-0.12)--(\i,0.12); \node[below,font=\small] at (\i,-0.12) {\i};}
\node[above,font=\small\bfseries,gblue] at (0,0.15) {Ο};
\node[above,font=\small\bfseries,gblue] at (1,0.15) {Α};
\draw[<->,gblue] (0,0.75)--(1,0.75) node[midway,above,font=\scriptsize] {μονάδα};
\end{tikzpicture}
\end{center}

\tmimasmall{Στρογγυλοποίηση}

\begin{kanonas}{στρογγυλοποίηση φυσικού}
\begin{enumerate}[leftmargin=1.6em,itemsep=0.2ex,topsep=2pt]
\item Βρίσκω την τάξη στην οποία θα στρογγυλοποιήσω (δεκάδες, εκατοντάδες, χιλιάδες, \dots).
\item Κοιτάζω το ψηφίο της \textbf{αμέσως μικρότερης} τάξης.
\item Αν είναι $0,\,1,\,2,\,3$ ή $4$, το ψηφίο της τάξης μου μένει ίδιο. Αν είναι $5,\,6,\,7,\,8$ ή $9$, αυξάνεται κατά 1.
\item Όλα τα ψηφία δεξιά από την τάξη μου γίνονται 0.
\end{enumerate}
\end{kanonas}

Το αποτέλεσμα λέγεται \textbf{προσέγγιση} του αριθμού και γράφεται με το σύμβολο $\approx$ (διαβάζεται «περίπου ίσο»). Για παράδειγμα, $2.917\approx 3.000$ στις χιλιάδες και $2.917\approx 2.900$ στις εκατοντάδες.

Η στρογγυλοποίηση ταιριάζει εκεί που μας ενδιαφέρει η τάξη μεγέθους, όπως το πλήθος των ανθρώπων σε μια πόλη. Δεν ταιριάζει εκεί που χρειάζεται ακρίβεια: ο αριθμός του τηλεφώνου, το ΑΦΜ, ο κωδικός ενός λουκέτου.

\begin{prosoxi}
\begin{itemize}[leftmargin=1.2em,itemsep=0.3ex,topsep=2pt]
\item Στρογγυλοποιούμε \textbf{μία φορά}, από τον αρχικό αριθμό. Το 1.449 στις εκατοντάδες είναι 1.400, γιατί το ψηφίο των δεκάδων είναι 4. Αν το κάνεις σε δύο βήματα, πρώτα στις δεκάδες (1.450) και μετά στις εκατοντάδες, βγάζεις 1.500, που είναι λάθος.
\item Όταν το ψηφίο που αυξάνεται είναι 9, η αύξηση «κουβαλάει»: $9.995\approx 10.000$ στις δεκάδες.
\item Όταν μετράς \textbf{από ως}, και τα δύο άκρα μετράνε. Από τη σελίδα 20 ως τη σελίδα 47 είναι $47-20+1=28$ σελίδες, όχι 27.
\end{itemize}
\end{prosoxi}

\tmima{Λυμένα παραδείγματα}

\paradeigma{Το ντουλάπι του Ηλία}
\tixero Το ντουλάπι του Ηλία ανοίγει με τριψήφιο κωδικό. Θυμάται ότι έχει τα ψηφία 3, 0 και 5, καθένα μία φορά, αλλά όχι με ποια σειρά.\par
\tithelo Όλους τους πιθανούς κωδικούς, σε αύξουσα σειρά.\par
\possindeo Δοκιμάζω όλες τις σειρές των τριών ψηφίων: 305, 350, 035, 053, 503, 530. Ο 035 και ο 053 δεν είναι τριψήφιοι αριθμοί, γιατί το 0 μπροστά δεν μετράει. Τους σβήνω.\par
\apantisi 305, 350, 503, 530.\par
\elegxos Είναι τέσσερις, και κανένας δεν ξεκινάει από 0. Ο μικρότερος είναι ο 305: ξεκινάει από το μικρότερο μη μηδενικό ψηφίο, το 3, και μετά βάζει το μικρότερο που μένει, το 0.

\paradeigma{Οι προβολές του Βασίλη}
\tixero Το βίντεο της σχολικής εκλογικής εκστρατείας του Βασίλη έχει 3.486.917 προβολές. (Ναι, ο Βασίλης υπερβάλλει, αλλά ας το δεχτούμε.)\par
\tithelo Να το στρογγυλοποιήσω (α) στις εκατοντάδες, (β) στις χιλιάδες, (γ) στις εκατοντάδες χιλιάδων.\par
\possindeo
(α) Τάξη: εκατοντάδες. Η αμέσως μικρότερη τάξη είναι οι δεκάδες, με ψηφίο 1. Επειδή $1<5$, το 9 μένει ίδιο και τα δεξιά γίνονται 0: $3.486.917\approx 3.486.900$.\par
(β) Τάξη: χιλιάδες, ψηφίο 6. Η εκατοντάδα έχει ψηφίο 9. Επειδή $9\ge 5$, το 6 γίνεται 7: $3.486.917\approx 3.487.000$.\par
(γ) Τάξη: εκατοντάδες χιλιάδων, ψηφίο 4. Οι δεκάδες χιλιάδων έχουν ψηφίο 8. Επειδή $8\ge 5$, το 4 γίνεται 5: $3.486.917\approx 3.500.000$.\par
\apantisi 3.486.900, 3.487.000 και 3.500.000.\par
\elegxos Στο (β) το 3.486.917 απέχει 83 από το 3.487.000 και 917 από το 3.486.000. Το πρώτο είναι πιο κοντά. Στο (γ) απέχει 13.083 από το 3.500.000 και 86.917 από το 3.400.000. Και πάλι το 3.500.000 είναι πιο κοντά.

\paradeigma{Πόσες σελίδες;}
\tixero Η Φανή διαβάζει από τη σελίδα 58 ως τη σελίδα 131 του βιβλίου της, και τις δύο μαζί.\par
\tithelo Πόσες σελίδες διαβάζει.\par
\possindeo Αν έβγαζα $131-58=73$, θα είχα αφήσει έξω την πρώτη σελίδα. Για να μετρήσω και τα δύο άκρα, προσθέτω 1: $131-58+1=74$.\par
\apantisi 74 σελίδες.\par
\elegxos Δοκιμή με μικρούς αριθμούς: από τη σελίδα 3 ως τη 5 είναι τρεις σελίδες (3, 4, 5). Ο τύπος δίνει $5-3+1=3$. Σωστά.

\tmima{Ασκήσεις}

\zoni{Το κατάλαβα}{εφαρμογή του κανόνα}

\ex Γράψε με ψηφία: α) τριακόσια δεκαοκτώ \quad β) εννιακόσια ένα \quad γ) δεκαπέντε χιλιάδες διακόσια τέσσερα \quad δ) δύο εκατομμύρια τριακόσιες χιλιάδες εβδομήντα.

\ex Γράψε με λέξεις: α) 64.307 \quad β) 1.050.009 \quad γ) 200.020.

\ex Ποιοι είναι οι τρεις προηγούμενοι αριθμοί του 1.000 και ποιοι οι δύο επόμενοι του 999.999; Ποιος είναι ο προηγούμενος του 0;

\ex Τα σκορ από το παιχνίδι του Βασίλη είναι 7.250, 7.205, 72.050, 7.520 και 5.270. Γράψε τα σε αύξουσα σειρά και μετά σε φθίνουσα.

\ex Βάλε το σωστό σύμβολο ($<$, $=$ ή $>$): α) $603\ \square\ 630$ \quad β) $9.099\ \square\ 9.100$ \quad γ) $45.012\ \square\ 45.021$ \quad δ) $100.000\ \square\ 99.999$ \quad ε) $7.008\ \square\ 7.008$ \quad στ) $12.345\ \square\ 1.234$.

\ex Σχεδίασε μια αριθμητική ευθεία με αρχή το Ο και μονάδα $\mathrm{ΟΑ}=3$ εκ. Τοποθέτησε τα σημεία Β, Γ, Δ και Ε σε αποστάσεις 9, 15, 21 και 27 εκ. από το Ο. Ποιοι αριθμοί αντιστοιχούν σε αυτά;

\ex Στρογγυλοποίησε στην πλησιέστερη εκατοντάδα: 248, 350, 1.249, 5.951, 19.999 και 70.049.

\ex Στρογγυλοποίησε τον αριθμό 2.846.593 (α) στις δεκάδες, (β) στις εκατοντάδες, (γ) στις χιλιάδες, (δ) στις δεκάδες χιλιάδων, (ε) στις εκατοντάδες χιλιάδων, (στ) στα εκατομμύρια.

\ex Σωστό ή λάθος;\par
α) Στον αριθμό 4.050.300 τα μηδενικά δείχνουν ότι δεν υπάρχουν εκατοντάδες χιλιάδων, χιλιάδες, δεκάδες και μονάδες.\quad β) Εκατό δεκάδες είναι χίλιες μονάδες.\quad γ) Από τον 40 ως τον 99 υπάρχουν 60 αριθμοί.\quad δ) Από τη σελίδα 20 ως τη σελίδα 29 είναι 9 σελίδες.\quad ε) Αν σήμερα είναι Δευτέρα, σε 15 ημέρες θα είναι πάλι Δευτέρα.\quad στ) Ανάμεσα στο 7 και στο 8 δεν υπάρχει άλλος φυσικός αριθμός.\quad ζ) Το 2.449 στρογγυλοποιημένο στις εκατοντάδες είναι 2.500.

\zoni{Το χρησιμοποιώ}{προβλήματα από τη ζωή}

\ex \textbf{Το λουκέτο της Φανής.} Το λουκέτο της ποδηλατοθήκης έχει τριψήφιο κωδικό με τα ψηφία 2, 0 και 8, καθένα μία φορά. Γράψε όλους τους δυνατούς κωδικούς από τον μικρότερο στον μεγαλύτερο. Πόσους θα δοκιμάσει η Φανή στη χειρότερη περίπτωση;

\ex \textbf{Το χωριό του παππού.} Η πινακίδα στην είσοδο του χωριού γράφει «Κάτοικοι: 1.846». Ο Νίκος λέει στους φίλους του «περίπου 2.000» και η Μυρτώ «περίπου 1.800». Σε ποια τάξη στρογγυλοποίησε ο καθένας; Ποιος είναι πιο κοντά; Αν ο δήμαρχος πρέπει να τυπώσει ψηφοδέλτια για όλους τους κατοίκους, μπορεί να χρησιμοποιήσει το «περίπου»; Γιατί;

\ex \textbf{Τα βήματα του Νίκου.} Η εφαρμογή στο κινητό του Νίκου γράφει ότι σήμερα έκανε 14.992 βήματα. Ο Νίκος λέει «Στρογγυλά, 15.000, πέτυχα τον στόχο!». Ο στόχος είναι τουλάχιστον 15.000 βήματα. Έφτασε τον στόχο; Πόσα του λείπουν; Γιατί εδώ η στρογγυλοποίηση δεν βοηθάει;

\ex \textbf{Τα πλακάκια του θείου.} Ο θείος Λευτέρης χρειάζεται 2.314 πλακάκια για τον κήπο. Πωλούνται σε κιβώτια των 100. Ο Στέλιος στρογγυλοποιεί στην πλησιέστερη εκατοντάδα και λέει «23 κιβώτια». Φτάνουν; Πόσα κιβώτια πρέπει να αγοράσει και πόσα πλακάκια θα περισσέψουν;

\ex \textbf{Το βιβλίο της Μυρτώς.} Η Μυρτώ διαβάζει από τη σελίδα 204 ως τη σελίδα 311, και τις δύο μαζί. Πόσες σελίδες διαβάζει; Αν διαβάζει 9 σελίδες την ημέρα, σε πόσες ημέρες θα τελειώσει;

\ex \textbf{Ο φράχτης.} Ο παππούς φράζει ένα κομμάτι του κήπου μήκους 36 μ. Βάζει έναν πάσσαλο κάθε 3 μ., και έναν στην αρχή και έναν στο τέλος. Πόσους πασσάλους θα χρειαστεί;

\ex \textbf{Οι κάρτες του αγώνα.} Οι κάρτες του σχολικού αγώνα είναι αριθμημένες από το 501 ως το 780. Πόσες κάρτες υπάρχουν; Πόσες από αυτές έχουν περιττό αριθμό;

\ex \textbf{Ο άξονας του Βασίλη.} Ο Βασίλης σχεδιάζει έναν άξονα για τις βαθμολογίες του παιχνιδιού του. Τα σημεία 0 και 1.000 απέχουν 4 εκ. (δες το σχήμα). Ποιοι αριθμοί αντιστοιχούν στα σημεία Α, Β, Γ και Δ;

\begin{center}
\begin{tikzpicture}[x=1.15cm]
\draw[-{Stealth}] (-0.3,0)--(9.6,0);
\draw (0,-0.12)--(0,0.12) node[above,font=\small] {0};
\draw (4,-0.12)--(4,0.12) node[above,font=\small] {1.000};
\foreach \x/\n in {1/Α,3/Β,5/Γ,7.5/Δ}{\draw[gblue,thick] (\x,-0.15)--(\x,0.15); \node[below,font=\small\bfseries,gblue] at (\x,-0.15) {\n};}
\end{tikzpicture}
\end{center}

\zoni{Το σκέφτομαι}{μια ερώτηση που πάει πιο μακριά}

\ex Πόσες φορές γράφεται το ψηφίο 3, αν γράψεις όλους τους αριθμούς από το 1 ως το 200;

\ex Ποιος είναι ο μικρότερος και ποιος ο μεγαλύτερος φυσικός αριθμός που, στρογγυλοποιημένος στις χιλιάδες, δίνει 8.000;

\ex Ποιος είναι ο μεγαλύτερος και ποιος ο μικρότερος πενταψήφιος αριθμός που έχει όλα τα ψηφία του διαφορετικά;

\ex Ο Νίκος λέει: «Θα βρω τον μεγαλύτερο φυσικό αριθμό και θα τον γράψω στον πίνακα». Η Ζωή λέει: «Δεν υπάρχει». Εξήγησε με λόγια γιατί έχει δίκιο η Ζωή.

\ex Ο Ηλίας δοκιμάζει στο κινητό: $7+5=12$, $3+9=12$. Ρωτάει την Πολυξένη: «Αν προσθέσω δύο φυσικούς, βγαίνει πάντα φυσικός. Αν αφαιρέσω;» Η Πολυξένη γράφει $3-5$. Μπορεί να υπάρχει αυτό το αποτέλεσμα στους φυσικούς; Τι νομίζεις ότι θα χρειαστούμε;

\zoni{Για το σπίτι}{πιο μακριά, με χαρτί και χάρακα}

\ex \textbf{Το θερμόμετρο της Αγγελικής.} Η Αγγελική φτιάχνει ένα χάρτινο θερμόμετρο για την τάξη. Σημειώνει το $0^\circ$ και το $40^\circ$ και η απόσταση ανάμεσά τους είναι 20 εκ. Πού θα βάλει τα σημάδια για τους $10^\circ$, $20^\circ$, $25^\circ$ και $35^\circ$;

\ex Δοκίμασε τις αφαιρέσεις $9-3$, $15-7$ και $21-5$. Τι είδους αριθμούς βρίσκεις; Εξήγησε γιατί η διαφορά δύο περιττών αριθμών είναι πάντα άρτιος.

\ex Πόσοι τριψήφιοι αριθμοί υπάρχουν; (Βοήθεια: ο μικρότερος είναι ο 100 και ο μεγαλύτερος ο 999, και τα δύο άκρα μετράνε.)

\paragrafos{Ευκλείδεια διαίρεση: το υπόλοιπο που μένει}{Α.1.2, «Ταυτότητα Ευκλείδειας Διαίρεσης»}

\begin{anagke}
Το σχολείο οργανώνει την πρώτη του εκδρομή. Θα πάνε και τα 387 παιδιά, με λεωφορεία των 49 θέσεων. Η κυρία Αλεξάνδρα ζητάει από την τάξη να βρει πόσα λεωφορεία χρειάζονται.

Ο Στέλιος, που αγαπάει το κομπιουτεράκι, κάνει $387:49$ και η οθόνη γράφει $7{,}897959\dots$ «Εφτά κόμμα εννιά!» λέει περήφανος. «Άρα εφτά λεωφορεία». Η Ρήνη τον κοιτάει. «Και τα υπόλοιπα παιδιά; Μένουν στην αυλή;» Ο Στέλιος ξαναδιαβάζει την οθόνη. «Μα το 0,897 τι είναι; Δεν υπάρχει 0,897 λεωφορείου». Η Ζωή κάνει τη διαίρεση στο χαρτί, όπως τη μάθαμε: πηλίκο 7 και κάτι περισσεύει. «Αυτό που περισσεύει είναι παιδιά, όχι κόμματα. Εκεί βρίσκεται η απάντηση».

Σε πολλά προβλήματα δεν ενδιαφέρει μόνο πόσες φορές χωράει ο ένας αριθμός στον άλλον, ενδιαφέρει και τι περισσεύει. Σε αυτή την παράγραφο μαθαίνουμε να γράφουμε τη διαίρεση με μία ισότητα, την ταυτότητα της Ευκλείδειας διαίρεσης, και να τη χρησιμοποιούμε για να λύνουμε προβλήματα.
\end{anagke}

\begin{lexi}{πηλίκο, υπόλοιπο, Ευκλείδεια}{πηλίκο, υπόλοιπο}
Το \lexiword{πηλίκο} έρχεται από το αρχαίο «πηλίκος», που σημαίνει «πόσο μεγάλος»: η διαίρεση απαντά στο «πόσο μεγάλο μερίδιο παίρνει ο καθένας». Το \lexiword{υπόλοιπο} είναι κυριολεκτικά αυτό που μένει, το «λοιπό» που δεν μοιράστηκε.

Η \lexiword{Ευκλείδεια διαίρεση} πήρε το όνομά της από τον Ευκλείδη, τον Έλληνα μαθηματικό που έζησε στην Αλεξάνδρεια γύρω στο 300 π.Χ. και έγραψε τα «Στοιχεία», το βιβλίο όπου συγκέντρωσε τα μαθηματικά της εποχής του.
\end{lexi}

\begin{metafrasi}
Μοιράζω 43 μπάλες σε ομάδες των 7: γίνονται 6 ομάδες και μένει 1 μπάλα. & $43=7\cdot 6+1$\\
Δεν περισσεύει τίποτα (τέλεια διαίρεση). & $\Delta=\delta\cdot\pi$, υπόλοιπο $0$\\
Το υπόλοιπο είναι πάντα μικρότερο από τον διαιρέτη. & $\upsilon<\delta$\\
Στη διαίρεση με το 8 τα πιθανά υπόλοιπα. & $0,1,2,3,4,5,6,7$\\
Δεν διαιρούμε ποτέ με το μηδέν. & $\delta\neq 0$\\
\end{metafrasi}

\tmima{Η θεωρία}

\begin{kanonas}{ταυτότητα της Ευκλείδειας διαίρεσης}
Αν $\Delta$ και $\delta$ είναι δύο φυσικοί αριθμοί με $\delta\neq 0$, τότε υπάρχουν δύο άλλοι φυσικοί $\pi$ και $\upsilon$ έτσι ώστε
$$\Delta=\delta\cdot\pi+\upsilon,\qquad\text{με}\quad \upsilon<\delta.$$
\end{kanonas}

Ο $\Delta$ λέγεται \textbf{διαιρετέος}, ο $\delta$ \textbf{διαιρέτης}, ο $\pi$ \textbf{πηλίκο} και ο $\upsilon$ \textbf{υπόλοιπο}. Για παράδειγμα, από τη διαίρεση $43:7$ έχουμε $\pi=6$ και $\upsilon=1$, δηλαδή $43=7\cdot 6+1$ με $1<7$.

Αν το υπόλοιπο είναι $0$, η διαίρεση λέγεται \textbf{τέλεια} και γράφεται $\Delta=\delta\cdot\pi$. Τότε ισχύει και $\Delta:\delta=\pi$ και $\Delta:\pi=\delta$. Η τέλεια διαίρεση είναι η πράξη αντίστροφη του πολλαπλασιασμού.

\begin{idiotita}{ειδικές περιπτώσεις}
\begin{itemize}[leftmargin=1.2em,itemsep=0.3ex,topsep=2pt]
\item Ο διαιρέτης δεν μπορεί να είναι $0$: $\delta\neq 0$.
\item $\alpha:\alpha=1$ (όταν $\Delta=\delta$ το πηλίκο είναι $1$).
\item $\alpha:1=\alpha$ (όταν $\delta=1$ το πηλίκο είναι ο διαιρετέος).
\item $0:\alpha=0$ (όταν $\Delta=0$ το πηλίκο είναι $0$).
\end{itemize}
\end{idiotita}

\begin{prosoxi}
Μια ισότητα της μορφής $\Delta=\delta\cdot\pi+\upsilon$ δεν είναι πάντα Ευκλείδεια διαίρεση. Πρέπει να ισχύει και $\upsilon<\delta$. Για παράδειγμα η $50=7\cdot 6+8$ είναι αληθινή ισότητα, αλλά δεν είναι Ευκλείδεια διαίρεση, γιατί $8>7$: στο $50$ το $7$ χωράει ακόμη μία φορά ($50=7\cdot 7+1$).
\end{prosoxi}

\tmimasmall{Τι υπόλοιπα μπορούν να βγουν}

Επειδή $\upsilon<\delta$, στη διαίρεση με το $5$ τα πιθανά υπόλοιπα είναι μόνο $0,1,2,3,4$. Στη διαίρεση με το $\delta$ υπάρχουν ακριβώς $\delta$ διαφορετικά υπόλοιπα, από το $0$ ως το $\delta-1$. Γι' αυτό τα προβλήματα που επαναλαμβάνονται κυκλικά (ημέρες της εβδομάδας, ώρες του ρολογιού) λύνονται με Ευκλείδεια διαίρεση.

\tmima{Λυμένα παραδείγματα}

\paradeigma{Ποιες ισότητες είναι Ευκλείδεια διαίρεση;}
\tixero Οι ισότητες α) $54=25\cdot 3+9$, β) $370=46\cdot 8+2$, γ) $680=51\cdot 12+68$.\par
\tithelo Σε ποιες ισχύει $\upsilon<\delta$ και ποιοι είναι οι $\Delta,\delta,\pi,\upsilon$;\par
\possindeo Στην α) είναι $9<25$: Ευκλείδεια διαίρεση με $\Delta=54$, $\delta=25$, $\pi=3$, $\upsilon=9$. Στη β) το $2$ είναι μικρότερο και από τον $46$ και από τον $8$, άρα μπορούμε να πούμε $\delta=46,\pi=8$ ή $\delta=8,\pi=46$ (και στις δύο $\upsilon=2$). Στη γ) το $68$ είναι μεγαλύτερο και από τον $51$ και από τον $12$, άρα δεν είναι διαίρεση.\par
\apantisi Οι α) και β) εκφράζουν Ευκλείδεια διαίρεση, η γ) όχι.

\paradeigma{Οι καρέκλες της αίθουσας}
\tixero Κάθε αίθουσα πρέπει να έχει $12$ θρανία. Σήμερα ήρθαν $3$ φορτηγά που το καθένα μετέφερε $35$ θρανία.\par
\tithelo Πόσες αίθουσες συμπληρώνονται και πόσα θρανία χρειάζονται ακόμη για να συμπληρωθεί μία ακόμη;\par
\possindeo Συνολικά $3\cdot 35=105$ θρανία. Κάνω τη διαίρεση $105:12$: $105=12\cdot 8+9$. Άρα $8$ αίθουσες συμπληρώνονται και περισσεύουν $9$ θρανία. Για την επόμενη χρειάζονται ακόμη $12-9=3$.\par
\apantisi $8$ αίθουσες, και χρειάζονται άλλα $3$ θρανία.\par
\elegxos $12\cdot 8+9=96+9=105$ και $9<12$.

\paradeigma{Η μέρα της εβδομάδας}
\tixero Σήμερα είναι Δευτέρα.\par
\tithelo Τι μέρα θα είναι σε $100$ ημέρες;\par
\possindeo Οι ημέρες επαναλαμβάνονται ανά $7$. Κάνω τη διαίρεση: $100=7\cdot 14+2$. Μετά από $14$ πλήρεις εβδομάδες ξαναγυρνάμε σε Δευτέρα, και μένουν $2$ ημέρες: Τρίτη, Τετάρτη.\par
\apantisi Τετάρτη.\par
\elegxos $7\cdot 14=98$ ημέρες είναι Δευτέρα, και $2$ ακόμη φέρνουν Τετάρτη.

\tmima{Ασκήσεις}

\zoni{Το κατάλαβα}{εφαρμογή της ταυτότητας}

\ex Συμπλήρωσε: α) Στη διαίρεση $47:6$ ο διαιρετέος είναι \dots, ο διαιρέτης \dots, το πηλίκο \dots και το υπόλοιπο \dots β) Ο διαιρέτης δεν μπορεί να είναι \dots γ) Αν ο διαιρετέος ισούται με τον διαιρέτη, το πηλίκο είναι \dots δ) Αν ο διαιρετέος είναι $0$, το πηλίκο είναι \dots

\ex Ποιες από τις ισότητες εκφράζουν Ευκλείδεια διαίρεση; α) $47=6\cdot 7+5$ \quad β) $50=7\cdot 6+8$ \quad γ) $29=4\cdot 7+1$ \quad δ) $100=9\cdot 10+10$ \quad ε) $63=9\cdot 7+0$.

\ex Κάνε τις διαιρέσεις και τις δοκιμές τους: α) $135:5$ \quad β) $95:4$ \quad γ) $276:12$ \quad δ) $1.000:7$. Ποιες είναι τέλειες;

\ex Ποια υπόλοιπα μπορεί να έχει μια διαίρεση με διαιρέτη το $5$; Ποια με διαιρέτη το $8$; Ποια με διαιρέτη το $1$;

\ex Σωστό ή λάθος; α) Η διαίρεση $452:2$ είναι τέλεια. \quad β) Ο διαιρέτης μπορεί να είναι $0$. \quad γ) Το υπόλοιπο είναι πάντα μικρότερο από τον διαιρέτη. \quad δ) Αν ο διαιρέτης είναι $1$, το πηλίκο ισούται με τον διαιρετέο. \quad ε) $0:7=7$. \quad στ) $\alpha:\alpha=\alpha$.

\ex Γράψε την ταυτότητα της Ευκλείδειας διαίρεσης για: α) $83:9$ \quad β) $250:25$ \quad γ) $7:10$ \quad δ) $0:5$.

\zoni{Το χρησιμοποιώ}{προβλήματα από τη ζωή}

\ex \textbf{Τα βαρέλια.} Ο παππούς Γιάννης έχει $350$ λίτρα κρασί και βάζει $12$ λίτρα σε κάθε βαρέλι. Πόσα βαρέλια γεμίζουν και πόσα λίτρα περισσεύουν;

\ex \textbf{Η εκδρομή.} Στην εκδρομή θέλουν να πάνε $250$ παιδιά και κάθε λεωφορείο έχει $45$ θέσεις. Πόσα λεωφορεία χρειάζονται; Πόσες θέσεις μένουν άδειες στο τελευταίο;

\ex \textbf{Οι ημέρες.} Σήμερα είναι Τετάρτη. Τι μέρα θα είναι μετά από $45$ ημέρες; Και μετά από $100$ ημέρες;

\ex \textbf{Τα βραχιόλια.} Η Ρήνη έχει $100$ χάντρες και φτιάχνει βραχιόλια με $8$ χάντρες το καθένα. Πόσα βραχιόλια φτιάχνει και πόσες χάντρες περισσεύουν; Πόσες χάντρες της λείπουν για να φτιάξει ένα ακόμη;

\ex \textbf{Το όνομα που επαναλαμβάνεται.} Η Αγγελική γράφει συνέχεια τη λέξη $\mathrm{ΑΓΓΕΛΙΚΗ}$ σε μια λωρίδα χαρτί: $\mathrm{ΑΓΓΕΛΙΚΗΑΓΓΕΛΙΚΗ}\dots$ Ποιο γράμμα βρίσκεται στην $100$ή θέση;

\zoni{Το σκέφτομαι}{ερωτήσεις με λίγη περισσότερη σκέψη}

\ex Ένας αριθμός διαιρούμενος με το $7$ δίνει πηλίκο $12$ και υπόλοιπο $6$. Ποιος είναι ο αριθμός;

\ex Ένας αριθμός διαιρείται τέλεια με το $14$ και δίνει πηλίκο $201$. Ποιος είναι ο αριθμός;

\ex Ένας αριθμός διαιρούμενος με το $9$ έχει πηλίκο $5$ και το μεγαλύτερο δυνατό υπόλοιπο. Ποιος είναι; Τι πηλίκο και τι υπόλοιπο δίνει ο επόμενός του;

\ex Ο Βασίλης λέει: «Αν ένας αριθμός διαιρούμενος με το $4$ δίνει υπόλοιπο $3$, τότε ο διπλάσιός του διαιρούμενος με το $4$ θα δώσει υπόλοιπο $6$». Εξήγησε γιατί αυτό δεν μπορεί να ισχύει και βρες το σωστό υπόλοιπο. (Δοκίμασε με τον αριθμό $7$.)

\zoni{Για το σπίτι}{εξάσκηση}

\ex Υπολόγισε τις τιμές: α) $35-12:3+3\cdot 3$ \quad β) $63:9+63:7-12:2$ \quad γ) $(5+12:3):3-24:(2\cdot 2+16:4)$.

\ex \textbf{Οι συνοδοί.} Στο καταφύγιο ζώων της περιοχής πρέπει να υπάρχει ένας εθελοντής ανά $15$ σκυλάκια, το πολύ. Αν στο καταφύγιο υπάρχουν $162$ σκυλάκια, ποιος είναι ο μικρότερος αριθμός εθελοντών που χρειάζεται;

\paragrafos{Δυνάμεις φυσικών αριθμών: ποιος πάει πρώτος}{Α.1.3, «Δυνάμεις φυσικών αριθμών»}

\begin{anagke}
Παρασκευή απόγευμα. Η Λίνα φτιάχνει κωδικό για το καινούργιο της κινητό: τέσσερα ψηφία. «Πόσοι διαφορετικοί κωδικοί υπάρχουν;» ρωτάει η Ζωή. Η Λίνα σκέφτεται δυνατά: «Για το πρώτο ψηφίο έχω 10 επιλογές, από το 0 ως το 9. Για το δεύτερο πάλι 10, και για το τρίτο, και για το τέταρτο». Ο Ηλίας γράφει $10\cdot 10\cdot 10\cdot 10$ και βρίσκει 10.000. «Και αν ο κωδικός είχε είκοσι ψηφία;» ρωτάει. Γράφει είκοσι δεκάρια στη σειρά και χάνει το μέτρημα στη μέση.

Ο Άρης, που βιάζεται πάντα, γράφει σε ένα χαρτάκι $2+3\cdot 4$ και ρωτάει τη Ζωή πόσο κάνει. «Δεκατέσσερα», λέει εκείνη. «Είκοσι!» φωνάζει ο Άρης. «Δύο και τρία πέντε, επί τέσσερα, είκοσι». Ο Θοδωρής, που περνάει από δίπλα, τους κοιτάει και λέει μόνο: «Και οι δύο έχετε ένα κομπιουτεράκι στην τσέπη. Δείτε τι βγάζει». Το κομπιουτεράκι δίνει 14.

Δύο προβλήματα. Το πρώτο: πώς γράφεται σύντομα ο ίδιος παράγοντας πολλές φορές; Το δεύτερο: όταν μια παράσταση έχει πολλές πράξεις, ποια γίνεται πρώτη, και ποιος το αποφασίζει;
\end{anagke}

\begin{lexi}{δύναμη, εκθέτης, βάση}{δύναμη, εκθέτης}
Η λέξη \lexiword{δύναμη} στα μαθηματικά δεν έχει σχέση με τα μπράτσα. Ο αριθμός «έχει τη δύναμη» να πολλαπλασιάζεται με τον εαυτό του πολλές φορές. Το $2^{10}$ διαβάζεται «το δύο στη δέκατη».

Ο \lexiword{εκθέτης} είναι το μικρό νούμερο πάνω δεξιά. Το «εκθέτω» σημαίνει «δείχνω»: ο εκθέτης δείχνει πόσες φορές μπαίνει ο παράγοντας. Ο αριθμός που πολλαπλασιάζεται λέγεται \lexiword{βάση}, γιατί πάνω του «στηρίζεται» το σύμβολο.
\end{lexi}

\begin{metafrasi}
Το 2 πολλαπλασιασμένο με τον εαυτό του 10 φορές. & $2^{10}=2\cdot 2\cdot\dots\cdot 2$ (10 παράγοντες)\\
Το τετράγωνο του 5. & $5^2=5\cdot 5=25$\\
Ο κύβος του 3. & $3^3=3\cdot 3\cdot 3=27$\\
Το 7 στην πρώτη. & $7^1=7$\\
Ένα εκατομμύριο, με δύναμη του 10. & $10^6=1.000.000$\\
Πρώτα οι δυνάμεις, μετά οι πολλαπλασιασμοί. & $2+3\cdot 4^2=2+3\cdot 16=50$\\
\end{metafrasi}

\tmima{Η θεωρία}

\begin{orismos}{δύναμη}
Το γινόμενο $\alpha\cdot\alpha\cdot\alpha\cdots\alpha$ με $\nu$ παράγοντες ίσους με το $\alpha$ λέγεται \textbf{δύναμη} του $\alpha$ στη $\nu$ και γράφεται $\alpha^{\nu}$. Ο $\alpha$ είναι η \textbf{βάση} και ο $\nu$ ο \textbf{εκθέτης}.
\end{orismos}

Η δεύτερη δύναμη $\alpha^2$ λέγεται και \textbf{τετράγωνο} του $\alpha$ (είναι το εμβαδό τετραγώνου με πλευρά $\alpha$). Η τρίτη δύναμη $\alpha^3$ λέγεται και \textbf{κύβος} του $\alpha$ (είναι ο όγκος κύβου με ακμή $\alpha$). Η πρώτη δύναμη $\alpha^1$ είναι ο ίδιος ο αριθμός $\alpha$. Οι δυνάμεις του 1 είναι όλες 1: $1^{\nu}=1$.

\begin{prosoxi}
Το $2^3$ δεν είναι $2\cdot 3=6$ αλλά $2\cdot 2\cdot 2=8$. Ο εκθέτης δείχνει πόσες φορές γράφεται η βάση, όχι επί πόσο την πολλαπλασιάζουμε.
\end{prosoxi}

\tmimasmall{Οι δυνάμεις του 10}

Οι δυνάμεις του 10 έχουν τόσα μηδενικά όσος ο εκθέτης:
$$10^2=100,\qquad 10^3=1.000,\qquad 10^6=1.000.000.$$
Γι' αυτό η δεκαδική μας γραφή αναλύεται εύκολα. Για παράδειγμα:
$$7.604=7\cdot 10^3+6\cdot 10^2+0\cdot 10+4.$$

\tmimasmall{Αριθμητική παράσταση και προτεραιότητα}

Μια \textbf{αριθμητική παράσταση} είναι μια σειρά αριθμών που συνδέονται με σύμβολα πράξεων. Ο αριθμός που βγαίνει όταν κάνουμε όλες τις πράξεις λέγεται \textbf{τιμή} της. Για να βγάζουμε όλοι την ίδια τιμή, όλος ο κόσμος έχει συμφωνήσει σε μια σειρά.

\begin{kanonas}{προτεραιότητα των πράξεων}
\begin{enumerate}[leftmargin=1.6em,itemsep=0.2ex,topsep=2pt]
\item Πρώτα οι πράξεις μέσα στις \textbf{παρενθέσεις}, με την ίδια σειρά που θα βρεις παρακάτω.
\item Μετά οι \textbf{δυνάμεις}.
\item Μετά οι \textbf{πολλαπλασιασμοί και οι διαιρέσεις}, από αριστερά προς τα δεξιά.
\item Τελευταίες οι \textbf{προσθέσεις και οι αφαιρέσεις}, από αριστερά προς τα δεξιά.
\end{enumerate}
\end{kanonas}

Η συμφωνία αυτή δεν είναι «φυσικός νόμος». Είναι συμφωνία, όπως ο κανόνας ότι στον δρόμο πάμε δεξιά. Αν δεν υπήρχε, το $2+3\cdot 4$ θα είχε δύο απαντήσεις (14 και 20) και δύο μηχανικοί δεν θα συμφωνούσαν ποτέ σε έναν υπολογισμό. Με τη συμφωνία η απάντηση είναι μία, 14, και η Ζωή έχει δίκιο. Αν θέλουμε το 20, γράφουμε $(2+3)\cdot 4$: η παρένθεση λέει ρητά «αυτό πρώτα».

\begin{prosoxi}
\begin{itemize}[leftmargin=1.2em,itemsep=0.3ex,topsep=2pt]
\item Στο $2\cdot 3^2$ ο εκθέτης αφορά μόνο το 3: $2\cdot 9=18$. Για να τετραγωνίσεις όλο το γινόμενο, γράφεις $(2\cdot 3)^2=36$.
\item Οι πράξεις του ίδιου επιπέδου γίνονται από αριστερά προς τα δεξιά: $24:4:2=6:2=3$, όχι $24:(4:2)=12$.
\end{itemize}
\end{prosoxi}


\tmimasmall{Ιστορική αναδρομή: πώς έγραφαν τους αριθμούς}

Το δεκαδικό σύστημα, που σου φαίνεται αυτονόητο, χρειάστηκε χιλιάδες χρόνια για να διαμορφωθεί. Οι Σουμέριοι και οι Βαβυλώνιοι (3η–2η χιλιετία π.Χ.) έγραφαν τους αριθμούς σε πήλινες πινακίδες. Οι Αιγύπτιοι είχαν ένα ιερογλυφικό για κάθε τάξη: 1, 10, 100, 1.000, \dots, 1.000.000. Οι Ίωνες, τον 5ο αιώνα π.Χ., έφτιαξαν το \textbf{αλφαβητικό σύστημα}: κάθε γράμμα του αλφαβήτου σήμαινε έναν αριθμό.

\begin{center}
\small\renewcommand{\arraystretch}{1.25}\setlength{\tabcolsep}{4pt}
\begin{tabular}{cccccccccc}
α & β & γ & δ & ε & \arch{ϛ} & ζ & η & θ & \\
1 & 2 & 3 & 4 & 5 & 6 & 7 & 8 & 9 & \\
ι & κ & λ & μ & ν & ξ & ο & π & \arch{ϙ} & \\
10 & 20 & 30 & 40 & 50 & 60 & 70 & 80 & 90 & \\
ρ & σ & τ & υ & φ & χ & ψ & ω & \arch{ϡ} & \\
100 & 200 & 300 & 400 & 500 & 600 & 700 & 800 & 900 & \\
\end{tabular}
\end{center}

Έτσι το 831 γράφεται ωλα ($800+30+1$). Οι \textbf{Ρωμαίοι} είχαν σύμβολα για τους αριθμούς $1,\,5,\,10,\,50,\,100,\,500,\,1.000$:

\begin{center}
\small\setlength{\tabcolsep}{6pt}
\begin{tabular}{ccccccc}
I & V & X & L & C & D & M\\
1 & 5 & 10 & 50 & 100 & 500 & 1.000\\
\end{tabular}
\end{center}

Πρόσθεταν από αριστερά προς τα δεξιά (II = 2, XVI = 16) και αφαιρούσαν όταν ένα μικρό σύμβολο προηγείται: IV = 4 ($5-1$), IX = 9, XL = 40, CM = 900. Τα σύμβολα που χρησιμοποιούμε σήμερα (0 ως 9) προέρχονται από τους Ινδούς, έφτασαν στην Ευρώπη μέσω των Αράβων και γι' αυτό λέγονται \textbf{ινδοαραβικά}. Εμφανίζονται πρώτη φορά στο έργο του Αλ-Χουαρίζμι (9ος αιώνας μ.Χ.).

Δεν είναι ανάγκη να έχουμε δέκα ψηφία. Ο υπολογιστής χρησιμοποιεί μόνο το 0 και το 1, στο \textbf{δυαδικό σύστημα}, όπου αντί για μονάδες, δεκάδες, εκατοντάδες έχουμε μονάδες, δυάδες, τετράδες, οκτάδες, δεκαεξάδες \dots

\begin{center}
\small\setlength{\tabcolsep}{4pt}
\begin{tabular}{l|cccccccccccc}
Δεκαδικό & 0 & 1 & 2 & 3 & 4 & 5 & 6 & 7 & 8 & 9 & 10 & 11\\
Δυαδικό & 0 & 1 & 10 & 11 & 100 & 101 & 110 & 111 & 1000 & 1001 & 1010 & 1011\\
\end{tabular}
\end{center}

Οι Πυθαγόρειοι πίστευαν ότι «τα πάντα είναι αριθμός». Έβαζαν τους αριθμούς σε σχήματα: οι \textbf{τρίγωνοι} αριθμοί $1,3,6,10,\dots$ φτιάχνουν τρίγωνα από τελείες και οι \textbf{τετράγωνοι} $1,4,9,16,\dots$ τετράγωνα. Οι τετράγωνοι είναι οι δεύτερες δυνάμεις.
\tmima{Λυμένα παραδείγματα}

\paradeigma{Οι κωδικοί της Λίνας}
\tixero Ο κωδικός έχει 4 ψηφία και κάθε ψηφίο έχει 10 επιλογές (0 ως 9).\par
\tithelo Πόσοι είναι οι δυνατοί κωδικοί, και πώς γράφεται σύντομα ο υπολογισμός;\par
\possindeo Για κάθε ψηφίο έχω 10 επιλογές, και οι επιλογές συνδυάζονται, άρα πολλαπλασιάζω: $10\cdot 10\cdot 10\cdot 10$. Είναι τέσσερα δεκάρια, δηλαδή $10^4$. Υπολογίζω βήμα βήμα: $10\cdot 10=100$, $100\cdot 10=1.000$, $1.000\cdot 10=10.000$.\par
\apantisi $10^4=10.000$ κωδικοί.\par
\elegxos Οι κωδικοί είναι οι αριθμοί από το 0000 ως το 9999. Από την παράγραφο 1.1: $9.999-0+1=10.000$. Βγήκε το ίδιο.

\paradeigma{Η παράσταση με πολλές πράξεις}
\tixero Η παράσταση $(2\cdot 5)^4+4\cdot(3+2)^2$.\par
\tithelo Την τιμή της.\par
\possindeo Πρώτα οι παρενθέσεις: $2\cdot 5=10$ και $3+2=5$. Έχω $10^4+4\cdot 5^2$. Μετά οι δυνάμεις: $10^4=10.000$ και $5^2=25$. Μετά ο πολλαπλασιασμός: $4\cdot 25=100$. Τέλος η πρόσθεση: $10.000+100=10.100$.\par
\apantisi $10.100$.\par
\elegxos Το $10^4$ έχει 4 μηδενικά, 10.000. Το άθροισμα με το 100 δίνει 10.100, όχι κάτι πιο μικρό από το 10.000. Σωστά.

\paradeigma{Ο αριθμός του Ηλία}
\tixero Ο Ηλίας γράφει ένα αριθμό σε ανάπτυγμα: $5\cdot 10^4+3\cdot 10^2+8$.\par
\tithelo Ποιος αριθμός είναι;\par
\possindeo $5\cdot 10^4=5\cdot 10.000=50.000$. $3\cdot 10^2=300$. Άρα $50.000+300+8=50.308$.\par
\apantisi $50.308$.\par
\elegxos Λείπουν οι χιλιάδες και οι δεκάδες, που σημαίνουν 0 στη θέση τους: 50.308. Τα μηδενικά έχουν μπει σωστά.

\tmima{Ασκήσεις}

\zoni{Το κατάλαβα}{εφαρμογή του κανόνα}

\ex Γράψε ως δύναμη: α) $5\cdot 5\cdot 5$ \quad β) $2\cdot 2\cdot 2\cdot 2\cdot 2$ \quad γ) $10\cdot 10\cdot 10\cdot 10$ \quad δ) $7\cdot 7$ \quad ε) $9$.

\ex Γράψε ως γινόμενο και υπολόγισε: α) $2^4$ \quad β) $3^3$ \quad γ) $5^2$ \quad δ) $10^5$ \quad ε) $1^{100}$ \quad στ) $6^1$.

\ex Ποια είναι η βάση και ποιος ο εκθέτης στα $4^5$, $12^2$, $1^{9}$ και $10^8$;

\ex Υπολόγισε το τετράγωνο και τον κύβο των αριθμών 2, 3, 4, 5 και 10. Φτιάξε πίνακα.

\ex Υπολόγισε την τιμή: α) $12+3\cdot 4$ \quad β) $(12+3)\cdot 4$ \quad γ) $24:4:2$ \quad δ) $100-5\cdot 4$ \quad ε) $2+3^2$ \quad στ) $(2+3)^2$.

\ex Βρες την τιμή: α) $(2\cdot 5)^4+4\cdot(3+2)^2$ \quad β) $(2+3)^3-8\cdot 3^2$ \quad γ) $3\cdot(4+2)^2-10\cdot 3$.

\ex Γράψε το ανάπτυγμα με δυνάμεις του 10: α) 7.604 \quad β) 53.280 \quad γ) 400.065.

\ex Γράψε με ψηφία τον αριθμό: α) $3\cdot 10^3+2\cdot 10^2+5$ \quad β) $9\cdot 10^5+4\cdot 10$ \quad γ) $10^6+10^2$.

\ex Σωστό ή λάθος; α) $2^3=6$ \quad β) $10^4=10.000$ \quad γ) $5^2=2\cdot 5$ \quad δ) $1^{57}=1$ \quad ε) $3+2\cdot 5=25$ \quad στ) $(1+1)^2=1^2+1^2$.

\zoni{Το χρησιμοποιώ}{προβλήματα από τη ζωή}

\ex \textbf{Το λουκέτο.} Ένα λουκέτο έχει τροχούς με ψηφία από το 0 ως το 9. Πόσους διαφορετικούς κωδικούς έχει αν έχει 3 τροχούς; Αν έχει 4; Αν έχει 5; Γράψε κάθε απάντηση και ως δύναμη του 10.

\ex \textbf{Οι διακόπτες.} Στην αποθήκη του θείου Λευτέρη ο πίνακας έχει διακόπτες που ο καθένας έχει δύο θέσεις: ανοιχτός ή κλειστός. Πόσους διαφορετικούς συνδυασμούς θέσεων έχουν 3 διακόπτες; Πόσους 5; Πόσους 10; Γράψε κάθε απάντηση και ως δύναμη του 2.

\ex \textbf{Ο κύβος του θείου.} Ο θείος Λευτέρης έχει δοχείο σε σχήμα κύβου με ακμή 4 δεκατόμετρα. Πόσα κυβικά δεκατόμετρα χωράει; (Θυμήσου ότι ο όγκος κύβου είναι $\alpha^3$.)

\ex \textbf{Ο κήπος.} Ο τετράγωνος κήπος της γιαγιάς Ρήνης έχει πλευρά 12 μ. Βρες το εμβαδό του. Αν διπλασιαστεί η πλευρά, ο κήπος γίνεται διπλάσιος σε εμβαδό; Υπολόγισε.

\ex \textbf{Τα σκακιστικά τετράγωνα.} Η σκακιέρα έχει 8 γραμμές και 8 στήλες. Πόσα τετράγωνα έχει; Γράψε τον υπολογισμό με δύναμη.

\ex \textbf{Ο λογαριασμός του Άρη.} Ο Άρης αγοράζει 3 μπάλες των 8 € και ένα φανελάκι των 12 €. Γράφει $3\cdot 8+12$ και βρίσκει 36. Η Μυρτώ γράφει $3\cdot(8+12)$. Ποια παράσταση ταιριάζει στο πρόβλημα και πόσο πληρώνει τελικά ο Άρης;

\ex \textbf{Τρεις απαντήσεις.} Τρεις φίλοι υπολογίζουν την παράσταση $4\cdot(7+7\cdot 9)+20$. Ο Άρης βρίσκει 335, η Ζωή 300 και ο Ηλίας 524. Ποιος έχει δίκιο; Βρες πού έκανε λάθος καθένας από τους άλλους δύο.

\zoni{Το σκέφτομαι}{μια ερώτηση που πάει πιο μακριά}

\ex Πόσο κάνει το $2^3\cdot 2^2$; Γράψε το ως γινόμενο παραγόντων (όλα τα δυάρια) και μέτρα. Ανακαλύπτεις έναν γρήγορο κανόνα για τον πολλαπλασιασμό δυνάμεων με την ίδια βάση;

\ex Συγκρίνε τα $2^5$ και $5^2$, και τα $2^4$ και $4^2$. Μετά βρες έναν αριθμό $\alpha$ διαφορετικό από το 2 για τον οποίο το $\alpha^2$ είναι ίσο με το $2\alpha$.

\ex Τοποθέτησε παρενθέσεις στην παράσταση $2+3\cdot 4+5$ με όσους τρόπους μπορείς και βρες όλες τις διαφορετικές τιμές που μπορεί να πάρει.

\ex Βάλε τα σύμβολα $+,-,\cdot$ (ή παρενθέσεις) ανάμεσα στα ψηφία $1\ \ 2\ \ 3\ \ 4$ ώστε το αποτέλεσμα να είναι 10. Βρες τουλάχιστον δύο τρόπους.

\ex Ο Ηλίας λέει: «Το $10^0$ πρέπει να είναι 0, γιατί δεν πολλαπλασιάζεις το 10 καμία φορά». Η Ζωή λέει: «Μάλλον είναι 1». Δες το μοτίβο $10^3=1.000$, $10^2=100$, $10^1=10$ και σκέψου τι πρέπει να συνεχίσει. Ποιος νομίζεις ότι έχει δίκιο; (Θα το μάθεις στο σχολείο σε λίγους μήνες.)

\zoni{Για το σπίτι}{εξάσκηση και ιστορία των αριθμών}

\ex Βρες τα τετράγωνα των 10, 20, 30, 40, 50, 60, 70, 80, 90 και τους κύβους των 10, 20, 30, 40, 50. Τι παρατηρείς στα μηδενικά;

\ex Υπολόγισε τις δυνάμεις $2^1, 2^2, 2^3, \dots, 2^{10}$ με τη σειρά, κάθε μία από την προηγούμενη.

\ex Γράψε πιο σύντομα: α) $\alpha+\alpha+\alpha$ \quad β) $\alpha\cdot\alpha\cdot\alpha$ \quad γ) $x+x+x+x$ \quad δ) $x\cdot x\cdot x\cdot x$ \quad ε) $2\cdot 2\cdot\alpha\cdot\alpha\cdot\alpha$.

\ex Βρες τις τιμές των παραστάσεων $(6+5)^2$ και $6^2+5^2$, και μετά των $(3+6)^2$ και $3^2+6^2$. Τι παρατηρείς;

\ex \textbf{Ρωμαϊκοί αριθμοί.} α) Γράψε με ρωμαϊκά τους 4, 9, 14, 40, 49, 90 και 1.994. β) Διάβασε τους XXXVII, CDXLIV και MMXXVI.

\ex \textbf{Αλφαβητικοί αριθμοί.} α) Γράψε με το αλφαβητικό σύστημα των Ιώνων τους 12, 25 και 831. β) Διάβασε τους ρκγ, σμ και τξε.

\ex \textbf{Δυαδικοί αριθμοί.} α) Γράψε στο δυαδικό σύστημα τους 11, 12 και 16. β) Διάβασε τους δυαδικούς 111, 1101 και 10101.

\ex \textbf{Οι αριθμοί των Πυθαγορείων.} Συνέχισε την ακολουθία των τριγώνων αριθμών $1,3,6,10$ για τρεις ακόμα όρους. Πόσο είναι ο δέκατος τρίγωνος αριθμός; Πόσο είναι ο δωδέκατος τετράγωνος;

\ex \textbf{Μαγικό τετράγωνο.} Βάλε τους αριθμούς από το 1 ως το 9, από μία φορά τον καθένα, ώστε κάθε γραμμή, κάθε στήλη και κάθε διαγώνιος να έχει άθροισμα 15. Έχουν μπει ήδη μερικοί:
\begin{center}
\renewcommand{\arraystretch}{1.6}\setlength{\tabcolsep}{12pt}
\begin{tabular}{|c|c|c|}\hline 2 & & 6\\\hline & 5 & \\\hline 4 & & 8\\\hline\end{tabular}
\end{center}

\paragrafos{ΕΚΠ και ΜΚΔ: πότε ξανασυναντιόμαστε, πόσο μεγάλα κομμάτια}{Α.1.4, «ΕΚΠ και ΜΚΔ φυσικών αριθμών»}

\begin{anagke}
Η γιαγιά Μυρτώ ποτίζει τα φυτά της αυλής κάθε 3 ημέρες και ο παππούς Γιάννης κλαδεύει τη λεμονιά κάθε 4 ημέρες. Σήμερα έκαναν και οι δύο τη δουλειά τους. «Πότε θα τα ξανακάνουμε και οι δύο την ίδια μέρα;» ρωτά η Ζωή, που έχει έρθει για επίσκεψη. Ο Ηλίας γράφει στο τετράδιο τις ημέρες της γιαγιάς: 3, 6, 9, 12, 15, 18, και του παππού: 4, 8, 12, 16, 20. «Στις 12!» φωνάζει.

Το απόγευμα η θεία Μαρία φέρνει 48 μανταρίνια και 60 πορτοκάλια και θέλει να τα βάλει σε ίσα καλάθια, με τον ίδιο αριθμό μανταρινιών και πορτοκαλιών σε κάθε καλάθι, χωρίς να περισσέψει τίποτα. «Πόσα το πολύ καλάθια μπορώ να φτιάξω;» ρωτά. Αυτό δεν είναι «πότε ξανασυναντιόμαστε»: είναι «σε πόσα ίσα κομμάτια χωρίζονται και τα δύο».

Δύο διαφορετικά προβλήματα: το ένα ψάχνει έναν αριθμό που τον φτάνουν και οι δύο (πολλαπλάσιο), το άλλο έναν αριθμό που χωράει σε όλους (διαιρέτη). Σε αυτή την παράγραφο μαθαίνουμε το ελάχιστο κοινό πολλαπλάσιο (ΕΚΠ) και τον μέγιστο κοινό διαιρέτη (ΜΚΔ), και πώς τα βρίσκουμε με την ανάλυση σε γινόμενο πρώτων παραγόντων και με τον αλγόριθμο του Ευκλείδη.
\end{anagke}

\begin{lexi}{πολλαπλάσιο, κοινό, ελάχιστο, μέγιστο}{πολλαπλάσια, ΕΚΠ, ΜΚΔ}
Το \lexiword{πολλαπλάσιο} ενός αριθμού είναι αυτός που προκύπτει όταν τον πάρουμε πολλές φορές, δηλαδή τον πολλαπλασιάσουμε. \lexiword{Κοινό} είναι αυτό που ανήκει και στους δύο. Με το ΕΚΠ ψάχνουμε το \lexiword{ελάχιστο}, το πιο μικρό, κοινό πολλαπλάσιο, και με το ΜΚΔ τον \lexiword{μέγιστο}, τον πιο μεγάλο, κοινό διαιρέτη. Το $0$ είναι πολλαπλάσιο κάθε αριθμού, γι' αυτό στο ΕΚΠ ζητάμε το μικρότερο \emph{μη μηδενικό} κοινό πολλαπλάσιο.
\end{lexi}

\begin{metafrasi}
Οι αριθμοί που προκύπτουν όταν πολλαπλασιάσω το 4. & $0,4,8,12,16,\dots$ (πολλαπλάσια του $4$)\\
Το πρώτο κοινό πολλαπλάσιο των 4 και 6. & $\EKP(4,6)=12$\\
Το μεγαλύτερο κοινό τους κομμάτι των 12 και 18. & $\MKD(12,18)=6$\\
Δύο αριθμοί δεν έχουν κοινό διαιρέτη εκτός από το 1. & $\MKD(\alpha,\beta)=1$, πρώτοι μεταξύ τους\\
\end{metafrasi}

\tmima{Η θεωρία}

\tmimasmall{Πολλαπλάσια και ΕΚΠ}

\begin{orismos}{πολλαπλάσια}
\textbf{Πολλαπλάσια} ενός φυσικού $\alpha$ είναι οι αριθμοί που προκύπτουν όταν τον πολλαπλασιάσουμε με όλους τους φυσικούς: $0,\ \alpha,\ 2\alpha,\ 3\alpha,\ 4\alpha,\dots$ Για παράδειγμα τα πολλαπλάσια του $9$ είναι τα $0,9,18,27,36,45,\dots$
\end{orismos}

\begin{orismos}{ελάχιστο κοινό πολλαπλάσιο}
Το \textbf{ελάχιστο κοινό πολλαπλάσιο} (ΕΚΠ) δύο ή περισσότερων φυσικών αριθμών που δεν είναι μηδέν είναι το μικρότερο (μη μηδενικό) από τα κοινά πολλαπλάσια των αριθμών αυτών.
\end{orismos}

Παράδειγμα: τα πολλαπλάσια του $4$ είναι $0,4,8,\mathbf{12},16,20,\mathbf{24},28,32,\mathbf{36},\dots$ και του $6$ είναι $0,6,\mathbf{12},18,\mathbf{24},30,\mathbf{36},42,\dots$ Τα κοινά είναι τα $12,24,36,\dots$ (άπειρα) και το μικρότερο είναι το $12$, άρα $\EKP(4,6)=12$.

\tmimasmall{Διαιρέτες και ΜΚΔ}

Από την προηγούμενη παράγραφο ξέρουμε ότι \textbf{διαιρέτες} ενός αριθμού είναι όσοι τον διαιρούν τέλεια. Ισχύουν τα εξής:
\begin{itemize}[leftmargin=1.2em,itemsep=0.2ex,topsep=2pt]
\item Κάθε φυσικός (εκτός από το $0$) διαιρεί τα πολλαπλάσιά του.
\item Κάθε φυσικός που διαιρείται από έναν άλλον είναι πολλαπλάσιο αυτού του αριθμού.
\item Αν ένας φυσικός διαιρεί έναν άλλον, θα διαιρεί και τα πολλαπλάσια του άλλου.
\end{itemize}

\begin{orismos}{μέγιστος κοινός διαιρέτης}
Ο \textbf{μέγιστος κοινός διαιρέτης} (ΜΚΔ) δύο ή περισσότερων φυσικών αριθμών είναι ο μεγαλύτερος από τους κοινούς διαιρέτες τους.
\end{orismos}

Παράδειγμα: οι διαιρέτες του $12$ είναι οι $1,2,3,4,6,12$ και του $18$ οι $1,2,3,6,9,18$. Οι κοινοί διαιρέτες είναι οι $1,2,3,6$ και ο μεγαλύτερος είναι ο $6$, άρα $\MKD(12,18)=6$.

Αν $\MKD(\alpha,\beta)=1$, οι αριθμοί $\alpha$ και $\beta$ λέγονται \textbf{πρώτοι μεταξύ τους}. Για παράδειγμα οι $6$ και $7$ έχουν κοινό διαιρέτη μόνο το $1$, άρα είναι πρώτοι μεταξύ τους (αν και ο $6$ δεν είναι πρώτος αριθμός).

\begin{prosoxi}
Το «πρώτοι μεταξύ τους» δεν σημαίνει ότι οι δύο αριθμοί είναι πρώτοι. Οι $8$ και $15$ δεν είναι πρώτοι αριθμοί, αλλά είναι πρώτοι μεταξύ τους.
\end{prosoxi}

\tmimasmall{ΕΚΠ και ΜΚΔ με ανάλυση σε πρώτους παράγοντες}

Η λίστα με τα πολλαπλάσια ή τους διαιρέτες είναι κουραστική για μεγάλους αριθμούς. Αναλύουμε πρώτα όλους τους αριθμούς σε γινόμενο πρώτων παραγόντων.

\begin{kanonas}{ΕΚΠ και ΜΚΔ με ανάλυση}
Το \textbf{ΕΚΠ} είναι το γινόμενο όλων των πρώτων παραγόντων (κοινών και μη κοινών) με τον \textbf{μεγαλύτερο} εκθέτη τους.\\
Ο \textbf{ΜΚΔ} είναι το γινόμενο μόνο των \textbf{κοινών} πρώτων παραγόντων με τον \textbf{μικρότερο} εκθέτη τους.
\end{kanonas}

\tmimasmall{Ο αλγόριθμος του Ευκλείδη για το ΜΚΔ}

Ο Ευκλείδης έδωσε μια μέθοδο που βρίσκει τον ΜΚΔ χωρίς καμία ανάλυση. Τα βήματα, για τους αριθμούς $125,170,280$:

\begin{center}
\renewcommand{\arraystretch}{1.3}\setlength{\tabcolsep}{9pt}
\begin{tabular}{ccc}
$125$ & $170$ & $280$\\
$125$ & $45$ & $30$\\
$5$ & $15$ & $30$\\
$5$ & $0$ & $0$
\end{tabular}
\end{center}

Γράφουμε τους αριθμούς και βρίσκουμε τον μικρότερο. Στην επόμενη γραμμή γράφουμε \emph{τον μικρότερο αριθμό ως έχει} και κάτω από κάθε άλλον το \emph{υπόλοιπο της διαίρεσής του με τον μικρότερο} ($170=125\cdot 1+45$, $280=125\cdot 2+30$). Επαναλαμβάνουμε με τον νέο μικρότερο μη μηδενικό αριθμό. Σταματάμε όταν όλοι οι άλλοι γίνουν $0$. Ο ΜΚΔ είναι ο μη μηδενικός αριθμός της τελευταίας γραμμής: $\MKD(125,170,280)=5$.

\tmima{Λυμένα παραδείγματα}

\paradeigma{Το λεωφορείο και το τρένο}
\tixero Από το χωριό φεύγει λεωφορείο κάθε $12$ λεπτά και τρένο κάθε $15$ λεπτά. Έφυγαν μαζί στις $8{:}00$.\par
\tithelo Πότε θα ξαναφύγουν μαζί;\par
\possindeo Ζητάω το μικρότερο κοινό πολλαπλάσιο των $12$ και $15$. Πολλαπλάσια του $12$: $12,24,36,48,60,\dots$ Πολλαπλάσια του $15$: $15,30,45,60,\dots$ Το πρώτο κοινό είναι το $60$.\par
\apantisi $\EKP(12,15)=60$, άρα μετά από $60$ λεπτά, στις $9{:}00$.\par
\elegxos Το $60=12\cdot 5=15\cdot 4$.

\paradeigma{Οι χάντρες της Αγγελικής}
\tixero Η Αγγελική έχει $36$ κόκκινες, $48$ κίτρινες και $60$ μπλε χάντρες. Φτιάχνει όσο το δυνατόν περισσότερα ίδια κολιέ, με τον ίδιο αριθμό χαντρών κάθε χρώματος, χωρίς να περισσέψει καμία.\par
\tithelo Πόσα κολιέ φτιάχνει και τι έχει το καθένα;\par
\possindeo Ο αριθμός των κολιέ πρέπει να διαιρεί και τα τρία πλήθη, και ζητάω τον μεγαλύτερο, άρα τον ΜΚΔ. Αναλύω: $36=2^2\cdot 3^2$, $48=2^4\cdot 3$, $60=2^2\cdot 3\cdot 5$. Κοινοί παράγοντες: $2$ και $3$ με μικρότερους εκθέτες $2$ και $1$. Άρα $\MKD=2^2\cdot 3=12$. Κάθε κολιέ έχει $36:12=3$ κόκκινες, $48:12=4$ κίτρινες και $60:12=5$ μπλε.\par
\apantisi $12$ κολιέ, με $3$ κόκκινες, $4$ κίτρινες και $5$ μπλε χάντρες.\par
\elegxos $12\cdot 3=36$, $12\cdot 4=48$, $12\cdot 5=60$.

\paradeigma{ΕΚΠ και ΜΚΔ με ανάλυση, και με Ευκλείδη}
\tixero Οι αριθμοί $84$, $90$ και $126$.\par
\tithelo Το ΕΚΠ και ο ΜΚΔ τους.\par
\possindeo Αναλύω: $84=2^2\cdot 3\cdot 7$, $90=2\cdot 3^2\cdot 5$, $126=2\cdot 3^2\cdot 7$. Όλοι οι πρώτοι παράγοντες είναι οι $2,3,5,7$, με μεγαλύτερους εκθέτες $2,2,1,1$: $\EKP=2^2\cdot 3^2\cdot 5\cdot 7=1.260$. Κοινοί παράγοντες είναι μόνο οι $2$ και $3$, με μικρότερους εκθέτες $1,1$: $\MKD=2\cdot 3=6$. Με τον αλγόριθμο του Ευκλείδη: $(84,90,126)\to(84,6,42)\to(0,6,0)$, άρα ΜΚΔ $=6$.\par
\apantisi $\EKP=1.260$ και $\MKD=6$.\par
\elegxos Το $1.260:84=15$, $1.260:90=14$, $1.260:126=10$ και ο $6$ διαιρεί τους $84,90,126$.

\tmima{Ασκήσεις}

\zoni{Το κατάλαβα}{ορισμοί και υπολογισμοί}

\ex Συμπλήρωσε: α) Αν $\MKD(\alpha,\beta)=1$, οι $\alpha$ και $\beta$ λέγονται \dots μεταξύ τους. β) Κάθε φυσικός \dots τα πολλαπλάσιά του. γ) Κάθε φυσικός που διαιρείται από έναν άλλον είναι \dots αυτού του αριθμού. δ) Ο ΜΚΔ δύο αριθμών είναι ο \dots από τους κοινούς διαιρέτες τους.

\ex Γράψε τα $8$ πρώτα πολλαπλάσια των $7$ και $12$ (ξεκινώντας από το $0$).

\ex Γράψε τα $8$ πρώτα μη μηδενικά πολλαπλάσια των $6$ και $9$. Ποια είναι κοινά; Ποιο είναι το $\EKP(6,9)$;

\ex Γράψε όλους τους διαιρέτες των $18$, $30$ και $36$. Βρες τους κοινούς διαιρέτες και τον $\MKD(18,30,36)$.

\ex Είναι πρώτοι μεταξύ τους οι αριθμοί; α) $8$ και $15$ \quad β) $14$ και $21$ \quad γ) $22$ και $45$ \quad δ) $9$ και $28$ \quad ε) $12$ και $35$.

\ex Συμπλήρωσε τον πίνακα με τον ΜΚΔ και το ΕΚΠ κάθε ζεύγους: $(6,8)$, $(10,15)$, $(4,9)$, $(14,21)$, $(16,24)$.

\ex Ανάλυσε τους $36$ και $54$ σε γινόμενο πρώτων παραγόντων και βρες τον ΜΚΔ και το ΕΚΠ τους.

\zoni{Το χρησιμοποιώ}{ΕΚΠ ή ΜΚΔ;}

\ex \textbf{Οι προπονήσεις.} Ο Νίκος προπονείται κάθε $3$ ημέρες, ο Βασίλης κάθε $5$ και η Πολυξένη κάθε $6$. Σήμερα προπονήθηκαν όλοι μαζί. Μετά από πόσες ημέρες θα ξανασυναντηθούν;

\ex \textbf{Τα καλάθια.} Η θεία Μαρία έχει $54$ μήλα και $72$ αχλάδια και τα βάζει σε ίσα καλάθια, χωρίς να περισσέψει τίποτα και με τον ίδιο αριθμό μήλων και αχλαδιών σε κάθε καλάθι. Πόσα το πολύ καλάθια μπορεί να φτιάξει και τι έχει το καθένα;

\ex \textbf{Τα πλακάκια.} Το δωμάτιο έχει διαστάσεις $240\,\mathrm{cm}\times 180\,\mathrm{cm}$ και θέλουμε να το στρώσουμε με τετράγωνα πλακάκια, ίδια, χωρίς κοψίματα. Ποια είναι η μεγαλύτερη δυνατή πλευρά του πλακιδίου και πόσα πλακάκια χρειάζονται;

\ex \textbf{Τα φωτάκια.} Στη γιορτή του χωριού ένα κόκκινο φωτάκι ανάβει κάθε $12$ δευτερόλεπτα και ένα πράσινο κάθε $18$. Άναψαν μαζί. Μετά από πόσα δευτερόλεπτα θα ανάψουν ξανά μαζί;

\ex \textbf{Οι σελιδοδείκτες.} Η Φανή έχει $40$ σελιδοδείκτες και $56$ αυτοκόλλητα. Φτιάχνει ίδια πακέτα, με τον ίδιο αριθμό από κάθε είδος, χωρίς να περισσέψει τίποτα. Πόσα το πολύ πακέτα φτιάχνει και τι έχει κάθε πακέτο;

\ex \textbf{Ο φούρνος.} Ένας φούρνος έχει $60$ κουλούρια, $84$ κρουασάν και $108$ τσουρέκια και τα βάζει σε ίδια κουτιά, με τον ίδιο αριθμό από κάθε είδος και χωρίς να περισσέψει τίποτα. Πόσα το πολύ κουτιά γεμίζει και τι έχει το καθένα;

\zoni{Το σκέφτομαι}{ερωτήσεις με λίγη περισσότερη σκέψη}

\ex \textbf{Μια ανακάλυψη.} Για τα ζεύγη $(4,6)$, $(12,18)$ και $(10,15)$ βρες τον ΜΚΔ και το ΕΚΠ. Μετά υπολόγισε το γινόμενο $\MKD\cdot\EKP$ και το γινόμενο των δύο αριθμών κάθε ζεύγους. Τι παρατηρείς;

\ex Βρες δύο διαφορετικούς φυσικούς που έχουν $\MKD=6$ και $\EKP=60$.

\ex α) Βρες το $\EKP(8,15)$ και το $\EKP(9,10)$. Πώς σχετίζονται με το γινόμενο των αριθμών, όταν οι αριθμοί είναι πρώτοι μεταξύ τους; β) Βρες τον $\MKD$ και το $\EKP$ των $6$ και $30$ (ο ένας διαιρεί τον άλλον). Τι παρατηρείς;

\zoni{Για το σπίτι}{εξάσκηση}

\ex Βρες με ανάλυση τον ΜΚΔ και το ΕΚΠ των $325$ και $338$ (από την παράγραφο 1.3) και των $24$, $42$, $54$.

\ex Βρες με τον αλγόριθμο του Ευκλείδη τον ΜΚΔ των $48$, $80$, $120$. Μετά εξέτασε αν οι $1.000$ και $1.303$ είναι πρώτοι μεταξύ τους.

\paragrafos{Χαρακτήρες διαιρετότητας, πρώτοι αριθμοί, ΜΚΔ και ΕΚΠ}{Α.1.5, «Χαρακτήρες διαιρετότητας – ΜΚΔ – ΕΚΠ – Ανάλυση αριθμού σε γινόμενο πρώτων παραγόντων»}
\ypotmima{Α}{Χαρακτήρες διαιρετότητας: ξέρω χωρίς να διαιρέσω}


\begin{anagke}
Το σχολικό παζάρι πήγε καλύτερα απ' όσο περίμενε κανείς. Η Φανή, που κρατάει το ταμείο, μετράει τα χρήματα: 2.358 €. Θέλουν να τα μοιράσουν ίσα σε τρεις ομάδες που θα φτιάξουν από ένα σχολικό έργο. Μετά η Ρήνη λέει: «Και αν οι ομάδες ήταν εννιά;» Ο Νίκος: «Και αν ήταν πέντε;» Ο Στέλιος ξεκινάει να κάνει τις διαιρέσεις με τη σειρά, μία μία. Τρεις διαιρέσεις με τετραψήφιο αριθμό είναι πολλή δουλειά.

Η Πολυξένη κοιτάζει τον αριθμό 2.358 και λέει: «Για το πέντε δεν χρειάζεται καν να κάνεις διαίρεση. Τελειώνει σε οκτώ. Δεν γίνεται». «Πώς το ξέρεις;» ρωτάει ο Ηλίας. «Το έχω δει σε όλους τους αριθμούς που διαιρούνται με το πέντε: τελειώνουν σε 0 ή σε 5». Η Ζωή ρωτάει: «Και για το τρία και για το εννιά υπάρχει κάτι τέτοιο;»

Υπάρχει. Και είναι από τα πιο χρήσιμα κόλπα στο γυμνάσιο: να καταλαβαίνεις αν ένας αριθμός διαιρείται με έναν άλλον κοιτάζοντας απλώς τα ψηφία του.
\end{anagke}

\begin{lexi}{κριτήριο, χαρακτήρας διαιρετότητας}{κριτήριο, χαρακτήρας}
Το \lexiword{κριτήριο} έρχεται από το «κρίνω», δηλαδή αποφασίζω. Κριτήριο είναι ο κανόνας με τον οποίο κρίνουμε κάτι. Ένας κριτής σε αγώνα έχει κριτήρια για να βγάλει τον νικητή. Ένα \lexiword{κριτήριο διαιρετότητας} μας βοηθάει να κρίνουμε αν ένας αριθμός διαιρείται με κάποιον άλλον, χωρίς να κάνουμε τη διαίρεση.

Το σχολικό βιβλίο τα λέει και \lexiword{χαρακτήρες διαιρετότητας}. Ο χαρακτήρας είναι το «σημάδι» που χαρακτηρίζει κάτι, όπως ο χαρακτήρας ενός ανθρώπου. Εδώ το σημάδι είναι στα ψηφία του αριθμού.
\end{lexi}

\begin{metafrasi}
Ο αριθμός είναι άρτιος. & Λήγει σε $0,\,2,4,\,6$ ή $8$ (διαιρείται με το 2)\\
Ο αριθμός διαιρείται με το 5. & Λήγει σε $0$ ή $5$\\
Ο αριθμός διαιρείται με το 10. & Λήγει σε $0$\\
Ο αριθμός διαιρείται με το 3. & Το άθροισμα των ψηφίων του διαιρείται με το 3\\
Ο αριθμός διαιρείται με το 9. & Το άθροισμα των ψηφίων του διαιρείται με το 9\\
Ο αριθμός διαιρείται με το 4. & Ο αριθμός από τα δύο τελευταία ψηφία διαιρείται με το 4\\
Ο αριθμός διαιρείται με το 25. & Λήγει σε $00$, $25$, $50$ ή $75$\\
Ο αριθμός διαιρείται με το 6. & Διαιρείται και με το 2 και με το 3\\
\end{metafrasi}

\tmima{Η θεωρία}

\begin{kanonas}{κριτήρια διαιρετότητας}
\begin{itemize}[leftmargin=1.2em,itemsep=0.3ex,topsep=2pt]
\item \textbf{Με το 2}: το τελευταίο ψηφίο είναι 0, 2, 4, 6 ή 8.
\item \textbf{Με το 5}: το τελευταίο ψηφίο είναι 0 ή 5.
\item \textbf{Με το 10}: το τελευταίο ψηφίο είναι 0. (Με το 100: τα δύο τελευταία ψηφία είναι 00, με το 1.000: τα τρία τελευταία, κ.ο.κ.)
\item \textbf{Με το 4}: ο αριθμός που σχηματίζουν τα δύο τελευταία ψηφία διαιρείται με το 4 (ή είναι 00).
\item \textbf{Με το 25}: τα δύο τελευταία ψηφία είναι 00, 25, 50 ή 75.
\item \textbf{Με το 3}: το άθροισμα των ψηφίων διαιρείται με το 3.
\item \textbf{Με το 9}: το άθροισμα των ψηφίων διαιρείται με το 9.
\end{itemize}
\end{kanonas}

Από τα κριτήρια αυτά φτιάχνουμε και άλλα. Για να διαιρείται ένας αριθμός με το 6 πρέπει να διαιρείται και με το 2 και με το 3. Για το 15, και με το 3 και με το 5. Για το 12, και με το 3 και με το 4.

\tmimasmall{Γιατί δουλεύουν}

Κάθε αριθμός ξεκινάει με πολλαπλάσια του 10, του 100, του 1.000. Για παράδειγμα $3.725=37\cdot 100+25$. Ο 100 διαιρείται με το 4 και με το 25, άρα το $37\cdot 100$ διαιρείται και με τους δύο, και το μόνο που μένει να ελέγξουμε είναι τα δύο τελευταία ψηφία, το 25. Το ίδιο και για τα 2 και 5, με το 10.

Για το 3 και το 9 δουλεύει ένα κόλπο. Οι αριθμοί 9, 99, 999, \dots\ διαιρούνται με το 9 (και με το 3). Γράφω $2.457=2\cdot 1.000+4\cdot 100+5\cdot 10+7$ και αντικαθιστώ το $1.000=999+1$, $100=99+1$, $10=9+1$:
$$2.457=2\cdot 999+4\cdot 99+5\cdot 9+(2+4+5+7).$$
Τα τρία πρώτα κομμάτια διαιρούνται με το 9. Άρα όλος ο αριθμός διαιρείται με το 9 ακριβώς όταν διαιρείται το $2+4+5+7=18$. Και διαιρείται.

\begin{prosoxi}
\begin{itemize}[leftmargin=1.2em,itemsep=0.3ex,topsep=2pt]
\item Αν ένας αριθμός τελειώνει σε 4, δεν σημαίνει ότι διαιρείται με το 4. Το 14 δεν διαιρείται. Για το 4 κοιτάμε τα δύο τελευταία ψηφία μαζί.
\item Για το 3 και το 9 μετράνε \textbf{όλα} τα ψηφία, όχι μόνο το τελευταίο.
\item Για να διαιρείται ένας αριθμός με το 24 δεν φτάνει να διαιρείται με το 4 και με το 6. Το 12 διαιρείται και με τα δύο, αλλά όχι με το 24. Μπορούμε να «σπάσουμε» έναν αριθμό σε δύο μόνο όταν οι δύο αριθμοί δεν έχουν κοινό παράγοντα (όπως το 3 και το 4, ή το 2 και το 5).
\item Το 0 διαιρείται με κάθε αριθμό (εκτός από το 0 ως διαιρέτη), γιατί $0=\delta\cdot 0$.
\end{itemize}
\end{prosoxi}

\tmima{Λυμένα παραδείγματα}

\paradeigma{Ο πίνακας των κριτηρίων}
\tixero Οι αριθμοί 4.860, 975, 1.236 και 7.500.\par
\tithelo Με ποιους από τους 2, 3, 4, 5, 9, 10, 25 και 100 διαιρείται ο καθένας.\par
\possindeo Εφαρμόζω κάθε κριτήριο. Στα 3 και 9 υπολογίζω το άθροισμα ψηφίων. Στα 4 και 25 κοιτάζω τα δύο τελευταία ψηφία.

\begin{center}
\small\renewcommand{\arraystretch}{1.3}\setlength{\tabcolsep}{4pt}
\begin{tabular}{@{}lcccccccc@{}}
\toprule
 & 2 & 3 & 4 & 5 & 9 & 10 & 25 & 100\\
\midrule
4.860 \ ($4+8+6+0=18$) & ναι & ναι & ναι & ναι & ναι & ναι & όχι & όχι\\
975 \ ($9+7+5=21$) & όχι & ναι & όχι & ναι & όχι & όχι & ναι & όχι\\
1.236 \ ($1+2+3+6=12$) & ναι & ναι & ναι & όχι & όχι & όχι & όχι & όχι\\
7.500 \ ($7+5=12$) & ναι & ναι & ναι & ναι & όχι & ναι & ναι & ναι\\
\bottomrule
\end{tabular}
\end{center}

\apantisi Δες τον πίνακα.\par
\elegxos Το 7.500 διαιρείται με το 100, άρα και με το 10, το 4, το 25, το 5 και το 2: κάθε διαιρέτης του 100 μπαίνει αυτόματα. Δοκιμή με το 4: $7.500:4=1.875$. Σωστά. Για το 1.236 με το 4: $36:4=9$, και πράγματι $1.236:4=309$.

\paradeigma{Το ψηφίο που λείπει}
\tixero Ο αριθμός $3.8\square 1$, όπου το $\square$ είναι ένα ψηφίο.\par
\tithelo Ποιο ψηφίο πρέπει να μπει ώστε ο αριθμός να διαιρείται με το 9.\par
\possindeo Το άθροισμα των γνωστών ψηφίων είναι $3+8+1=12$. Πρέπει $12+\square$ να διαιρείται με το 9. Το επόμενο πολλαπλάσιο του 9 μετά το 12 είναι το 18, άρα $\square=6$. (Το επόμενο, 27, θα απαιτούσε ψηφίο 15, που δεν υπάρχει.)\par
\apantisi $3.861$.\par
\elegxos $3.861:9=429$, γιατί $9\cdot 429=3.861$. Σωστά.

\paradeigma{Η μοιρασιά του παζαριού}
\tixero Η Φανή έχει 2.358 €.\par
\tithelo Σε ποιες από τις 3, 9, 5, 4 και 6 ισόποσες ομάδες μοιράζονται χωρίς να περισσεύει τίποτα.\par
\possindeo Άθροισμα ψηφίων: $2+3+5+8=18$. Διαιρείται με το 3 και με το 9. Τελευταίο ψηφίο 8, άρα όχι με το 5. Τελευταία δύο ψηφία 58: το $58:4$ δεν είναι ακέραιο, άρα όχι με το 4. Είναι άρτιος και διαιρείται με το 3, άρα διαιρείται με το 6.\par
\apantisi Σε 3, 6 και 9 ομάδες.\par
\elegxos $2.358:3=786$, $2.358:6=393$, $2.358:9=262$. Και οι τρεις είναι τέλειες διαιρέσεις.

\tmima{Ασκήσεις}

\zoni{Το κατάλαβα}{εφαρμογή του κανόνα}

\ex Ποιοι από τους αριθμούς 2.480, 3.115, 6.002, 1.050 και 819 διαιρούνται με το 2; Ποιοι με το 5; Ποιοι με το 10;

\ex Βρες ποιοι από τους 252, 3.614, 8.001, 55.551 και 1.370 διαιρούνται με το 3. Έλεγξε μία από τις περιπτώσεις με διαίρεση.

\ex Βρες ποιοι από τους 738, 5.427, 92.808, 6.661 και 14.445 διαιρούνται με το 9.

\ex Βρες ποιοι από τους 524, 1.018, 6.300, 3.772 και 90 διαιρούνται με το 4.

\ex Βρες ποιοι από τους 450, 2.125, 3.300, 1.080 και 7.675 διαιρούνται με το 25.

\ex Συμπλήρωσε έναν πίνακα με τους αριθμούς 360, 1.125, 5.724 και 40.500 και τους διαιρέτες 2, 3, 4, 5, 9, 10, 25. Για κάθε αριθμό γράψε ναι ή όχι.

\ex Ποιοι από τους 78, 250, 3.402, 1.113 και 9.996 διαιρούνται με το 6;

\ex Βρες το ψηφίο (ή τα ψηφία) που πρέπει να μπουν στη θέση του $\square$: α) $6.3\square 8$ να διαιρείται με το 9 \quad β) $4\square 2$ να διαιρείται με το 3 \quad γ) $2\square 6$ να διαιρείται με το 4 \quad δ) $29\square$ να διαιρείται και με το 2 και με το 5.

\ex Σωστό ή λάθος; α) Κάθε αριθμός που τελειώνει σε 6 διαιρείται με το 6. \quad β) Αν ένας αριθμός διαιρείται με το 9, διαιρείται και με το 3. \quad γ) Το 0 διαιρείται με το 13. \quad δ) Αν ένας αριθμός διαιρείται με το 6 και με το 8, διαιρείται και με το 48. \quad ε) Αν διαιρείται με το 3 και με το 5, διαιρείται με το 15. \quad στ) Αν διαιρείται με το 4 και με το 25, διαιρείται με το 100.

\zoni{Το χρησιμοποιώ}{προβλήματα από τη ζωή}

\ex \textbf{Τα έσοδα της λαχειοφόρου.} Η λαχειοφόρος του σχολείου έφερε 4.752 €. Σε ποιες από τις 3, 4, 5, 6 και 9 ομάδες μοιράζονται ακριβώς; Βρες πόσα παίρνει η κάθε ομάδα, όπου γίνεται.

\ex \textbf{Τα μπλουζάκια.} Ο προπονητής του Νίκου έχει 150 μπλουζάκια και θέλει να τα βάλει σε ίσες στοίβες. Με ποιους από τους αριθμούς 2, 3, 4, 5, 6 και 10 γίνεται, χωρίς να περισσεύει κανένα; Χρησιμοποίησε κριτήρια.

\ex \textbf{Ο κωδικός της Πολυξένης.} Ο νέος κωδικός της Πολυξένης είναι τετραψήφιος. Ξεκινάει με 5, το δεύτερο ψηφίο είναι το 4, τελειώνει σε 0, διαιρείται με το 9 και το τρίτο ψηφίο είναι μεγαλύτερο από το δεύτερο. Ποιος είναι ο κωδικός;

\ex \textbf{Οι θέσεις του γηπέδου.} Οι θέσεις ενός μικρού γηπέδου είναι αριθμημένες από το 1 ως το 90. Οι θέσεις που διαιρούνται με το 5 έχουν σκιά. Πόσες είναι; Πόσες διαιρούνται με το 10; Πόσες διαιρούνται και με το 3 και με το 5;

\ex \textbf{Τα δίσεκτα έτη.} Ένα έτος είναι δίσεκτο αν διαιρείται με το 4. Εξαίρεση οι αιώνες (1.900, 2.000, \dots), που είναι δίσεκτοι μόνο όταν διαιρούνται με το 400. Ποια από τα έτη 1.900, 2.000, 2.024, 2.030 και 2.100 είναι δίσεκτα;

\zoni{Το σκέφτομαι}{μια ερώτηση που πάει πιο μακριά}

\ex Γιατί ένας αριθμός που τελειώνει σε 0 διαιρείται και με το 2 και με το 5; (Βοήθεια: $10=2\cdot 5$.)

\ex Γράψε το 2.763 ως $2\cdot 999+7\cdot 99+6\cdot 9+(2+7+6+3)$ και εξήγησε με τη βοήθεια αυτής της γραφής αν διαιρείται με το 3 και αν διαιρείται με το 9.

\ex Ποιος είναι ο μικρότερος τετραψήφιος αριθμός που διαιρείται ταυτόχρονα με τα 4, 5 και 6;

\ex Οι αριθμοί 121, 132, 143 και 154 διαιρούνται όλοι με το 11. Τι παρατηρείς στα ψηφία τους; Δοκίμασε το ίδιο μοτίβο στο 165 και στο 175 και έλεγξε με διαίρεση.

\ex Ο Στέλιος ισχυρίζεται ότι ο 5.034 διαιρείται με το 4, γιατί τελειώνει σε 4. Ποιο είναι το λάθος του; Τι θα έπρεπε να είχε κάνει;

\ypotmima{Β}{Πρώτοι αριθμοί και ανάλυση: τα δομικά υλικά}

\begin{anagke}
Στο διάλειμμα η Πολυξένη, που κρατάει ένα τετράδιο με κωδικούς και μοτίβα, δείχνει στην παρέα κάτι που της εξήγησε η μαμά της, που δουλεύει στην ασφάλεια ενός δικτύου υπολογιστών. «Σκέφτομαι δύο πρώτους αριθμούς», λέει στη Ρήνη. «Το γινόμενό τους είναι 143. Ποιοι είναι;»

Η Ρήνη δοκιμάζει. Δεν διαιρείται με το 2, γιατί τελειώνει σε 3. Το άθροισμα ψηφίων $1+4+3=8$, άρα ούτε με το 3. Ούτε με το 5. Με το 7, $7\cdot 20=140$ και περισσεύουν 3, όχι. Με το 11: $11\cdot 13=143$. «Έντεκα και δεκατρία!» φωνάζει. «Μπράβο», λέει η Πολυξένη. «Τώρα σκέψου ότι ο αριθμός έχει 600 ψηφία. Ένας υπολογιστής πολλαπλασιάζει τους δύο πρώτους σε ένα κλάσμα του δευτερολέπτου. Για να τους βρει από το γινόμενο, θέλει απίστευτα πολύ χρόνο. Πάνω σε αυτό στηρίζονται οι κωδικοί των τραπεζών».

Η Ζωή ρωτάει κάτι πιο απλό. «Δηλαδή κάθε αριθμός είναι φτιαγμένος από πρώτους, όπως ο τοίχος από τούβλα;» Ναι. Και το να βρίσκεις αυτά τα τούβλα είναι δουλειά με μέθοδο.
\end{anagke}

\begin{lexi}{πρώτος, σύνθετος, παράγοντας}{πρώτος, σύνθετος}
Ο \lexiword{πρώτος} αριθμός δεν λέγεται έτσι επειδή έρχεται πρώτος στη σειρά, αλλά επειδή είναι «πρωταρχικός», πρωτογενής: δεν φτιάχνεται από μικρότερους αριθμούς με πολλαπλασιασμό. Είναι το δομικό υλικό.

Ο \lexiword{σύνθετος} αριθμός είναι αυτός που συντίθεται, δηλαδή φτιάχνεται, από άλλους: το 12 είναι $2\cdot 2\cdot 3$. Ο \lexiword{παράγοντας} είναι ο αριθμός που παίρνει μέρος σε έναν πολλαπλασιασμό και «παράγει» το γινόμενο.
\end{lexi}

\begin{metafrasi}
Το 7 έχει διαιρέτες μόνο το 1 και το 7. & Το 7 είναι πρώτος\\
Το 12 έχει διαιρέτες τους 1, 2, 3, 4, 6, 12. & Το 12 είναι σύνθετος\\
Το 60 σαν γινόμενο πρώτων. & $60=2\cdot 2\cdot 3\cdot 5=2^2\cdot 3\cdot 5$\\
Ο $\alpha$ διαιρεί τον $\beta$. & $\beta=\alpha\cdot\kappa$, με $\kappa$ φυσικό\\
Πολλαπλάσια του 4. & $0,\,4,\,8,\,12,\,16,\,20,\,\dots$\\
\end{metafrasi}

\tmima{Η θεωρία}

\tmimasmall{Πολλαπλάσια και διαιρέτες}

\textbf{Πολλαπλάσια} ενός φυσικού $\alpha$ είναι οι αριθμοί που βγαίνουν αν τον πολλαπλασιάσουμε με φυσικούς: $0,\ \alpha,\ 2\alpha,\ 3\alpha,\ 4\alpha,\ \dots$ \textbf{Διαιρέτες} του $\alpha$ είναι όλοι οι αριθμοί που τον διαιρούν ακριβώς. Κάθε αριθμός $\alpha$ έχει διαιρέτες τουλάχιστον τον 1 και τον εαυτό του. Αν ο $\delta$ διαιρεί έναν αριθμό, τότε ο αριθμός είναι πολλαπλάσιο του $\delta$. Είναι το ίδιο πράγμα, από δύο πλευρές.

\begin{orismos}{πρώτος και σύνθετος}
\textbf{Πρώτος} λέγεται ένας φυσικός αριθμός μεγαλύτερος από το 1 που έχει διαιρέτες μόνο τον 1 και τον εαυτό του. Ένας φυσικός μεγαλύτερος από το 1 που έχει και άλλους διαιρέτες λέγεται \textbf{σύνθετος}. Το 1 δεν είναι ούτε πρώτος ούτε σύνθετος.
\end{orismos}

Οι πρώτοι μέχρι το 50 είναι οι: \blue{2, 3, 5, 7, 11, 13, 17, 19, 23, 29, 31, 37, 41, 43, 47}. Ο 2 είναι ο μοναδικός άρτιος πρώτος, γιατί κάθε άλλος άρτιος διαιρείται και με το 2. Οι πρώτοι είναι άπειροι: το απέδειξε ο Ευκλείδης πριν από 2.300 χρόνια.

\tmimasmall{Το κόσκινο του Ερατοσθένη}

Ο Ερατοσθένης, που έζησε στην Αλεξάνδρεια τον 3ο αιώνα π.Χ., βρήκε μια μέθοδο για να «κοσκινίζουμε» τους πρώτους. Γράφεις τους αριθμούς, σβήνεις το 1, μετά κρατάς το 2 και σβήνεις όλα τα πολλαπλάσιά του, μετά κρατάς το επόμενο που δεν έσβησες, το 3, και σβήνεις τα πολλαπλάσιά του, και ούτω καθεξής. Αυτά που μένουν είναι οι πρώτοι.

\begin{center}
\renewcommand{\arraystretch}{1.25}
\setlength{\tabcolsep}{5pt}
\newcommand{\pr}[1]{\textcolor{gblue}{\bfseries #1}}
\begin{tabular}{cccccccccc}
1 & \pr{2} & \pr{3} & 4 & \pr{5} & 6 & \pr{7} & 8 & 9 & 10\\
\pr{11} & 12 & \pr{13} & 14 & 15 & 16 & \pr{17} & 18 & \pr{19} & 20\\
21 & 22 & \pr{23} & 24 & 25 & 26 & 27 & 28 & \pr{29} & 30\\
\pr{31} & 32 & 33 & 34 & 35 & 36 & \pr{37} & 38 & 39 & 40\\
\pr{41} & 42 & \pr{43} & 44 & 45 & 46 & \pr{47} & 48 & 49 & 50\\
\end{tabular}
\end{center}

\tmimasmall{Ανάλυση σε γινόμενο πρώτων παραγόντων}

Κάθε σύνθετος αριθμός γράφεται ως γινόμενο πρώτων, και \textbf{με έναν μόνο τρόπο} (αν αγνοήσουμε τη σειρά των παραγόντων). Γι' αυτό οι πρώτοι είναι τα δομικά υλικά: κάθε αριθμός έχει το δικό του «σύνολο υλικών».

\begin{kanonas}{ανάλυση σε πρώτους παράγοντες}
\begin{enumerate}[leftmargin=1.6em,itemsep=0.2ex,topsep=2pt]
\item Βρίσκω τον μικρότερο πρώτο που διαιρεί τον αριθμό (δοκιμάζω με τη σειρά 2, 3, 5, 7, 11, \dots με τα κριτήρια διαιρετότητας).
\item Διαιρώ και γράφω τον πρώτο στα δεξιά και το πηλίκο από κάτω.
\item Επαναλαμβάνω με το πηλίκο, ώσπου να φτάσω στο 1.
\item Γράφω το γινόμενο των πρώτων και τους ίδιους παράγοντες τους ενώνω σε δυνάμεις.
\end{enumerate}
\end{kanonas}

Αν ο αριθμός είναι πρώτος, δεν χρειάζεται να δοκιμάσεις όλους τους πρώτους. Δοκιμάζεις με τη σειρά και σταματάς όταν το \textbf{πηλίκο γίνει μικρότερο από τον διαιρέτη}. Αν μέχρι τότε δεν βρήκες διαιρέτη, ο αριθμός είναι πρώτος.

\begin{prosoxi}
\begin{itemize}[leftmargin=1.2em,itemsep=0.3ex,topsep=2pt]
\item Το 1 \textbf{δεν είναι πρώτος}. Αν ήταν, το 6 θα γραφόταν και $2\cdot 3$ και $1\cdot 2\cdot 3$ και $1\cdot 1\cdot 2\cdot 3$, και η ανάλυση δεν θα ήταν μοναδική.
\item Ο 2 είναι πρώτος, και μάλιστα ο μόνος άρτιος. Όχι «όλοι οι πρώτοι είναι περιττοί».
\item Ο αριθμός που μοιάζει πρώτος δεν είναι πάντα. Το 91 φαίνεται πρώτος, αλλά $91=7\cdot 13$.
\end{itemize}
\end{prosoxi}

\tmima{Λυμένα παραδείγματα}

\paradeigma{Πρώτος ή σύνθετος;}
\tixero Οι αριθμοί 91 και 97.\par
\tithelo Να αποφασίσω αν είναι πρώτοι.\par
\possindeo Για το 91: δεν διαιρείται με το 2 (τελειώνει σε 1), ούτε με το 3 ($9+1=10$), ούτε με το 5. Με το 7: $7\cdot 13=91$. Βρήκα διαιρέτη, άρα είναι σύνθετος.\par
Για το 97: ούτε με 2, ούτε με 3 ($9+7=16$), ούτε με 5. Με το 7: $97=7\cdot 13+6$, δεν διαιρείται, και το πηλίκο 13 είναι ακόμα μεγαλύτερο από το 7. Με το 11: $97=11\cdot 8+9$, και το πηλίκο 8 είναι μικρότερο από το 11. Σταματάω.\par
\apantisi Το 91 είναι σύνθετος ($7\cdot 13$), το 97 είναι πρώτος.\par
\elegxos Οι πρώτοι κάτω από το 100 είναι 25 και ο τελευταίος από αυτούς είναι το 97.

\paradeigma{Ανάλυση του 1.980}
\tixero Ο αριθμός 1.980.\par
\tithelo Να τον αναλύσω σε γινόμενο πρώτων παραγόντων.\par
\possindeo Διαιρώ διαδοχικά: $1.980:2=990$, $990:2=495$. Το 495 είναι περιττός, άρα δοκιμάζω με το 3: $495:3=165$, $165:3=55$. Μετά με το 5: $55:5=11$. Τέλος $11:11=1$.\par
Οι παράγοντες είναι δύο δυάρια, δύο τριάρια, ένα πεντάρι και ένα ενδεκάρι.\par
\apantisi $1.980=2^2\cdot 3^2\cdot 5\cdot 11$.\par
\elegxos $4\cdot 9=36$, $36\cdot 5=180$, $180\cdot 11=1.980$. Σωστά.

\paradeigma{Το μυστικό της Πολυξένης}
\tixero Το 143 είναι το γινόμενο δύο πρώτων.\par
\tithelo Ποιοι είναι αυτοί οι πρώτοι.\par
\possindeo Δοκιμάζω τους πρώτους με τη σειρά. Το 2, το 3, το 5 και το 7 δεν διαιρούν το 143. Με το 11: $143:11=13$, και το 13 είναι πρώτος.\par
\apantisi $143=11\cdot 13$.\par
\elegxos $11\cdot 13=143$. Και οι δύο παράγοντες είναι πρώτοι, άρα αυτή είναι η ανάλυση.

\tmima{Ασκήσεις}

\zoni{Το κατάλαβα}{εφαρμογή του κανόνα}

\ex Γράψε όλους τους διαιρέτες των αριθμών 16, 20, 28, 42 και 49.

\ex Γράψε τα έξι πρώτα πολλαπλάσια (μαζί με το 0) των αριθμών 3, 8 και 15.

\ex Ποιοι από τους 1, 3, 15, 23, 27, 31, 39, 43, 49, 51, 59, 63, 77 και 83 είναι πρώτοι; Για τους σύνθετους, γράψε έναν διαιρέτη τους.

\ex Βρες όλους τους πρώτους αριθμούς ανάμεσα στο 80 και στο 110.

\ex Ανάλυσε σε γινόμενο πρώτων παραγόντων (με δυνάμεις): α) 36 \quad β) 90 \quad γ) 126 \quad δ) 200 \quad ε) 360 \quad στ) 462.

\ex Ποιοι αριθμοί έχουν ανάλυση: α) $2^4\cdot 3$ \quad β) $2\cdot 5^2\cdot 7$ \quad γ) $3^2\cdot 5\cdot 7$ \quad δ) $2^3\cdot 3\cdot 11$;

\ex Σωστό ή λάθος; α) Ο 2 είναι ο μοναδικός άρτιος πρώτος. \quad β) Το 1 είναι πρώτος. \quad γ) Το 77 είναι πρώτος. \quad δ) Το άθροισμα δύο πρώτων είναι πάντα πρώτος. \quad ε) Το γινόμενο δύο πρώτων είναι πάντα σύνθετος. \quad στ) Το 0 είναι πολλαπλάσιο κάθε αριθμού. \quad ζ) Το 57 είναι πρώτος.

\zoni{Το χρησιμοποιώ}{προβλήματα από τη ζωή}

\ex \textbf{Το μυστικό της Φανής.} Η Φανή σκέφτεται δύο πρώτους με γινόμενο 323. Ποιοι είναι; Δοκίμασε με τη σειρά τους πρώτους μέχρι να βρεις διαιρέτη.

\ex \textbf{Οι καραμέλες της Μυρτώς.} Η Ρήνη έχει 54 καραμέλες και θέλει να τις βάλει σε ίσα σακουλάκια, χωρίς να περισσέψει καμία. Πόσες καραμέλες μπορεί να βάλει σε κάθε σακουλάκι; Η Ρήνη έχει και ένα άλλο βάζο με 53 καραμέλες. Με πόσους τρόπους μπορεί να τις μοιράσει εκεί; Γιατί τόσο λίγους;

\ex \textbf{Οι ομάδες.} Ο γυμναστής θέλει να χωρίσει 72 παιδιά σε ίσες ομάδες. Θέλει τουλάχιστον 2 ομάδες και τουλάχιστον 2 παιδιά σε κάθε ομάδα. Με πόσους διαφορετικούς τρόπους γίνεται; Αν τα παιδιά ήταν 73, πόσοι τρόποι θα υπήρχαν;

\ex \textbf{Η φανέλα του Ηλία.} Ο αριθμός στη φανέλα του Ηλία είναι πρώτος αριθμός ανάμεσα στο 30 και στο 50, και το άθροισμα των ψηφίων του είναι 7. Ποιος είναι ο αριθμός;

\ex \textbf{Το κόσκινο.} Φτιάξε το κόσκινο του Ερατοσθένη για τους αριθμούς από το 1 ως το 60. Πόσοι πρώτοι υπάρχουν; Πόσους αριθμούς σβήνει το 2, και πόσους ακόμα το 3;

\zoni{Το σκέφτομαι}{μια ερώτηση που πάει πιο μακριά}

\ex Βρες τους διαιρέτες του 6 και του 28, εκτός από τον ίδιο τον αριθμό, και πρόσθεσέ τους. Τι παρατηρείς; Οι μαθηματικοί λένε αυτούς τους αριθμούς «τέλειους».

\ex Δύο πρώτοι που διαφέρουν κατά 2 λέγονται δίδυμοι, όπως οι 3 και 5. Βρες όλα τα ζευγάρια διδύμων πρώτων μέχρι το 50.

\ex Ο 72 είναι $2^3\cdot 3^2$. Γράψε όλους τους αριθμούς της μορφής $2^{\alpha}\cdot 3^{\beta}$, όπου $\alpha=0,1,2,3$ και $\beta=0,1,2$. Τι είναι αυτοί οι αριθμοί για τον 72; Πόσοι είναι;

\ex Η Ζωή ρωτάει «γιατί το 1 δεν είναι πρώτος, αφού διαιρείται μόνο με τον εαυτό του και με το 1;» Εξήγησε με την ανάλυση του 6 γιατί, αν το 1 ήταν πρώτος, δεν θα είχε κάθε αριθμός μία μοναδική ανάλυση.

\ex Η Πολυξένη ισχυρίζεται ότι $10.403=101\cdot 103$. Έλεγξέ το με πολλαπλασιασμό. Πόσο δύσκολο θα ήταν να βρεις τους δύο παράγοντες χωρίς να τους ξέρεις, δοκιμάζοντας τους πρώτους με τη σειρά;

\ypotmima{Γ}{ΜΚΔ και ΕΚΠ: μοιράζω και συναντιέμαι}

\begin{anagke}
Αυτή την εβδομάδα η κυρία Αλεξάνδρα βάζει δύο προβλήματα που μοιάζουν να έχουν αντίθετη ιδέα.

Πρώτο: το σχολείο φτιάχνει κήπο. Η Α’ τάξη έχει 72 φυτά, 48 γλάστρες και 120 σακούλες χώμα. Θέλουν να φτιάξουν ίδια κιτ για τα τμήματα, ώστε κάθε κιτ να έχει και από τα τρία είδη και να μη μείνει τίποτα. «Πόσα το πολύ ίδια κιτ γίνονται;» Ο Στέλιος λέει αμέσως «δώδεκα». Η Φανή ελέγχει: δώδεκα κιτ θέλουν 6 φυτά, 4 γλάστρες και 10 σακούλες το καθένα. Βγαίνει. «Μα δεν είπαμε πόσα το πολύ;» λέει η Ζωή.

Δεύτερο: η θεία Μαρία και ο Ηλίας κάθονται το βράδυ στην ακτή και βλέπουν δύο φάρους. Ο ένας αναβοσβήνει κάθε 12 δευτερόλεπτα, ο άλλος κάθε 18. Μόλις άναψαν και οι δύο μαζί. «Σε πόσο ξανά θα ανάψουν μαζί;» ρωτάει η θεία.

Στο πρώτο θέλουμε το μεγαλύτερο κοινό κομμάτι: μοιράζουμε. Στο δεύτερο θέλουμε την πρώτη φορά που συμπίπτουν: συναντιόμαστε. Και τα δύο λύνονται με τους ίδιους αριθμούς, διαιρέτες και πολλαπλάσια.
\end{anagke}

\begin{lexi}{ΜΚΔ και ΕΚΠ}{μέγιστος, ελάχιστος, κοινός}
Οι δύο συντομογραφίες κρύβουν τρεις λέξεις η καθεμία. \lexiword{Κοινός} είναι αυτός που ανήκει και στους δύο. \lexiword{Μέγιστος} είναι ο μεγαλύτερος που υπάρχει (το «μέγας» στον υπερθετικό), και \lexiword{ελάχιστος} ο μικρότερος που υπάρχει (το «ελαχύς»).

Το \lexiword{ΜΚΔ} είναι ο Μέγιστος Κοινός Διαιρέτης: ο μεγαλύτερος αριθμός που διαιρεί και τους δύο. Το \lexiword{ΕΚΠ} είναι το Ελάχιστο Κοινό Πολλαπλάσιο: το μικρότερο (μη μηδενικό) πολλαπλάσιο που είναι και των δύο.
\end{lexi}

\begin{metafrasi}
Ο μεγαλύτερος αριθμός που διαιρεί και το 12 και το 18. & $\MKD(12,\,18)=6$\\
Ο μικρότερος αριθμός (όχι 0) που είναι πολλαπλάσιο και του 12 και του 18. & $\EKP(12,\,18)=36$\\
Οι $\alpha$ και $\beta$ δεν έχουν κοινό διαιρέτη εκτός από το 1. & $\MKD(\alpha,\beta)=1$ (πρώτοι μεταξύ τους)\\
Ο $\alpha$ διαιρεί τον $\beta$. & $\MKD(\alpha,\beta)=\alpha$, $\EKP(\alpha,\beta)=\beta$\\
Το γινόμενο των δύο αριθμών. & $\MKD\cdot\EKP=\alpha\cdot\beta$\\
\end{metafrasi}

\tmima{Η θεωρία}

\begin{orismos}{ΜΚΔ}
Ο \textbf{Μέγιστος Κοινός Διαιρέτης} δύο ή περισσότερων φυσικών είναι ο μεγαλύτερος από τους κοινούς τους διαιρέτες. Γράφουμε $\MKD(\alpha,\beta)$.
\end{orismos}

\begin{orismos}{ΕΚΠ}
Το \textbf{Ελάχιστο Κοινό Πολλαπλάσιο} δύο ή περισσότερων φυσικών είναι το μικρότερο από τα κοινά τους πολλαπλάσια, εκτός από το 0 (που είναι πολλαπλάσιο όλων). Γράφουμε $\EKP(\alpha,\beta)$.
\end{orismos}

\tmimasmall{Δύο μέθοδοι}

\textbf{Με λίστες.} Γράφω τους διαιρέτες (για ΜΚΔ) ή τα πρώτα πολλαπλάσια (για ΕΚΠ) κάθε αριθμού και διαλέγω τον κοινό που ψάχνω. Είναι εύκολο για μικρούς αριθμούς.

\begin{kanonas}{ΜΚΔ και ΕΚΠ με ανάλυση}
\begin{enumerate}[leftmargin=1.6em,itemsep=0.2ex,topsep=2pt]
\item Αναλύω κάθε αριθμό σε γινόμενο πρώτων παραγόντων (με δυνάμεις).
\item \textbf{ΜΚΔ}: παίρνω μόνο τους \emph{κοινούς} πρώτους παράγοντες, καθένα με τον \emph{μικρότερο} εκθέτη που έχει.
\item \textbf{ΕΚΠ}: παίρνω \emph{όλους} τους πρώτους παράγοντες (κοινούς και μη), καθένα με τον \emph{μεγαλύτερο} εκθέτη που έχει.
\item Πολλαπλασιάζω τους παράγοντες που πήρα.
\end{enumerate}
\end{kanonas}

Η ιδέα: το ΜΚΔ πρέπει να χωράει και στους δύο αριθμούς, γι' αυτό παίρνει όσα έχουν και οι δύο, το λιγότερο. Το ΕΚΠ πρέπει να χωράει και τους δύο, γι' αυτό παίρνει όλα όσα χρειάζεται ο καθένας, το περισσότερο.

\tmimasmall{Πρώτοι μεταξύ τους και μια σχέση}

Δύο αριθμοί λέγονται \textbf{πρώτοι μεταξύ τους} αν ο μοναδικός κοινός διαιρέτης τους είναι το 1, δηλαδή αν $\MKD(\alpha,\beta)=1$. Δεν χρειάζεται να είναι οι ίδιοι πρώτοι: οι 8 και 15 είναι σύνθετοι αλλά πρώτοι μεταξύ τους. Τότε το ΕΚΠ τους είναι το γινόμενό τους: $\EKP(8,\,15)=120$.

Για δύο φυσικούς ισχύει πάντα
$$\MKD(\alpha,\beta)\cdot\EKP(\alpha,\beta)=\alpha\cdot\beta.$$
Για τα 12 και 18: $6\cdot 36=216=12\cdot 18$.

\begin{prosoxi}
\begin{itemize}[leftmargin=1.2em,itemsep=0.3ex,topsep=2pt]
\item Το ΜΚΔ είναι πάντα μικρότερο ή ίσο από τον μικρότερο αριθμό. Το ΕΚΠ είναι πάντα μεγαλύτερο ή ίσο από τον μεγαλύτερο. Αν βγάλεις ΜΚΔ μεγαλύτερο από έναν από τους αριθμούς, κάτι πήγε στραβά.
\item Στο ΜΚΔ προσοχή: παίρνω μόνο τους \emph{κοινούς} παράγοντες. Στο ΕΚΠ παίρνω \emph{όλους}.
\item Όταν ένας αριθμός διαιρεί τον άλλον, τα πράγματα είναι εύκολα: $\MKD(7,42)=7$ και $\EKP(7,42)=42$.
\end{itemize}
\end{prosoxi}

\tmima{Λυμένα παραδείγματα}

\paradeigma{Οι φάροι}
\tixero Ο ένας φάρος αναβοσβήνει κάθε 12 δευτερόλεπτα, ο άλλος κάθε 18. Άναψαν μαζί τώρα.\par
\tithelo Σε πόσα δευτερόλεπτα θα ανάψουν ξανά μαζί.\par
\possindeo Ο πρώτος ανάβει στα δευτερόλεπτα $12, 24, 36, 48, \dots$ και ο δεύτερος στα $18, 36, 54, \dots$ Η πρώτη φορά που συμπίπτουν είναι το 36. Δηλαδή ψάχνω το $\EKP(12,\,18)$.\par
Με ανάλυση: $12=2^2\cdot 3$ και $18=2\cdot 3^2$. Παίρνω όλους τους παράγοντες με τον μεγαλύτερο εκθέτη: $2^2\cdot 3^2=36$.\par
\apantisi Σε 36 δευτερόλεπτα.\par
\elegxos Το 36 είναι πολλαπλάσιο και του 12 ($12\cdot 3$) και του 18 ($18\cdot 2$), και κανένας μικρότερος αριθμός δεν είναι (δες τις λίστες).

\paradeigma{Τα κιτ του κήπου}
\tixero 72 φυτά, 48 γλάστρες και 120 σακούλες χώμα. Τα κιτ πρέπει να είναι ίδια και να μη μείνει τίποτα.\par
\tithelo Πόσα το πολύ κιτ γίνονται και τι έχει το καθένα.\par
\possindeo Ο αριθμός των κιτ πρέπει να διαιρεί και το 72 και το 48 και το 120. Ψάχνω το μεγαλύτερο κοινό διαιρέτη. $72=2^3\cdot 3^2$, $48=2^4\cdot 3$, $120=2^3\cdot 3\cdot 5$. Κοινοί παράγοντες: το 2 (μικρότερος εκθέτης 3) και το 3 (μικρότερος εκθέτης 1). Το 5 δεν είναι κοινό. Άρα $\MKD=2^3\cdot 3=24$.\par
Κάθε κιτ έχει: $72:24=3$ φυτά, $48:24=2$ γλάστρες, $120:24=5$ σακούλες.\par
\apantisi 24 κιτ, με 3 φυτά, 2 γλάστρες και 5 σακούλες το καθένα.\par
\elegxos Τα 12 κιτ του Στέλιου βγαίνουν, αλλά το 24 είναι κοινός διαιρέτης και μεγαλύτερος από το 12. Αν δοκιμάσουμε μεγαλύτερο αριθμό, π.χ. 48 κιτ, το 72 δεν διαιρείται με το 48 ($72:48$ δεν είναι ακέραιο).

\paradeigma{ΜΚΔ και ΕΚΠ τριών αριθμών}
\tixero Οι αριθμοί 1.800, 2.100 και 3.000.\par
\tithelo Το ΜΚΔ και το ΕΚΠ τους.\par
\possindeo Αναλύω: $1.800=2^3\cdot 3^2\cdot 5^2$, $2.100=2^2\cdot 3\cdot 5^2\cdot 7$ και $3.000=2^3\cdot 3\cdot 5^3$.\par
ΜΚΔ: οι παράγοντες 2, 3 και 5 υπάρχουν και στους τρεις (το 7 όχι). Παίρνω τον μικρότερο εκθέτη του καθενός: $2^2\cdot 3\cdot 5^2=300$.\par
ΕΚΠ: παίρνω όλους τους παράγοντες με τον μεγαλύτερο εκθέτη: $2^3\cdot 3^2\cdot 5^3\cdot 7=8\cdot 9\cdot 125\cdot 7=63.000$.\par
\apantisi $\MKD=300$ και $\EKP=63.000$.\par
\elegxos Το 300 διαιρεί τον καθένα: $1.800:300=6$, $2.100:300=7$, $3.000:300=10$. Και το 63.000 είναι πολλαπλάσιο του καθενός: $63.000:1.800=35$, $63.000:2.100=30$, $63.000:3.000=21$.

\tmima{Ασκήσεις}

\zoni{Το κατάλαβα}{εφαρμογή του κανόνα}

\ex Βρες το ΜΚΔ γράφοντας τους διαιρέτες: α) 14 και 35 \quad β) 24 και 36 \quad γ) 16 και 45 \quad δ) 28 και 42.

\ex Βρες το ΕΚΠ γράφοντας τα πρώτα πολλαπλάσια: α) 6 και 8 \quad β) 9 και 12 \quad γ) 10 και 15 \quad δ) 4 και 9.

\ex Βρες το ΜΚΔ και το ΕΚΠ με ανάλυση σε πρώτους παράγοντες: α) 30 και 45 \quad β) 72 και 108 \quad γ) 84 και 126 \quad δ) 150 και 225.

\ex Βρες το ΜΚΔ και το ΕΚΠ των τριών αριθμών 20, 30 και 50.

\ex Ποια από τα ζευγάρια είναι πρώτοι μεταξύ τους; α) 9 και 16 \quad β) 15 και 28 \quad γ) 18 και 27 \quad δ) 21 και 22 \quad ε) 35 και 49 \quad στ) 1 και 50.

\ex Βρες το ΜΚΔ και το ΕΚΠ των: α) 6 και 48 \quad β) 13 και 52 \quad γ) 20 και 100. Τι παρατηρείς όταν ο ένας αριθμός διαιρεί τον άλλον;

\ex Σωστό ή λάθος; α) $\MKD(8,\,12)=24$. \quad β) $\EKP(6,\,9)=54$. \quad γ) Το ΜΚΔ δύο διαφορετικών πρώτων είναι 1. \quad δ) Το ΕΚΠ δύο αριθμών δεν είναι ποτέ μικρότερο από τον μεγαλύτερο. \quad ε) Το ΜΚΔ δεν είναι ποτέ μεγαλύτερο από τον μικρότερο. \quad στ) Αν δύο αριθμοί είναι πρώτοι μεταξύ τους, το ΕΚΠ τους είναι το γινόμενό τους.

\zoni{Το χρησιμοποιώ}{προβλήματα από τη ζωή}

\ex \textbf{Οι φάροι.} Δύο άλλοι φάροι αναβοσβήνουν ο ένας κάθε 10 δευτερόλεπτα και ο άλλος κάθε 25. Άναψαν μαζί τώρα. Πότε θα ξανανάψουν μαζί; Πόσες φορές θα συμπέσουν στα επόμενα 3 λεπτά;

\ex \textbf{Τα καλάθια της γιαγιάς.} Η γιαγιά Μυρτώ έχει 90 μήλα και 150 πορτοκάλια και θέλει να φτιάξει ίδια καλάθια, με και από τα δύο φρούτα και χωρίς να περισσέψει τίποτα. Πόσα το πολύ καλάθια μπορεί να φτιάξει; Τι έχει το καθένα;

\ex \textbf{Τα πλακάκια της θείας.} Η αυλή της θείας Μαρίας είναι ορθογώνια, 240 εκ. επί 180 εκ. Θέλει να την καλύψει με τετράγωνα πλακάκια ίδιου μεγέθους, χωρίς να κόψει κανένα. Ποια είναι η μεγαλύτερη δυνατή πλευρά του πλακιδίου; Πόσα πλακάκια θα χρειαστεί;

\ex \textbf{Τα λεωφορεία.} Δύο γραμμές λεωφορείων ξεκινούν μαζί από τη στάση στις 7:00. Η μία περνάει κάθε 12 λεπτά και η άλλη κάθε 16. Πότε θα ξαναπεράσουν μαζί; Πόσες φορές θα βρεθούν μαζί από τις 7:00 ως τις 10:00 (μετρώντας και την πρώτη φορά στις 7:00);

\ex \textbf{Το πότισμα.} Η Ρήνη ποτίζει τις γλάστρες κάθε 4 μέρες, ο Ηλίας κάθε 6 και η Ζωή κάθε 9. Σήμερα τις πότισαν και οι τρεις. Σε πόσες μέρες θα τις ποτίσουν ξανά και οι τρεις την ίδια μέρα;

\zoni{Το σκέφτομαι}{μια ερώτηση που πάει πιο μακριά}

\ex Για τα ζευγάρια 10 και 25, 16 και 24, 21 και 35, 18 και 30 βρες το ΜΚΔ, το ΕΚΠ και το γινόμενο των δύο αριθμών. Σύγκρινε το γινόμενο ΜΚΔ·ΕΚΠ με το γινόμενο των αριθμών. Τι ανακαλύπτεις;

\ex Η Ζωή παρατηρεί ότι $90=36\cdot 2+18$ και ότι $\MKD(90,\,36)=18=\MKD(36,\,18)$. Δοκίμασε το ίδιο με τα 180 και 84: γράψε την Ευκλείδεια διαίρεση και σύγκρινε το $\MKD(180,\,84)$ με το $\MKD(84,\,\upsilon)$. Αν συνεχίσεις μέχρι να βρεις υπόλοιπο 0, τι δείχνει ο τελευταίος διαιρέτης; (Αυτή είναι η μέθοδος του Ευκλείδη.)

\ex Γιατί, όταν δύο αριθμοί είναι πρώτοι μεταξύ τους, το ΕΚΠ είναι το γινόμενό τους; Εξήγησε με τα 9 και 20, χρησιμοποιώντας τις αναλύσεις τους.

\ex Βρες δύο αριθμούς που να έχουν ΜΚΔ 4 και ΕΚΠ 48. Υπάρχει και δεύτερο ζευγάρι; (Βοήθεια: το γινόμενο των δύο πρέπει να είναι $4\cdot 48$.)

\ex Μπορεί το ΜΚΔ δύο \emph{διαφορετικών} αριθμών να είναι ίσο με το ΕΚΠ τους; Εξήγησε.

\zoni{Για το σπίτι}{και λίγα ακόμη}

\ex \textbf{Ο φούρνος του χωριού.} Η γιαγιά Μυρτώ ψήνει ψωμί κάθε 5 ημέρες, ο φούρνος του χωριού καθαρίζεται κάθε 6 και το φορτηγό φέρνει αλεύρι κάθε 8. Σήμερα συνέπεσαν και τα τρία. Σε πόσες ημέρες θα συμπέσουν ξανά;

\ex \textbf{Τα κουλουράκια της Μυρτώς.} Όταν η Ρήνη βάζει τα κουλουράκια της σε σακουλάκια των 4, των 6 ή των 9, δεν περισσεύει κανένα. Πόσα κουλουράκια έχει το λιγότερο; Αν ξέρουμε ότι έχει ανάμεσα στα 150 και στα 250, πόσα μπορεί να είναι;

\ex \textbf{Μουσική και προπόνηση.} Η Αγγελική έχει μάθημα μουσικής κάθε 6 ημέρες και ο Νίκος προπόνηση κάθε 8. Μια μέρα συνέπεσαν. Σε πόσες ημέρες θα συμπέσουν ξανά; Πόσες φορές πήγε ο καθένας ξεχωριστά στο διάστημα ανάμεσα στις δύο αυτές ημέρες, χωρίς να μετρήσεις τις ημέρες που συμπίπτουν;

\ex Δύο αριθμοί έχουν ΜΚΔ το 18. Εξήγησε γιατί έχουν και άλλους κοινούς διαιρέτες εκτός από τη μονάδα.

\ex Το τριπλάσιο ενός πρώτου αριθμού είναι πρώτος ή σύνθετος; Εξήγησε.

\ex Βρες όλους τους διαιρέτες των αριθμών: α) 40 \quad β) 56 \quad γ) 66 \quad δ) 91 \quad ε) 150 \quad στ) 221.

\ex Ανάλυσε σε γινόμενο πρώτων παραγόντων τους: α) 66 \quad β) 275 \quad γ) 840 \quad δ) 1.001.

\ex Αν ένας αριθμός διαιρεί δύο αριθμούς, διαιρεί και το άθροισμα και τη διαφορά τους. Δοκίμασέ το με το 6 και τους 42 και 18. Μετά εξήγησε γιατί ο αριθμός $5\alpha+15$ διαιρείται με το 5 για κάθε φυσικό $\alpha$.

\ex Η Πολυξένη παρατηρεί ότι $37\cdot 101=3.737$. Δοκίμασε τα $24\cdot 101$ και $58\cdot 101$. Τι παρατηρείς στο αποτέλεσμα; Εξήγησε με την επιμεριστική ιδιότητα γιατί γίνεται αυτό.

\epanalipsi{1}

\tmima{Λεξικό του κεφαλαίου}

\renewcommand{\arraystretch}{1.5}
\begin{tabular*}{\linewidth}{@{\extracolsep{\fill}} >{\raggedright\arraybackslash}p{0.26\linewidth} >{\raggedright\arraybackslash}p{0.64\linewidth}}
\toprule
\textbf{Λέξη ή τύπος} & \textbf{Με μία πρόταση}\\
\midrule
φυσικός & Ένας από τους $0,\,1,2,\,3,\dots$: μετράει πράγματα που δεν σπάνε.\\
στρογγυλοποίηση & Αλλάζω έναν αριθμό με κοντινό του, με συμφωνημένα βήματα, μία φορά.\\
$\alpha^{\nu}$ & Ο $\alpha$ γραμμένος $\nu$ φορές σε γινόμενο.\\
προτεραιότητα & Παρενθέσεις, δυνάμεις, $\cdot$ και $:$, και τέλος $+$ και $-$.\\
$\Delta=\delta\cdot\pi+\upsilon$ & Ευκλείδεια διαίρεση, με $\upsilon<\delta$.\\
κριτήρια & Κανόνες στα ψηφία που δείχνουν αν ένας αριθμός διαιρείται με το 2, 3, 4, 5, 9, 10, 25.\\
πρώτος & Ο αριθμός μεγαλύτερος από το 1 με διαιρέτες μόνο τον 1 και τον εαυτό του.\\
ΜΚΔ & Ο μεγαλύτερος κοινός διαιρέτης: μοιράζω.\\
ΕΚΠ & Το μικρότερο (μη μηδενικό) κοινό πολλαπλάσιο: συναντιέμαι.\\
\bottomrule
\end{tabular*}

\tmima{Ασκήσεις επανάληψης}

\ex Γράψε με ψηφία τον αριθμό «τριάντα επτά εκατομμύρια, δύο χιλιάδες, πέντε». Μετά στρογγυλοποίησέ τον στα εκατομμύρια.

\ex Βάλε σε αύξουσα σειρά τους 5.012, 5.102, 5.021, 4.999 και 10.000.

\ex Στρογγυλοποίησε το 86.549 στις δεκάδες, στις εκατοντάδες και στις χιλιάδες.

\ex Υπολόγισε με τις ιδιότητες των πράξεων: α) $25\cdot 17\cdot 4$ \quad β) $99\cdot 8$ \quad γ) $36\cdot 15+36\cdot 85$ \quad δ) $47+68+53+32$.

\ex Βρες την τιμή: α) $3\cdot(8-2)^2-10$ \quad β) $2^5-5^2$ \quad γ) $(2+3)\cdot 4-2^3$ \quad δ) $100-4\cdot 5^2$.

\ex Κάνε την Ευκλείδεια διαίρεση $1.000:48$ και γράψε την ισότητα $\Delta=\delta\cdot\pi+\upsilon$. Ποιος αριθμός δίνει πηλίκο 7 και υπόλοιπο 6 όταν διαιρεθεί με το 15;

\ex Σήμερα είναι Πέμπτη. Τι μέρα θα είναι μετά από 100 ημέρες;

\ex Με ποιους από τους 2, 3, 4, 5, 9, 10 και 25 διαιρείται ο αριθμός 3.240;

\ex Ποιο ψηφίο πρέπει να μπει στη θέση του $\square$ ώστε ο $5\square 7$ να διαιρείται με το 9;

\ex Ποιοι από τους 57, 59, 63 και 67 είναι πρώτοι; Για τους σύνθετους, βρες την ανάλυσή τους.

\ex Ανάλυσε τον 360 σε γινόμενο πρώτων παραγόντων. Πόσους πρώτους διαφορετικούς παράγοντες έχει;

\ex Βρες το ΜΚΔ και το ΕΚΠ των αριθμών 36 και 60.

\ex \textbf{Οι κάρτες του Στέλιου.} Ο Στέλιος έχει 36 κάρτες με ζώα και 60 κάρτες με αυτοκίνητα. Θέλει να φτιάξει όσο το δυνατόν περισσότερες ίδιες στοίβες, με και από τα δύο είδη, χωρίς να περισσέψει καμία. Πόσες στοίβες θα γίνουν και τι έχει η κάθε μία;

\ex \textbf{Η πισίνα.} Ο Ηλίας πάει στην πισίνα κάθε 8 μέρες και η Ζωή κάθε 12. Σήμερα πήγαν μαζί. Πότε θα ξαναπάνε μαζί;

\ex Σωστό ή λάθος; α) Κάθε άρτιος αριθμός διαιρείται με το 4. \quad β) $\MKD(12,\,18)=6$. \quad γ) Το 91 είναι πρώτος. \quad δ) $2^4=8$. \quad ε) $17-(9-3)=17-9-3$. \quad στ) $0:5=0$.

\ex Ο Νίκος λέει ότι $6+4\cdot 5=50$ και η Ζωή ότι $6+4\cdot 5=26$. Ποιος έχει δίκιο και γιατί; Πού πρέπει να βάλεις παρένθεση ώστε να βγει 50;

\ex Πόσους διαιρέτες έχει το 100; Πόσοι από αυτούς είναι πρώτοι;

\ex \textbf{Ερωτήσεις αυτοαξιολόγησης (πράξεις και δυνάμεις).} Σωστό ή λάθος; α) $(50-20)-5=50-(20-5)$. \quad β) $5\cdot 5\cdot 5=5^3$. \quad γ) $3^4=12$. \quad δ) $\beta+\beta+\beta=3\beta$. \quad ε) $2^4+4=20$. \quad στ) $2\cdot 10^2+5\cdot 10+1=251$. \quad ζ) $30-18:6=2$. \quad η) $6\cdot 4-3+7=28$.

\ex \textbf{Ερωτήσεις αυτοαξιολόγησης (διαιρετότητα).} Σωστό ή λάθος; α) Το άθροισμα δύο περιττών αριθμών είναι πάντα άρτιος. \quad β) Αν ο $\alpha$ είναι πολλαπλάσιο του $\beta$ (με $\beta\neq 0$), τότε ο $\beta$ διαιρεί τον $\alpha$. \quad γ) Το 51 είναι πολλαπλάσιο του 3 και του 5. \quad δ) Ο 630 διαιρείται με το 9. \quad ε) Η διαίρεση $365:20$ δίνει πηλίκο 18 και υπόλοιπο 5. \quad στ) Η σχέση $250=12\cdot 20+10$ είναι Ευκλείδεια διαίρεση. \quad ζ) Ο αριθμός $4\alpha+8$ διαιρείται με το 4. \quad η) Ο 312 διαιρείται με το 4 και με το 8.

\ex \textbf{Ερωτήσεις αυτοαξιολόγησης (ΜΚΔ, ΕΚΠ, πρώτοι).} Σωστό ή λάθος; α) Ο 14 και ο 42 έχουν ΜΚΔ το 7. \quad β) $\EKP(5,\,15)=75$. \quad γ) Η γραφή $150=5\cdot 30$ είναι ανάλυση σε γινόμενο πρώτων παραγόντων. \quad δ) Δύο διαδοχικοί φυσικοί (π.χ. 14 και 15) είναι πάντα πρώτοι μεταξύ τους. \quad ε) Αν ο $\alpha$ διαιρεί τον $\beta$, τότε $\MKD(\alpha,\,\beta)=\alpha$. \quad στ) Ο 2 είναι ο μόνος άρτιος πρώτος.

\tmima{Ο χάρτης της επανάληψης}

Μετά τις λύσεις, δες σε ποιες ασκήσεις δυσκολεύτηκες. Ο πίνακας σου λέει ποιο κεφάλαιο να ξαναδιαβάσεις.

\renewcommand{\arraystretch}{1.5}
\begin{tabular*}{\linewidth}{@{\extracolsep{\fill}} >{\raggedright\arraybackslash}p{0.38\linewidth} >{\raggedright\arraybackslash}p{0.52\linewidth}}
\toprule
\textbf{Αν δυσκολεύτηκες στις} & \textbf{ξαναδιάβασε τις παραγράφους}\\
\midrule
Ε1.1, Ε1.2, Ε1.3 & 1.1 (φυσικοί, διάταξη, στρογγυλοποίηση)\\
Ε1.4 & 1.2 (ιδιότητες των πράξεων)\\
Ε1.5, Ε1.16, Ε1.18 & 1.3 (δυνάμεις, προτεραιότητα)\\
Ε1.6, Ε1.7, Ε1.19 & 1.4 (Ευκλείδεια διαίρεση)\\
Ε1.8, Ε1.9 & 1.5, μέρος Α (κριτήρια διαιρετότητας)\\
Ε1.10, Ε1.11, Ε1.17 & 1.5, μέρος Β (πρώτοι, ανάλυση)\\
Ε1.12, Ε1.13, Ε1.14 & 1.5, μέρος Γ (ΜΚΔ, ΕΚΠ)\\
Ε1.15, Ε1.18, Ε1.19, Ε1.20 & 1.1 έως 1.5\\
\bottomrule
\end{tabular*}


\kefalaio{Τα κλάσματα}{Το ένα έχει κομμάτια, και τα κομμάτια έχουν ονόματα}{%
\begin{tabular}{@{}l@{\hspace{1em}}p{0.82\linewidth}@{}}
2.1 & Η έννοια του κλάσματος\\[2pt]
2.2 & Ισοδύναμα κλάσματα\\[2pt]
2.3 & Σύγκριση κλασμάτων\\[2pt]
2.4 & Πρόσθεση και αφαίρεση κλασμάτων\\[2pt]
2.5 & Πολλαπλασιασμός κλασμάτων\\[2pt]
2.6 & Διαίρεση κλασμάτων\\[2pt]
 & Επανάληψη του κεφαλαίου\\
\end{tabular}}
\paragrafos{Θετικοί και αρνητικοί ακέραιοι: ο κόσμος κάτω από το μηδέν}{Α.2.1, «Θετικοί, αρνητικοί, ομόσημοι, ετερόσημοι ακέραιοι»}

\begin{anagke}
Η Α3 κάνει εκδρομή στο Μέτσοβο. Στο πούλμαν το ψηφιακό θερμόμετρο γράφει $-4$. «Τι είναι αυτό το μείον;» ρωτάει η Πολυξένη. «Τέσσερις βαθμοί κάτω από το μηδέν», απαντάει ο Νίκος. Ο Στέλιος ταράζεται: «Κάτω από το μηδέν; Το μηδέν δεν είναι το τίποτα; Πώς μπορεί να υπάρχει κάτι μικρότερο από το τίποτα;».

Στο ξενοδοχείο ο ανελκυστήρας έχει κουμπιά $-2$, $-1$, $0$, $1$, $2$, $3$. Ο Ηλίας πατάει το $-1$ και κατεβαίνει στο υπόγειο με το πάρκινγκ. Η Αγγελική βλέπει στον πίνακα του καιρού: Αθήνα $+7$, Ιωάννινα $-3$, Φλώρινα $-10$, Χανιά $+11$. Η Ρήνη θυμάται ότι στο ταμείο της παρέας κάποτε έλειπαν $8$ ευρώ, ένα χρέος. Και στο ντοκιμαντέρ που είδαν, ο καρχαρίας βρισκόταν $200$ μέτρα κάτω από την επιφάνεια ενώ το αεροπλάνο πετούσε $200$ μέτρα πάνω από αυτήν.

Η κυρία Αλεξάνδρα το συνοψίζει: «Οι φυσικοί αριθμοί που ξέρετε μετράνε πόσα. Όταν όμως χρειάζεται να πούμε «πάνω ή κάτω», «κέρδος ή ζημιά», «δεξιά ή αριστερά», οι φυσικοί δεν φτάνουν. Χρειαζόμαστε αριθμούς με κατεύθυνση». Κάθε «δεξιά» προϋποθέτει ένα «αριστερά» και κάθε «πάνω» ένα «κάτω»: ο κόσμος των αριθμών έχει δύο πλευρές. Σε αυτή την παράγραφο γνωρίζουμε τους θετικούς και τους αρνητικούς ακέραιους, τους βάζουμε πάνω στον άξονα και μαθαίνουμε να λέμε ποιοι έχουν το ίδιο πρόσημο και ποιοι όχι.
\end{anagke}

\begin{lexi}{θετικός, αρνητικός, ακέραιος}{θετικός, αρνητικός, ακέραιος}
\lexiword{Θετικός} είναι αυτός που «τίθεται», που βάζουμε ως βάση. \lexiword{Αρνητικός} είναι αυτός που «αρνείται», που πηγαίνει ανάποδα από τον θετικό. Το \lexiword{πρόσημο} είναι το «προ-σήμα», το σημάδι που μπαίνει πριν από τον αριθμό. Ο \lexiword{ακέραιος} είναι ο «ακέραιος», ο ολόκληρος: χωρίς κομμάτια. \lexiword{Ομόσημοι} είναι οι αριθμοί με το «ίδιο σήμα» και \lexiword{ετερόσημοι} αυτοί με «άλλο» (έτερο) σήμα. Η \lexiword{τετμημένη} είναι αυτή που «τέμνεται», που κόβουμε από την αρχή του άξονα ως το σημείο.
\end{lexi}

\begin{metafrasi}
Τέσσερις βαθμοί κάτω από το μηδέν. & $-4$\\
Αεροπλάνο $200$ μ. πάνω από τη θάλασσα. & $+200$\\
Καρχαρίας $200$ μ. κάτω από την επιφάνεια. & $-200$\\
Χρέος $8$ ευρώ. & $-8$\\
Κέρδος $15$ ευρώ. & $+15$\\
Δεύτερο υπόγειο. & $-2$\\
Δύο αριθμοί με το ίδιο πρόσημο. & ομόσημοι\\
Ο αριθμός που αντιστοιχεί σε σημείο του άξονα. & τετμημένη, $\mathrm{Α}(-3)$\\
\end{metafrasi}

\tmima{Η θεωρία}

\tmimasmall{Πού χρησιμοποιούμε αρνητικούς αριθμούς}

Οι αρνητικοί αριθμοί εκφράζουν το αντίθετο από κάτι που ήδη μετράμε με φυσικούς:
\begin{itemize}[leftmargin=1.2em,itemsep=0.3ex,topsep=2pt]
\item \textbf{Κέρδος και ζημιά.} Έσοδο, κέρδος, είσπραξη: θετικό ποσό. Έξοδο, ζημιά, χρέος: αρνητικό ποσό.
\item \textbf{Κατεύθυνση.} Δεξιά (θετική) και αριστερά (αρνητική). Πάνω (θετική) και κάτω (αρνητική). Αύξηση και μείωση.
\item \textbf{Στροφή.} Αντίθετα από τους δείκτες του ρολογιού (θετική στροφή) και κατά τη φορά των δεικτών (αρνητική στροφή).
\end{itemize}
Την ένδειξη $-2$ στο θερμόμετρο τη διαβάζουμε «πλην δύο» ή «δύο υπό το μηδέν».

\begin{orismos}{πρόσημο, θετικός, αρνητικός}
Τα σύμβολα «$+$» και «$-$», που γράφονται πριν από τους αριθμούς, λέγονται \textbf{πρόσημα}. Οι αριθμοί με πρόσημο $+$ (ή χωρίς πρόσημο) λέγονται \textbf{θετικοί}: $+5$, $+2$, $7$. Οι αριθμοί με πρόσημο $-$ λέγονται \textbf{αρνητικοί}: $-1$, $-2$, $-6$.
\end{orismos}

Το $+$ μπροστά από έναν θετικό συχνά παραλείπεται: γράφουμε $2$ αντί για $+2$. Το $-$ μπροστά από έναν αρνητικό \textbf{δεν παραλείπεται ποτέ}. Το $0$ \textbf{δεν είναι ούτε θετικός ούτε αρνητικός}.

\begin{orismos}{ομόσημοι, ετερόσημοι}
Δύο αριθμοί είναι \textbf{ομόσημοι} όταν έχουν το ίδιο πρόσημο και \textbf{ετερόσημοι} όταν έχουν διαφορετικό πρόσημο.
\end{orismos}

Για παράδειγμα οι $+2$ και $+5$ είναι ομόσημοι, οι $-1$ και $-3$ επίσης. Οι $+2$ και $-5$ είναι ετερόσημοι.

\begin{orismos}{ακέραιοι αριθμοί}
\textbf{Ακέραιοι} είναι οι φυσικοί αριθμοί μαζί με τους αντίστοιχους αρνητικούς: $$\dots,\,-4,\,-3,\,-2,\,-1,\,0,\,+1,\,+2,\,+3,\,+4,\dots$$
\end{orismos}

Οι φυσικοί $0,\,1,\,2,\,3,\dots$ είναι λοιπόν και ακέραιοι. Αν στους φυσικούς προσθέσουμε και τους αρνητικούς $-1,\,-2,\,-3,\dots$ παίρνουμε όλους τους ακέραιους. Ο $-3$ είναι ακέραιος αλλά \emph{όχι} φυσικός.

\tmimasmall{Μια ιστορική σημείωση}

Οι αρνητικοί αριθμοί εμφανίζονται για πρώτη φορά σε ένα κινέζικο μαθηματικό βιβλίο, τα «Μαθηματικά σε εννέα βιβλία», που τοποθετείται στην εποχή της δυναστείας των Χαν (206 π.Χ. -- 220 μ.Χ.). Στην ευρωπαϊκή παράδοση ο Διόφαντος (γύρω στο 250 μ.Χ.) τους χρησιμοποιεί στους ενδιάμεσους υπολογισμούς του, αλλά ως λύση ενός προβλήματος δέχεται πάντα θετικό αριθμό. Για πολλούς αιώνες οι μαθηματικοί τους θεωρούσαν «παράλογους» ή «ψεύτικους» αριθμούς. Σήμερα τους χρησιμοποιούμε στο θερμόμετρο, στον λογαριασμό της τράπεζας, στο υψόμετρο.

\tmimasmall{Οι ακέραιοι στην αριθμογραμμή: η τετμημένη}

Στον ημιάξονα $Ox$ τοποθετούσαμε τους φυσικούς δεξιά από την αρχή $\mathrm{Ο}$. Αν τον προεκτείνουμε προς τα αριστερά (στον αντικείμενο ημιάξονα $Ox'$) με την ίδια μονάδα, παίρνουμε τον \textbf{άξονα $x'x$}. Οι αρνητικοί παριστάνονται με σημεία συμμετρικά, ως προς το $\mathrm{Ο}$, των σημείων που παριστάνουν τους αντίστοιχους θετικούς. Έτσι οι θετικοί βρίσκονται δεξιά του $\mathrm{Ο}$ και οι αρνητικοί αριστερά του.

\begin{center}
\akgr[0.6]{-10}{10}{}
\end{center}

Η θέση ενός σημείου πάνω στον άξονα $x'x$ ορίζεται από έναν ακέραιο που λέγεται \textbf{τετμημένη} του σημείου. Γράφουμε $\mathrm{Α}(4)$, $\mathrm{Β}(-2)$, $\mathrm{Ο}(0)$.

\begin{center}
\akgr[1]{-5}{5}{-2/\mathrm{Β},0/\mathrm{Ο},4/\mathrm{Α}}
\end{center}

Το σημείο $\mathrm{Α}$ έχει τετμημένη $4$, το $\mathrm{Β}$ έχει $-2$ και το $\mathrm{Ο}$ έχει $0$. Όσο πηγαίνουμε προς τα αριστερά οι αριθμοί μικραίνουν.

Ο ίδιος άξονας μπορεί να σταθεί και κατακόρυφα, όπως στον ανελκυστήρα ή στο θερμόμετρο:

\begin{center}
\begin{tikzpicture}[y=0.55cm]
  \draw[-{Stealth},thick] (0,-2.7)--(0,3.9);
  \foreach \j in {-2,-1,0,1,2,3}{\draw (-0.1,\j)--(0.1,\j) node[right=3pt,font=\small]{$\j$};}
  \node[left=4pt,font=\small] at (0,3) {τρίτος όροφος};
  \node[left=4pt,font=\small] at (0,0) {ισόγειο};
  \node[left=4pt,font=\small] at (0,-2) {δεύτερο υπόγειο};
\end{tikzpicture}
\end{center}

\begin{prosoxi}
\begin{itemize}[leftmargin=1.2em,itemsep=0.3ex,topsep=2pt]
\item Το $-3$ δεν είναι «τρία ευρώ λιγότερο από το τίποτα». Είναι μια θέση, τρεις μονάδες αριστερά του $0$, και η μονάδα μπορεί να είναι βαθμός, ευρώ ή μέτρο.
\item Το $0$ δεν έχει πρόσημο. Δεν γράφουμε $+0$ ή $-0$.
\item Όλοι οι αρνητικοί είναι μικρότεροι από το $0$. Το $-100$ δεν είναι «μεγάλος αριθμός»: είναι πολύ αριστερά.
\item Μην μπερδεύεις το πρόσημο με την πράξη: στο $5-3$ το μείον είναι αφαίρεση, στο $-3$ είναι πρόσημο.
\item Οι αριθμοί $+5$ και $5$ είναι ο ίδιος αριθμός. Αντίθετα, οι $5$ και $-5$ είναι διαφορετικοί.
\end{itemize}
\end{prosoxi}

\tmima{Λυμένα παραδείγματα}

\paradeigma{Αριθμοί για μεγέθη}
\tixero Αθήνα $7$° πάνω από το $0$, Ιωάννινα $3$° κάτω, Φλώρινα $10$° κάτω. Ένα χρέος $12$ ευρώ. Μια αύξηση $200$ ευρώ και μια μείωση $50$ ευρώ.\par
\tithelo Να τα γράψω ως ακέραιους.\par
\possindeo Πάνω από το μηδέν, αύξηση, κέρδος είναι θετικό. Κάτω από το μηδέν, μείωση, χρέος είναι αρνητικό.\par
\apantisi Αθήνα $+7$, Ιωάννινα $-3$, Φλώρινα $-10$, χρέος $-12$, αύξηση $+200$, μείωση $-50$.\par
\elegxos Η Φλώρινα ($-10$) είναι πιο κρύα από τα Ιωάννινα ($-3$): θα το συγκρίνουμε στην επόμενη παράγραφο.

\paradeigma{Θετικοί, αρνητικοί, ομόσημοι, ετερόσημοι}
\tixero Οι αριθμοί $+4$, $-9$, $-1$, $+7$.\par
\tithelo Ποιοι είναι θετικοί, ποιοι αρνητικοί και ποια ζεύγη είναι ομόσημα ή ετερόσημα.\par
\possindeo Κοιτάζω μόνο τα πρόσημα. Θετικοί: $+4$ και $+7$. Αρνητικοί: $-9$ και $-1$. Ομόσημα ζεύγη είναι αυτά που έχουν το ίδιο πρόσημο, όπως $(+4,\,+7)$ και $(-9,\,-1)$. Τα υπόλοιπα, π.χ. $(+4,\,-9)$, είναι ετερόσημα.\par
\apantisi Ομόσημοι: $+4,\,+7$ και $-9,\,-1$. Ετερόσημοι: κάθε θετικός με κάθε αρνητικό.\par
\elegxos Δύο αριθμοί με το ίδιο πρόσημο βρίσκονται στην ίδια πλευρά του $0$ στον άξονα.

\paradeigma{Φυσικοί ή ακέραιοι;}
\tixero Οι αριθμοί $-4$, $0$, $15$, $-1$, $3$, $-100$.\par
\tithelo Ποιοι είναι φυσικοί και ποιοι ακέραιοι.\par
\possindeo Φυσικοί είναι οι $0,\,1,\,2,\dots$ Οι αρνητικοί δεν είναι φυσικοί. Όλοι οι αριθμοί της λίστας όμως είναι ακέραιοι.\par
\apantisi Φυσικοί: $0$, $15$, $3$. Ακέραιοι: και οι έξι.\par
\elegxos Κάθε φυσικός είναι ακέραιος. Το αντίστροφο δεν ισχύει: ο $-4$ είναι ακέραιος αλλά όχι φυσικός.

\paradeigma{Σημεία στον άξονα}
\tixero Ο άξονας του σχήματος και τα σημεία $\mathrm{Δ}$, $\mathrm{Ε}$, $\mathrm{Ζ}$.\par
\tithelo Να βρω τις τετμημένες τους και να τοποθετήσω τον $-6$.\par
\possindeo Μετράω μονάδες από το $0$, δεξιά με $+$ και αριστερά με $-$.\par
\begin{center}
\akgr[0.7]{-8}{8}{-5/\mathrm{Δ},1/\mathrm{Ε},6/\mathrm{Ζ}}
\end{center}
\apantisi $\mathrm{Δ}(-5)$, $\mathrm{Ε}(1)$, $\mathrm{Ζ}(6)$. Ο $-6$ βρίσκεται $6$ μονάδες αριστερά του $0$, δηλαδή μία μονάδα πιο αριστερά από το $\mathrm{Δ}$.\par
\elegxos Τα $\mathrm{Δ}$ και $\mathrm{Ε}$ είναι σε διαφορετικές πλευρές του $0$: η τετμημένη του $\mathrm{Δ}$ είναι αρνητική και του $\mathrm{Ε}$ θετική.

\tmima{Ασκήσεις}

\zoni{Το κατάλαβα}{πρόσημα, ακέραιοι, άξονας}

\ex Εξέφρασε με ακέραιο: α) $5$ βαθμοί κάτω από το μηδέν \quad β) $12$ μ. πάνω από τη θάλασσα \quad γ) χρέος $8$ € \quad δ) κέρδος $15$ € \quad ε) το δεύτερο υπόγειο \quad στ) $30$ μ. κάτω από την επιφάνεια της θάλασσας.

\ex Ποιοι από τους αριθμούς $-7$, $+3$, $0$, $-12$, $8$, $+1$ είναι θετικοί, ποιοι αρνητικοί και ποιος δεν είναι ούτε το ένα ούτε το άλλο;

\ex Πες αν είναι ομόσημοι ή ετερόσημοι: α) $-3$ και $-8$ \quad β) $+4$ και $-4$ \quad γ) $5$ και $+9$ \quad δ) $-2$ και $6$ \quad ε) $-1$ και $-100$ \quad στ) $+8$ και $-8$.

\ex Η Ρήνη λέει «Αντί για $+2$ μπορώ να γράψω $2$», και ο Στέλιος «Αντί για $-2$ μπορώ να γράψω $2$». Ποιος έχει δίκιο;

\ex Σχεδίασε άξονα με μονάδα $1$ cm και σημείωσε τους $-4$, $-1$, $0$, $2$, $5$.

\ex Στον άξονα βρες τις τετμημένες των $\mathrm{Α},\,\mathrm{Β},\,\Gamma,\,\mathrm{Δ}$.
\begin{center}
\akgr{-6}{6}{-4/\mathrm{Α},-1/\mathrm{Β},3/\Gamma,5/\mathrm{Δ}}
\end{center}

\ex Σωστό ή λάθος; α) Το $0$ είναι θετικός αριθμός. \quad β) Κάθε φυσικός αριθμός είναι ακέραιος. \quad γ) Κάθε ακέραιος είναι φυσικός. \quad δ) Οι $-3$ και $-5$ είναι ετερόσημοι. \quad ε) Ο $-2$ είναι φυσικός αριθμός. \quad στ) Στον άξονα οι αρνητικοί βρίσκονται αριστερά του $0$.

\zoni{Το χρησιμοποιώ}{θερμόμετρα, ασανσέρ, θάλασσες}

\ex \textbf{Ο καιρός.} Οι προβλέψεις: Αθήνα $7$° πάνω από το $0$, Θεσσαλονίκη $3$° κάτω, Ιωάννινα $5$° κάτω, Χανιά $11$° πάνω, Φλώρινα $10$° κάτω. Γράψε τις θερμοκρασίες ως ακέραιους και τοποθέτησέ τες σε έναν κατακόρυφο άξονα.

\ex \textbf{Ο ανελκυστήρας.} Οι όροφοι του ξενοδοχείου είναι $-2,\,-1,\,0,\,1,\,2,\,3$. Ο Βασίλης είναι στον $3$ και θέλει το δεύτερο υπόγειο. Πόσους ορόφους κατεβαίνει;

\ex \textbf{Βουνό και θάλασσα.} Το αεροπλάνο πετάει στα $+300$ μ., το ψάρι στα $-45$ μ. και το υποβρύχιο στα $-120$ μ. Ποιο από τα τρία είναι πιο κοντά στην επιφάνεια της θάλασσας;

\ex \textbf{Το ταμείο.} Η Ρήνη έβαλε $20$ €, ο Στέλιος πήρε $7$ €, η Φανή έβαλε $5$ € και ο Βασίλης πήρε $12$ €. Γράψε κάθε κίνηση ως ακέραιο.

\ex \textbf{Ο καιρός αλλάζει.} Το πρωί στο Μέτσοβο είναι $-3$° και μέχρι το μεσημέρι ανεβαίνει $5$ βαθμούς. Βρες τη νέα θερμοκρασία μετρώντας πάνω στον άξονα.

\zoni{Το σκέφτομαι}{είναι το μηδέν τίποτα;}

\ex Ο Στέλιος ισχυρίζεται ότι δεν υπάρχει αριθμός μικρότερος από το $0$. Φέρε αντιπαράδειγμα από τη ζωή.

\ex Γιατί το $0$ δεν είναι ούτε θετικός ούτε αρνητικός αριθμός;

\ex Τα $0$°C είναι η θερμοκρασία πήξης του νερού και τα $32$°F είναι η ίδια θερμοκρασία σε βαθμούς Φαρενάιτ. Τι δείχνει αυτό για το «μηδέν» στο θερμόμετρο;

\ex Πόσοι ακέραιοι υπάρχουν από το $-5$ ως το $5$, με τα άκρα; Πόσοι ανάμεσα στο $-5$ και στο $5$ χωρίς τα άκρα;

\zoni{Για το σπίτι}{ψάξε τους αρνητικούς γύρω σου}

\ex Βρες πέντε περιπτώσεις της καθημερινής ζωής που χρησιμοποιούνται αρνητικοί αριθμοί.

\ex Μέτρησε τη θερμοκρασία κάθε πρωί για μια εβδομάδα και παράστησέ τη σε άξονα.

\paragrafos{Ισοδύναμα κλάσματα: το ίδιο κομμάτι με διαφορετικό κόψιμο}{Α.2.2, «Ισοδύναμα κλάσματα»}

\begin{anagke}
Παρασκευή βράδυ στο σπίτι του Βασίλη, που γιορτάζει τη νίκη του στις εκλογές της τάξης με πίτσες. Οι πίτσες είναι τέσσερις, ίδιου μεγέθους, αλλά κάθε παιδί κόβει τη δική του όπως θέλει. Ο Νίκος τη χωρίζει σε 4 κομμάτια, η Αγγελική σε 6, ο Βασίλης σε 8 και η Ζωή σε 10. Και οι τέσσερις τρώνε κάποια κομμάτια και ανακοινώνουν: ο Νίκος 2, η Αγγελική 3, ο Βασίλης 4 και η Ζωή 5.

«Εγώ έφαγα δύο, εσύ τρία, εγώ τέσσερα, εσύ πέντε. Ο Βασίλης έφαγε τα πιο πολλά!» λέει θριαμβευτικά ο Στέλιος, που μετράει μόνο κομμάτια. Η Πολυξένη κοιτάζει τις τέσσερις πίτσες πάνω στο τραπέζι και λέει ήσυχα: «Έφαγαν όλοι το ίδιο. Όλοι έφαγαν τη μισή πίτσα».

«Μα ο ένας έφαγε δύο και ο άλλος πέντε!» διαμαρτύρεται ο Στέλιος. Η Μυρτώ σηκώνει το κομμάτι της Ζωής και το βάζει δίπλα στης Αγγελικής. «Δεν είναι ίδια κομμάτια. Τα κομμάτια της Ζωής είναι μικρότερα». Πώς γράφουμε με αριθμούς το ότι τα $\frac{2}{4}$, τα $\frac{3}{6}$, τα $\frac{4}{8}$ και τα $\frac{5}{10}$ είναι το ίδιο κομμάτι;
\end{anagke}

\begin{lexi}{ισοδύναμα, απλοποίηση, ανάγωγο}{ισοδύναμα, απλοποίηση, ανάγωγο}
Δύο κλάσματα είναι \lexiword{ισοδύναμα} όταν έχουν την ίδια «δύναμη», δηλαδή την ίδια αξία. Το $\frac{2}{4}$ και το $\frac{3}{6}$ δεν είναι ίδια σύμβολα, αλλά είναι ίσα, γιατί δείχνουν το ίδιο κομμάτι.

\lexiword{Απλοποιώ} ένα κλάσμα σημαίνει το κάνω πιο απλό, με μικρότερους αριθμούς, χωρίς να αλλάξει η αξία του. Το \lexiword{ανάγωγο} κλάσμα είναι αυτό που δεν μπορεί να γίνει πιο απλό: δεν μπορεί να «αναχθεί», δηλαδή να φερθεί πίσω σε μικρότερους όρους.

Και δύο λέξεις για τα ονόματα: \lexiword{ομώνυμα} είναι τα κλάσματα που έχουν το ίδιο όνομα (παρονομαστή), και \lexiword{ετερώνυμα} εκείνα που έχουν διαφορετικό.
\end{lexi}

\begin{metafrasi}
Η μισή πίτσα, κομμένη σε 4 ή σε 6 κομμάτια. & $\dfrac{2}{4}=\dfrac{3}{6}$\\
Τα χιαστί γινόμενα είναι ίσα. & $2\cdot 6=4\cdot 3$\\
Πολλαπλασιάζω και τους δύο όρους επί 3. & $\dfrac{2}{5}=\dfrac{2\cdot 3}{5\cdot 3}=\dfrac{6}{15}$\\
Διαιρώ και τους δύο όρους με το 14. & $\dfrac{42}{56}=\dfrac{42:14}{56:14}=\dfrac{3}{4}$\\
Δεν απλοποιείται άλλο (ΜΚΔ $=1$). & $\dfrac{7}{12}$ είναι ανάγωγο\\
\end{metafrasi}

\tmima{Η θεωρία}

\begin{orismos}{ισοδύναμα κλάσματα}
Δύο κλάσματα $\dfrac{\alpha}{\beta}$ και $\dfrac{\gamma}{\delta}$ λέγονται \textbf{ισοδύναμα} ή \textbf{ίσα} όταν εκφράζουν το ίδιο μέρος ενός μεγέθους. Γράφουμε $\dfrac{\alpha}{\beta}=\dfrac{\gamma}{\delta}$.
\end{orismos}

\begin{idiotita}{χιαστί γινόμενα}
Αν $\dfrac{\alpha}{\beta}=\dfrac{\gamma}{\delta}$, τότε τα «χιαστί γινόμενα» είναι ίσα: $\alpha\cdot\delta=\beta\cdot\gamma$. Και αντιστρόφως.
\end{idiotita}

Έτσι $\frac{3}{4}=\frac{9}{12}$, γιατί $3\cdot 12=36=4\cdot 9$. Και τα $\frac{3}{5}$ και $\frac{4}{7}$ δεν είναι ισοδύναμα, γιατί $3\cdot 7=21$ ενώ $5\cdot 4=20$.

\begin{kanonas}{ισοδύναμα κλάσματα}
\begin{itemize}[leftmargin=1.2em,itemsep=0.3ex,topsep=2pt]
\item Αν πολλαπλασιάσω \textbf{και τους δύο} όρους ενός κλάσματος με τον ίδιο φυσικό αριθμό (όχι 0), βγαίνει ισοδύναμο κλάσμα: $\dfrac{2}{3}=\dfrac{2\cdot 4}{3\cdot 4}=\dfrac{8}{12}$.
\item Αν διαιρέσω \textbf{και τους δύο} όρους με τον ίδιο φυσικό αριθμό (όχι 0, που διαιρεί και τους δύο), βγαίνει ισοδύναμο κλάσμα: $\dfrac{18}{30}=\dfrac{18:6}{30:6}=\dfrac{3}{5}$.
\end{itemize}
\end{kanonas}

Η δεύτερη διαδικασία λέγεται \textbf{απλοποίηση}. Αν διαιρέσουμε και τους δύο όρους με τον \textbf{ΜΚΔ} τους, το κλάσμα γίνεται ανάγωγο σε ένα βήμα. Ένα κλάσμα λέγεται \textbf{ανάγωγο} όταν δεν έχει κοινό διαιρέτη ο αριθμητής και ο παρονομαστής εκτός από το 1, δηλαδή όταν είναι πρώτοι μεταξύ τους. Το $\frac{7}{12}$ είναι ανάγωγο, γιατί $\MKD(7,\,12)=1$.

\tmimasmall{Ομώνυμα και ετερώνυμα}

Τα κλάσματα που έχουν τον ίδιο παρονομαστή λέγονται \textbf{ομώνυμα}, όπως τα $\frac{2}{7}$ και $\frac{5}{7}$. Αυτά που έχουν διαφορετικούς παρονομαστές λέγονται \textbf{ετερώνυμα}, όπως τα $\frac{2}{7}$ και $\frac{5}{3}$.

\begin{kanonas}{πώς κάνω ετερώνυμα κλάσματα ομώνυμα}
\begin{enumerate}[leftmargin=1.6em,itemsep=0.2ex,topsep=2pt]
\item Απλοποιώ όσα κλάσματα απλοποιούνται.
\item Βρίσκω το \textbf{ΕΚΠ} των παρονομαστών: θα είναι ο κοινός παρονομαστής.
\item Διαιρώ το ΕΚΠ με τον παρονομαστή κάθε κλάσματος.
\item Πολλαπλασιάζω και τους δύο όρους κάθε κλάσματος με τον αριθμό που βρήκα για εκείνο το κλάσμα.
\end{enumerate}
\end{kanonas}

\begin{prosoxi}
\begin{itemize}[leftmargin=1.2em,itemsep=0.3ex,topsep=2pt]
\item Πολλαπλασιάζουμε ή διαιρούμε \textbf{και τους δύο} όρους με τον \textbf{ίδιο} αριθμό. Αν προσθέσεις τον ίδιο αριθμό και στους δύο όρους, αλλάζει η αξία: το $\frac{1}{2}$ δεν είναι ίσο με το $\frac{1+3}{2+3}=\frac{4}{5}$.
\item Στην απλοποίηση ο αριθμός που διαιρεί πρέπει να διαιρεί \textbf{και τους δύο} όρους ακριβώς.
\item Αν πολλαπλασιάσεις με το 0, χάνεις το κλάσμα: δεν επιτρέπεται.
\end{itemize}
\end{prosoxi}

\tmima{Λυμένα παραδείγματα}

\paradeigma{Οι πίτσες του Βασίλη}
\tixero Ο Νίκος έφαγε $\frac{2}{4}$ της πίτσας του, η Αγγελική $\frac{3}{6}$ και ο Βασίλης $\frac{4}{8}$.\par
\tithelo Να εξετάσω αν έφαγαν το ίδιο.\par
\possindeo Υπολογίζω τα χιαστί γινόμενα. Για το $\frac{2}{4}$ και το $\frac{3}{6}$: $2\cdot 6=12$ και $4\cdot 3=12$. Ίσα, άρα ισοδύναμα. Για το $\frac{3}{6}$ και το $\frac{4}{8}$: $3\cdot 8=24$ και $6\cdot 4=24$. Ίσα, άρα ισοδύναμα.\par
\apantisi Έφαγαν το ίδιο: όλοι τη μισή πίτσα.\par
\elegxos Και τα τρία απλοποιούνται στο $\frac{1}{2}$: $\frac{2}{4}=\frac{2:2}{4:2}=\frac{1}{2}$, $\frac{3}{6}=\frac{3:3}{6:3}=\frac{1}{2}$, $\frac{4}{8}=\frac{4:4}{8:4}=\frac{1}{2}$.

\paradeigma{Απλοποίηση του $\frac{84}{126}$}
\tixero Το κλάσμα $\frac{84}{126}$.\par
\tithelo Να το κάνω ανάγωγο.\par
\possindeo Βρίσκω το ΜΚΔ των όρων. $84=2^2\cdot 3\cdot 7$ και $126=2\cdot 3^2\cdot 7$, άρα $\MKD(84,\,126)=2\cdot 3\cdot 7=42$. Διαιρώ και τους δύο όρους με το 42: $\frac{84:42}{126:42}=\frac{2}{3}$.\par
\apantisi $\dfrac{84}{126}=\dfrac{2}{3}$.\par
\elegxos Το $\frac{2}{3}$ είναι ανάγωγο, γιατί το 2 και το 3 είναι πρώτοι μεταξύ τους. Και τα χιαστί γινόμενα: $84\cdot 3=252=126\cdot 2$. Σωστά.

\paradeigma{Τρία κλάσματα, ένας παρονομαστής}
\tixero Τα κλάσματα $\frac{3}{4}$, $\frac{5}{6}$ και $\frac{7}{10}$.\par
\tithelo Να τα κάνω ομώνυμα.\par
\possindeo Κανένα δεν απλοποιείται, γιατί είναι όλα ανάγωγα. Βρίσκω το $\EKP(4,\,6,\,10)$: $4=2^2$, $6=2\cdot 3$, $10=2\cdot 5$, άρα $2^2\cdot 3\cdot 5=60$. Διαιρώ: $60:4=15$, $60:6=10$, $60:10=6$. Πολλαπλασιάζω τους όρους: $\frac{3}{4}=\frac{3\cdot 15}{4\cdot 15}=\frac{45}{60}$, $\frac{5}{6}=\frac{5\cdot 10}{6\cdot 10}=\frac{50}{60}$, $\frac{7}{10}=\frac{7\cdot 6}{10\cdot 6}=\frac{42}{60}$.\par
\apantisi $\dfrac{45}{60}$, $\dfrac{50}{60}$ και $\dfrac{42}{60}$.\par
\elegxos Απλοποιώντας το $\frac{45}{60}$ με το 15 ξαναβρίσκω το $\frac{3}{4}$. Και το $\frac{50}{60}=\frac{5}{6}$ (με το 10) και το $\frac{42}{60}=\frac{7}{10}$ (με το 6). Σωστά.

\tmima{Ασκήσεις}

\zoni{Το κατάλαβα}{εφαρμογή του κανόνα}

\ex Γράψε τρία κλάσματα ισοδύναμα με το $\frac{2}{3}$, πολλαπλασιάζοντας τους όρους του επί 2, επί 3 και επί 10.

\ex Βρες τον αριθμό που λείπει: α) $\frac{3}{4}=\frac{\square}{20}$ \quad β) $\frac{5}{7}=\frac{35}{\square}$ \quad γ) $\frac{18}{24}=\frac{3}{\square}$ \quad δ) $\frac{\square}{9}=\frac{20}{36}$.

\ex Έλεγξε με τα χιαστί γινόμενα ποια ζευγάρια είναι ισοδύναμα: α) $\frac{4}{6}$ και $\frac{10}{15}$ \quad β) $\frac{5}{8}$ και $\frac{15}{32}$ \quad γ) $\frac{7}{9}$ και $\frac{21}{27}$ \quad δ) $\frac{12}{18}$ και $\frac{14}{21}$ \quad ε) $\frac{3}{5}$ και $\frac{4}{7}$.

\ex Απλοποίησε μέχρι να γίνουν ανάγωγα: α) $\frac{18}{24}$ \quad β) $\frac{35}{42}$ \quad γ) $\frac{48}{72}$ \quad δ) $\frac{100}{250}$ \quad ε) $\frac{56}{98}$.

\ex Ποια από τα κλάσματα $\frac{9}{16}$, $\frac{12}{18}$, $\frac{15}{28}$, $\frac{21}{35}$, $\frac{17}{51}$ και $\frac{25}{36}$ είναι ανάγωγα; Για τα υπόλοιπα, γράψε το ανάγωγο.

\ex Μετάτρεψε σε ισοδύναμα κλάσματα με παρονομαστή 100: α) $\frac{3}{4}$ \quad β) $\frac{7}{20}$ \quad γ) $\frac{2}{5}$ \quad δ) $\frac{9}{50}$ \quad ε) $\frac{13}{25}$.

\ex Μετάτρεψε σε ισοδύναμα κλάσματα με παρονομαστή 12: α) $\frac{5}{6}$ \quad β) $\frac{3}{4}$ \quad γ) $\frac{2}{3}$ \quad δ) $\frac{1}{2}$.

\ex Κάνε ομώνυμα τα κλάσματα: α) $\frac{2}{3}$ και $\frac{3}{4}$ \quad β) $\frac{5}{6}$ και $\frac{7}{8}$ \quad γ) $\frac{4}{5}$, $\frac{1}{2}$ και $\frac{7}{10}$ \quad δ) $\frac{1}{6}$, $\frac{3}{8}$ και $\frac{5}{12}$.

\ex Σωστό ή λάθος; α) $\frac{2}{3}=\frac{4}{6}$ \quad β) $\frac{3}{5}=\frac{6}{12}$ \quad γ) Αν διαιρέσω τον αριθμητή με το 3 και τον παρονομαστή με το 4, βγαίνει ισοδύναμο κλάσμα. \quad δ) $\frac{7}{7}=1$ \quad ε) Το $\frac{6}{9}$ είναι ανάγωγο. \quad στ) Τα $\frac{4}{9}$ και $\frac{8}{18}$ είναι ισοδύναμα.

\zoni{Το χρησιμοποιώ}{προβλήματα από τη ζωή}

\ex \textbf{Οι μερίδες.} Τέσσερα παιδιά έφαγαν από μια ίδια πίτσα. Η Μυρτώ έφαγε $\frac{3}{4}$, η Φανή $\frac{6}{8}$, ο Ηλίας $\frac{9}{12}$ και ο Στέλιος $\frac{4}{6}$. Ποιος έφαγε περισσότερο και ποιος λιγότερο; Κάνε τα κλάσματα ομώνυμα για να βεβαιωθείς.

\ex \textbf{Η σοκολάτα.} Η σοκολάτα της Αγγελικής έχει 24 τετράγωνα και έφαγε τα 18. Τι κλάσμα της σοκολάτας έφαγε; Γράψε το ως ανάγωγο.

\ex \textbf{Τα λεπτά.} Η ώρα έχει 60 λεπτά. Γράψε ως ανάγωγο κλάσμα της ώρας τα 45, τα 20, τα 40 και τα 12 λεπτά.

\ex \textbf{Το τυρί της γιαγιάς.} Η γιαγιά Ρήνη αγοράζει τυρί 250, 750, 400 και 125 γραμμάρια σε διαφορετικές μέρες. Ποιο μέρος του κιλού (1.000 γραμμάρια) είναι κάθε ποσότητα; Γράψε το ανάγωγο κλάσμα.

\ex \textbf{Οι νότες της Αγγελικής.} Στο μάθημα της μουσικής η Αγγελική μαθαίνει ότι μια νότα «τέταρτο» διαρκεί το $\frac{1}{4}$ ενός μέτρου, μια «όγδοο» το $\frac{1}{8}$ και μια «δέκατο έκτο» το $\frac{1}{16}$. Πόσα όγδοα χωράνε σε ένα τέταρτο; Πόσα δέκατα έκτα; Γράψε τις ισότητες $\frac{1}{4}=\frac{\square}{8}=\frac{\square}{16}$.

\zoni{Το σκέφτομαι}{μια ερώτηση που πάει πιο μακριά}

\ex Ο Στέλιος έγραψε $\frac{1}{2}=\frac{1+3}{2+3}=\frac{4}{5}$. Δείξε με μια πίτσα ότι αυτό είναι λάθος. Ποιος είναι ο σωστός τρόπος να φτιάξεις ισοδύναμο κλάσμα με παρονομαστή 6;

\ex Πόσα διαφορετικά κλάσματα ισοδύναμα με το $\frac{3}{5}$ έχουν παρονομαστή μικρότερο από το 50; Γράψε τα.

\ex Βρες ένα κλάσμα ισοδύναμο με το $\frac{5}{8}$, του οποίου οι δύο όροι έχουν άθροισμα 91.

\ex Ένα κλάσμα είναι ισοδύναμο με το $\frac{3}{7}$ και οι όροι του διαφέρουν κατά 20. Ποιο είναι;

\ex Γιατί, όταν απλοποιούμε και τους δύο όρους με τον ΜΚΔ τους, το κλάσμα που βγαίνει είναι ανάγωγο; Σκέψου το $\frac{84}{126}$ και το $\frac{2}{3}$: αν τα 2 και 3 είχαν κοινό διαιρέτη μεγαλύτερο από το 1, τι θα συνέβαινε με το 42;

\ex Γράψε τρία κλάσματα που είναι ισοδύναμα με το 0 και τρία που είναι ισοδύναμα με το 1.

\zoni{Για το σπίτι}{με χαρτί και ψαλίδι}

\ex Δες τις τέσσερις ράβδους. Γράψε το κλάσμα του χρωματισμένου μέρους καθεμιάς και παρατήρησε τι συμβαίνει.\par\smallskip
\barfig{3}{2}\quad\barfig{6}{4}\quad\barfig{9}{6}\quad\barfig{12}{8}

\ex Πάρε ένα φύλλο χαρτί και δίπλωσέ το στη μέση, μετά πάλι στη μέση, και πάλι. Γράψε με κλάσματα το μέρος του φύλλου που έχει ένα στρώμα, δύο στρώματα και τέσσερα στρώματα όταν το ξεδιπλώσεις. Ποια ισοδύναμα κλάσματα βρήκες για το $\frac{1}{2}$;

\paragrafos{Πρόσθεση και αφαίρεση ακεραίων: κέρδη, ζημιές και διαφορές}{Α.2.3, «Πρόσθεση και αφαίρεση ακεραίων»}

\begin{anagke}
Στο πούλμαν η παρέα παίζει ένα παιχνίδι με κάρτες: κάθε κάρτα γράφει έναν πόντο, θετικό για κέρδος και αρνητικό για ζημιά. Ο Βασίλης τραβάει $+5$ και μετά $-3$. «Πόσους πόντους έχω;» ρωτάει. Ο Νίκος υπολογίζει: «Πέντε μείον τρία κάνει δύο». Η Πολυξένη τραβάει $-4$ και μετά $-2$ και ο Στέλιος λέει «έξι», μα η Ρήνη τον ρωτάει: «Έξι κέρδος ή έξι ζημιά;». Ο Ηλίας κάνει μια πιο δύσκολη ερώτηση: «Αν τραβήξω $-7$ και μετά $+3$, είμαι ακόμα πίσω;».

Η Αγγελική το βλέπει στον άξονα: «Κάθε κάρτα είναι μια μετακίνηση. Το $+5$ σημαίνει «πέντε βήματα δεξιά», το $-3$ «τρία βήματα αριστερά». Ξεκινάω από το μηδέν, πηγαίνω πέντε δεξιά και μετά τρία αριστερά και καταλήγω στο $+2$». Για το $-4$ και το $-2$ είναι έξι βήματα αριστερά: $-6$. Και για τα $-7$ και $+3$ καταλήγουμε στο $-4$: ακόμα πίσω, αλλά λιγότερο. Ο Νίκος παρατηρεί και κάτι ωραίο: αν μια κάρτα $+3$ και μια $-3$ βρεθούν μαζί, «ακυρώνονται».

Το βράδυ ο μετεωρολόγος στην τηλεόραση λέει: «Η μέγιστη θερμοκρασία σήμερα είναι $+5$ βαθμοί και η ελάχιστη $-3$. Η διαφορά είναι $8$ βαθμοί». Ο Στέλιος κάνει τον λογαριασμό: «Πέντε μείον τρία κάνει δύο. Ο μετεωρολόγος κάνει λάθος». Η Πολυξένη τον διορθώνει: «Δεν αφαιρούμε το τρία από το πέντε. Αφαιρούμε το μείον τρία!». Η Αγγελική μετράει στον άξονα: από το $-3$ ως το $+5$ είναι τρία βήματα ως το μηδέν και πέντε ακόμη, οκτώ. Ο Ηλίας το γενικεύει: «Μήπως η αφαίρεση του $-3$ είναι το ίδιο με την πρόσθεση του $+3$;».

Σε αυτή την παράγραφο η παρέα ανακαλύπτει τους κανόνες της πρόσθεσης και της αφαίρεσης ακεραίων, μαθαίνει γιατί ισχύουν και πώς γίνονται σύντομοι οι υπολογισμοί όταν έχουμε πολλούς όρους και παρενθέσεις.
\end{anagke}

\begin{lexi}{πρόσθεση, άθροισμα, αφαίρεση, διαφορά}{πρόσθεση, άθροισμα, αφαίρεση, διαφορά}
Η \lexiword{πρόσθεση} είναι το «προσ-τίθημι», δηλαδή «βάζω επιπλέον». Οι \lexiword{προσθετέοι} είναι αυτά που προσθέτουμε και το \lexiword{άθροισμα} είναι αυτό που βγαίνει όταν τα «αθροίζουμε», τα βάζουμε σε ένα σωρό. Η \lexiword{αφαίρεση} είναι το «αφ-αιρώ», «βγάζω από κάτι». Η \lexiword{διαφορά} είναι αυτό που μένει και δείχνει «πόσο διαφέρουν» δύο αριθμοί. Στους ακέραιους, το «προσθέτω» μπορεί να σημαίνει και «ζημιώνω».
\end{lexi}

\begin{metafrasi}
$5$ βήματα δεξιά και $3$ αριστερά. & $(+5)+(-3)=+2$\\
$4$ βήματα αριστερά και $2$ αριστερά ακόμα. & $(-4)+(-2)=-6$\\
Κέρδος και ζημιά που ακυρώνονται. & $(+3)+(-3)=0$\\
Δεν αλλάζει η σειρά των προσθετέων. & $\alpha+\beta=\beta+\alpha$\\
Αφαίρεση σημαίνει πρόσθεση του αντίθετου. & $\alpha-\beta=\alpha+(-\beta)$\\
Η διαφορά θερμοκρασιών $+5$ και $-3$. & $(+5)-(-3)=8$\\
\end{metafrasi}

\tmima{Η θεωρία}

\tmimasmall{Πρόσθεση ακεραίων}

Θεωρούμε ένα ρομπότ που ξεκινά από την αρχή $\mathrm{Ο}$ του άξονα και κινείται μόνο πάνω στην ευθεία. Μια εντολή $+4$ το στέλνει $4$ μονάδες δεξιά και μια εντολή $-3$ το στέλνει $3$ μονάδες αριστερά. Αν δώσουμε δύο εντολές, η θέση που καταλήγει είναι το άθροισμα των δύο αριθμών.

\begin{center}
\begin{tikzpicture}[x=0.65cm]
  \draw[-{Stealth},thick] (-7.5,0)--(8.5,0);
  \foreach \j in {-7,...,8}{\draw (\j,0.1)--(\j,-0.1) node[below,font=\scriptsize]{$\j$};}
  \draw[thick] (0,0.17)--(0,-0.17);
  \draw[-{Stealth},gblue,thick] (0,0.45) to[bend left=25] node[above,font=\small]{$+4$} (4,0.45);
  \draw[-{Stealth},gwarmD,thick] (4,0.45) to[bend right=25] node[above,font=\small]{$-3$} (1,0.45);
  \fill[gblue] (1,0) circle (2.3pt);
\end{tikzpicture}
\end{center}

Το ρομπότ πήγε $4$ δεξιά και μετά $3$ αριστερά και έφτασε στο $+1$: $(+4)+(-3)=+1$. Αν κινηθεί $3$ αριστερά και μετά $1$ αριστερά, φτάνει στο $-4$: $(-3)+(-1)=-4$. Από αυτά βγαίνουν οι κανόνες:

\begin{kanonas}{πρόσθεση ακεραίων}
\begin{itemize}[leftmargin=1.2em,itemsep=0.3ex,topsep=2pt]
\item \textbf{Ομόσημοι:} προσθέτω τις απόλυτες τιμές και κρατάω το κοινό πρόσημο. $(-8)+(-5)=-(8+5)=-13$ και $(+9)+(+4)=+13$.
\item \textbf{Ετερόσημοι:} αφαιρώ τη μικρότερη απόλυτη τιμή από τη μεγαλύτερη και βάζω το πρόσημο του προσθετέου με τη μεγαλύτερη απόλυτη τιμή. $(+9)+(-4)=+5$ και $(-9)+(+4)=-5$.
\item \textbf{Αντίθετοι:} το άθροισμα είναι $0$. $(-6)+(+6)=0$.
\item Με το $0$: $\alpha+0=\alpha$.
\end{itemize}
\end{kanonas}

\begin{idiotita}{ιδιότητες της πρόσθεσης}
Για όλους τους ακέραιους $\alpha,\,\beta,\,\gamma$ ισχύουν: \textbf{αντιμεταθετική} $\alpha+\beta=\beta+\alpha$, \textbf{προσεταιριστική} $(\alpha+\beta)+\gamma=\alpha+(\beta+\gamma)$, $\alpha+0=\alpha$ και $\alpha+(-\alpha)=0$.
\end{idiotita}

Χάρη στις ιδιότητες, όταν έχω πολλούς προσθετέους μπορώ να τους ομαδοποιήσω όπως θέλω: προσθέτω χωριστά όλους τους θετικούς και χωριστά όλους τους αρνητικούς και μετά τα δύο αποτελέσματα. Έτσι $(-7)+12+(-5)+9+(-3)=(12+9)+\big((-7)+(-5)+(-3)\big)=21+(-15)=6$.

\tmimasmall{Αφαίρεση ακεραίων}

\begin{kanonas}{αφαίρεση ακεραίων}
Για να αφαιρέσω από τον $\alpha$ τον $\beta$, προσθέτω στον $\alpha$ τον \textbf{αντίθετο} του $\beta$: $$\alpha-\beta=\alpha+(-\beta).$$ Έτσι $(+2)-(+3)=(+2)+(-3)=-1$ και $(+4)-(-1)=(+4)+(+1)=+5$.
\end{kanonas}

Στους ακέραιους η αφαίρεση γίνεται πάντα: ο μειωτέος δεν χρειάζεται να είναι μεγαλύτερος από τον αφαιρετέο. Η διαφορά $\alpha-\beta$ δείχνει και τον δρόμο από το $\beta$ προς το $\alpha$ πάνω στον άξονα: από το $-3$ ως το $+5$ είναι $8$ βήματα δεξιά, και πράγματι $(+5)-(-3)=8$.

\tmimasmall{Απαλοιφή παρενθέσεων και αλγεβρικό άθροισμα}

Μια παράσταση με προσθέσεις και αφαιρέσεις γράφεται ως πρόσθεση ακεραίων: $12-15+8-20+3=12+(-15)+8+(-20)+3$. Τέτοια παράσταση ονομάζεται \textbf{αλγεβρικό άθροισμα} και μπορώ να τη λύσω με όποια σειρά θέλω: $(12+8+3)+(-15-20)=23-35=-12$. Για να σβήσω παρενθέσεις:
\begin{itemize}[leftmargin=1.2em,itemsep=0.3ex,topsep=2pt]
\item Όταν η παρένθεση έχει μπροστά της $+$ (ή δεν έχει πρόσημο), τη σβήνω μαζί με αυτό το $+$ και γράφω τους όρους της \textbf{με τα ίδια πρόσημα}: $(+5)+(-4)+(-2+6)=5-4-2+6$.
\item Όταν η παρένθεση έχει μπροστά της $-$, τη σβήνω μαζί με το $-$ και γράφω τους όρους της \textbf{με αντίθετα πρόσημα}: $-(+3)-(-4)-(-2+5)=-3+4+2-5$.
\end{itemize}

\begin{prosoxi}
\begin{itemize}[leftmargin=1.2em,itemsep=0.3ex,topsep=2pt]
\item Στους ακέραιους η πρόσθεση δεν σημαίνει πάντα αύξηση: $5+(-3)=2<5$. Και η αφαίρεση δεν σημαίνει πάντα μείωση: $3-(-2)=5>3$.
\item Στους ετερόσημους δεν προσθέτω τις απόλυτες τιμές: $(+9)+(-4)\neq13$, αλλά $9-4=5$.
\item Το πρόσημο του αθροίσματος ετεροσήμων είναι του μεγαλύτερου κατ' απόλυτη τιμή: $(-9)+4=-5$, όχι $+5$.
\item Η διαφορά δεν είναι αντιμεταθετική: $5-3=2$ αλλά $3-5=-2$.
\item Με το $-$ μπροστά από παρένθεση αλλάζω το πρόσημο \emph{όλων} των όρων της παρένθεσης, όχι μόνο του πρώτου: $10-(3-8+2)=10-3+8-2$.
\item Η διαφορά θερμοκρασιών $+5$ και $-3$ είναι $8$, όχι $2$. Αφαιρώ τον αριθμό μαζί με το πρόσημό του.
\end{itemize}
\end{prosoxi}

\tmima{Λυμένα παραδείγματα}

\paradeigma{Πρόσθεση στον άξονα}
\tixero Οι προσθέσεις $(-2)+5$, $3+(-7)$, $(-4)+(-3)$.\par
\tithelo Να τις κάνω στον άξονα.\par
\possindeo Ξεκινάω από τον πρώτο προσθετέο και κινούμαι: από το $-2$ πέντε δεξιά: $3$. Από το $3$ επτά αριστερά: $-4$. Από το $-4$ τρία αριστερά: $-7$.\par
\apantisi $3$, $-4$, $-7$.\par
\elegxos Με τους κανόνες: $5-2=3$, $-(7-3)=-4$, $-(4+3)=-7$.

\paradeigma{Με τους κανόνες}
\tixero Οι προσθέσεις $(-8)+(-5)$, $(+9)+(-4)$, $(-9)+(+4)$, $(-6)+(+6)$.\par
\tithelo Να υπολογίσω τα αθροίσματα.\par
\possindeo Ομόσημοι: $-(8+5)=-13$. Ετερόσημοι με μεγαλύτερη απόλυτη τιμή το $+9$: $+(9-4)=5$. Ετερόσημοι με μεγαλύτερη το $-9$: $-(9-4)=-5$. Αντίθετοι: $0$.\par
\apantisi $-13$, $5$, $-5$, $0$.\par
\elegxos Κάθε πρόσημο το βρίσκω από τον προσθετέο με τη μεγαλύτερη απόλυτη τιμή.

\paradeigma{Αφαιρέσεις}
\tixero Οι διαφορές $(+5)-(-3)$, $(+2)-(+3)$, $(-4)-(-6)$, $(-7)-(+5)$.\par
\tithelo Να τις υπολογίσω.\par
\possindeo Τις γράφω ως προσθέσεις του αντίθετου: $5+3$, $2+(-3)$, $(-4)+6$, $(-7)+(-5)$.\par
\apantisi $8$, $-1$, $2$, $-12$.\par
\elegxos Για την πρώτη: $(-3)+8=5$. Για την τελευταία: $(-12)+5=-7$.

\paradeigma{Το ταμείο της παρέας}
\tixero Το ταμείο έχει $+30$ €. Κινήσεις: $-12$, $+8$, $-35$, $+20$.\par
\tithelo Το υπόλοιπο.\par
\possindeo Θετικοί: $30+8+20=58$. Αρνητικοί: $-12-35=-47$. $58+(-47)=11$.\par
\apantisi $11$ €.\par
\elegxos Βήμα βήμα: $30,\ 18,\ 26,\ -9,\ 11$. Συμφωνεί.

\paradeigma{Παρενθέσεις με δύο τρόπους}
\tixero Η παράσταση $-(6-9+2)+(-5+1)-(4-10-3)$.\par
\tithelo Να την υπολογίσω.\par
\possindeo \textbf{1ος τρόπος:} κάνω πρώτα τις πράξεις μέσα στις παρενθέσεις: $6-9+2=-1$, $-5+1=-4$, $4-10-3=-9$. Άρα $-(-1)+(-4)-(-9)=1-4+9=6$. \textbf{2ος τρόπος:} σβήνω τις παρενθέσεις: $-6+9-2-5+1-4+10+3$. Θετικοί: $9+1+10+3=23$. Αρνητικοί: $-6-2-5-4=-17$. Άρα $23-17=6$.\par
\apantisi $6$.\par
\elegxos Και οι δύο τρόποι δίνουν το ίδιο αποτέλεσμα.

\paradeigma{Παραστάσεις με γράμματα}
\tixero Οι αριθμοί $\alpha=-2$, $\beta=+4$, $\gamma=-6$.\par
\tithelo Να υπολογίσω τις παραστάσεις $\alpha+\beta+\gamma$, $\beta-\alpha-\gamma$ και $\alpha-(\beta-\gamma)$.\par
\possindeo Βάζω τον αριθμό στη θέση του γράμματος, μέσα σε παρένθεση για να φαίνεται το πρόσημο. $\alpha+\beta+\gamma=(-2)+4+(-6)=-4$. $\beta-\alpha-\gamma=4-(-2)-(-6)=4+2+6=12$. Για την τρίτη, πρώτα $\beta-\gamma=4-(-6)=10$ και μετά $\alpha-10=(-2)-10=-12$.\par
\apantisi $-4$, $12$ και $-12$.\par
\elegxos Για την πρώτη: $(-2)+(-6)=-8$ και $-8+4=-4$.

\tmima{Ασκήσεις}

\zoni{Το κατάλαβα}{κανόνες, ιδιότητες, παρενθέσεις}

\ex Στον άξονα βρες τα αθροίσματα: α) $(-3)+5$ \quad β) $2+(-6)$ \quad γ) $(-5)+(-2)$ \quad δ) $7+(-7)$ \quad ε) $0+(-4)$.

\ex Υπολόγισε: α) $(-8)+(-5)$ \quad β) $(+9)+(-4)$ \quad γ) $(-9)+(+4)$ \quad δ) $(+7)+(+5)$ \quad ε) $(-12)+0$ \quad στ) $(-15)+(+15)$ \quad ζ) $(+20)+(-35)$ \quad η) $(-17)+(+6)$.

\ex Υπολόγισε τις διαφορές: α) $(+3)-(+1)$ \quad β) $(+10)-(+20)$ \quad γ) $(+15)-(-5)$ \quad δ) $(-6)-(+2)$ \quad ε) $(-11)-(-6)$ \quad στ) $(+8)-(-8)$ \quad ζ) $(-14)-(-14)$ \quad η) $(-25)-(+8)$ \quad θ) $(-9)-(-17)$.

\ex Συμπλήρωσε τον αριθμό που λείπει ώστε το άθροισμα να είναι $0$: α) $(-7)+\kouti$ \quad β) $(+13)+\kouti$ \quad γ) $\kouti+(-4)$ \quad δ) $(-21)+\kouti$.

\ex Για κάθε ζεύγος $(\alpha,\beta)$ βρες τα $\alpha+\beta$, $\alpha-\beta$ και $\beta-\alpha$: α) $(+6,\,+2)$ \quad β) $(-12,\,-7)$ \quad γ) $(+9,\,-14)$ \quad δ) $(-2,\,+5)$. Τι παρατηρείς για τους αριθμούς $\alpha-\beta$ και $\beta-\alpha$;

\ex Υπολόγισε: α) $(-7)+12+(-5)+9+(-3)$ \quad β) $(-15)+4+(-25)+11$ \quad γ) $10+(-3)+(-4)+(-5)$.

\ex Σβήσε τις παρενθέσεις και υπολόγισε: α) $5+(-4)+(-2+6)$ \quad β) $-(+3)-(-4)-(-2+5)$ \quad γ) $(14-9+1)-(13-9)+3$ \quad δ) $12-(5-8+3)-(-4+9)$.

\ex Σωστό ή λάθος; α) Αν το άθροισμα δύο ακεραίων είναι θετικό, τότε και οι δύο είναι θετικοί. \quad β) Αν το άθροισμα δύο ακεραίων είναι αρνητικό, τότε και οι δύο είναι αρνητικοί. \quad γ) Η πρόσθεση στους ακέραιους σημαίνει πάντα αύξηση. \quad δ) Η αφαίρεση στους ακέραιους σημαίνει πάντα μείωση. \quad ε) Αν η διαφορά δύο ακεραίων είναι $0$, οι αριθμοί είναι ίσοι. \quad στ) $-(-1)=1$. \quad ζ) Το άθροισμα δύο αντίθετων αριθμών είναι $0$. \quad η) $5-7=7-5$.

\zoni{Το χρησιμοποιώ}{θερμόμετρα, ανελκυστήρες, ταμεία, παιχνίδια}

\ex \textbf{Η θερμοκρασία.} Το πρωί στο Μέτσοβο είναι $-6$°, το μεσημέρι ανεβαίνει $9$ βαθμούς και το βράδυ πέφτει $5$ βαθμούς. Πόση είναι η θερμοκρασία το μεσημέρι και το βράδυ;

\ex \textbf{Ο ανελκυστήρας.} Ο Ηλίας είναι στον όροφο $-2$, ανεβαίνει $5$ ορόφους και μετά κατεβαίνει $4$. Σε ποιον όροφο φτάνει;

\ex \textbf{Το παιχνίδι.} Οι πόντοι της Πολυξένης στους τέσσερις γύρους είναι $+8$, $-5$, $-7$, $+12$. Πόσους έχει συνολικά;

\ex \textbf{Ο λογαριασμός.} Η Ρήνη έχει υπόλοιπο $-45$ € στον λογαριασμό, καταθέτει $80$ € και μετά αναλαμβάνει $20$ €. Πόσα έχει;

\ex \textbf{Θερμική διαφορά.} Στο Μέτσοβο η μέγιστη θερμοκρασία ήταν $+4$° και η ελάχιστη $-6$°. Στη Φλώρινα η μέγιστη ήταν $-1$° και η ελάχιστη $-9$°. Σε ποια πόλη ήταν μεγαλύτερη η διαφορά μέγιστης και ελάχιστης;

\ex \textbf{Από την περιγραφή στην πράξη.} Γράψε μια παράσταση για καθεμιά από τις καταστάσεις και υπολόγισέ την. α) Ένα σημείο ξεκινά από το $-4$, κινείται $6$ μονάδες δεξιά και μετά $3$ αριστερά. β) Ο Βασίλης βρίσκεται στα $15$ μ. υψόμετρο, κατεβαίνει $25$ μ. και μετά ανεβαίνει $6$ μ. γ) Ο λογαριασμός της Φανής έχει $120$ €, γίνεται κατάθεση $40$ € και ανάληψη $90$ €.

\zoni{Το σκέφτομαι}{γιατί ισχύουν οι κανόνες}

\ex Πότε το άθροισμα δύο ακεραίων είναι $0$;

\ex Αν οι $\alpha,\,\beta$ είναι ετερόσημοι, ποιο πρόσημο έχει το άθροισμά τους;

\ex Είναι πάντα το άθροισμα μεγαλύτερο από κάθε προσθετέο; Δώσε αντιπαράδειγμα.

\ex Ο Στέλιος γράφει $5-(-3)=2$, γιατί «πέντε μείον τρία κάνει δύο». Ποιο είναι το λάθος του και ποιο το σωστό αποτέλεσμα;

\zoni{Για το σπίτι}{ο προϋπολογισμός σου}

\ex Κράτησε για μια εβδομάδα τα έσοδα και τα έξοδα του χαρτζιλικιού σου ως ακέραιους και υπολόγισε το τελικό υπόλοιπο.

\ex Φτιάξε ένα παιχνίδι με κάρτες $+$ και $-$ και παίξε με έναν φίλο τρεις γύρους. Ποιος έχει το μεγαλύτερο άθροισμα;

\paragrafos{Πολλαπλασιασμός ακεραίων: το μείον επί μείον}{Α.2.4, «Πολλαπλασιασμός ακεραίων»}

\begin{anagke}
Το βράδυ στο Μέτσοβο η θερμοκρασία πέφτει $2$ βαθμούς κάθε ώρα. Η Πολυξένη παρατηρεί το θερμόμετρο στις $20{:}00$, που δείχνει $0$°. «Τι θα δείχνει σε τρεις ώρες;». Ο Νίκος σκέφτεται: «Κάθε ώρα $-2$, τρεις ώρες: $-2-2-2=-6$». Η Αγγελική το γράφει ως πολλαπλασιασμό: $3\cdot(-2)=-6$. «Και πριν από τρεις ώρες, στις $17{:}00$;» ρωτάει η Ρήνη. «Αφού τώρα πέφτει δύο βαθμούς ανά ώρα, πριν από τρεις ώρες είχε $6$ βαθμούς περισσότερο», απαντάει η Πολυξένη. Ο ρυθμός είναι $-2$ και ο χρόνος «τρεις ώρες πριν» είναι $-3$, και το αποτέλεσμα είναι $+6$. Άρα $(-2)\cdot(-3)=+6$.

Ο Στέλιος ξαφνιάζεται: «Το μείον επί μείον κάνει συν; Σαν τη διπλή άρνηση της γλώσσας που καταλήγει σε κατάφαση;». Η κυρία Αλεξάνδρα του απαντά: «Η γλώσσα είναι ωραία αναλογία, αλλά τα μαθηματικά θέλουν να δουν γιατί». Το ίδιο το μοτίβο δείχνει το γιατί: $3\cdot(-2)=-6$, $2\cdot(-2)=-4$, $1\cdot(-2)=-2$, $0\cdot(-2)=0$. Κάθε φορά που το πρώτο μικραίνει κατά $1$ το γινόμενο μεγαλώνει κατά $2$. Συνεχίζοντας: $(-1)\cdot(-2)=2$, $(-2)\cdot(-2)=4$.

Σε αυτή την παράγραφο η παρέα μαθαίνει τον κανόνα των προσήμων στον πολλαπλασιασμό, τον εφαρμόζει σε γινόμενα με πολλούς παράγοντες και στις παραστάσεις με γράμματα, και προσέχει να μη μπερδεύει τον κανόνα του πολλαπλασιασμού με αυτόν της πρόσθεσης.
\end{anagke}

\begin{lexi}{παράγοντας, γινόμενο}{παράγοντας, γινόμενο}
\lexiword{Παράγοντες} είναι οι αριθμοί που πολλαπλασιάζονται και \lexiword{γινόμενο} αυτό που «γίνεται». Ο \lexiword{πολλαπλασιασμός} έρχεται από το «πολλαπλάσιο» και το «πλάσσω»: φτιάχνω πολλές φορές το ίδιο.
\end{lexi}

\begin{metafrasi}
Πέφτει $2$ βαθμούς την ώρα, για $3$ ώρες. & $3\cdot(-2)=-6$\\
Πριν από $3$ ώρες. & $(-2)\cdot(-3)=+6$\\
Ομόσημοι παράγοντες. & γινόμενο θετικό\\
Ετερόσημοι παράγοντες. & γινόμενο αρνητικό\\
Ένας παράγοντας μηδέν. & γινόμενο $0$\\
Το γινόμενο των $\alpha$ και $\beta$. & $\alpha\beta$ ή $\alpha\cdot\beta$\\
\end{metafrasi}

\tmima{Η θεωρία}

\tmimasmall{Γινόμενο δύο ακεραίων}

Σε ένα παιχνίδι γνώσεων, ο παίκτης κερδίζει $5$ πόντους ($+5$) για κάθε σωστή απάντηση και χάνει $2$ πόντους ($-2$) για κάθε λάθος. Αν απαντήσει σωστά σε $3$ ερωτήσεις έχει $(+3)\cdot(+5)=(+5)+(+5)+(+5)=+15$ πόντους. Αν απαντήσει λάθος σε $3$ ερωτήσεις έχει $(+3)\cdot(-2)=(-2)+(-2)+(-2)=-6$ πόντους. Παρατηρούμε ότι:

\begin{itemize}[leftmargin=1.2em,itemsep=0.3ex,topsep=2pt]
\item Το γινόμενο δύο θετικών είναι θετικός αριθμός.
\item Το γινόμενο ενός θετικού και ενός αρνητικού είναι αρνητικός αριθμός.
\end{itemize}

Για τους δύο αρνητικούς, ας δούμε το μοτίβο:
$$\begin{array}{rcl}
(-4)\cdot(+3)&=&-12\\
(-4)\cdot(+2)&=&-8\\
(-4)\cdot(+1)&=&-4\\
(-4)\cdot\phantom{(}0&=&0\\
(-4)\cdot(-1)&=&\ \ldots\\
(-4)\cdot(-2)&=&\ \ldots
\end{array}$$
Κάθε γραμμή δίνει αποτέλεσμα κατά $4$ μεγαλύτερο από την προηγούμενη. Αν συνεχίσουμε με τον ίδιο τρόπο, πρέπει να είναι $(-4)\cdot(-1)=+4$ και $(-4)\cdot(-2)=+8$. Διαπιστώνουμε ότι το γινόμενο δύο αρνητικών είναι θετικός αριθμός.

\begin{kanonas}{γινόμενο ακεραίων}
Για να πολλαπλασιάσω δύο ακέραιους, πολλαπλασιάζω τις απόλυτες τιμές τους και στο γινόμενο βάζω:
\begin{itemize}[leftmargin=1.2em,itemsep=0.3ex,topsep=2pt]
\item πρόσημο $+$, αν είναι \textbf{ομόσημοι}: $(+2)\cdot(+3)=+6$ και $(-2)\cdot(-3)=+6$.
\item πρόσημο $-$, αν είναι \textbf{ετερόσημοι}: $(+4)\cdot(-5)=-20$ και $(-3)\cdot(+3)=-9$.
\end{itemize}
Αν ένας παράγοντας είναι $0$, το γινόμενο είναι $0$.
\end{kanonas}

Σύντομα: $(+)\cdot(+)=+$, $(+)\cdot(-)=-$, $(-)\cdot(+)=-$, $(-)\cdot(-)=+$. Το $+$ δίνεται όταν τα πρόσημα είναι ίδια. Όταν χρησιμοποιούμε γράμματα, γράφουμε το γινόμενο χωρίς τελεία: $\alpha\beta$ αντί $\alpha\cdot\beta$ και $3\gamma$ αντί $3\cdot\gamma$.

\begin{center}
\begin{tikzpicture}[font=\small]
  \draw[rounded corners=3pt,fill=gblueL,draw=gblue] (0,0) rectangle (9.6,3.2);
  \node[anchor=west] at (0.3,2.65) {Ο φίλος του φίλου μου είναι φίλος μου.};
  \node[anchor=east] at (9.3,2.65) {$+\cdot+=+$};
  \node[anchor=west] at (0.3,1.95) {Ο εχθρός του φίλου μου είναι εχθρός μου.};
  \node[anchor=east] at (9.3,1.95) {$+\cdot-=-$};
  \node[anchor=west] at (0.3,1.25) {Ο φίλος του εχθρού μου είναι εχθρός μου.};
  \node[anchor=east] at (9.3,1.25) {$-\cdot+=-$};
  \node[anchor=west] at (0.3,0.55) {Ο εχθρός του εχθρού μου είναι φίλος μου.};
  \node[anchor=east] at (9.3,0.55) {$-\cdot-=+$};
\end{tikzpicture}
\end{center}

\tmimasmall{Γινόμενο πολλών παραγόντων}

Για να υπολογίσω ένα γινόμενο πολλών παραγόντων, πολλαπλασιάζω τις απόλυτες τιμές τους και στο γινόμενο βάζω:
\begin{itemize}[leftmargin=1.2em,itemsep=0.3ex,topsep=2pt]
\item πρόσημο $+$, αν το πλήθος των \textbf{αρνητικών} παραγόντων είναι \textbf{άρτιο},
\item πρόσημο $-$, αν το πλήθος των \textbf{αρνητικών} παραγόντων είναι \textbf{περιττό}.
\end{itemize}
Το πρόσημο λοιπόν το καθορίζει μόνο το πλήθος των αρνητικών παραγόντων. Αν τουλάχιστον ένας παράγοντας είναι $0$, το γινόμενο είναι $0$. Για παράδειγμα $(-2)\cdot3\cdot(-1)\cdot(-5)\cdot2=-60$ (τρεις αρνητικοί) και $(-1)\cdot(-1)\cdot(-1)\cdot(-1)=+1$ (τέσσερις αρνητικοί).

Ισχύουν οι γνωστές ιδιότητες: $\alpha\beta=\beta\alpha$, $(\alpha\beta)\gamma=\alpha(\beta\gamma)$, $\alpha\cdot1=\alpha$, $\alpha\cdot0=0$ και $\alpha\cdot(-1)=-\alpha$.

\begin{prosoxi}
\begin{itemize}[leftmargin=1.2em,itemsep=0.3ex,topsep=2pt]
\item Στον πολλαπλασιασμό δύο αρνητικοί δίνουν θετικό, ενώ στην πρόσθεση δύο αρνητικοί δίνουν αρνητικό: $(-3)\cdot(-4)=+12$ αλλά $(-3)+(-4)=-7$.
\item Το πρόσημο του γινομένου μετράει το πλήθος των \emph{αρνητικών} παραγόντων, όχι το πλήθος των θετικών.
\item Ο πολλαπλασιασμός δεν «μεγαλώνει» πάντα: $5\cdot(-2)=-10<5$.
\item Στις παραστάσεις με πράξεις, ο πολλαπλασιασμός προηγείται της πρόσθεσης και της αφαίρεσης: $5-3\cdot(-2)=5-(-6)=11$, όχι $2\cdot(-2)=-4$.
\item Αν στο γινόμενο υπάρχει παράγοντας $0$, δεν χρειάζεται να μετρήσω πρόσημα: το γινόμενο είναι $0$.
\end{itemize}
\end{prosoxi}

\tmima{Λυμένα παραδείγματα}

\paradeigma{Γινόμενα δύο παραγόντων}
\tixero Τα γινόμενα $(+6)\cdot(+7)$, $(+4)\cdot(-5)$, $(-8)\cdot(+3)$, $(-6)\cdot(-7)$, $(-15)\cdot0$.\par
\tithelo Να τα υπολογίσω.\par
\possindeo Πολλαπλασιάζω τις απόλυτες τιμές και βάζω $+$ αν τα πρόσημα είναι ίδια, $-$ αν είναι διαφορετικά.\par
\apantisi $42$, $-20$, $-24$, $42$, $0$.\par
\elegxos Στα ετερόσημα ($-20$ και $-24$) το αποτέλεσμα είναι αρνητικό και στα ομόσημα θετικό.

\paradeigma{Πολλοί παράγοντες}
\tixero Τα γινόμενα $(-2)\cdot3\cdot(-1)\cdot(-5)\cdot2$ και $2\cdot(-3)\cdot0\cdot(-7)$.\par
\tithelo Να τα υπολογίσω.\par
\possindeo Στο πρώτο μετράω τους αρνητικούς: $-2$, $-1$, $-5$, δηλαδή τρεις (περιττό πλήθος), άρα το γινόμενο είναι αρνητικό. Η απόλυτη τιμή είναι $2\cdot3\cdot1\cdot5\cdot2=60$. Το δεύτερο έχει παράγοντα $0$.\par
\apantisi $-60$ και $0$.\par
\elegxos Αν άλλαζα έναν παράγοντα σε θετικό, θα είχα δύο αρνητικούς και το αποτέλεσμα θα γινόταν θετικό.

\paradeigma{Πόντοι στο παιχνίδι}
\tixero Κάθε πράσινη κάρτα μετράει $+3$ πόντους και κάθε μαύρη $-2$. Η Φανή έχει $4$ πράσινες και $3$ μαύρες.\par
\tithelo Πόσους πόντους έχει;\par
\possindeo $4\cdot(+3)+3\cdot(-2)=12+(-6)=6$.\par
\apantisi $6$ πόντους.\par
\elegxos Οι πράσινες δίνουν $12$ και οι μαύρες αφαιρούν $6$: $12-6=6$.

\paradeigma{Παράσταση με πράξεις και παρενθέσεις}
\tixero Η παράσταση $5\cdot(-2+6)-3\cdot(-7+4)$.\par
\tithelo Να την υπολογίσω.\par
\possindeo Πρώτα οι παρενθέσεις: $-2+6=4$ και $-7+4=-3$. Μετά οι πολλαπλασιασμοί: $5\cdot4=20$ και $3\cdot(-3)=-9$. Τέλος η αφαίρεση: $20-(-9)=29$.\par
\apantisi $29$.\par
\elegxos Το δεύτερο γινόμενο είναι αρνητικό και το αφαιρούμε: ένα «μείον» μπροστά σε αρνητικό δίνει $+$.

\paradeigma{Αριθμητικές τιμές παραστάσεων}
\tixero Οι τιμές $\alpha=-2$, $\beta=+5$, $\gamma=-3$.\par
\tithelo Να βρω τις τιμές των $\alpha\beta$, $\alpha\gamma$, $\alpha\beta\gamma$ και $3\beta-\alpha\gamma$.\par
\possindeo $\alpha\beta=(-2)\cdot5=-10$. $\alpha\gamma=(-2)\cdot(-3)=6$. $\alpha\beta\gamma=(-10)\cdot(-3)=30$. $3\beta-\alpha\gamma=3\cdot5-6=15-6=9$.\par
\apantisi $-10$, $6$, $30$, $9$.\par
\elegxos Στο $\alpha\beta\gamma$ οι αρνητικοί παράγοντες είναι δύο (άρτιο πλήθος): θετικό γινόμενο.

\paradeigma{Ποιοι δύο ακέραιοι;}
\tixero Δύο ακέραιοι έχουν γινόμενο $+12$ και άθροισμα $-7$.\par
\tithelo Να τους βρω.\par
\possindeo Αφού το γινόμενο είναι θετικό, είναι ομόσημοι. Αφού το άθροισμα είναι αρνητικό, είναι και οι δύο αρνητικοί. Ζεύγη με γινόμενο $12$: $(-1,\,-12)$, $(-2,\,-6)$, $(-3,\,-4)$. Το άθροισμα $-7$ το δίνει το τελευταίο.\par
\apantisi $-3$ και $-4$.\par
\elegxos $(-3)\cdot(-4)=12$ και $(-3)+(-4)=-7$.

\tmima{Ασκήσεις}

\zoni{Το κατάλαβα}{κανόνας προσήμων, πολλοί παράγοντες}

\ex Υπολόγισε τα γινόμενα: α) $(+2)\cdot(+3)$ \quad β) $(-4)\cdot(+6)$ \quad γ) $(-6)\cdot(-7)$ \quad δ) $(+5)\cdot(-5)$ \quad ε) $(-15)\cdot0$ \quad στ) $(-1)\cdot(+207)$.

\ex Συμπλήρωσε με τη σωστή λέξη. α) Το γινόμενο δύο ομοσήμων ακεραίων είναι πάντα \ldots\ αριθμός. β) Το γινόμενο δύο ετεροσήμων ακεραίων είναι πάντα \ldots\ αριθμός. γ) Το πρόσημο ενός γινομένου καθορίζεται από το πλήθος των \ldots\ παραγόντων. δ) Αν το γινόμενο δύο ακεραίων είναι αρνητικό, οι αριθμοί είναι \ldots.

\ex Υπολόγισε: α) $(-1)\cdot(-1)\cdot(-1)\cdot(-1)\cdot(-1)$ \quad β) $(-1)\cdot(-2)\cdot(+3)\cdot(-4)\cdot(-5)$ \quad γ) $(+5)\cdot(-10)\cdot(+5)\cdot(-10)\cdot(-2)$.

\ex Συμπλήρωσε τον πίνακα, όπου κάθε κελί είναι το γινόμενο της γραμμής με τη στήλη.
\begin{center}
\renewcommand{\arraystretch}{1.3}
\begin{tabular}{c|c|c|c|c}
$\cdot$ & $+3$ & $-4$ & $-6$ & $0$\\\hline
$-5$ & & & &\\\hline
$-1$ & & & &\\\hline
$+2$ & & & &\\\hline
$7$ & & & &\\
\end{tabular}
\end{center}

\ex Βρες τον αριθμό που λείπει ώστε να αληθεύουν οι ισότητες: α) $(+3)\cdot(-4)=(-2)\cdot(\ldots)$ \quad β) $(-5)\cdot(-6)=(+3)\cdot(\ldots)$ \quad γ) $(+9)\cdot(+8)=(-6)\cdot(\ldots)$ \quad δ) $(-3)\cdot8=12\cdot(\ldots)$.

\ex Υπολόγισε: α) $2\cdot(-5+3)+4\cdot(-1+6)$ \quad β) $(-7+11-13)\cdot(-5+4)$ \quad γ) $(+6)\cdot(-3)+(-1)\cdot(+4)-2\cdot(-5)\cdot(-4)$ \quad δ) $3\cdot\big[\,2\cdot(-4)-(-5)\cdot2\,\big]$.

\ex Σωστό ή λάθος; α) Το γινόμενο ενός θετικού και ενός αρνητικού είναι αρνητικό. \quad β) $(-3)\cdot(-4)=-12$. \quad γ) Αν ένας παράγοντας είναι $0$, το γινόμενο είναι $0$. \quad δ) Το γινόμενο $(-1)\cdot(-1)\cdot\ldots\cdot(-1)$ με $100$ παράγοντες είναι $-1$. \quad ε) $3\cdot(-2)=(-2)\cdot3$. \quad στ) Το γινόμενο αρνητικών ακεραίων είναι πάντα αρνητικό. \quad ζ) $\alpha\cdot(-1)=-\alpha$.

\ex Για $\alpha=-1$, $\beta=+4$, $\gamma=-5$ υπολόγισε τα $\alpha\beta$, $\beta\gamma$, $\alpha\beta\gamma$, $\alpha\gamma-\beta$ και $\alpha-\beta\gamma$.

\zoni{Το χρησιμοποιώ}{θερμόμετρα, ασανσέρ, παιχνίδια}

\ex \textbf{Η νύχτα.} Στις $20{:}00$ το θερμόμετρο δείχνει $0$° και πέφτει $3$° κάθε ώρα. Η Φανή ρωτά: α) Τι θα δείχνει σε $4$ ώρες; β) Τι έδειχνε $2$ ώρες πριν; Γράψε τις πράξεις με τον ρυθμό $-3$.

\ex \textbf{Το παιχνίδι ερωτήσεων.} Κάθε σωστή απάντηση δίνει $+5$ πόντους και κάθε λάθος $-2$. Ο Ηλίας απαντά σωστά σε $7$ ερωτήσεις και λάθος σε $4$. Ο Στέλιος σωστά σε $3$ και λάθος σε $9$. Πόσους πόντους έχει ο καθένας;

\ex \textbf{Ο ανελκυστήρας.} Ο ανελκυστήρας του ξενοδοχείου κατεβαίνει $2$ ορόφους κάθε φορά που πατάμε «κάτω». Ο Νίκος ξεκινά από τον $6$ο όροφο και πατάει «κάτω» $5$ φορές. Σε ποιον όροφο βρίσκεται;

\ex \textbf{Το νερό του χωριού.} Στη δεξαμενή του χωριού, η γιαγιά Μυρτώ βλέπει ότι το νερό μειώνεται κατά $15$ λίτρα κάθε μέρα (ρυθμός $-15$). Πόσα λίτρα λιγότερα υπάρχουν μετά από $6$ μέρες; Πόσα περισσότερα υπήρχαν πριν από $4$ μέρες;

\ex \textbf{Το μοτίβο της Αγγελικής.} Η Αγγελική γράφει: $(-6)\cdot3=-18$, $(-6)\cdot2=-12$, $(-6)\cdot1=-6$, $(-6)\cdot0=0$. Συνέχισε το μοτίβο με $(-6)\cdot(-1)$ και $(-6)\cdot(-2)$. Τι παρατηρείς;

\zoni{Το σκέφτομαι}{ομόσημοι, ετερόσημοι, πρόσημο γινομένου}

\ex Αν $\alpha\beta>0$, τι ξέρεις για τα πρόσημα των $\alpha$ και $\beta$; Και αν $\alpha\beta<0$;

\ex Πόσοι από τους $5$ παράγοντες ενός γινομένου ακεραίων μπορεί να είναι αρνητικοί, ώστε το γινόμενο να είναι θετικό (όλοι οι παράγοντες μη μηδενικοί);

\ex Ο Στέλιος λέει: «Μείον επί μείον κάνει συν, άρα και μείον συν μείον κάνει συν: $(-3)+(-4)=+7$». Τι μπερδεύει;

\ex Βρες δύο ακέραιους με γινόμενο $-12$ και άθροισμα $1$.

\zoni{Για το σπίτι}{πίνακες και παραδείγματα}

\ex Φτιάξε έναν πίνακα πολλαπλασιασμού με τους ακέραιους από το $-3$ ως το $3$ σε γραμμές και στήλες. Τι παρατηρείς στη γραμμή και στη στήλη του $0$;

\ex Βρες ένα δικό σου παράδειγμα από την καθημερινή ζωή που εξηγεί γιατί το γινόμενο δύο αρνητικών είναι θετικό (π.χ. μια ταινία που παίζει ανάποδα).

\paragrafos{Πολλαπλασιασμός κλασμάτων: το «του» είναι «επί»}{Α.2.5, «Πολλαπλασιασμός κλασμάτων»}

\begin{anagke}
Η Αγγελική ετοιμάζει ένα πίνακα για τη σχολική έκθεση. Παίρνει ένα ορθογώνιο χαρτόνι και ζωγραφίζει θάλασσα στα τρία τέταρτα του. Από τη θάλασσα, τα δύο τρίτα τα βάφει σκούρο μπλε, για το βάθος. «Τι μέρος του χαρτονιού είναι σκούρο μπλε;» ρωτάει η Ζωή.

Ο Ηλίας, που δοκιμάζει πάντα πρώτα κάτι με τα χέρια του, κόβει ένα χαρτί και το χωρίζει σε 4 ίσες στήλες. Σκιάζει τις 3. Μετά χωρίζει όλο το σχήμα σε 3 ίσες σειρές και σκουραίνει τις 2 από τις 3 σειρές πάνω στη σκιασμένη περιοχή. Μετράει: από 12 ίσα τετράγωνα, τα 6 είναι σκούρα. «Τα μισά!» λέει.

Η Μυρτώ σκέφτεται δυνατά: «Τα δύο τρίτα των τριών τετάρτων είναι το μισό του χαρτονιού. Αλλά όταν πολλαπλασιάζεις, δεν βγαίνει πάντα κάτι μεγαλύτερο; Εδώ βγήκε μικρότερο και από τα δύο!» Και η Πολυξένη προσθέτει: «Και αν βρούμε έναν κανόνα, ώστε να μην χρειάζεται να ζωγραφίζουμε κάθε φορά;»
\end{anagke}

\begin{lexi}{αντίστροφος}{αντίστροφος}
Η λέξη \lexiword{αντίστροφος} έρχεται από το «αντί» και το «στρέφω», δηλαδή «γυρίζω ανάποδα». Ο αντίστροφος ενός κλάσματος είναι το ίδιο κλάσμα γυρισμένο ανάποδα: ο αριθμητής γίνεται παρονομαστής και ο παρονομαστής αριθμητής. Δύο αριθμοί είναι αντίστροφοι όταν το γινόμενό τους είναι 1.

Και μια μικρή λέξη που κρύβει τα πάντα: το «\lexiword{του}» σε φράσεις όπως «τα $\frac{2}{3}$ του $\frac{3}{4}$» σημαίνει \emph{επί}. «Τα δύο τρίτα του» είναι «δύο τρίτα επί».
\end{lexi}

\begin{metafrasi}
Τα $\frac{2}{3}$ των $\frac{3}{4}$. & $\dfrac{2}{3}\cdot\dfrac{3}{4}=\dfrac{6}{12}=\dfrac{1}{2}$\\
Τα $\frac{3}{5}$ του 20. & $\dfrac{3}{5}\cdot 20=12$\\
Ο αντίστροφος του $\frac{3}{7}$. & $\dfrac{7}{3}$, γιατί $\dfrac{3}{7}\cdot\dfrac{7}{3}=1$\\
Ο αντίστροφος του 5. & $\dfrac{1}{5}$, γιατί $5\cdot\dfrac{1}{5}=1$\\
Δύο ολόκληρα και μισό, τέσσερις φορές. & $2\dfrac{1}{2}\cdot 4=\dfrac{5}{2}\cdot 4=10$\\
\end{metafrasi}

\tmima{Η θεωρία}

\begin{kanonas}{γινόμενο κλασμάτων}
Το γινόμενο δύο κλασμάτων είναι κλάσμα με αριθμητή το γινόμενο των αριθμητών και παρονομαστή το γινόμενο των παρονομαστών:
$$\dfrac{\alpha}{\beta}\cdot\dfrac{\gamma}{\delta}=\dfrac{\alpha\cdot\gamma}{\beta\cdot\delta}.$$
Το γινόμενο φυσικού επί κλάσμα είναι κλάσμα με αριθμητή το γινόμενο του αριθμητή επί τον φυσικό και τον ίδιο παρονομαστή: $\lambda\cdot\dfrac{\alpha}{\beta}=\dfrac{\lambda\cdot\alpha}{\beta}$.
\end{kanonas}

Γιατί δουλεύει; Πάρε το χαρτόνι της Αγγελικής. Χωρίζεται σε $4\cdot 3=12$ ίσα τετράγωνα, και τα σκούρα είναι $3\cdot 2=6$. Δηλαδή $\frac{3}{4}\cdot\frac{2}{3}=\frac{3\cdot 2}{4\cdot 3}=\frac{6}{12}$. Κάθε φορά οι παρονομαστές δίνουν πόσα τετράγωνα έχει το όλο και οι αριθμητές πόσα χρωματίστηκαν.

\begin{center}
\begin{tikzpicture}[x=0.9cm,y=0.9cm]
\fill[gblue!25] (0,0) rectangle (3,3);
\fill[gblue!70] (0,1) rectangle (3,3);
\foreach \i in {0,...,4}{\draw (\i,0)--(\i,3);}
\foreach \j in {0,...,3}{\draw (0,\j)--(4,\j);}
\draw[very thick] (0,0) rectangle (4,3);
\node[below] at (2,-0.1) {\small 4 στήλες, σκιασμένες 3};
\node[right] at (4.2,1.5) {\small 3 σειρές, σκούρες 2};
\end{tikzpicture}
\end{center}

\tmimasmall{Αντίστροφοι}

Δύο κλάσματα λέγονται \textbf{αντίστροφα} όταν το γινόμενό τους είναι 1: $\dfrac{\alpha}{\beta}\cdot\dfrac{\beta}{\alpha}=\dfrac{\alpha\beta}{\beta\alpha}=1$. Ο αντίστροφος του $\frac{3}{7}$ είναι ο $\frac{7}{3}$. Ο αντίστροφος του φυσικού $5=\frac{5}{1}$ είναι ο $\frac{1}{5}$. Ο αντίστροφος του 1 είναι ο ίδιος ο 1. Το 0 \textbf{δεν έχει αντίστροφο}, γιατί δεν υπάρχει αριθμός που επί 0 να κάνει 1.

Όλες οι ιδιότητες του πολλαπλασιασμού των φυσικών ισχύουν και στα κλάσματα: αντιμεταθετική, προσεταιριστική, επιμεριστική και ο αριθμός 1 που δεν αλλάζει το γινόμενο.

\tmimasmall{Μεικτοί και απλοποίηση}

Για να πολλαπλασιάσω μεικτούς αριθμούς, τους γράφω πρώτα ως κλάσματα: $2\frac{1}{3}\cdot 1\frac{1}{2}=\frac{7}{3}\cdot\frac{3}{2}$. Πριν πολλαπλασιάσω, μπορώ να απλοποιήσω «χιαστί»: το 3 του παρονομαστή και το 3 του αριθμητή απλοποιούνται, και μένει $\frac{7}{2}=3\frac{1}{2}$. Έτσι δουλεύω με μικρότερους αριθμούς.

\begin{prosoxi}
\begin{itemize}[leftmargin=1.2em,itemsep=0.3ex,topsep=2pt]
\item Ο πολλαπλασιασμός με κλάσμα μικρότερο από το 1 \textbf{μικραίνει}. Τα $\frac{2}{3}$ του $\frac{3}{4}$ είναι μικρότερα και από τα δύο: $\frac{1}{2}<\frac{2}{3}$ και $\frac{1}{2}<\frac{3}{4}$. Ο πολλαπλασιασμός με κλάσμα μεγαλύτερο από το 1 μεγαλώνει.
\item Δεν πολλαπλασιάζουμε χωριστά τα ολόκληρα και τα κλάσματα ενός μεικτού: $2\frac{1}{2}\cdot 3\frac{1}{2}$ δεν είναι $(2\cdot 3)+\left(\frac{1}{2}\cdot\frac{1}{2}\right)=6\frac{1}{4}$. Σωστά: $\frac{5}{2}\cdot\frac{7}{2}=\frac{35}{4}=8\frac{3}{4}$.
\item Δεν πολλαπλασιάζουμε «και παρονομαστές» όταν προσθέτουμε: στο άθροισμα κρατάμε τον ίδιο παρονομαστή, στο γινόμενο τους πολλαπλασιάζουμε.
\end{itemize}
\end{prosoxi}

\tmima{Λυμένα παραδείγματα}

\paradeigma{Το χαρτόνι της Αγγελικής}
\tixero Τα $\frac{3}{4}$ του χαρτονιού είναι θάλασσα, και τα $\frac{2}{3}$ της θάλασσας είναι σκούρα.\par
\tithelo Τι μέρος του χαρτονιού είναι σκούρο.\par
\possindeo Το «τα $\frac{2}{3}$ των $\frac{3}{4}$» σημαίνει $\frac{2}{3}\cdot\frac{3}{4}=\frac{2\cdot 3}{3\cdot 4}=\frac{6}{12}$. Το $\frac{6}{12}$ απλοποιείται με το 6: $\frac{1}{2}$. Μπορώ να απλοποιήσω και πριν: $\frac{2\cdot\cancel{3}}{\cancel{3}\cdot 4}=\frac{2}{4}=\frac{1}{2}$.\par
\apantisi Το $\dfrac{1}{2}$ του χαρτονιού.\par
\elegxos Το σχήμα της θεωρίας δείχνει 6 σκούρα τετράγωνα από 12. Σωστά.

\paradeigma{Οι μαθητές της Α’}
\tixero Η Α’ τάξη έχει 96 μαθητές και τα $\frac{5}{8}$ έχουν αδέλφια.\par
\tithelo Πόσοι έχουν αδέλφια και πόσοι όχι.\par
\possindeo Τα $\frac{5}{8}$ του 96 είναι $\frac{5}{8}\cdot 96=\frac{5\cdot 96}{8}=\frac{480}{8}=60$. Οι υπόλοιποι είναι $96-60=36$.\par
\apantisi 60 έχουν αδέλφια και 36 δεν έχουν.\par
\elegxos Τα υπόλοιπα $\frac{3}{8}$ του 96 είναι $\frac{3\cdot 96}{8}=36$. Συμφωνεί.

\paradeigma{Μεικτοί αριθμοί}
\tixero Οι αριθμοί $2\frac{1}{3}$ και $1\frac{1}{2}$.\par
\tithelo Το γινόμενό τους.\par
\possindeo $2\frac{1}{3}=\frac{7}{3}$ και $1\frac{1}{2}=\frac{3}{2}$. Άρα $\frac{7}{3}\cdot\frac{3}{2}=\frac{7\cdot\cancel{3}}{\cancel{3}\cdot 2}=\frac{7}{2}=3\frac{1}{2}$.\par
\apantisi $3\dfrac{1}{2}$.\par
\elegxos Το $2\frac{1}{3}$ είναι λίγο πάνω από 2, και το $1\frac{1}{2}$ είναι το ενάμισι. Ενάμισι φορά το 2 και κάτι είναι λίγο πάνω από 3. Το $3\frac{1}{2}$ βγάζει νόημα.

\paradeigma{Οι αντίστροφοι}
\tixero Οι αριθμοί $\frac{4}{9}$, 6 και $1\frac{1}{2}$.\par
\tithelo Οι αντίστροφοί τους.\par
\possindeo Γυρίζω το κλάσμα ανάποδα. Για το $\frac{4}{9}$: $\frac{9}{4}$. Για το $6=\frac{6}{1}$: $\frac{1}{6}$. Για το $1\frac{1}{2}=\frac{3}{2}$: $\frac{2}{3}$.\par
\apantisi $\dfrac{9}{4}$, $\dfrac{1}{6}$ και $\dfrac{2}{3}$.\par
\elegxos $\frac{4}{9}\cdot\frac{9}{4}=1$, $6\cdot\frac{1}{6}=1$, $\frac{3}{2}\cdot\frac{2}{3}=1$. Σωστά.

\tmima{Ασκήσεις}

\zoni{Το κατάλαβα}{εφαρμογή του κανόνα}

\ex Βρες τα γινόμενα και απλοποίησε: α) $\frac{2}{3}\cdot\frac{5}{7}$ \quad β) $\frac{3}{4}\cdot\frac{8}{9}$ \quad γ) $\frac{5}{6}\cdot\frac{12}{25}$ \quad δ) $\frac{7}{10}\cdot\frac{20}{21}$ \quad ε) $\frac{4}{9}\cdot\frac{3}{8}$ \quad στ) $\frac{9}{14}\cdot\frac{7}{12}$.

\ex Υπολόγισε: α) $4\cdot\frac{2}{5}$ \quad β) $\frac{3}{7}\cdot 14$ \quad γ) $6\cdot\frac{5}{9}$ \quad δ) $12\cdot\frac{3}{8}$ \quad ε) $\frac{5}{6}\cdot 18$ \quad στ) $100\cdot\frac{7}{10}$.

\ex Βρες: α) το $\frac{1}{2}$ του $\frac{1}{2}$ \quad β) το $\frac{1}{3}$ του $\frac{1}{5}$ \quad γ) τα $\frac{2}{3}$ του $\frac{3}{4}$ \quad δ) τα $\frac{3}{4}$ του $\frac{4}{5}$.

\ex Υπολόγισε τα γινόμενα με μεικτούς: α) $2\frac{1}{3}\cdot 3$ \quad β) $1\frac{1}{2}\cdot 2\frac{2}{3}$ \quad γ) $3\frac{1}{4}\cdot 2$ \quad δ) $2\frac{1}{2}\cdot 1\frac{3}{5}$.

\ex Βρες τον αντίστροφο των αριθμών: $\frac{7}{5}$, 3, $\frac{1}{8}$, 1, $\frac{11}{4}$ και $2\frac{1}{2}$.

\ex Συμπλήρωσε ώστε το γινόμενο να είναι 1: α) $\frac{3}{4}\cdot\square$ \quad β) $\frac{5}{9}\cdot\square$ \quad γ) $6\cdot\square$ \quad δ) $\square\cdot\frac{2}{7}$.

\ex Απλοποίησε πριν πολλαπλασιάσεις: α) $\frac{14}{15}\cdot\frac{25}{28}$ \quad β) $\frac{18}{35}\cdot\frac{21}{27}$ \quad γ) $\frac{8}{9}\cdot\frac{15}{16}$ \quad δ) $\frac{4}{21}\cdot\frac{35}{12}$.

\ex Βρες την τιμή, με προτεραιότητα των πράξεων: α) $\frac{2}{3}+\frac{1}{2}\cdot\frac{3}{4}$ \quad β) $\left(\frac{2}{3}+\frac{1}{2}\right)\cdot\frac{3}{4}$ \quad γ) $\frac{1}{2}\cdot\left(\frac{3}{4}-\frac{1}{6}\right)$.

\ex Σωστό ή λάθος; α) $\frac{2}{3}\cdot\frac{3}{2}=1$ \quad β) $\frac{3}{5}\cdot 4=\frac{12}{20}$ \quad γ) $\frac{1}{2}\cdot\frac{1}{3}=\frac{1}{5}$ \quad δ) Ο αντίστροφος του 0 είναι το 0. \quad ε) $4\cdot\frac{1}{4}=1$ \quad στ) $2\frac{1}{2}\cdot 2=4\frac{1}{2}$.

\zoni{Το χρησιμοποιώ}{προβλήματα από τη ζωή}

\ex \textbf{Η συνταγή.} Η συνταγή της γιαγιάς Ρήνης για 6 άτομα θέλει $\frac{3}{4}$ του κιλού αλεύρι. Η Μυρτώ θέλει να μαγειρέψει για 15 άτομα, δηλαδή για $2\frac{1}{2}$ φορές περισσότερα. Πόσο αλεύρι χρειάζεται;

\ex \textbf{Ο χυμός.} Ένα μπουκάλι έχει $1\frac{1}{2}$ λίτρα χυμό. Ο Νίκος πίνει τα $\frac{2}{5}$ του μπουκαλιού. Πόσα λίτρα ήπιε;

\ex \textbf{Ο κήπος.} Τα $\frac{3}{5}$ του κήπου της γιαγιάς είναι λαχανικά και τα $\frac{2}{3}$ από αυτά είναι ντομάτες. Τι μέρος ολόκληρου του κήπου είναι οι ντομάτες; Αν ο κήπος έχει 300 τ.μ., πόσα τ.μ. έχουν ντομάτες;

\ex \textbf{Το εμβαδό.} Ένα ορθογώνιο έχει μήκος $3\frac{1}{2}$ μ. και πλάτος $2\frac{2}{5}$ μ. Βρες το εμβαδό του.

\ex \textbf{Ο περίπατος.} Ο Βασίλης περπατάει $4\frac{1}{2}$ χλμ. την ώρα. Πόσα χιλιόμετρα κάνει σε $\frac{2}{3}$ της ώρας, δηλαδή σε 40 λεπτά;

\ex \textbf{Η έκπτωση.} Ο Νίκος βρίσκει αθλητικά παπούτσια που κοστίζουν 80 € και έχουν έκπτωση τα $\frac{3}{10}$ της τιμής τους. Πόσα ευρώ είναι η έκπτωση; Πόσο πληρώνει; (Δοκίμασε να βρεις και τα $\frac{7}{10}$ της τιμής.)

\zoni{Το σκέφτομαι}{μια ερώτηση που πάει πιο μακριά}

\ex Ο Στέλιος λέει: «Ο πολλαπλασιασμός κάνει τα πράγματα μεγαλύτερα, άρα το $\frac{1}{2}\cdot\frac{1}{3}$ είναι πάνω από το $\frac{1}{2}$». Εξήγησε με το «μισό του ενός τρίτου» ότι βγαίνει μικρότερο.

\ex Πολλαπλασίασε το 12 επί $\frac{3}{4}$ και επί $\frac{5}{4}$. Πότε ο πολλαπλασιασμός με κλάσμα μικραίνει τον αριθμό και πότε τον μεγαλώνει;

\ex Ποιο κλάσμα πρέπει να πολλαπλασιάσεις με το $\frac{3}{8}$ για να βρεις $\frac{1}{2}$; (Σκέψου τον αντίστροφο του $\frac{3}{8}$.)

\ex Υπολόγισε $\frac{1}{2}\cdot\frac{2}{3}\cdot\frac{3}{4}\cdot\frac{4}{5}\cdot\frac{5}{6}$. Τι παρατηρείς; Μάντεψε πόσο θα βγει το $\frac{1}{2}\cdot\frac{2}{3}\cdot\dots\cdot\frac{9}{10}$.

\ex Γιατί το 0 δεν έχει αντίστροφο; Εξήγησε με την εξίσωση $0\cdot x=1$.

\zoni{Για το σπίτι}{με χαρτί και με το μυαλό}

\ex Πάρε ένα φύλλο χαρτί, δίπλωσέ το σε 3 ίσες σειρές και σε 4 ίσες στήλες. Σκίασε τις 3 από τις 4 στήλες και μετά σκούρυνε 2 από τις 3 σειρές πάνω στη σκιασμένη περιοχή. Δείξε με το φύλλο ότι $\frac{2}{3}\cdot\frac{3}{4}=\frac{1}{2}$.

\ex Βρες το $\frac{1}{2}$ του $\frac{1}{2}$ του $\frac{1}{2}$ του $\frac{1}{2}$ του 80.

\paragrafos{Διαίρεση κλασμάτων: πόσες φορές χωράει}{Α.2.6, «Διαίρεση κλασμάτων»}

\begin{anagke}
Ο παππούς Γιάννης ετοιμάζει δείγματα ελαιόλαδου για το παζάρι. Έχει $\frac{3}{4}$ του λίτρου λάδι και θέλει να το βάλει σε μικρά μπουκαλάκια που χωράνε το καθένα $\frac{1}{8}$ του λίτρου. «Πόσα μπουκαλάκια θα γεμίσουν;» ρωτάει, και η Ζωή και ο Ηλίας αρχίζουν να σκέφτονται.

Ο Ηλίας γράφει τα $\frac{3}{4}$ ως $\frac{6}{8}$. «Το κάθε μπουκαλάκι είναι ένα όγδοο. Έχω έξι όγδοα. Έξι μπουκαλάκια». Η Ζωή ρωτάει: «Αυτό είναι διαίρεση. Πώς γράφεται;» Γράφουν $\frac{3}{4}:\frac{1}{8}$. Ο Στέλιος, που θέλει πάντα έναν κανόνα, προσπαθεί: «Διαιρώ αριθμητή με αριθμητή και παρονομαστή με παρονομαστή. Το $3:1=3$ και το $4:8$... δεν βγαίνει φυσικός». Η Ρήνη σηκώνει τους ώμους: «Βγάζεις περίεργα νούμερα».

Η Πολυξένη παρατηρεί κάτι στην απάντηση του Ηλία: έξι είναι ακριβώς $\frac{3}{4}\cdot 8$. «Μήπως η διαίρεση με το $\frac{1}{8}$ είναι πολλαπλασιασμός με το 8, δηλαδή με τον αντίστροφό του;» Αν είναι έτσι, ο κανόνας υπάρχει και δεν χρειάζεται καν τα όγδοα.
\end{anagke}

\begin{lexi}{σύνθετο κλάσμα}{σύνθετο κλάσμα}
Η λέξη \lexiword{σύνθετο} σημαίνει «φτιαγμένο από πολλά κομμάτια». Σύνθετο κλάσμα είναι ένα κλάσμα που έχει μέσα του άλλα κλάσματα, στον αριθμητή ή στον παρονομαστή ή και στα δύο, για παράδειγμα $\dfrac{\frac{2}{3}}{\frac{4}{9}}$. Η κλασματική γραμμή του σημαίνει διαίρεση: $\frac{2}{3}:\frac{4}{9}$.

Και η διαίρεση, στη γλώσσα της ζωής, απαντάει στο \emph{πόσες φορές χωράει}: πόσες φορές χωράει το μικρό μπουκαλάκι στο λάδι που έχω.
\end{lexi}

\begin{metafrasi}
Πόσες φορές χωράει το $\frac{1}{8}$ στα $\frac{3}{4}$; & $\dfrac{3}{4}:\dfrac{1}{8}=\dfrac{3}{4}\cdot\dfrac{8}{1}=6$\\
Διαίρεση σημαίνει πολλαπλασιασμός με τον αντίστροφο. & $\dfrac{\alpha}{\beta}:\dfrac{\gamma}{\delta}=\dfrac{\alpha}{\beta}\cdot\dfrac{\delta}{\gamma}$\\
Ο φυσικός 6 διαιρούμενος με το $\frac{2}{3}$. & $6:\dfrac{2}{3}=6\cdot\dfrac{3}{2}=9$\\
Το μισό από τα $\frac{3}{4}$. & $\dfrac{3}{4}:2=\dfrac{3}{4}\cdot\dfrac{1}{2}=\dfrac{3}{8}$\\
Σύνθετο κλάσμα. & $\dfrac{\frac{2}{3}}{\frac{5}{7}}=\dfrac{2}{3}:\dfrac{5}{7}=\dfrac{14}{15}$\\
\end{metafrasi}

\tmima{Η θεωρία}

Στους φυσικούς, για να διαιρέσουμε το 6 με το 3, ρωτάμε πόσες φορές χωράει το 3 στο 6. Το ίδιο ισχύει και στα κλάσματα. Και υπάρχει ένας πολύ βολικός κανόνας.

\begin{kanonas}{διαίρεση}
\begin{itemize}[leftmargin=1.2em,itemsep=0.3ex,topsep=2pt]
\item Για να διαιρέσω δύο φυσικούς: πολλαπλασιάζω τον διαιρετέο με τον αντίστροφο του διαιρέτη. $\alpha:\beta=\alpha\cdot\dfrac{1}{\beta}$.
\item Για να διαιρέσω δύο κλάσματα: πολλαπλασιάζω τον διαιρετέο με τον \textbf{αντίστροφο του διαιρέτη}. $\dfrac{\alpha}{\beta}:\dfrac{\gamma}{\delta}=\dfrac{\alpha}{\beta}\cdot\dfrac{\delta}{\gamma}$.
\end{itemize}
\end{kanonas}

Γιατί; Στο λάδι του παππού: $\frac{3}{4}:\frac{1}{8}$ ρωτάει πόσες φορές χωράει το $\frac{1}{8}$ στα $\frac{3}{4}$. Σε ένα ολόκληρο λίτρο χωράνε 8 τέτοια μπουκαλάκια (γιατί $8\cdot\frac{1}{8}=1$). Στα $\frac{3}{4}$ του λίτρου χωράνε τα $\frac{3}{4}$ από το 8, δηλαδή $\frac{3}{4}\cdot 8=6$. Ο αριθμός 8 είναι ο αντίστροφος του $\frac{1}{8}$.

\tmimasmall{Σύνθετο κλάσμα}

Ένα κλάσμα που έχει κλάσμα στον αριθμητή ή στον παρονομαστή λέγεται \textbf{σύνθετο}. Η κλασματική του γραμμή σημαίνει διαίρεση, άρα το μετατρέπω σε απλό με τον κανόνα της διαίρεσης:
$$\dfrac{\frac{\alpha}{\beta}}{\frac{\gamma}{\delta}}=\dfrac{\alpha}{\beta}:\dfrac{\gamma}{\delta}=\dfrac{\alpha\cdot\delta}{\beta\cdot\gamma}.$$
Αν ο διαιρετέος ή ο διαιρέτης είναι μεικτός αριθμός, τον γράφω πρώτα ως κλάσμα.

\begin{prosoxi}
\begin{itemize}[leftmargin=1.2em,itemsep=0.3ex,topsep=2pt]
\item Αντιστρέφουμε \textbf{μόνο τον διαιρέτη}, όχι τον διαιρετέο. $\frac{3}{4}:\frac{1}{8}=\frac{3}{4}\cdot 8=6$, ενώ ο «αντίστροφος και των δύο», $\frac{4}{3}\cdot 8=\frac{32}{3}$, είναι λάθος.
\item Η διαίρεση με κλάσμα μικρότερο από το 1 \textbf{μεγαλώνει} τον αριθμό: $6:\frac{1}{2}=12$, γιατί στο 6 χωράνε 12 μισά. Η διαίρεση με φυσικό μεγαλύτερο από το 1 τον μικραίνει.
\item Δεν διαιρούμε με το 0, όπως πάντα. Ούτε με το κλάσμα $\frac{0}{5}$, που είναι 0.
\end{itemize}
\end{prosoxi}

\tmima{Λυμένα παραδείγματα}

\paradeigma{Τα δείγματα του παππού}
\tixero Ο παππούς Γιάννης έχει $\frac{3}{4}$ του λίτρου λάδι και το βάζει σε μπουκαλάκια των $\frac{1}{8}$ του λίτρου.\par
\tithelo Πόσα μπουκαλάκια γεμίζουν.\par
\possindeo Ζητάω πόσες φορές χωράει το $\frac{1}{8}$ στο $\frac{3}{4}$, δηλαδή $\frac{3}{4}:\frac{1}{8}$. Ο αντίστροφος του $\frac{1}{8}$ είναι το 8. Άρα $\frac{3}{4}\cdot 8=\frac{24}{4}=6$.\par
\apantisi 6 μπουκαλάκια.\par
\elegxos $6\cdot\frac{1}{8}=\frac{6}{8}=\frac{3}{4}$. Βγαίνει το λάδι που είχαμε.

\paradeigma{Μπουκαλάκια που δεν χωράνε ακριβώς}
\tixero Ο παππούς Γιάννης έχει $\frac{2}{3}$ του λίτρου και μπουκαλάκια των $\frac{1}{4}$ του λίτρου.\par
\tithelo Πόσα μπουκαλάκια θα γεμίσουν.\par
\possindeo $\frac{2}{3}:\frac{1}{4}=\frac{2}{3}\cdot 4=\frac{8}{3}=2\frac{2}{3}$.\par
\apantisi Γεμίζουν 2 ολόκληρα μπουκαλάκια και τα $\frac{2}{3}$ ενός τρίτου.\par
\elegxos $2\cdot\frac{1}{4}=\frac{1}{2}$ και μένουν $\frac{2}{3}-\frac{1}{2}=\frac{1}{6}$ του λίτρου. Το $\frac{1}{6}$ είναι τα $\frac{1}{6}:\frac{1}{4}=\frac{4}{6}=\frac{2}{3}$ του μπουκαλιού. Συμφωνεί.

\paradeigma{Σύνθετα κλάσματα}
\tixero Τα κλάσματα $\dfrac{\frac{5}{6}}{\frac{10}{9}}$ και $\dfrac{4}{\frac{2}{5}}$.\par
\tithelo Να τα κάνω απλά.\par
\possindeo Η κλασματική γραμμή σημαίνει διαίρεση. $\dfrac{\frac{5}{6}}{\frac{10}{9}}=\dfrac{5}{6}:\dfrac{10}{9}=\dfrac{5}{6}\cdot\dfrac{9}{10}=\dfrac{45}{60}=\dfrac{3}{4}$. Και $\dfrac{4}{\frac{2}{5}}=4:\dfrac{2}{5}=4\cdot\dfrac{5}{2}=10$.\par
\apantisi $\dfrac{3}{4}$ και $10$.\par
\elegxos Το $\frac{3}{4}\cdot\frac{10}{9}=\frac{30}{36}=\frac{5}{6}$ και το $10\cdot\frac{2}{5}=4$. Σωστά.

\paradeigma{Διαίρεση μεικτών}
\tixero Οι αριθμοί $3\frac{1}{3}$ και $1\frac{1}{9}$.\par
\tithelo Το πηλίκο $3\frac{1}{3}:1\frac{1}{9}$.\par
\possindeo $3\frac{1}{3}=\frac{10}{3}$ και $1\frac{1}{9}=\frac{10}{9}$. Άρα $\frac{10}{3}:\frac{10}{9}=\frac{10}{3}\cdot\frac{9}{10}=\frac{\cancel{10}\cdot 9}{3\cdot\cancel{10}}=3$.\par
\apantisi 3.\par
\elegxos $3\cdot 1\frac{1}{9}=3\cdot\frac{10}{9}=\frac{10}{3}=3\frac{1}{3}$. Σωστά.

\tmima{Ασκήσεις}

\zoni{Το κατάλαβα}{εφαρμογή του κανόνα}

\ex Κάνε τις διαιρέσεις: α) $\frac{3}{5}:\frac{1}{10}$ \quad β) $\frac{5}{6}:\frac{5}{12}$ \quad γ) $\frac{7}{8}:\frac{3}{4}$ \quad δ) $\frac{4}{9}:\frac{8}{27}$ \quad ε) $\frac{2}{3}:2$ \quad στ) $6:\frac{3}{4}$.

\ex Κάνε τις διαιρέσεις: α) $12:\frac{2}{3}$ \quad β) $5:\frac{1}{4}$ \quad γ) $\frac{3}{7}:3$ \quad δ) $10:\frac{5}{2}$ \quad ε) $\frac{1}{2}:\frac{1}{6}$ \quad στ) $\frac{9}{10}:\frac{3}{5}$.

\ex Πόσες φορές χωράει: α) το $\frac{1}{4}$ στο 1 \quad β) το $\frac{1}{3}$ στα $\frac{2}{3}$ \quad γ) το $\frac{1}{6}$ στα $\frac{5}{12}$ \quad δ) το $\frac{2}{5}$ στο 2.

\ex Υπολόγισε με μεικτούς: α) $2\frac{1}{2}:\frac{5}{6}$ \quad β) $3\frac{1}{3}:2$ \quad γ) $4:1\frac{1}{3}$ \quad δ) $1\frac{3}{4}:2\frac{1}{2}$.

\ex Μετάτρεψε σε απλά τα σύνθετα κλάσματα: α) $\dfrac{\frac{2}{3}}{\frac{4}{9}}$ \quad β) $\dfrac{5}{\frac{5}{6}}$ \quad γ) $\dfrac{\frac{7}{8}}{7}$ \quad δ) $\dfrac{\frac{3}{4}}{\frac{9}{10}}$.

\ex Ποιο κλάσμα πολλαπλασιασμένο με το $\frac{3}{5}$ δίνει $\frac{9}{10}$;

\ex Βρες τον αριθμό, αν: α) το $\frac{1}{3}$ του είναι $\frac{7}{9}$ \quad β) τα $\frac{2}{5}$ του είναι $\frac{6}{7}$.

\ex Υπολόγισε με προτεραιότητα πράξεων: α) $\frac{1}{2}:\frac{1}{4}\cdot\frac{1}{3}$ \quad β) $\frac{1}{2}:\left(\frac{1}{4}\cdot\frac{1}{3}\right)$ \quad γ) $\left(\frac{3}{4}+\frac{1}{4}\right):\frac{1}{2}$ \quad δ) $\frac{3}{4}+\frac{1}{4}:\frac{1}{2}$.

\ex Σωστό ή λάθος; α) $\frac{3}{4}:\frac{1}{2}=\frac{3}{8}$ \quad β) $6:\frac{1}{3}=18$ \quad γ) Η διαίρεση με κλάσμα μικρότερο από το 1 μεγαλώνει τον αριθμό. \quad δ) $\frac{2}{5}:\frac{2}{5}=1$ \quad ε) $\frac{1}{2}:3=6$ \quad στ) Για να διαιρέσω αντιστρέφω και τα δύο κλάσματα.

\zoni{Το χρησιμοποιώ}{προβλήματα από τη ζωή}

\ex \textbf{Το λάδι.} Ο παππούς Γιάννης έχει $\frac{5}{6}$ του λίτρου λάδι και μπουκαλάκια των $\frac{1}{12}$ του λίτρου. Πόσα γεμίζει; Αν είχε μπουκαλάκια των $\frac{1}{5}$ του λίτρου, πόσα θα γέμιζε και πόσο μέρος του τελευταίου θα έμενε γεμάτο;

\ex \textbf{Η κορδέλα.} Η Αγγελική έχει μια κορδέλα μήκους $\frac{9}{2}$ μ. και την κόβει σε κομμάτια μήκους $\frac{3}{4}$ μ. Πόσα κομμάτια βγάζει;

\ex \textbf{Η ταχύτητα.} Ο Νίκος τρέχει $2\frac{1}{2}$ χλμ. σε $\frac{1}{4}$ της ώρας (15 λεπτά). Πόσα χιλιόμετρα θα έτρεχε σε μια ώρα με την ίδια ταχύτητα;

\ex \textbf{Οι βάρδιες.} Η θεία Μαρία έχει $3\frac{1}{2}$ ώρες δουλειά στην αποθήκη και κάθε βάρδια διαρκεί $\frac{7}{4}$ της ώρας. Πόσες βάρδιες χρειάζεται;

\ex \textbf{Οι μερίδες.} Η γιαγιά Μυρτώ έχει 3 πίτες και κόβει κάθε μερίδα στα $\frac{3}{8}$ μιας πίτας. Πόσες μερίδες βγάζει;

\ex \textbf{Η συνταγή.} Για 6 άτομα η συνταγή θέλει $\frac{3}{4}$ του κιλού αλεύρι. Πόσο αλεύρι αντιστοιχεί σε ένα άτομο;

\zoni{Το σκέφτομαι}{μια ερώτηση που πάει πιο μακριά}

\ex Ο Στέλιος γράφει $\frac{3}{4}:\frac{1}{8}=\frac{4}{3}\cdot 8=\frac{32}{3}$. Βρες το λάθος του και το σωστό αποτέλεσμα.

\ex Τι βγαίνει αν διαιρέσεις ένα κλάσμα με τον αντίστροφό του; Δοκίμασε με τα $\frac{2}{3}$ και με τα $\frac{3}{5}$.

\ex Γιατί το $6:\frac{1}{2}=12$ είναι μεγαλύτερο από το 6; Εξήγησε με το «πόσα μισά χωράνε στο 6».

\ex Οι Αιγύπτιοι έγραφαν τα $\frac{2}{3}$ του $\frac{1}{7}$ ως $\frac{1}{14}+\frac{1}{42}$. Έλεγξε ότι αυτό ισούται με $\frac{2}{3}\cdot\frac{1}{7}$. Βρες με τον ίδιο τρόπο το $\frac{2}{3}$ του $\frac{1}{5}$ ως άθροισμα δύο κλασματικών μονάδων.

\ex Ποιος αριθμός πολλαπλασιασμένος επί $\frac{4}{5}$ κάνει 12; Ποιος αριθμός διαιρεμένος με το $\frac{4}{5}$ κάνει 12;

\zoni{Για το σπίτι}{με το μυαλό και με το μπουκάλι}

\ex Πόσα ποτήρια των $\frac{1}{4}$ του λίτρου γεμίζουν από ένα μπουκάλι $1\frac{1}{2}$ λίτρου; Έλεγξε με πραγματικό μπουκάλι και ποτήρι.

\ex Υπολόγισε το σύνθετο κλάσμα $\dfrac{\frac{1}{2}+\frac{1}{3}}{\frac{1}{2}-\frac{1}{3}}$.

\epanalipsi{2}

\tmima{Λεξικό του κεφαλαίου}

\renewcommand{\arraystretch}{1.5}
\begin{tabular*}{\linewidth}{@{\extracolsep{\fill}} >{\raggedright\arraybackslash}p{0.28\linewidth} >{\raggedright\arraybackslash}p{0.62\linewidth}}
\toprule
\textbf{Λέξη ή τύπος} & \textbf{Με μία πρόταση}\\
\midrule
$\dfrac{\kappa}{\nu}$ & $\kappa$ μέρη από $\nu$ ίσα μέρη, και ταυτόχρονα το $\kappa:\nu$ (με $\nu\neq 0$).\\
κλασματική μονάδα & Το $\dfrac{1}{\nu}$: ένα από τα $\nu$ ίσα μέρη.\\
ισοδύναμα & Ίδιο κομμάτι, άλλο κόψιμο: $\dfrac{\alpha}{\beta}=\dfrac{\gamma}{\delta}$ όταν $\alpha\delta=\beta\gamma$.\\
ανάγωγο & Δεν απλοποιείται άλλο: οι όροι είναι πρώτοι μεταξύ τους.\\
ομώνυμα, ετερώνυμα & Ίδιος ή διαφορετικός παρονομαστής. Για πράξεις και συγκρίσεις κάνω ετερώνυμα ομώνυμα με το ΕΚΠ.\\
μεικτός & Ολόκληρο και κλάσμα μαζί: $2\dfrac{3}{4}=2+\dfrac{3}{4}=\dfrac{11}{4}$.\\
αντίστροφος & Το κλάσμα γυρισμένο ανάποδα, με γινόμενο 1. Το 0 δεν έχει.\\
διαίρεση & Πολλαπλασιάζω με τον αντίστροφο του διαιρέτη.\\
σύνθετο κλάσμα & Κλάσμα με κλάσμα στους όρους του: η γραμμή σημαίνει διαίρεση.\\
\bottomrule
\end{tabular*}

\tmima{Ασκήσεις επανάληψης}

\ex Γράψε ως κλάσμα: α) τα πέντε από τα δώδεκα ίσα μέρη \quad β) το πηλίκο $7:9$ \quad γ) τον φυσικό 4 με παρονομαστή 6.

\ex Βρες τα $\frac{3}{8}$ του 240. Βρες και τον αριθμό που τα $\frac{5}{6}$ του είναι 75.

\ex Απλοποίησε μέχρι να γίνουν ανάγωγα: α) $\frac{36}{48}$ \quad β) $\frac{90}{150}$ \quad γ) $\frac{66}{99}$.

\ex Κάνε ομώνυμα και σύγκρινε: α) $\frac{5}{12}$ και $\frac{3}{8}$ \quad β) $\frac{7}{9}$ και $\frac{5}{6}$.

\ex Γράψε σε αύξουσα σειρά τα $\frac{1}{2}$, $\frac{5}{8}$, $\frac{3}{5}$ και $\frac{7}{10}$.

\ex Υπολόγισε: α) $\frac{3}{4}+\frac{5}{6}$ \quad β) $\frac{7}{8}-\frac{1}{3}$ \quad γ) $2\frac{1}{4}+1\frac{2}{3}$ \quad δ) $5-2\frac{3}{5}$.

\ex Υπολόγισε: α) $\frac{5}{6}\cdot\frac{9}{10}$ \quad β) $3\frac{1}{3}\cdot 2\frac{1}{4}$ \quad γ) $\frac{8}{9}:\frac{4}{3}$ \quad δ) $2\frac{1}{2}:1\frac{1}{4}$.

\ex Βρες την τιμή: α) $\left(\frac{1}{2}+\frac{1}{3}\right)\cdot\frac{6}{5}$ \quad β) $\frac{3}{4}-\frac{1}{2}:2$.

\ex Κάνε απλά τα σύνθετα κλάσματα: α) $\dfrac{\frac{3}{4}}{\frac{5}{8}}$ \quad β) $\dfrac{2+\frac{1}{2}}{1-\frac{1}{2}}$.

\ex Γράψε το $\frac{31}{6}$ ως μεικτό και τον $4\frac{3}{7}$ ως κλάσμα.

\ex \textbf{Το κέικ της Μυρτώς.} Η Μυρτώ κόβει το κέικ σε 12 ίσα κομμάτια. Η παρέα των εννέα παιδιών τρώει από ένα κομμάτι. Τι μέρος του κέικ έφαγαν και τι μέρος έμεινε; Η γιαγιά Ρήνη τρώει το μισό από αυτό που έμεινε. Τι μέρος ολόκληρου του κέικ είναι αυτό;

\ex \textbf{Οι δραστηριότητες.} Η Α3 έχει 27 μαθητές. Τα $\frac{4}{9}$ παίζουν μπάσκετ και το $\frac{1}{3}$ μαθαίνουν μουσική (κανείς δεν κάνει και τα δύο). Πόσα παιδιά κάνουν το καθένα και τι μέρος της τάξης δεν κάνει καμία από τις δύο;

\ex \textbf{Η προπόνηση.} Ο Νίκος έτρεξε το πρωί τα $\frac{2}{5}$ μιας διαδρομής και το απόγευμα τα $\frac{3}{10}$. Τι μέρος έτρεξε συνολικά και τι μέρος του μένει; Αν η διαδρομή είναι 20 χλμ., πόσα έτρεξε;

\ex \textbf{Το διάβασμα.} Ο Βασίλης χρειάζεται $\frac{3}{4}$ της ώρας για κάθε μάθημα. Πόσα μαθήματα προλαβαίνει σε 3 ώρες;

\ex \textbf{Ερωτήσεις αυτοαξιολόγησης (έννοια, ισοδύναμα).} Σωστό ή λάθος; α) Το $\frac{5}{5}$ είναι ίσο με 1. \quad β) Κάθε κλάσμα έχει αριθμητή μικρότερο από τον παρονομαστή. \quad γ) $\frac{4}{6}=\frac{2}{3}$. \quad δ) Το $\frac{3}{7}$ είναι ανάγωγο. \quad ε) Αν προσθέσω 2 και στον αριθμητή και στον παρονομαστή, το κλάσμα μένει ίδιο. \quad στ) $\frac{0}{9}=0$.

\ex \textbf{Ερωτήσεις αυτοαξιολόγησης (σύγκριση, πράξεις).} Σωστό ή λάθος; α) $\frac{3}{8}>\frac{3}{5}$. \quad β) $\frac{1}{2}+\frac{1}{2}=\frac{2}{4}$. \quad γ) $\frac{2}{3}\cdot\frac{3}{2}=1$. \quad δ) $6:\frac{1}{3}=2$. \quad ε) $1\frac{3}{4}=\frac{7}{4}$. \quad στ) Το άθροισμα δύο κλασμάτων είναι μεγαλύτερο από καθένα τους. \quad ζ) Το γινόμενο δύο κλασμάτων μικρότερων από το 1 είναι μικρότερο από καθένα τους.

\ex \textbf{Ερωτήσεις αυτοαξιολόγησης (αντίστροφοι, διαίρεση).} Σωστό ή λάθος; α) Ο αντίστροφος του $\frac{2}{5}$ είναι ο $\frac{5}{2}$. \quad β) Ο αντίστροφος του 0 είναι το 0. \quad γ) $\frac{1}{2}:\frac{1}{4}=2$. \quad δ) Η διαίρεση με κλάσμα μικρότερο από το 1 μικραίνει τον αριθμό. \quad ε) $3:\frac{3}{4}=4$. \quad στ) $\frac{5}{6}:\frac{5}{6}=1$.

\tmima{Ο χάρτης της επανάληψης}

Μετά τις λύσεις, δες σε ποιες ασκήσεις δυσκολεύτηκες. Ο πίνακας σου λέει ποια παράγραφο να ξαναδιαβάσεις.

\renewcommand{\arraystretch}{1.5}
\begin{tabular*}{\linewidth}{@{\extracolsep{\fill}} >{\raggedright\arraybackslash}p{0.38\linewidth} >{\raggedright\arraybackslash}p{0.52\linewidth}}
\toprule
\textbf{Αν δυσκολεύτηκες στις} & \textbf{ξαναδιάβασε τις παραγράφους}\\
\midrule
Ε2.1, Ε2.2, Ε2.12 & 2.1 (έννοια του κλάσματος)\\
Ε2.3 & 2.2 (ισοδύναμα κλάσματα)\\
Ε2.4, Ε2.5 & 2.2 και 2.3 (ομώνυμα, σύγκριση)\\
Ε2.6, Ε2.10, Ε2.13 & 2.4 (πρόσθεση, αφαίρεση, μεικτοί)\\
Ε2.7 (α, β), Ε2.8, Ε2.11 & 2.5 (πολλαπλασιασμός)\\
Ε2.7 (γ, δ), Ε2.9, Ε2.14 & 2.6 (διαίρεση)\\
Ε2.15, Ε2.16, Ε2.17 & 2.1 έως 2.6\\
\bottomrule
\end{tabular*}


\providecommand{\binpie}[2]{\begin{tikzpicture}[baseline=-0.4ex]%
  \fill[gblue!35] (0,0) -- (90:0.5) arc (90:{90-360*#2/#1}:0.5) -- cycle;%
  \foreach \i in {1,...,#1}{\draw[thin] (0,0) -- ({90-360*\i/#1}:0.5);}%
  \draw[thick] (0,0) circle (0.5);\end{tikzpicture}}
\kefalaio{Ρητοί αριθμοί}{Κάθε αριθμός έχει θέση, απόσταση και κατεύθυνση}{%
\begin{tabular}{@{}l@{\hspace{1em}}p{0.82\linewidth}@{}}
3.1 & Η έννοια του ρητού\\[2pt]
3.2 & Απόλυτη τιμή ρητού, σύγκριση και διάταξη ρητών, αντίθετοι ρητοί\\[2pt]
3.3 & Πρόσθεση και αφαίρεση ρητών\\[2pt]
3.4 & Πολλαπλασιασμός ρητών, ιδιότητες του πολλαπλασιασμού, αντίστροφοι ρητοί\\[2pt]
3.5 & Διαίρεση ρητών\\[2pt]
3.6 & Δυνάμεις ρητών με εκθέτη μη μηδενικό φυσικό\\[2pt]
3.7 & Τυποποιημένη μορφή μεγάλων αριθμών\\[2pt]
 & Επανάληψη του κεφαλαίου\\
\end{tabular}}
\paragrafos{Δεκαδικά κλάσματα και δεκαδικοί: ανάμεσα στους ακέραιους}{Α.3.1, «Δεκαδικά κλάσματα – Δεκαδικοί αριθμοί – Διάταξη – Στρογγυλοποίηση»}

\begin{anagke}
Αθλητική μέρα στο καινούργιο σχολείο, και για πρώτη φορά οι δύο γειτονιές τρέχουν στην ίδια πίστα. Η κυρία Αλεξάνδρα κρατάει το χρονόμετρο στα 60 μέτρα και η Πολυξένη γράφει τους χρόνους στο τετράδιο, γιατί στη Θεσσαλονίκη είχε τρέξει σε σύλλογο.

Ο Νίκος τερματίζει πρώτος και η καθηγήτρια φωνάζει: «Εννιά και ογδόντα!». Λίγο μετά τερματίζει ο Βασίλης: «Εννιά και εβδομήντα πέντε». Ο Βασίλης σκύβει το κεφάλι: «Το εβδομήντα πέντε είναι πολύ μεγαλύτερο από το οκτώ. Άρα έκανα πολύ περισσότερο χρόνο από σένα, έχασα με μεγάλη διαφορά». Ο Νίκος διαμαρτύρεται: «Μα έφτασες σχεδόν μαζί μου, με είχες πάνω από την πλάτη μου!». Η Πολυξένη γράφει τους δύο χρόνους, $9{,}8$ και $9{,}75$, και παρατηρεί ότι δεν είναι ίδιου μήκους: «Το ένα έχει ένα ψηφίο μετά την υποδιαστολή και το άλλο δύο. Πώς τα συγκρίνεις;»

Ο Στέλιος προσθέτει το δικό του πρόβλημα: «Με τι μονάδα μετράμε τέτοιους χρόνους; Το δευτερόλεπτο είναι πολύ χοντρή μονάδα. Αν ο χρόνος μου είναι ανάμεσα στο 9 και στο 10, δεν μπορώ να πω ούτε 9 ούτε 10». Η Μυρτώ το είχε πάθει και στο θρανίο: το πλάτος του ήταν ανάμεσα στα 62 και στα 63 εκατοστά.

Οι φυσικοί δεν φτάνουν. Χρειαζόμαστε αριθμούς \emph{ανάμεσα} στους φυσικούς, που να τους κόβουν σε δέκατα, εκατοστά, χιλιοστά, όσο λεπτά χρειάζεται η μέτρηση. Τέτοιους αριθμούς τους ξέρουμε ήδη από τα κλάσματα, και τώρα θα τους γράψουμε με τρόπο που συγκρίνεται εύκολα.
\end{anagke}

\begin{lexi}{δεκαδικός, υποδιαστολή}{δεκαδικός, υποδιαστολή}
Η λέξη \lexiword{δεκαδικός} έρχεται από το «δέκα»: ο αριθμός είναι γραμμένος με βάση το 10, και μετά την υποδιαστολή συνεχίζονται οι θέσεις με δέκατα, εκατοστά, χιλιοστά. Η \lexiword{υποδιαστολή} είναι το κόμμα που χωρίζει τον ακέραιο από το δεκαδικό μέρος· «υπό-διαστέλλω» σημαίνει «χωρίζω από κάτω», και παλιότερα ήταν στο ύψος της βάσης των ψηφίων. Το \lexiword{δεκαδικό κλάσμα} είναι το κλάσμα που έχει παρονομαστή δύναμη του 10: $10,\,100,\,1\,000,\,10\,000,\dots$.
\end{lexi}

\begin{metafrasi}
Τρία δέκατα. & $\dfrac{3}{10}=0{,}3$\\
Εκατόν τριάντα πέντε εκατοστά. & $\dfrac{135}{100}=1{,}35$\\
Δύο ακέραια και επτά χιλιοστά. & $2{,}007=2+\dfrac{7}{1\,000}$\\
Ένα κλάσμα $\dfrac{\kappa}{\nu}$ γίνεται δεκαδικός. & $\kappa:\nu$\\
Οι ίδιοι χρόνοι, με τα ίδια ψηφία. & $9{,}8=9{,}80$\\
Ποιος δεκαδικός είναι μεγαλύτερος; & συγκρίνω ακέραιο, δέκατα, εκατοστά, …\\
Προσεγγίζω στα εκατοστά. & $9{,}7468\approx 9{,}75$\\
\end{metafrasi}

\tmima{Η θεωρία}

\tmimasmall{Δεκαδικά κλάσματα και δεκαδικοί αριθμοί}

Τα κλάσματα $\frac{3}{10}$, $\frac{825}{100}$, $\frac{53}{1\,000}$ και $\frac{1004}{10\,000}$ έχουν παρονομαστές τους $10$, $10^2$, $10^3$ και $10^4$. Είναι δεκαδικά κλάσματα.

\begin{orismos}{δεκαδικό κλάσμα και δεκαδικός αριθμός}
Δεκαδικό κλάσμα λέγεται κάθε κλάσμα με παρονομαστή δύναμη του 10. Κάθε δεκαδικό κλάσμα γράφεται ως \textbf{δεκαδικός αριθμός} με τόσα δεκαδικά ψηφία όσα μηδενικά έχει ο παρονομαστής του:
$$\dfrac{3}{10}=0{,}3,\qquad \dfrac{825}{100}=8{,}25,\qquad \dfrac{53}{1\,000}=0{,}053,\qquad \dfrac{1004}{10\,000}=0{,}1004.$$
\end{orismos}

Για να πάω από το δεκαδικό στο κλάσμα, κάνω το αντίστροφο: γράφω στον αριθμητή όλα τα ψηφία χωρίς την υποδιαστολή και στον παρονομαστή το $1$ με τόσα μηδενικά όσα τα δεκαδικά ψηφία. Άρα $3{,}25=\frac{325}{100}$.

Για να μετατρέψω ένα \emph{οποιοδήποτε} κλάσμα $\frac{\kappa}{\nu}$ σε δεκαδικό, κάνω τη διαίρεση $\kappa:\nu$, βάζοντας υποδιαστολή και μηδενικά στον διαιρετέο όταν τελειώσει το ακέραιο μέρος. Αν η διαίρεση τελειώνει, ο δεκαδικός είναι ακριβής, π.χ. $\frac{3}{8}=0{,}375$. Αν δεν τελειώνει ποτέ, όπως στο $\frac{1}{3}=0{,}333\ldots$, τον γράφουμε κατά προσέγγιση (θα το ξαναδούμε στο κεφάλαιο των ρητών).

\tmimasmall{Οι τάξεις των δεκαδικών ψηφίων}

Δεξιά από την υποδιαστολή οι θέσεις μικραίνουν κατά δέκα, όπως αριστερά μεγαλώνουν:

\begin{center}\small
\renewcommand{\arraystretch}{1.4}\setlength{\tabcolsep}{3pt}
\begin{tabular}{@{}cccccccc@{}}
\toprule
εκατοντάδες & δεκάδες & μονάδες & \textbf{,} & δέκατα & εκατοστά & χιλιοστά & δεκάκις χιλιοστά\\
$100$ & $10$ & $1$ & & $\frac{1}{10}$ & $\frac{1}{100}$ & $\frac{1}{1\,000}$ & $\frac{1}{10\,000}$\\
\bottomrule
\end{tabular}
\end{center}

Το επόμενο ψηφίο λέγεται εκατοντάκις χιλιοστό ($\frac{1}{100\,000}$), μετά εκατομμυριοστό, και ούτω καθεξής. Ο $305{,}0408$ έχει δεκαδικό μέρος $0{,}0408$ και διαβάζεται «τριακόσια πέντε ακέραια και τετρακόσια οκτώ δεκάκις χιλιοστά», γιατί το δεκαδικό μέρος $0{,}0408=\frac{408}{10\,000}$. Μπορούμε να τον αναλύσουμε:
$$305{,}0408=3\cdot100+0\cdot10+5\cdot1+\dfrac{0}{10}+\dfrac{4}{100}+\dfrac{0}{1\,000}+\dfrac{8}{10\,000}.$$

Αν προσθέσω ή σβήσω μηδενικά στο τέλος του δεκαδικού μέρους, ο αριθμός δεν αλλάζει: $9{,}8=9{,}80=9{,}800$, γιατί $\frac{8}{10}=\frac{80}{100}=\frac{800}{1\,000}$. Δεν αλλάζει όμως αν βάλω μηδενικό \emph{πριν} από τα ψηφία του δεκαδικού μέρους: $0{,}5\neq 0{,}05$.

\tmimasmall{Μια ιστορική σημείωση: ο Στέβιν και ο Νάπιερ}

Οι Βαβυλώνιοι και οι αρχαίοι Έλληνες αστρονόμοι χρησιμοποιούσαν κλάσματα με παρονομαστή 60, και οι έμποροι έβρισκαν δύσκολα τους υπολογισμούς με κλάσματα. Ο Φλαμανδός μαθηματικός \textbf{Σιμόν Στέβιν} (1548--1620) κατάλαβε ότι οι υπολογισμοί γίνονται πολύ πιο εύκολα με κλάσματα που έχουν παρονομαστή δύναμη του 10. Το 1585 έγραψε ένα μικρό βιβλίο, τη «Δεκάτη» (\emph{La disme}), για τους τοκογλύφους, τους τοπογράφους και τους εμπόρους, όπου εισήγαγε τον δεκαδικό υπολογισμό. Ο Στέβιν έγραφε τον $34{,}25$ ως $34_{(0)}\,2_{(1)}\,5_{(2)}$: ο δείκτης έδειχνε πόσες θέσεις μετά τη μονάδα βρίσκεται το κάθε ψηφίο. Το κόμμα (ή τελεία) ανάμεσα στον ακέραιο και στο δεκαδικό μέρος το πρότεινε αργότερα ο Σκοτσέζος \textbf{Τζον Νάπιερ} (1550--1617), και η απλότητά του επικράτησε.

\tmimasmall{Σύγκριση δεκαδικών}

\begin{kanonas}{σύγκριση δεκαδικών}
\begin{enumerate}[leftmargin=1.6em,itemsep=0.3ex,topsep=2pt]
\item Συγκρίνω πρώτα τα \textbf{ακέραια μέρη}. Μεγαλύτερο ακέραιο μέρος σημαίνει μεγαλύτερο αριθμό.
\item Αν τα ακέραια μέρη είναι ίσα, συγκρίνω τα δέκατα, μετά τα εκατοστά, μετά τα χιλιοστά, και σταματάω στο πρώτο ψηφίο που διαφέρει.
\item Για να μην μπερδεύομαι, γράφω στον αριθμό με τα λιγότερα δεκαδικά ψηφία μηδενικά στο τέλος, ώστε να έχουν το ίδιο πλήθος.
\end{enumerate}
\end{kanonas}

Έτσι ο χρόνος του Νίκου $9{,}8=9{,}80$ και του Βασίλη $9{,}75$: τα ακέραια ίσα, τα δέκατα $8>7$, άρα $9{,}80>9{,}75$. Ο Βασίλης συνέκρινε τα δεκαδικά μέρη $75$ και $8$ σαν να ήταν φυσικοί αριθμοί. Στην πραγματικότητα έκανε \emph{μικρότερο} χρόνο και κέρδισε, με διαφορά πέντε εκατοστά του δευτερολέπτου.

\tmimasmall{Στρογγυλοποίηση δεκαδικών}

Όταν δεν χρειαζόμαστε όλα τα ψηφία, τα κόβουμε και κρατάμε μια προσέγγιση. Για να στρογγυλοποιήσω σε μια δεδομένη τάξη (π.χ. στα εκατοστά):
\begin{itemize}[leftmargin=1.2em,itemsep=0.3ex,topsep=2pt]
\item κοιτάζω το \textbf{επόμενο} ψηφίο (στα χιλιοστά, αν στρογγυλοποιώ στα εκατοστά)·
\item αν είναι $0,1,2,3,4$ αφήνω το ψηφίο της τάξης όπως είναι και κόβω τα υπόλοιπα·
\item αν είναι $5,6,7,8,9$ αυξάνω το ψηφίο της τάξης κατά $1$ και κόβω τα υπόλοιπα.
\end{itemize}

\tmimasmall{Οι δεκαδικοί στην ευθεία}

Για να βάλω δεκαδικούς στην ευθεία των αριθμών, χωρίζω κάθε μονάδα σε $10$ ίσα μέρη (δέκατα) και, αν χρειάζεται, κάθε δέκατο σε $10$ ίσα μέρη (εκατοστά), και πάει λέγοντας. Δεν υπάρχει όριο: ανάμεσα σε δύο οποιουσδήποτε διαφορετικούς δεκαδικούς υπάρχει πάντα άλλος δεκαδικός, π.χ. ο μέσος τους, και άρα υπάρχουν άπειροι.

\begin{center}
\begin{tikzpicture}[x=1cm,y=1cm]
\draw[thick,-{Stealth}] (-0.3,0)--(10.7,0);
\foreach \i in {0,...,10}{\draw (\i,-0.12)--(\i,0.12);}
\node[below] at (0,-0.15) {$2$};
\node[below] at (5,-0.15) {$2{,}5$};
\node[below] at (10,-0.15) {$3$};
\foreach \i/\t in {1/2{,}1,3/2{,}3,8/2{,}8}{\node[below,font=\small] at (\i,-0.15) {$\t$};}
\fill[gblue] (7.5,0) circle (2.5pt);
\node[above,gblue] at (7.5,0.2) {$2{,}75$};
\fill[gblue] (3,0) circle (2.5pt);
\node[above,gblue] at (3,0.2) {$2{,}3$};
\end{tikzpicture}
\end{center}

Το $2{,}75$ βρίσκεται στη μέση του διαστήματος από το $2{,}7$ ως το $2{,}8$, που αντιστοιχεί στο σημείο $7{,}5$ των δέκατων.

\begin{prosoxi}
\begin{itemize}[leftmargin=1.2em,itemsep=0.3ex,topsep=2pt]
\item Δεν συγκρίνουμε τα δεκαδικά μέρη \emph{σαν φυσικούς}. Το $0{,}9$ είναι μεγαλύτερο από το $0{,}12$, γιατί $0{,}90>0{,}12$, αν και το $9<12$.
\item Το $0{,}50$ ισούται με το $0{,}5$, γιατί το μηδενικό στο τέλος δεν αλλάζει τίποτα, αλλά το $0{,}05$ είναι δέκα φορές μικρότερο.
\item Στρογγυλοποιούμε \emph{μια φορά}, από τον αρχικό αριθμό. Αν στρογγυλοποιήσουμε πρώτα στα εκατοστά και μετά στα δέκατα, μπορεί να βγει λάθος: $2{,}449\approx2{,}4$, αλλά $2{,}449\to2{,}45\to2{,}5$.
\item Το $\frac{2}{5}$ δεν έχει παρονομαστή δύναμη του 10, δεν είναι δηλαδή δεκαδικό κλάσμα, αλλά είναι ισοδύναμο με το δεκαδικό κλάσμα $\frac{4}{10}=0{,}4$.
\end{itemize}
\end{prosoxi}

\tmima{Λυμένα παραδείγματα}

\paradeigma{Ποιος κέρδισε;}
\tixero Οι χρόνοι στα 60 μ. του Νίκου και του Βασίλη είναι $9{,}8$ s και $9{,}75$ s.\par
\tithelo Ποιος έκανε τον μικρότερο χρόνο.\par
\possindeo Ισοσκελίζω τα δεκαδικά ψηφία: $9{,}8=9{,}80$. Τα ακέραια είναι ίσα (9) και στα δέκατα $8>7$, άρα $9{,}80>9{,}75$.\par
\apantisi Ο Βασίλης έκανε μικρότερο χρόνο ($9{,}75$ s) και κέρδισε με διαφορά πέντε εκατοστά του δευτερολέπτου.\par
\elegxos Σε εκατοστά, $980$ εκατοστά και $975$ εκατοστά: το $975$ είναι μικρότερο. Συμφωνεί.

\paradeigma{Από κλάσμα σε δεκαδικό και ανάποδα}
\tixero Τα κλάσματα $\frac{3}{8}$ και $\frac{7}{20}$ και ο δεκαδικός $0{,}625$.\par
\tithelo Να τα γράψω στην άλλη μορφή.\par
\possindeo Για το $\frac{3}{8}$ κάνω τη διαίρεση $3:8$: $3=3{,}000$ και $3000:8=375$, άρα $0{,}375$. Το $\frac{7}{20}$ γίνεται δεκαδικό κλάσμα αν πολλαπλασιάσω αριθμητή και παρονομαστή με $5$: $\frac{35}{100}=0{,}35$. Και ο $0{,}625$ έχει τρία δεκαδικά ψηφία, άρα $\frac{625}{1\,000}$, που απλοποιείται με το $125$ σε $\frac{5}{8}$.\par
\apantisi $\frac{3}{8}=0{,}375$, $\frac{7}{20}=0{,}35$ και $0{,}625=\frac{5}{8}$.\par
\elegxos $0{,}375\cdot 8=3$ και $\frac{5}{8}:1=0{,}625$ (αφού $5:8=0{,}625$). Σωστά.

\paradeigma{Διαβάζω και αναλύω}
\tixero Ο αριθμός $305{,}0408$.\par
\tithelo Να πω τι ψηφίο βρίσκεται σε κάθε θέση και να τον διαβάσω.\par
\possindeo Πριν από την υποδιαστολή: $3$ εκατοντάδες, $0$ δεκάδες, $5$ μονάδες. Μετά: $0$ δέκατα, $4$ εκατοστά, $0$ χιλιοστά, $8$ δεκάκις χιλιοστά. Το δεκαδικό μέρος είναι $0{,}0408=\frac{408}{10\,000}$.\par
\apantisi «Τριακόσια πέντε ακέραια και τετρακόσια οκτώ δεκάκις χιλιοστά».\par
\elegxos $\frac{4}{100}+\frac{8}{10\,000}=\frac{400}{10\,000}+\frac{8}{10\,000}=\frac{408}{10\,000}$. Σωστά.

\paradeigma{Στρογγυλοποιώ σε τάξεις}
\tixero Ο αριθμός $16{,}4963$.\par
\tithelo Να τον στρογγυλοποιήσω στην μονάδα, στα δέκατα, στα εκατοστά και στα χιλιοστά.\par
\possindeo Κοιτάζω κάθε φορά το επόμενο ψηφίο. Μονάδα: επόμενο το $4$ (<5) άρα $16$. Δέκατα: επόμενο το $9$ (≥5) άρα $4+1=5$, δηλαδή $16{,}5$. Εκατοστά: επόμενο το $6$ άρα $9+1=10$, που δίνει $16{,}50$. Χιλιοστά: επόμενο το $3$ άρα $16{,}496$.\par
\apantisi $16$, $16{,}5$, $16{,}50$, $16{,}496$.\par
\elegxos Το σφάλμα στα εκατοστά είναι $16{,}50-16{,}4963=0{,}0037<0{,}005$, δηλαδή μικρότερο από μισό εκατοστό. Σωστό. Το μηδενικό στο $16{,}50$ μένει για να δείχνει ότι μετρήσαμε μέχρι τα εκατοστά.

\paradeigma{Ένας δεκαδικός ανάμεσα}
\tixero Οι αριθμοί $3{,}14$ και $3{,}15$.\par
\tithelo Να βρω έναν δεκαδικό ανάμεσά τους.\par
\possindeo Γράφω τα δύο με τρία δεκαδικά: $3{,}140$ και $3{,}150$. Ανάμεσα είναι όλοι οι $3{,}141,\,3{,}142,\dots,3{,}149$.\par
\apantisi Για παράδειγμα ο $3{,}145$, που είναι ο μέσος τους.\par
\elegxos $3{,}14<3{,}145<3{,}15$. Και αφού μπορώ να ξαναχωρίσω το διάστημα, υπάρχουν άπειροι.

\tmima{Ασκήσεις}

\zoni{Το κατάλαβα}{δεκαδικά κλάσματα, ανάγνωση, σύγκριση}

\ex Γράψε ως δεκαδικούς: α) $\frac{7}{10}$ \quad β) $\frac{43}{100}$ \quad γ) $\frac{9}{1\,000}$ \quad δ) $\frac{2051}{1\,000}$ \quad ε) $\frac{38}{10\,000}$ \quad στ) $\frac{5}{100\,000}$.

\ex Γράψε ως δεκαδικά κλάσματα (με έναν αριθμητή και έναν παρονομαστή): α) $0{,}6$ \quad β) $0{,}07$ \quad γ) $3{,}25$ \quad δ) $0{,}0019$ \quad ε) $12{,}5$ \quad στ) $0{,}400$.

\ex Μετάτρεψε σε δεκαδικούς με διαίρεση ή με ισοδύναμο κλάσμα: α) $\frac{3}{4}$ \quad β) $\frac{7}{8}$ \quad γ) $\frac{9}{20}$ \quad δ) $\frac{11}{25}$ \quad ε) $\frac{5}{16}$ \quad στ) $\frac{3}{5}$.

\ex Γράψε με ψηφία: α) τρία ακέραια και είκοσι πέντε χιλιοστά \quad β) τετρακόσια εβδομήντα δύο εκατοστά \quad γ) δύο δέκατα \quad δ) δεκαπέντε δεκάκις χιλιοστά \quad ε) έξι ακέραια, πέντε εκατοστά και ένα εκατοντάκις χιλιοστό.

\ex Ποιο ψηφίο βρίσκεται στα εκατοστά και ποιο στα δεκάκις χιλιοστά του $24{,}0713$; Πώς διαβάζεται ο αριθμός;

\ex Ανάλυσε τους αριθμούς όπως στη θεωρία: α) $4{,}073$ \quad β) $12{,}408$ \quad γ) $0{,}0905$.

\ex Σωστό ή λάθος; α) $0{,}5=0{,}50$ \quad β) $0{,}5=0{,}05$ \quad γ) $3{,}40=3{,}4$ \quad δ) $0{,}0070=0{,}007$ \quad ε) Ο $12{,}5$ διαβάζεται «δώδεκα και πέντε» \quad στ) $\frac{1}{4}=0{,}4$ \quad ζ) $\frac{25}{100}=0{,}25$.

\ex Σύγκρινε τα ζευγάρια με τα σύμβολα $<,\,=,\,>$: α) $3{,}07$ και $3{,}7$ \quad β) $0{,}58$ και $0{,}6$ \quad γ) $12{,}09$ και $12{,}090$ \quad δ) $5{,}999$ και $6{,}001$ \quad ε) $0{,}3$ και $0{,}29$ \quad στ) $7{,}45$ και $7{,}5$.

\ex Στρογγυλοποίησε τους $7{,}486$ και $0{,}0951$: α) στη μονάδα \quad β) στα δέκατα \quad γ) στα εκατοστά.

\zoni{Το χρησιμοποιώ}{ο στίβος, η ζυγαριά, η πινακίδα}

\ex \textbf{Ο στίβος.} Στα 60 μ. οι χρόνοι ήταν: Νίκος $9{,}38$ s, Βασίλης $9{,}4$ s, Στέλιος $9{,}29$ s, Πολυξένη $9{,}35$ s, Αγγελική $9{,}3$ s, Ηλίας $9{,}41$ s. Γράψε τους μαθητές με τη σειρά που τερμάτισαν, από τον πρώτο.

\ex \textbf{Το άλμα.} Στο άλμα εις μήκος η Ζωή πήδηξε $3{,}08$ μ., η Μυρτώ $2{,}98$ μ., η Φανή $3{,}1$ μ. και η Πολυξένη $3{,}085$ μ. Γράψε τα μήκη από το μεγαλύτερο στο μικρότερο και πες ποια παιδιά πήραν το χρυσό και το ασημένιο μετάλλιο.

\ex \textbf{Η ζυγαριά.} Η Αγγελική ζυγίζει τα μήλα του παππού και βλέπει $2{,}35$ κιλά. Πόσα είναι στρογγυλοποιημένα στο κιλό και πόσα στα δέκατα του κιλού;

\ex \textbf{Η πινακίδα.} Στο βενζινάδικο η πινακίδα γράφει $1{,}879$ € το λίτρο. Ο θείος Λευτέρης λέει ότι «κοντά στα $1{,}9$», και η Πολυξένη ότι «κοντά στα $1{,}88$». Ποια στρογγυλοποίηση είναι ποια, και ποια είναι πιο κοντινή στην πραγματική τιμή;

\ex \textbf{Ανάμεσα.} Βρες τρεις διαφορετικούς δεκαδικούς ανάμεσα στους $2{,}4$ και $2{,}5$. Πόσοι ακόμη υπάρχουν;

\ex \textbf{Στην ευθεία.} Σχεδίασε ένα ευθύγραμμο τμήμα με μήκος $10$ εκ., που θα παριστάνει την απόσταση από το $0$ ως το $1$, και σημείωσε τους $0{,}3$, $0{,}75$, $0{,}9$ και $\frac{1}{2}$.

\ex \textbf{Τελειώνει ή όχι;} Κάνε τις διαιρέσεις (μέχρι να τελειώσουν ή να δεις ότι το ίδιο ψηφίο επαναλαμβάνεται): α) $\frac{3}{20}$ \quad β) $\frac{7}{30}$ \quad γ) $\frac{11}{50}$ \quad δ) $\frac{5}{12}$. Ποια τελειώνουν και ποια όχι;

\zoni{Το σκέφτομαι}{όταν το κόμμα ξεγελάει}

\ex Ο Βασίλης ισχυρίζεται ότι $0{,}9<0{,}12$, γιατί το $9$ είναι μικρότερο από το $12$. Εξήγησε το λάθος του και βρες ένα παράδειγμα με λεφτά ($0{,}90$ € και $0{,}12$ €) που το δείχνει.

\ex Πόσα δεκαδικά ψηφία έχει ο $\frac{7}{16}$; Πώς θα το προέβλεπες χωρίς να κάνεις τη διαίρεση, αν σκεφτείς ότι ο παρονομαστής $16$ είναι $2^4$ και ότι το $10^4=2^4\cdot5^4$;

\ex Με τα ψηφία $3$, $0$ και $8$ (το καθένα μία φορά στο δεκαδικό μέρος) και ακέραιο μέρος το $1$, γράψε τον μικρότερο και τον μεγαλύτερο δεκαδικό.

\ex α) Πόσοι δεκαδικοί, με το πολύ δύο δεκαδικά ψηφία, υπάρχουν ανάμεσα στο $2$ και στο $3$ (χωρίς αυτά); β) Πόσοι από αυτούς έχουν ακριβώς δύο δεκαδικά ψηφία (το τελευταίο όχι μηδέν);

\ex Ο Ηλίας στρογγυλοποιεί τον $2{,}449$ στα δέκατα με δύο βήματα: πρώτα στα εκατοστά ($2{,}45$) και μετά στα δέκατα ($2{,}5$). Ο Νίκος τον στρογγυλοποιεί κατευθείαν. Ποιος έχει δίκιο;

\ex Ο Στέβιν έγραφε τον $7{,}206$ ως $7_{(0)}\,2_{(1)}\,0_{(2)}\,6_{(3)}$. Πώς θα έγραφε τον $18{,}053$; Ποιος αριθμός είναι ο $8_{(0)}\,0_{(1)}\,4_{(2)}\,9_{(3)}$;

\zoni{Για το σπίτι}{μετρώ και καταγράφω}

\ex Μέτρησε το πλάτος του θρανίου σου με ένα μολύβι, με ένα χάρακα και με μετροταινία. Γράψε τα αποτελέσματα και εξήγησε ποια μέτρηση είναι η πιο ακριβής και γιατί χρειάζεται δεκαδικός αριθμός.

\ex Στο σούπερ μάρκετ βρες τρία προϊόντα που η τιμή ανά κιλό ή ανά λίτρο έχει τρία δεκαδικά ψηφία. Στρογγυλοποίησέ τις στο λεπτό του ευρώ και βάλε τες σε σειρά από την πιο ακριβή.

\paragrafos{Απόλυτη τιμή και διάταξη: ποιος είναι δεξιότερα;}{Α.3.2, «Απόλυτη τιμή ρητού, σύγκριση και διάταξη ρητών, αντίθετοι ρητοί»}

\begin{anagke}
Στην ώρα των μαθηματικών η κυρία Αλεξάνδρα ρωτάει: «Ποιος είναι ο επόμενος ρητός μετά το $\frac{2}{5}$;». Ο Βασίλης απαντάει με σιγουριά: «Το $\frac{3}{5}$». Η Ρήνη σηκώνει το χέρι: «Στον χάρακα το $\frac{2}{5}$ είναι στο $0{,}4$ και το $\frac{3}{5}$ στο $0{,}6$. Ανάμεσα υπάρχει το $0{,}5$. Και το $0{,}45$. Και το $0{,}41$!». Ο Νίκος ρωτάει: «Αν δεν υπάρχει επόμενος, πώς τους βάζουμε σε σειρά;».

Λίγο αργότερα ο Ηλίας διαβάζει την πρόγνωση: Ιωάννινα $-3^\circ$, Φλώρινα $-10^\circ$. «Το δέκα είναι μεγαλύτερο από το τρία, άρα η Φλώρινα έχει μεγαλύτερη θερμοκρασία», λέει ο Βασίλης. Η Ζωή διαφωνεί: «Δεν μπορεί να είναι πιο ζεστό το δέκα κάτω από το μηδέν!». Η Αγγελική ζωγραφίζει τον άξονα: «Όσο πιο αριστερά, τόσο πιο κρύο. Το δέκα μετράει πόσο μακριά είμαστε από το μηδέν, και το πρόσημο λέει από ποια μεριά».

Η Πολυξένη το διατυπώνει με δύο ερωτήματα: «Πόσο μακριά είναι κάθε αριθμός από το μηδέν; Και ποιος είναι πιο δεξιά;» Η πρώτη ιδέα θα λέγεται απόλυτη τιμή και η δεύτερη διάταξη. Σε αυτή την παράγραφο θα μάθουμε να συγκρίνουμε οποιουσδήποτε δύο ρητούς: κλάσματα, δεκαδικούς, θετικούς και αρνητικούς.
\end{anagke}

\begin{lexi}{απόλυτη τιμή, αντίθετοι, διάταξη}{απόλυτη τιμή, αντίθετοι, διάταξη}
\lexiword{Απόλυτη} σημαίνει «λυμένη από δεσμούς» (από το «απολύω»): η απόλυτη τιμή είναι ο αριθμός ελεύθερος από το πρόσημό του, δηλαδή μόνο η απόσταση από το $0$. Οι \lexiword{αντίθετοι} βρίσκονται σε «αντίθετη» κατεύθυνση: ίδια απόσταση από το $0$, ο ένας δεξιά και ο άλλος αριστερά. Η \lexiword{διάταξη} είναι η «διά-ταξη», το να βάζεις τους αριθμούς «σε τάξη», δηλαδή σε σειρά.
\end{lexi}

\begin{metafrasi}
Η απόσταση του $-2$ από το $0$ είναι $2$. & $|-2|=2$\\
Η απόσταση του $+\frac{1}{2}$ από το $0$ είναι $\frac{1}{2}$. & $\left|+\frac{1}{2}\right|=\frac{1}{2}$\\
Ο αντίθετος του $+2{,}8$ είναι ο $-2{,}8$. & $-(+2{,}8)=-2{,}8$\\
Το $-3{,}7$ είναι αριστερά του $-2{,}5$. & $-3{,}7<-2{,}5$\\
Το $1{,}6$ είναι δεξιά του $0{,}5$. & $1{,}6>0{,}5$\\
\end{metafrasi}

\tmima{Η θεωρία}

\tmimasmall{Απόλυτη τιμή και αντίθετοι}

Ο ορισμός της απόλυτης τιμής που ξέρουμε από τους ακεραίους ισχύει και για τους ρητούς.

\begin{orismos}{απόλυτη τιμή ρητού}
\textbf{Απόλυτη τιμή} ενός ρητού $\alpha$ είναι η \textbf{απόσταση} του σημείου με τετμημένη $\alpha$ από την αρχή $\mathrm{Ο}$ του άξονα. Συμβολίζεται $|\alpha|$.
\end{orismos}

Για παράδειγμα $|-2|=2$ και $\left|+\frac{1}{2}\right|=\frac{1}{2}$. Η απόλυτη τιμή ενός θετικού είναι ο ίδιος ο αριθμός, ενός αρνητικού ο αντίθετός του και $|0|=0$. Δεν είναι ποτέ αρνητική.

\begin{center}
\akgr[1.8]{-2}{2}{-2/-2,0.5/\frac{1}{2}}
\end{center}

\begin{orismos}{αντίθετοι ρητοί}
\textbf{Αντίθετοι} λέγονται δύο ρητοί που είναι \textbf{ετερόσημοι} και έχουν την \textbf{ίδια απόλυτη τιμή}. Ο αντίθετος του $\alpha$ είναι ο $-\alpha$. Ο αντίθετος του $0$ είναι το $0$ και ισχύει $-(-\alpha)=\alpha$.
\end{orismos}

Οι $+2{,}8$ και $-2{,}8$ είναι αντίθετοι. Το ίδιο και οι $-\frac{1}{7}$ και $+\frac{1}{7}$. Στον άξονα, οι αντίθετοι είναι σημεία συμμετρικά ως προς το $0$.

\tmimasmall{Σύγκριση και διάταξη ρητών}

Από δύο ρητούς \textbf{μεγαλύτερος} είναι αυτός που βρίσκεται \textbf{δεξιότερα} στον άξονα.

\begin{center}
\akgr[1.8]{-4}{2}{-3.7/-3{,}7,-2.5/-2{,}5,0.5/0{,}5,1.6/1{,}6}
\end{center}

Στο σχήμα $1{,}6>0{,}5$ και $-2{,}5>-3{,}7$.

\begin{kanonas}{σύγκριση ρητών με απόλυτες τιμές}
\begin{itemize}[leftmargin=1.2em,itemsep=0.3ex,topsep=2pt]
\item Κάθε θετικός είναι μεγαλύτερος από κάθε αρνητικό. Το $0$ είναι μεγαλύτερο από κάθε αρνητικό και μικρότερο από κάθε θετικό.
\item Από δύο \textbf{θετικούς} μεγαλύτερος είναι αυτός που έχει τη \textbf{μεγαλύτερη} απόλυτη τιμή: $+4{,}3>+4{,}2$, γιατί $|{+4{,}3}|>|{+4{,}2}|$.
\item Από δύο \textbf{αρνητικούς} μεγαλύτερος είναι αυτός που έχει τη \textbf{μικρότερη} απόλυτη τιμή: $-0{,}1>-0{,}2$, γιατί $|{-0{,}1}|=0{,}1<0{,}2=|{-0{,}2}|$.
\end{itemize}
\end{kanonas}

\tmimasmall{Σύγκριση κλασμάτων}

Για κλάσματα με \textbf{θετικούς} όρους ισχύουν τα εξής.

\begin{kanonas}{σύγκριση κλασμάτων}
\begin{itemize}[leftmargin=1.2em,itemsep=0.3ex,topsep=2pt]
\item Από δύο \textbf{ομώνυμα} κλάσματα μεγαλύτερο είναι αυτό με τον \textbf{μεγαλύτερο αριθμητή}: $\frac{2}{5}<\frac{3}{5}$.
\item Από δύο κλάσματα με τον \textbf{ίδιο αριθμητή} μεγαλύτερο είναι αυτό με τον \textbf{μικρότερο παρονομαστή}: $\frac{1}{4}<\frac{1}{2}$.
\item Για να συγκρίνουμε \textbf{ετερώνυμα} κλάσματα, τα κάνουμε \textbf{ομώνυμα} και συγκρίνουμε τους αριθμητές τους. Διαφορετικά, εκτελούμε τις διαιρέσεις και συγκρίνουμε τους δεκαδικούς.
\end{itemize}
\end{kanonas}

\begin{center}
\(\dfrac{2}{5}\)\ \binpie{5}{2}\ \ $<$\ \ \(\dfrac{3}{5}\)\ \binpie{5}{3}\qquad\qquad\(\dfrac{1}{4}\)\ \binpie{4}{1}\ \ $<$\ \ \(\dfrac{1}{2}\)\ \binpie{2}{1}
\end{center}

\tmimasmall{Σύγκριση δεκαδικών}

\begin{kanonas}{σύγκριση δεκαδικών}
\begin{itemize}[leftmargin=1.2em,itemsep=0.3ex,topsep=2pt]
\item Αν έχουν διαφορετικό ακέραιο μέρος, μεγαλύτερος είναι αυτός με το μεγαλύτερο ακέραιο μέρος: $8{,}43>5{,}96$.
\item Αν έχουν το ίδιο ακέραιο μέρος, συγκρίνουμε τα δεκαδικά τους ψηφία \textbf{ένα προς ένα από αριστερά προς τα δεξιά} και μεγαλύτερος είναι αυτός που έχει μεγαλύτερο το πρώτο ψηφίο στο οποίο διαφέρουν: $3{,}132>3{,}1294$, γιατί στα χιλιοστά $3>2$.
\end{itemize}
\end{kanonas}

\tmimasmall{Άπειροι ρητοί ανάμεσα σε δύο ρητούς}

Ανάμεσα σε δύο οποιουσδήποτε διαφορετικούς ρητούς υπάρχει ένα \textbf{άπειρο πλήθος} ρητών. Λέμε ότι οι ρητοί είναι «\textbf{πυκνοί}» στην ευθεία: δεν υπάρχει «επόμενος» ενός ρητού, σε αντίθεση με τους ακεραίους, όπου κάθε ακέραιος έχει επόμενο. Για να βρούμε έναν ρητό ανάμεσα σε δύο άλλους, τους γράφουμε ως ομώνυμα (και αν δεν υπάρχει ακέραιος ανάμεσα στους αριθμητές, πολλαπλασιάζουμε τους όρους με έναν αρκετά μεγάλο αριθμό, π.χ. $10$) ή ως δεκαδικούς.

\begin{prosoxi}
\begin{itemize}[leftmargin=1.2em,itemsep=0.3ex,topsep=2pt]
\item Στους αρνητικούς το «μεγάλο νούμερο» είναι μικρός αριθμός: $-10<-3$ παρά το $10>3$.
\item Ο αντίθετος του $-5$ είναι το $+5$ και όχι το $-5$. Και $|-5|=5$, όχι $-5$.
\item Στα δεκαδικά, ο αριθμός με τα περισσότερα ψηφία δεν είναι μεγαλύτερος: $0{,}3>0{,}2999$. Και ο $-1{,}36$ είναι \emph{μικρότερος} από τον $-1{,}3$.
\item Στα κλάσματα, ο μεγαλύτερος παρονομαστής δεν δίνει μεγαλύτερο κλάσμα: $\frac{1}{8}<\frac{1}{3}$. Τα κομμάτια είναι πιο μικρά.
\item Δεν συγκρίνουμε ετερώνυμα κλάσματα κοιτάζοντας μόνο τους αριθμητές ή μόνο τους παρονομαστές.
\end{itemize}
\end{prosoxi}

\tmima{Λυμένα παραδείγματα}

\paradeigma{Ποια πόλη είναι πιο κρύα}
\tixero Ιωάννινα $-3^\circ$, Φλώρινα $-10^\circ$.\par
\tithelo Ποια θερμοκρασία είναι μικρότερη;\par
\possindeo Και οι δύο είναι αρνητικοί, και μεγαλύτερος είναι αυτός με τη μικρότερη απόλυτη τιμή. Ισχύει $|-10|=10>3=|-3|$.\par
\apantisi $-10<-3$: η Φλώρινα είναι πιο κρύα.\par
\elegxos Στον άξονα το $-10$ είναι αριστερότερα από το $-3$.

\paradeigma{Αντίθετοι και απόλυτες τιμές}
\tixero Οι ρητοί $+4{,}5$, $-\frac{2}{3}$ και $-7$.\par
\tithelo Να βρω τις απόλυτες τιμές και τους αντίθετούς τους.\par
\possindeo Η απόλυτη τιμή είναι η απόσταση από το $0$: $4{,}5$, $\frac{2}{3}$, $7$. Ο αντίθετος του $x$ είναι ο $-x$: $-4{,}5$, $+\frac{2}{3}$, $+7$.\par
\apantisi $|{+4{,}5}|=4{,}5$, $\left|-\frac{2}{3}\right|=\frac{2}{3}$, $|-7|=7$ και αντίθετοι $-4{,}5$, $+\frac{2}{3}$, $+7$.\par
\elegxos Κάθε αριθμός και ο αντίθετός του έχουν την ίδια απόλυτη τιμή.

\paradeigma{Ετερώνυμα κλάσματα}
\tixero Τα κλάσματα $\frac{3}{5}$, $\frac{5}{8}$ και $\frac{7}{12}$.\par
\tithelo Να τα βάλω σε αύξουσα σειρά.\par
\possindeo Τα κάνω ομώνυμα. $\EKP(5,8,12)=120$. $\frac{3}{5}=\frac{72}{120}$, $\frac{5}{8}=\frac{75}{120}$, $\frac{7}{12}=\frac{70}{120}$. Συγκρίνω τους αριθμητές: $70<72<75$.\par
\apantisi $\dfrac{7}{12}<\dfrac{3}{5}<\dfrac{5}{8}$.\par
\elegxos Με δεκαδικούς: $0{,}583\ldots<0{,}6<0{,}625$.

\paradeigma{Αρνητικά κλάσματα και δεκαδικοί}
\tixero Οι αριθμοί $-\frac{5}{2}$, $-0{,}9$, $-4{,}5$, $0{,}1$, $\frac{1}{2}$, $3$.\par
\tithelo Να τους τοποθετήσω στην αριθμογραμμή και να τους συγκρίνω.\par
\possindeo Γράφω τα κλάσματα ως δεκαδικούς: $-2{,}5$ και $0{,}5$. Οι αρνητικοί είναι αριστερά του $0$ και ο μικρότερος είναι αυτός με τη μεγαλύτερη απόλυτη τιμή: $-4{,}5<-2{,}5<-0{,}9$. Μετά οι θετικοί: $0{,}1<0{,}5<3$.\par
\apantisi $-4{,}5<-\dfrac{5}{2}<-0{,}9<0{,}1<\dfrac{1}{2}<3$.\par
\elegxos Η σειρά συμφωνεί με τη σειρά των σημείων στον άξονα, από αριστερά προς τα δεξιά.

\begin{center}
\akgr[1.1]{-5}{4}{-4.5/-4{,}5,-2.5/-\frac{5}{2},-0.9/-0{,}9,0.1/0{,}1,0.5/\frac{1}{2},3/3}
\end{center}

\paradeigma{Ένας ρητός ανάμεσα}
\tixero Οι αριθμοί $\frac{1}{3}$ και $\frac{1}{2}$.\par
\tithelo Να βρω έναν ρητό ανάμεσά τους.\par
\possindeo \emph{1ος τρόπος.} Κάνω τα κλάσματα ομώνυμα: $\frac{1}{3}=\frac{2}{6}$ και $\frac{1}{2}=\frac{3}{6}$. Δεν υπάρχει ακέραιος ανάμεσα στο $2$ και στο $3$, γι' αυτό πολλαπλασιάζω τους όρους επί $10$: $\frac{20}{60}$ και $\frac{30}{60}$. Ανάμεσα βρίσκεται το $\frac{25}{60}$. \emph{2ος τρόπος.} Με δεκαδικούς: $\frac{1}{3}=0{,}333\ldots$ και $\frac{1}{2}=0{,}5$, άρα ανάμεσα βρίσκεται ο $0{,}4$.\par
\apantisi $\dfrac{1}{3}<\dfrac{25}{60}<\dfrac{1}{2}$ και $\dfrac{1}{3}<0{,}4<\dfrac{1}{2}$.\par
\elegxos $\frac{25}{60}=\frac{5}{12}=0{,}41\overline{6}$. Βρίσκεται ανάμεσα στο $0{,}333\ldots$ και στο $0{,}5$.

\paradeigma{Σύγκριση δεκαδικών}
\tixero Οι αριθμοί $9{,}8$ και $9{,}75$, και οι $-1{,}36$ και $-1{,}3$.\par
\tithelo Ποιος είναι μεγαλύτερος σε κάθε ζεύγος;\par
\possindeo Ίδιο ακέραιο μέρος στο πρώτο ζεύγος. Γράφω $9{,}8=9{,}80$ και συγκρίνω τα δέκατα: $8>7$, άρα $9{,}80>9{,}75$. Στο δεύτερο οι αριθμοί είναι αρνητικοί, και $|-1{,}36|=1{,}36>1{,}30=|-1{,}3|$, άρα $-1{,}36<-1{,}3$.\par
\apantisi $9{,}8>9{,}75$ και $-1{,}3>-1{,}36$.\par
\elegxos Αν $9{,}8$ και $9{,}75$ είναι χρόνοι τρεξίματος, ο $9{,}75$ είναι ο καλύτερος.

\tmima{Ασκήσεις}

\zoni{Το κατάλαβα}{απόλυτη τιμή, αντίθετοι, σύγκριση}

\ex Συμπλήρωσε τα κενά. α) Ο αντίθετος του $\alpha$ είναι ο \ldots \quad β) Από δύο θετικούς ρητούς μεγαλύτερος είναι αυτός με τη \ldots απόλυτη τιμή. \quad γ) Από δύο αρνητικούς ρητούς μεγαλύτερος είναι αυτός με τη \ldots απόλυτη τιμή. \quad δ) Από δύο ομώνυμα κλάσματα μεγαλύτερο είναι αυτό με τον μεγαλύτερο \ldots \quad ε) Από δύο κλάσματα με ίδιο αριθμητή μεγαλύτερο είναι αυτό με τον \ldots παρονομαστή.

\ex Συμπλήρωσε τον πίνακα.

\begin{center}
\renewcommand{\arraystretch}{1.6}
\begin{tabular}{|l|c|c|c|c|c|c|}
\hline
Αριθμός & $-3$ & $0{,}35$ & $\frac{1}{9}$ & & & $-\frac{2}{5}$\\
\hline
Αντίθετος & & & & $-0{,}1$ & $\frac{7}{4}$ & \\
\hline
Απόλυτη τιμή & & & & & & \\
\hline
\end{tabular}
\end{center}

\ex Σύγκρινε με ένα από τα $<,\ >,\ =$: α) $\frac{21}{6}\ \ldots\ \frac{12}{6}$ \quad β) $\frac{4}{9}\ \ldots\ \frac{4}{6}$ \quad γ) $\frac{2}{3}\ \ldots\ \frac{7}{8}$ \quad δ) $\frac{1}{100}\ \ldots\ -\frac{2}{5}$ \quad ε) $-\frac{8}{10}\ \ldots\ -1$ \quad στ) $-\frac{8}{3}\ \ldots\ -3$ \quad ζ) $\frac{9}{8}\ \ldots\ 1{,}2$.

\ex Σύγκρινε τους δεκαδικούς: α) $3{,}16\ \ldots\ 3{,}2$ \quad β) $0{,}092\ \ldots\ 0{,}087$ \quad γ) $-1{,}36\ \ldots\ -0{,}01$ \quad δ) $-2{,}69\ \ldots\ 1{,}4$ \quad ε) $-1{,}3\ \ldots\ -1{,}33$ \quad στ) $-1{,}10\ \ldots\ -1{,}1$.

\ex Βάλε σε αύξουσα σειρά τους δεκαδικούς: $1,\ -3,\ 0{,}2,\ 1{,}37,\ -0{,}03,\ 0$.

\ex Βάλε σε αύξουσα σειρά τα κλάσματα: α) $\frac{5}{4},\ \frac{1}{4},\ -\frac{3}{4},\ \frac{3}{4},\ -\frac{1}{4},\ -\frac{5}{4}$ \quad β) $\frac{5}{3},\ \frac{5}{4},\ -\frac{5}{2},\ \frac{5}{8},\ -\frac{5}{7},\ -\frac{5}{3}$.

\ex Δίνονται τα κλάσματα $\frac{9}{4},\ \frac{3}{16},\ \frac{9}{8},\ \frac{15}{16},\ \frac{5}{4},\ \frac{7}{16}$. α) Βρες ποια είναι μικρότερα του $1$ και γράψε τα σε αύξουσα σειρά. β) Βρες ποια είναι μεγαλύτερα του $1$ και γράψε τα σε αύξουσα σειρά. γ) Γράψε και τα έξι σε αύξουσα σειρά.

\ex Σωστό ή λάθος; α) $|\alpha|\ge0$ για κάθε ρητό $\alpha$. \quad β) $|-5|=-5$. \quad γ) Ο αντίθετος του $-4$ είναι ο $4$. \quad δ) $-9>-2$. \quad ε) $\frac{3}{7}<\frac{3}{5}$. \quad στ) Ο αντίθετος του $0$ είναι το $0$.

\zoni{Το χρησιμοποιώ}{θερμόμετρα, ώρες, αθλήματα}

\ex \textbf{Οι θερμοκρασίες.} Στον πίνακα του καιρού: Μέτσοβο $-5{,}5^\circ$, Ιωάννινα $-3^\circ$, Φλώρινα $-10^\circ$, Αθήνα $7^\circ$, Θεσσαλονίκη $-1{,}5^\circ$. Γράψε τις πόλεις από την πιο κρύα στην πιο ζεστή.

\ex \textbf{Οι δραστηριότητες.} Στο κυκλικό διάγραμμα της Αγγελικής ο Ύπνος παίρνει το $\frac{1}{3}$ του εικοσιτετραώρου, το Σχολείο το $\frac{1}{4}$, το Διάβασμα το $\frac{1}{8}$ και η Διασκέδαση το $\frac{1}{6}$. α) Πόσες ώρες διαρκεί η κάθε δραστηριότητα; β) Γράψε τις δραστηριότητες από αυτή που παίρνει τον περισσότερο χρόνο ως αυτή που παίρνει τον λιγότερο. γ) Τι μέρος του εικοσιτετραώρου μένει για όλα τα άλλα;

\ex \textbf{Τα εκατοστά του δευτερολέπτου.} Ο Νίκος τρέχει τα $60$ μέτρα σε $9{,}8$ δευτερόλεπτα και ο Βασίλης σε $9{,}75$. Ο Βασίλης λέει «το $75$ είναι μεγαλύτερο από το $8$, άρα έκανα περισσότερο χρόνο». Ποιος κέρδισε και ποιο είναι το λάθος του Βασίλη;

\ex \textbf{Οι αποστάσεις.} Στον άξονα της οδού το σχολείο είναι στο $0$ (μονάδα $100$ μ.), ο φούρνος στο $+3$ και το βιβλιοπωλείο στο $-2$. Ποιο μαγαζί είναι πιο κοντά στο σχολείο; Πόσο απέχουν μεταξύ τους τα δύο μαγαζιά;

\ex \textbf{Τα υψόμετρα.} Η Νεκρά Θάλασσα είναι στα $-430$ μ., ο Όλυμπος στα $+2\,917$ μ., η επιφάνεια της θάλασσας στο $0$ και ο βυθός της τάφρου των Μαριανών στα $-11\,034$ μ. Γράψε τα σε αύξουσα σειρά υψομέτρου.

\zoni{Το σκέφτομαι}{γιατί δεν υπάρχει επόμενος;}

\ex Ο Βασίλης λέει ότι ο επόμενος ρητός του $\frac{2}{5}$ είναι ο $\frac{3}{5}$. Βρες τρεις ρητούς ανάμεσά τους.

\ex Βρες δύο ρητούς ανάμεσα στους: α) $-0{,}7$ και $-0{,}8$ \quad β) $\frac{3}{10}$ και $\frac{2}{5}$ \quad γ) $\frac{7}{9}$ και $\frac{5}{6}$.

\ex Για ποιους ρητούς ισχύει $|\alpha|=\alpha$ και για ποιους $|\alpha|=-\alpha$; Γιατί ισχύει πάντα $|\alpha|=|-\alpha|$;

\ex Βρες τους δύο διαδοχικούς ακεραίους ανάμεσα στους οποίους βρίσκεται κάθε αριθμός: α) $0{,}3$ \quad β) $-1{,}25$ \quad γ) $\frac{7}{2}$ \quad δ) $-\frac{11}{3}$ \quad ε) $\frac{2024}{2025}$.

\zoni{Για το σπίτι}{ψάξε και σύγκρινε}

\ex Βρες τις θερμοκρασίες πέντε πόλεων του κόσμου και βάλε τις σε σειρά από την πιο κρύα στην πιο ζεστή. Σχεδίασε και τον άξονα.

\ex Ζήτησε από τρία μέλη της οικογένειάς σου να σου πουν τον χρόνο που κοιμούνται (σε ώρες) και γράψε το κλάσμα του εικοσιτετραώρου. Ποιος κοιμάται περισσότερο;

\providecommand{\key}[1]{\mbox{\fbox{\rule{0pt}{1.9ex}\hspace{0.15em}#1\hspace{0.15em}}}\hspace{0.1em}}
\providecommand{\kpl}{\key{$+$}}\providecommand{\kmi}{\key{$-$}}\providecommand{\kpo}{\key{$\times$}}\providecommand{\kdi}{\key{$\div$}}\providecommand{\kis}{\key{$=$}}

\paragrafos{Υπολογιστής τσέπης: ένα εργαλείο, όχι ένα μυαλό}{Α.3.3, «Υπολογισμοί με τη βοήθεια υπολογιστή τσέπης»}

\begin{anagke}
Μετά το παζάρι, το ταμείο έχει γεμίσει αποδείξεις. Ο Νίκος, που κρατάει το ταμείο, αρχίζει να τις προσθέτει με το χέρι και βγάζει $68{,}79$ €. Η Πολυξένη τις ξαναπροσθέτει με το κομπιουτεράκι του κινητού της και βγάζει $77{,}79$ €. «Ένας από τους δύο έχει λάθος», λέει ο Ηλίας. «Το κομπιουτεράκι δεν κάνει λάθος», απαντά ο Βασίλης. «Άρα ο Νίκος».

Η Πολυξένη κοιτάζει την οθόνη της και βλέπει ότι πάτησε κατά λάθος το $23{,}45$ ως $32{,}45$: τα ψηφία αντιστράφηκαν, και η διαφορά στο σύνολο είναι ακριβώς $9$ ευρώ. «Το κομπιουτεράκι έκανε σωστή πράξη με λάθος δεδομένα», παρατηρεί η κυρία Αλεξάνδρα. «Δεν σκέφτεται: κάνει ό,τι του πεις». Και ο Στέλιος, που θέλει να δοκιμάσει το εργαλείο, πατάει $1:3$ και βλέπει $0{,}33333333$ με οκτώ ψηφία. «Χάλασε; Το $\frac{1}{3}$ δεν είναι ακριβώς αυτό!».

Στο κεφάλαιο αυτό η παρέα μαθαίνει να χρησιμοποιεί το κομπιουτεράκι σωστά: πώς πατάμε τις πράξεις, πώς φυλάμε ενδιάμεσα αποτελέσματα στη μνήμη, πώς ξέρουμε αν η οθόνη λέει αλήθεια.
\end{anagke}

\begin{lexi}{υπολογιστής, μνήμη, οθόνη}{υπολογιστής, μνήμη}
Ο \lexiword{υπολογιστής} είναι αυτός που «υπολογίζει», από το «λόγος, λογαριασμός»: το να υπολογίζεις σημαίνει να κάνεις λογαριασμό. Παλιά οι υπολογιστές ήταν άνθρωποι με άβακα. Ο \lexiword{άβακας} ήταν ένας πίνακας με χάντρες. Σήμερα το μηχάνημα έχει δική του \lexiword{μνήμη}, ένα «συρτάρι» όπου φυλάμε έναν αριθμό για να τον ξαναβρούμε αργότερα.
\end{lexi}

\begin{metafrasi}
Πρόσθεσε $12{,}35$ και $7{,}8$. & \key{12,35} \kpl\ \key{7,8} \kis\\
Φύλαξε τον αριθμό της οθόνης στη μνήμη. & \key{M+}\\
Αφαίρεσε από τη μνήμη τον αριθμό της οθόνης. & \key{M$-$}\\
Δείξε μου τι υπάρχει στη μνήμη. & \key{MR}\\
Άδειασε τη μνήμη. & \key{MC}\\
Πρώτα η παρένθεση, μετά το υπόλοιπο. & \key{(}\ldots\key{)} ή ενδιάμεσα αποτελέσματα\\
\end{metafrasi}

\tmima{Η θεωρία}

\tmimasmall{Οι τέσσερις πράξεις}

Όλα τα κομπιουτεράκια εκτελούν τις τέσσερις πράξεις. Για να κάνουμε μια πράξη μεταξύ δύο αριθμών, πληκτρολογούμε τον πρώτο, το σύμβολο της πράξης, τον δεύτερο και μετά το \key{$=$}. Τα σύμβολα είναι \kpl, \kmi, \kpo{} και \kdi. Το κόμμα των δεκαδικών είναι συνήθως τελεία στην οθόνη και στο πλήκτρο. Για παράδειγμα, για το $7{,}28:5{,}2$ πατάμε \key{7.28} \kdi\ \key{5.2} \kis{} και βλέπουμε $1{,}4$.

\tmimasmall{Η μνήμη}

Αν το κομπιουτεράκι έχει μνήμη, έχει τέσσερα πλήκτρα:
\begin{itemize}[leftmargin=1.2em,itemsep=0.3ex,topsep=2pt]
\item \key{MR} (ή \key{MRC}): εμφανίζει στην οθόνη τον αριθμό που βρίσκεται στη μνήμη.
\item \key{MC}: σβήνει το περιεχόμενο της μνήμης.
\item \key{M+}: προσθέτει στη μνήμη τον αριθμό της οθόνης.
\item \key{M$-$}: αφαιρεί από τη μνήμη τον αριθμό της οθόνης.
\end{itemize}
Η μνήμη είναι ό,τι χρειάζεται η παρένθεση του χαρτιού: φυλάμε ένα αποτέλεσμα, κάνουμε άλλο υπολογισμό και μετά τα συνδυάζουμε.

\tmimasmall{Παραστάσεις με πολλές πράξεις}

Δεν εμπιστευόμαστε τυφλά τη σειρά των πλήκτρων. Μερικά απλά κομπιουτεράκια εκτελούν τις πράξεις με τη σειρά που τις πληκτρολογούμε, χωρίς να υπολογίζουν πρώτα τους πολλαπλασιασμούς (τα επιστημονικά την σέβονται). Για παραστάσεις με παρενθέσεις ή μικτές πράξεις, είναι ασφαλές να κάνουμε:
\begin{enumerate}[leftmargin=1.6em,itemsep=0.3ex,topsep=2pt]
\item Μια πράξη τη φορά, με την προτεραιότητα του χαρτιού (παρενθέσεις, δυνάμεις, πολλαπλασιασμός/διαίρεση, μετά πρόσθεση/αφαίρεση).
\item Φύλαξη των αποτελεσμάτων που θα χρειαστούμε αργότερα με \key{M+} ή \key{M$-$}, ή απλή γραφή τους στο τετράδιο.
\end{enumerate}

\tmimasmall{Η οθόνη δείχνει προσέγγιση}

Η οθόνη έχει συγκεκριμένο πλήθος ψηφίων (συνήθως $8$ ή $10$). Όταν το αποτέλεσμα έχει περισσότερα ψηφία, το κομπιουτεράκι το κόβει ή το στρογγυλοποιεί. Γι' αυτό $1:3$ δείχνει $0{,}33333333$ και $2:3$ δείχνει $0{,}66666667$. Άρα $3\cdot0{,}33333333=0{,}99999999$ και όχι $1$. Δεν σημαίνει ότι το κομπιουτεράκι χάλασε, αλλά ότι δείχνει μια προσέγγιση.

\begin{prosoxi}
\begin{itemize}[leftmargin=1.2em,itemsep=0.3ex,topsep=2pt]
\item Το κομπιουτεράκι κάνει σωστά ό,τι του πεις, ακόμα και αν του πεις λάθος. Πριν από κάθε χρήση κάνω μια \textbf{εκτίμηση} με στρογγυλά νούμερα: τόσο περίπου πρέπει να βγει.
\item Δεν ξεχνάμε να σβήσουμε τη μνήμη (\key{MC}) πριν από νέο υπολογισμό, αλλιώς προσθέτουμε σε παλιό αριθμό.
\item Η τιμή που δείχνει η οθόνη μπορεί να είναι στρογγυλοποιημένη. Όταν χρειαζόμαστε ακρίβεια, δεν στρογγυλοποιούμε ενδιάμεσα αποτελέσματα, μόνο το τελικό.
\end{itemize}
\end{prosoxi}

\tmima{Λυμένα παραδείγματα}

\paradeigma{Το σύνολο του ταμείου με τη μνήμη}
\tixero Οι αποδείξεις είναι $12{,}35$, $7{,}8$, $23{,}45$, $9{,}99$ και $15{,}2$ ευρώ.\par
\tithelo Το σύνολο, και ο μέσος όρος ανά απόδειξη.\par
\possindeo Εκτίμηση: $12+8+23+10+15=68$. Πατάω \key{MC}, και για κάθε απόδειξη πληκτρολογώ τον αριθμό και \key{M+}: \key{12.35}\key{M+}\ \key{7.8}\key{M+}\ \key{23.45}\key{M+}\ \key{9.99}\key{M+}\ \key{15.2}\key{M+}. Μετά \key{MR}. Για τον μέσο όρο διαιρώ με $5$: $68{,}79:5=13{,}758$.\par
\apantisi Σύνολο $68{,}79$ €, κατά μέσο όρο $13{,}76$ € (στρογγυλοποιημένο στο λεπτό).\par
\elegxos Η εκτίμηση ($68$) είναι πολύ κοντά στο $68{,}79$, ενώ το $77{,}79$ της Πολυξένης είναι εκτός (η εκτίμηση είναι $68$).

\paradeigma{Σύνθετη παράσταση με μνήμη}
\tixero Η παράσταση $(2{,}4:0{,}6+1{,}5\cdot4)\cdot2{,}5-18{,}9:(0{,}7\cdot5+3{,}5:0{,}5)$.\par
\tithelo Να την υπολογίσω με κομπιουτεράκι που έχει μνήμη.\par
\possindeo Πρώτα η πρώτη παρένθεση: $2{,}4:0{,}6=4$, $1{,}5\cdot4=6$, $4+6=10$, $10\cdot2{,}5=25$. Πατάω \key{M+} (η μνήμη έχει $25$). Μετά η δεύτερη: $0{,}7\cdot5=3{,}5$, $3{,}5:0{,}5=7$, $3{,}5+7=10{,}5$ και $18{,}9:10{,}5=1{,}8$. Πατάω \key{M$-$} (η μνήμη έχει $25-1{,}8$). Τέλος \key{MR}.\par
\apantisi $23{,}2$.\par
\elegxos Εκτίμηση: $(4+6)\cdot2{,}5=25$ και $18{,}9:10{,}5\approx2$, άρα $\approx23$. Συμφωνεί.

\paradeigma{Ο έλεγχος με εκτίμηση}
\tixero Ο Ηλίας πληκτρολογεί $38{,}5\cdot4{,}2$ και η οθόνη δείχνει $16{,}17$.\par
\tithelo Να δω αν είναι σωστό.\par
\possindeo Εκτίμηση: $40\cdot4=160$. Το $16{,}17$ είναι δέκα φορές μικρότερο. Ο Ηλίας πάτησε $3{,}85$ αντί για $38{,}5$ (η υποδιαστολή σε λάθος θέση). Διορθώνω: $38{,}5\cdot4{,}2=161{,}7$.\par
\apantisi $161{,}7$.\par
\elegxos $161{,}7\approx160$. Και $3{,}85\cdot4{,}2=16{,}17$, το λάθος που είδαμε.

\paradeigma{Η οθόνη δείχνει προσέγγιση}
\tixero Το κομπιουτεράκι έχει οθόνη οκτώ ψηφίων.\par
\tithelo Τι δείχνει για το $2:3$ και πόσο βγαίνει το $3\cdot(2:3)$ αν χρησιμοποιήσω την οθόνη;\par
\possindeo $2:3=0{,}666\ldots$ και η οθόνη κρατάει οκτώ ψηφία, στρογγυλοποιώντας το τελευταίο: $0{,}66666667$. Τριπλασιάζω: $3\cdot0{,}66666667=2{,}00000001$.\par
\apantisi Η οθόνη δείχνει $0{,}66666667$, και το γινόμενο με το $3$ βγαίνει $2{,}00000001$ αντί για $2$.\par
\elegxos Η διαφορά $0{,}00000001$ είναι αμελητέα, αλλά δείχνει ότι δουλεύουμε με προσέγγιση.

\paradeigma{Πολλαπλασιασμοί που προστίθενται}
\tixero Στο παζάρι πουλήθηκαν $12$ κέικ των $2{,}8$ €, $8$ μπουκάλια των $1{,}5$ € και $20$ μπισκότα των $0{,}6$ €.\par
\tithelo Τα συνολικά έσοδα.\par
\possindeo Κάθε γινόμενο το προσθέτω στη μνήμη: \key{12}\kpo\key{2.8}\kis\key{M+}\ \key{8}\kpo\key{1.5}\kis\key{M+}\ \key{20}\kpo\key{0.6}\kis\key{M+}\ \key{MR}.\par
\apantisi $33{,}6+12+12=57{,}6$ €.\par
\elegxos Εκτίμηση: $12\cdot3=36$ και $8\cdot1{,}5=12$ και $20\cdot0{,}6=12$, άρα $\approx60$. Συμφωνεί.

\tmima{Ασκήσεις}

\zoni{Το κατάλαβα}{πλήκτρα, μνήμη, εκτίμηση}

\ex Υπολόγισε με κομπιουτεράκι: α) $318{,}76+47{,}092$ \quad β) $1\,000-87{,}654$ \quad γ) $24{,}6\cdot3{,}05$ \quad δ) $91{,}2:0{,}16$.

\ex Με εκτίμηση βρες ποιο αποτέλεσμα είναι σωστό χωρίς να κάνεις την πράξη: α) $48{,}7\cdot5{,}9$ είναι $28{,}733$ ή $287{,}33$ ή $2873{,}3$; β) $19{,}8:4{,}5$ είναι $0{,}44$ ή $4{,}4$ ή $44$;

\ex Γράψε με ποια πλήκτρα θα υπολογίσεις: α) $(4{,}5+2{,}7)\cdot1{,}25$ \quad β) $36:(0{,}9+2{,}1)$. Πόσο βγαίνουν;

\ex Ποια παράσταση υπολογίζεται και πόσο βγαίνει; α) \key{MC} \key{6.5}\kpo\key{4}\kis\key{M+} \key{3.2}\kpo\key{5}\kis\key{M$-$} \key{MR}. \quad β) \key{1.2}\kpl\key{0.8}\kis\kdi\key{0.5}\kis.

\ex Γράψε τα πλήκτρα που χρειάζονται για τον υπολογισμό της $2{,}5\cdot8+1{,}5\cdot6-4{,}25\cdot2$ με χρήση μνήμης. Πόσο είναι το αποτέλεσμα;

\ex Υπολόγισε με κομπιουτεράκι: α) $(3{,}6+4{,}4):0{,}8-2{,}5$ \quad β) $12{,}6:(4{,}2\cdot0{,}5)$ \quad γ) $0{,}25\cdot(10-3{,}2)+1{,}2$.

\ex Σωστό ή λάθος; α) Το κομπιουτεράκι δεν κάνει ποτέ λάθος. \quad β) Η οθόνη μετά το $1:3$ δείχνει την ακριβή τιμή του $\frac{1}{3}$. \quad γ) Το \key{M+} προσθέτει στη μνήμη τον αριθμό της οθόνης. \quad δ) Όλα τα κομπιουτεράκια εκτελούν πρώτα τους πολλαπλασιασμούς. \quad ε) Όταν έχω κομπιουτεράκι δεν χρειάζεται εκτίμηση.

\zoni{Το χρησιμοποιώ}{ταμείο, καύσιμα, μοιρασιά}

\ex \textbf{Το ταμείο.} Στο παζάρι έσοδα ήταν $57{,}6$ € και τα έξοδα (υλικά, χαρτικά) $18{,}35$ €. Πόσα καθαρά κέρδη έχει το ταμείο της εκδρομής; Έλεγξε με εκτίμηση.

\ex \textbf{Η απόδειξη.} Ο Στέλιος πληκτρολόγησε $6{,}75\cdot1{,}2$ και είδε $81$. Εξήγησε τι πάτησε λάθος και βρες το σωστό αποτέλεσμα.

\ex \textbf{Το ταξίδι.} Το αυτοκίνητο του θείου Λευτέρη καταναλώνει $7{,}8$ λίτρα ανά $100$ χλμ. Πόση βενζίνη θα χρειαστεί για μια διαδρομή $312$ χλμ.; Στρογγυλοποίησε στο δέκατο του λίτρου.

\ex \textbf{Η μοιρασιά.} Το ταμείο των $100$ € μοιράζεται εξίσου σε $7$ παιδιά. Τι δείχνει η οθόνη; Πόσα ευρώ παίρνει το καθένα στρογγυλοποιημένα στο λεπτό και πόσα θα έδινε συνολικά το ταμείο με αυτή τη στρογγυλοποίηση;

\ex \textbf{Το ρεύμα.} Μια λάμπα ισχύος $0{,}06$ kW ανάβει $4{,}5$ ώρες την ημέρα για $30$ ημέρες. Με τιμή $0{,}18$ € ανά kWh, πόσο κοστίζει στον μήνα; (Κόστος = ισχύς × ώρες × τιμή.)

\zoni{Το σκέφτομαι}{ο έλεγχος είναι δική μας δουλειά}

\ex Εξήγησε γιατί στο κομπιουτεράκι των $8$ ψηφίων το $3\cdot(1:3)$ δεν δίνει ακριβώς $1$.

\ex Με ένα απλό κομπιουτεράκι, αν πατήσεις \key{2}\kpl\key{3}\kpo\key{4}\kis{} βλέπεις συχνά $20$ και όχι $14$. Ποια σειρά εκτέλεσης ακολούθησε και πώς θα την πατήσεις για να βγει το σωστό;

\ex Χωρίς να κάνεις την πράξη, εξήγησε γιατί το $487\cdot0{,}98$ δεν μπορεί να είναι $4\,773{,}3$ και βρες τη σωστή τάξη μεγέθους.

\ex Πόσες φορές πρέπει να πατήσεις \kpl\ \key{0.1} \kis{} ξεκινώντας από το $0$, για να φτάσεις το $2{,}5$; Με ποια διαίρεση το βρίσκεις αμέσως;

\ex Κάνε τη διαίρεση $1:7$ στο κομπιουτεράκι σου. Γράψε όσα ψηφία φαίνονται και στρογγυλοποίησέ τον στα εκατοστά. Τι μοτίβο βλέπεις στα ψηφία;

\zoni{Για το σπίτι}{δοκίμασε το δικό σου κομπιουτεράκι}

\ex Πάρε τον λογαριασμό ρεύματος ή νερού του σπιτιού, και με το κομπιουτεράκι έλεγξε ένα από τα γινόμενα ή τα αθροίσματα που αναγράφονται. Βγαίνει το ίδιο;

\ex Βρες τις θερμοκρασίες της τελευταίας εβδομάδας και με το κομπιουτεράκι υπολόγισε τον μέσο όρο τους. Στρογγυλοποίησέ τον στο δέκατο.

\paragrafos{Τυποποιημένη μορφή: οι μεγάλοι αριθμοί χωράνε σε μια γραμμή}{Α.3.4, «Τυποποιημένη μορφή μεγάλων αριθμών»}

\begin{anagke}
Η κυρία Αλεξάνδρα ζητάει από την Α3 μια εργασία για το ηλιακό σύστημα. Η Αγγελική γράφει στο τετράδιό της την απόσταση της Γης από τον Ήλιο: $149\,600\,000$ χιλιόμετρα. Ο Ηλίας γράφει τη μάζα της Γης: $5\,976\,000\,000\,000\,000\,000\,000\,000$ κιλά, και το γράφει τρεις φορές, γιατί δεν είναι σίγουρος για τα μηδενικά. Ο Στέλιος δοκιμάζει με τον Ήλιο και χάνεται στο δεύτερο μέτρημα. «Αυτό είναι αδύνατο να διαβαστεί», παραπονιέται η Μυρτώ.

Η Πολυξένη, που έχει το κομπιουτεράκι της, πατάει $149\,600\,000\cdot10$ και η οθόνη δείχνει $1.496\mathrm{E}9$. «Τι είναι αυτό το E9;». Η κυρία Αλεξάνδρα χαμογελάει: «Το κομπιουτεράκι δεν χωράει όλα τα μηδενικά, και τα γράφει με συντομία: $1{,}496$ επί δέκα στην ένατη. Κάνουν ό,τι κάνουν και οι επιστήμονες». Και η Φανή ρωτάει: «Μπορούμε να γράφουμε κι εμείς έτσι, ώστε να μη χάνουμε τα μηδενικά;».
\end{anagke}

\begin{lexi}{τυποποιημένη μορφή}{τυποποιημένη μορφή}
Η λέξη \lexiword{τύπος} σημαίνει καλούπι, πρότυπο. \lexiword{Τυποποιημένη} είναι η μορφή που έχει το ίδιο καλούπι για όλους τους αριθμούς: ένας αριθμός από $1$ μέχρι $9{,}99\ldots$ επί μια δύναμη του $10$. Οι επιστήμονες την χρησιμοποιούν γιατί δείχνει αμέσως το \emph{μέγεθος} του αριθμού (πόσα μηδενικά) και τα σημαντικά του ψηφία.
\end{lexi}

\begin{metafrasi}
Εκατόν σαράντα εννέα εκατομμύρια εξακόσιες χιλιάδες. & $149\,600\,000=1{,}496\cdot10^8$\\
Τόσα μηδενικά όσο ο εκθέτης. & $10^8=100\,000\,000$ (το $1$ και $8$ μηδενικά)\\
Ο $4{,}56\cdot10^{10}$ γραμμένος κανονικά. & $45\,600\,000\,000$\\
Στην οθόνη του κομπιουτεράκι. & $1.496\mathrm{E}8$ σημαίνει $1{,}496\cdot10^8$\\
Σύγκριση: μεγαλύτερος εκθέτης. & $8{,}2\cdot10^5>9{,}1\cdot10^4$\\
\end{metafrasi}

\tmima{Η θεωρία}

\begin{orismos}{τυποποιημένη μορφή}
Ένας (μεγάλος) αριθμός είναι γραμμένος σε \textbf{τυποποιημένη μορφή} όταν είναι γραμμένος ως $\alpha\cdot10^{\nu}$, όπου $\nu$ φυσικός αριθμός και ο $\alpha$ είναι ένας δεκαδικός (ή φυσικός) με ακέραιο μέρος ένα ψηφίο, δηλαδή $1\leq\alpha<10$.
\end{orismos}

Για παράδειγμα
$$3\,140\,000\,000\,000\,000\,000=3{,}14\cdot10^{18},\qquad 234\,000\,000\,000\,000\,000=2{,}34\cdot10^{17}.$$
Και αντίστροφα, ο $5{,}21\cdot10^5$ είναι ο $521\,000$ και ο $2\cdot10^3$ ο $2\,000$.

\tmimasmall{Πώς γράφω έναν αριθμό σε τυποποιημένη μορφή}

\begin{kanonas}{τυποποίηση}
\begin{enumerate}[leftmargin=1.6em,itemsep=0.3ex,topsep=2pt]
\item Μετακινώ την υποδιαστολή (στους ακεραίους είναι στο τέλος) προς τα αριστερά, μέχρι να μείνει \textbf{ένα μη μηδενικό ψηφίο} πριν από αυτήν. Έτσι βρίσκω τον $\alpha$ (με $1\leq\alpha<10$).
\item Ο εκθέτης $\nu$ είναι το πλήθος των θέσεων που μετακίνησα.
\item Γράφω $\alpha\cdot10^{\nu}$. Αφαιρώ τα μηδενικά που περισσεύουν στο τέλος του $\alpha$.
\end{enumerate}
\end{kanonas}

Για να πάω πίσω, από την τυποποιημένη στη δεκαδική μορφή, πολλαπλασιάζω τον $\alpha$ επί $10^{\nu}$: μετακινώ την υποδιαστολή του $\alpha$ δεξιά κατά $\nu$ θέσεις, βάζοντας μηδενικά όπου χρειάζεται. Έτσι $7{,}4\cdot10^3=7\,400$.

\tmimasmall{Σύγκριση και πράξεις}

Για να συγκρίνω δύο αριθμούς σε τυποποιημένη μορφή, κοιτάζω πρώτα τον \textbf{εκθέτη}: μεγαλύτερος εκθέτης σημαίνει μεγαλύτερο αριθμό, επειδή έχει περισσότερα ψηφία. Αν οι εκθέτες είναι ίσοι, συγκρίνω τους $\alpha$.

Για να πολλαπλασιάσω, πολλαπλασιάζω τους $\alpha$ μεταξύ τους και προσθέτω τους εκθέτες του $10$: $(2{,}5\cdot10^4)\cdot(4\cdot10^3)=10\cdot10^7=1\cdot10^8$. Προσοχή: αν το γινόμενο των $\alpha$ βγει $10$ ή περισσότερο, πρέπει να ξαναγράψω το αποτέλεσμα σε τυποποιημένη μορφή.

\begin{prosoxi}
\begin{itemize}[leftmargin=1.2em,itemsep=0.3ex,topsep=2pt]
\item Ο $\alpha$ είναι πάντα μικρότερος από το $10$ και τουλάχιστον $1$. Οι $52\cdot10^2$, $0{,}52\cdot10^4$ και $10\cdot10^5$ δεν είναι τυποποιημένοι (αν και ισούνται με $5{,}2\cdot10^3$, $5{,}2\cdot10^3$ και $1\cdot10^6$ αντίστοιχα).
\item Ο εκθέτης δεν είναι το πλήθος των ψηφίων αλλά το πλήθος των ψηφίων μείον ένα. Ο $45\,600\,000\,000$ έχει $11$ ψηφία και γράφεται $4{,}56\cdot10^{10}$.
\item Δεν συγκρίνω τους $\alpha$ όταν διαφέρουν οι εκθέτες: $9{,}9\cdot10^3<1\cdot10^4$, αν και το $9{,}9>1$.
\end{itemize}
\end{prosoxi}

\tmima{Λυμένα παραδείγματα}

\paradeigma{Η απόσταση Γη -- Ήλιος}
\tixero Η Γη απέχει από τον Ήλιο $149\,600\,000$ χλμ.\par
\tithelo Να το γράψω σε τυποποιημένη μορφή.\par
\possindeo Μετακινώ την υποδιαστολή (που είναι στο τέλος) προς τα αριστερά ώσπου να μείνει ένα ψηφίο πριν από αυτήν: $1{,}49600000$. Μετακίνησα $8$ θέσεις. Αφαιρώ τα τελευταία μηδενικά.\par
\apantisi $1{,}496\cdot10^8$ χλμ.\par
\elegxos $1{,}496\cdot100\,000\,000=149\,600\,000$. Και ο αριθμός έχει $9$ ψηφία, εκθέτης $9-1=8$.

\paradeigma{Από τυποποιημένη σε δεκαδική}
\tixero Οι αριθμοί $7{,}03\cdot10^6$ και $4{,}56\cdot10^{10}$.\par
\tithelo Να τους γράψω κανονικά.\par
\possindeo Ο $7{,}03\cdot10^6$: μετακινώ την υποδιαστολή $6$ θέσεις δεξιά, $7\,030\,000$. Ο $4{,}56\cdot10^{10}$: $10$ θέσεις δεξιά, $45\,600\,000\,000$.\par
\apantisi $7\,030\,000$ και $45\,600\,000\,000$.\par
\elegxos Στον πρώτο υπάρχουν $7$ ψηφία (εκθέτης $6$) και στον δεύτερο $11$ ψηφία (εκθέτης $10$).

\paradeigma{Ποιος είναι μεγαλύτερος;}
\tixero Οι αριθμοί $8{,}2\cdot10^5$, $9{,}1\cdot10^4$, $3{,}5\cdot10^7$ και $2{,}9\cdot10^7$.\par
\tithelo Να τους βάλω σε αύξουσα σειρά.\par
\possindeo Κοιτάζω πρώτα τους εκθέτες: $4<5<7$. Ο $9{,}1\cdot10^4$ είναι ο μικρότερος και ο $8{,}2\cdot10^5$ ακολουθεί. Οι άλλοι δύο έχουν εκθέτη $7$, και $2{,}9<3{,}5$.\par
\apantisi $9{,}1\cdot10^4<8{,}2\cdot10^5<2{,}9\cdot10^7<3{,}5\cdot10^7$.\par
\elegxos Δεκαδικά: $91\,000<820\,000<29\,000\,000<35\,000\,000$.

\paradeigma{Πόση ώρα ταξιδεύει το φως;}
\tixero Το φως ταξιδεύει με $3\cdot10^5$ χλμ. ανά δευτερόλεπτο και ο Ήλιος απέχει $1{,}496\cdot10^8$ χλμ.\par
\tithelo Σε πόσο χρόνο το φως φτάνει από τον Ήλιο στη Γη.\par
\possindeo Χρόνος = απόσταση : ταχύτητα $=\dfrac{1{,}496\cdot10^8}{3\cdot10^5}=\dfrac{1{,}496}{3}\cdot10^3\approx0{,}4987\cdot10^3\approx499$ δευτερόλεπτα. Στα λεπτά: $499:60\approx8{,}3$.\par
\apantisi Περίπου $8{,}3$ λεπτά.\par
\elegxos Σε $8{,}3$ λεπτά, δηλαδή περίπου $499$ δευτερόλεπτα, το φως διανύει $499\cdot300\,000\approx1{,}5\cdot10^8$ χλμ. Συμφωνεί.

\paradeigma{Γινόμενο σε τυποποιημένη μορφή}
\tixero Οι αριθμοί $2{,}5\cdot10^4$ και $4\cdot10^3$.\par
\tithelo Το γινόμενό τους σε τυποποιημένη μορφή.\par
\possindeo Πολλαπλασιάζω τους $\alpha$: $2{,}5\cdot4=10$. Προσθέτω τους εκθέτες: $4+3=7$. Έχω $10\cdot10^7$. Το $10$ δεν είναι μικρότερο του $10$, γι' αυτό γράφω $10\cdot10^7=10^8$.\par
\apantisi $1\cdot10^8$.\par
\elegxos $25\,000\cdot4\,000=100\,000\,000=10^8$.

\tmima{Ασκήσεις}

\zoni{Το κατάλαβα}{τυποποίηση, αποτυποποίηση, σύγκριση}

\ex Γράψε σε τυποποιημένη μορφή: α) $72\,000$ \quad β) $1\,250\,000$ \quad γ) $903\,000\,000$ \quad δ) $6\,000\,000\,000$ \quad ε) $48\,700$ \quad στ) $15\,000\,000\,000\,000$.

\ex Γράψε με δεκαδική μορφή: α) $2{,}7\cdot10^5$ \quad β) $5{,}06\cdot10^3$ \quad γ) $1{,}9\cdot10^7$ \quad δ) $8\cdot10^4$ \quad ε) $3{,}125\cdot10^6$ \quad στ) $4{,}008\cdot10^9$.

\ex Ποιοι από τους αριθμούς είναι σε τυποποιημένη μορφή; α) $5{,}2\cdot10^3$ \quad β) $52\cdot10^2$ \quad γ) $0{,}52\cdot10^4$ \quad δ) $10\cdot10^5$ \quad ε) $9{,}99\cdot10^1$.

\ex Γράψε σε τυποποιημένη μορφή: α) $24\cdot10^5$ \quad β) $0{,}35\cdot10^7$ \quad γ) $150\cdot10^3$ \quad δ) $0{,}008\cdot10^9$.

\ex Σύγκρινε με τα $<$, $=$, $>$: α) $8{,}2\cdot10^5$ και $9{,}1\cdot10^4$ \quad β) $3{,}5\cdot10^7$ και $2{,}9\cdot10^7$ \quad γ) $6\cdot10^9$ και $5{,}99\cdot10^9$ \quad δ) $4{,}4\cdot10^3$ και $4{,}4\cdot10^5$.

\ex Γράψε σε αύξουσα σειρά: $3{,}2\cdot10^6$, $9\cdot10^5$, $1{,}1\cdot10^7$, $4{,}5\cdot10^6$, $7{,}9\cdot10^5$.

\ex Η Πολυξένη βλέπει στην οθόνη του κομπιουτεράκι: α) $4.5\mathrm{E}7$ \quad β) $2.5\mathrm{E}10$. Ποιοι είναι οι αριθμοί; Πόσο δείχνει η οθόνη αν πολλαπλασιάσει τον $12\,000\,000$ με τον $3\,000\,000$;

\zoni{Το χρησιμοποιώ}{ηλιακό σύστημα και ανθρώπινο σώμα}

\ex \textbf{Σελήνη και Ήλιος.} Η Σελήνη απέχει $384\,400$ χλμ. από τη Γη και ο Ήλιος $149\,600\,000$ χλμ. Γράψε τις αποστάσεις σε τυποποιημένη μορφή και πες περίπου πόσες φορές μακρύτερα είναι ο Ήλιος.

\ex \textbf{Το έτος φωτός.} Ένα έτος φωτός είναι $9\,461\,000\,000\,000$ χλμ. Γράψε το σε τυποποιημένη μορφή. Ο κοντινότερος αστέρας, ο Εγγύτατος του Κενταύρου, απέχει περίπου $4{,}3$ έτη φωτός. Πόσα χιλιόμετρα είναι αυτά, σε τυποποιημένη μορφή, με δύο δεκαδικά;

\ex \textbf{Το φως του Ήλιου.} Η Φανή ρωτάει σε πόσο χρόνο φτάνει στη Γη το φως, αν ταξιδεύει με $3\cdot10^5$ χλμ. το δευτερόλεπτο και ο Ήλιος απέχει $1{,}496\cdot10^8$ χλμ.; Απάντησε σε δευτερόλεπτα και σε λεπτά.

\ex \textbf{Πληθυσμοί.} Ο πληθυσμός της Γης είναι περίπου $8\,000\,000\,000$ άνθρωποι και της Ελλάδας περίπου $10\,400\,000$. Γράψε τους σε τυποποιημένη μορφή και βρες περίπου πόσες φορές μεγαλύτερος είναι ο πρώτος.

\ex \textbf{Οι μάζες.} Η μάζα της Γης είναι $5{,}976\cdot10^{24}$ κιλά και του Ήλιου $1{,}983\cdot10^{30}$ κιλά. Γράψε τη μάζα της Γης σε δεκαδική μορφή και βρες περίπου πόσες φορές είναι ο Ήλιος βαρύτερος.

\ex \textbf{Η μνήμη του κινητού.} Η μνήμη του κινητού της Μυρτώς είναι $128$ GB, δηλαδή $1{,}28\cdot10^{11}$ bytes. Μια φωτογραφία έχει περίπου $4$ MB, δηλαδή $4\cdot10^6$ bytes. Πόσες φωτογραφίες χωράνε;

\zoni{Το σκέφτομαι}{οι κανόνες, όταν σπάνε}

\ex Γιατί ο $10\cdot10^5$ δεν είναι σε τυποποιημένη μορφή; Πώς γράφεται σωστά;

\ex Πόσα ψηφία έχει ο $4{,}56\cdot10^{10}$ όταν γραφεί κανονικά; Και ο $1\cdot10^{\nu}$; Εξήγησε πώς το βρίσκεις από τον εκθέτη.

\ex Υπολόγισε τα $(3\cdot10^4)\cdot(2\cdot10^5)$ και $(4\cdot10^3)\cdot(5\cdot10^6)$. Βάλε τα αποτελέσματα σε τυποποιημένη μορφή και διατύπωσε τον κανόνα του πολλαπλασιασμού.

\ex Ο Βασίλης λέει ότι $9{,}9\cdot10^3>1\cdot10^4$, γιατί $9{,}9>1$. Έχει δίκιο;

\ex Αν μετράμε αριθμούς φωναχτά, ένα νούμερο ανά δευτερόλεπτο, πόσα χρόνια χρειαζόμαστε για να φτάσουμε στο $10^{15}$; (Ένα έτος έχει περίπου $3{,}15\cdot10^7$ δευτερόλεπτα.)

\zoni{Για το σπίτι}{ψάξε τους μεγάλους αριθμούς του κόσμου}

\ex Βρες από αξιόπιστη πηγή τη μάζα του Ηλίου, το πλήθος των αστεριών του γαλαξία μας και την ακτίνα της Γης. Γράψε τους σε δεκαδική και σε τυποποιημένη μορφή.

\ex Φτιάξε πέντε δικούς σου αριθμούς, γράψε τους σε τυποποιημένη μορφή και ζήτησε από έναν συμμαθητή να βρει τους αρχικούς. Έλεγξε τις απαντήσεις του.

\paragrafos{Διαίρεση ρητών: πόσες φορές χωράει;}{Α.3.5, «Διαίρεση ρητών»}

\begin{anagke}
Ο παππούς Γιάννης έχει $450$ λίτρα ελαιόλαδο και θέλει να το βάλει σε δοχεία που χωράνε από $1\frac{1}{2}$ λίτρα. «Πόσα δοχεία θα χρειαστώ;» ρωτάει. Ο Ηλίας γράφει $450:1\frac{1}{2}$ και κολλάει: «Πώς διαιρώ με μεικτό;». Ο Στέλιος προσπαθεί να εφαρμόσει έναν κανόνα που μόλις έμαθε: «Διαιρώ τους αριθμητές μεταξύ τους και τους παρονομαστές μεταξύ τους». Τον σταματά η Ρήνη: «Σε ένα δοχείο χωράνε λίγο παραπάνω από ένα λίτρο. Τα δοχεία δεν μπορεί να είναι περισσότερα από $450$». Η Ζωή ρωτάει «πόσες φορές χωράει;». Μετράνε: κάθε $3$ λίτρα γεμίζουν $2$ δοχεία, άρα τα $450$ λίτρα γεμίζουν $300$ δοχεία.

Η Πολυξένη παρατηρεί ότι το $300$ προκύπτει και ως $450\cdot\frac{2}{3}$. «Το να διαιρέσεις με το $\frac{3}{2}$ είναι το ίδιο με το να πολλαπλασιάσεις με το $\frac{2}{3}$, τον αντίστροφό του;» Η Αγγελική προσθέτει ένα δεύτερο πρόβλημα: «Στο ταμείο της εκδρομής έλειπαν $24$ ευρώ. Αν χωριστούν σε $6$ παιδιά, πόσο έλειπε σε καθένα;» Και ο Νίκος: «Και αν ο διαιρέτης είναι δεκαδικός, π.χ. $3{,}6:0{,}45$;».

Σε αυτή την παράγραφο θα μάθουμε να διαιρούμε κάθε είδους ρητούς: ακεραίους, κλάσματα, δεκαδικούς και θετικούς και αρνητικούς.
\end{anagke}

\begin{lexi}{διαίρεση, πηλίκο, λόγος}{διαίρεση, πηλίκο, λόγος}
Η \lexiword{διαίρεση} είναι το «δι-αιρώ», «χωρίζω σε μέρη». Ο \lexiword{διαιρετέος} είναι αυτός που μοιράζεται, ο \lexiword{διαιρέτης} αυτός που μοιράζει και το \lexiword{πηλίκο} (από το «πηλίκος», «πόσος;») το αποτέλεσμα. Το πηλίκο $\alpha:\beta$ ή $\frac{\alpha}{\beta}$ λέγεται και \lexiword{λόγος} του $\alpha$ προς το $\beta$. Ένα κλάσμα που έχει μέσα του κλάσμα λέγεται \lexiword{σύνθετο}, «φτιαγμένο από πολλά κομμάτια».
\end{lexi}

\begin{metafrasi}
Πόσες φορές χωράει το $2$ στο $6$; & $6:2=3$ γιατί $3\cdot2=6$\\
Ομόσημοι. & πηλίκο θετικό\\
Ετερόσημοι. & πηλίκο αρνητικό\\
Διαιρώ και πολλαπλασιάζω τον αντίστροφο. & $\alpha:\beta=\alpha\cdot\dfrac{1}{\beta}$\\
Κλάσμα που έχει κλάσμα μέσα του. & $\dfrac{\frac{\alpha}{\beta}}{\frac{\gamma}{\delta}}$ (σύνθετο)\\
\end{metafrasi}

\tmima{Η θεωρία}

\begin{orismos}{πηλίκο ρητών}
Το πηλίκο της διαίρεσης $\alpha:\beta$ ή $\dfrac{\alpha}{\beta}$ ($\beta\neq0$) λέγεται και \textbf{λόγος} του $\alpha$ προς το $\beta$. Αν $\alpha:\beta=\gamma$, τότε ο $\gamma$ είναι ο μοναδικός αριθμός για τον οποίο ισχύει $\gamma\cdot\beta=\alpha$.
\end{orismos}

\begin{kanonas}{πρόσημο του πηλίκου}
Για να διαιρέσουμε δύο ρητούς, διαιρούμε τις \textbf{απόλυτες τιμές} τους και στο πηλίκο βάζουμε το πρόσημο $+$, αν είναι \textbf{ομόσημοι}, και το πρόσημο $-$, αν είναι \textbf{ετερόσημοι}.
\end{kanonas}

Για παράδειγμα $(+6):(+2)=+3$, γιατί $(+3)\cdot(+2)=+6$. Επίσης $(-8):(+4)=-2$, γιατί $(-2)\cdot(+4)=-8$, και $(-12):(-3)=+4$, γιατί $(+4)\cdot(-3)=-12$.

\begin{kanonas}{διαίρεση με τον αντίστροφο}
Για να \textbf{διαιρέσουμε} δύο ρητούς αρκεί να \textbf{πολλαπλασιάσουμε} τον διαιρετέο με τον \textbf{αντίστροφο} του διαιρέτη:
$$\alpha:\beta=\alpha\cdot\dfrac{1}{\beta}=\dfrac{\alpha}{\beta}\qquad\text{και}\qquad\dfrac{\alpha}{\beta}:\dfrac{\gamma}{\delta}=\dfrac{\alpha}{\beta}\cdot\dfrac{\delta}{\gamma}.$$
\end{kanonas}

Για παράδειγμα $\frac{2}{3}:\frac{4}{5}=\frac{2}{3}\cdot\frac{5}{4}=\frac{10}{12}=\frac{5}{6}$ και $\left(+\frac{2}{5}\right):\left(-\frac{1}{2}\right)=-\frac{2}{5}\cdot\frac{2}{1}=-\frac{4}{5}$.

\tmimasmall{Σύνθετα κλάσματα}

Ένα κλάσμα, του οποίου ένας τουλάχιστον όρος είναι κλάσμα, λέγεται \textbf{σύνθετο}. Η γραμμή του κλάσματος σημαίνει διαίρεση, άρα για να το κάνουμε \textbf{απλό} εκτελούμε τη διαίρεση:
$$\dfrac{\frac{\alpha}{\beta}}{\frac{\gamma}{\delta}}=\dfrac{\alpha}{\beta}:\dfrac{\gamma}{\delta}=\dfrac{\alpha\cdot\delta}{\beta\cdot\gamma}.$$
Αν ο παρονομαστής ή ο αριθμητής είναι ακέραιος, τον γράφουμε με παρονομαστή $1$. Αν υπάρχουν πρόσημα στους όρους, τα βγάζουμε έξω από τα κλάσματα: $\frac{-1}{2}$, $\frac{3}{-4}$ και $-\frac{3}{4}$ είναι ο ίδιος τύπος αριθμών.

\tmimasmall{Διαίρεση δεκαδικών}

Για να διαιρέσουμε δεκαδικό με δεκαδικό, πολλαπλασιάζουμε \textbf{και τον διαιρετέο και τον διαιρέτη} με την ίδια κατάλληλη δύναμη του $10$ ($10$, $100$, \dots) ώστε ο διαιρέτης να γίνει φυσικός αριθμός. Το πηλίκο δεν αλλάζει. Έτσι $4{,}2:1{,}75=420:175=2{,}4$.

\begin{kanonas}{διαίρεση με $10,\ 100,\ 1\,000,\dots$ και με $0{,}1,\ 0{,}01,\dots$}
Όταν διαιρούμε έναν δεκαδικό με $10,\ 100,\ 1\,000,\dots$ μεταφέρουμε την υποδιαστολή τόσες θέσεις \textbf{προς τα αριστερά} όσα είναι τα μηδενικά. Όταν τον διαιρούμε με $0{,}1,\ 0{,}01,\ 0{,}001,\dots$ μεταφέρουμε την υποδιαστολή τόσες θέσεις \textbf{προς τα δεξιά} όσα είναι τα δεκαδικά ψηφία του διαιρέτη.
\end{kanonas}

Για παράδειγμα $12{,}68:100=0{,}1268$ και $12{,}68:0{,}01=1\,268$.

\begin{prosoxi}
\begin{itemize}[leftmargin=1.2em,itemsep=0.3ex,topsep=2pt]
\item Δεν διαιρούμε ποτέ με το $0$. Αν $5:0=x$ τότε θα έπρεπε $x\cdot0=5$, που δεν ισχύει για κανέναν αριθμό. Το $0:5=0$ είναι σωστό.
\item Η διαίρεση δεν «μικραίνει» πάντα: $3:\frac{1}{2}=6>3$. Διαιρώντας με αριθμό μικρότερο του $1$, το πηλίκο μεγαλώνει.
\item Δεν αντιστρέφουμε και τα δύο κλάσματα: $\frac{2}{3}:\frac{4}{5}=\frac{2}{3}\cdot\frac{5}{4}$ και όχι $\frac{3}{2}\cdot\frac{5}{4}$. Αντιστρέφουμε μόνο τον \emph{διαιρέτη}.
\item Η διαίρεση δεν είναι αντιμεταθετική: $6:2=3$ ενώ $2:6=\frac{1}{3}$.
\item Στη διαίρεση με $0{,}1$ η υποδιαστολή πάει \emph{δεξιά} (το πηλίκο μεγαλώνει) και όχι αριστερά.
\end{itemize}
\end{prosoxi}

\tmima{Λυμένα παραδείγματα}

\paradeigma{Τα δοχεία του παππού Γιάννη}
\tixero $450$ λίτρα λάδι και δοχεία των $1\frac{1}{2}=\frac{3}{2}$ λίτρων.\par
\tithelo Πόσα δοχεία θα γεμίσουν;\par
\possindeo Διαιρώ: $450:\frac{3}{2}=450\cdot\frac{2}{3}=\frac{900}{3}=300$.\par
\apantisi $300$ δοχεία.\par
\elegxos $300\cdot\frac{3}{2}=450$.

\paradeigma{Πρόσημα στη διαίρεση}
\tixero Οι διαιρέσεις $(+6):(+2)$, $(-8):(+4)$, $(-12):(-3)$ και $\left(-\frac{2}{5}\right):\left(-\frac{3}{10}\right)$.\par
\tithelo Να βρω τα πηλίκα.\par
\possindeo Διαιρώ τις απόλυτες τιμές και βάζω πρόσημο: ομόσημοι $+$, ετερόσημοι $-$. Για το τελευταίο: $\frac{2}{5}\cdot\frac{10}{3}=\frac{20}{15}=\frac{4}{3}$.\par
\apantisi $+3$, $-2$, $+4$, $+\dfrac{4}{3}$.\par
\elegxos Για κάθε πηλίκο ο πολλαπλασιασμός με τον διαιρέτη δίνει τον διαιρετέο: $(-2)\cdot(+4)=-8$.

\paradeigma{Διαίρεση κλασμάτων}
\tixero Οι $\frac{2}{3}:\frac{4}{5}$ και $\left(+\frac{2}{5}\right):\left(-\frac{1}{2}\right)$.\par
\tithelo Να τις υπολογίσω.\par
\possindeo Πολλαπλασιάζω με τον αντίστροφο του διαιρέτη. Για το πρώτο ο αντίστροφος του $\frac{4}{5}$ είναι ο $\frac{5}{4}$. Για το δεύτερο ο αντίστροφος του $-\frac{1}{2}$ είναι ο $-2$.\par
\apantisi $\dfrac{2}{3}\cdot\dfrac{5}{4}=\dfrac{10}{12}=\dfrac{5}{6}$ και $\dfrac{2}{5}\cdot(-2)=-\dfrac{4}{5}$.\par
\elegxos $\frac{5}{6}\cdot\frac{4}{5}=\frac{20}{30}=\frac{2}{3}$.

\paradeigma{Σύνθετα κλάσματα}
\tixero Τα $\dfrac{\frac{3}{4}}{\frac{9}{8}}$, $\dfrac{\frac{5}{6}}{10}$ και $\dfrac{4}{\frac{2}{3}}$.\par
\tithelo Να τα γράψω ως απλά κλάσματα.\par
\possindeo Γράφω κάθε ένα ως διαίρεση: $\frac{3}{4}:\frac{9}{8}=\frac{3}{4}\cdot\frac{8}{9}=\frac{24}{36}=\frac{2}{3}$. Το δεύτερο: $\frac{5}{6}:10=\frac{5}{6}\cdot\frac{1}{10}=\frac{5}{60}=\frac{1}{12}$. Το τρίτο: $4:\frac{2}{3}=4\cdot\frac{3}{2}=6$.\par
\apantisi $\dfrac{2}{3}$, $\dfrac{1}{12}$, $6$.\par
\elegxos Στο τρίτο: $6\cdot\frac{2}{3}=4$.

\paradeigma{Παράσταση με σύνθετο κλάσμα}
\tixero Η παράσταση $\dfrac{\frac{3}{4}+\frac{1}{2}}{\frac{5}{6}-\frac{1}{3}}$.\par
\tithelo Να βρω την τιμή της.\par
\possindeo Εκτελώ τις πράξεις στον αριθμητή και στον παρονομαστή: $\frac{3}{4}+\frac{1}{2}=\frac{5}{4}$ και $\frac{5}{6}-\frac{1}{3}=\frac{5}{6}-\frac{2}{6}=\frac{3}{6}=\frac{1}{2}$. Τότε $\frac{5}{4}:\frac{1}{2}=\frac{5}{4}\cdot2=\frac{10}{4}$.\par
\apantisi $\dfrac{5}{2}$.\par
\elegxos Με δεκαδικούς: $1{,}25:0{,}5=2{,}5$.

\paradeigma{Διαίρεση δεκαδικών}
\tixero Η διαίρεση $3{,}6:0{,}45$ (κιλά αλεύρι προς κιλά ανά κέικ).\par
\tithelo Πόσα κέικ βγαίνουν;\par
\possindeo Πολλαπλασιάζω και τους δύο αριθμούς επί $100$, ώστε ο διαιρέτης να γίνει φυσικός: $360:45=8$.\par
\apantisi $8$ κέικ.\par
\elegxos $8\cdot0{,}45=3{,}6$.

\paradeigma{Διαίρεση με $10$, $100$, $0{,}1$, $0{,}01$}
\tixero Ο αριθμός $12{,}68$.\par
\tithelo Να τον διαιρέσω με $10$, $100$, $0{,}1$, $0{,}01$.\par
\possindeo Με $10$ και $100$ η υποδιαστολή πάει αριστερά κατά $1$ και $2$ θέσεις. Με $0{,}1$ και $0{,}01$ πάει δεξιά κατά $1$ και $2$ θέσεις.\par
\apantisi $1{,}268$, \ $0{,}1268$, \ $126{,}8$, \ $1\,268$.\par
\elegxos Το $12{,}68:0{,}1$ είναι «πόσα δέκατα χωράνε στο $12{,}68$», δηλαδή $126{,}8$.

\tmima{Ασκήσεις}

\zoni{Το κατάλαβα}{κανόνες και πράξεις}

\ex Κάνε τις διαιρέσεις: α) $15:(-5)$ \quad β) $(-60):(+15)$ \quad γ) $(+81):(-9)$ \quad δ) $(-490):(-7)$.

\ex Βρες τα πηλίκα: α) $\frac{2}{5}:\frac{4}{15}$ \quad β) $\left(+\frac{1}{2}\right):\left(-\frac{1}{3}\right)$ \quad γ) $4:\frac{1}{2}$ \quad δ) $\left(-\frac{4}{3}\right):(-8)$ \quad ε) $\frac{18}{5}:\frac{6}{5}$ \quad στ) $\frac{12}{21}:\frac{7}{4}$.

\ex Μετάτρεψε τα σύνθετα κλάσματα σε απλά: α) $\dfrac{\frac{7}{11}}{\frac{3}{22}}$ \quad β) $\dfrac{\frac{27}{49}}{\frac{9}{14}}$ \quad γ) $\dfrac{4}{\frac{56}{9}}$ \quad δ) $\dfrac{\frac{11}{2}}{121}$ \quad ε) $\dfrac{\frac{3}{4}}{2\frac{1}{8}}$.

\ex Υπολόγισε τα πηλίκα: α) $4{,}5:0{,}15$ \quad β) $28{,}42:49$ \quad γ) $35{,}76:1{,}5$ \quad δ) $292{,}62:29{,}262$ \quad ε) $0{,}48:0{,}6$.

\ex Υπολόγισε χωρίς κάθετη πράξη: α) $45{,}7:10$ \quad β) $3{,}2:100$ \quad γ) $0{,}5:1\,000$ \quad δ) $6{,}4:0{,}1$ \quad ε) $0{,}35:0{,}01$ \quad στ) $12:0{,}001$.

\ex Σωστό ή λάθος; α) $\alpha:\beta=\alpha\cdot\frac{1}{\beta}$. \quad β) Επιτρέπεται η διαίρεση με το $0$. \quad γ) Το πηλίκο ετεροσήμων είναι αρνητικό. \quad δ) $(-12):(-4)=-3$. \quad ε) $3:\frac{1}{3}=1$. \quad στ) Το πηλίκο δύο ρητών είναι πάντα μικρότερο από τον διαιρετέο.

\ex Υπολόγισε: α) $\dfrac{-3}{4}-\dfrac{5}{-8}+\dfrac{-1}{-2}$ \quad β) $1+\dfrac{5}{-4}-\dfrac{-1}{8}$.

\ex Βρες τον αριθμό $x$ ώστε: α) $(-3)\cdot x=12$ \quad β) $x\cdot\frac{2}{3}=4$ \quad γ) $(-5)\cdot x=-\frac{1}{2}$.

\zoni{Το χρησιμοποιώ}{μοιράζω και μετράω}

\ex \textbf{Το λάδι.} Ο παππούς Γιάννης έχει $6$ λίτρα λάδι και θέλει να το βάλει σε μπουκάλια των $\frac{3}{4}$ του λίτρου. Πόσα μπουκάλια θα γεμίσουν;

\ex \textbf{Η κορδέλα.} Η Αγγελική κόβει μια κορδέλα $4{,}5$ μ. σε κομμάτια των $0{,}15$ μ. Πόσα κομμάτια βγαίνουν;

\ex \textbf{Η τιμή.} Η γιαγιά Μυρτώ πλήρωσε $6$ € για $0{,}75$ κιλά τυρί. Πόσο κοστίζει το κιλό;

\ex \textbf{Το χρέος του ταμείου.} Το ταμείο της εκδρομής έχει υπόλοιπο $-24{,}60$ €. Τα $6$ παιδιά της παρέας μοιράζονται εξίσου το χρέος. Τι αριθμός εκφράζει το μερίδιο του καθενός;

\ex \textbf{Η ταχύτητα.} Ο Νίκος τρέχει $\frac{3}{5}$ του χιλιομέτρου σε $\frac{1}{10}$ της ώρας. Πόσα χιλιόμετρα τρέχει σε μία ώρα (αν διατηρεί τον ίδιο ρυθμό);

\zoni{Το σκέφτομαι}{γιατί το πηλίκο μεγαλώνει;}

\ex Ο Βασίλης λέει ότι η διαίρεση πάντα μικραίνει τον αριθμό. Βρες αντιπαράδειγμα και εξήγησε πότε το πηλίκο είναι μεγαλύτερο από τον διαιρετέο.

\ex Γιατί δεν επιτρέπεται η διαίρεση με το $0$; Σκέψου τι θα σήμαινε $5:0=x$ αν υπήρχε τέτοιος αριθμός.

\ex Ποιο είναι μεγαλύτερο: το $3:\frac{1}{2}$ ή το $3:\frac{1}{3}$; Εξήγησε με τη φράση «πόσες φορές χωράει».

\ex Αν $\frac{\alpha}{\beta}>0$, τι πρόσημα μπορούν να έχουν οι $\alpha$ και $\beta$; Αν $\frac{\alpha}{\beta}<0$;

\zoni{Για το σπίτι}{κόβω και μετράω}

\ex Πάρε ένα σπάγκο μήκους $1{,}2$ μ. και κόψε τον σε κομμάτια των $0{,}15$ μ. Πρόβλεψε πρώτα με διαίρεση πόσα θα βγουν και μετά έλεγξε.

\ex Γράψε τρία σύνθετα κλάσματα της επιλογής σου και μετάτρεψέ τα σε απλά.

\paragrafos{Δυνάμεις ρητών: οι πολλαπλασιασμοί που επαναλαμβάνονται}{Α.3.6, «Δυνάμεις ρητών με εκθέτη μη μηδενικό φυσικό»}

\begin{anagke}
Στην αυλή του σχολείου ο Νίκος αφήνει μια μπάλα του μπάσκετ από ύψος $2$ μέτρων. Κάθε φορά που αναπηδά, φτάνει στο μισό του ύψους της προηγούμενης αναπήδησης. «Σε τι ύψος θα φτάσει στην τρίτη αναπήδηση;» ρωτάει η Πολυξένη. Ο Ηλίας γράφει: μετά την πρώτη $2\cdot\frac{1}{2}$, μετά τη δεύτερη $2\cdot\frac{1}{2}\cdot\frac{1}{2}$, μετά την τρίτη $2\cdot\frac{1}{2}\cdot\frac{1}{2}\cdot\frac{1}{2}$. «Αυτό γράφεται πιο σύντομα», λέει η Ρήνη: «το $\frac{1}{2}$ τρεις φορές». Στους φυσικούς το ξέρουν ήδη: $2\cdot2\cdot2=2^3$. Μπορεί να γίνει και με κλάσματα;

Η κυρία Αλεξάνδρα βάζει στον πίνακα δύο παραστάσεις: $\left(-\frac{1}{2}\right)^3$ και $-\frac{1}{2^3}$. «Είναι ίδιες;» ρωτά. Ο Στέλιος λέει «ναι, όλα είναι μείον ένα όγδοο» και η Φανή γράφει στο τετράδιό της: «ίδια, αλλά γιατί;». Ο Βασίλης προτείνει $(-3)^2=-9$ και η Ζωή τον ρωτάει «πώς βγαίνει αρνητικό το γινόμενο δύο αρνητικών;»

Σε αυτή την παράγραφο θα επεκτείνουμε τις δυνάμεις στους ρητούς: θα μάθουμε να προσδιορίζουμε το πρόσημο μιας δύναμης και να υπολογίζουμε παραστάσεις με δυνάμεις.
\end{anagke}

\begin{lexi}{δύναμη, βάση, εκθέτης}{δύναμη, βάση, εκθέτης}
Η \lexiword{δύναμη} είναι η «ύψωση», το γινόμενο ίσων παραγόντων. Στα αρχαία ελληνικά «δύναμις» ήταν και το τετράγωνο. Η \lexiword{βάση} είναι ο αριθμός που πολλαπλασιάζεται με τον εαυτό του και ο \lexiword{εκθέτης} (από το «εκ-τίθημι», «βάζω έξω») ο αριθμός που γράφουμε πάνω δεξιά και δείχνει πόσους παράγοντες έχουμε.
\end{lexi}

\begin{metafrasi}
Το $\frac{2}{3}$ τρεις φορές. & $\left(\dfrac{2}{3}\right)^{3}=\dfrac{2}{3}\cdot\dfrac{2}{3}\cdot\dfrac{2}{3}=\dfrac{8}{27}$\\
Δύο αρνητικοί παράγοντες. & $(-1{,}5)^2=+2{,}25$\\
Τρεις αρνητικοί παράγοντες. & $(-0{,}1)^3=-0{,}001$\\
Η δύναμη με εκθέτη $1$. & $\alpha^1=\alpha$\\
Το μείον μόνο, χωρίς παρένθεση. & $-2^4=-16$\\
\end{metafrasi}

\tmima{Η θεωρία}

\begin{orismos}{δύναμη ρητού}
Η \textbf{δύναμη} του $\alpha$ στην $\nu$ ($\nu$ θετικός φυσικός) είναι το γινόμενο $\nu$ παραγόντων ίσων με $\alpha$:
$$\alpha^{\nu}=\underbrace{\alpha\cdot\alpha\cdot\ldots\cdot\alpha}_{\nu\ \text{παράγοντες}}.$$
Ο $\alpha$ λέγεται \textbf{βάση} (είναι ρητός) και ο $\nu$ \textbf{εκθέτης}. Ισχύει $\alpha^1=\alpha$.
\end{orismos}

Για παράδειγμα $\left(\frac{2}{3}\right)^3=\frac{2}{3}\cdot\frac{2}{3}\cdot\frac{2}{3}=\frac{2^3}{3^3}=\frac{8}{27}$. Γενικά, αφού το γινόμενο κλασμάτων έχει αριθμητή το γινόμενο των αριθμητών και παρονομαστή το γινόμενο των παρονομαστών, ισχύει $\left(\dfrac{\alpha}{\beta}\right)^{\nu}=\dfrac{\alpha^{\nu}}{\beta^{\nu}}$.

\begin{kanonas}{το πρόσημο μιας δύναμης}
\begin{itemize}[leftmargin=1.2em,itemsep=0.3ex,topsep=2pt]
\item Δύναμη με βάση \textbf{θετικό} αριθμό είναι πάντα \textbf{θετική}: αν $\alpha>0$, τότε $\alpha^{\nu}>0$.
\item Δύναμη με βάση \textbf{αρνητικό} αριθμό και εκθέτη \textbf{άρτιο} είναι \textbf{θετική}: αν $\alpha<0$ και $\nu$ άρτιος, τότε $\alpha^{\nu}>0$. Π.χ. $(-1{,}5)^2=(-1{,}5)\cdot(-1{,}5)=+2{,}25$.
\item Δύναμη με βάση \textbf{αρνητικό} αριθμό και εκθέτη \textbf{περιττό} είναι \textbf{αρνητική}: αν $\alpha<0$ και $\nu$ περιττός, τότε $\alpha^{\nu}<0$. Π.χ. $(-0{,}1)^3=-0{,}001$.
\end{itemize}
\end{kanonas}

Ο λόγος είναι ο κανόνας των προσήμων του γινομένου: κάθε ζευγάρι αρνητικών παραγόντων δίνει θετικό, άρα το πρόσημο εξαρτάται από το αν ο αριθμός των αρνητικών παραγόντων είναι άρτιος ή περιττός. Ειδικά $(-1)^{\nu}=1$ αν ο $\nu$ είναι άρτιος και $-1$ αν είναι περιττός. Επίσης $1^{\nu}=1$ και $0^{\nu}=0$.

\tmimasmall{Προτεραιότητα των πράξεων}

Στις παραστάσεις με πράξεις εκτελούμε με τη σειρά: πρώτα τις πράξεις μέσα σε \textbf{παρενθέσεις}, μετά τις \textbf{δυνάμεις}, μετά τους \textbf{πολλαπλασιασμούς και τις διαιρέσεις} (από αριστερά προς τα δεξιά) και τέλος τις \textbf{προσθέσεις και αφαιρέσεις}.

\begin{prosoxi}
\begin{itemize}[leftmargin=1.2em,itemsep=0.3ex,topsep=2pt]
\item Το $\alpha^{3}$ δεν είναι $3\cdot\alpha$, αλλά $\alpha\cdot\alpha\cdot\alpha$. Για παράδειγμα $0{,}2^4\neq4\cdot0{,}2$: $0{,}2^4=0{,}0016$.
\item Το $(-3)^2$ και το $-3^2$ δεν είναι ίδια. Το πρώτο είναι $(-3)(-3)=9$ και το δεύτερο είναι ο αντίθετος του $3^2$, δηλαδή $-9$. Η βάση είναι ό,τι βρίσκεται μέσα στην παρένθεση.
\item Στο $\left(\frac{2}{3}\right)^3$ ο εκθέτης δουλεύει και στον αριθμητή και στον παρονομαστή: $\frac{2^3}{3^3}=\frac{8}{27}$ και όχι $\frac{2^3}{3}$.
\item Ο εκθέτης δείχνει πλήθος παραγόντων, όχι πρόσθεση: $2^3=2\cdot2\cdot2=8$, όχι $2+2+2=6$.
\item Η δύναμη με άρτιο εκθέτη δεν είναι ποτέ αρνητική, ακόμη και με αρνητική βάση.
\end{itemize}
\end{prosoxi}

\tmima{Λυμένα παραδείγματα}

\paradeigma{Η μπάλα του Νίκου}
\tixero Η μπάλα αφήνεται από ύψος $2$ μ. και κάθε αναπήδηση φτάνει στο μισό της προηγούμενης.\par
\tithelo Σε τι ύψος φτάνει στην τρίτη αναπήδηση;\par
\possindeo Κάθε αναπήδηση πολλαπλασιάζει το ύψος επί $\frac{1}{2}$. Μετά από τρεις: $2\cdot\left(\frac{1}{2}\right)^3=2\cdot\frac{1}{8}=\frac{2}{8}=\frac{1}{4}$.\par
\apantisi $\frac{1}{4}$ μ., δηλαδή $25$ εκατοστά.\par
\elegxos Βήμα βήμα: $2\to1\to\frac{1}{2}\to\frac{1}{4}$.

\paradeigma{Δυνάμεις με πρόσημο}
\tixero Οι δυνάμεις $\left(\frac{2}{3}\right)^3$, $(-1{,}5)^2$, $(-0{,}1)^3$.\par
\tithelo Να τις υπολογίσω.\par
\possindeo Πρώτα το πρόσημο: $\frac{2}{3}$ θετική βάση, άρα θετικό. Το $(-1{,}5)^2$ έχει αρνητική βάση και άρτιο εκθέτη, άρα θετικό. Το $(-0{,}1)^3$ έχει αρνητική βάση και περιττό εκθέτη, άρα αρνητικό. Μετά οι απόλυτες τιμές: $\frac{8}{27}$, $1{,}5\cdot1{,}5=2{,}25$, $0{,}1\cdot0{,}1\cdot0{,}1=0{,}001$.\par
\apantisi $\dfrac{8}{27}$, $\ 2{,}25$, $\ -0{,}001$.\par
\elegxos Το $0{,}001$ έχει τρία δεκαδικά ψηφία: $1+1+1$.

\paradeigma{Η παρένθεση μετράει}
\tixero Οι $(-2)^4$, $-2^4$, $\left(-\frac{1}{2}\right)^3$ και $-\frac{1}{2^3}$.\par
\tithelo Να τις υπολογίσω και να τις συγκρίνω.\par
\possindeo Το $(-2)^4$ έχει αρνητική βάση και άρτιο εκθέτη: $(-2)\cdot(-2)\cdot(-2)\cdot(-2)=16$. Το $-2^4$ είναι ο αντίθετος του $2^4$: $-16$. Το $\left(-\frac{1}{2}\right)^3=-\frac{1}{8}$. Το $-\frac{1}{2^3}=-\frac{1}{8}$.\par
\apantisi $16$, $-16$, $-\dfrac{1}{8}$, $-\dfrac{1}{8}$. Οι δύο τελευταίες είναι ίσες.\par
\elegxos Το $(-2)^4$ είναι θετικό, ενώ το $-2^4$ αρνητικό: διαφέρουν.

\paradeigma{Η ειδική περίπτωση του $-1$}
\tixero Οι δυνάμεις $(-1)^{100}$, $(-1)^{101}$, $1^{50}$, $0^3$.\par
\tithelo Να τις υπολογίσω χωρίς να γράψω όλους τους παράγοντες.\par
\possindeo Το $(-1)^{\nu}$ είναι $1$ για άρτιο και $-1$ για περιττό $\nu$. Το $1^{\nu}=1$ και $0^{\nu}=0$.\par
\apantisi $1$, $-1$, $1$, $0$.\par
\elegxos Σε ζεύγη: $(-1)^{100}=\big((-1)(-1)\big)^{50}=1^{50}=1$.

\paradeigma{Παράσταση με προτεραιότητα}
\tixero Η παράσταση $3\cdot\left(\frac{1}{2}-\frac{5}{6}\right)+\left(\frac{1}{3}-\frac{2}{3}\right)^2$.\par
\tithelo Να βρω την τιμή της.\par
\possindeo Πρώτα οι παρενθέσεις: $\frac{1}{2}-\frac{5}{6}=\frac{3}{6}-\frac{5}{6}=-\frac{1}{3}$ και $\frac{1}{3}-\frac{2}{3}=-\frac{1}{3}$. Μετά η δύναμη: $\left(-\frac{1}{3}\right)^2=\frac{1}{9}$. Μετά ο πολλαπλασιασμός: $3\cdot\left(-\frac{1}{3}\right)=-1$. Τέλος η πρόσθεση: $-1+\frac{1}{9}=-\frac{9}{9}+\frac{1}{9}=-\frac{8}{9}$.\par
\apantisi $-\dfrac{8}{9}$.\par
\elegxos Με δεκαδικούς: $3\cdot(-0{,}333\ldots)+0{,}111\ldots=-0{,}888\ldots$

\tmima{Ασκήσεις}

\zoni{Το κατάλαβα}{ορισμός, πρόσημο, υπολογισμός}

\ex Σωστό ή λάθος; α) $\left(\frac{5}{2}\right)^2=\frac{5^2}{2^2}$. \quad β) $\left(\frac{2}{3}\right)^3=\frac{2^3}{3}$. \quad γ) $0{,}2^4=4\cdot0{,}2$. \quad δ) $\left(-\frac{3}{10}\right)^2=0{,}3^2$. \quad ε) $\left(-\frac{1}{3}\right)^3=-\left(\frac{1}{3}\right)^3$.

\ex Υπολόγισε: α) $\left(\frac{1}{2}\right)^2$ \quad β) $\left(-\frac{1}{3}\right)^3$ \quad γ) $\left(\frac{2}{5}\right)^2$ \quad δ) $(-0{,}2)^2$ \quad ε) $0{,}3^3$ \quad στ) $(-1)^7$.

\ex Βρες το πρόσημο χωρίς να υπολογίσεις: α) $\left(-\frac{3}{4}\right)^5$ \quad β) $(-0{,}2)^4$ \quad γ) $\left(-\frac{1}{2}\right)^{10}$ \quad δ) $-\left(\frac{3}{5}\right)^3$ \quad ε) $(-7)^{11}$.

\ex Γράψε ως δύναμη: α) $\frac{3}{5}\cdot\frac{3}{5}\cdot\frac{3}{5}$ \quad β) $(-2)(-2)(-2)(-2)$ \quad γ) $0{,}1\cdot0{,}1$ \quad δ) $\left(-\frac{1}{2}\right)\cdot\left(-\frac{1}{2}\right)$.

\ex Υπολόγισε και σύγκρινε: α) $(-2)^4$ και $-2^4$ \quad β) $(-3)^3$ και $-3^3$ \quad γ) $\left(-\frac{1}{2}\right)^2$ και $-\left(\frac{1}{2}\right)^2$.

\ex Υπολόγισε: α) $1^{100}$ \quad β) $0^5$ \quad γ) $(-1)^{100}$ \quad δ) $(-1)^{101}$ \quad ε) $10^3$ \quad στ) $(0{,}1)^2$.

\ex Κάνε τις πράξεις: α) $\left(\frac{1}{2}\right)^2+\left(-\frac{1}{2}\right)^3+\frac{1}{4^2}$ \quad β) $0{,}1^2+\left(\frac{1}{10}\right)^2-0{,}2^2$ \quad γ) $-1^{10}+1^{100}+(-1)^2+(-1)^3$.

\zoni{Το χρησιμοποιώ}{μπάλες, διπλώσεις, εμβαδά}

\ex \textbf{Η μπάλα.} Ο Νίκος αφήνει μια μπάλα από ύψος $2$ μ. και κάθε αναπήδηση φτάνει στο μισό της προηγούμενης. Σε τι ύψος φτάνει μετά την τρίτη και μετά την τέταρτη αναπήδηση;

\ex \textbf{Οι διπλώσεις.} Η Αγγελική διπλώνει ένα φύλλο χαρτί στη μέση πέντε φορές. Τι μέρος του αρχικού εμβαδού έχει το διπλωμένο φύλλο; Γράψε το ως δύναμη και ως κλάσμα.

\ex \textbf{Το τετράγωνο.} Μια αυλή είναι τετράγωνο με πλευρά $1{,}5$ μ. Πόσο είναι το εμβαδόν της; Ο Ηλίας λέει $1{,}5\cdot2=3$. Τι λάθος έκανε;

\ex \textbf{Η παράσταση.} Υπολόγισε $8\cdot\left(-\frac{1}{2}\right)^3+(0{,}5+0{,}3)^2$.

\ex \textbf{Το μοτίβο.} Υπολόγισε τα $\left(\frac{1}{10}\right)^1$, $\left(\frac{1}{10}\right)^2$, $\left(\frac{1}{10}\right)^3$, $\left(\frac{1}{10}\right)^4$ ως δεκαδικούς. Τι παρατηρείς;

\zoni{Το σκέφτομαι}{γιατί βγαίνει έτσι;}

\ex Ο Βασίλης υπολόγισε $-3^2=9$. Εξήγησε το λάθος του και δείξε ότι $(-3)^2=9$.

\ex Πότε είναι θετική μια δύναμη με αρνητική βάση; Εξήγησε με τον κανόνα των προσήμων.

\ex Ποιο είναι μεγαλύτερο, το $\left(\frac{1}{2}\right)^3$ ή το $\left(\frac{1}{2}\right)^2$; Και στη βάση $2$; Τι παρατηρείς για βάσεις μικρότερες και μεγαλύτερες από το $1$;

\ex Αν $\alpha=-\frac{1}{2}$, ποιο είναι μεγαλύτερο, το $\alpha^2$ ή το $\alpha^3$; Εξήγησε χωρίς πράξεις.

\zoni{Για το σπίτι}{πίνακες δυνάμεων}

\ex Φτιάξε πίνακα με τις δυνάμεις $\left(\frac{1}{2}\right)^{\nu}$ για $\nu=1,2,\dots,6$, με κλάσμα και δεκαδικό. Τι παρατηρείς;

\ex Φτιάξε τρεις παραστάσεις με ρητούς και δυνάμεις, που να χρειάζονται και παρενθέσεις, και υπολόγισέ τες.

\paragrafos{Τυποποιημένη μορφή: οι μεγάλοι αριθμοί χωράνε σε μια γραμμή}{Α.3.7, «Τυποποιημένη μορφή μεγάλων αριθμών»}

\begin{anagke}
Η κυρία Αλεξάνδρα ζητάει από την Α3 μια εργασία για το ηλιακό σύστημα. Η Αγγελική γράφει στο τετράδιό της την απόσταση της Γης από τον Ήλιο: $149\,600\,000\,000$ μέτρα. Ο Ηλίας γράφει την ηλικία της Γης: $4\,500\,000\,000$ έτη, και την ξαναγράφει τρεις φορές, γιατί δεν είναι σίγουρος για τα μηδενικά. Ο Στέλιος δοκιμάζει την ηλικία του Γαλαξία και χάνεται στο δεύτερο μέτρημα. «Αυτό δεν διαβάζεται», παραπονιέται η Ρήνη.

Η Πολυξένη, που έχει το κομπιουτεράκι της, πατάει $149\,600\,000\cdot1\,000$ και η οθόνη δείχνει $1{,}496\mathrm{E}11$. «Τι είναι αυτό το E11;». Η κυρία Αλεξάνδρα χαμογελάει: «Το κομπιουτεράκι δεν χωράει όλα τα μηδενικά και τα γράφει με συντομία: $1{,}496$ επί δέκα στην ενδέκατη. Κάνουν ό,τι κάνουν και οι επιστήμονες». Η Φανή ρωτάει: «Μπορούμε να γράφουμε κι εμείς έτσι, ώστε να μη χάνουμε τα μηδενικά;».

Σε αυτή την παράγραφο θα μάθουμε να γράφουμε τους πολύ μεγάλους αριθμούς με σύντομο και ενιαίο τρόπο, τον οποίο χρησιμοποιούν οι επιστήμονες.
\end{anagke}

\begin{lexi}{τυποποιημένη μορφή}{τυποποιημένη μορφή}
Ο \lexiword{τύπος} σημαίνει καλούπι, πρότυπο. \lexiword{Τυποποιημένη} είναι η μορφή που έχει το ίδιο καλούπι για όλους τους αριθμούς: ένας αριθμός από το $1$ μέχρι κάτω από το $10$, επί μια δύναμη του $10$. Οι επιστήμονες την προτιμούν γιατί φαίνεται αμέσως το \emph{μέγεθος} του αριθμού (πόσα μηδενικά) και τα σημαντικά του ψηφία.
\end{lexi}

\begin{metafrasi}
Ένα ακολουθούμενο από $7$ μηδενικά. & $10^7=10\,000\,000$\\
Τριάντα εκατομμύρια. & $30\,000\,000=3\cdot10^7$\\
Η απόσταση Γης -- Ήλιου σε χιλιόμετρα. & $149\,600\,000=1{,}496\cdot10^8$\\
Στην οθόνη του κομπιουτεράκι. & $1{,}496\mathrm{E}8$\\
\end{metafrasi}

\tmima{Η θεωρία}

Θυμόμαστε ότι η δύναμη $10^{\nu}$ είναι ο αριθμός που γράφεται με το $1$ και μετά $\nu$ μηδενικά: $10^3=1\,000$, $10^6=1\,000\,000$.

\begin{orismos}{τυποποιημένη μορφή}
Ένας «μεγάλος» αριθμός μπορεί να γραφεί στη μορφή $\alpha\cdot10^{\nu}$, δηλαδή ως γινόμενο ενός αριθμού $\alpha$ επί μια δύναμη του $10$, όπου ο $\nu$ είναι ένας μη μηδενικός φυσικός αριθμός. Η μορφή αυτή λέγεται \textbf{τυποποιημένη}. Ο $\alpha$ είναι ένας δεκαδικός αριθμός με ακέραιο μέρος ένα ψηφίο μεγαλύτερο του $0$, δηλαδή $1\le\alpha<10$.
\end{orismos}

Η τυποποιημένη μορφή μας δίνει τη δυνατότητα να γράψουμε έναν «μεγάλο» αριθμό με απλούστερο τρόπο.

\begin{kanonas}{από τον αριθμό στην τυποποιημένη μορφή}
\begin{enumerate}[leftmargin=1.6em,itemsep=0.2ex,topsep=2pt]
\item Βάζω την υποδιαστολή μετά το \emph{πρώτο ψηφίο} (αριστερά), ώστε να προκύψει αριθμός $\alpha$ από $1$ έως $9{,}99\ldots$
\item Μετράω πόσες θέσεις μετακίνησα την υποδιαστολή προς τα αριστερά. Αυτός είναι ο εκθέτης $\nu$.
\item Γράφω $\alpha\cdot10^{\nu}$ (και αφαιρώ τα μηδενικά που περισσεύουν στο τέλος του $\alpha$).
\end{enumerate}
Για παράδειγμα $815\,000\,000\,000=8{,}15\cdot10^{11}$, γιατί η υποδιαστολή μετακινήθηκε $11$ θέσεις προς τα αριστερά. Ο εκθέτης είναι πάντα κατά $1$ μικρότερος από το πλήθος των ψηφίων του αριθμού.
\end{kanonas}

\begin{kanonas}{από την τυποποιημένη μορφή στον αριθμό}
Για να γράψω τον $\alpha\cdot10^{\nu}$ ως συνηθισμένο αριθμό, μεταφέρω την υποδιαστολή του $\alpha$ \textbf{$\nu$ θέσεις προς τα δεξιά}, συμπληρώνοντας με μηδενικά όπου χρειάζεται. Έτσι $3{,}42\cdot10^9=3\,420\,000\,000$.
\end{kanonas}

Η τυποποιημένη μορφή βοηθά και στη σύγκριση. Από δύο αριθμούς σε τυποποιημένη μορφή μεγαλύτερος είναι αυτός με τον \textbf{μεγαλύτερο εκθέτη}. Αν έχουν τον ίδιο εκθέτη, συγκρίνουμε τους $\alpha$. Έτσι $2\cdot10^{5}>9{,}9\cdot10^{4}$ και $3{,}2\cdot10^{6}<4{,}1\cdot10^{6}$.

\begin{prosoxi}
\begin{itemize}[leftmargin=1.2em,itemsep=0.3ex,topsep=2pt]
\item Το $12\cdot10^{5}$ δεν είναι τυποποιημένη μορφή, γιατί το $12$ δεν είναι μικρότερο από το $10$. Σωστά: $1{,}2\cdot10^{6}$.
\item Το $0{,}5\cdot10^{3}$ δεν είναι τυποποιημένη, γιατί το ακέραιο μέρος είναι $0$. Σωστά: $5\cdot10^{2}$.
\item Το $10^{\nu}$ δεν είναι $10\cdot\nu$. Το $10^{3}=1\,000$ και όχι $30$.
\item Ο εκθέτης δεν είναι το πλήθος των ψηφίων αλλά το πλήθος των μηδενικών μετά το $1$, δηλαδή τα ψηφία μείον $1$: ο $5\,000\,000$ έχει $7$ ψηφία και είναι $5\cdot10^{6}$.
\item Τα μηδενικά στο τέλος του $\alpha$ δεν γράφονται: $2{,}50\cdot10^4=2{,}5\cdot10^4$.
\end{itemize}
\end{prosoxi}

\tmima{Λυμένα παραδείγματα}

\paradeigma{Στην τυποποιημένη μορφή}
\tixero Οι αριθμοί $30\,000\,000$ και $815\,000\,000\,000$.\par
\tithelo Να τους γράψω σε τυποποιημένη μορφή.\par
\possindeo Ο πρώτος έχει $8$ ψηφία: μετακινώ την υποδιαστολή κατά $7$ θέσεις και παίρνω $3$. Ο δεύτερος έχει $12$ ψηφία: κατά $11$ θέσεις, και παίρνω $8{,}15$.\par
\apantisi $30\,000\,000=3\cdot10^{7}$ και $815\,000\,000\,000=8{,}15\cdot10^{11}$.\par
\elegxos $3\cdot10^{7}=3\cdot10\,000\,000$ και $8{,}15\cdot10^{11}$: η υποδιαστολή πάει $11$ θέσεις δεξιά, και παίρνουμε το $815$ και $9$ μηδενικά.

\paradeigma{Από την τυποποιημένη μορφή}
\tixero Ο αριθμός $3{,}42\cdot10^{9}$.\par
\tithelo Να τον γράψω ως συνηθισμένο αριθμό.\par
\possindeo Μετακινώ την υποδιαστολή κατά $9$ θέσεις δεξιά. Το $3{,}42$ έχει δύο δεκαδικά ψηφία, οπότε προσθέτω $7$ μηδενικά.\par
\apantisi $3\,420\,000\,000$.\par
\elegxos Ο αριθμός έχει $10$ ψηφία, ενώ ο εκθέτης είναι $9$.

\paradeigma{Σύγκριση}
\tixero Οι αριθμοί $2\cdot10^{3}$ και $3\cdot10^{2}$, και οι $0{,}1\cdot10^{2}$ και $10$.\par
\tithelo Να τους συγκρίνω.\par
\possindeo Στο πρώτο ζεύγος οι εκθέτες διαφέρουν: $2\,000>300$. Το δεύτερο ζεύγος: το $0{,}1\cdot10^{2}$ δεν είναι τυποποιημένο. Είναι $0{,}1\cdot100=10$.\par
\apantisi $2\cdot10^{3}>3\cdot10^{2}$ και $0{,}1\cdot10^{2}=10$.\par
\elegxos $2\cdot10^3=2\,000$ και $3\cdot10^2=300$.

\paradeigma{Χρόνος σε λεπτά}
\tixero Οι καρχαρίες της Γροιλανδίας ζουν περίπου $400$ χρόνια. Ένα έτος έχει περίπου $365$ ημέρες.\par
\tithelo Πόσα λεπτά είναι τα $400$ χρόνια; Να δώσω την απάντηση σε τυποποιημένη μορφή.\par
\possindeo Ένα έτος έχει $365\cdot24\cdot60=525\,600$ λεπτά. Τα $400$ έτη έχουν $400\cdot525\,600=210\,240\,000$ λεπτά. Ο αριθμός έχει $9$ ψηφία, άρα ο εκθέτης είναι $8$.\par
\apantisi $2{,}1024\cdot10^{8}$ λεπτά.\par
\elegxos $2{,}1024\cdot10^8$: η υποδιαστολή πάει $8$ θέσεις δεξιά και δίνει $210\,240\,000$.

\paradeigma{Ποιοι αριθμοί είναι τυποποιημένοι;}
\tixero Οι $4{,}5\cdot10^{3}$, $45\cdot10^{2}$, $0{,}45\cdot10^{4}$, $1\cdot10^{6}$.\par
\tithelo Να βρω ποιοι είναι σε τυποποιημένη μορφή και να διορθώσω τους άλλους.\par
\possindeo Ελέγχω αν το ακέραιο μέρος του πρώτου παράγοντα είναι ένα ψηφίο από το $1$ ως το $9$. Στο $4{,}5$ και στο $1$ ναι. Το $45$ έχει δύο ψηφία και το $0{,}45$ έχει ακέραιο μέρος $0$.\par
\apantisi Τυποποιημένοι είναι οι $4{,}5\cdot10^{3}$ και $1\cdot10^{6}$. Οι άλλοι δύο είναι ίσοι με τον πρώτο, $45\cdot10^2=4\,500=0{,}45\cdot10^4=4{,}5\cdot10^3$.\par
\elegxos Και οι τρεις δίνουν $4\,500$.

\tmima{Ασκήσεις}

\zoni{Το κατάλαβα}{από και προς την τυποποιημένη μορφή}

\ex Γράψε σε τυποποιημένη μορφή: α) $150\,000$ \quad β) $21\,000$ \quad γ) $1\,300\,000$ \quad δ) $2\,030\,000$.

\ex Γράψε σε πλήρη μορφή: α) $2{,}31\cdot10^{4}$ \quad β) $1{,}36\cdot10^{9}$ \quad γ) $1{,}37\cdot10^{2}$ \quad δ) $10^{6}$.

\ex Σύγκρινε τους αριθμούς με ένα από τα $<,\ >,\ =$: α) $2\cdot10^{3}\ \ldots\ 3\cdot10^{2}$ \quad β) $0{,}1\cdot10^{2}\ \ldots\ 10$ \quad γ) $1{,}7\cdot10^{4}\ \ldots\ 3\cdot10^{3}$ \quad δ) $6\cdot10^{5}\ \ldots\ 0{,}6\cdot10^{6}$.

\ex Ποιοι από τους αριθμούς $4{,}5\cdot10^{3}$, $45\cdot10^{2}$, $0{,}45\cdot10^{4}$, $9{,}99\cdot10^{8}$, $10\cdot10^{5}$, $1\cdot10^{6}$ είναι σε τυποποιημένη μορφή;

\ex Γράψε σε τυποποιημένη μορφή τους αριθμούς που δεν είναι, από την προηγούμενη άσκηση.

\ex Γράψε στη μορφή $\alpha\cdot10^{\nu}$ τους αριθμούς: α) $7\,000$ \quad β) $480\,000\,000$ \quad γ) $12\,500\,000\,000$ \quad δ) $6\,021\,000$.

\ex Σωστό ή λάθος; α) $3\cdot10^{5}=300\,000$. \quad β) $2{,}5\cdot10^{4}=250\,000$. \quad γ) Ο $10^{7}$ γράφεται με το $1$ και επτά μηδενικά. \quad δ) $5{,}2\cdot10^{3}>4{,}9\cdot10^{3}$. \quad ε) $3\cdot10^{2}=6$. \quad στ) $7\cdot10^{1}=70$.

\zoni{Το χρησιμοποιώ}{διάστημα, χρόνος, πληθυσμοί}

\ex \textbf{Τα αστρονομικά μεγέθη.} Γράψε σε τυποποιημένη μορφή: α) Απόσταση Γης -- Σελήνης: $384\,400$ χιλιόμετρα. \quad β) Απόσταση Γης -- Ήλιου: $149\,600\,000$ χιλιόμετρα. \quad γ) Ηλικία της Γης: $4\,500\,000\,000$ έτη. \quad δ) Ηλικία του Γαλαξία: $13\,600\,000\,000$ έτη. \quad ε) Απόσταση Πλούτωνα -- Ήλιου: $5\,900\,000\,000$ χιλιόμετρα.

\ex \textbf{Η χελώνα.} Μια γιγάντια χελώνα ζει περίπου $150$ χρόνια. Σε πόσα λεπτά αντιστοιχούν (με έτος $365$ ημερών); Γράψε την απάντηση σε τυποποιημένη μορφή.

\ex \textbf{Οι πλανήτες.} Οι αποστάσεις από τον Ήλιο: Δίας $7{,}8\cdot10^{8}$ km, Άρης $2{,}3\cdot10^{8}$ km, Ποσειδώνας $4{,}5\cdot10^{9}$ km, Ερμής $5{,}8\cdot10^{7}$ km. Γράψε τους πλανήτες από τον πιο κοντινό στον πιο μακρινό.

\ex \textbf{Το κομπιουτεράκι.} Το κομπιουτεράκι της Πολυξένης δείχνει $3{,}2\mathrm{E}8$ και μετά $4{,}05\mathrm{E}6$. Ποιοι είναι οι αριθμοί;

\ex \textbf{Οι πληθυσμοί.} Η Κίνα έχει περίπου $1\,400\,000\,000$ κατοίκους και η Ελλάδα περίπου $10\,400\,000$. Γράψε τους αριθμούς σε τυποποιημένη μορφή και πες ποιος είναι μεγαλύτερος και πώς το βλέπεις από τους εκθέτες.

\zoni{Το σκέφτομαι}{γιατί; και τι θα συνέβαινε αν...;}

\ex Εξήγησε γιατί το $12\cdot10^{5}$ δεν είναι σε τυποποιημένη μορφή και γράψε το σωστά.

\ex Πόσα ψηφία έχει ένας αριθμός που γράφεται στη μορφή $\alpha\cdot10^{\nu}$; Πόσα ψηφία έχουν οι $4\cdot10^{9}$ και $4{,}5\cdot10^{9}$;

\ex Ο Στέλιος έγραψε $5\,400\,000=5{,}4\cdot10^{5}$. Βρες το λάθος του.

\ex Ανάμεσα σε ποιες δύο διαδοχικές δυνάμεις του $10$ βρίσκεται ο $4{,}5\cdot10^{6}$;

\zoni{Για το σπίτι}{ψάξε και γράψε}

\ex Ψάξε πέντε μεγάλα μεγέθη από την εγκυκλοπαίδεια ή το διαδίκτυο (π.χ. πληθυσμός της Γης, μάζα ή ακτίνα πλανήτη) και γράψε τα σε τυποποιημένη μορφή.

\ex Υπολόγισε πόσα δευτερόλεπτα έχει ένα έτος ($365$ ημέρες) και γράψε το αποτέλεσμα σε τυποποιημένη μορφή.

\epanalipsi{Α3}

\tmima{Λεξικό του κεφαλαίου}

\renewcommand{\arraystretch}{1.5}
\begin{longtable}{@{}>{\raggedright\arraybackslash}p{0.28\linewidth} >{\raggedright\arraybackslash}p{0.62\linewidth}@{}}
\toprule
\textbf{Λέξη ή τύπος} & \textbf{Με μία πρόταση}\\
\midrule
\endhead
ρητοί & Οι αριθμοί που γράφονται ως κλάσμα $\frac{\alpha}{\beta}$ με $\alpha,\beta$ ακέραιους και $\beta\neq0$.\\
ισοδύναμα κλάσματα & Ίδια αξία: $\frac{\alpha}{\beta}=\frac{\gamma}{\delta}$ όταν $\alpha\delta=\beta\gamma$.\\
ανάγωγο & Κλάσμα που δεν απλοποιείται: $\MKD$ των όρων ίσος με $1$.\\
$|\alpha|$ & Η απόσταση του $\alpha$ από το $0$: ποτέ αρνητική.\\
αντίθετοι & Ετερόσημοι με ίδια απόλυτη τιμή: $\alpha$ και $-\alpha$, με άθροισμα $0$.\\
σύγκριση & Μεγαλύτερος είναι ο δεξιότερος στον άξονα.\\
$\alpha-\beta$ & $=\alpha+(-\beta)$: προσθέτω τον αντίθετο.\\
γινόμενο, πηλίκο & Ομόσημοι: θετικό. Ετερόσημοι: αρνητικό. Ποτέ διαίρεση με το $0$.\\
αντίστροφος & Ο $\frac{\beta}{\alpha}$ για τον $\frac{\alpha}{\beta}$: γινόμενο $1$. Το $0$ δεν έχει.\\
$\alpha^{\nu}$ & $\nu$ παράγοντες ίσοι με $\alpha$. Ίδιο πρόσημο με το γινόμενο.\\
$\alpha\cdot10^{\nu}$ & Τυποποιημένη μορφή: $1\le\alpha<10$ και $\nu$ μη μηδενικός φυσικός.\\
\bottomrule
\end{longtable}

\tmima{Ασκήσεις επανάληψης}

\ex Για καθέναν από τους $-\frac{6}{3}$, $\frac{4}{5}$, $0$, $-2{,}5$, $\frac{7}{7}$ πες αν είναι φυσικός, ακέραιος, ρητός.

\ex Απλοποίησε το $\frac{36}{48}$ μέχρι να γίνει ανάγωγο και γράψε ένα κλάσμα ισοδύναμο με το $\frac{5}{6}$ που έχει παρονομαστή $42$.

\ex Γράψε το $\frac{7}{8}$ ως δεκαδικό και τον $0{,}36$ ως ανάγωγο κλάσμα.

\ex Βρες: α) $|-7{,}5|$ \quad β) $\left|\frac{3}{4}\right|$ \quad γ) τον αντίθετο του $-\frac{2}{3}$ \quad δ) όλους τους $\alpha$ με $|\alpha|=4$.

\ex Γράψε σε αύξουσα σειρά τους αριθμούς $-2{,}5$, $0$, $-\frac{1}{2}$, $\frac{3}{2}$, $-3$, $1{,}2$.

\ex Βρες δύο ρητούς ανάμεσα στους $\frac{2}{7}$ και $\frac{3}{7}$.

\ex Υπολόγισε: α) $-\frac{3}{4}+\frac{5}{6}$ \quad β) $2-\frac{7}{5}$ \quad γ) $-1{,}5+0{,}75$ \quad δ) $\frac{1}{2}-\left(-\frac{1}{3}\right)$.

\ex Υπολόγισε την παράσταση $-\frac{1}{2}+\frac{2}{3}-\frac{1}{6}+\frac{1}{4}$.

\ex Υπολόγισε: α) $\left(-\frac{3}{4}\right)\cdot\frac{8}{9}$ \quad β) $(-0{,}5)\cdot(-1{,}2)$ \quad γ) $12\cdot\left(\frac{1}{3}+\frac{1}{4}\right)$.

\ex Υπολόγισε: α) $\left(-\frac{3}{5}\right):\frac{9}{10}$ \quad β) $6:\left(-\frac{3}{4}\right)$ \quad γ) $7{,}2:0{,}09$.

\ex Μετάτρεψε σε απλό κλάσμα το $\dfrac{\frac{2}{3}+\frac{1}{6}}{\frac{5}{4}}$.

\ex Υπολόγισε: α) $\left(-\frac{2}{3}\right)^3$ \quad β) $-3^2$ \quad γ) $(-0{,}1)^2$ \quad δ) $\frac{1}{2}+\left(-\frac{1}{2}\right)^2$.

\ex Γράψε σε τυποποιημένη μορφή τον $45\,000\,000$, σε πλήρη μορφή τον $3{,}02\cdot10^{5}$ και σύγκρινε τους $8\cdot10^{6}$ και $9{,}5\cdot10^{5}$.

\ex \textbf{Ο αγώνας.} Ένα σχολείο έχει $240$ μαθητές. Τα $\frac{3}{8}$ παίζουν μπάσκετ και το $\frac{1}{3}$ αυτών είναι κορίτσια. Πόσα κορίτσια παίζουν μπάσκετ;

\ex \textbf{Σωστό ή λάθος (πράξεις).} α) $-\frac{3}{4}<-\frac{2}{3}$. \quad β) $\frac{1}{3}+\frac{1}{3}=\frac{2}{6}$. \quad γ) $(-2)\cdot(-3)\cdot(-4)=-24$. \quad δ) $5:\frac{1}{5}=1$. \quad ε) Το $0$ έχει αντίστροφο. \quad στ) Ο αντίθετος του $0$ είναι το $0$.

\ex \textbf{Σωστό ή λάθος (δυνάμεις, τυποποιημένη μορφή).} α) $(-2)^3=-6$. \quad β) $-2^2=-4$. \quad γ) $\left(\frac{2}{3}\right)^2=\frac{4}{3}$. \quad δ) $5\cdot10^{4}=50\,000$. \quad ε) Το $12\cdot10^{3}$ είναι σε τυποποιημένη μορφή. \quad στ) Τα $\frac{3}{4}$ και $\frac{9}{12}$ είναι ισοδύναμα.

\Needspace{24\baselineskip}\tmima{Ο χάρτης της επανάληψης}

Μετά τις λύσεις, δες σε ποιες ασκήσεις δυσκολεύτηκες. Ο πίνακας σου λέει ποια παράγραφο να ξαναδιαβάσεις.

\renewcommand{\arraystretch}{1.5}
\begin{tabular*}{\linewidth}{@{\extracolsep{\fill}} >{\raggedright\arraybackslash}p{0.38\linewidth} >{\raggedright\arraybackslash}p{0.52\linewidth}}
\toprule
\textbf{Αν δυσκολεύτηκες στις} & \textbf{ξαναδιάβασε την παράγραφο}\\
\midrule
ΕΑ3.1, ΕΑ3.2, ΕΑ3.3 & Α3.1 (ρητοί, ισοδύναμα, δεκαδικοί)\\
ΕΑ3.4, ΕΑ3.5, ΕΑ3.6 & Α3.2 (απόλυτη τιμή, διάταξη, ενδιάμεσοι ρητοί)\\
ΕΑ3.7, ΕΑ3.8 & Α3.3 (πρόσθεση και αφαίρεση)\\
ΕΑ3.9, ΕΑ3.14 & Α3.4 (πολλαπλασιασμός, κλάσμα ποσότητας)\\
ΕΑ3.10, ΕΑ3.11 & Α3.5 (διαίρεση, σύνθετα κλάσματα)\\
ΕΑ3.12 & Α3.6 (δυνάμεις)\\
ΕΑ3.13 & Α3.7 (τυποποιημένη μορφή)\\
ΕΑ3.15, ΕΑ3.16 & Α3.1 έως Α3.7\\
\bottomrule
\end{tabular*}


% Βοηθητικά σχήματα του κεφαλαίου Α4 (μοναδικά ονόματα \bn...)
\providecommand{\bnMatch}[1]{\begin{tikzpicture}[scale=0.5,baseline=(current bounding box.center)]
  \foreach \i in {1,...,#1}{\draw[brown!80!black,line width=1.6pt] (\i-1,0) rectangle (\i,1);}
\end{tikzpicture}}
\providecommand{\bnStair}[1]{\begin{tikzpicture}[scale=0.38,baseline=(current bounding box.center)]
  \foreach \r in {1,...,#1}{\foreach \c in {1,...,\r}{\draw[fill=gblue!35] (\c-1,-\r) rectangle (\c,-\r+1);}}
\end{tikzpicture}}
\providecommand{\bnL}[1]{\begin{tikzpicture}[scale=0.4,baseline=(current bounding box.center)]
  \foreach \c in {1,...,#1}{\draw[fill=gwarmD!35] (\c-1,0) rectangle (\c,1);}
  \foreach \r in {2,...,#1}{\draw[fill=gwarmD!35] (0,\r-1) rectangle (1,\r);}
\end{tikzpicture}}
\providecommand{\bnBW}[1]{\begin{tikzpicture}[scale=0.4,baseline=(current bounding box.center)]
  \foreach \c in {1,...,#1}{\draw[fill=gblue!55] (\c-1,0) rectangle (\c,1);}
  \pgfmathtruncatemacro{\bnnext}{#1+1}
  \draw (#1,0) rectangle (\bnnext,1);
\end{tikzpicture}}
\providecommand{\bnShapes}[2]{\begin{tikzpicture}[scale=0.45,baseline=(current bounding box.center)]
  \foreach \i in {0,...,#1}{\pgfmathtruncatemacro{\bnm}{mod(\i,#2)}
    \ifnum\bnm=0 \draw[fill=gblue!50] (\i*1.3,0) circle (0.42);\fi
    \ifnum\bnm=1 \draw[fill=gwarmD!45] (\i*1.3-0.42,-0.42) rectangle (\i*1.3+0.42,0.42);\fi
    \ifnum\bnm=2 \draw[fill=green!50!black!35] (\i*1.3-0.5,-0.4)--(\i*1.3+0.5,-0.4)--(\i*1.3,0.5)--cycle;\fi}
\end{tikzpicture}}
\kefalaio{Κανονικότητες}{Όποιος βρει τον κανόνα, προβλέπει το επόμενο}{%
\begin{tabular}{@{}l@{\hspace{1em}}p{0.82\linewidth}@{}}
4.1 & Κανονικότητα, έννοια και αναπαραστάσεις\\[2pt]
4.2 & Κανονικότητες της μορφής $\alpha\cdot\nu$ με $\alpha$ ρητό\\[2pt]
 & Επανάληψη του κεφαλαίου\\
\end{tabular}}
\paragrafos{Η έννοια της εξίσωσης: ο αριθμός που κρύβεται}{Α.4.1, «Η έννοια της εξίσωσης – Οι εξισώσεις $\alpha+x=\beta$, $x-\alpha=\beta$, $\alpha-x=\beta$, $\alpha x=\beta$, $\alpha:x=\beta$, $x:\alpha=\beta$»}

\begin{anagke}
Στο διάλειμμα η Πολυξένη, που αγαπάει τους κωδικούς, προτείνει ένα παιχνίδι: «Σκέψου έναν αριθμό, πολλαπλασίασέ τον επί τρία, πρόσθεσε πέντε και πες μου μόνο το αποτέλεσμα. Εγώ θα βρω τον αριθμό σου». Ο Νίκος σκέφτεται και λέει «είκοσι έξι». Η Πολυξένη απαντάει αμέσως: «Επτά!». Ο Νίκος ανοίγει το στόμα: ήταν πράγματι το επτά.

Η Μυρτώ δοκιμάζει με το $41$ και η Πολυξένη βρίσκει πάλι τον αριθμό. «Μαγεία;» ρωτάει ο Στέλιος. «Όχι», απαντάει η Πολυξένη. «Αναποδογυρίζω τα βήματα: αφαιρώ το πέντε και διαιρώ με το τρία». Ο Βασίλης αμφιβάλλει: «Μα εσύ δεν ξέρεις τον αριθμό. Πώς γράφεις κάτι που δεν ξέρεις;». Η Αγγελική παίρνει το μολύβι και ζωγραφίζει στο τετράδιο ένα κουτάκι: «Βάζω ένα κουτάκι εκεί που κρύβεται ο αριθμός. Τρία κουτάκια συν πέντε κάνουν είκοσι έξι».

Στο μάθημα της κυρίας Αλεξάνδρας η Φανή ρωτάει: «Αν αντί για κουτάκι γράψουμε ένα γράμμα, το ίδιο παιχνίδι λειτουργεί για κάθε αριθμό;». Εκεί αρχίζουν τα γράμματα στα μαθηματικά. Το γράμμα είναι ένα κουτάκι για τον αριθμό που δεν ξέρουμε ακόμα, και με τη βοήθειά του μπορούμε να γράψουμε το παιχνίδι ως ισότητα και να το λύσουμε.
\end{anagke}

\begin{lexi}{εξίσωση, άγνωστος, άλγεβρα}{εξίσωση, άγνωστος, ρίζα}
Η λέξη \lexiword{εξίσωση} βγαίνει από το «ίσος»: εξισώνω σημαίνει «κάνω ίσα» δύο ποσά, όπως οι δύο δίσκοι μιας ζυγαριάς. \lexiword{Άγνωστος} είναι ο αριθμός που δεν ξέρουμε και ψάχνουμε. Η \lexiword{ρίζα} της εξίσωσης είναι η λύση της, και λέγεται έτσι γιατί «βγαίνει από το βάθος», όπως η ρίζα ενός φυτού. Η λέξη \lexiword{άλγεβρα} έρχεται από το αραβικό «αλ-τζαμπρ» (αποκατάσταση), από τον τίτλο ενός βιβλίου του Αλ-Χουαρίζμι τον 9ο αιώνα μ.Χ., όπου λύνονταν εξισώσεις.
\end{lexi}

\begin{metafrasi}
Ο επόμενος και ο προηγούμενος του $\nu$. & $\nu+1$ και $\nu-1$\\
Ένας άρτιος και ένας περιττός. & $2\nu$ και $2\nu+1$\\
Τα πολλαπλάσια του $3$. & $3\nu$\\
Το διπλάσιο του $x$. & $2x$\\
Ένας αριθμός αυξημένος κατά $8$, ελαττωμένος κατά $4$. & $x+8$ και $x-4$\\
Το τετραπλάσιο ενός αριθμού, αυξημένο κατά $2$, κάνει $22$. & $4x+2=22$\\
Αν σε έναν αριθμό προσθέσω $5$ βγαίνει $8$. & $x+5=8$\\
\end{metafrasi}

\tmima{Η θεωρία}

\tmimasmall{Γράμματα στη θέση αριθμών}

Με ένα γράμμα (π.χ. $x$, $\alpha$, $\nu$) παριστάνουμε έναν αριθμό, γνωστό ή άγνωστο, και με γράμματα και πράξεις γράφουμε προτάσεις συνοπτικά. Αν η πλευρά ενός τετραγώνου είναι $\alpha$, η περίμετρός του είναι $4\alpha$ και το εμβαδόν του $\alpha\cdot\alpha$. Μια έκφραση με γράμματα και αριθμούς λέγεται \textbf{αλγεβρική παράσταση}.

Όταν ο ίδιος άγνωστος εμφανίζεται πολλές φορές, συντομεύουμε με την επιμεριστική ιδιότητα: $x+x+x+x=4x$, $3\alpha+5\alpha=(3+5)\alpha=8\alpha$, $15\beta-9\beta=6\beta$, $18x+7x+4x=29x$. Μπορούμε να προσθέτουμε ή να αφαιρούμε μόνο τους όρους που έχουν το \emph{ίδιο γράμμα}: $4\alpha+2\beta+\alpha+3\beta=5\alpha+5\beta$. Δεν συνενώνουμε το $5\alpha$ με το $5\beta$.

\tmimasmall{Τι είναι εξίσωση}

Η ισότητα $2\cdot6=12$ μπορούμε μόνο να τη βεβαιώσουμε. Η ισότητα $2\cdot x=12$ δεν είναι ίδια: κρύβει έναν αριθμό που, αν τον βάλουμε στη θέση του $x$, την κάνει αληθινή, ενώ κάθε άλλος αριθμός την κάνει ψευδή.

\begin{orismos}{εξίσωση με έναν άγνωστο}
\textbf{Εξίσωση} με έναν άγνωστο λέγεται μια ισότητα που περιέχει αριθμούς και ένα γράμμα (τον άγνωστο), π.χ. $x+5=12$, $10-z=1$, $7\varphi=14$, $24:\psi=6$.\\
\textbf{Λύση ή ρίζα} της εξίσωσης είναι ο αριθμός που, όταν αντικαταστήσει τον άγνωστο, \textbf{επαληθεύει} την ισότητα. Η διαδικασία εύρεσης της λύσης λέγεται \textbf{επίλυση}.
\end{orismos}

Για παράδειγμα η ρίζα της $x-7=5$ είναι το $12$, γιατί $12-7=5$. Γράφουμε $x=12$.

\tmimasmall{Ταυτότητα και αδύνατη εξίσωση}

\begin{itemize}[leftmargin=1.2em,itemsep=0.3ex,topsep=2pt]
\item Μια εξίσωση λέγεται \textbf{ταυτότητα} (ή αόριστη) όταν \emph{κάθε} αριθμός είναι λύση της, π.χ. $x=x$ ή $0\cdot z=0$.
\item Λέγεται \textbf{αδύνατη} όταν \emph{κανένας} αριθμός δεν την επαληθεύει, π.χ. $x+2=x+6$ ή $0\cdot\omega=5$.
\end{itemize}

\tmimasmall{Οι έξι βασικές εξισώσεις}

Από τον ορισμό των πράξεων (η αφαίρεση είναι η αντίστροφη πράξη της πρόσθεσης και η διαίρεση του πολλαπλασιασμού) βρίσκουμε τη λύση κάθε μίας:

\begin{center}
\renewcommand{\arraystretch}{1.5}
\begin{tabular}{@{}lll@{}}
\toprule
\textbf{Εξίσωση} & \textbf{Λύση} & \textbf{Παράδειγμα}\\
\midrule
$x+\alpha=\beta$ & $x=\beta-\alpha$ & $x+5=12\ \Rightarrow\ x=12-5=7$\\
$x-\alpha=\beta$ & $x=\beta+\alpha$ & $y-2=3\ \Rightarrow\ y=3+2=5$\\
$\alpha-x=\beta$ & $x=\alpha-\beta$ & $10-z=1\ \Rightarrow\ z=10-1=9$\\
$\alpha\cdot x=\beta$ & $x=\beta:\alpha$ & $7\varphi=14\ \Rightarrow\ \varphi=14:7=2$\\
$x:\alpha=\beta$ & $x=\beta\cdot\alpha$ & $\omega:5=4\ \Rightarrow\ \omega=4\cdot5=20$\\
$\alpha:x=\beta$ & $x=\alpha:\beta$ & $24:\psi=6\ \Rightarrow\ \psi=24:6=4$\\
\bottomrule
\end{tabular}
\end{center}

(Στις δύο τελευταίες και στην τέταρτη προϋποθέτουμε ότι δεν διαιρούμε με το $0$.) Μετά τη λύση \textbf{επαληθεύουμε}: βάζουμε την τιμή στην αρχική εξίσωση και ελέγχουμε ότι βγαίνει ισότητα. Στις εξισώσεις δεν γράφουμε τις μονάδες, αλλά πρέπει να ξέρουμε σε ποιες μονάδες δουλεύουμε.

\begin{prosoxi}
\begin{itemize}[leftmargin=1.2em,itemsep=0.3ex,topsep=2pt]
\item Στην $\alpha-x=\beta$ η λύση είναι $\alpha-\beta$, όχι $\beta-\alpha$. Η $12-x=5$ δίνει $x=7$, όχι $17$.
\item Στις $\alpha:x=\beta$ και $x:\alpha=\beta$ οι λύσεις διαφέρουν: $24:\psi=6$ δίνει $\psi=24:6=4$, ενώ $\omega:6=24$ δίνει $\omega=24\cdot6=144$.
\item Το $3\alpha+5\beta$ δεν γίνεται $8\alpha\beta$ ούτε $8\alpha$: έχουμε διαφορετικά γράμματα και δεν συντομεύεται.
\item Δεν αντικαθιστούμε τον άγνωστο και μετά «επαληθεύουμε» με την ίδια πράξη που κάναμε: επαληθεύω βάζοντας τη λύση στην \emph{αρχική} εξίσωση.
\end{itemize}
\end{prosoxi}

\tmima{Λυμένα παραδείγματα}

\paradeigma{Μεταφράζω σε γλώσσα γραμμάτων}
\tixero Ο Νίκος έχει $x$ ευρώ. Η Αγγελική έχει $7$ ευρώ περισσότερα και μαζί έχουν $25$ ευρώ.\par
\tithelo Να γράψω με γράμματα τα χρήματα της Αγγελικής και την ισότητα για το σύνολο.\par
\possindeo Η Αγγελική έχει $x+7$. Μαζί έχουν $x+(x+7)=2x+7$.\par
\apantisi Αγγελική: $x+7$. Σύνολο: $2x+7=25$.\par
\elegxos Αν ο Νίκος είχε $9$ €, η Αγγελική θα είχε $16$ € και μαζί $25$. Άρα $x=9$ ταιριάζει (με την $2x+7=25$ θα το δούμε στην επόμενη παράγραφο).

\paradeigma{Συντομεύω παραστάσεις}
\tixero Οι παραστάσεις $5x+3x-2x$ και $4\alpha+2\beta+\alpha+3\beta$.\par
\tithelo Να τις γράψω απλούστερα.\par
\possindeo Προσθέτω ή αφαιρώ τους όρους που έχουν το ίδιο γράμμα. $5x+3x-2x=(5+3-2)x=6x$. Και $4\alpha+\alpha=5\alpha$, $2\beta+3\beta=5\beta$.\par
\apantisi $6x$ και $5\alpha+5\beta$.\par
\elegxos Με $x=2$: $10+6-4=12=6\cdot2$. Με $\alpha=1$, $\beta=2$: $4+4+1+6=15=5+10$. Σωστά.

\paradeigma{Ποιος αριθμός επαληθεύει;}
\tixero Η εξίσωση $x-4=1$ και οι αριθμοί $3$, $5$, $8$.\par
\tithelo Ποιος είναι ρίζα.\par
\possindeo Αντικαθιστώ: για $x=3$: $3-4$, δεν γίνεται στους φυσικούς. Για $x=5$: $5-4=1$ σωστό. Για $x=8$: $8-4=4\neq1$.\par
\apantisi Ρίζα είναι ο $5$.\par
\elegxos Από τον κανόνα: $x=1+4=5$. Συμφωνεί.

\paradeigma{Λύνω τις έξι βασικές}
\tixero Οι εξισώσεις $x+7{,}5=12$, $x-3{,}2=5$, $15-x=6{,}4$, $4x=18$, $x:6=2{,}5$, $36:x=9$.\par
\tithelo Να τις λύσω και να επαληθεύσω.\par
\possindeo $x=12-7{,}5=4{,}5$. $x=5+3{,}2=8{,}2$. $x=15-6{,}4=8{,}6$. $x=18:4=4{,}5$. $x=2{,}5\cdot6=15$. $x=36:9=4$.\par
\apantisi $4{,}5$, $8{,}2$, $8{,}6$, $4{,}5$, $15$, $4$.\par
\elegxos $4{,}5+7{,}5=12$, $8{,}2-3{,}2=5$, $15-8{,}6=6{,}4$, $4\cdot4{,}5=18$, $15:6=2{,}5$, $36:4=9$. Όλα σωστά.

\paradeigma{Ταυτότητα ή αδύνατη;}
\tixero Οι εξισώσεις $x+3=3+x$, $x+2=x+6$ και $2x-2x=5$.\par
\tithelo Τι είδους είναι.\par
\possindeo Η πρώτη ισχύει για κάθε αριθμό (αντιμεταθετική ιδιότητα): ταυτότητα. Στη δεύτερη, όποιον αριθμό και να βάλω, το αριστερό μέλος είναι κατά $4$ μικρότερο από το δεξί: αδύνατη. Στην τρίτη το $2x-2x=0$ και δεν ισούται με $5$: αδύνατη.\par
\apantisi Ταυτότητα, αδύνατη, αδύνατη.\par
\elegxos Με $x=10$: $13=13$, αλλά $12\neq16$ και $0\neq5$.

\paradeigma{Το ύψος της δεξαμενής}
\tixero Η νέα δεξαμενή του χωριού έχει όγκο $7{,}2$ m$^3$, μήκος $2$ μ. και πλάτος $1{,}2$ μ.\par
\tithelo Το ύψος της.\par
\possindeo Αν $x$ το ύψος σε μέτρα, ο όγκος είναι $2\cdot1{,}2\cdot x=2{,}4x$. Άρα $2{,}4x=7{,}2$ και $x=7{,}2:2{,}4=3$.\par
\apantisi $3$ μ.\par
\elegxos $2\cdot1{,}2\cdot3=7{,}2$ m$^3$. Σωστό.

\tmima{Ασκήσεις}

\zoni{Το κατάλαβα}{γράμματα, παραστάσεις, έξι βασικές εξισώσεις}

\ex Γράψε με γράμματα (με $x$ έναν αριθμό): α) ο επόμενός του \quad β) ο προηγούμενός του \quad γ) το τριπλάσιό του \quad δ) αυξημένος κατά $12$ \quad ε) ελαττωμένος κατά $5$ \quad στ) το μισό του \quad ζ) το άθροισμά του με το διπλάσιό του.

\ex Γράψε με γράμματα: α) ένας άρτιος φυσικός \quad β) ένας περιττός φυσικός \quad γ) ένα πολλαπλάσιο του $5$ \quad δ) τρεις διαδοχικοί φυσικοί.

\ex Διατύπωσε με λόγια: α) $5x-3$ \quad β) $2(x+1)=14$ \quad γ) $\alpha+\beta$ \quad δ) $6\kappa-\kappa=30$.

\ex Η πλευρά ενός τετραγώνου είναι $\alpha$. Γράψε την περίμετρο και το εμβαδόν του. Κάνε το ίδιο για ένα ορθογώνιο με πλευρές $\alpha$ και $\beta$.

\ex Γράψε απλούστερα: α) $x+x+x$ \quad β) $3\alpha+5\alpha$ \quad γ) $7\beta-2\beta$ \quad δ) $2x+5y+4x-y$ \quad ε) $9\omega-4\omega+\omega$ \quad στ) $10\kappa+3\lambda-10\kappa$.

\ex Εξέτασε αν ο $12$ είναι λύση της $x-5=7$. Ποιος από τους $2,\,4,\,6$ είναι λύση της $3x=12$;

\ex Βρες τη λύση: α) $x+6=10$ \quad β) $x-2=7$ \quad γ) $15-x=9$ \quad δ) $4x=20$ \quad ε) $x:3=5$ \quad στ) $18:x=6$.

\ex Λύσε: α) $x+2{,}75=6$ \quad β) $x-0{,}8=1{,}15$ \quad γ) $12-x=4{,}35$ \quad δ) $0{,}5x=3{,}2$ \quad ε) $x:0{,}4=7$ \quad στ) $9:x=0{,}6$.

\ex Λύσε: α) $x+\frac{1}{4}=\frac{3}{4}$ \quad β) $\frac{2}{3}x=\frac{4}{9}$ \quad γ) $\frac{5}{6}-x=\frac{1}{3}$ \quad δ) $x:3=\frac{2}{9}$.

\ex Σωστό ή λάθος; α) Η $3+x=10$ έχει λύση το $7$. \quad β) Η $x-x=0$ είναι ταυτότητα. \quad γ) Η $x+1=x+2$ είναι αδύνατη. \quad δ) Η $0\cdot x=5$ είναι αδύνατη. \quad ε) Η $5x=0$ έχει λύση το $0$. \quad στ) Η $6:x=2$ έχει λύση το $12$.

\zoni{Το χρησιμοποιώ}{ζυγαριές, δεξαμενές, μυστικοί αριθμοί}

\ex \textbf{Η ζυγαριά.} Μια ζυγαριά ισορροπεί όταν στον έναν δίσκο υπάρχει μια σοκολάτα βάρους $x$ g και στον άλλον $150$ g και μισή σοκολάτα. Γράψε την ισότητα και βρες το βάρος της σοκολάτας με σκέψη.

\ex \textbf{Η δεξαμενή.} Μια δεξαμενή έχει όγκο $12$ m$^3$, μήκος $2$ μ. και πλάτος $1{,}5$ μ. Βρες το ύψος της με εξίσωση.

\ex \textbf{Ο μυστικός αριθμός.} Η Πολυξένη λέει: «Σκέφτηκα έναν αριθμό και βρήκα...». Γράψε εξίσωση και βρες τον αριθμό όταν: α) του πρόσθεσε $17$ και βρήκε $52$ \quad β) αφαίρεσε $9$ και βρήκε $14$ \quad γ) τον αφαίρεσε από το $50$ και βρήκε $18$ \quad δ) τον πολλαπλασίασε επί $6$ και βρήκε $84$ \quad ε) τον διαίρεσε με $8$ και βρήκε $7$.

\ex \textbf{Τα εισιτήρια.} Τέσσερα ίδια εισιτήρια πλοίου κοστίζουν $54$ €. Πόσο κάνει το καθένα; Γράψε την εξίσωση.

\ex \textbf{Ο κουμπαράς.} Ο Ηλίας έχει $36$ € και ξοδεύει κάποιο ποσό. Του μένουν $11{,}5$ €. Πόσα ξόδεψε;

\ex \textbf{Το μοτίβο.} Βάζουμε σε σειρά τετράγωνα πλευράς $1$ cm: το πρώτο σχήμα έχει $1$ τετράγωνο (περίμετρος $4$ cm), το δεύτερο $2$ (περίμετρος $6$ cm), το τρίτο $3$ (περίμετρος $8$ cm). α) Πόση είναι η περίμετρος του πέμπτου σχήματος; β) Γράψε έναν τύπο για το $\nu$-οστό σχήμα. γ) Ποιο σχήμα έχει περίμετρο $98$ cm;

\zoni{Το σκέφτομαι}{όταν το γράμμα κρύβει κάτι}

\ex \textbf{Το κόλπο της Πολυξένης.} Σκέψου έναν αριθμό $x$, πρόσθεσε $4$, διπλασίασε και αφαίρεσε $8$. Γράψε με παράσταση τι κάνεις και δείξε ότι το αποτέλεσμα είναι πάντα το διπλάσιο του αριθμού που σκέφτηκες.

\ex Η εξίσωση $x+8=5$ δεν έχει λύση στους φυσικούς αριθμούς. Γιατί; Τι είδους αριθμός θα την έλυνε;

\ex Ο Στέλιος λύνει την $12-x=5$ και γράφει $x=12+5=17$. Βρες το λάθος, το σωστό και επαλήθευσε.

\ex Γιατί η $0\cdot x=0$ έχει άπειρες λύσεις και η $0\cdot x=5$ καμία; Εξήγησε με το «ποιος αριθμός επί $0$ δίνει...».

\ex Γράψε μια εξίσωση με λύση το $6$, μια αδύνατη και μια ταυτότητα.

\zoni{Για το σπίτι}{φτιάξε το δικό σου}

\ex Ο παππούς λέει: «Σκέφτηκα έναν αριθμό και το τριπλάσιό του είναι $2{,}4$». Ποιος είναι; Γράψε την εξίσωση.

\ex Φτιάξε μια καθημερινή ιστορία που να οδηγεί στην εξίσωση $x+12=30$ και ζήτησε από ένα μέλος της οικογένειας να τη λύσει.

\paragrafos{Ο γενικός όρος α·ν: προβλέπω οποιονδήποτε όρο}{Α.4.2, «Κανονικότητες της μορφής α·ν με α ρητό»}

\begin{anagke}
Η Πολυξένη αγοράζει από το περίπτερο πακέτα με κάρτες για το άλμπουμ της. Κάθε πακέτο έχει $5$ κάρτες. Με $1$ πακέτο έχει $5$ κάρτες, με $2$ πακέτα $10$, με $3$ πακέτα $15$. «Πόσες κάρτες θα έχω αν αγοράσω $100$ πακέτα;» ρωτάει η Αγγελική. Ο Στέλιος αρχίζει να γράφει: $5,\ 10,\ 15,\ 20,\ 25,\ \ldots$ και σταματάει στο $60$ να παραπονιέται: «Δεν θα γράφω εκατό αριθμούς!».

Ο Ηλίας φτιάχνει πίνακα με τη σειρά και τον αριθμό των καρτών και βλέπει κάτι: $5=5\cdot1$, $10=5\cdot2$, $15=5\cdot3$, $20=5\cdot4$. «Κάθε αριθμός είναι το $5$ επί τη θέση του! Ο εκατοστός είναι ο $5\cdot100$». Η Ρήνη ρωτάει το αντίστροφο: «Αν έχω $225$ κάρτες, πόσα πακέτα αγόρασα;». Η Αγγελική προσθέτει: «Και αν ένα ρομποτάκι προχωρά $\frac{3}{4}$ του μέτρου κάθε δευτερόλεπτο; Ο κανόνας λειτουργεί και τότε;».

Η κυρία Αλεξάνδρα χαμογελά: «Αν βρείτε έναν τύπο που δίνει τον όρο για οποιαδήποτε θέση, μπορείτε να τον υπολογίσετε χωρίς να γράψετε τους προηγούμενους». Σε αυτή την παράγραφο μαθαίνουμε τι είναι ο γενικός όρος μιας κανονικότητας, τον βρίσκουμε όταν είναι της μορφής $\alpha\cdot\nu$ (και όταν ο $\alpha$ είναι ρητός, όχι μόνο φυσικός) και τον χρησιμοποιούμε για να βρίσκουμε όρους και θέσεις.
\end{anagke}

\begin{lexi}{ακολουθία, ν-οστός όρος, γενικός όρος}{ακολουθία, ν-οστός όρος, γενικός όρος}
Η \lexiword{ακολουθία} είναι αυτό που «ακολουθεί»: ένα στοιχείο μετά το άλλο, με τη σειρά τους. Ο \lexiword{ν-οστός όρος} είναι ο όρος που βρίσκεται στη θέση $\nu$, όπως λέμε «πέμπτος», «εκατοστός». Ο \lexiword{γενικός όρος} είναι «γενικός» επειδή ισχύει για κάθε θέση: ένας τύπος που περιγράφει όλους τους όρους μαζί.
\end{lexi}

\begin{metafrasi}
Ο όρος που βρίσκεται στη θέση $\nu$. & $\nu$-οστός όρος\\
Πέντε επί τη θέση του όρου. & $5\cdot\nu$ ή $5\nu$\\
Ο $10$ος όρος της κανονικότητας $5\nu$. & $5\cdot10=50$\\
Ποιος όρος ισούται με $225$; & $5\nu=225$, άρα $\nu=225:5=45$\\
Κάθε όρος είναι $\alpha$ φορές η θέση του. & $\alpha\cdot\nu$\\
\end{metafrasi}

\tmima{Η θεωρία}

\tmimasmall{Ακολουθίες και κανονικότητες}

Οι αριθμοί $5,\ 10,\ 15,\ 20,\ 25,\dots$ αποτελούν μια \textbf{ακολουθία}. Ακολουθία είναι μια διαδοχή στοιχείων, που ονομάζονται \textbf{όροι}, και έχουν μια ορισμένη σειρά: $1$ος όρος, $2$ος όρος, $3$ος όρος κ.ο.κ.

\begin{center}
\renewcommand{\arraystretch}{1.3}
\begin{tabular}{rcccccc}
$5,$ & $10,$ & $15,$ & $20,$ & $25,$ & $\dots$\\
$\uparrow$ & $\uparrow$ & $\uparrow$ & $\uparrow$ & $\uparrow$ & \\
{\small $1$ος όρος} & {\small $2$ος όρος} & {\small $3$ος όρος} & {\small $4$ος όρος} & {\small $5$ος όρος} &
\end{tabular}
\end{center}

Αν κάθε όρος μιας ακολουθίας προκύπτει με έναν συγκεκριμένο κανόνα, τότε οι όροι της αποτελούν μια \textbf{κανονικότητα}.

\begin{orismos}{ν-οστός όρος ή γενικός όρος}
Ο όρος που βρίσκεται στη θέση $\nu$ ($\nu=1,\,2,\,3,\dots$) λέγεται \textbf{ν-οστός όρος} ή \textbf{γενικός όρος} της ακολουθίας.
\end{orismos}

\tmimasmall{Η κανονικότητα $\alpha\cdot\nu$}

Στην κανονικότητα $5,\ 10,\ 15,\ 20,\ 25,\dots$ κάθε όρος είναι κατά $5$ μεγαλύτερος από τον προηγούμενο (κανόνας: $+5$). Παρατηρούμε ότι:
\begin{itemize}[leftmargin=1.2em,itemsep=0.3ex,topsep=2pt]
\item $1$ος όρος ($\nu=1$): $5=5\cdot1$
\item $2$ος όρος ($\nu=2$): $10=5\cdot2$
\item $3$ος όρος ($\nu=3$): $15=5\cdot3$
\item $4$ος όρος ($\nu=4$): $20=5\cdot4$
\item $\nu$-οστός όρος: $5\cdot\nu$
\end{itemize}

Ο γενικός όρος είναι λοιπόν $5\cdot\nu$ ή, πιο σύντομα, $5\nu$. Αν γνωρίζουμε τον γενικό όρο, μπορούμε να βρούμε όποιον όρο θέλουμε:
\begin{itemize}[leftmargin=1.2em,itemsep=0.3ex,topsep=2pt]
\item Ο $10$ος όρος είναι $5\cdot10=50$ και ο $23$ος είναι $5\cdot23=115$.
\item Για να βρούμε ποιος όρος ισούται με $225$, γράφουμε $5\nu=225$, άρα $\nu=225:5=45$. Ο $45$ος όρος ισούται με $225$.
\end{itemize}

\begin{center}
\begin{tikzpicture}[x=1.1cm,y=0.13cm]
  \draw[-{Stealth},thick] (0,0)--(5.6,0) node[right,font=\small]{$\nu$};
  \draw[-{Stealth},thick] (0,0)--(0,28) node[above,font=\small]{όρος};
  \foreach \x in {1,...,5}{\draw (\x,0.6)--(\x,-0.6) node[below,font=\scriptsize]{$\x$};}
  \foreach \y in {5,10,15,20,25}{\draw (0.07,\y)--(-0.07,\y) node[left,font=\scriptsize]{$\y$};}
  \draw[dashed,gray] (0,0)--(5.4,27);
  \foreach \x in {1,...,5}{\pgfmathtruncatemacro{\yy}{5*\x}\fill[gblue] (\x,\yy) circle (2.4pt);}
\end{tikzpicture}
\end{center}

Τα σημεία $(\nu,\ \text{όρος})$ της κανονικότητας $5\nu$ βρίσκονται σε μία ευθεία γραμμή, που περνά από την αρχή των αξόνων.

\begin{kanonas}{κανονικότητες της μορφής $\alpha\cdot\nu$}
Αν ο $1$ος όρος μιας κανονικότητας είναι ο $\alpha$ και κάθε όρος είναι κατά $\alpha$ μεγαλύτερος από τον προηγούμενο, τότε ο γενικός όρος είναι $\alpha\cdot\nu$. Ο $\alpha$ είναι ταυτόχρονα ο πρώτος όρος και το «βήμα». Αν γνωρίζω έναν όρο και τη θέση του, βρίσκω τον $\alpha$ διαιρώντας: $\alpha=\text{όρος}:\nu$.
\end{kanonas}

Ο $\alpha$ μπορεί να είναι οποιοσδήποτε ρητός: φυσικός ($7\nu$), αρνητικός ($-3\nu$), δεκαδικός ($0{,}5\nu$) ή κλάσμα ($\frac{3}{4}\nu$). Η σειρά $\nu$ είναι πάντα φυσικός αριθμός ($1,\,2,\,3,\dots$), γι' αυτό ένας αριθμός είναι όρος της κανονικότητας $\alpha\nu$ μόνο όταν το πηλίκο «αριθμός $:\ \alpha$» είναι φυσικός.

\tmimasmall{Γεωμετρικά μοτίβα}

Ένα ισόπλευρο τρίγωνο με πλευρά από $\nu$ οδοντογλυφίδες έχει συνολικά $3\nu$ οδοντογλυφίδες. Μια πυραμίδα με βάση πολύγωνο $\nu$ πλευρών έχει $\nu$ ακμές στη βάση και άλλες $\nu$ που ενώνουν την κορυφή με τη βάση: $2\nu$ ακμές. Σε κάθε περίπτωση ο γενικός όρος μας λέει πόσα στοιχεία χρειάζομαι για οποιοδήποτε σχήμα της σειράς.

\begin{prosoxi}
\begin{itemize}[leftmargin=1.2em,itemsep=0.3ex,topsep=2pt]
\item Το $7\nu$ δεν είναι το $7+\nu$: το πρώτο δίνει $7,\,14,\,21,\dots$ και το δεύτερο $8,\,9,\,10,\dots$
\item Το $\nu$ είναι η \emph{θέση} και όχι ο όρος. Ο $225$ ως όρος της $5\nu$ είναι ο $45$ος όρος, όχι ο $225$ος.
\item Για να βρω τη θέση διαιρώ: $\nu=\text{όρος}:\alpha$. Δεν πολλαπλασιάζω.
\item Αν το πηλίκο «όρος $:\ \alpha$» δεν βγαίνει φυσικός, ο αριθμός δεν είναι όρος της κανονικότητας: ο $100$ δεν είναι όρος της $7\nu$, γιατί $100:7$ δεν είναι φυσικός.
\item Δεν είναι κάθε προσθετική κανονικότητα της μορφής $\alpha\nu$: η $4,\,7,\,10,\,13,\dots$ έχει βήμα $3$, αλλά ο πρώτος όρος είναι $4$, όχι $3$. Ο γενικός της όρος είναι άλλος ($3\nu+1$).
\end{itemize}
\end{prosoxi}

\tmima{Λυμένα παραδείγματα}

\paradeigma{Βρίσκω τον γενικό όρο}
\tixero Η κανονικότητα $6,\ 12,\ 18,\ 24,\dots$\par
\tithelo Ο γενικός όρος, ο $7$ος και ο $20$ός όρος.\par
\possindeo Κάθε όρος είναι κατά $6$ μεγαλύτερος από τον προηγούμενο και ο πρώτος είναι $6$: $6=6\cdot1$, $12=6\cdot2$, $18=6\cdot3$.\par
\apantisi Γενικός όρος $6\nu$. Ο $7$ος είναι $6\cdot7=42$ και ο $20$ός $6\cdot20=120$.\par
\elegxos Ο $4$ος όρος από τον τύπο είναι $6\cdot4=24$, όπως στη σειρά.

\paradeigma{Από τη θέση στον όρο και αντίστροφα}
\tixero Ο γενικός όρος είναι $7\nu$.\par
\tithelo Ο $20$ός όρος. Ποιος όρος ισούται με $511$; Υπάρχει όρος ίσος με $100$;\par
\possindeo Ο $20$ός όρος: $\nu=20$. Για το $511$ λύνω $7\nu=511$, άρα $\nu=511:7=73$. Για το $100$: $100:7=14{,}28\ldots$ που δεν είναι φυσικός.\par
\apantisi $7\cdot20=140$. Ο $73$ος όρος ισούται με $511$. Όχι, δεν υπάρχει όρος ίσος με $100$.\par
\elegxos $7\cdot73=511$. Και τα πολλαπλάσια του $7$ δίπλα στο $100$ είναι τα $98$ και $105$.

\paradeigma{Αρνητικός συντελεστής}
\tixero Η κανονικότητα $-4,\ -8,\ -12,\ -16,\dots$\par
\tithelo Ο γενικός όρος, ο $25$ος όρος και η θέση του όρου $-96$.\par
\possindeo Κάθε όρος είναι κατά $4$ μικρότερος (βήμα $-4$) και ο πρώτος είναι $-4$: $-4=(-4)\cdot1$, $-8=(-4)\cdot2$. Για τη θέση: $(-4)\nu=-96$, άρα $\nu=(-96):(-4)=24$.\par
\apantisi Γενικός όρος $-4\nu$. Ο $25$ος είναι $-100$. Ο $-96$ είναι ο $24$ος όρος.\par
\elegxos $(-4)\cdot24=-96$.

\paradeigma{Δεκαδικός συντελεστής}
\tixero Η κανονικότητα $0{,}5,\ 1,\ 1{,}5,\ 2,\ \ldots$\par
\tithelo Ο γενικός όρος, ο $30$ός όρος και η θέση του όρου $12{,}5$.\par
\possindeo Ο πρώτος όρος είναι $0{,}5$ και κάθε όρος αυξάνεται κατά $0{,}5$: $0{,}5\nu$. Ο $30$ός όρος: $0{,}5\cdot30=15$. Για το $12{,}5$: $0{,}5\nu=12{,}5$, άρα $\nu=12{,}5:0{,}5=25$.\par
\apantisi Γενικός όρος $0{,}5\nu$, ο $30$ός είναι $15$ και ο $12{,}5$ είναι ο $25$ος όρος.\par
\elegxos Το $0{,}5\nu$ σημαίνει «το μισό του $\nu$»: το μισό του $30$ είναι $15$ και το μισό του $25$ είναι $12{,}5$.

\paradeigma{Κλασματικός συντελεστής}
\tixero Η κανονικότητα $\frac{3}{4},\ \frac{3}{2},\ \frac{9}{4},\ 3,\ \ldots$\par
\tithelo Ο γενικός όρος, ο $8$ος όρος και η θέση του όρου $15$.\par
\possindeo Ο πρώτος όρος είναι $\frac{3}{4}$ και κάθε όρος αυξάνεται κατά $\frac{3}{4}$: $\frac{3}{4}\nu$. Ο $8$ος: $\frac{3}{4}\cdot8=6$. Για το $15$: $\frac{3}{4}\nu=15$, άρα $\nu=15:\frac{3}{4}=15\cdot\frac{4}{3}=20$.\par
\apantisi Γενικός όρος $\frac{3}{4}\nu$, ο $8$ος είναι $6$ και ο $15$ είναι ο $20$ός όρος.\par
\elegxos $\frac{3}{4}\cdot20=15$.

\paradeigma{Πυραμίδες}
\tixero Μια πυραμίδα με βάση τρίγωνο έχει $6$ ακμές, με τετράπλευρη βάση $8$, με πενταγωνική $10$.\par
\tithelo Ο γενικός όρος και η πυραμίδα που έχει $50$ ακμές.\par
\possindeo Οι ακμές είναι $6,\,8,\,10,\dots$ για βάσεις με $3,\,4,\,5,\dots$ πλευρές, δηλαδή $6=2\cdot3$, $8=2\cdot4$, $10=2\cdot5$. Άρα $2\nu$ ακμές για βάση με $\nu$ πλευρές. Για τις $50$ ακμές: $2\nu=50$, $\nu=25$.\par
\apantisi Γενικός όρος $2\nu$ και η βάση έχει $25$ πλευρές.\par
\elegxos Η βάση με $25$ πλευρές χρειάζεται $25$ ακμές στη βάση και $25$ ακμές προς την κορυφή: $50$.

\paradeigma{Πίνακας με άγνωστους όρους}
\tixero Οι όροι της κανονικότητας με γενικό όρο $\alpha\nu$ είναι: $1{,}2,\ 2{,}4,\ x,\ 4{,}8,\ y,\ 7{,}2$.\par
\tithelo Οι $x$ και $y$ και ο $100$ός όρος.\par
\possindeo Ο πρώτος όρος είναι $1{,}2$, άρα $\alpha=1{,}2$ και ο γενικός όρος είναι $1{,}2\nu$. Ο $x$ είναι ο $3$ος και ο $y$ ο $5$ος όρος. Ο $6$ος είναι $1{,}2\cdot6=7{,}2$, όπως δίνεται.\par
\apantisi $x=1{,}2\cdot3=3{,}6$, $y=1{,}2\cdot5=6$ και ο $100$ός όρος $1{,}2\cdot100=120$.\par
\elegxos Η διαφορά διαδοχικών όρων είναι $1{,}2$: $2{,}4+1{,}2=3{,}6$ και $3{,}6+1{,}2=4{,}8$.

\tmima{Ασκήσεις}

\zoni{Το κατάλαβα}{γενικός όρος, ν-οστός όρος, θέση}

\ex Συμπλήρωσε τον πίνακα με τη βοήθεια του κανόνα. Σειρά $\nu$: $1,\ 2,\ 3,\ 4,\ 5,\ 6$ και όροι: $9,\ 18,\ 27,\ 36,\ \_\_,\ \_\_$. Ποιος από τους τύπους α) $9+\nu$ \ β) $9\nu$ \ γ) $9+9\nu$ είναι ο γενικός όρος; Ποιος από τους α) $4+\nu$ \ β) $4\nu$ \ γ) $4+4\nu$ \ δ) $8\nu-4$ είναι ο γενικός όρος της σειράς $4,\ 8,\ 12,\ 16,\dots$;

\ex Βρες τους $5$ πρώτους όρους της κανονικότητας με γενικό όρο: α) $8\nu$ \quad β) $11\nu$ \quad γ) $-3\nu$ \quad δ) $0{,}5\nu$.

\ex Συμπλήρωσε τον πίνακα και βρες τον γενικό όρο: α) $6,\ 12,\ 18,\ 24,\ \_\_,\ \_\_$ \quad β) $-2,\ -4,\ -6,\ -8,\ \_\_,\ \_\_$ \quad γ) $0{,}25,\ 0{,}5,\ 0{,}75,\ 1,\ \_\_,\ \_\_$.

\ex Η κανονικότητα έχει γενικό όρο $9\nu$. Βρες τον $10$ο, τον $25$ο και τον $100$ό όρο.

\ex Η κανονικότητα έχει γενικό όρο $6\nu$. Ποιος όρος ισούται με $72$; Με $240$; Με $1\,500$;

\ex Είναι όρος της κανονικότητας $7\nu$ ο αριθμός: α) $91$ \quad β) $150$ \quad γ) $700$ \quad δ) $5$; Αν ναι, ποια είναι η θέση του;

\ex Σωστό ή λάθος; α) Ο γενικός όρος της σειράς $5,\,10,\,15,\dots$ είναι $5+\nu$. \quad β) Ο $4$ος όρος της $3\nu$ είναι $12$. \quad γ) Στην κανονικότητα $2\nu$ ο $1$ος όρος είναι $2$. \quad δ) Ο γενικός όρος της $-3,\,-6,\,-9,\dots$ είναι $-3\nu$. \quad ε) Στον τύπο $4\nu$ το $\nu$ μπορεί να είναι $0{,}5$. \quad στ) Αν ο γενικός όρος είναι $\alpha\nu$, ο $1$ος όρος είναι ο $\alpha$. \quad ζ) Ο $9$ος όρος της $2\nu$ είναι $18$.

\zoni{Το χρησιμοποιώ}{οδοντογλυφίδες, πυραμίδες, παγωτά, ρομπότ}

\ex \textbf{Οι οδοντογλυφίδες.} Ο Βασίλης φτιάχνει ισόπλευρα τρίγωνα με πλευρά $\nu$ οδοντογλυφίδες. α) Πόσες οδοντογλυφίδες χρειάζεται για τρίγωνο με πλευρά $7$; β) Πόση πλευρά έχει το τρίγωνο που χρειάζεται $120$ οδοντογλυφίδες; γ) Υπάρχει τρίγωνο της σειράς που χρειάζεται ακριβώς $100$ οδοντογλυφίδες;

\ex \textbf{Οι πυραμίδες.} Μια πυραμίδα με βάση πολύγωνο $\nu$ πλευρών έχει $\nu$ ακμές στη βάση και $\nu$ ακμές προς την κορυφή. Συμπλήρωσε τον πίνακα για βάση με $3,\,4,\,5,\,6$ πλευρές. Βρες τις ακμές πυραμίδας με βάση $8$-γωνο, $9$-γωνο και $20$-γωνο και τον γενικό όρο.

\ex \textbf{Τα παγωτά.} Ο Νίκος αγοράζει παγωτά των $1{,}5$ € το ένα. Πόσο κοστίζουν $\nu$ παγωτά; Πόσο κοστίζουν $12$; Πόσα παγωτά αγοράζει με $27$ €;

\ex \textbf{Το ρομποτάκι.} Το ρομποτάκι της Ζωής προχωρά $\frac{3}{4}$ του μέτρου κάθε δευτερόλεπτο. Πόσα μέτρα έχει διανύσει μετά από $12$ δευτερόλεπτα; Σε πόσα δευτερόλεπτα θα διανύσει $15$ μέτρα;

\ex \textbf{Το θερμόμετρο.} Σε μια παγωμένη νύχτα η θερμοκρασία ξεκινά από $0$° και πέφτει $1{,}5$° κάθε ώρα. Γράψε τον γενικό όρο της θερμοκρασίας μετά από $\nu$ ώρες. Πόση θα είναι μετά από $8$ ώρες; Πότε θα φτάσει τους $-27$°;

\zoni{Το σκέφτομαι}{ποιος όρος; ποιο βήμα;}

\ex Ο Στέλιος γράφει ότι ο γενικός όρος της $6,\ 12,\ 18,\ 24,\dots$ είναι $6+6\nu$. Πώς μπορείς να ελέγξεις με τον $1$ο όρο ότι κάνει λάθος; Ποιος είναι ο σωστός τύπος;

\ex Μια κανονικότητα έχει γενικό όρο $\alpha\nu$ και ο $4$ος όρος της είναι $10$. Βρες τον $\alpha$, τους δύο πρώτους όρους και τον $50$ό όρο.

\ex Βρες όλους τους φυσικούς $\alpha$ για τους οποίους ο αριθμός $30$ είναι όρος της κανονικότητας $\alpha\nu$.

\ex Η Ζωή γράφει τη σειρά $3,\ 6,\ 9,\ 12,\dots$ και ο Ηλίας τη σειρά $4,\ 7,\ 10,\ 13,\dots$ Και οι δύο έχουν βήμα $3$. Ποια από τις δύο έχει γενικό όρο $3\nu$; Γιατί η άλλη όχι;

\zoni{Για το σπίτι}{πίνακες και καταστάσεις}

\ex Βρες δύο καθημερινές καταστάσεις που περιγράφονται από τον γενικό όρο $12\nu$ και από τον $0{,}5\nu$ (π.χ. τα αυγά σε $\nu$ δωδεκάδες).

\ex Για την κανονικότητα με γενικό όρο $2{,}5\nu$, φτιάξε πίνακα με τους $8$ πρώτους όρους και σημείωσε τα σημεία σε σύστημα αξόνων. Τι παρατηρείς για τη θέση των σημείων;

\epanalipsi{4}

\tmima{Λεξικό του κεφαλαίου}

\renewcommand{\arraystretch}{1.5}
\begin{tabular*}{\linewidth}{@{\extracolsep{\fill}} >{\raggedright\arraybackslash}p{0.28\linewidth} >{\raggedright\arraybackslash}p{0.62\linewidth}}
\toprule
\textbf{Λέξη ή τύπος} & \textbf{Με μία πρόταση}\\
\midrule
αλγεβρική παράσταση & Έκφραση με αριθμούς, γράμματα και πράξεις, π.χ. $3x+5$.\\
εξίσωση & Ισότητα με άγνωστο γράμμα, π.χ. $2x+3=11$.\\
ρίζα, λύση & Ο αριθμός που βάζω στη θέση του αγνώστου και η ισότητα ισχύει.\\
επαλήθευση & Βάζω τη λύση στην αρχική εξίσωση (και στο πρόβλημα) και ελέγχω.\\
ταυτότητα & Ισχύει για κάθε αριθμό, π.χ. $x+3=3+x$.\\
αδύνατη & Δεν ισχύει για κανέναν αριθμό, π.χ. $x+2=x+6$.\\
$x+\alpha=\beta$ και οι άλλες & $x=\beta-\alpha$, $x=\beta+\alpha$, $x=\alpha-\beta$, $x=\beta:\alpha$, $x=\beta\cdot\alpha$, $x=\alpha:\beta$.\\
βήματα επίλυσης προβλήματος & Άγνωστος · εκφράζω · εξίσωση · λύνω · επαληθεύω.\\
\bottomrule
\end{tabular*}

\tmima{Ασκήσεις επανάληψης}

\ex Γράψε με γράμματα το πενταπλάσιο ενός αριθμού ελαττωμένο κατά $3$ και τον επόμενο του διπλάσιου ενός αριθμού $\nu$.

\ex Γράψε απλούστερα την παράσταση $7x+2y-3x+5y$.

\ex Λύσε: α) $x+3{,}4=10$ \quad β) $2{,}5x=15$ \quad γ) $20-x=7{,}5$ \quad δ) $x:4=2{,}5$ \quad ε) $18:x=4$.

\ex Λύσε: α) $3x+7=34$ \quad β) $5x-2=18$ \quad γ) $40-3x=7$ \quad δ) $\frac{x}{4}+2=6$.

\ex Χαρακτήρισε ως ταυτότητα, αδύνατη ή εξίσωση με μία λύση: α) $x+5=5+x$ \quad β) $x+5=x+7$ \quad γ) $2x=0$ \quad δ) $0\cdot x=3$.

\ex \textbf{Η δεξαμενή.} Μια δεξαμενή έχει όγκο $9$ m$^3$, μήκος $2{,}5$ μ. και πλάτος $1{,}2$ μ. Βρες το ύψος της.

\ex \textbf{Το πακέτο.} Το πακέτο κινητής έχει $5$ € πάγιο και $0{,}03$ € ανά λεπτό. Ο λογαριασμός ήταν $9{,}2$ €. Πόσα λεπτά μίλησε;

\ex \textbf{Η τούρτα.} Τα $\frac{3}{4}$ των κομματιών μιας τούρτας είναι $18$. Πόσα κομμάτια έχει η τούρτα;

\ex \textbf{Οι διαδοχικοί.} Βρες τρεις διαδοχικούς φυσικούς με άθροισμα $87$.

\ex \textbf{Η μοιρασιά.} Η Αγγελική και ο Νίκος μοιράζονται $100$ €, ο Νίκος παίρνει $12$ € περισσότερα. Πόσα παίρνει ο καθένας;

\ex \textbf{Οι χειραψίες.} Σε μια συνάντηση $15$ ατόμων όλοι δίνουν το χέρι σε όλους. Πόσες χειραψίες γίνονται;

\ex \textbf{Η αμοιβή.} Ο θείος Λευτέρης δούλεψε $5$ ημέρες και πήρε προκαταβολή τα $\frac{3}{5}$ της αμοιβής του, $240$ €. Πόσο ήταν το μεροκάματο;

\ex \textbf{Ερωτήσεις αυτοαξιολόγησης (εξισώσεις).} Σωστό ή λάθος; α) Η $2x=12$ είναι εξίσωση. \quad β) Η ρίζα της $x-7=5$ είναι το $12$. \quad γ) $3x+4x=7x^2$. \quad δ) Η $x+1=x+1$ είναι ταυτότητα. \quad ε) Η $12-x=5$ δίνει $x=17$. \quad στ) Η $0\cdot x=0$ έχει μία μόνο λύση.

\ex \textbf{Ερωτήσεις αυτοαξιολόγησης (προβλήματα).} Σωστό ή λάθος; α) Το πρώτο βήμα είναι να ονομάσω τον άγνωστο. \quad β) Αρκεί να επαληθευτεί η εξίσωση, όχι το πρόβλημα. \quad γ) Κάθε πρόβλημα λύνεται με εξίσωση. \quad δ) Η $2x+3=11$ έχει λύση το $7$. \quad ε) $\nu$ άτομα κάνουν $\frac{\nu(\nu-1)}{2}$ χειραψίες.

\tmima{Ο χάρτης της επανάληψης}

Μετά τις λύσεις, δες σε ποιες ασκήσεις δυσκολεύτηκες. Ο πίνακας σου λέει ποια παράγραφο να ξαναδιαβάσεις.

\renewcommand{\arraystretch}{1.5}
\begin{tabular*}{\linewidth}{@{\extracolsep{\fill}} >{\raggedright\arraybackslash}p{0.38\linewidth} >{\raggedright\arraybackslash}p{0.52\linewidth}}
\toprule
\textbf{Αν δυσκολεύτηκες στις} & \textbf{ξαναδιάβασε τις παραγράφους}\\
\midrule
Ε4.1, Ε4.2, Ε4.3, Ε4.5, Ε4.13 & 4.1 (έννοια της εξίσωσης)\\
Ε4.4, Ε4.6, Ε4.7, Ε4.8, Ε4.14 & 4.2 (επίλυση προβλημάτων)\\
Ε4.9, Ε4.10, Ε4.11, Ε4.12 & 4.3 (παραδείγματα επίλυσης προβλημάτων)\\
\bottomrule
\end{tabular*}


\kefalaio{Αλγεβρικές παραστάσεις}{Ένα γράμμα που μπορεί να παριστάνει όλους τους αριθμούς}{%
\begin{tabular}{@{}l@{\hspace{1em}}p{0.82\linewidth}@{}}
5.1 & Αλγεβρική παράσταση, αριθμητική τιμή, επιμεριστική ιδιότητα\\[2pt]
 & Επανάληψη του κεφαλαίου\\
\end{tabular}}
\providecommand{\permil}{\mathord{\text{‰}}}
\paragrafos{Ποσοστά: το κλάσμα με παρονομαστή το εκατό}{Α.5.1, «Ποσοστά»}

\begin{anagke}
Στο σχολείο γίνεται ψηφοφορία για την επιτροπή της εκδρομής και η Φανή είναι υποψήφια. Τα τρία τμήματα της πρώτης τάξης ψηφίζουν χωριστά και τα αποτελέσματα, γραμμένα στον πίνακα, είναι: στο Α1 η Φανή πήρε $12$ ψήφους από τους $20$ που ψήφισαν, στο Α2 πήρε $17$ από τους $25$ και στο Α3 πήρε $21$ από τους $35$.

«Στο Α3 με ψήφισαν περισσότεροι», λέει η Φανή περήφανα, «είκοσι ένας!». Ο Νίκος σηκώνει το φρύδι: «Στο Α3 όμως ψήφισαν πολύ περισσότερα παιδιά. Δεν συγκρίνεις το ίδιο πράγμα». Ο Βασίλης προτείνει να τα συγκρίνουν σαν κλάσματα, και η Πολυξένη ρωτάει πώς θα το έλεγαν πιο γρήγορα, με μία ματιά: «Στις ειδήσεις λένε πάντα "το 60\%" ή "το 12\%". Μήπως μπορούμε να κάνουμε όλα τα κλάσματα να έχουν παρονομαστή το εκατό, για να συγκρίνονται εύκολα;».

Το ίδιο ισχύει στη ζωή. Ο Ηλίας είδε στη βιτρίνα «έκπτωση 30\%» και η θεία Μαρία διάβασε σε μια απόδειξη «ΦΠΑ 24\%». Αν είναι να μιλάμε για μέρη ενός συνόλου, χρειαζόμαστε μια κοινή μονάδα σύγκρισης, και η πιο συνηθισμένη είναι το «τοις εκατό».
\end{anagke}

\begin{lexi}{ποσοστό, τοις εκατό}{ποσοστό, τοις εκατό}
Η λέξη \lexiword{ποσοστό} έρχεται από το «πόσο»: ένα \emph{πόσο} μέρος του συνόλου. Το «\lexiword{τοις εκατό}» σημαίνει «ανά εκατό»: το $12\%$ σημαίνει «$12$ για κάθε $100$». Το σύμβολο $\%$ θυμίζει ένα κλάσμα, με το $1$ και το $0$ και $0$ να έχουν κατέβει στη θέση του παρονομαστή. Το \lexiword{τοις χιλίοις} ($\permil$) σημαίνει «ανά χίλια».
\end{lexi}

\begin{metafrasi}
Δώδεκα για κάθε εκατό. & $12\%=\dfrac{12}{100}$\\
Το $35\%$ των $80$ ευρώ. & $\dfrac{35}{100}\cdot80=28$\\
Τρία από τα πέντε ίσα μέρη. & $\dfrac{3}{5}=\dfrac{60}{100}=60\%$\\
Το ποσοστό ως δεκαδικός. & $37\%=0{,}37$\\
Τέσσερα για κάθε χίλια. & $4\permil=\dfrac{4}{1\,000}$\\
\end{metafrasi}

\tmima{Η θεωρία}

\begin{orismos}{ποσοστό επί τοις εκατό}
Το \textbf{ποσοστό $\alpha\%$} (διαβάζεται «$\alpha$ τοις εκατό») είναι ο αριθμός $\dfrac{\alpha}{100}$. Το \textbf{$\alpha\permil$} («$\alpha$ τοις χιλίοις») είναι ο αριθμός $\dfrac{\alpha}{1\,000}$. Το \textbf{$\alpha\%$ του $\beta$} είναι ο αριθμός $\dfrac{\alpha}{100}\cdot\beta$.
\end{orismos}

Άρα κάθε ποσοστό είναι ένα κλάσμα (ή δεκαδικός), και οι κανόνες των κλασμάτων και των δεκαδικών ισχύουν. Μόνο που τώρα, αντί να λέμε «τα $\frac{3}{4}$ των $80$», λέμε «το $75\%$ των $80$».

\tmimasmall{Από κλάσμα σε ποσοστό και ανάποδα}

\begin{itemize}[leftmargin=1.2em,itemsep=0.3ex,topsep=2pt]
\item \textbf{Κλάσμα σε ποσοστό:} κάνω το κλάσμα να έχει παρονομαστή $100$ (ή το διαιρώ και βρίσκω δεκαδικό), και μετά γράφω το αριθμητή με το σύμβολο $\%$. Για παράδειγμα $\frac{4}{5}=\frac{80}{100}=80\%$ και $\frac{3}{8}=0{,}375=37{,}5\%$.
\item \textbf{Δεκαδικός σε ποσοστό:} πολλαπλασιάζω με το $100$ (η υποδιαστολή πάει δύο θέσεις δεξιά): $0{,}52=52\%$, $3{,}41=341\%$, $0{,}07=7\%$.
\item \textbf{Ποσοστό σε κλάσμα:} γράφω $\dfrac{\alpha}{100}$ και απλοποιώ: $12\%=\frac{12}{100}=\frac{3}{25}$.
\item \textbf{Ποιο ποσοστό είναι το $\mu$ του $\nu$:} διαιρώ $\mu:\nu$ και μετατρέπω σε ποσοστό: $15$ € από τα $60$ € είναι $\frac{15}{60}=0{,}25=25\%$.
\end{itemize}

Ένα ποσοστό μπορεί να είναι και μεγαλύτερο από $100\%$: το $150\%$ ενός αριθμού είναι ο ίδιος αριθμός μαζί με το μισό του. Και το $100\%$ ενός αριθμού είναι όλος ο αριθμός.

\tmimasmall{Το ποσοστό στην καθημερινή ζωή}

Πολλές εκφράσεις της καθημερινότητας είναι ποσοστά: «Αύξηση $14\%$» (στον μισθό $100$ € προστίθενται $14$ €), «έκπτωση $60\%$» (από κάθε $100$ € δεν πληρώνουμε $60$), «ΦΠΑ $24\%$» (στην αξία $100$ € προστίθενται $24$ € φόρος), «επιτόκιο $2\%$», «είμαι $100\%$ σίγουρος» (χωρίς καμία αμφιβολία), «πιθανότητα $1\%$» (μια φορά στις εκατό). Για να βρούμε το ποσοστό μιας ποσότητας, την πολλαπλασιάζουμε με το κλάσμα του ποσοστού.

\begin{prosoxi}
\begin{itemize}[leftmargin=1.2em,itemsep=0.3ex,topsep=2pt]
\item Το $\%$ δεν είναι αριθμός από μόνο του: το $30\%$ είναι πάντα το $30\%$ \emph{κάποιου} συνόλου. Το $30\%$ των $50$ δεν είναι ίσο με το $30\%$ των $200$.
\item Όταν συγκρίνω ποσοστά, πρέπει να είναι του ίδιου συνόλου ή να μιλάω για ποσοστά ως τιμές: τα $21$ από τα $35$ δεν είναι περισσότερα «σαν ποσοστό» από τα $17$ από τα $25$, παρότι $21>17$.
\item Στη μετατροπή δεκαδικού σε ποσοστό μετακινώ την υποδιαστολή \emph{δύο} θέσεις δεξιά και όχι μία: $0{,}3=30\%$, όχι $3\%$.
\item Το $\permil$ είναι δέκα φορές μικρότερο από το $\%$: $1\%=10\permil$.
\end{itemize}
\end{prosoxi}

\tmima{Λυμένα παραδείγματα}

\paradeigma{Τα ποσοστά της Φανής}
\tixero Η Φανή πήρε $12$ από $20$ ψήφους στο Α1, $17$ από $25$ στο Α2 και $21$ από $35$ στο Α3.\par
\tithelo Σε ποιο τμήμα την ψήφισαν «περισσότερο».\par
\possindeo Βρίσκω το ποσοστό κάθε τμήματος: $\frac{12}{20}=\frac{60}{100}=60\%$, $\frac{17}{25}=\frac{68}{100}=68\%$ και $\frac{21}{35}=\frac{3}{5}=\frac{60}{100}=60\%$.\par
\apantisi Στο Α2, με $68\%$. Στα Α1 και Α3 πήρε από $60\%$.\par
\elegxos Στο Α3 πήρε περισσότερες ψήφους ($21$), αλλά ψήφισαν και περισσότεροι. Το ποσοστό δίνει το σωστό μέτρο.

\paradeigma{Μετατροπές}
\tixero Τα κλάσματα $\frac{4}{5}$, $\frac{3}{8}$, $\frac{5}{7}$ και τα ποσοστά $15\%$, $64\%$, $2{,}5\%$.\par
\tithelo Τα κλάσματα ως ποσοστά (το $\frac{5}{7}$ στο δέκατο του ποσοστού) και τα ποσοστά ως ανάγωγα κλάσματα.\par
\possindeo $\frac{4}{5}=\frac{80}{100}=80\%$. $\frac{3}{8}=0{,}375=37{,}5\%$. $\frac{5}{7}=0{,}7142\ldots\approx0{,}714=71{,}4\%$. Και $15\%=\frac{15}{100}=\frac{3}{20}$, $64\%=\frac{64}{100}=\frac{16}{25}$, $2{,}5\%=\frac{2{,}5}{100}=\frac{25}{1\,000}=\frac{1}{40}$.\par
\apantisi $80\%$, $37{,}5\%$, περίπου $71{,}4\%$ και $\frac{3}{20}$, $\frac{16}{25}$, $\frac{1}{40}$.\par
\elegxos $\frac{1}{40}=0{,}025=2{,}5\%$.

\paradeigma{Το ποσοστό ενός ποσού}
\tixero Η Ρήνη βλέπει ένα μπλουζάκι των $80$ € με έκπτωση $35\%$ και ένα βιβλίο των $50$ € με ΦΠΑ $24\%$.\par
\tithelo Το ποσό της έκπτωσης και η νέα τιμή, και ο ΦΠΑ και η τιμή με ΦΠΑ.\par
\possindeo Έκπτωση $35\%$ του $80$: $\frac{35}{100}\cdot80=28$ €, νέα τιμή $80-28=52$ €. ΦΠΑ $24\%$ του $50$: $\frac{24}{100}\cdot50=12$ €, τιμή $50+12=62$ €.\par
\apantisi $52$ € και $62$ €.\par
\elegxos Το $35\%$ είναι λίγο περισσότερο από το $\frac{1}{3}$, και το $\frac{1}{3}$ του $80$ είναι περίπου $27$. Συμφωνεί.

\paradeigma{Ποιο ποσοστό είναι;}
\tixero Ο Νίκος έβαλε $14$ από $20$ βολές, ο Βασίλης $9$ από $12$ και ο Στέλιος $5$ από $8$.\par
\tithelo Ποιος έχει το καλύτερο ποσοστό επιτυχίας.\par
\possindeo $\frac{14}{20}=0{,}7=70\%$, $\frac{9}{12}=0{,}75=75\%$, $\frac{5}{8}=0{,}625=62{,}5\%$.\par
\apantisi Ο Βασίλης με $75\%$.\par
\elegxos Ο Νίκος έβαλε τις περισσότερες βολές (14), αλλά ο Βασίλης έχει το καλύτερο ποσοστό.

\paradeigma{Το τοις χιλίοις}
\tixero Ένα εργοστάσιο βγάζει $25\,000$ τεμάχια και τα $4\permil$ είναι ελαττωματικά.\par
\tithelo Πόσα ελαττωματικά είναι.\par
\possindeo $4\permil=\frac{4}{1\,000}$. Άρα $\frac{4}{1\,000}\cdot25\,000=4\cdot25=100$.\par
\apantisi $100$ τεμάχια.\par
\elegxos $100$ από $25\,000$ είναι $\frac{100}{25\,000}=0{,}004=4\permil$.

\tmima{Ασκήσεις}

\zoni{Το κατάλαβα}{μετατροπές και πράξεις με ποσοστά}

\ex Γράψε ως ποσοστά επί τοις εκατό τα κλάσματα: α) $\frac{1}{5}$ \quad β) $\frac{3}{2}$ \quad γ) $\frac{1}{4}$ \quad δ) $\frac{3}{4}$ \quad ε) $\frac{7}{20}$.

\ex Γράψε ως ποσοστά τους δεκαδικούς: α) $0{,}37$ \quad β) $2{,}5$ \quad γ) $0{,}08$ \quad δ) $0{,}6$ \quad ε) $0{,}005$.

\ex Γράψε ως ανάγωγα κλάσματα τα ποσοστά: α) $15\%$ \quad β) $8\%$ \quad γ) $45\%$ \quad δ) $60\%$ \quad ε) $125\%$ \quad στ) $0{,}5\%$.

\ex Υπολόγισε: α) το $10\%$ των $3\,500$ € \quad β) το $25\%$ των $64$ κιλών \quad γ) το $40\%$ της ώρας \quad δ) το $75\%$ του λίτρου \quad ε) το $5\%$ των $240$ \quad στ) το $150\%$ των $30$.

\ex Τι ποσοστό είναι: α) τα $15$ € από τα $60$ € \quad β) τα $45$ λεπτά από την ώρα \quad γ) τα $7$ από τα $28$ \quad δ) τα $9$ παιδιά από τα $24$ \quad ε) τα $330$ από τα $1\,100$.

\ex α) Γράψε το $7\permil$ ως δεκαδικό και ως κλάσμα. β) Βρες το $25\permil$ του $4\,000$. γ) Πόσα τοις χιλίοις είναι το $1\%$;

\ex Σωστό ή λάθος; α) $50\%=\frac{1}{2}$ \quad β) Το $100\%$ ενός αριθμού είναι ο ίδιος ο αριθμός. \quad γ) Το $200\%$ του $30$ είναι $60$. \quad δ) Το $10\%$ του $50$ ισούται με το $50\%$ του $10$. \quad ε) $\frac{1}{4}=14\%$. \quad στ) $0{,}3\%=0{,}3$.

\zoni{Το χρησιμοποιώ}{ψηφοφορίες, μπάσκετ, εκπτώσεις και φόροι}

\ex \textbf{Η ψηφοφορία.} Στη δεύτερη ψηφοφορία ο Βασίλης πήρε $15$ από $25$ ψήφους στο Α1, $24$ από $40$ στο Α2 και $27$ από $36$ στο Α3. Βρες τα ποσοστά και πες σε ποιο τμήμα τον ψήφισαν περισσότερο.

\ex \textbf{Το μπάσκετ.} Η Αγγελική έβαλε $18$ από $24$ βολές και η Ρήνη $13$ από $20$. Ποια έχει το καλύτερο ποσοστό;

\ex \textbf{Τα αθλητικά.} Ένα ζευγάρι αθλητικά κοστίζει $80$ € και έχει έκπτωση $25\%$. Πόσο είναι η έκπτωση και πόσο η νέα τιμή;

\ex \textbf{Το ΦΠΑ.} Το ποδήλατο κοστίζει $150$ € χωρίς ΦΠΑ. Πόσο πληρώνει ο Ηλίας αν ο ΦΠΑ είναι $24\%$;

\ex \textbf{Το μπουκάλι.} Η Ζωή ήπιε το $40\%$ από ένα μπουκάλι $1{,}5$ λίτρων. Πόσο ήπιε και πόσο έμεινε;

\ex \textbf{Το κέρδος του κυλικείου.} Το κυλικείο έβγαλε $450$ € και το $20\%$ ήταν καθαρό κέρδος. Πόσο ήταν το κέρδος;

\ex \textbf{Η τάξη.} Στην Α3 υπάρχουν $25$ παιδιά και τα $11$ είναι αγόρια. Τι ποσοστό της τάξης είναι τα αγόρια και τι τα κορίτσια;

\ex \textbf{Το νερό.} Το νερό είναι περίπου το $60\%$ του βάρους ενός ανθρώπου. Πόσα κιλά νερό υπάρχουν σε ένα παιδί βάρους $45$ κιλών;

\zoni{Το σκέφτομαι}{ποσοστά που ξεγελάνε}

\ex Δώσε ένα παράδειγμα ποσοστού μεγαλύτερου από το $100\%$. Μια τιμή ανέβηκε από $50$ € σε $110$ €. Η νέα τιμή είναι τι ποσοστό της παλιάς;

\ex Ποιο είναι μεγαλύτερο, το $40\%$ του $50$ ή το $50\%$ του $40$; Εξήγησε γιατί.

\ex Μια τιμή των $100$ € αυξάνεται κατά $10\%$ και μετά μειώνεται κατά $10\%$. Επιστρέφει στα $100$ €;

\ex Πώς βρίσκεις γρήγορα με το μυαλό το $10\%$ και το $1\%$ ενός αριθμού; Βρες το $10\%$ και το $1\%$ του $3\,450$ και μετά το $11\%$.

\ex Ο Στέλιος λέει «είμαι $80\%$ σίγουρος ότι θα βρέξει». Τι εννοεί; Και τι θα σήμαινε «$100\%$ σίγουρος», «$0\%$ σίγουρος»;

\zoni{Για το σπίτι}{ποσοστά στα νέα}

\ex Βρες πέντε φράσεις με ποσοστά από εφημερίδα ή ειδήσεις και εξήγησε τι σημαίνει η καθεμία.

\ex Πάρε μια απόδειξη σούπερ μάρκετ και βρες το ποσό του ΦΠΑ που αναγράφεται. Τι ποσοστό της συνολικής αξίας είναι;

\epanalipsi{5}

\tmima{Λεξικό του κεφαλαίου}

\renewcommand{\arraystretch}{1.5}
\begin{tabular*}{\linewidth}{@{\extracolsep{\fill}} >{\raggedright\arraybackslash}p{0.28\linewidth} >{\raggedright\arraybackslash}p{0.62\linewidth}}
\toprule
\textbf{Λέξη ή τύπος} & \textbf{Με μία πρόταση}\\
\midrule
$\alpha\%$ & Ο αριθμός $\dfrac{\alpha}{100}$: «$\alpha$ για κάθε $100$».\\
$\alpha\%$ του $\beta$ & $\dfrac{\alpha}{100}\cdot\beta$.\\
$\alpha\permil$ & Ο αριθμός $\dfrac{\alpha}{1\,000}$.\\
έκπτωση, αύξηση & Πολλαπλασιάζω με $1-\frac{\alpha}{100}$ ή $1+\frac{\alpha}{100}$.\\
ποσοστιαία μεταβολή & $\dfrac{\text{νέα}-\text{παλιά}}{\text{παλιά}}$, πάντα επί της αρχικής.\\
ΦΠΑ & Φόρος που προστίθεται στην αξία· μέσα σε τιμή με $24\%$ είναι τα $\frac{24}{124}$.\\
τόκος & $\text{κεφάλαιο}\cdot\text{επιτόκιο}\cdot\text{χρόνος}$· με κεφαλαιοποίηση ο τόκος προστίθεται στο κεφάλαιο.\\
\bottomrule
\end{tabular*}

\tmima{Ασκήσεις επανάληψης}

\ex Γράψε ως ποσοστά τα $\frac{3}{5}$, $\frac{7}{8}$ και τον $0{,}45$.

\ex Γράψε ως ανάγωγα κλάσματα τα $35\%$ και $12{,}5\%$.

\ex Βρες το $15\%$ των $240$ και το $60\%$ του $35$.

\ex Τι ποσοστό είναι τα $18$ από τα $72$;

\ex Μια τιμή των $140$ € έχει έκπτωση $30\%$. Πόσο κάνει;

\ex Μια τιμή των $250$ € αυξάνεται κατά $12\%$. Πόσο γίνεται;

\ex Πόσο \% αυξήθηκε μια τιμή από $80$ σε $100$;

\ex Μετά από έκπτωση $20\%$ ένα προϊόν κοστίζει $64$ €. Ποια ήταν η αρχική τιμή;

\ex Η τελική τιμή ενός προϊόντος με ΦΠΑ $24\%$ είναι $186$ €. Πόση είναι η αξία χωρίς ΦΠΑ;

\ex Κατάθεση $3\,000$ € με επιτόκιο $2\%$ τον χρόνο. Πόσος είναι ο τόκος ενός έτους και πόσος εννέα μηνών;

\ex Σε κυκλικό διάγραμμα $400$ ψήφων οι υποψήφιοι πήραν $45\%$, $30\%$, $15\%$ και ο τέταρτος το υπόλοιπο. Πόσο ποσοστό και πόσες ψήφους πήρε ο καθένας;

\ex Δύο διαδοχικές εκπτώσεις $10\%$ στα $200$ €. Πόσο κάνει και ποια είναι η συνολική έκπτωση;

\ex \textbf{Ερωτήσεις αυτοαξιολόγησης (ποσοστά).} Σωστό ή λάθος; α) $25\%=\frac{1}{4}$. \quad β) $0{,}07=7\%$. \quad γ) Το $150\%$ του $20$ είναι $30$. \quad δ) $10\permil=10\%$. \quad ε) Το $20\%$ του $50$ ισούται με το $50\%$ του $20$. \quad στ) $0{,}5\%=0{,}5$.

\ex \textbf{Ερωτήσεις αυτοαξιολόγησης (προβλήματα).} Σωστό ή λάθος; α) Έκπτωση $20\%$ και μετά $30\%$ κάνουν $50\%$. \quad β) Αν μια τιμή αυξηθεί κατά $25\%$ και μετά μειωθεί κατά $25\%$, επιστρέφει στην ίδια. \quad γ) Ο ΦΠΑ μέσα στην τιμή με $24\%$ είναι το $24\%$ της τελικής τιμής. \quad δ) Τα ποσοστά ενός πλήρους κυκλικού διαγράμματος κάνουν $100\%$. \quad ε) Ο τόκος ενός έτους είναι κεφάλαιο επί επιτόκιο.

\tmima{Ο χάρτης της επανάληψης}

Μετά τις λύσεις, δες σε ποιες ασκήσεις δυσκολεύτηκες. Ο πίνακας σου λέει ποια παράγραφο να ξαναδιαβάσεις.

\renewcommand{\arraystretch}{1.5}
\begin{tabular*}{\linewidth}{@{\extracolsep{\fill}} >{\raggedright\arraybackslash}p{0.38\linewidth} >{\raggedright\arraybackslash}p{0.52\linewidth}}
\toprule
\textbf{Αν δυσκολεύτηκες στις} & \textbf{ξαναδιάβασε τις παραγράφους}\\
\midrule
Ε5.1, Ε5.2, Ε5.3, Ε5.4, Ε5.13 & 5.1 (ποσοστά)\\
Ε5.5 έως Ε5.12, Ε5.14 & 5.2 (προβλήματα με ποσοστά)\\
\bottomrule
\end{tabular*}


\kefalaio{Ανάλογα ποσά – Αντιστρόφως ανάλογα ποσά}{Όταν ένα ποσό αλλάζει, ποιο άλλο το ακολουθεί;}{%
\begin{tabular}{@{}l@{\hspace{1em}}p{0.82\linewidth}@{}}
6.1 & Παράσταση σημείων στο επίπεδο\\[2pt]
6.2 & Λόγος δύο αριθμών – Αναλογία\\[2pt]
6.3 & Ανάλογα ποσά – Ιδιότητες αναλόγων ποσών\\[2pt]
6.4 & Γραφική παράσταση σχέσης αναλογίας\\[2pt]
6.5 & Προβλήματα αναλογιών\\[2pt]
6.6 & Αντιστρόφως ανάλογα ποσά\\[2pt]
 & Επανάληψη του κεφαλαίου\\
\end{tabular}}
\paragrafos{Ιδιότητες της ισότητας και η έννοια της εξίσωσης}{Α.6.1, «Ιδιότητες της ισότητας (ισότητα και πράξεις) και η έννοια της εξίσωσης»}

\begin{anagke}
Στο εργαστήριο του σχολείου ο Ηλίας βρίσκει μια παλιά ζυγαριά με δύο δίσκους. Στον αριστερό δίσκο βάζει ένα κλειστό σακουλάκι με φουντούκια, που δεν ξέρει πόσο ζυγίζει, και ένα βαρίδι των $100$ g. Στον δεξί δίσκο βάζει βαρίδια που ζυγίζουν συνολικά $350$ g. Η ζυγαριά ισορροπεί. «Πόσο ζυγίζουν τα φουντούκια;» ρωτάει η Ρήνη.

Ο Στέλιος αρπάζει το βαρίδι των $100$ g από τον αριστερό δίσκο. Η ζυγαριά γέρνει δεξιά. «Δεν πειράζει», λέει η Φανή, «αν αφαιρέσουμε $100$ g και από τον άλλο δίσκο, θα ξαναισορροπήσει». Έτσι κάνουν. Στον αριστερό μένει το σακουλάκι και στον δεξιό $250$ g. Η Πολυξένη το γράφει στο τετράδιό της και το κάνει ισότητα: «Πρόσθεσα ή αφαίρεσα το ίδιο και από τις δύο πλευρές, και η ισότητα έμεινε αληθινή». Ο Βασίλης αμφιβάλλει: «Και αν ήταν πολλαπλασιασμός; Ή διαίρεση;».

Στην ισότητα αυτή υπάρχει ένας αριθμός που δεν ξέρουμε, το βάρος του σακουλιού. Σε αυτή την παράγραφο θα μάθουμε ποιες πράξεις μπορούμε να κάνουμε χωρίς να χαλάσουμε μια ισότητα, και θα γνωρίσουμε την εξίσωση, δηλαδή την ισότητα που κρύβει έναν άγνωστο αριθμό.
\end{anagke}

\begin{lexi}{ισότητα, εξίσωση, άγνωστος, ρίζα, επαλήθευση}{ισότητα, εξίσωση, ρίζα}
Η λέξη \lexiword{εξίσωση} βγαίνει από το «ίσος»: εξισώνω σημαίνει «κάνω ίσα» τα δύο μέλη, όπως οι δύο δίσκοι μιας ζυγαριάς. \lexiword{Άγνωστος} είναι ο αριθμός που δεν ξέρουμε και ψάχνουμε. Η \lexiword{ρίζα} της εξίσωσης είναι η λύση της, και λέγεται έτσι γιατί «βγαίνει από το βάθος», όπως η ρίζα ενός φυτού. Η \lexiword{επαλήθευση} είναι ο έλεγχος ότι κάτι είναι αληθινό («επί» + «αληθεύω»).
\end{lexi}

\begin{metafrasi}
Η ζυγαριά ισορροπεί. & $\alpha=\beta$\\
Βάζω το ίδιο βάρος και στους δύο δίσκους. & $\alpha+\gamma=\beta+\gamma$\\
Βγάζω το ίδιο βάρος και από τους δύο δίσκους. & $\alpha-\gamma=\beta-\gamma$\\
Διπλασιάζω ή τριπλασιάζω ό,τι υπάρχει και στους δύο δίσκους. & $\alpha\cdot\gamma=\beta\cdot\gamma$\\
Χωρίζω και τους δύο δίσκους σε ίσα μέρη. & $\alpha:\gamma=\beta:\gamma$, με $\gamma\neq0$\\
Σακουλάκι και $100$ g ισορροπούν με $350$ g. & $x+100=350$\\
\end{metafrasi}

\tmima{Η θεωρία}

\tmimasmall{Ισότητα και ισορροπία}

Μια ισότητα $\alpha=\beta$ λέει ότι οι δύο πλευρές της, που λέγονται \textbf{μέλη}, έχουν την ίδια τιμή. Η ζυγαριά που ισορροπεί είναι η καλύτερη εικόνα για αυτό.

\begin{center}
\scalebox{0.78}{\binzyg{\binsx\ \binsw{100 g}}{\binsw{350 g}}}
\quad\raisebox{6mm}{$\Rightarrow$}\quad
\scalebox{0.78}{\binzyg{\binsx}{\binsw{250 g}}}
\end{center}

Αν βγάλουμε $100$ g από τον έναν δίσκο και $100$ g από τον άλλον, η ζυγαριά ισορροπεί ακόμη και φαίνεται ότι $x=250$ g.

\begin{idiotita}{ισότητα και πράξεις}
Για οποιουσδήποτε αριθμούς $\alpha,\beta,\gamma$ ισχύει:
\begin{itemize}[leftmargin=1.2em,itemsep=0.3ex,topsep=2pt]
\item Αν $\alpha=\beta$, τότε $\alpha+\gamma=\beta+\gamma$ (προσθέτουμε τον ίδιο αριθμό και στα δύο μέλη).
\item Αν $\alpha=\beta$, τότε $\alpha-\gamma=\beta-\gamma$ (αφαιρούμε τον ίδιο αριθμό και από τα δύο μέλη).
\item Αν $\alpha=\beta$, τότε $\alpha\cdot\gamma=\beta\cdot\gamma$ (πολλαπλασιάζουμε και τα δύο μέλη με τον ίδιο αριθμό).
\item Αν $\alpha=\beta$, τότε $\alpha:\gamma=\beta:\gamma$ (διαιρούμε και τα δύο μέλη με τον ίδιο αριθμό \textbf{που δεν είναι μηδέν}).
\end{itemize}
\end{idiotita}

Για παράδειγμα, από την $8=8$ προκύπτουν οι $8+5=8+5$, $8\cdot3=8\cdot3$ και $8:2=8:2$, δηλαδή $13=13$, $24=24$ και $4=4$. Οι ιδιότητες αυτές είναι το εργαλείο μας για όλη την Άλγεβρα.

\begin{prosoxi}
Δεν διαιρούμε με το $0$. Από την $0\cdot3=0\cdot5$ (που ισχύει, γιατί $0=0$) δεν μπορούμε να «διαιρέσουμε με το $0$» και να καταλήξουμε στο $3=5$. Για αυτό στην τελευταία ιδιότητα απαιτείται $\gamma\neq0$.
\end{prosoxi}

\tmimasmall{Η εξίσωση}

\begin{orismos}{εξίσωση με έναν άγνωστο}
\textbf{Εξίσωση} με έναν άγνωστο λέγεται μια ισότητα που περιέχει αριθμούς και ένα γράμμα (τον άγνωστο), π.χ. $x+5=15$ ή $2\cdot y=16$.\\
\textbf{Λύση ή ρίζα} της εξίσωσης είναι ο αριθμός που, όταν αντικαταστήσει τον άγνωστο, επαληθεύει την ισότητα. Η διαδικασία εύρεσης της λύσης λέγεται \textbf{επίλυση} της εξίσωσης.
\end{orismos}

Ο αριθμός $10$ είναι λύση της $x+5=15$, γιατί όταν αντικαταστήσουμε το $x$ με $10$ προκύπτει $10+5=15$, που ισχύει. Ο αριθμός $7$ δεν είναι λύση, γιατί $7+5=12\neq15$. Μετά από κάθε επίλυση κάνουμε \textbf{επαλήθευση}: βάζουμε τη λύση στην \emph{αρχική} εξίσωση και ελέγχουμε ότι βγαίνει σωστή ισότητα.

\tmimasmall{Ειδικές περιπτώσεις}

\begin{itemize}[leftmargin=1.2em,itemsep=0.3ex,topsep=2pt]
\item Η εξίσωση $0\cdot x=2$ δεν επαληθεύεται από κανέναν αριθμό, γιατί κάθε αριθμός επί $0$ κάνει $0$. Μια εξίσωση χωρίς καμία λύση λέγεται \textbf{αδύνατη}.
\item Η εξίσωση $0\cdot x=0$ επαληθεύεται από όλους τους αριθμούς. Μια εξίσωση με λύσεις όλους τους αριθμούς λέγεται \textbf{ταυτότητα} ή \textbf{αόριστη}.
\end{itemize}

\tmimasmall{Ισοδύναμες εξισώσεις και οι έξι βασικές μορφές}

Δύο εξισώσεις που έχουν τις ίδιες λύσεις λέγονται \textbf{ισοδύναμες}. Οι ιδιότητες της ισότητας μας επιτρέπουν να αντικαθιστούμε μια εξίσωση με μια ισοδύναμη, πιο απλή. Από τον ορισμό των πράξεων (η αφαίρεση είναι η αντίστροφη της πρόσθεσης και η διαίρεση του πολλαπλασιασμού) βρίσκουμε τη λύση κάθε μίας από τις παρακάτω εξισώσεις.

\begin{center}
\renewcommand{\arraystretch}{1.5}
\begin{tabular}{@{}lll@{}}
\toprule
\textbf{Εξίσωση} & \textbf{Λύση} & \textbf{Παράδειγμα}\\
\midrule
$x+\alpha=\beta$ & $x=\beta-\alpha$ & $x+2=5\ \Rightarrow\ x=5-2=3$\\
$x-\alpha=\beta$ & $x=\beta+\alpha$ & $x-6=2\ \Rightarrow\ x=2+6=8$\\
$\alpha-x=\beta$ & $x=\alpha-\beta$ & $10-x=3\ \Rightarrow\ x=10-3=7$\\
$\alpha\cdot x=\beta$ & $x=\beta:\alpha$ & $2x=40\ \Rightarrow\ x=40:2=20$\\
$x:\alpha=\beta$ & $x=\beta\cdot\alpha$ & $x:5=6\ \Rightarrow\ x=6\cdot5=30$\\
$\alpha:x=\beta$ & $x=\alpha:\beta$ & $20:x=5\ \Rightarrow\ x=20:5=4$\\
\bottomrule
\end{tabular}
\end{center}

(Στις εξισώσεις με διαίρεση υποθέτουμε ότι δεν διαιρούμε με το $0$.) Οι αριθμοί των εξισώσεων μπορεί να είναι και ρητοί, με πρόσημο, δεκαδικοί ή κλάσματα.

\begin{prosoxi}
\begin{itemize}[leftmargin=1.2em,itemsep=0.3ex,topsep=2pt]
\item Στην $\alpha-x=\beta$ η λύση είναι $\alpha-\beta$, όχι $\beta-\alpha$ ή $\alpha+\beta$. Η $12-x=5$ δίνει $x=7$, όχι $17$.
\item Στις $\alpha:x=\beta$ και $x:\alpha=\beta$ οι λύσεις διαφέρουν: $24:x=6$ δίνει $x=24:6=4$, ενώ $x:6=24$ δίνει $x=24\cdot6=144$.
\item Στην επαλήθευση βάζουμε τη λύση στην \emph{αρχική} εξίσωση, όχι σε κάποια που προέκυψε στη μέση της λύσης μας.
\end{itemize}
\end{prosoxi}

\tmima{Λυμένα παραδείγματα}

\paradeigma{Η ζυγαριά του Ηλία}
\tixero Σακουλάκι με $x$ g φουντούκια και βαρίδι $100$ g ισορροπούν με $350$ g.\par
\tithelo Το βάρος των φουντουκιών.\par
\possindeo Η ισότητα είναι $x+100=350$. Αφαιρώ $100$ και από τα δύο μέλη (ιδιότητα της αφαίρεσης): $x+100-100=350-100$, δηλαδή $x=250$.\par
\apantisi $250$ g.\par
\elegxos $250+100=350$. Σωστό.

\paradeigma{Είναι ο αριθμός λύση;}
\tixero Η εξίσωση $2x=16$ και οι αριθμοί $6$ και $8$. Η εξίσωση $x-4=1$ και οι αριθμοί $3$ και $5$.\par
\tithelo Ποιοι από αυτούς είναι ρίζες.\par
\possindeo Αντικαθιστώ. Για $x=6$: $2\cdot6=12\neq16$. Για $x=8$: $2\cdot8=16$ σωστό. Για $x=3$: $3-4=-1\neq1$. Για $x=5$: $5-4=1$ σωστό.\par
\apantisi Ρίζες είναι ο $8$ και ο $5$.\par
\elegxos Από τους κανόνες: $x=16:2=8$ και $x=1+4=5$.

\paradeigma{Λύνω τις έξι βασικές εξισώσεις}
\tixero Οι εξισώσεις $x+2{,}5=7$, $x-\frac{3}{4}=\frac{1}{4}$, $15-x=6{,}4$, $4x=18$, $x:6=2{,}5$, $36:x=9$.\par
\tithelo Να τις λύσω.\par
\possindeo $x=7-2{,}5=4{,}5$. $x=\frac{1}{4}+\frac{3}{4}=1$. $x=15-6{,}4=8{,}6$. $x=18:4=4{,}5$. $x=2{,}5\cdot6=15$. $x=36:9=4$.\par
\apantisi $4{,}5$, $1$, $8{,}6$, $4{,}5$, $15$, $4$.\par
\elegxos $4{,}5+2{,}5=7$, $1-\frac{3}{4}=\frac{1}{4}$, $15-8{,}6=6{,}4$, $4\cdot4{,}5=18$, $15:6=2{,}5$, $36:4=9$. Όλα σωστά.

\paradeigma{Ρητοί αριθμοί με πρόσημο}
\tixero Οι εξισώσεις $-5x=25$ και $x+(-8)=-12$ και $-8-x=11$.\par
\tithelo Να τις λύσω και να επαληθεύσω.\par
\possindeo Στην πρώτη διαιρώ και τα δύο μέλη με $-5$: $x=25:(-5)=-5$. Στη δεύτερη προσθέτω $8$: $x=-12+8=-4$. Στην τρίτη, αν προσθέσω $8$ και στα δύο μέλη, παίρνω $-x=19$, άρα $x=-19$.\par
\apantisi $-5$, $-4$, $-19$.\par
\elegxos $-5\cdot(-5)=25$. $-4+(-8)=-12$. $-8-(-19)=-8+19=11$. Σωστά.

\paradeigma{Ταυτότητα ή αδύνατη;}
\tixero Οι εξισώσεις $x+3=3+x$, $x+2=x+6$ και $5x-5x=7$.\par
\tithelo Τι είδους είναι.\par
\possindeo Η πρώτη ισχύει για κάθε αριθμό (αντιμεταθετική ιδιότητα): ταυτότητα. Στη δεύτερη το αριστερό μέλος είναι πάντα κατά $4$ μικρότερο από το δεξί: αδύνατη. Στην τρίτη το $5x-5x=0$ και δεν ισούται με $7$: αδύνατη.\par
\apantisi Ταυτότητα, αδύνατη, αδύνατη.\par
\elegxos Με $x=10$: $13=13$, αλλά $12\neq16$ και $0\neq7$.

\paradeigma{Η δεξαμενή του χωριού}
\tixero Μια δεξαμενή έχει όγκο $7{,}2$ m$^3$, μήκος $2$ μ. και πλάτος $1{,}2$ μ.\par
\tithelo Το ύψος της.\par
\possindeo Αν $x$ είναι το ύψος σε μέτρα, ο όγκος είναι $2\cdot1{,}2\cdot x=2{,}4x$. Άρα $2{,}4x=7{,}2$ και $x=7{,}2:2{,}4=3$.\par
\apantisi $3$ μ.\par
\elegxos $2\cdot1{,}2\cdot3=7{,}2$. Σωστό.

\tmima{Ασκήσεις}

\zoni{Το κατάλαβα}{ιδιότητες, ρίζες, βασικές εξισώσεις}

\ex Σωστό ή λάθος; α) Αν $\alpha=\beta$, τότε $\alpha+5=\beta+5$. \quad β) Αν $x-3=7$, τότε $x=4$. \quad γ) Η εξίσωση $0\cdot x=3$ είναι αδύνατη. \quad δ) Αν $3\alpha=0$, τότε $\alpha=0$. \quad ε) Αν $\dfrac{\alpha}{3}=\dfrac{\beta}{3}$, τότε $\alpha=\beta$. \quad στ) Η εξίσωση $2x=6$ έχει λύση το $6$.

\ex \textbf{Η ζυγαριά.} Στην εικόνα η ζυγαριά ισορροπεί. Γράψε την εξίσωση και βρες πόσο ζυγίζει το σακουλάκι.
\begin{center}
\binzyg{\binsx\ \binsw{200 g}}{\binsw{500 g}}
\end{center}

\ex Αξιοποίησε τις ιδιότητες της ισότητας και συμπλήρωσε τα κενά. α) Αν $\alpha+2=7$, τότε $\alpha=\dots$ \quad β) Αν $\kappa-3=10$, τότε $\kappa=\dots$ \quad γ) Αν $3=\lambda$, τότε $\dots=3$ \quad δ) Αν $9\rho=45$, τότε $\rho=\dots$

\ex Εξέτασε αν οι αριθμοί είναι λύσεις των εξισώσεων: α) $x+3=9$, $x=6$ \quad β) $x-2=7$, $x=9$ \quad γ) $2x=10$, $x=20$ \quad δ) $12-x=10$, $x=22$ \quad ε) $x:4=3$, $x=12$.

\ex Λύσε τις εξισώσεις: α) $x+12=45$ \quad β) $9+y=17$ \quad γ) $x-48=5$ \quad δ) $65-x=17$ \quad ε) $15y=45$ \quad στ) $z:5=40$ \quad ζ) $26:x=13$.

\ex Λύσε τις εξισώσεις και επαλήθευσε: α) $-5x=25$ \quad β) $x+(-8)=-12$ \quad γ) $-6x=-18$ \quad δ) $-8-x=11$.

\ex Λύσε: α) $x+0{,}75=3$ \quad β) $\dfrac{1}{2}x=\dfrac{9}{4}$ \quad γ) $x-\dfrac{1}{3}=\dfrac{5}{6}$ \quad δ) $x:\dfrac{3}{5}=10$.

\zoni{Το χρησιμοποιώ}{ζυγαριές, μυστικοί αριθμοί, σχήματα}

\ex \textbf{Ο μυστικός αριθμός.} Η Πολυξένη λέει: «Σκέφτηκα έναν αριθμό». Γράψε εξίσωση και βρες τον, όταν: α) πρόσθεσε $17$ και βρήκε $52$ \quad β) τον πολλαπλασίασε επί $3$ και βρήκε $84$ \quad γ) τον αφαίρεσε από το $50$ και βρήκε $18$ \quad δ) τον διαίρεσε με $2$ και βρήκε $7$.

\ex \textbf{Τα βαρίδια.} Στη ζυγαριά της θείας Μαρίας δύο ίδια μπουκάλια ζυγίζουν όσο ένα βαρίδι $1\,200$ g. Γράψε την εξίσωση και βρες το βάρος κάθε μπουκαλιού.

\ex \textbf{Τα εισιτήρια και ο κουμπαράς.} α) Τέσσερα ίδια εισιτήρια πλοίου κοστίζουν $54$ €. Πόσο κάνει το καθένα; β) Ο Ηλίας είχε στον κουμπαρά $36$ €, ξόδεψε κάποιο ποσό και του έμειναν $11{,}5$ €. Πόσα ξόδεψε;

\ex \textbf{Οι περίμετροι.} α) Η Αγγελική φτιάχνει μια τετράγωνη κορνίζα με περίμετρο $98$ cm. Πόση είναι η πλευρά της; β) Ένα ισόπλευρο τρίγωνο έχει περίμετρο $4{,}5$ m. Πόση είναι η πλευρά του;

\ex Βρες το είδος κάθε εξίσωσης (ένας αριθμός λύση, αδύνατη ή ταυτότητα): α) $x+3=3+x$ \quad β) $0\cdot x=7$ \quad γ) $x+2=x+6$ \quad δ) $5x-5x=0$.

\zoni{Το σκέφτομαι}{γιατί δουλεύει, πού κάνουμε λάθος}

\ex Γιατί η $0\cdot x=0$ έχει άπειρες λύσεις και η $0\cdot x=5$ καμία; Εξήγησε με το «ποιος αριθμός επί $0$ δίνει...».

\ex Ο Στέλιος λύνει την $12-x=5$ και γράφει $x=12+5=17$. Λύνει και την $24:x=6$ και γράφει $x=24\cdot6$. Βρες τα λάθη του, γράψε τις σωστές λύσεις και επαλήθευσε.

\ex Γράψε μια εξίσωση με λύση το $6$, μια αδύνατη και μια ταυτότητα.

\ex Ποια ζεύγη εξισώσεων είναι ισοδύναμα; α) $x+3=7$ και $2x=8$ \quad β) $x-1=2$ και $x+1=2$ \quad γ) $x+1=x+2$ και $0\cdot x=5$. Εξήγησε.

\zoni{Για το σπίτι}{φτιάξε το δικό σου}

\ex Ο παππούς Γιάννης λέει: «Σκέφτηκα έναν αριθμό και το τριπλάσιό του είναι $2{,}4$». Ποιος είναι; Γράψε την εξίσωση.

\ex Φτιάξε μια καθημερινή ιστορία που να οδηγεί στην εξίσωση $x+12=30$ και ζήτησε από ένα μέλος της οικογένειάς σου να τη λύσει.

\paragrafos{Επίλυση εξίσωσης της μορφής αx+β=γ}{Α.6.2, «Επίλυση εξίσωσης της μορφής $\alpha x+\beta=\gamma$»}

\begin{anagke}
Η θεία Μαρία δείχνει στην Ζωή τον τελευταίο λογαριασμό του κινητού της: $15$ ευρώ. «Το πακέτο μου έχει πάγιο $9$ ευρώ τον μήνα και κάθε λεπτό ομιλίας εκτός πακέτου κοστίζει $0{,}03$ ευρώ», λέει. «Πόσα λεπτά μίλησα;». Η Ζωή δοκιμάζει: «Εκατό λεπτά, τριάντα εκατό... όχι, τρία ευρώ... και πάλι λάθος». Ο Ηλίας γράφει: αν $x$ είναι τα λεπτά, ο λογαριασμός είναι $9+0{,}03x$, και αυτός ισούται με $15$.

Η Φανή ρωτάει: «Η $9+0{,}03x=15$ δεν είναι σαν τις εξισώσεις της ζυγαριάς. Έχει δύο πράξεις γύρω από το $x$». Η κυρία Αλεξάνδρα απαντά με ερώτηση: «Ποια πράξη έγινε στο $x$ πρώτη και ποια δεύτερη; Και αν γυρίσουμε τη ζυγαριά πίσω, με ποια σειρά θα τις αναιρέσουμε;».

Σε αυτή την παράγραφο θα μάθουμε να λύνουμε εξισώσεις της μορφής $\alpha x+\beta=\gamma$ με τις ιδιότητες της ισότητας, να ελέγχουμε τη λύση με επαλήθευση και να μεταφράζουμε προβλήματα της ζωής σε τέτοιες εξισώσεις.
\end{anagke}

\begin{lexi}{επίλυση, ισοδύναμες, πρόβλημα}{επίλυση, πρόβλημα}
Η \lexiword{επίλυση} είναι η διαδικασία με την οποία «λύνουμε» τον κόμπο, δηλαδή βρίσκουμε τη λύση. Το \lexiword{πρόβλημα} έρχεται από το «προ-βάλλω», δηλαδή «βάζω μπροστά σου» ένα εμπόδιο. Οι εξισώσεις που προκύπτουν σε κάθε βήμα της επίλυσης λέγονται \lexiword{ισοδύναμες}, γιατί έχουν την ίδια λύση.
\end{lexi}

\begin{metafrasi}
Ο λογαριασμός είναι $15$ ευρώ. & $9+0{,}03x=15$\\
Αφαιρώ και από τα δύο μέλη το $9$. & $0{,}03x=6$\\
Διαιρώ και τα δύο μέλη με το $0{,}03$. & $x=200$\\
Το διπλάσιο ενός αριθμού αυξημένο κατά $3$ κάνει $15$. & $2x+3=15$\\
Πέντε ίσα μέρη και $10$ ακόμη κάνουν $70$. & $5x+10=70$\\
\end{metafrasi}

\tmima{Η θεωρία}

\tmimasmall{Η μέθοδος}

Στην εξίσωση $\alpha x+\beta=\gamma$ ο άγνωστος πολλαπλασιάστηκε με το $\alpha$ και μετά προστέθηκε το $\beta$. Για να τον ελευθερώσουμε αναιρούμε τις πράξεις με την αντίστροφη σειρά, με τις ιδιότητες της ισότητας:
\begin{gather*}
\alpha x+\beta=\gamma\\
\alpha x+\beta-\beta=\gamma-\beta\qquad(\text{αφαιρώ και από τα δύο μέλη το }\beta)\\
\alpha x=\gamma-\beta\qquad(\text{κάνω τις πράξεις})\\
\frac{\alpha x}{\alpha}=\frac{\gamma-\beta}{\alpha}\qquad(\text{διαιρώ και τα δύο μέλη με το }\alpha,\ \alpha\neq0)\\
x=\frac{\gamma-\beta}{\alpha}.
\end{gather*}
Όλες οι εξισώσεις που γράψαμε είναι \textbf{ισοδύναμες}: έχουν την ίδια λύση. Η τελευταία δείχνει καθαρά τη λύση.

\begin{center}
\binzyg{\binsx\ \binsx\ \binsw{6}}{\binsw{14}}
\quad$\Longrightarrow$\quad
\binzyg{\binsx\ \binsx}{\binsw{8}}
\quad$\Longrightarrow$\quad
\binzyg{\binsx}{\binsw{4}}
\end{center}
Η ζυγαριά δείχνει την εξίσωση $2x+6=14$. Βγάζω $6$ από κάθε δίσκο και μένει $2x=8$. Χωρίζω κάθε δίσκο σε δύο ίσα μέρη και μένει $x=4$.

\begin{kanonas}{λύση της $\alpha x+\beta=\gamma$}
1. Αφαιρώ και από τα δύο μέλη το $\beta$ (ή προσθέτω, αν το $\beta$ είναι αρνητικό).\\
2. Κάνω τις πράξεις.\\
3. Διαιρώ και τα δύο μέλη με τον συντελεστή $\alpha$ του αγνώστου.\\
4. \textbf{Επαληθεύω} στην αρχική εξίσωση.
\end{kanonas}

\tmimasmall{Όταν υπάρχουν παρενθέσεις ή όμοιοι όροι}

Αν στην εξίσωση υπάρχουν παρενθέσεις, κάνουμε πρώτα την επιμεριστική ιδιότητα και την αναγωγή ομοίων όρων σε κάθε μέλος (Α5), και μετά την αντιμετωπίζουμε ως $\alpha x+\beta=\gamma$. Για παράδειγμα η $3(x-2)+2(x+1)=6$ γίνεται $3x-6+2x+2=6$, δηλαδή $5x-4=6$.

\tmimasmall{Όταν ο άγνωστος χάνεται}

Αν ο συντελεστής $\alpha$ γίνει $0$, δεν μπορούμε να διαιρέσουμε. Η εξίσωση $0\cdot x=\delta$ είναι τότε αδύνατη (αν $\delta\neq0$) ή ταυτότητα (αν $\delta=0$), όπως είδαμε στην προηγούμενη παράγραφο.

\tmimasmall{Προβλήματα με εξίσωση}

\begin{enumerate}[leftmargin=1.6em,itemsep=0.3ex,topsep=2pt]
\item Διαβάζω και βρίσκω τι ζητάω. Τον ονομάζω με ένα γράμμα και γράφω τι σημαίνει, με τις μονάδες του: «Έστω $x$ η τιμή ενός εισιτηρίου, σε ευρώ».
\item Εκφράζω τα άλλα στοιχεία με τη βοήθεια του αγνώστου. Εδώ βοηθάει ένα σχέδιο, π.χ. ένα ορθογώνιο \emph{μοντέλο ράβδου} (δες το παράδειγμα με το ταμείο της εκδρομής).
\item Γράφω την εξίσωση, δηλαδή τη φράση «το ένα ποσό ισούται με το άλλο».
\item Λύνω την εξίσωση.
\item Επαληθεύω στο \emph{πρόβλημα}, όχι μόνο στην εξίσωση, και απαντώ με λόγια και μονάδες.
\end{enumerate}

\begin{prosoxi}
\begin{itemize}[leftmargin=1.2em,itemsep=0.3ex,topsep=2pt]
\item Στην $2x+3=11$ αφαιρώ πρώτα το $3$ και μετά διαιρώ με το $2$: $x=(11-3):2=4$. Δεν γράφω $x=11+3=14$ και δεν διαιρώ το $11$ με το $2$ πριν από την αφαίρεση.
\item Όταν διαιρώ με τον συντελεστή, τον διαιρώ και με το πρόσημό του: από την $-3x=12$ βγαίνει $x=-4$, όχι $4$.
\item Μια λύση που δεν βγάζει νόημα στο πρόβλημα (π.χ. $-5$ μαθητές ή $2{,}5$ παιδιά) σημαίνει ότι κάτι δεν πήγε καλά.
\end{itemize}
\end{prosoxi}

\tmima{Λυμένα παραδείγματα}

\paradeigma{Δύο βήματα}
\tixero Η εξίσωση $2x+6=14$.\par
\tithelo Να τη λύσω και να επαληθεύσω.\par
\possindeo $2x+6=14$, άρα $2x+6-6=14-6$, δηλαδή $2x=8$. Διαιρώ με το $2$: $x=4$.\par
\apantisi $x=4$.\par
\elegxos $2\cdot4+6=14$. Σωστό.

\paradeigma{Αρνητικά, δεκαδικά, κλάσματα}
\tixero Οι εξισώσεις $7-3x=-5$, $\dfrac{1}{2}x-\dfrac{3}{4}=\dfrac{1}{4}$ και $1{,}5x-2{,}25=3$.\par
\tithelo Να τις λύσω.\par
\possindeo Πρώτη: $-3x=-5-7=-12$, άρα $x=-12:(-3)=4$. Δεύτερη: $\frac{1}{2}x=\frac{1}{4}+\frac{3}{4}=1$, άρα $x=2$. Τρίτη: $1{,}5x=3+2{,}25=5{,}25$, άρα $x=5{,}25:1{,}5=3{,}5$.\par
\apantisi $4$, $2$ και $3{,}5$.\par
\elegxos $7-12=-5$. $\frac{1}{2}\cdot2-\frac{3}{4}=\frac{1}{4}$. $1{,}5\cdot3{,}5-2{,}25=5{,}25-2{,}25=3$.

\paradeigma{Με παρενθέσεις}
\tixero Η εξίσωση $3(x-2)+2(x+1)=6$.\par
\tithelo Να τη λύσω.\par
\possindeo Επιμεριστική ιδιότητα: $3x-6+2x+2=6$. Αναγωγή ομοίων όρων: $5x-4=6$. Προσθέτω $4$: $5x=10$. Διαιρώ με $5$: $x=2$.\par
\apantisi $x=2$.\par
\elegxos $3\cdot0+2\cdot3=6$. Σωστό.

\paradeigma{Ο άγνωστος χάνεται}
\tixero Οι εξισώσεις $2(x+3)=2x+5$ και $3(x+1)-3x=3$.\par
\tithelo Να τις λύσω.\par
\possindeo Η πρώτη γίνεται $2x+6=2x+5$ και, αφαιρώντας $2x$ και από τα δύο μέλη, $6=5$, που δεν ισχύει ποτέ. Η δεύτερη γίνεται $3x+3-3x=3$, δηλαδή $3=3$, που ισχύει πάντα.\par
\apantisi Η πρώτη είναι αδύνατη και η δεύτερη ταυτότητα.\par
\elegxos Με $x=1$ η πρώτη δίνει $8\neq7$ και η δεύτερη $3=3$. Με $x=5$ η πρώτη $16\neq15$ και η δεύτερη $3=3$.

\paradeigma{Ο λογαριασμός του κινητού}
\tixero Το πακέτο της θείας Μαρίας κοστίζει $9$ € πάγιο και $0{,}03$ € το λεπτό. Ο λογαριασμός ήταν $15$ €.\par
\tithelo Πόσα λεπτά μίλησε.\par
\possindeo Έστω $x$ τα λεπτά. Ο λογαριασμός είναι $9+0{,}03x=15$. Άρα $0{,}03x=6$ και $x=6:0{,}03=200$.\par
\apantisi $200$ λεπτά.\par
\elegxos $9+0{,}03\cdot200=9+6=15$ €. Σωστό.

\paradeigma{Το ταμείο της εκδρομής}
\tixero Τα $24$ παιδιά της Α3 πληρώνουν από $x$ ευρώ το καθένα και το σχολείο δίνει ακόμη $60$ €. Το σύνολο είναι $420$ €.\par
\tithelo Πόσο πληρώνει κάθε παιδί.\par
\possindeo Το μοντέλο ράβδου:
\begin{center}
\begin{tikzpicture}[x=0.19cm]
\foreach \i in {0,...,23}{\draw[fill=gblue!25] (\i*1.0,0) rectangle (\i*1.0+1,0.5);}
\draw[fill=black!15] (24,0) rectangle (30,0.5);
\node at (27,0.25) {\scriptsize 60};
\draw[<->] (0,-0.35)--(24,-0.35) node[midway,below,font=\small] {$24$ φορές το $x$};
\draw[<->] (0,0.95)--(30,0.95) node[midway,above,font=\small] {$420$};
\end{tikzpicture}
\end{center}
Άρα $24x+60=420$, $24x=360$ και $x=15$.\par
\apantisi $15$ € το παιδί.\par
\elegxos $24\cdot15=360$ και $360+60=420$. Σωστό.

\paradeigma{Οι σελίδες του βιβλίου}
\tixero Η Ρήνη διάβασε τα $\frac{2}{5}$ ενός βιβλίου και άλλες $30$ σελίδες. Έχει διαβάσει συνολικά $150$ σελίδες.\par
\tithelo Πόσες σελίδες έχει το βιβλίο.\par
\possindeo Έστω $x$ οι σελίδες. Τότε $\frac{2}{5}x+30=150$. Άρα $\frac{2}{5}x=120$ και $x=120\cdot\frac{5}{2}=300$.\par
\apantisi $300$ σελίδες.\par
\elegxos Τα $\frac{2}{5}$ των $300$ είναι $120$ και $120+30=150$.

\paradeigma{Η ηλικία της γιαγιάς}
\tixero Η γιαγιά Μυρτώ είναι $70$ ετών, δέκα χρόνια περισσότερα από το πενταπλάσιο της ηλικίας ενός εγγονού της.\par
\tithelo Την ηλικία του εγγονού.\par
\possindeo Έστω $x$ η ηλικία του εγγονού. Τότε $5x+10=70$, $5x=60$ και $x=12$.\par
\apantisi $12$ ετών.\par
\elegxos $5\cdot12+10=70$.

\tmima{Ασκήσεις}

\zoni{Το κατάλαβα}{τα βήματα και η επαλήθευση}

\ex Λύσε τις εξισώσεις: α) $3x+4=19$ \quad β) $6x+8=62$ \quad γ) $5+12\alpha=77$ \quad δ) $9\alpha-15=111$.

\ex Λύσε τις εξισώσεις: α) $-3x-5=10$ \quad β) $2-5x=-13$ \quad γ) $12x+31=7$ \quad δ) $8+6x-2=7-7$.

\ex Βάλε στη σωστή σειρά τα βήματα επίλυσης της $3x-7=11$: α) Διαιρώ και τα δύο μέλη με το $3$. \quad β) Προσθέτω $7$ και στα δύο μέλη. \quad γ) Επαληθεύω: $3\cdot6-7=11$. \quad δ) Έχω $3x=18$. \quad ε) Βρίσκω $x=6$.

\ex Συμπλήρωσε τη λύση: $5x+10=45$, άρα $5x=45-\dots=\dots$ και $x=\dots:5=\dots$.

\ex Εξέτασε αν ο αριθμός είναι λύση της εξίσωσης: α) $4x-5=7$, $x=3$ \quad β) $3x+8=2$, $x=-2$ \quad γ) $2x-1=9$, $x=4$ \quad δ) $4x+3=5$, $x=\dfrac{1}{2}$.

\ex Αξιοποίησε την επιμεριστική ιδιότητα και λύσε: α) $2(x-1)=18$ \quad β) $3-(x-2)=7$ \quad γ) $4(x+1)-2(x-3)=20$ \quad δ) $x-2(x-4)=10$.

\ex Λύσε: α) $\dfrac{1}{2}x+3=7$ \quad β) $\dfrac{2}{3}x-1=5$ \quad γ) $0{,}5x+1{,}2=4$ \quad δ) $\dfrac{x}{8}-1=0$.

\zoni{Το χρησιμοποιώ}{ιστορίες που γίνονται εξισώσεις}

\ex Γράψε εξίσωση και λύσε: α) το διπλάσιο ενός αριθμού αυξημένο κατά $3$ κάνει $15$ \quad β) το τριπλάσιο ενός αριθμού ελαττωμένο κατά $7$ κάνει $20$ \quad γ) το $56$ ελαττωμένο κατά το τετραπλάσιο ενός αριθμού κάνει $8$ \quad δ) το μισό ενός αριθμού αυξημένο κατά $3$ κάνει $10$.

\ex \textbf{Το ποδήλατο.} Ένα ποδήλατο κοστίζει $180$ €. Ο Ηλίας έχει $45$ € και αποταμιεύει $15$ € την εβδομάδα. Σε πόσες εβδομάδες θα το αγοράσει;

\ex \textbf{Η εκδρομή.} Τα $25$ παιδιά μιας τάξης πληρώνουν από $x$ ευρώ και το σχολείο δίνει ακόμη $50$ €. Το σύνολο είναι $300$ €. Σχεδίασε το μοντέλο ράβδου, γράψε την εξίσωση και βρες το $x$.

\ex \textbf{Τα καφάσια.} Ο παππούς Γιάννης έχει δύο ίδια καφάσια με ελιές και ακόμη $3$ κιλά ελιές ξεχωριστά. Συνολικά έχει $31$ κιλά. Πόσα κιλά έχει κάθε καφάσι;

\ex \textbf{Οι ηλικίες.} Η θεία Μαρία είναι $38$ ετών, δύο χρόνια περισσότερα από το τριπλάσιο της ηλικίας του Νίκου. Πόσο ετών είναι ο Νίκος;

\zoni{Το σκέφτομαι}{ανάποδη πορεία, λάθη, δικά σου προβλήματα}

\ex Ο Στέλιος λύνει την $2x+3=11$ γράφοντας $x=11+3=14$ και μετά $x=14:2=7$. Βρες το λάθος, γράψε τη σωστή λύση και επαλήθευσε.

\ex Σκέφτηκα έναν αριθμό, τον πολλαπλασίασα επί $3$, αφαίρεσα $4$, διαίρεσα με το $2$ και βρήκα $10$. Βρες τον αριθμό πηγαίνοντας ανάποδα και μετά γράψε την εξίσωση.

\ex Γράψε ένα πρόβλημα της καθημερινής ζωής που λύνεται με την εξίσωση $2x+4=20$ και λύσε το.

\ex \textbf{Η μυστική συνταγή της γιαγιάς.} Η γιαγιά Μυρτώ φτιάχνει κουλουράκια με ζάχαρη, βούτυρο και αμύγδαλα. Το βούτυρο είναι διπλάσιο από τη ζάχαρη και τα αμύγδαλα είναι $10$ g λιγότερα από τη ζάχαρη. Όλα τα υλικά ζυγίζουν $390$ g. Πόσα γραμμάρια είναι το καθένα;

\zoni{Για το σπίτι}{φτιάξε και λύσε}

\ex Σκέψου ένα πρόβλημα από το σπίτι σου που λύνεται με την εξίσωση $4x+5=45$ και λύσε το.

\ex Λύσε την εξίσωση $2x-3(0{,}5x+1)-x=1$ και κάνε την επαλήθευση.

\paragrafos{Ανάλογα ποσά: όταν το ένα διπλασιάζεται, διπλασιάζεται και το άλλο}{Α.6.3, «Ανάλογα ποσά – Ιδιότητες αναλόγων ποσών»}

\begin{anagke}
Στο παζάρι ο παππούς πουλάει το μέλι του: $8$ ευρώ το κιλό. Ένας πελάτης παίρνει ένα κιλό και πληρώνει $8$ €. Άλλος παίρνει δύο κιλά και πληρώνει $16$ €. Η Φανή κρατάει το βιβλίο και παρατηρεί ότι «όσο μεγαλώνουν τα κιλά, μεγαλώνουν και τα λεφτά, και μάλιστα πάντα το ίδιο: τα διπλά κιλά, διπλά λεφτά, τα τριπλά κιλά, τριπλά λεφτά». Η Πολυξένη το γράφει σε έναν πίνακα και βρίσκει ότι η τιμή διά τα κιλά κάνει πάντα $8$.

Ο Στέλιος ενθουσιάζεται και προτείνει: «Αυτό ισχύει παντού! Ο Νίκος είναι $12$ ετών και $1{,}55$ μ. Αν διπλασιάσουμε την ηλικία, στα $24$ θα έχει διπλάσιο ύψος, $3{,}10$ μ.». Όλοι γελούν και ο Νίκος φωνάζει: «Δεν είμαι καμηλοπάρδαλη!». Εδώ η Πολυξένη πάει στο τετράδιο: «Δύο μεγέθη που αλλάζουν μαζί δεν είναι αυτόματα ανάλογα. Πρέπει να ελέγχουμε».

Η παρέα θα μάθει σήμερα να ξεχωρίζει τα ποσά που μεταβάλλονται με τον ίδιο τρόπο (τα ανάλογα) από εκείνα που απλώς μεγαλώνουν μαζί.
\end{anagke}

\begin{lexi}{ανάλογα, συντελεστής αναλογίας}{ανάλογα, συντελεστής αναλογίας}
Τα \lexiword{ανάλογα} ποσά είναι αυτά που αλλάζουν «κατά τον ίδιο λόγο» (ανά + λόγος). Ο \lexiword{συντελεστής αναλογίας} είναι ο σταθερός λόγος $\dfrac{y}{x}$ των δύο ποσών· «συν-τελεστής» είναι αυτός που «τελεί» μαζί με το $x$, δηλαδή που το πολλαπλασιάζει.
\end{lexi}

\begin{metafrasi}
Κιλά μέλι και κόστος, $8$ € το κιλό. & $y=8\cdot x$\\
Διπλά κιλά, διπλά λεφτά. & $x\to2x\ \Rightarrow\ y\to2y$\\
Ο λόγος των ποσών είναι σταθερός. & $\dfrac{y}{x}=\alpha$\\
Το $12\%$ ενός ποσού $x$. & $y=0{,}12\cdot x$ (ανάλογα με $\alpha=0{,}12$)\\
\end{metafrasi}

\tmima{Η θεωρία}

\begin{orismos}{ανάλογα ποσά}
Δύο ποσά λέγονται \textbf{ανάλογα} όταν μεταβάλλονται με τέτοιον τρόπο ώστε, όταν οι τιμές του ενός πολλαπλασιάζονται με έναν αριθμό, και οι αντίστοιχες τιμές του άλλου να πολλαπλασιάζονται με τον ίδιο αριθμό.
\end{orismos}

Ισοδύναμα, δύο ποσά $x$ και $y$ είναι ανάλογα όταν οι αντίστοιχες τιμές τους δίνουν πάντα το ίδιο πηλίκο:
$$\dfrac{y}{x}=\alpha\quad\text{ή}\quad y=\alpha\cdot x,$$
όπου ο σταθερός αριθμός $\alpha$ (με $\alpha\neq0$) λέγεται \textbf{συντελεστής αναλογίας}. Στο μέλι ο $\alpha=8$ και $y=8x$.

\tmimasmall{Πώς ελέγχω αν δύο ποσά είναι ανάλογα}

\begin{enumerate}[leftmargin=1.6em,itemsep=0.3ex,topsep=2pt]
\item \textbf{Με τον ορισμό:} αν πολλαπλασιάζω το $x$ επί έναν αριθμό, το $y$ πολλαπλασιάζεται επί τον ίδιο;
\item \textbf{Με τον λόγο:} βρίσκω τους λόγους $\frac{y}{x}$ για κάθε ζεύγος του πίνακα. Αν είναι όλοι ίσοι, τα ποσά είναι ανάλογα.
\item \textbf{Με τον τύπο:} αν από το πρόβλημα έχω $y=\alpha x$, είναι ανάλογα.
\end{enumerate}

Αν δεν ισχύει έστω και μία φορά, τα ποσά δεν είναι ανάλογα. Επιπλέον, αν τα ποσά είναι ανάλογα και για $x=0$ έχουμε $y=0$.

\tmimasmall{Το ποσοστό είναι αναλογία}

Όταν το $y$ είναι ένα ποσοστό του $x$, δηλαδή $y=\frac{\alpha}{100}x$, τα δύο ποσά είναι ανάλογα με συντελεστή αναλογίας $\frac{\alpha}{100}$ ή $\alpha\%$. Έτσι το $12\%$ του $x$ είναι $y=0{,}12x$.

\tmimasmall{Ιδιότητες των αναλόγων ποσών}

Αν τα $x,\,y$ είναι ανάλογα, τότε:
\begin{itemize}[leftmargin=1.2em,itemsep=0.3ex,topsep=2pt]
\item Αν $(x_1,\,y_1)$ και $(x_2,\,y_2)$ είναι ζεύγη τιμών, τότε $\dfrac{y_1}{x_1}=\dfrac{y_2}{x_2}$, άρα έχουμε αναλογία και μπορούμε να βρούμε έναν άγνωστο όρο.
\item Το άθροισμα των τιμών ενός ποσού αντιστοιχεί στο άθροισμα των τιμών του άλλου: $(x_1+x_2,\ y_1+y_2)$ είναι κι αυτό ζεύγος τιμών.
\end{itemize}

\begin{prosoxi}
\begin{itemize}[leftmargin=1.2em,itemsep=0.3ex,topsep=2pt]
\item Δύο ποσά που μεγαλώνουν μαζί δεν είναι απαραίτητα ανάλογα: η ηλικία και το ύψος μεγαλώνουν μαζί, αλλά ο λόγος τους δεν μένει σταθερός.
\item Η σχέση $y=x+2$ δεν είναι αναλογία: για $x=0$ δίνει $y=2\neq0$. Στην αναλογία το $0$ πηγαίνει στο $0$.
\item Για να είναι ανάλογα, πρέπει να ισχύει για \emph{όλα} τα ζεύγη του πίνακα, όχι μόνο για τα δύο πρώτα.
\item Το εμβαδόν τετραγώνου δεν είναι ανάλογο της πλευράς του: πλευρά $1$ δίνει $1$, πλευρά $2$ δίνει $4$ (όχι $2$).
\end{itemize}
\end{prosoxi}

\tmima{Λυμένα παραδείγματα}

\paradeigma{Το μέλι του παππού}
\tixero Το μέλι κοστίζει $8$ € το κιλό.\par
\tithelo Να συμπληρώσω τον πίνακα για $0{,}5$, $1$, $2$, $3$, $5$ κιλά και να βρω πόσα κιλά αγοράζονται με $20$ €.\par
\possindeo Τα ποσά είναι ανάλογα με $y=8x$. Για $x=0{,}5$: $4$. $x=1$: $8$. $x=2$: $16$. $x=3$: $24$. $x=5$: $40$. Για $y=20$: $x=20:8=2{,}5$.\par
\apantisi $4,\,8,\,16,\,24,\,40$ € και $2{,}5$ κιλά.\par
\elegxos $2{,}5\cdot8=20$. Και $\frac{y}{x}=8$ σε κάθε ζεύγος.

\paradeigma{Είναι ανάλογα;}
\tixero Δύο πίνακες. (Α) $x$: $1,\,2,\,3,\,4$ και $y$: $2,\,4,\,6,\,8$. (Β) $x$: $1,\,2,\,3$ και $y$: $3,\,5,\,7$.\par
\tithelo Να δω ποιος περιγράφει ανάλογα ποσά.\par
\possindeo Ο λόγος $\frac{y}{x}$. Πίνακας Α: $\frac{2}{1}=\frac{4}{2}=\frac{6}{3}=\frac{8}{4}=2$, ίδιος. Πίνακας Β: $\frac{3}{1}=3$, $\frac{5}{2}=2{,}5$, $\frac{7}{3}\approx2{,}33$, δεν είναι ίσοι.\par
\apantisi Ο Α είναι ανάλογα ($\alpha=2$), ο Β όχι.\par
\elegxos Στον Β όταν το $x$ διπλασιάζεται από $1$ σε $2$, το $y$ γίνεται $5$, όχι το διπλάσιο του $3$.

\paradeigma{Βρίσκω το συντελεστή}
\tixero Τα $x,\,y$ είναι ανάλογα: για $x=4$ είναι $y=6$ και για $x=10$ είναι $y=15$.\par
\tithelo Ο συντελεστής και το $y$ για $x=14$.\par
\possindeo $\alpha=\frac{6}{4}=1{,}5$ και $\frac{15}{10}=1{,}5$. Για $x=14$: $y=1{,}5\cdot14=21$.\par
\apantisi $\alpha=1{,}5$ και $y=21$.\par
\elegxos $4+10=14$ και $6+15=21$ (ιδιότητα του αθροίσματος).

\paradeigma{Το διάλυμα}
\tixero Ένα διάλυμα έχει περιεκτικότητα $12\%$ σε αλάτι.\par
\tithelo Πόσο αλάτι υπάρχει σε $250$ g διαλύματος και πόσο διάλυμα περιέχει $9$ g αλάτι.\par
\possindeo Αλάτι $y=0{,}12x$, όπου $x$ το διάλυμα. Για $x=250$: $y=0{,}12\cdot250=30$ g. Για $y=9$: $x=9:0{,}12=75$ g.\par
\apantisi $30$ g αλάτι και $75$ g διαλύματος.\par
\elegxos $\frac{30}{250}=0{,}12$ και $\frac{9}{75}=0{,}12$.

\paradeigma{Ύψος και ηλικία}
\tixero Ο Νίκος στα $8$ ήταν $1{,}28$ μ. και στα $12$ είναι $1{,}55$ μ.\par
\tithelo Είναι ύψος και ηλικία ανάλογα;\par
\possindeo $\frac{1{,}28}{8}=0{,}16$ και $\frac{1{,}55}{12}\approx0{,}129$. Οι λόγοι δεν είναι ίσοι. Επίσης, η ηλικία αυξήθηκε $1{,}5$ φορές αλλά το ύψος μόνο $1{,}2$ φορές.\par
\apantisi Όχι, δεν είναι ανάλογα.\par
\elegxos Αν ήταν ανάλογα, στα $24$ ο Νίκος θα ήταν $3{,}1$ μ. Αδύνατο.

\tmima{Ασκήσεις}

\zoni{Το κατάλαβα}{πίνακες, συντελεστής, έλεγχος}

\ex Τα $x,\,y$ είναι ανάλογα με συντελεστή $\alpha=3$. Συμπλήρωσε: $x=0,\ 1,\ 2,\ 5,\ 0{,}4$ \ και βρες το $x$ όταν $y=9$ και όταν $y=4{,}5$.

\ex Τα $x,\,y$ είναι ανάλογα με $\alpha=\frac{2}{3}$. Βρες το $y$ για $x=0,\ 3,\ 6,\ 1$ και το $x$ για $y=10$.

\ex Εξέτασε αν τα $x,\,y$ είναι ανάλογα: α) $x$: $2,4,6,8$, $y$: $5,10,15,20$ \quad β) $x$: $1,2,3$, $y$: $2,3,4$ \quad γ) $x$: $3,6,9$, $y$: $1,2,3$ \quad δ) $x$: $1,2,4$, $y$: $5,10,21$.

\ex Βρες τον τύπο και πες αν τα ποσά είναι ανάλογα: α) το κόστος $y$ για $x$ τεμάχια των $1{,}5$ € \quad β) η περίμετρος $y$ τετραγώνου με πλευρά $x$ \quad γ) το εμβαδόν $y$ τετραγώνου με πλευρά $x$.

\ex Τα $x,\,y$ είναι ανάλογα με $x$: $5,\,8,\,12$ και $y$: $15,\,24,\,36$. Βρες τον συντελεστή.

\ex Σωστό ή λάθος; α) Αν διπλασιάσω το $x$ και το $y$ διπλασιάζεται, τα ποσά είναι ανάλογα. \quad β) Σε ανάλογα ποσά το $y$ είναι $0$ όταν το $x$ είναι $0$. \quad γ) Η $y=x+2$ είναι αναλογία. \quad δ) Το $y=0{,}2x$ είναι αναλογία με συντελεστή $20\%$. \quad ε) Το εμβαδόν ενός τετραγώνου είναι ανάλογο της πλευράς του.

\zoni{Το χρησιμοποιώ}{μέλι, βενζίνη, ύφασμα, ταχύτητα}

\ex \textbf{Το μέλι.} Συμπλήρωσε τα κόστη για $0{,}5,\ 1,\ 1{,}5,\ 2{,}5,\ 4$ κιλά μέλι των $8$ €. Πόσα κιλά αγοράζονται με $52$ €;

\ex \textbf{Η βενζίνη.} Ο θείος Λευτέρης αγοράζει βενζίνη προς $1{,}85$ € το λίτρο. Πόσο πληρώνει για $20$ και για $35$ λίτρα; Πόσα λίτρα παίρνει με $50$ €;

\ex \textbf{Το ύφασμα.} Η Αγγελική αγοράζει ύφασμα προς $4{,}8$ € το μέτρο. Πόσο κοστίζουν $2{,}5$ μ. και $0{,}75$ μ.; Πόσα μέτρα παίρνει με $33{,}6$ €;

\ex \textbf{Το διάλυμα.} Ένα διάλυμα έχει περιεκτικότητα $12\%$ σε αλάτι. Πόσο αλάτι υπάρχει σε $250$ g; Πόσο διάλυμα περιέχει $9$ g αλάτι;

\ex \textbf{Ο ΦΠΑ.} Ο ΦΠΑ είναι $24\%$ της αξίας. Βρες τον ΦΠΑ σε αξίες $50$ € και $200$ €. Για ποια αξία ο ΦΠΑ είναι $30$ €;

\ex \textbf{Ο ποδηλάτης.} Ο Νίκος ποδηλατεί με σταθερή ταχύτητα $18$ km την ώρα. Πόση απόσταση κάνει σε $0{,}5$, $2$ και $3{,}5$ ώρες; Σε πόσο χρόνο κάνει $27$ km;

\ex \textbf{Η μπογιά.} Ένα λίτρο μπογιάς καλύπτει $8$ m$^2$. Πόση επιφάνεια καλύπτουν $2{,}5$ λίτρα; Πόσα λίτρα χρειάζονται για $60$ m$^2$;

\zoni{Το σκέφτομαι}{ανάλογα ή όχι}

\ex Ο Στέλιος λέει ότι αφού ο Νίκος στα $12$ είναι $1{,}55$ μ., στα $24$ θα είναι $3{,}10$ μ. Γιατί σφάλλει;

\ex Ποια από τα παρακάτω ζεύγη ποσών είναι ανάλογα; α) τα κιλά και η τιμή \quad β) η ηλικία παιδιού και το βάρος του \quad γ) ο χρόνος και η απόσταση με σταθερή ταχύτητα \quad δ) η πλευρά και το εμβαδόν τετραγώνου \quad ε) η πλευρά και η περίμετρος τετραγώνου.

\ex Αν $y=2x$ και $z=5x$, δείξε ότι τα $y$ και $z$ είναι ανάλογα και βρες τον συντελεστή.

\ex Για το $y=3x$ τα $(2,\,6)$ και $(5,\,15)$ είναι ζεύγη τιμών. Έλεγξε ότι και το $(7,\,21)$ (το άθροισμα) είναι ζεύγος τιμών.

\zoni{Για το σπίτι}{ανάλογα στην κουζίνα και στο σούπερ μάρκετ}

\ex Βρες από τη συσκευασία ενός προϊόντος την τιμή ανά κιλό και φτιάξε πίνακα με τις τιμές για διάφορες ποσότητες.

\ex Βρες δύο ποσά της καθημερινής ζωής που είναι ανάλογα και δύο που μεγαλώνουν μαζί χωρίς να είναι ανάλογα.

\paragrafos{Γραφική παράσταση της αναλογίας: μια ημιευθεία από την αρχή}{Α.6.4, «Γραφική παράσταση σχέσης αναλογίας»}

\begin{anagke}
Ο Νίκος και ο Βασίλης κάνουν ποδήλατο στον παραλιακό δρόμο. Ο Νίκος κάνει $18$ χιλιόμετρα την ώρα και ο Βασίλης $12$. Η Πολυξένη γράφει τις αποστάσεις τους σε πίνακες: ώρα $0,\ 1,\ 2,\ 3$ και χιλιόμετρα $0,\ 18,\ 36,\ 54$ για τον Νίκο, και $0,\ 12,\ 24,\ 36$ για τον Βασίλη. Ο Ηλίας κοιτάζει τους αριθμούς και λέει: «Βλέπω ότι ο Νίκος είναι γρηγορότερος, αλλά δεν μπορώ να δω \emph{πόσο}».

Η Αγγελική παίρνει ένα τετραγωνισμένο φύλλο και, όπως στη ναυμαχία, σημειώνει τα ζευγάρια ώρα--χιλιόμετρα ως σημεία. Όλα τα σημεία του Νίκου πέφτουν σε μια ευθεία γραμμή που ξεκινά από την αρχή. Το ίδιο και του Βασίλη, αλλά η δική του γραμμή είναι πιο «ξαπλωτή». «Όσο πιο γρήγορα, τόσο πιο απότομη η γραμμή», παρατηρεί η Ρήνη. Ο Στέλιος ρωτάει αν αυτό ισχύει για κάθε αναλογία, και η Φανή δοκιμάζει με το μέλι του παππού και μια δική της σχέση, που δεν πέφτει σε ευθεία.

Ένα σχήμα λέει με μια ματιά αυτό που ένας πίνακας λέει με πολλά νούμερα. Σήμερα βρίσκουμε τι είδους σχήμα έχει η αναλογία και πώς το χρησιμοποιούμε για να διαβάζουμε τιμές.
\end{anagke}

\begin{lexi}{γραφική παράσταση, ημιευθεία}{γραφική παράσταση}
Η \lexiword{γραφική παράσταση} είναι η «εικόνα που γράφεται» με σημεία και γραμμές. Η \lexiword{ημιευθεία} είναι η μισή ευθεία: ξεκινά από ένα σημείο και συνεχίζει στο άπειρο προς μία κατεύθυνση. Στις αναλογίες το σημείο εκκίνησης είναι η αρχή $O$ των αξόνων.
\end{lexi}

\begin{metafrasi}
Ένα ζεύγος τιμών $(x,\,y)$ ενός πίνακα. & ένα σημείο του επιπέδου\\
Ανάλογα ποσά με συντελεστή $\alpha$. & σημεία πάνω στην ημιευθεία $y=\alpha x$ από το $O$\\
Μεγαλύτερος συντελεστής. & πιο απότομη ημιευθεία\\
Για $x=1$ η τεταγμένη. & είναι ο συντελεστής $\alpha$\\
\end{metafrasi}

\tmima{Η θεωρία}

Αν δύο ποσά είναι ανάλογα με συντελεστή $\alpha$, δηλαδή $y=\alpha x$, ένα σημείο του επιπέδου αντιστοιχεί σε κάθε ζεύγος τιμών $(x,\,y)$. Για παράδειγμα με $\alpha=1{,}5$ έχουμε τον πίνακα
\begin{center}
\renewcommand{\arraystretch}{1.3}
\begin{tabular}{c|cccccc}
$x$ & $0$ & $1$ & $2$ & $4$ & $5$ & $6$\\\hline
$y=1{,}5x$ & $0$ & $1{,}5$ & $3$ & $6$ & $7{,}5$ & $9$
\end{tabular}
\end{center}
και τα σημεία $(0,\,0)$, $(1,\,1{,}5)$, $(2,\,3)$, $(4,\,6)$, $(5,\,7{,}5)$, $(6,\,9)$.

\begin{center}
\begin{tikzpicture}[x=0.9cm,y=0.5cm]
\draw[gray!35,very thin] (0,0) grid (6.4,9.4);
\draw[thick,-{Stealth}] (0,0)--(6.8,0) node[below] {$x$};
\draw[thick,-{Stealth}] (0,0)--(0,10) node[left] {$y$};
\foreach \i in {1,...,6}{\node[below,font=\scriptsize] at (\i,0) {\i};}
\foreach \j in {1,3,5,7,9}{\node[left,font=\scriptsize] at (0,\j) {\j};}
\draw[gblue,thick] (0,0)--(6,9);
\foreach \x/\y in {0/0,1/1.5,2/3,4/6,5/7.5,6/9}{\fill[gblue] (\x,\y) circle (2.3pt);}
\node[gblue,right] at (6,9) {$y=1{,}5x$};
\end{tikzpicture}
\end{center}

Παρατηρούμε ότι τα σημεία βρίσκονται πάνω στην ίδια ημιευθεία με αρχή το $O(0,\,0)$.

\begin{idiotita}{γραφική παράσταση της αναλογίας}
Τα σημεία που αντιστοιχούν στα ζεύγη τιμών $(x,\,y)$ δύο \textbf{ανάλογων} ποσών βρίσκονται πάνω σε μία \textbf{ημιευθεία που ξεκινά από την αρχή $O(0,\,0)$ των ημιαξόνων}. Και αντίστροφα, αν τα σημεία ενός πίνακα βρίσκονται σε ημιευθεία από το $O$, τα ποσά είναι ανάλογα.
\end{idiotita}

Ο συντελεστής $\alpha$ βρίσκεται από τη γραφική παράσταση: είναι η τεταγμένη του σημείου που έχει τετμημένη $1$, ή γενικά το πηλίκο $\frac{y}{x}$ οποιουδήποτε σημείου της. Όσο μεγαλύτερος ο $\alpha$, τόσο πιο «απότομη» η ημιευθεία. Από τη γραφική παράσταση διαβάζουμε και τιμές που δεν υπάρχουν στον πίνακα: ανεβαίνουμε από ένα $x$ ως την ημιευθεία και διαβάζουμε το $y$.

\tmimasmall{Πώς διαπιστώνω γραφικά ότι τα ποσά \emph{δεν} είναι ανάλογα}

Αν τα σημεία δεν είναι σε ευθεία (έχουν «τσακίσματα» ή κάμπτονται), ή αν είναι σε ευθεία που δεν περνά από το $O$, τα ποσά δεν είναι ανάλογα. Π.χ. η $y=x+2$ έχει σημεία $(0,\,2)$, $(1,\,3)$, $(2,\,4)$ που βρίσκονται σε ευθεία που κόβει τον $Oy$ στο $2$.

\begin{prosoxi}
\begin{itemize}[leftmargin=1.2em,itemsep=0.3ex,topsep=2pt]
\item Δεν αρκεί να είναι τα σημεία σε ευθεία: πρέπει η ευθεία να περνά από την αρχή $O$. Η $y=x+2$ είναι ευθεία χωρίς να είναι αναλογία.
\item Δεν ενώνουμε τα σημεία «όπως μας βολεύει»: ενώνουμε με ημιευθεία μόνο αν ξέρουμε ότι υπάρχουν και ενδιάμεσες τιμές (π.χ. μπορεί να πάρω $2{,}5$ κιλά μέλι).
\item Προσοχή στη μονάδα στους άξονες: κάθε διαίρεση πρέπει να είναι ίση παντού, αλλά δεν χρειάζεται να είναι ίδια στον $x$ και στον $y$.
\item Σε ανάλογα ποσά, η μεγάλη τιμή του $y$ δεν σημαίνει αυτόματα μεγάλο συντελεστή: μετράει ο λόγος $\frac{y}{x}$.
\end{itemize}
\end{prosoxi}

\tmima{Λυμένα παραδείγματα}

\paradeigma{Σχεδιάζω την αναλογία}
\tixero Τα $x,\,y$ είναι ανάλογα με $y=1{,}5x$.\par
\tithelo Να τα παραστήσω γραφικά.\par
\possindeo Κάνω πίνακα ($x=0,\,2,\,4,\,6$ και $y=0,\,3,\,6,\,9$), τοποθετώ τα σημεία και φέρνω ημιευθεία από το $O$ που τα περιέχει.\par
\apantisi Η ημιευθεία από το $O$ που περνά από το $(2,\,3)$ (σχήμα θεωρίας).\par
\elegxos Το $(4,\,6)$ είναι πάνω της και $\frac{6}{4}=1{,}5$.

\paradeigma{Ποιος πίνακας είναι αναλογία;}
\tixero Πίνακες με $x=0,\,1,\,2,\,3$: (Α) $y=0,\,2,\,3,\,5$. (Β) $y=2,\,4,\,6,\,8$. (Γ) $y=0,\,2,\,4,\,6$. (Δ) $y=0,\,1,\,4,\,9$.\par
\tithelo Να διαπιστώσω ποιοι δίνουν ανάλογα ποσά.\par
\possindeo Σε κάθε πίνακα κοιτάζω αν τα σημεία είναι σε ημιευθεία από το $O$. Α: οι λόγοι $\frac{2}{1}=2$, $\frac{3}{2}=1{,}5$: όχι. Β: έχει $y=2$ όταν $x=0$, άρα δεν ξεκινά από το $O$: όχι. Γ: $\frac{y}{x}=2$ παντού: ναι. Δ: $\frac{1}{1}=1$, $\frac{4}{2}=2$: όχι, είναι καμπύλη.\par
\apantisi Μόνο ο Γ ($y=2x$).\par
\elegxos Το Β είναι ευθεία $y=x\cdot2+2$ που περνά από το $(0,\,2)$.

\paradeigma{Οι δύο ποδηλάτες}
\tixero Ο Νίκος κάνει $18$ km/h και ο Βασίλης $12$ km/h.\par
\tithelo Γραφική παράσταση για κάθε ποδηλάτη, απόσταση στις $2$ ώρες, και πότε ο Νίκος κάνει $54$ km.\par
\possindeo Οι αποστάσεις είναι $y=18x$ και $y=12x$ ($x$ ώρες, $y$ χλμ.). Τα σημεία: Νίκος $(0,\,0)$, $(1,\,18)$, $(2,\,36)$, $(3,\,54)$. Βασίλης $(0,\,0)$, $(1,\,12)$, $(2,\,24)$, $(3,\,36)$. Στις $2$ ώρες: $36$ και $24$ km. Για τα $54$ km ο Νίκος: $54:18=3$ ώρες.\par
\apantisi Η ημιευθεία του Νίκου είναι πιο απότομη. Στις $2$ ώρες: $36$ και $24$ km. $54$ km σε $3$ ώρες. Μετά από $3$ ώρες η διαφορά τους είναι $54-36=18$ km.\par
\elegxos Ο συντελεστής $18>12$ κάνει την ημιευθεία του Νίκου πιο απότομη.

\begin{center}
\begin{tikzpicture}[x=1.5cm,y=0.06cm]
\draw[gray!35,very thin] (0,0) grid[xstep=0.5,ystep=10] (3.2,60);
\draw[thick,-{Stealth}] (0,0)--(3.5,0) node[below] {ώρες};
\draw[thick,-{Stealth}] (0,0)--(0,66) node[left] {km};
\foreach \i in {1,2,3}{\node[below,font=\scriptsize] at (\i,0) {\i};}
\foreach \j in {20,40,60}{\node[left,font=\scriptsize] at (0,\j) {\j};}
\draw[gblue,thick] (0,0)--(3,54) node[right] {Νίκος};
\draw[gblue!60!black,thick,dashed] (0,0)--(3,36) node[right] {Βασίλης};
\end{tikzpicture}
\end{center}

\paradeigma{Βρίσκω τον συντελεστή από το γράφημα}
\tixero Δύο ανάλογα ποσά έχουν γραφική παράσταση ημιευθεία από το $O$ που περνά από το σημείο $(4,\,10)$.\par
\tithelo Ο συντελεστής και το $y$ όταν $x=7$.\par
\possindeo $\alpha=\frac{10}{4}=2{,}5$. Τότε $y=2{,}5\cdot7=17{,}5$.\par
\apantisi $\alpha=2{,}5$ και $y=17{,}5$.\par
\elegxos Το σημείο $(1,\,2{,}5)$ είναι κι αυτό στην ημιευθεία.

\tmima{Ασκήσεις}

\zoni{Το κατάλαβα}{σημεία, ημιευθείες, πίνακες}

\ex Φτιάξε πίνακα για την $y=2x$ με $x=0,\,1,\,2,\,3,\,4,\,5$ και σχεδίασε τη γραφική παράσταση.

\ex Σχεδίασε στο ίδιο σύστημα αξόνων τις $y=x$, $y=2x$ και $y=0{,}5x$ για $x$ από $0$ ως $6$. Ποια είναι πιο απότομη και γιατί;

\ex Ποιοι από τους πίνακες παριστάνουν ανάλογα ποσά; α) $x$: $0,2,4,6$, $y$: $0,3,6,9$ \quad β) $x$: $0,1,2,3$, $y$: $1,2,3,4$ \quad γ) $x$: $0,1,2,3$, $y$: $0,1,3,6$ \quad δ) $x$: $0,2,4,5$, $y$: $0,5,10,12{,}5$.

\ex Τα σημεία $(2,\,3)$, $(4,\,6)$, $(6,\,9)$ και $(8,\,12)$ βρίσκονται σε ημιευθεία από το $O$; Βρες τον συντελεστή. Και τα $(2,\,3)$, $(4,\,6)$, $(5,\,8)$;

\ex Η ημιευθεία από το $O$ περνά από το $(5,\,12)$. Βρες τον συντελεστή αναλογίας, το $y$ για $x=2$ και το $x$ για $y=30$.

\ex Για την ευθεία $y=x+2$ φτιάξε πίνακα για $x=0,\,1,\,2,\,3$ και σχεδίασε. Γιατί δεν είναι αναλογία;

\zoni{Το χρησιμοποιώ}{μέλι, ποδήλατα, βενζίνη}

\ex \textbf{Το μέλι.} Σχεδίασε τη γραφική παράσταση της τιμής για μέλι $8$ €/κιλό. Διάβασε από αυτή το κόστος για $2{,}5$ κιλά και τα κιλά που παίρνεις με $36$ €.

\ex \textbf{Οι ποδηλάτες.} Ο Νίκος κάνει $18$ km/h και ο Βασίλης $12$ km/h. Σχεδίασε τις αποστάσεις για $0$ ως $4$ ώρες. Πόση είναι η διαφορά τους μετά από $3$ ώρες; Σε πόση ώρα ο Βασίλης κάνει $30$ km;

\ex \textbf{Η βενζίνη.} Σχεδίασε την τιμή της βενζίνης προς $1{,}85$ €/λίτρο για $0$ ως $40$ λίτρα (χρησιμοποίησε κατάλληλη μονάδα στον κάθε άξονα). Πόσο κοστίζουν $12$ λίτρα περίπου;

\ex \textbf{Το σχοινί.} Ένα σχοινί βάρους $0{,}4$ kg ανά μέτρο. Σχεδίασε το βάρος για $0$ ως $5$ μ. και πες πόσο ζυγίζουν $3{,}5$ μ.

\zoni{Το σκέφτομαι}{γιατί ημιευθεία από το $O$}

\ex Γιατί η γραφική παράσταση μιας αναλογίας περνά πάντα από το $O$;

\ex Η γραφική παράσταση ενός πίνακα είναι ευθεία που δεν περνά από το $O$. Μπορούν τα ποσά να είναι ανάλογα; Εξήγησε.

\ex Δύο ανάλογα ποσά έχουν σημεία $(3,\,4)$ και $(6,\,8)$. Ποιο είναι το $y$ για $x=9$ και ποιο το $x$ για $y=2$;

\ex Ο Στέλιος ένωσε τα σημεία της θερμοκρασίας του πυρετού της Φανής με ημιευθεία από το $O$. Γιατί είναι λάθος;

\zoni{Για το σπίτι}{μέτρησε και σχεδίασε}

\ex Μέτρησε το μήκος $1,\,2,\,3,\,4$ βημάτων σου και σχεδίασε το γράφημα βήματα--μέτρα. Είναι αναλογία;

\ex Πάρε τα κιλά και την τιμή τριών προϊόντων από το σούπερ μάρκετ και σχεδίασε την τιμή σε σχέση με το βάρος για ένα από αυτά.

\paragrafos{Προβλήματα αναλογιών: από το ένα στο πολλά}{Α.6.5, «Προβλήματα αναλογιών»}

\begin{anagke}
Η Α3 ετοιμάζει ένα μεγάλο φαγητό για το παζάρι. Η συνταγή της γιαγιάς Μυρτούς γράφει «για $4$ άτομα: $300$ γραμμάρια ρύζι». Θα έρθουν $10$ άτομα. «Πόσο ρύζι θέλουμε;» ρωτάει η Ρήνη. Ο Βασίλης προτείνει: «Για τέσσερα $300$, για άλλα τέσσερα $300$, και για τα υπόλοιπα δύο κάτι λιγότερο. Να βάλουμε $800$». Η Πολυξένη τον σταματάει: «Μην μαντεύεις. Πόσο ρύζι τρώει ένα άτομο;». Η Φανή λέει ότι ένα άτομο θέλει $300:4=75$ γραμμάρια, άρα τα $10$ άτομα $750$. «Αυτό δεν είναι ο κανόνας των τριών που έλεγαν στο δημοτικό;» ρωτάει ο Στέλιος. Η Αγγελική παρατηρεί κάτι ακόμη: «Δεν είναι όλα τα προβλήματα έτσι. Αν η Ζωή είναι $5$ χρονών όταν η αδελφή της είναι $2$, δεν σημαίνει ότι όταν η Ζωή γίνει $10$ η αδελφή της θα γίνει $4$». Η παρέα καταλαβαίνει ότι πριν από οποιαδήποτε μέθοδο πρέπει να ελέγξει αν τα ποσά είναι ανάλογα.

Σε αυτή την παράγραφο η παρέα μαζεύει όλες τις τεχνικές που ξέρει: πίνακα, αναλογία, αναγωγή στη μονάδα, γράφημα, και χρησιμοποιεί και τον έλεγχο της αναλογίας ώστε να μην εφαρμόζει κανόνες όπου δεν ισχύουν.
\end{anagke}

\begin{lexi}{αναγωγή στη μονάδα}{αναγωγή στη μονάδα}
\lexiword{Αναγωγή} σημαίνει «φέρνω πίσω»: φέρνω το πρόβλημα πίσω στη μονάδα, δηλαδή βρίσκω πρώτα πόσο αντιστοιχεί στο ένα, και μετά πολλαπλασιάζω. Το \lexiword{ανά} σημαίνει «για κάθε»: «ανά άτομο» είναι «για κάθε άτομο».
\end{lexi}

\begin{metafrasi}
Για $4$ άτομα $300$ g, για $10$ άτομα $x$ g. & $\dfrac{4}{10}=\dfrac{300}{x}$\\
Πόσο για ένα άτομο; & $300:4=75$ g\\
Για $10$ άτομα. & $10\cdot75=750$ g\\
Ο πίνακας των δύο ποσών. & $\begin{array}{c|cc}\text{άτομα}&4&10\\\hline\text{ρύζι}&300&x\end{array}$\\
\end{metafrasi}

\tmima{Η θεωρία}

\tmimasmall{Πρώτα ελέγχω αν είναι ανάλογα}

Για να δω αν δύο ποσά είναι ανάλογα χρησιμοποιώ έναν από τους τρεις τρόπους της προηγούμενης παραγράφου: τον ορισμό (πολλαπλασιάζω το ένα, πολλαπλασιάζεται και το άλλο), τον τύπο $y=\alpha x$ ή τον σταθερό λόγο $\frac{y}{x}$. Στα προβλήματα της ζωής σκέφτομαι: \emph{αν διπλασιάσω το ένα, διπλασιάζεται και το άλλο;} Αν η απάντηση είναι ναι, τα ποσά είναι ανάλογα.

\tmimasmall{Τρεις τρόποι επίλυσης}

\begin{enumerate}[leftmargin=1.6em,itemsep=0.3ex,topsep=2pt]
\item \textbf{Αναγωγή στη μονάδα.} Βρίσκω πρώτα την τιμή για το $1$. $300:4=75$ g το άτομο, και $10\cdot75=750$ g.
\item \textbf{Αναλογία (τέταρτος ανάλογος).} Οργανώνω πίνακα και γράφω ισότητα λόγων: $\dfrac{4}{10}=\dfrac{300}{x}$, άρα $4x=3\,000$, $x=750$.
\item \textbf{Συντελεστής αναλογίας.} Βρίσκω τον συντελεστή $\alpha=\frac{300}{4}=75$ και γράφω $y=75x$. Για $x=10$: $y=750$. Ή από τη γραφική παράσταση.
\end{enumerate}

Και οι τρεις δίνουν το ίδιο. Διαλέγω αυτόν που μου φαίνεται πιο εύκολος και, αν μπορώ, ελέγχω με έναν δεύτερο.

\tmimasmall{Διαιρώ ένα ποσό ανάλογα με δοσμένους αριθμούς}

Όταν μοιράζω ένα ποσό ανάλογα με τους αριθμούς $3$ και $5$, σημαίνει ότι το κάθε μέρος είναι $3$ και $5$ «μονάδες» αντίστοιχα. Αν το σύνολο είναι $120$, οι μονάδες είναι $3+5=8$, η μία μονάδα $120:8=15$ και τα μέρη $3\cdot15=45$ και $5\cdot15=75$.

\begin{prosoxi}
\begin{itemize}[leftmargin=1.2em,itemsep=0.3ex,topsep=2pt]
\item Δεν είναι όλα ανάλογα. Αν η Ζωή είναι $5$ χρονών και η αδελφή της $2$, οι ηλικίες διαφέρουν πάντα κατά $3$: όχι ανάλογα.
\item Ο ίδιος αριθμός ($4$) δεν είναι πάντα το $x$: διάβασε προσεκτικά ποιο ποσό ποιον αριθμό δίνει.
\item Στον πίνακα βάζω τα ομοειδή μεγέθη στην ίδια γραμμή: άτομα με άτομα, γραμμάρια με γραμμάρια.
\item Αν βγει αποτέλεσμα που δεν έχει νόημα (π.χ. $3{,}5$ άνθρωποι), το κρατάω ως τιμή αλλά αναθεωρώ την απάντηση για τη ζωή.
\end{itemize}
\end{prosoxi}

\tmima{Λυμένα παραδείγματα}

\paradeigma{Το ρύζι της συνταγής}
\tixero Για $4$ άτομα η συνταγή θέλει $300$ g ρύζι και $2$ κρεμμύδια.\par
\tithelo Η ποσότητα για $10$ άτομα.\par
\possindeo Αναγωγή στη μονάδα: ρύζι ανά άτομο $300:4=75$ g, άρα $750$ g. Κρεμμύδια ανά άτομο $2:4=0{,}5$, άρα $5$.\par
\apantisi $750$ g ρύζι και $5$ κρεμμύδια.\par
\elegxos Με αναλογία: $\frac{4}{10}=\frac{300}{x}$ δίνει $x=750$.

\paradeigma{Το ύφασμα}
\tixero Τα $5$ μ. ύφασμα κοστίζουν $36$ €.\par
\tithelo Πόσο κοστίζουν $8$ μ.\par
\possindeo Ένα μέτρο κοστίζει $36:5=7{,}2$ €. Τα $8$ μ. κοστίζουν $8\cdot7{,}2=57{,}6$ €.\par
\apantisi $57{,}60$ €.\par
\elegxos Με αναλογία $\frac{5}{8}=\frac{36}{x}$, $5x=288$, $x=57{,}6$.

\paradeigma{Το ταξίδι}
\tixero Ένα αυτοκίνητο με σταθερή ταχύτητα κάνει $120$ km σε $1{,}5$ ώρες.\par
\tithelo Σε πόση ώρα θα κάνει $280$ km.\par
\possindeo Η απόσταση και ο χρόνος είναι ανάλογα. Ταχύτητα $120:1{,}5=80$ km/h. Χρόνος $280:80=3{,}5$ ώρες.\par
\apantisi $3{,}5$ ώρες.\par
\elegxos $80\cdot3{,}5=280$.

\paradeigma{Δεν είναι ανάλογα}
\tixero Η Ζωή είναι $5$ ετών όταν η αδελφή της είναι $2$.\par
\tithelo Πόσο θα είναι η αδελφή όταν η Ζωή είναι $15$;\par
\possindeo Αν ήταν ανάλογα θα είχαμε $6$. Αλλά οι ηλικίες διαφέρουν πάντα κατά $3$ χρόνια: $15-3=12$.\par
\apantisi $12$ ετών.\par
\elegxos Οι λόγοι: $\frac{2}{5}=0{,}4$ και $\frac{12}{15}=0{,}8$. Δεν είναι ίσοι.

\paradeigma{Ανάλογη μοιρασιά}
\tixero Ο Νίκος δούλεψε $3$ ώρες και ο Βασίλης $5$ ώρες στο παζάρι και μοιράζονται $120$ € ανάλογα με τις ώρες.\par
\tithelo Πόσα παίρνει ο καθένας.\par
\possindeo Οι ώρες είναι $3+5=8$. Η κάθε ώρα αντιστοιχεί σε $120:8=15$ €. Νίκος $3\cdot15=45$ €, Βασίλης $5\cdot15=75$ €.\par
\apantisi $45$ € και $75$ €.\par
\elegxos $45+75=120$ και $\frac{45}{3}=\frac{75}{5}=15$.

\tmima{Ασκήσεις}

\zoni{Το κατάλαβα}{αναγωγή στη μονάδα και αναλογία}

\ex Τα $3$ μπουκάλια νερό κοστίζουν $2{,}4$ €. Πόσο κοστίζουν τα $8$;

\ex Τα $5$ κιλά πορτοκάλια κοστίζουν $6{,}5$ €. Πόσο κοστίζουν $8$ κιλά; Πόσα κιλά παίρνεις με $20$ €;

\ex Τα $7$ τεμάχια ενός προϊόντος κοστίζουν $31{,}5$ €. Πόσο κοστίζουν τα $12$;

\ex Συμπλήρωσε τον πίνακα αναλόγων ποσών: $x$: $3,\ 5,\ 8,\ 12$ και $y$: $9,\ ?,\ ?,\ ?$.

\ex Ποια ποσά είναι ανάλογα; α) ηλικία και ύψος \quad β) κιλά και κόστος \quad γ) πλευρά και περίμετρος ισόπλευρου τριγώνου \quad δ) ώρες εργασίας και αμοιβή με σταθερό μεροκάματο \quad ε) πλευρά και εμβαδόν τετραγώνου \quad στ) ηλικία πατέρα και ηλικία γιου.

\zoni{Το χρησιμοποιώ}{κουζίνα, ταξίδια, δουλειές}

\ex \textbf{Το ρύζι.} Για $4$ άτομα η συνταγή θέλει $300$ g ρύζι και $2$ κρεμμύδια. Βρες τις ποσότητες για $10$ άτομα και για $6$ άτομα.

\ex \textbf{Οι πίτες.} Η γιαγιά Μυρτώ φτιάχνει $6$ πίτες από $1{,}2$ κιλά αλεύρι. Πόσο αλεύρι θέλει για $15$ πίτες;

\ex \textbf{Το ταξίδι.} Ένα αυτοκίνητο κάνει $120$ km σε $1{,}5$ ώρες. Πόσα km θα κάνει σε $3{,}5$ ώρες με την ίδια ταχύτητα;

\ex \textbf{Η κατανάλωση.} Ένα αυτοκίνητο καταναλώνει $8$ λίτρα ανά $100$ km. Πόση βενζίνη θέλει για $340$ km;

\ex \textbf{Οι φωτοτυπίες.} Οι $35$ φωτοτυπίες κοστίζουν $2{,}45$ €. Πόσο κοστίζουν οι $120$;

\ex \textbf{Οι ελιές.} Τα $15$ ελαιόδεντρα δίνουν $240$ κιλά ελιές. Πόσα κιλά δίνουν $40$ δέντρα;

\ex \textbf{Το νόμισμα.} Με ισοτιμία $1$ € $=1{,}1$ δολάρια, πόσα δολάρια είναι τα $50$ €; Πόσα ευρώ είναι τα $132$ δολάρια;

\ex \textbf{Η μοιρασιά.} Ο Νίκος δούλεψε $3$ ώρες και ο Βασίλης $5$ ώρες στο παζάρι, και μοιράζονται $120$ € ανάλογα με τις ώρες. Πόσα παίρνει ο καθένας;

\zoni{Το σκέφτομαι}{όταν ο κανόνας δεν ισχύει}

\ex Η Φανή είναι $5$ ετών όταν η αδελφή της είναι $2$. Όταν η Φανή είναι $15$, πόσο είναι η αδελφή της; Γιατί δεν είναι ανάλογα;

\ex Ο Ηλίας τρέχει $5$ km σε $25$ λεπτά. Θα τρέξει $10$ km σε $50$ λεπτά; Συζήτησε πότε το συμπέρασμα ισχύει και πότε όχι.

\ex Λύσε με δύο τρόπους: τα $9$ κιλά κοστίζουν $22{,}5$ €. Πόσο κοστίζουν $14$ κιλά;

\ex Ο Στέλιος λέει: «$3$ εργάτες χτίζουν έναν τοίχο σε $6$ ημέρες, άρα $6$ εργάτες σε $12$ ημέρες». Έχει δίκιο; Πόσες ημέρες περιμένεις να χρειαστούν $6$ εργάτες;

\ex Φτιάξε ένα δικό σου πρόβλημα αναλογίας με τιμή ανά μέτρο και λύσε το με τρεις τρόπους.

\zoni{Για το σπίτι}{στην κουζίνα}

\ex Πάρε μια συνταγή από το σπίτι σου και μετέτρεψέ την για διπλάσιο αριθμό ατόμων και για μισό αριθμό ατόμων.

\ex Βρες την τιμή ενός προϊόντος ανά κιλό και ανά $100$ γραμμάρια και σύγκρινε δύο συσκευασίες.

\paragrafos{Αντιστρόφως ανάλογα ποσά: περισσότερο το ένα, λιγότερο το άλλο}{Α.6.6, «Αντιστρόφως ανάλογα ποσά»}

\begin{anagke}
Το Σάββατο η Α3 πρέπει να στήσει τα τραπέζια του παζαριού στην αυλή. Ο Στέλιος μετράει: «Πέντε παιδιά, στήσιμο σε $36$ λεπτά». Τότε ο Νίκος λέει «Θα φέρω κι άλλα παιδιά» και ο Στέλιος αναρωτιέται: «Αν διπλασιάσουμε τα παιδιά, θα διπλασιαστεί ο χρόνος;». Όλοι γελούν: με περισσότερα χέρια θα τελειώσουν πιο γρήγορα. Η Μυρτώ παρατηρεί: «Δύο φορές περισσότερα παιδιά, μισός χρόνος».

Η Πολυξένη το γράφει. «Τα πέντε παιδιά κάνουν $36$ λεπτά, δηλαδή $180$ «λεπτά δουλειάς παιδιού». Τα $10$ παιδιά πρέπει να κάνουν κι αυτά $180$ συνολικά, άρα $18$ λεπτά». Ο Βασίλης δοκιμάζει με $9$ παιδιά και βρίσκει $20$, και παρατηρεί ότι το γινόμενο παίδες επί λεπτά βγαίνει πάντα $180$.

Το ίδιο συμβαίνει στο ταξίδι: αν διπλασιάσεις την ταχύτητα, ο χρόνος για την ίδια διαδρομή γίνεται μισός. Και στη μοιρασιά: όσο περισσότερα παιδιά, τόσο λιγότερα μπισκότα το καθένα. Εδώ το ένα ποσό μεγαλώνει και το άλλο μικραίνει με τέτοιον τρόπο, ώστε το γινόμενό τους να μένει ίδιο.
\end{anagke}

\begin{lexi}{αντιστρόφως ανάλογα}{αντιστρόφως ανάλογα}
Τα \lexiword{αντιστρόφως ανάλογα} ποσά είναι αυτά που αλλάζουν «ανάλογα, αλλά ανάποδα» (αντίστροφα): όσες φορές μεγαλώνει το ένα, τόσες μικραίνει το άλλο. Ο \lexiword{αντίστροφος} ενός αριθμού είναι αυτός που, πολλαπλασιασμένος με αυτόν, δίνει $1$.
\end{lexi}

\begin{metafrasi}
Διπλάσια παιδιά, μισός χρόνος. & $x\to2x\ \Rightarrow\ y\to\dfrac{y}{2}$\\
Το γινόμενο των δύο ποσών μένει σταθερό. & $x\cdot y=\alpha$\\
$5$ παιδιά και $36$ λεπτά. & $5\cdot36=180$\\
Τα $10$ παιδιά θα κάνουν. & $y=\dfrac{180}{10}=18$ λεπτά\\
\end{metafrasi}

\tmima{Η θεωρία}

\begin{orismos}{αντιστρόφως ανάλογα ποσά}
Δύο ποσά λέγονται \textbf{αντιστρόφως ανάλογα} όταν μεταβάλλονται με τέτοιον τρόπο ώστε, όταν το ένα πολλαπλασιάζεται επί έναν αριθμό, το άλλο διαιρείται με τον ίδιο αριθμό.
\end{orismos}

Σε αντιστρόφως ανάλογα ποσά $x,\,y$ το \textbf{γινόμενο των αντίστοιχων τιμών} είναι σταθερό:
$$x\cdot y=\alpha\quad(\alpha\neq0)\qquad\text{ή}\qquad y=\dfrac{\alpha}{x}.$$
Στην ειδική περίπτωση $\alpha=1$ τα $x,\,y$ είναι αντίστροφοι αριθμοί ($x\cdot y=1$).

\tmimasmall{Πίνακας και γραφική παράσταση}

Για την ταχύτητα $x$ και τον χρόνο $y$ μιας διαδρομής $120$ km είναι $x\cdot y=120$:
\begin{center}
\renewcommand{\arraystretch}{1.3}
\begin{tabular}{c|cccccc}
ταχύτητα $x$ (km/h) & $30$ & $40$ & $60$ & $80$ & $120$ & $240$\\\hline
χρόνος $y$ (h) & $4$ & $3$ & $2$ & $1{,}5$ & $1$ & $0{,}5$
\end{tabular}
\end{center}
Αν τοποθετήσουμε τα σημεία $(x,\,y)$ σε σύστημα ημιαξόνων, \textbf{δεν} βρίσκονται σε ημιευθεία αλλά πάνω σε μια \textbf{καμπύλη γραμμή}, που πλησιάζει τους άξονες χωρίς να τους ακουμπά.

\begin{center}
\begin{tikzpicture}[x=0.05cm,y=0.095cm]
\draw[gray!35,very thin] (0,0) grid[xstep=30,ystep=1] (250,4.4);
\draw[thick,-{Stealth}] (0,0)--(262,0) node[below] {$x$};
\draw[thick,-{Stealth}] (0,0)--(0,4.9) node[left] {$y$};
\foreach \i in {30,60,120,240}{\node[below,font=\scriptsize] at (\i,0) {\i};}
\foreach \j in {1,2,3,4}{\node[left,font=\scriptsize] at (0,\j) {\j};}
\draw[gblue,thick,domain=28:240,samples=60,smooth] plot (\x,{120/\x});
\foreach \x/\y in {30/4,40/3,60/2,80/1.5,120/1,240/0.5}{\fill[gblue] (\x,\y) circle (2.3pt);}
\end{tikzpicture}
\end{center}

\begin{prosoxi}
\begin{itemize}[leftmargin=1.2em,itemsep=0.3ex,topsep=2pt]
\item Στα αντιστρόφως ανάλογα ποσά, αν το ένα διπλασιάζεται, το άλλο \emph{γίνεται το μισό} (δεν διπλασιάζεται).
\item Δεν ενώνουμε τα σημεία με ημιευθεία και δεν περνά η γραμμή από το $O$: η καμπύλη δεν κόβει τους άξονες (δεν μπορώ να έχω $x=0$ ή $y=0$).
\item Δεν είναι κάθε ζευγάρι ποσών που το ένα μεγαλώνει και το άλλο μικραίνει αντιστρόφως ανάλογο: πρέπει το γινόμενο να μένει σταθερό.
\item Για να δω το είδος της σχέσης, ρωτάω: με διπλάσιο $x$, το $y$ γίνεται διπλάσιο (ανάλογα), μισό (αντιστρόφως ανάλογα) ή τίποτα από τα δύο;
\end{itemize}
\end{prosoxi}

\tmima{Λυμένα παραδείγματα}

\paradeigma{Τα τραπέζια του παζαριού}
\tixero Τα $5$ παιδιά στήνουν τα τραπέζια σε $36$ λεπτά.\par
\tithelo Πόση ώρα θα κάνουν $10$, $9$ και $12$ παιδιά;\par
\possindeo Η δουλειά είναι $5\cdot36=180$ «λεπτά παιδιού». Με $10$: $180:10=18$. Με $9$: $20$. Με $12$: $15$.\par
\apantisi $18$, $20$ και $15$ λεπτά.\par
\elegxos $10\cdot18=9\cdot20=12\cdot15=5\cdot36=180$.

\paradeigma{Η ταχύτητα και ο χρόνος}
\tixero Η διαδρομή είναι $120$ km.\par
\tithelo Ο χρόνος για ταχύτητες $60$, $80$, $48$ km/h.\par
\possindeo $x\cdot y=120$. Για $60$: $2$ ώρες. Για $80$: $1{,}5$. Για $48$: $120:48=2{,}5$.\par
\apantisi $2$, $1{,}5$ και $2{,}5$ ώρες.\par
\elegxos $60\cdot2=80\cdot1{,}5=48\cdot2{,}5=120$.

\paradeigma{Το ορθογώνιο με εμβαδόν $48$}
\tixero Ορθογώνιο με εμβαδόν $48$ m$^2$ και διαστάσεις $x,\,y$.\par
\tithelo Οι ακέραιες διαστάσεις και ποιο έχει τη μικρότερη περίμετρο.\par
\possindeo $x\cdot y=48$: $1\times48$, $2\times24$, $3\times16$, $4\times12$, $6\times8$. Περίμετροι $2(x+y)$: $98$, $52$, $38$, $32$, $28$.\par
\apantisi Η μικρότερη περίμετρος είναι η $28$ m για το $6\times8$ (το πιο «τετράγωνο»).\par
\elegxos Όσο περισσότερο μοιάζει με τετράγωνο, τόσο μικρότερη η περίμετρος για το ίδιο εμβαδόν.

\paradeigma{Είναι αντιστρόφως ανάλογα;}
\tixero Πίνακες: (Α) $x=2,4,6$ και $y=12,6,4$. (Β) $x=2,4,8$ και $y=10,6,3$.\par
\tithelo Ποιος δίνει αντιστρόφως ανάλογα ποσά.\par
\possindeo Γινόμενο: Α: $24$, $24$, $24$. Β: $20$, $24$, $24$.\par
\apantisi Ο Α ναι ($\alpha=24$). Ο Β όχι.\par
\elegxos Στον Α το $x$ διπλασιάζεται από $2$ σε $4$ και το $y$ γίνεται το μισό ($12\to6$).

\paradeigma{Αντίστροφοι αριθμοί}
\tixero Οι αριθμοί $5$, $0{,}5$, $\frac{3}{4}$ και $1{,}25$.\par
\tithelo Οι αντίστροφοί τους.\par
\possindeo Ο αντίστροφος του $x$ είναι ο $y$ με $xy=1$, δηλαδή $y=\frac{1}{x}$. Για το $5$: $\frac{1}{5}$. Για το $0{,}5=\frac{1}{2}$: $2$. Για το $\frac{3}{4}$: $\frac{4}{3}$. Για το $1{,}25=\frac{5}{4}$: $\frac{4}{5}=0{,}8$.\par
\apantisi $\frac{1}{5}$, $2$, $\frac{4}{3}$, $0{,}8$.\par
\elegxos $1{,}25\cdot0{,}8=1$.

\tmima{Ασκήσεις}

\zoni{Το κατάλαβα}{γινόμενα, πίνακες, αντίστροφοι}

\ex Ποιοι πίνακες δίνουν αντιστρόφως ανάλογα ποσά; α) $x$: $2,4,6,12$, $y$: $12,6,4,2$ \quad β) $x$: $1,2,3$, $y$: $5,10,15$ \quad γ) $x$: $3,6,9$, $y$: $8,4,3$ \quad δ) $x$: $5,10,20$, $y$: $8,4,2$.

\ex Τα $x,\,y$ είναι αντιστρόφως ανάλογα με $x\cdot y=36$. Συμπλήρωσε τον πίνακα για $x=1,\,2,\,3,\,4,\,6,\,9$.

\ex Τα $x,\,y$ έχουν $x\cdot y=60$. Βρες το $y$ για $x=5,\ 15,\ 0{,}5$ και το $x$ για $y=30$.

\ex Βρες τους αντίστροφους των $5$, $0{,}5$, $\frac{3}{4}$, $1{,}25$.

\ex Σωστό ή λάθος; α) Σε αντιστρόφως ανάλογα ποσά, όταν το $x$ διπλασιάζεται το $y$ γίνεται το μισό. \quad β) Το γινόμενο $x\cdot y$ είναι σταθερό. \quad γ) Όταν $x=0$, είναι $y=0$. \quad δ) Τα σημεία βρίσκονται σε ημιευθεία από το $O$. \quad ε) Αν $8$ εργάτες κάνουν ένα έργο σε $6$ ημέρες, $12$ εργάτες θέλουν $9$ ημέρες.

\zoni{Το χρησιμοποιώ}{παιδιά, ταχύτητες, δοχεία}

\ex \textbf{Το στήσιμο.} Τα $5$ παιδιά στήνουν το παζάρι σε $36$ λεπτά. Συμπλήρωσε τον πίνακα για $1,\,2,\,3,\,4,\,6,\,9,\,10,\,12,\,15,\,18$ παιδιά.

\ex \textbf{Η ταχύτητα.} Η διαδρομή είναι $120$ km. Βρες τον χρόνο για ταχύτητες $30,\,40,\,48,\,60,\,80,\,120$ km/h.

\ex \textbf{Το ορθογώνιο.} Ένα ορθογώνιο έχει εμβαδόν $48$ m$^2$. Βρες τις ακέραιες διαστάσεις του και πες ποιο έχει τη μικρότερη περίμετρο.

\ex \textbf{Οι βρύσες.} Τέσσερις ίδιες βρύσες γεμίζουν μια δεξαμενή σε $6$ ώρες. Σε πόσες ώρες θα τη γεμίσουν $3$ βρύσες και σε πόσες $8$;

\ex \textbf{Τα μπισκότα.} Μοιράζουμε $72$ μπισκότα εξίσου σε $x$ παιδιά. Βρες πόσα παίρνει το καθένα για $x=8,\ 12,\ 18$.

\ex \textbf{Το διάβασμα.} Το βιβλίο έχει $60$ σελίδες. Πόσες ώρες θέλει η Ζωή να το διαβάσει με $20$ σελίδες ανά ώρα και πόσες με $30$ σελίδες ανά ώρα; Πόσες σελίδες την ώρα πρέπει να διαβάζει για να το τελειώσει σε $1{,}5$ ώρες;

\ex \textbf{Το ταμείο.} Τα $40$ € του ταμείου μοιράζονται εξίσου σε $8$ παιδιά και σε $10$ παιδιά. Πόσο παίρνει το καθένα στην κάθε περίπτωση;

\zoni{Το σκέφτομαι}{ανάλογα ή αντιστρόφως;}

\ex Ο Στέλιος λέει «διπλάσια παιδιά, διπλάσιος χρόνος». Πού κάνει λάθος;

\ex Τοποθέτησε στο ίδιο σύστημα αξόνων τα σημεία $(1,\,48)$, $(2,\,24)$, $(3,\,16)$, $(4,\,12)$, $(6,\,8)$, $(8,\,6)$, $(12,\,4)$ και ένωσέ τα. Είναι ημιευθεία; Φτάνει η γραμμή στους άξονες;

\ex Αν η ταχύτητα διπλασιαστεί, τι γίνεται ο χρόνος για την ίδια διαδρομή; Και αν τριπλασιαστεί;

\ex Πες για κάθε ζευγάρι αν τα ποσά είναι ανάλογα, αντιστρόφως ανάλογα ή τίποτα από τα δύο: α) κιλά και τιμή \quad β) ταχύτητα και χρόνος \quad γ) ηλικία και ύψος \quad δ) αριθμός ατόμων και μερίδιο από σταθερό ποσό \quad ε) πλευρά και εμβαδόν τετραγώνου.

\ex Αν τα $x,\,y$ έχουν $xy=36$, δείξε ότι τα $y$ και $\frac{1}{x}$ είναι ανάλογα και βρες τον συντελεστή.

\zoni{Για το σπίτι}{πόσο γρήγορα φτάνω;}

\ex Μέτρησε την απόσταση από το σπίτι ως το σχολείο και υπολόγισε τον χρόνο για ταχύτητες περπατήματος, ποδηλάτου και αυτοκινήτου.

\ex Βρες δύο παραδείγματα αντιστρόφως αναλόγων ποσών από την καθημερινή ζωή.

\epanalipsi{Α6}

\tmima{Λεξικό του κεφαλαίου}

\renewcommand{\arraystretch}{1.5}
\begin{tabular*}{\linewidth}{@{\extracolsep{\fill}} >{\raggedright\arraybackslash}p{0.28\linewidth} >{\raggedright\arraybackslash}p{0.62\linewidth}}
\toprule
\textbf{Λέξη ή σύμβολο} & \textbf{Με μία πρόταση}\\
\midrule
ισότητα $\alpha=\beta$ & Τα δύο μέλη έχουν την ίδια τιμή, όπως μια ζυγαριά που ισορροπεί.\\
ιδιότητες της ισότητας & Αν $\alpha=\beta$, τότε $\alpha\pm\gamma=\beta\pm\gamma$, $\alpha\gamma=\beta\gamma$ και $\alpha:\gamma=\beta:\gamma$ (με $\gamma\neq0$).\\
εξίσωση & Ισότητα με έναν άγνωστο αριθμό, π.χ. $2x+6=14$.\\
λύση ή ρίζα & Ο αριθμός που, στη θέση του αγνώστου, κάνει την ισότητα αληθινή.\\
επίλυση, επαλήθευση & Η εύρεση της λύσης· ο έλεγχος της λύσης στην αρχική εξίσωση.\\
ισοδύναμες εξισώσεις & Εξισώσεις με τις ίδιες λύσεις.\\
αδύνατη, ταυτότητα & Καμία λύση (π.χ. $0\cdot x=2$)· όλοι οι αριθμοί είναι λύση ($0\cdot x=0$).\\
$x=\dfrac{\gamma-\beta}{\alpha}$ & Η λύση της $\alpha x+\beta=\gamma$, με $\alpha\neq0$.\\
\bottomrule
\end{tabular*}

\tmima{Ασκήσεις επανάληψης}

\ex Αν $\alpha=\beta$, συμπλήρωσε με το σωστό: $\alpha+7=\dots$, \quad $\alpha-2=\dots$, \quad $5\alpha=\dots$, \quad $\alpha:4=\dots$

\ex Εξέτασε αν το $4$ είναι λύση των εξισώσεων: α) $x+9=13$ \quad β) $3x=15$ \quad γ) $20-x=16$ \quad δ) $2x+1=10$.

\ex Λύσε: α) $x+19=40$ \quad β) $x-2{,}5=4$ \quad γ) $30-x=12$ \quad δ) $6x=42$ \quad ε) $x:7=3$ \quad στ) $56:x=8$.

\ex Λύσε και επαλήθευσε: α) $-4x=28$ \quad β) $x+(-3)=-10$ \quad γ) $5-x=-2$.

\ex Λύσε: α) $5x+3=28$ \quad β) $4x-9=15$ \quad γ) $2-3x=14$ \quad δ) $0{,}5x+2=3{,}5$.

\ex Λύσε: α) $\dfrac{2}{3}x+1=5$ \quad β) $\dfrac{x}{4}-2=1$.

\ex Λύσε: $2(x+5)-3(x-1)=10$.

\ex Λύσε: $3(2x-1)-(x+4)=12$.

\ex Εξέτασε αν η εξίσωση $4(x+2)=4x+8$ είναι αδύνατη ή ταυτότητα και αν η $3x+1=3x-4$ είναι αδύνατη ή ταυτότητα.

\ex Γράψε εξίσωση και λύσε: το τριπλάσιο ενός αριθμού αυξημένο κατά $8$ κάνει $29$.

\ex Η Αγγελική αγόρασε $3$ ίδια τετράδια και μια κασετίνα των $6$ €. Πλήρωσε συνολικά $15$ €. Πόσο κοστίζει κάθε τετράδιο;

\ex Η Ζωή έχει $x$ ευρώ και ο Νίκος $7$ ευρώ περισσότερα. Μαζί έχουν $33$ ευρώ. Πόσα έχει η Ζωή;

\ex Ένα ορθογώνιο έχει πλάτος $8$ cm και περίμετρο $54$ cm. Βρες το μήκος του με εξίσωση.

\ex Ο παππούς Γιάννης λέει: «Αν διπλασιάσεις την ηλικία μου και αφαιρέσεις $20$, βρίσκεις $130$». Πόσο ετών είναι;

\ex \textbf{Σωστό ή λάθος (ισότητα).} α) Αν $\alpha=\beta$, τότε $\alpha-5=\beta-5$. \quad β) Αν $3\alpha=3\beta$, τότε $\alpha=\beta$. \quad γ) Η $0\cdot x=5$ είναι ταυτότητα. \quad δ) Αν $\alpha\cdot0=\beta\cdot0$, τότε $\alpha=\beta$.

\ex \textbf{Σωστό ή λάθος (εξισώσεις).} α) Η $2x=2$ έχει λύση το $2$. \quad β) Η $x+1=x+2$ είναι αδύνατη. \quad γ) Η λύση της $\alpha x+\beta=\gamma$ ($\alpha\neq0$) είναι $x=\dfrac{\gamma-\beta}{\alpha}$. \quad δ) Οι $x+3=7$ και $2x=8$ είναι ισοδύναμες.

\tmima{Ο χάρτης της επανάληψης}

Μετά τις λύσεις, δες σε ποιες ασκήσεις δυσκολεύτηκες. Ο πίνακας σου λέει ποια παράγραφο να ξαναδιαβάσεις.

\renewcommand{\arraystretch}{1.5}
\begin{tabular*}{\linewidth}{@{\extracolsep{\fill}} >{\raggedright\arraybackslash}p{0.38\linewidth} >{\raggedright\arraybackslash}p{0.52\linewidth}}
\toprule
\textbf{Αν δυσκολεύτηκες στις} & \textbf{ξαναδιάβασε τις παραγράφους}\\
\midrule
ΕΑ6.1, ΕΑ6.14 & 6.1 (ιδιότητες της ισότητας)\\
ΕΑ6.2, ΕΑ6.3, ΕΑ6.4, ΕΑ6.9 & 6.1 (εξίσωση, βασικές μορφές, ειδικές περιπτώσεις)\\
ΕΑ6.5, ΕΑ6.6, ΕΑ6.7, ΕΑ6.8 & 6.2 (η μέθοδος $\alpha x+\beta=\gamma$)\\
ΕΑ6.10, ΕΑ6.11, ΕΑ6.12, ΕΑ6.13, ΕΑ6.15 & 6.2 (προβλήματα με εξίσωση)\\
ΕΑ6.16 & 6.1 και 6.2\\
\bottomrule
\end{tabular*}


\kefalaio{Θετικοί και αρνητικοί αριθμοί}{Ο κόσμος των αριθμών έχει δύο πλευρές}{%
\begin{tabular}{@{}l@{\hspace{1em}}p{0.82\linewidth}@{}}
7.1 & Θετικοί και αρνητικοί αριθμοί (ρητοί) – Η ευθεία των ρητών – Τετμημένη σημείου\\[2pt]
7.2 & Απόλυτη τιμή ρητού – Αντίθετοι ρητοί – Σύγκριση ρητών\\[2pt]
7.3 & Πρόσθεση ρητών αριθμών\\[2pt]
7.4 & Αφαίρεση ρητών αριθμών\\[2pt]
7.5 & Πολλαπλασιασμός ρητών αριθμών\\[2pt]
7.6 & Διαίρεση ρητών αριθμών\\[2pt]
7.7 & Δεκαδική μορφή ρητών αριθμών\\[2pt]
7.8 & Δυνάμεις ρητών αριθμών με εκθέτη φυσικό\\[2pt]
7.9 & Δυνάμεις ρητών αριθμών με εκθέτη ακέραιο\\[2pt]
7.10 & Τυποποιημένη μορφή μεγάλων και μικρών αριθμών\\[2pt]
 & Επανάληψη του κεφαλαίου\\
\end{tabular}}
\paragrafos{Θετικοί και αρνητικοί αριθμοί: ο κόσμος κάτω από το μηδέν}{Α.7.1, «Θετικοί και αρνητικοί αριθμοί (ρητοί αριθμοί) – Η ευθεία των ρητών – Τετμημένη σημείου»}

\begin{anagke}
Η Α3 κάνει εκδρομή στο Μέτσοβο. Στο πούλμαν το ψηφιακό θερμόμετρο γράφει $-4$. «Τι είναι αυτό το μείον;» ρωτάει η Πολυξένη. «Τέσσερις βαθμούς κάτω από το μηδέν», απαντάει ο Νίκος. Ο Στέλιος ταράζεται: «Κάτω από το μηδέν; Το μηδέν δεν είναι το τίποτα; Πώς μπορεί να υπάρχει κάτι μικρότερο από το τίποτα;».

Στο ξενοδοχείο ο ανελκυστήρας έχει κουμπιά $-2$, $-1$, $0$, $1$, $2$, $3$. Ο Ηλίας πατάει το $-1$ και κατεβαίνει στο υπόγειο με το πάρκινγκ. Η Αγγελική βλέπει στον πίνακα του καιρού: Αθήνα $+7$, Ιωάννινα $-3$, Φλώρινα $-10$, Χανιά $+11$. Η Μυρτώ θυμάται ότι στο ταμείο της παρέας κάποτε έλειπαν $8$ ευρώ, ένα χρέος. Και στο ντοκιμαντέρ που είδαν, ο καρχαρίας βρισκόταν $200$ μέτρα κάτω από την επιφάνεια ενώ το αεροπλάνο πετούσε $200$ μέτρα πάνω από αυτήν.

Η κυρία Αλεξάνδρα το συνοψίζει: «Οι φυσικοί, οι δεκαδικοί και τα κλάσματα που ξέρετε μετράνε πόσα. Όταν όμως χρειάζεται να πούμε «πάνω ή κάτω», «κέρδος ή ζημιά», «πάνω ή κάτω από το μηδέν», χρειαζόμαστε και αριθμούς με κατεύθυνση». Κάθε «δεξιά» προϋποθέτει ένα «αριστερά» και κάθε «πάνω» ένα «κάτω»: ο κόσμος των αριθμών έχει δύο πλευρές.
\end{anagke}

\begin{lexi}{θετικός, αρνητικός, ρητός}{θετικός, αρνητικός, ρητός}
\lexiword{Θετικός} είναι αυτός που «τίθεται», που βάζουμε ως βάση. \lexiword{Αρνητικός} είναι αυτός που «αρνείται», που πηγαίνει ανάποδα από τον θετικό. Το \lexiword{πρόσημο} είναι το «προ-σήμα», το σημάδι που μπαίνει πριν από τον αριθμό. Ο \lexiword{ακέραιος} είναι «ακέραιος», ολόκληρος. Ο \lexiword{ρητός} (από το «ρητός», αυτός που μπορεί να ειπωθεί) είναι ο αριθμός που εκφράζεται ως κλάσμα, δηλαδή ως λόγος δύο ακεραίων.
\end{lexi}

\begin{metafrasi}
Τέσσερις βαθμοί κάτω από το μηδέν. & $-4$\\
Αεροπλάνο $200$ μ. πάνω από τη θάλασσα. & $+200$\\
Καρχαρίας $200$ μ. κάτω από την επιφάνεια. & $-200$\\
Χρέος $8$ ευρώ. & $-8$\\
Κέρδος $15$ ευρώ. & $+15$\\
Δεύτερο υπόγειο. & $-2$\\
\end{metafrasi}

\tmima{Η θεωρία}

Τα σύμβολα «$+$» και «$-$» λέγονται \textbf{πρόσημα}. Γράφονται πριν από τους αριθμούς και τους χαρακτηρίζουν, αντίστοιχα, ως \textbf{θετικούς} ή \textbf{αρνητικούς}.
$$+3,\ +\tfrac{3}{4},\ +15{,}7,\ +0{,}352\ \text{ είναι θετικοί αριθμοί},\qquad -2,\ -10,\ -28{,}95,\ -0{,}098\ \text{ είναι αρνητικοί}.$$
Όταν μιλάμε μόνο για θετικούς, μπορούμε να παραλείπουμε το $+$: γράφουμε $7$ αντί για $+7$. Το $0$ \textbf{δεν είναι ούτε θετικός ούτε αρνητικός}.

\begin{orismos}{ομόσημοι, ετερόσημοι, ακέραιοι, ρητοί}
Δύο αριθμοί είναι \textbf{ομόσημοι} όταν έχουν το ίδιο πρόσημο και \textbf{ετερόσημοι} όταν έχουν διαφορετικό. \textbf{Ακέραιοι} είναι οι φυσικοί μαζί με τους αντίστοιχους αρνητικούς: $\dots,\,-3,\,-2,\,-1,\,0,\,1,\,2,\,3,\dots$ \textbf{Ρητοί} είναι όλοι οι αριθμοί που ξέρουμε (φυσικοί, κλάσματα, δεκαδικοί) μαζί με τους αντίστοιχους αρνητικούς.
\end{orismos}

Για παράδειγμα οι $-7$, $-0{,}58$, $-\frac{3}{4}$ είναι ομόσημοι (όλοι αρνητικοί), οι $+1{,}25$, $+10$, $5$ επίσης, ενώ οι $-7$ και $+0{,}58$ είναι ετερόσημοι.

\tmimasmall{Μια ιστορική σημείωση}

Οι αρνητικοί αριθμοί εμφανίζονται για πρώτη φορά σε ένα κινέζικο μαθηματικό βιβλίο, τα «Μαθηματικά σε εννέα βιβλία», που τοποθετείται στην εποχή της δυναστείας των Χαν (206 π.Χ. -- 220 μ.Χ.). Στην ευρωπαϊκή παράδοση ο Διόφαντος (γύρω στο 250 μ.Χ.) τους χρησιμοποιεί στους ενδιάμεσους υπολογισμούς του, αλλά ως λύση ενός προβλήματος δέχεται πάντα θετικό αριθμό. Για πολλούς αιώνες οι μαθηματικοί τους θεωρούσαν «παράλογους» ή «ψεύτικους» αριθμούς. Σήμερα τους χρησιμοποιούμε στο θερμόμετρο, στον λογαριασμό της τράπεζας, στο υψόμετρο.

\tmimasmall{Η ευθεία των ρητών και η τετμημένη}

Στον ημιάξονα $Ox$ των αριθμών τοποθετούσαμε τους φυσικούς, δεξιά από την αρχή $O$. Αν προεκτείνουμε προς τα αριστερά (στον αντικείμενο ημιάξονα $Ox'$) με την ίδια μονάδα, παίρνουμε τον \textbf{άξονα $x'x$}, την ευθεία των ρητών. Τα σημεία δεξιά του $O$ παριστάνουν θετικούς και τα αριστερά αρνητικούς. Ο αριθμός που αντιστοιχεί σε ένα σημείο λέγεται \textbf{τετμημένη} του σημείου και γράφουμε $A(-3)$.

\begin{center}
\akgr{-5}{5}{-3/A,2.5/B,-1.5/\Gamma,4/\Delta}
\end{center}

Στο σχήμα: $A(-3)$, $B(2{,}5)$, $\Gamma(-1{,}5)$, $\Delta(4)$. Στην ευθεία των ρητών οι αρνητικοί βρίσκονται αριστερά του $0$ και, όσο πηγαίνουμε αριστερά, οι αριθμοί μικραίνουν.

\begin{prosoxi}
\begin{itemize}[leftmargin=1.2em,itemsep=0.3ex,topsep=2pt]
\item Το $-3$ δεν είναι «τρία ευρώ λιγότερο από το τίποτα». Είναι μια θέση, τρεις μονάδες αριστερά του $0$, και η μονάδα μπορεί να είναι βαθμός, ευρώ ή μέτρο.
\item Το $0$ δεν έχει πρόσημο. Δεν γράφουμε $+0$ ή $-0$.
\item Στη μονάδα του άξονα πρέπει οι αποστάσεις να είναι ίσες: το $-1{,}5$ είναι στη μέση του $-2$ και του $-1$.
\item Μην μπερδεύεις το πρόσημο με την πράξη: στο $5-3$ το μείον είναι αφαίρεση, στο $-3$ είναι πρόσημο.
\end{itemize}
\end{prosoxi}

\tmima{Λυμένα παραδείγματα}

\paradeigma{Αριθμοί για μεγέθη}
\tixero Οι πόλεις και οι θερμοκρασίες: Αθήνα $7$° πάνω από το $0$, Ιωάννινα $3$° κάτω, Φλώρινα $10$° κάτω. Επίσης ένα ντεπόζιτο είναι $12$ ευρώ χρέος.\par
\tithelo Να τα γράψω ως ρητούς.\par
\possindeo Πάνω από το μηδέν είναι θετικό και κάτω αρνητικό. Αθήνα $+7$, Ιωάννινα $-3$, Φλώρινα $-10$. Το χρέος είναι $-12$.\par
\apantisi $+7$, $-3$, $-10$ και $-12$.\par
\elegxos Η Φλώρινα ($-10$) είναι πιο κρύα από τα Ιωάννινα ($-3$): θα το συγκρίνουμε στην επόμενη παράγραφο.

\paradeigma{Ομόσημοι ή ετερόσημοι;}
\tixero Τα ζεύγη $(-3,\,-8)$, $(+4,\,-4)$, $(0{,}5,\,\frac{7}{2})$, $(-2,\,2{,}5)$.\par
\tithelo Να πω τι είναι.\par
\possindeo Κοιτάζω τα πρόσημα. $-3$ και $-8$ είναι αρνητικοί: ομόσημοι. $+4$ και $-4$: ετερόσημοι. $0{,}5$ και $\frac{7}{2}$ θετικοί: ομόσημοι. $-2$ και $2{,}5$: ετερόσημοι.\par
\apantisi Ομόσημοι: το πρώτο και το τρίτο. Ετερόσημοι: το δεύτερο και το τέταρτο.\par
\elegxos Δύο αριθμοί με το ίδιο πρόσημο βρίσκονται στην ίδια πλευρά του $0$ στον άξονα.

\paradeigma{Ακέραιοι και ρητοί}
\tixero Οι αριθμοί $-4$, $2{,}5$, $0$, $\frac{9}{3}$, $-\frac{1}{2}$, $10$.\par
\tithelo Ποιοι είναι ακέραιοι.\par
\possindeo Ακέραιοι είναι οι αριθμοί που δεν έχουν δεκαδικό ή κλασματικό μέρος. $\frac{9}{3}=3$ είναι ακέραιος. Όλοι είναι ρητοί.\par
\apantisi Ακέραιοι: $-4$, $0$, $\frac{9}{3}$, $10$. Οι $2{,}5$ και $-\frac{1}{2}$ είναι ρητοί αλλά όχι ακέραιοι.\par
\elegxos Κάθε ακέραιος είναι ρητός: π.χ. $-4=-\frac{4}{1}$.

\paradeigma{Σημεία στον άξονα}
\tixero Στο σχήμα της θεωρίας.\par
\tithelo Να βρω τετμημένες και να τοποθετήσω τον $-\frac{1}{2}$.\par
\possindeo Το $A$ είναι $3$ μονάδες αριστερά του $0$: $-3$. Το $B$ είναι στη μέση του $2$ και του $3$: $2{,}5$. Το $\Gamma$ στη μέση του $-2$ και $-1$: $-1{,}5$. Ο $-\frac{1}{2}=-0{,}5$ είναι στη μέση του $-1$ και του $0$.\par
\apantisi $A(-3)$, $B(2{,}5)$, $\Gamma(-1{,}5)$, $\Delta(4)$ και ο $-0{,}5$ ανάμεσα στο $-1$ και στο $0$.\par
\elegxos Τα $-3$ και $-1{,}5$ είναι αριστερά του $0$, άρα αρνητικοί.

\tmima{Ασκήσεις}

\zoni{Το κατάλαβα}{πρόσημα, ακέραιοι, άξονας}

\ex Εξέφρασε με αριθμό: α) $5$ βαθμοί κάτω από το μηδέν \quad β) $12$ μ. πάνω από τη θάλασσα \quad γ) χρέος $8$ € \quad δ) κέρδος $15$ € \quad ε) το δεύτερο υπόγειο \quad στ) $30$ μ. κάτω από την επιφάνεια της θάλασσας.

\ex Ποιοι από τους αριθμούς $-7$, $+3{,}5$, $0$, $-\frac{1}{2}$, $8$, $+0{,}01$ είναι θετικοί, ποιοι αρνητικοί και ποιος δεν είναι ούτε το ένα ούτε το άλλο;

\ex Πες αν είναι ομόσημοι ή ετερόσημοι: α) $-3$ και $-8$ \quad β) $+4$ και $-4$ \quad γ) $0{,}5$ και $\frac{7}{2}$ \quad δ) $-2$ και $2{,}5$.

\ex Ποιοι από τους $-4$, $2{,}5$, $0$, $\frac{9}{3}$, $-\frac{1}{2}$, $10$ είναι ακέραιοι;

\ex Σχεδίασε άξονα με μονάδα $1$ cm και σημείωσε τους $-4$, $-1{,}5$, $0$, $2$, $3{,}5$.

\ex Στον άξονα βρες τις τετμημένες των $A,\,B,\,\Gamma,\,\Delta$.
\begin{center}
\akgr{-5}{5}{-4/A,-1.5/B,1/\Gamma,3.5/\Delta}
\end{center}

\ex Σωστό ή λάθος; α) Το $0$ είναι θετικός αριθμός. \quad β) Κάθε φυσικός αριθμός είναι ακέραιος. \quad γ) Κάθε ακέραιος είναι ρητός. \quad δ) Οι $-3$ και $-5$ είναι ετερόσημοι. \quad ε) Το $\frac{2}{3}$ είναι ρητός. \quad στ) Στον άξονα οι αρνητικοί βρίσκονται αριστερά του $0$.

\zoni{Το χρησιμοποιώ}{θερμόμετρα, ασανσέρ, θάλασσες}

\ex \textbf{Ο καιρός.} Οι προβλέψεις: Αθήνα $7$° πάνω από το $0$, Θεσσαλονίκη $3$° κάτω, Ιωάννινα $5$° κάτω, Χανιά $11$° πάνω, Φλώρινα $10$° κάτω. Γράψε τις θερμοκρασίες ως ρητούς και τοποθέτησέ τες σε έναν κατακόρυφο άξονα.

\ex \textbf{Ο ανελκυστήρας.} Οι όροφοι του ξενοδοχείου είναι $-2,\,-1,\,0,\,1,\,2,\,3$. Ο Νίκος είναι στον $3$ και θέλει το δεύτερο υπόγειο. Πόσους ορόφους κατεβαίνει;

\ex \textbf{Βουνό και θάλασσα.} Το αεροπλάνο πετάει στα $+300$ μ., το ψάρι στα $-45$ μ. και το υποβρύχιο στα $-120$ μ. Ποιο από τα τρία είναι πιο κοντά στην επιφάνεια της θάλασσας;

\ex \textbf{Το ταμείο.} Η Μυρτώ έβαλε $20$ €, ο Στέλιος πήρε $7$ €, η Φανή έβαλε $5$ € και ο Βασίλης πήρε $12$ €. Γράψε κάθε κίνηση ως ρητό.

\ex \textbf{Ο καιρός αλλάζει.} Το πρωί στο Μέτσοβο είναι $-3$° και μέχρι το μεσημέρι ανεβαίνει $5$ βαθμούς. Βρες τη νέα θερμοκρασία μετρώντας πάνω στον άξονα.

\zoni{Το σκέφτομαι}{είναι το μηδέν τίποτα;}

\ex Ο Στέλιος ισχυρίζεται ότι δεν υπάρχει αριθμός μικρότερος από το $0$. Φέρε αντιπαράδειγμα από τη ζωή.

\ex Γιατί το $0$ δεν είναι ούτε θετικός ούτε αρνητικός αριθμός;

\ex Τα $0$°C είναι η θερμοκρασία πήξης του νερού και τα $32$°F είναι η ίδια θερμοκρασία σε βαθμούς Φαρενάιτ. Τι δείχνει αυτό για το «μηδέν» στο θερμόμετρο;

\ex Πόσοι ακέραιοι υπάρχουν από το $-5$ ως το $5$, με τα άκρα; Πόσοι ανάμεσα στο $-5$ και στο $5$ χωρίς τα άκρα;

\zoni{Για το σπίτι}{ψάξε τους αρνητικούς γύρω σου}

\ex Βρες πέντε περιπτώσεις της καθημερινής ζωής που χρησιμοποιούνται αρνητικοί αριθμοί.

\ex Μέτρησε τη θερμοκρασία κάθε πρωί για μια εβδομάδα και παράστησέ τη σε άξονα.

\paragrafos{Απόλυτη τιμή, αντίθετοι και σύγκριση: ποιος είναι πιο κρύος;}{Α.7.2, «Απόλυτη τιμή ρητού – Αντίθετοι ρητοί – Σύγκριση ρητών»}

\begin{anagke}
Στο πούλμαν της επιστροφής η παρέα συγκρίνει τις θερμοκρασίες του πίνακα του καιρού: Ιωάννινα $-3$, Φλώρινα $-10$. «Ποια πόλη είναι πιο κρύα;» ρωτάει η Ρήνη. Ο Βασίλης απαντάει αμέσως: «Το δέκα είναι μεγαλύτερο από το τρία, άρα το $-10$ είναι μεγαλύτερο από το $-3$. Άρα η Φλώρινα είναι πιο ζεστή». Ο Νίκος διαμαρτύρεται: «Δεν μπορεί να είναι πιο ζεστό το δέκα βαθμούς κάτω από το μηδέν από το τρεις βαθμούς κάτω!». Η Αγγελική παίρνει ένα χαρτί και σχεδιάζει τον άξονα: «Όσο πιο αριστερά, τόσο πιο κρύο. Το $-10$ είναι πολύ πιο αριστερά από το $-3$».

Ο Ηλίας προσθέτει ένα δεύτερο ερώτημα: «Αν η Ζωή μένει $3$ στάσεις αριστερά από το σχολείο και ο Στέλιος $3$ στάσεις δεξιά, ποιος μένει πιο μακριά;». Οι δύο μένουν το ίδιο μακριά! Το «πόσο μακριά» δεν εξαρτάται από την κατεύθυνση, αλλά το «πού ακριβώς» εξαρτάται. Έτσι έχουμε δύο έννοιες: την απόσταση από το $0$, που λέγεται απόλυτη τιμή, και τη σχετική θέση δεξιά ή αριστερά, που μας βοηθάει να συγκρίνουμε.
\end{anagke}

\begin{lexi}{απόλυτη τιμή, αντίθετος}{απόλυτη τιμή, αντίθετος}
\lexiword{Απόλυτη} σημαίνει «λυμένη από τους δεσμούς» (από το «απολύω»): η απόλυτη τιμή είναι ο αριθμός «απελευθερωμένος» από το πρόσημό του, δηλαδή μόνο η απόσταση από το $0$. Οι \lexiword{αντίθετοι} είναι αριθμοί σε «αντίθετη» κατεύθυνση: ίδια απόσταση από το $0$, αλλά η μία δεξιά και η άλλη αριστερά.
\end{lexi}

\begin{metafrasi}
Η απόσταση του $-3$ από το $0$. & $|-3|=3$\\
Ο αντίθετος του $5$. & $-5$\\
Ο αντίθετος του $-2{,}7$. & $-(-2{,}7)=+2{,}7$\\
Το $-10$ είναι πιο αριστερά από το $-3$. & $-10<-3$\\
Δύο αριθμοί με απόλυτη τιμή $2$. & $+2$ και $-2$\\
\end{metafrasi}

\tmima{Η θεωρία}

\begin{orismos}{απόλυτη τιμή}
\textbf{Απόλυτη τιμή} ενός ρητού $\alpha$ είναι η απόσταση του σημείου με τετμημένη $\alpha$ από την αρχή $O$ του άξονα. Συμβολίζεται $|\alpha|$.
\end{orismos}

\begin{itemize}[leftmargin=1.2em,itemsep=0.3ex,topsep=2pt]
\item Η απόλυτη τιμή ενός θετικού είναι ο ίδιος ο αριθμός: $|+9{,}63|=9{,}63$.
\item Η απόλυτη τιμή ενός αρνητικού είναι ο αντίθετός του: $|-8{,}4|=8{,}4=-(-8{,}4)$.
\item $|0|=0$. Η απόλυτη τιμή δεν είναι ποτέ αρνητική.
\end{itemize}

\begin{orismos}{αντίθετοι αριθμοί}
\textbf{Αντίθετοι} είναι δύο ετερόσημοι αριθμοί που έχουν την ίδια απόλυτη τιμή. Ο αντίθετος του $x$ είναι ο $-x$. Για το $0$ ο αντίθετος είναι το $0$. Ισχύει $-(-x)=x$.
\end{orismos}

Δύο σημεία σε ίση απόσταση από το $O$, δεξιά και αριστερά, έχουν αντίθετες τετμημένες: το $+3$ και το $-3$ απέχουν από το $0$ απόσταση $3$ και $|+3|=|-3|=3$.

\tmimasmall{Σύγκριση ρητών}

\begin{kanonas}{σύγκριση ρητών}
Από δύο ρητούς μεγαλύτερος είναι αυτός που βρίσκεται \textbf{δεξιότερα} στον άξονα. Έτσι:
\begin{itemize}[leftmargin=1.2em,itemsep=0.3ex,topsep=2pt]
\item Κάθε θετικός είναι μεγαλύτερος από κάθε αρνητικό. Το $0$ είναι μικρότερο από κάθε θετικό και μεγαλύτερο από κάθε αρνητικό.
\item Από δύο \textbf{θετικούς} μεγαλύτερος είναι αυτός με τη μεγαλύτερη απόλυτη τιμή.
\item Από δύο \textbf{αρνητικούς} μεγαλύτερος είναι αυτός με τη μικρότερη απόλυτη τιμή: $-6{,}8<-3{,}7$, γιατί $|-6{,}8|=6{,}8>3{,}7=|-3{,}7|$.
\end{itemize}
\end{kanonas}

Έτσι $-4{,}75<-3{,}5<-\frac{3}{4}<-0{,}6<0<0{,}6<\frac{3}{4}<3{,}5<4{,}75$.

\begin{prosoxi}
\begin{itemize}[leftmargin=1.2em,itemsep=0.3ex,topsep=2pt]
\item Στους αρνητικούς το «μεγάλο νούμερο» σημαίνει «μικρός αριθμός»: $-10<-3$ παρά το $10>3$.
\item Η απόλυτη τιμή δεν «σβήνει το μείον» από μόνο του: $|-x|=x$ μόνο αν το $x$ είναι θετικό. Η $|-(-3)|=|3|=3$.
\item Η $|x|=3$ έχει δύο λύσεις: $x=3$ και $x=-3$. Η $|x|=-1$ δεν έχει καμία, γιατί η απόλυτη τιμή δεν είναι αρνητική.
\item Ο αντίθετος του $-5$ είναι το $+5$, όχι το $-5$.
\end{itemize}
\end{prosoxi}

\tmima{Λυμένα παραδείγματα}

\paradeigma{Ποια πόλη είναι πιο κρύα}
\tixero Ιωάννινα $-3$°, Φλώρινα $-10$°.\par
\tithelo Ποια θερμοκρασία είναι μικρότερη.\par
\possindeo Συγκρίνω στον άξονα: το $-10$ είναι αριστερότερα από το $-3$. Ή με απόλυτες τιμές: και οι δύο αρνητικοί και $|-10|=10>3=|-3|$, άρα $-10<-3$.\par
\apantisi $-10<-3$: η Φλώρινα είναι πιο κρύα.\par
\elegxos Από το $-10$ μέχρι το $0$ έχουμε $10$ βαθμούς, από το $-3$ μόνο $3$. Συμφωνεί.

\paradeigma{Τοποθετώ και διατάσσω}
\tixero Οι αριθμοί $-2{,}5$, $-0{,}75$, $\frac{1}{2}$, $-\frac{7}{4}$, $3{,}2$.\par
\tithelo Να τους βάλω σε αύξουσα σειρά.\par
\possindeo Γράφω τους δεκαδικούς: $-2{,}5$, $-0{,}75$, $0{,}5$, $-1{,}75$, $3{,}2$. Οι αρνητικοί είναι αριστερά του $0$: από αυτούς μικρότερος είναι αυτός με τη μεγαλύτερη απόλυτη τιμή: $-2{,}5<-1{,}75<-0{,}75$. Μετά οι θετικοί $0{,}5<3{,}2$.\par
\apantisi $-2{,}5<-\frac{7}{4}<-0{,}75<\frac{1}{2}<3{,}2$.\par
\elegxos Στον άξονα, με τη σειρά από αριστερά προς τα δεξιά, εμφανίζονται οι ίδιοι.

\begin{center}
\akgr[1.3]{-3}{4}{-2.5/{-2{,}5},-1.75/{-\frac{7}{4}},-0.75/{-0{,}75},0.5/{\frac{1}{2}},3.2/{3{,}2}}
\end{center}

\paradeigma{Αντίθετος}
\tixero Το σημείο $M$ έχει τετμημένη $+4{,}5$ και το $N$ τετμημένη $-\frac{2}{3}$.\par
\tithelo Οι αντίθετοι.\par
\possindeo Ο αντίθετος του $x$ είναι ο $-x$. Του $+4{,}5$ είναι ο $-4{,}5$. Του $-\frac{2}{3}$ είναι ο $-\left(-\frac{2}{3}\right)=+\frac{2}{3}$.\par
\apantisi $-4{,}5$ και $+\frac{2}{3}$.\par
\elegxos Το άθροισμα απόστασης: $|{+4{,}5}|=|{-4{,}5}|$ και $|-\frac{2}{3}|=|+\frac{2}{3}|$.

\paradeigma{Εξισώσεις με απόλυτη τιμή}
\tixero Οι εξισώσεις $|x|=2$, $|x|=0$, $|x|=-1$.\par
\tithelo Να βρω τους $x$.\par
\possindeo Ψάχνω αριθμούς με απόσταση $2$ από το $0$: $+2$ και $-2$. Με απόσταση $0$: μόνο το $0$. Απόσταση $-1$ δεν υπάρχει.\par
\apantisi $x=2$ ή $x=-2$, $x=0$, αδύνατη.\par
\elegxos $|2|=|-2|=2$.

\paradeigma{Πράξεις με απόλυτες τιμές και αντίθετους}
\tixero Οι παραστάσεις $|-6|\cdot|-2|$, $-|-5|$ και $-(-(-2))$.\par
\tithelo Να τις υπολογίσω.\par
\possindeo $|-6|=6$ και $|-2|=2$, άρα $6\cdot2=12$. Το $|-5|=5$ και ο αντίθετός του είναι $-5$. Ο αντίθετος του $-2$ είναι $2$, και ο αντίθετος του $2$ είναι $-2$.\par
\apantisi $12$, $-5$, $-2$.\par
\elegxos Κάθε δεύτερο «μείον» αναιρεί το προηγούμενο: δύο φορές $-$ επιστρέφει το αρχικό.

\tmima{Ασκήσεις}

\zoni{Το κατάλαβα}{απόλυτη τιμή, αντίθετοι, σύγκριση}

\ex Βρες τις απόλυτες τιμές: α) $|+7|$ \quad β) $|-7|$ \quad γ) $|0|$ \quad δ) $|-2{,}5|$ \quad ε) $\left|-\frac{3}{4}\right|$ \quad στ) $|+0{,}01|$.

\ex Βρες τους αντίθετους των: $5$, $-8$, $0$, $2{,}5$, $-\frac{1}{3}$, $+0{,}6$.

\ex Απλοποίησε: α) $-(-4)$ \quad β) $-(+3)$ \quad γ) $-(-(-2))$ \quad δ) $-|-6|$.

\ex Βρες τον $\alpha$: α) $|\alpha|=5$ \quad β) $|\alpha|=0$ \quad γ) $|\alpha|=\frac{1}{2}$ \quad δ) $|\alpha|=-3$.

\ex Σύγκρινε: α) $-3$ και $-8$ \quad β) $+2$ και $-9$ \quad γ) $-0{,}5$ και $-0{,}05$ \quad δ) $-\frac{7}{2}$ και $-3$ \quad ε) $0$ και $-0{,}001$ \quad στ) $-4$ και $4$.

\ex Γράψε σε αύξουσα σειρά: $-3$, $5$, $-7$, $0$, $2{,}5$, $-0{,}5$, $4$.

\ex Ποιος είναι ο μεγαλύτερος από τους $-1{,}2$, $-1{,}02$, $-1{,}22$ και ποιος ο μικρότερος;

\ex Σωστό ή λάθος; α) $|\alpha|\ge0$ για κάθε $\alpha$. \quad β) $|-5|=-5$. \quad γ) Ο αντίθετος του $-4$ είναι $4$. \quad δ) $-9>-2$. \quad ε) Η $|x|=3$ έχει δύο λύσεις. \quad στ) Ο αντίθετος του $0$ είναι το $0$.

\zoni{Το χρησιμοποιώ}{θερμόμετρα, αποστάσεις, υψόμετρα}

\ex \textbf{Οι θερμοκρασίες.} Μέτσοβο $-5$, Ιωάννινα $-3$, Φλώρινα $-10$, Αθήνα $7$, Θεσσαλονίκη $-1$. Γράψε τις πόλεις από την πιο κρύα στην πιο ζεστή.

\ex \textbf{Οι αποστάσεις.} Στον άξονα της οδού το σχολείο είναι στο $0$ (μονάδα $100$ μ.), ο φούρνος στο $+3$ και το βιβλιοπωλείο στο $-2$. Ποιο είναι πιο κοντά στο σχολείο; Πόση είναι η απόσταση των δύο μαγαζιών μεταξύ τους;

\ex \textbf{Υψόμετρα.} Η Νεκρά Θάλασσα είναι στα $-430$ μ., ο Όλυμπος στα $+2\,917$ μ., η επιφάνεια της θάλασσας στο $0$ και η τάφρος των Μαριανών στα $-11\,034$ μ. Βάλε τα σε αύξουσα σειρά.

\ex \textbf{Το ταμείο.} Η Ζωή έχει $-15$ €, ο Ηλίας $-8$ € και η Ρήνη $+12$ € (αρνητικό σημαίνει ότι χρωστάει). Ποιος χρωστάει περισσότερο; Ποιος έχει το μεγαλύτερο ποσό;

\ex \textbf{Αντίθετες θέσεις.} Το σημείο $K$ έχει τετμημένη $-7$. Βρες το σημείο $\Lambda$ με αντίθετη τετμημένη και πες πόσο απέχουν τα $K$ και $\Lambda$.

\zoni{Το σκέφτομαι}{γιατί το μεγάλο νούμερο είναι μικρός αριθμός}

\ex Ο Βασίλης λέει ότι $-10>-3$, γιατί $10>3$. Εξήγησε το λάθος του με τον άξονα και με το θερμόμετρο.

\ex Γιατί ισχύει πάντα $|\alpha|=|-\alpha|$;

\ex Για ποιους αριθμούς ισχύει $|\alpha|=\alpha$ και για ποιους $|\alpha|=-\alpha$;

\ex Βρες δύο διαφορετικούς αριθμούς $x\neq y$ με $|x|=|y|$. Πόσοι ακέραιοι έχουν απόλυτη τιμή μικρότερη από $2$;

\ex Βρες έναν ρητό ανάμεσα στο $-2$ και στο $-1$ και δύο ρητούς ανάμεσα στο $-0{,}3$ και στο $-0{,}2$. Υπάρχουν άλλοι;

\zoni{Για το σπίτι}{ψάξε θερμοκρασίες}

\ex Βρες τις θερμοκρασίες πέντε πόλεων του κόσμου από τον καιρό και βάλε τες σε σειρά από την πιο κρύα στην πιο ζεστή.

\paragrafos{Πρόσθεση ρητών: κέρδη, ζημιές και ισορροπία}{Α.7.3, «Πρόσθεση ρητών αριθμών»}

\begin{anagke}
Στο πούλμαν η παρέα παίζει ένα παιχνίδι με κάρτες: κάθε κάρτα γράφει έναν πόντο, θετικό για κέρδος και αρνητικό για ζημιά. Ο Βασίλης τραβάει $+5$ και μετά $-3$. «Πόσους πόντους έχω;» ρωτάει. Ο Νίκος υπολογίζει: «Πέντε μείον τρία κάνει δύο». Η Πολυξένη τραβάει $-4$ και μετά $-2$ και ο Στέλιος λέει «έξι», μα η Μυρτώ τον ρωτάει: «Έξι κέρδος ή έξι ζημιά;». Ο Ηλίας κάνει μια πιο δύσκολη ερώτηση: «Αν τραβήξω $-7$ και μετά $+3$, είμαι ακόμα πίσω;».

Η Αγγελική το βλέπει στον άξονα: «Κάθε κάρτα είναι μια μετακίνηση. Το $+5$ σημαίνει «πέντε βήματα δεξιά», το $-3$ «τρία βήματα αριστερά». Ξεκινάω από το μηδέν, πηγαίνω πέντε δεξιά και μετά τρία αριστερά και καταλήγω στο $+2$». Για το $-4$ και το $-2$ είναι έξι βήματα αριστερά: $-6$. Και για τα $-7$ και $+3$ καταλήγουμε στο $-4$: ακόμα πίσω, αλλά λιγότερο.

Με αυτόν τον τρόπο η παρέα ανακαλύπτει τους κανόνες της πρόσθεσης ρητών και μαθαίνει γιατί ισχύουν. Μάλιστα στο τέλος του παιχνιδιού ο Νίκος παρατηρεί κάτι ωραίο: αν μια κάρτα $+3$ και μια $-3$ βρεθούν μαζί, «ακυρώνονται».
\end{anagke}

\begin{lexi}{πρόσθεση, προσθετέος, άθροισμα}{πρόσθεση, άθροισμα}
Η \lexiword{πρόσθεση} είναι το «προσ-τίθημι», δηλαδή «βάζω επιπλέον». Οι \lexiword{προσθετέοι} είναι αυτά που προσθέτουμε και το \lexiword{άθροισμα} είναι αυτό που βγαίνει όταν τα «αθροίζουμε», τα βάζουμε σε ένα σωρό. Στους ρητούς, το «προσθέτω» μπορεί να σημαίνει και «ζημιώνω».
\end{lexi}

\begin{metafrasi}
$5$ βήματα δεξιά και $3$ αριστερά. & $(+5)+(-3)=+2$\\
$4$ βήματα αριστερά και $2$ αριστερά ακόμα. & $(-4)+(-2)=-6$\\
Κέρδος και ζημιά που ακυρώνονται. & $(+3)+(-3)=0$\\
Δεν αλλάζει η σειρά των προσθετέων. & $\alpha+\beta=\beta+\alpha$\\
Το $0$ δεν αλλάζει τίποτα. & $\alpha+0=\alpha$\\
\end{metafrasi}

\tmima{Η θεωρία}

Στον άξονα, προσθέτω έναν θετικό όταν πηγαίνω προς τα δεξιά και έναν αρνητικό όταν πηγαίνω προς τα αριστερά, ξεκινώντας από τη θέση του πρώτου προσθετέου. Από αυτό προκύπτουν οι κανόνες:

\begin{kanonas}{πρόσθεση ρητών}
\begin{itemize}[leftmargin=1.2em,itemsep=0.3ex,topsep=2pt]
\item \textbf{Ομόσημοι:} προσθέτω τις απόλυτες τιμές και κρατάω το κοινό πρόσημο. $(-8)+(-5)=-(8+5)=-13$ και $(+9)+(+4)=+13$.
\item \textbf{Ετερόσημοι:} αφαιρώ τη μικρότερη απόλυτη τιμή από τη μεγαλύτερη και βάζω το πρόσημο του προσθετέου με τη μεγαλύτερη απόλυτη τιμή. $(+9)+(-4)=+5$ και $(-9)+(+4)=-5$.
\item \textbf{Αντίθετοι:} το άθροισμα είναι $0$. $(-3{,}5)+(+3{,}5)=0$.
\item Με το $0$: $\alpha+0=\alpha$.
\end{itemize}
\end{kanonas}

\begin{idiotita}{ιδιότητες της πρόσθεσης}
Για όλους τους ρητούς $\alpha,\,\beta,\,\gamma$ ισχύουν: \textbf{αντιμεταθετική} $\alpha+\beta=\beta+\alpha$, \textbf{προσεταιριστική} $(\alpha+\beta)+\gamma=\alpha+(\beta+\gamma)$, $\alpha+0=\alpha$ και $\alpha+(-\alpha)=0$.
\end{idiotita}

Χάρη στις ιδιότητες, όταν έχω πολλούς προσθετέους μπορώ να τους ομαδοποιήσω όπως θέλω: προσθέτω χωριστά όλους τους θετικούς και χωριστά όλους τους αρνητικούς και μετά τα δύο αποτελέσματα. Έτσι $(-7)+12+(-5)+9+(-3)=(12+9)+\big((-7)+(-5)+(-3)\big)=21+(-15)=6$.

Η πρόσθεση κλασμάτων και δεκαδικών ακολουθεί τους ίδιους κανόνες: $\left(-\frac{1}{2}\right)+\left(-\frac{1}{4}\right)=-\frac{3}{4}$ και $\left(+\frac{2}{3}\right)+\left(-\frac{1}{2}\right)=\frac{4}{6}-\frac{3}{6}=\frac{1}{6}$.

\begin{prosoxi}
\begin{itemize}[leftmargin=1.2em,itemsep=0.3ex,topsep=2pt]
\item Το άθροισμα δύο αριθμών δεν είναι πάντα μεγαλύτερο από τους προσθετέους: $(-3)+(-2)=-5<-3$.
\item Στους ετερόσημους δεν προσθέτω τις απόλυτες τιμές: $(+9)+(-4)\neq13$, αλλά $9-4=5$.
\item Το πρόσημο του αθροίσματος ετεροσήμων είναι του μεγαλύτερου κατ' απόλυτη τιμή: $(-9)+4=-5$, όχι $+5$.
\item Στο $5+(-3)$ οι παρενθέσεις δείχνουν ότι το $-3$ είναι ο αριθμός που προσθέτω, όχι δύο πράξεις.
\end{itemize}
\end{prosoxi}

\tmima{Λυμένα παραδείγματα}

\paradeigma{Πρόσθεση στον άξονα}
\tixero Οι προσθέσεις $(-2)+5$, $3+(-7)$, $(-4)+(-3)$.\par
\tithelo Να τις κάνω στον άξονα.\par
\possindeo Ξεκινάω από τον πρώτο προσθετέο και κινούμαι: από το $-2$ πέντε δεξιά: $3$. Από το $3$ επτά αριστερά: $-4$. Από το $-4$ τρία αριστερά: $-7$.\par
\apantisi $3$, $-4$, $-7$.\par
\elegxos Με τους κανόνες: $5-2=3$, $-(7-3)=-4$, $-(4+3)=-7$.

\paradeigma{Με τους κανόνες}
\tixero Οι αριθμοί $(-8)+(-5)$, $(+9)+(-4)$, $(-9)+(+4)$, $\left(-\frac{1}{2}\right)+\left(-\frac{1}{4}\right)$.\par
\tithelo Να υπολογίσω τα αθροίσματα.\par
\possindeo Ομόσημοι: $-(8+5)=-13$. Ετερόσημοι με μεγαλύτερη απόλυτη τιμή το $+9$: $+(9-4)=5$. Ετερόσημοι με μεγαλύτερη το $-9$: $-(9-4)=-5$. Κλάσματα ομόσημα: $-\left(\frac{2}{4}+\frac{1}{4}\right)=-\frac{3}{4}$.\par
\apantisi $-13$, $5$, $-5$, $-\frac{3}{4}$.\par
\elegxos Κάθε πρόσημο το βρίσκω από τον προσθετέο με τη μεγαλύτερη απόλυτη τιμή.

\paradeigma{Πολλοί προσθετέοι}
\tixero Το άθροισμα $(-7)+12+(-5)+9+(-3)$.\par
\tithelo Να το υπολογίσω.\par
\possindeo Ομαδοποιώ: θετικοί $12+9=21$. Αρνητικοί $(-7)+(-5)+(-3)=-15$. Τότε $21+(-15)=6$.\par
\apantisi $6$.\par
\elegxos Με τη σειρά: $-7+12=5$, $5-5=0$, $0+9=9$, $9-3=6$.

\paradeigma{Οι ιδιότητες βοηθούν}
\tixero Το άθροισμα $(-4{,}6)+(+3{,}9)+(+4{,}6)+(-3{,}9)$.\par
\tithelo Να το βρω εύκολα.\par
\possindeo Με την αντιμεταθετική και την προσεταιριστική ομαδοποιώ τους αντίθετους: $\big((-4{,}6)+(+4{,}6)\big)+\big((+3{,}9)+(-3{,}9)\big)=0+0=0$.\par
\apantisi $0$.\par
\elegxos Ο καθένας έχει τον αντίθετό του και ακυρώνεται.

\paradeigma{Το ταμείο της παρέας}
\tixero Το ταμείο έχει $+30$ €. Κινήσεις: $-12$, $+8$, $-35$, $+20$.\par
\tithelo Το υπόλοιπο.\par
\possindeo Θετικοί: $30+8+20=58$. Αρνητικοί: $-12-35=-47$. $58+(-47)=11$.\par
\apantisi $11$ €.\par
\elegxos Βήμα βήμα: $30,\ 18,\ 26,\ -9,\ 11$. Συμφωνεί.

\tmima{Ασκήσεις}

\zoni{Το κατάλαβα}{κανόνες, ιδιότητες, πολλοί προσθετέοι}

\ex Στον άξονα βρες τα αθροίσματα: α) $(-2)+5$ \quad β) $3+(-7)$ \quad γ) $(-4)+(-3)$ \quad δ) $6+(-6)$ \quad ε) $(-1)+(+1)$.

\ex Υπολόγισε: α) $(-8)+(-5)$ \quad β) $(+9)+(-4)$ \quad γ) $(-9)+(+4)$ \quad δ) $(+7)+(+5)$ \quad ε) $(-12)+0$ \quad στ) $(-2{,}5)+(+1{,}5)$.

\ex Υπολόγισε: α) $(-3{,}5)+(+3{,}5)$ \quad β) $(-0{,}7)+(-1{,}8)$ \quad γ) $(+2{,}4)+(-5{,}6)$ \quad δ) $\left(-\frac{1}{2}\right)+\left(-\frac{1}{4}\right)$ \quad ε) $\left(+\frac{2}{3}\right)+\left(-\frac{1}{2}\right)$ \quad στ) $\left(-\frac{5}{6}\right)+\left(+\frac{1}{3}\right)$.

\ex Συμπλήρωσε τον αριθμό που λείπει ώστε το άθροισμα να είναι $0$: α) $(-7)+\kouti$ \quad β) $(+3{,}5)+\kouti$ \quad γ) $\kouti+(-4)$ \quad δ) $\left(-\frac{1}{3}\right)+\kouti$.

\ex Υπολόγισε: α) $(-7)+12+(-5)+9+(-3)$ \quad β) $(-1{,}5)+4+(-2{,}5)+1$ \quad γ) $10+(-3)+(-4)+(-5)$.

\ex Υπολόγισε με τις ιδιότητες: α) $(-4{,}6)+(+3{,}9)+(+4{,}6)+(-3{,}9)$ \quad β) $37+(-18)+(-37)+8$.

\ex Σωστό ή λάθος; α) Το άθροισμα δύο αρνητικών είναι αρνητικό. \quad β) Το άθροισμα δύο ετεροσήμων είναι πάντα αρνητικό. \quad γ) $(-5)+5=0$. \quad δ) $(-3)+(-4)=+7$. \quad ε) $\alpha+0=\alpha$. \quad στ) Το άθροισμα δύο αντίθετων είναι $0$. \quad ζ) $(-2)+(+5)>0$.

\zoni{Το χρησιμοποιώ}{θερμόμετρα, ταμεία, παιχνίδια}

\ex \textbf{Η θερμοκρασία.} Το πρωί στο Μέτσοβο είναι $-6$°, το μεσημέρι ανεβαίνει $9$ βαθμούς και το βράδυ πέφτει $5$ βαθμούς. Πόση είναι η θερμοκρασία το μεσημέρι και το βράδυ;

\ex \textbf{Το ταμείο.} Το ταμείο έχει $+30$ €. Γίνονται οι κινήσεις $-12$, $+8$, $-35$, $+20$. Πόσο είναι το υπόλοιπο μετά από κάθε κίνηση και στο τέλος;

\ex \textbf{Ο ανελκυστήρας.} Ο Νίκος είναι στον όροφο $-2$, ανεβαίνει $5$ ορόφους και μετά κατεβαίνει $4$. Σε ποιον όροφο φτάνει;

\ex \textbf{Ο δύτης.} Ο δύτης είναι στα $-18$ μ., ανεβαίνει $7$ μ. και μετά κατεβαίνει $4$ μ. Πού βρίσκεται;

\ex \textbf{Το παιχνίδι.} Οι πόντοι της Πολυξένης στους τέσσερις γύρους είναι $+8$, $-5$, $-7$, $+12$. Πόσους έχει συνολικά;

\ex \textbf{Ο λογαριασμός.} Η Μυρτώ έχει υπόλοιπο $-45$ € στον λογαριασμό, καταθέτει $80$ € και μετά αναλαμβάνει $20$ €. Πόσα έχει;

\ex \textbf{Η διαφορά καλαθιών.} Η ομάδα μπάσκετ έχει στους πέντε αγώνες διαφορά $+12$, $-7$, $-9$, $+4$, $+6$. Ποια είναι η συνολική διαφορά;

\zoni{Το σκέφτομαι}{γιατί ισχύουν οι κανόνες}

\ex Πότε το άθροισμα δύο ρητών είναι $0$;

\ex Αν οι $\alpha,\,\beta$ είναι ετερόσημοι, ποιο πρόσημο έχει το άθροισμά τους;

\ex Γράψε δύο διαφορετικά ζεύγη ρητών με άθροισμα $-5$.

\ex Είναι πάντα το άθροισμα μεγαλύτερο από κάθε προσθετέο; Δώσε αντιπαράδειγμα.

\ex Δείξε με τα $\alpha=3$ και $\beta=4$ ότι $(-\alpha)+(-\beta)=-(\alpha+\beta)$.

\zoni{Για το σπίτι}{ο προϋπολογισμός σου}

\ex Κράτησε για μια εβδομάδα τα έσοδα και έξοδα του χαρτζιλικιού σου ως ρητούς και υπολόγισε το τελικό υπόλοιπο.

\ex Φτιάξε ένα παιχνίδι με κάρτες $+$ και $-$ και παίξε με έναν φίλο τρεις γύρους. Ποιος έχει το μεγαλύτερο άθροισμα;

\paragrafos{Αφαίρεση ρητών: η απόσταση ανάμεσα σε δύο θέσεις}{Α.7.4, «Αφαίρεση ρητών αριθμών»}

\begin{anagke}
Ο μετεωρολόγος στην τηλεόραση του ξενοδοχείου λέει: «Η μέγιστη θερμοκρασία σήμερα είναι $+5$ βαθμοί και η ελάχιστη $-3$. Η διαφορά είναι $8$ βαθμοί». Ο Στέλιος κάνει τον λογαριασμό: «Πέντε μείον τρία κάνει δύο. Ο μετεωρολόγος κάνει λάθος». Η Πολυξένη τον διορθώνει: «Δεν αφαιρούμε το τρία από το πέντε. Αφαιρούμε το μείον τρία!».

Η Αγγελική σχεδιάζει τον άξονα: «Από το $-3$ ως το $+5$ πόσα βήματα υπάρχουν;». Μετράνε: τρία βήματα ως το μηδέν και πέντε ακόμη. Οκτώ. Ο Νίκος συμπληρώνει: «Η αφαίρεση είναι αναποδογύρισμα της πρόσθεσης: ψάχνω αριθμό που προστιθέμενος στο $-3$ να δώσει $+5$. Δηλαδή $(-3)+x=5$. Αφού $(-3)+8=5$, το $5-(-3)$ είναι $8$».

Ο Ηλίας το γενικεύει: «Αν η αφαίρεση ψάχνει «πόσο πρέπει να προσθέσω», μήπως $\alpha-\beta$ είναι το ίδιο με το να προσθέσω τον αντίθετο του $\beta$;». Έτσι η παρέα ανακαλύπτει ότι κάθε αφαίρεση μπορεί να γίνει πρόσθεση και ότι όλα όσα ξέρουν για την πρόσθεση λειτουργούν και για την αφαίρεση.
\end{anagke}

\begin{lexi}{αφαίρεση, διαφορά}{αφαίρεση, διαφορά}
Η \lexiword{αφαίρεση} είναι το «αφ-αιρώ», «βγάζω από κάτι». Η \lexiword{διαφορά} είναι αυτό που μένει και δείχνει «πόσο διαφέρουν» δύο αριθμοί. Ο \lexiword{μειωτέος} είναι αυτός που «μειώνεται» και ο \lexiword{αφαιρετέος} αυτός που αφαιρείται.
\end{lexi}

\begin{metafrasi}
Η διαφορά $\alpha-\beta$ είναι ο αριθμός $x$ που... & $\beta+x=\alpha$\\
Αφαίρεση σημαίνει πρόσθεση του αντίθετου. & $\alpha-\beta=\alpha+(-\beta)$\\
Η διαφορά θερμοκρασιών $+5$ και $-3$. & $(+5)-(-3)=8$\\
Αφαιρώ ολόκληρη παρένθεση. & $\alpha-(\beta+\gamma)=\alpha-\beta-\gamma$\\
\end{metafrasi}

\tmima{Η θεωρία}

\begin{orismos}{διαφορά}
\textbf{Διαφορά} $\alpha-\beta$ δύο ρητών ορίζεται ως η μοναδική λύση της εξίσωσης $\beta+x=\alpha$.
\end{orismos}

Για παράδειγμα η διαφορά $(+5)-(-3)$ είναι η λύση της $(-3)+x=5$, δηλαδή $x=8$. Από τον ορισμό προκύπτει ο κανόνας:

\begin{kanonas}{αφαίρεση ρητών}
Για να αφαιρέσω έναν αριθμό, προσθέτω τον \textbf{αντίθετό} του: $$\alpha-\beta=\alpha+(-\beta).$$ Έτσι $7-(+9)=7+(-9)=-2$ και $(-4)-(-6)=(-4)+(+6)=2$.
\end{kanonas}

Η διαφορά $\alpha-\beta$ δείχνει και τον δρόμο από το $\beta$ στο $\alpha$ στον άξονα: αν $\alpha>\beta$ είναι θετική, αλλιώς αρνητική.

\tmimasmall{Αλγεβρικό άθροισμα και απαλοιφή παρενθέσεων}

Μια παράσταση με προσθέσεις και αφαιρέσεις γράφεται ως πρόσθεση ρητών. Για παράδειγμα $12-15+8-20+3=12+(-15)+8+(-20)+3$, και μπορώ να τα αθροίσω με όποια σειρά θέλω: $(12+8+3)+(-15-20)=23-35=-12$. Τέτοιο άθροισμα ονομάζεται \textbf{αλγεβρικό άθροισμα}. Για να σβήσω παρενθέσεις:
\begin{itemize}[leftmargin=1.2em,itemsep=0.3ex,topsep=2pt]
\item Αν μπροστά από την παρένθεση υπάρχει $+$, τη σβήνω χωρίς να αλλάξω τα πρόσημα: $\alpha+(\beta-\gamma)=\alpha+\beta-\gamma$.
\item Αν μπροστά υπάρχει $-$, τη σβήνω και \textbf{αλλάζω το πρόσημο όλων} των όρων μέσα: $\alpha-(\beta-\gamma+\delta)=\alpha-\beta+\gamma-\delta$.
\end{itemize}

\begin{prosoxi}
\begin{itemize}[leftmargin=1.2em,itemsep=0.3ex,topsep=2pt]
\item Η διαφορά δεν είναι αντιμεταθετική: $5-3=2$ αλλά $3-5=-2$. Δεν ισχύει $\alpha-\beta=\beta-\alpha$.
\item Η αφαίρεση δεν «μικραίνει» πάντα: $3-(-2)=5$, μεγαλύτερο από το $3$.
\item Με το $-$ μπροστά από παρένθεση αλλάζω το πρόσημο \emph{όλων} των όρων της παρένθεσης, όχι μόνο του πρώτου: $10-(3-8+2)=10-3+8-2$.
\item Η διαφορά θερμοκρασιών $+5$ και $-3$ είναι $8$, όχι $2$. Αφαιρώ τον αριθμό μαζί με το πρόσημό του.
\end{itemize}
\end{prosoxi}

\tmima{Λυμένα παραδείγματα}

\paradeigma{Η διαφορά των θερμοκρασιών}
\tixero Μέγιστη $+5$°, ελάχιστη $-3$°.\par
\tithelo Η θερμική διακύμανση, δηλαδή η διαφορά μέγιστης και ελάχιστης.\par
\possindeo $(+5)-(-3)=(+5)+(+3)=8$. Έλεγχος στον άξονα: από το $-3$ ως το $+5$ υπάρχουν $8$ μονάδες.\par
\apantisi $8$ βαθμοί.\par
\elegxos $(-3)+8=5$.

\paradeigma{Με τον κανόνα}
\tixero Οι διαφορές $7-(+9)$, $(-4)-(-6)$, $(+3{,}5)-(-1{,}5)$, $\left(-\frac{1}{2}\right)-\left(+\frac{1}{3}\right)$.\par
\tithelo Να τις υπολογίσω.\par
\possindeo Προσθέτω τον αντίθετο: $7+(-9)=-2$. $(-4)+(+6)=2$. $3{,}5+1{,}5=5$. $-\frac{1}{2}-\frac{1}{3}=-\frac{3}{6}-\frac{2}{6}=-\frac{5}{6}$.\par
\apantisi $-2$, $2$, $5$, $-\frac{5}{6}$.\par
\elegxos Για την πρώτη: $9+(-2)=7$. Σωστό.

\paradeigma{Απαλοιφή παρενθέσεων}
\tixero Η παράσταση $10-(3-8+2)$.\par
\tithelo Να την υπολογίσω με δύο τρόπους.\par
\possindeo \emph{1ος τρόπος:} πρώτα η παρένθεση: $3-8+2=-3$, και $10-(-3)=13$. \emph{2ος τρόπος:} σβήνω την παρένθεση αλλάζοντας πρόσημα: $10-3+8-2=13$.\par
\apantisi $13$.\par
\elegxos Και οι δύο δίνουν $13$.

\paradeigma{Αλγεβρικό άθροισμα}
\tixero Η παράσταση $12-15+8-20+3$.\par
\tithelo Να την υπολογίσω.\par
\possindeo Ομαδοποιώ: θετικοί $12+8+3=23$. Αρνητικοί $-15-20=-35$. $23+(-35)=-12$.\par
\apantisi $-12$.\par
\elegxos Με τη σειρά: $12-15=-3$, $-3+8=5$, $5-20=-15$, $-15+3=-12$.

\paradeigma{Εξισώσεις με αφαίρεση}
\tixero Οι εξισώσεις $x+(-3)=5$, $x-7=-2$ και $4-x=9$.\par
\tithelo Να τις λύσω.\par
\possindeo $x=5-(-3)=8$. $x=-2+7=5$. $x=4-9=-5$ (από την $\alpha-x=\beta\Rightarrow x=\alpha-\beta$).\par
\apantisi $8$, $5$, $-5$.\par
\elegxos $8+(-3)=5$, $5-7=-2$, $4-(-5)=9$.

\tmima{Ασκήσεις}

\zoni{Το κατάλαβα}{διαφορές, παρενθέσεις, αλγεβρικά αθροίσματα}

\ex Υπολόγισε: α) $(+5)-(+3)$ \quad β) $(+3)-(+5)$ \quad γ) $(-4)-(-1)$ \quad δ) $(-1)-(-4)$ \quad ε) $0-(+6)$ \quad στ) $0-(-6)$.

\ex Μετάτρεψε σε πρόσθεση και υπολόγισε: α) $7-(+9)$ \quad β) $(-4)-(-6)$ \quad γ) $(+3{,}5)-(-1{,}5)$ \quad δ) $\left(-\frac{1}{2}\right)-\left(+\frac{1}{3}\right)$ \quad ε) $(-2{,}4)-(+0{,}6)$ \quad στ) $8-(-8)$.

\ex Λύσε: α) $x+7=3$ \quad β) $x-5=-2$ \quad γ) $4-x=9$ \quad δ) $x+(-6)=-6$ \quad ε) $-3-x=2$.

\ex Κάνε απαλοιφή παρενθέσεων και υπολόγισε: α) $10-(3-8+2)$ \quad β) $5+(-3+7)-(2-9)$ \quad γ) $-(-4+5)-(7-10)$ \quad δ) $20-(-5)-(-(3))$.

\ex Υπολόγισε τα αλγεβρικά αθροίσματα: α) $12-15+8-20+3$ \quad β) $-2{,}5+7-4{,}5-1$ \quad γ) $\frac{1}{2}-\frac{3}{4}+\frac{1}{4}-1$.

\ex Σωστό ή λάθος; α) $\alpha-\beta=\beta-\alpha$. \quad β) $\alpha-0=\alpha$. \quad γ) $0-\alpha=-\alpha$. \quad δ) $(-3)-(-3)=0$. \quad ε) $(+4)-(-2)=2$. \quad στ) Η διαφορά δύο θετικών μπορεί να είναι αρνητική.

\zoni{Το χρησιμοποιώ}{θερμοκρασίες, υψόμετρα, χρόνια}

\ex \textbf{Η διακύμανση.} Βρες τη διαφορά μέγιστης και ελάχιστης θερμοκρασίας όταν α) μέγιστη $+5$, ελάχιστη $-3$ \quad β) μέγιστη $-2$, ελάχιστη $-11$.

\ex \textbf{Το υψόμετρο.} Ο Όλυμπος είναι στα $+2\,917$ μ. και η Νεκρά Θάλασσα στα $-430$ μ. Πόση είναι η υψομετρική διαφορά;

\ex \textbf{Η ιστορία.} Η δυναστεία των Χαν ξεκίνησε το $-206$ και ο Διόφαντος έζησε γύρω στο $250$ μ.Χ. Πόσα χρόνια μεσολαβούν (αγνοώντας ότι δεν υπήρξε έτος $0$);

\ex \textbf{Το ταμείο.} Η Ζωή έχει $-15$ € και ο Ηλίας $+8$ €. Πόση είναι η διαφορά των δύο ποσών;

\ex \textbf{Ο δύτης.} Ο δύτης ανεβαίνει από το $-22$ μ. στο $-9$ μ. Πόσα μέτρα ανέβηκε;

\ex \textbf{Ο ανελκυστήρας.} Από τον όροφο $3$ στον $-2$ η μεταβολή είναι $(-2)-3$. Πόσο είναι και τι σημαίνει;

\ex \textbf{Η πτώση.} Η θερμοκρασία από $+4$° έπεσε στα $-7$°. Πόση είναι η μεταβολή (νέα $-$ παλιά);

\zoni{Το σκέφτομαι}{όταν η αφαίρεση μεγαλώνει}

\ex Ο Στέλιος γράφει $5-(-3)=2$. Βρες το λάθος και δείξε το σωστό με τον άξονα.

\ex Μικραίνει πάντα ένας αριθμός όταν του αφαιρώ κάτι; Δώσε παράδειγμα.

\ex Αν $\alpha>\beta$, ποιο πρόσημο έχει η διαφορά $\alpha-\beta$; Και αν $\alpha<\beta$;

\ex Δείξε με $\alpha=10$, $\beta=3$, $\gamma=4$ ότι $\alpha-(\beta+\gamma)=\alpha-\beta-\gamma$.

\ex Βρες τον αριθμό που, αν του αφαιρέσουμε το $-5$, δίνει $2$.

\zoni{Για το σπίτι}{διαφορές του κόσμου}

\ex Βρες τη διαφορά θερμοκρασίας ανάμεσα στο πιο ζεστό και στο πιο κρύο μέρος του κόσμου σήμερα.

\ex Φτιάξε μια παράσταση με έξι όρους και παρενθέσεις και υπολόγισέ την με δύο τρόπους.

\paragrafos{Πολλαπλασιασμός ρητών: το μείον επί μείον}{Α.7.5, «Πολλαπλασιασμός ρητών αριθμών»}

\begin{anagke}
Το βράδυ στο Μέτσοβο η θερμοκρασία πέφτει $2$ βαθμούς κάθε ώρα. Η Πολυξένη παρατηρεί το θερμόμετρο στις $20{:}00$, που δείχνει $0$°. «Τι θα δείχνει σε τρεις ώρες;». Ο Νίκος σκέφτεται: «Κάθε ώρα $-2$, τρεις ώρες: $-2-2-2=-6$». Η Αγγελική το γράφει ως πολλαπλασιασμό: $3\cdot(-2)=-6$. «Και πριν από τρεις ώρες, στις $17{:}00$;» ρωτάει η Μυρτώ. «Αφού τώρα πέφτει δύο βαθμούς ανά ώρα, πριν από τρεις ώρες είχε $6$ βαθμούς περισσότερο», απαντάει η Πολυξένη. Ο ρυθμός είναι $-2$ και ο χρόνος «τρεις ώρες πριν» είναι $-3$, και το αποτέλεσμα είναι $+6$. Άρα $(-2)\cdot(-3)=+6$.

Ο Στέλιος ξαφνιάζεται: «Το μείον επί μείον κάνει συν; Σαν τη διπλή άρνηση της γλώσσας που καταλήγει σε κατάφαση;». Η κυρία Αλεξάνδρα του απαντά: «Η γλώσσα είναι ωραία αναλογία, αλλά τα μαθηματικά θέλουν να δουν γιατί». Το ίδιο το μοτίβο δείχνει το γιατί: $3\cdot(-2)=-6$, $2\cdot(-2)=-4$, $1\cdot(-2)=-2$, $0\cdot(-2)=0$. Κάθε φορά που το πρώτο μικραίνει κατά $1$ το γινόμενο μεγαλώνει κατά $2$. Συνεχίζοντας: $(-1)\cdot(-2)=2$, $(-2)\cdot(-2)=4$.
\end{anagke}

\begin{lexi}{παράγοντας, γινόμενο, ομόσημοι}{παράγοντας, γινόμενο}
\lexiword{Παράγοντες} είναι οι αριθμοί που πολλαπλασιάζονται και \lexiword{γινόμενο} αυτό που «γίνεται». Η \lexiword{επιμεριστική} ιδιότητα είναι αυτή που «επιμερίζει» (μοιράζει σε μέρη) τον πολλαπλασιασμό πάνω σε ένα άθροισμα.
\end{lexi}

\begin{metafrasi}
Πέφτει $2$ βαθμούς την ώρα, για $3$ ώρες. & $3\cdot(-2)=-6$\\
Πριν από $3$ ώρες. & $(-2)\cdot(-3)=+6$\\
Ετερόσημοι. & γινόμενο αρνητικό\\
Ομόσημοι. & γινόμενο θετικό\\
Ο πολλαπλασιασμός πάνω σε άθροισμα. & $\alpha(\beta+\gamma)=\alpha\beta+\alpha\gamma$\\
\end{metafrasi}

\tmima{Η θεωρία}

\begin{kanonas}{γινόμενο ρητών}
Το γινόμενο δύο ρητών έχει απόλυτη τιμή το γινόμενο των απόλυτων τιμών τους και:
\begin{itemize}[leftmargin=1.2em,itemsep=0.3ex,topsep=2pt]
\item πρόσημο \textbf{θετικό} αν οι παράγοντες είναι ομόσημοι: $(-7)\cdot(-8)=+56$, $(+5)(+4)=+20$.
\item πρόσημο \textbf{αρνητικό} αν οι παράγοντες είναι ετερόσημοι: $(-6)\cdot(+4)=-24$.
\item Αν ένας παράγοντας είναι $0$, το γινόμενο είναι $0$.
\end{itemize}
\end{kanonas}

Για παράδειγμα $\left(-\frac{3}{4}\right)\cdot\left(-\frac{8}{9}\right)=+\frac{3\cdot8}{4\cdot9}=\frac{2}{3}$ και $(+0{,}5)\cdot(-12)=-6$.

\begin{idiotita}{ιδιότητες του πολλαπλασιασμού}
Για κάθε τρεις ρητούς ισχύουν: \textbf{αντιμεταθετική} $\alpha\beta=\beta\alpha$, \textbf{προσεταιριστική} $(\alpha\beta)\gamma=\alpha(\beta\gamma)$, $\alpha\cdot1=\alpha$, $\alpha\cdot0=0$, $\alpha\cdot(-1)=-\alpha$ και \textbf{επιμεριστική} $\alpha(\beta+\gamma)=\alpha\beta+\alpha\gamma$ (και $\alpha(\beta-\gamma)=\alpha\beta-\alpha\gamma$).
\end{idiotita}

\tmimasmall{Γινόμενο πολλών παραγόντων}

Το γινόμενο πολλών ρητών (όλων μη μηδενικών) είναι θετικό αν το πλήθος των αρνητικών παραγόντων είναι \emph{άρτιο} και αρνητικό αν είναι \emph{περιττό}. Η απόλυτη τιμή του είναι το γινόμενο των απόλυτων τιμών. Έτσι $(-2)\cdot3\cdot(-5)\cdot(-1)=-30$ (τρεις αρνητικοί).

Η επιμεριστική ιδιότητα δίνει συντομίες: $99\cdot(-6)=(100-1)\cdot(-6)=-600+6=-594$.

\begin{prosoxi}
\begin{itemize}[leftmargin=1.2em,itemsep=0.3ex,topsep=2pt]
\item Στον πολλαπλασιασμό, δύο αρνητικοί δίνουν θετικό, ενώ στην πρόσθεση δύο αρνητικοί δίνουν αρνητικό: $(-3)\cdot(-4)=+12$ και $(-3)+(-4)=-7$.
\item Το πρόσημο του γινομένου μετράει το πλήθος των αρνητικών παραγόντων, όχι το πλήθος των θετικών.
\item Στην επιμεριστική πολλαπλασιάζω τον αριθμό με \emph{κάθε} όρο της παρένθεσης: $(-2)(5-9)=-10+18$.
\item $(-3)^2$ και $-3^2$ δεν είναι ίδια, αλλά θα το δούμε στις δυνάμεις.
\end{itemize}
\end{prosoxi}

\tmima{Λυμένα παραδείγματα}

\paradeigma{Το μοτίβο}
\tixero Τα γινόμενα $3\cdot(-4)$, $2\cdot(-4)$, $1\cdot(-4)$, $0\cdot(-4)$.\par
\tithelo Να συνεχίσω το μοτίβο.\par
\possindeo Τα γινόμενα είναι $-12$, $-8$, $-4$, $0$ και αυξάνουν κατά $4$ κάθε φορά που ο πρώτος παράγοντας μειώνεται κατά $1$. Άρα $(-1)\cdot(-4)=4$, $(-2)\cdot(-4)=8$, $(-3)\cdot(-4)=12$.\par
\apantisi $4$, $8$, $12$.\par
\elegxos Ο κανόνας ομόσημοι δίνει θετικό: $(-3)\cdot(-4)=+12$.

\paradeigma{Γινόμενα με τον κανόνα}
\tixero Τα $(-6)(+4)$, $(-7)(-8)$, $(+0{,}5)(-12)$, $\left(-\frac{3}{4}\right)\left(-\frac{8}{9}\right)$.\par
\tithelo Να τα υπολογίσω.\par
\possindeo Πρώτα το πρόσημο, μετά οι απόλυτες τιμές. Ετερόσημοι: $-24$. Ομόσημοι: $+56$. Ετερόσημοι: $-6$. Ομόσημοι: $+\frac{24}{36}=\frac{2}{3}$.\par
\apantisi $-24$, $56$, $-6$, $\frac{2}{3}$.\par
\elegxos Το $7\cdot8=56$ και $6\cdot4=24$.

\paradeigma{Πολλοί παράγοντες}
\tixero Το γινόμενο $(-2)\cdot3\cdot(-5)\cdot(-1)$.\par
\tithelo Να βρω το πρόσημο και την τιμή.\par
\possindeo Αρνητικοί παράγοντες: $3$ (περιττό πλήθος), άρα αρνητικό. Απόλυτη τιμή: $2\cdot3\cdot5\cdot1=30$.\par
\apantisi $-30$.\par
\elegxos Βήμα βήμα: $(-2)\cdot3=-6$, $-6\cdot(-5)=30$, $30\cdot(-1)=-30$.

\paradeigma{Η επιμεριστική}
\tixero Οι $(-4)\cdot(7+3)$ και $99\cdot(-6)$.\par
\tithelo Να τις υπολογίσω εύκολα.\par
\possindeo $(-4)\cdot7+(-4)\cdot3=-28-12=-40$ (ή $-4\cdot10$). Για το δεύτερο: $99=100-1$, $(100-1)\cdot(-6)=-600+6=-594$.\par
\apantisi $-40$ και $-594$.\par
\elegxos $99\cdot6=594$ και ο αριθμός είναι αρνητικός.

\paradeigma{Παράσταση}
\tixero Η παράσταση $3\cdot(-4)-(-5)\cdot2+(-2)\cdot(-3)$.\par
\tithelo Να την υπολογίσω.\par
\possindeo Πρώτα οι πολλαπλασιασμοί: $3\cdot(-4)=-12$, $(-5)\cdot2=-10$, $(-2)(-3)=6$. Άρα $-12-(-10)+6=-12+10+6=4$.\par
\apantisi $4$.\par
\elegxos Προτεραιότητα: πολλαπλασιασμοί πριν από προσθέσεις και αφαιρέσεις.

\tmima{Ασκήσεις}

\zoni{Το κατάλαβα}{πρόσημα και τιμές}

\ex Συμπλήρωσε το μοτίβο: $3\cdot(-4)$, $2\cdot(-4)$, $1\cdot(-4)$, $0\cdot(-4)$, $(-1)\cdot(-4)$, $(-2)\cdot(-4)$, $(-3)\cdot(-4)$.

\ex Υπολόγισε: α) $(-6)\cdot(+4)$ \quad β) $(+7)\cdot(-3)$ \quad γ) $(-5)\cdot(-9)$ \quad δ) $(-12)\cdot0$ \quad ε) $(+8)\cdot(+5)$ \quad στ) $(-1)\cdot(-1)$.

\ex Υπολόγισε: α) $(+0{,}5)\cdot(-12)$ \quad β) $(-1{,}2)\cdot(-3)$ \quad γ) $(-2{,}5)\cdot(+0{,}4)$ \quad δ) $\left(-\frac{3}{4}\right)\cdot\left(-\frac{8}{9}\right)$ \quad ε) $\left(+\frac{2}{5}\right)\cdot(-15)$ \quad στ) $\left(-\frac{1}{2}\right)\cdot\left(-\frac{2}{3}\right)\cdot(-3)$.

\ex Βρες το πρόσημο του γινομένου χωρίς να το υπολογίσεις: α) $(-2)\cdot3\cdot(-5)\cdot(-1)$ \quad β) $(-1)(-1)(-1)(-1)$ \quad γ) $7\cdot(-3)\cdot(-2)\cdot4\cdot(-1)$ \quad δ) $(-5)\cdot0\cdot(-8)$.

\ex Υπολόγισε: α) $(-8)\cdot1$ \quad β) $(-8)\cdot(-1)$ \quad γ) $0\cdot(-3)$ \quad δ) $\alpha\cdot(-1)$ για $\alpha=-7$.

\ex Με την επιμεριστική: α) $(-4)\cdot(7+3)$ \quad β) $5\cdot(-3+8)$ \quad γ) $(-2)\cdot(5-9)$ \quad δ) $99\cdot(-6)$ \quad ε) $(-12)\cdot\left(\frac{1}{3}+\frac{1}{4}\right)$.

\ex Υπολόγισε: α) $3\cdot(-4)-(-5)\cdot2+(-2)\cdot(-3)$ \quad β) $(-2)\cdot(4-7)+3\cdot(-1-2)$ \quad γ) $10-2\cdot(-3)\cdot(-5)$.

\ex Σωστό ή λάθος; α) $(-3)(-4)=-12$. \quad β) Το γινόμενο ετεροσήμων είναι αρνητικό. \quad γ) $\alpha\cdot0=0$. \quad δ) $(-1)\cdot\alpha=-\alpha$. \quad ε) Το γινόμενο τριών αρνητικών είναι θετικό. \quad στ) $(-2)(-3)(-4)=-24$.

\zoni{Το χρησιμοποιώ}{θερμοκρασίες, χρέη, βουτιές}

\ex \textbf{Η θερμοκρασία.} Πέφτει $2$ βαθμούς κάθε ώρα. Πόσο θα έχει μεταβληθεί σε $3$ ώρες και πόσο έχει μεταβληθεί από πριν $3$ ώρες;

\ex \textbf{Το χρέος.} Κάθε εβδομάδα το χρέος της Μυρτώς αυξάνεται κατά $4$ €, δηλαδή η μεταβολή είναι $-4$ €. Ποια είναι η μεταβολή σε $5$ εβδομάδες;

\ex \textbf{Ο δύτης.} Ο δύτης κατεβαίνει $3$ μέτρα το λεπτό. Πού θα είναι σε $8$ λεπτά; Πού ήταν πριν από $5$ λεπτά;

\ex \textbf{Ο ανελκυστήρας.} Ο ανελκυστήρας κατεβαίνει $2$ ορόφους κάθε φορά, $4$ φορές. Πόση είναι η συνολική μεταβολή;

\ex \textbf{Το ταμείο.} Αν τα $9$ παιδιά της παρέας χρωστάνε από $4$ € το καθένα, τι αριθμός εκφράζει το συνολικό χρέος;

\ex \textbf{Η μετοχή.} Η τιμή μιας μετοχής πέφτει $0{,}25$ € κάθε ημέρα. Πόση είναι η συνολική μεταβολή σε $6$ ημέρες;

\zoni{Το σκέφτομαι}{γιατί το μείον επί μείον κάνει συν}

\ex Εξήγησε με το μοτίβο ή με τη θερμοκρασία γιατί $(-2)\cdot(-3)=+6$.

\ex Γιατί το τετράγωνο ενός ρητού δεν είναι ποτέ αρνητικό;

\ex Πόσα αρνητικά πρέπει να έχει ένα γινόμενο $10$ παραγόντων για να είναι θετικό; Δώσε δύο διαφορετικά πλήθη.

\ex Αν $\alpha\cdot\beta=0$, τι συμπεραίνεις για τους $\alpha,\,\beta$;

\ex Υπολόγισε το $(-1)\cdot(-1)\cdot\ldots\cdot(-1)$ με $100$ παράγοντες και με $101$ παράγοντες.

\zoni{Για το σπίτι}{φτιάξε το μοτίβο}

\ex Φτιάξε το μοτίβο του πολλαπλασιασμού με το $-5$ από $3$ ως $-3$ και πες τι παρατηρείς.

\paragrafos{Διαίρεση ρητών: μοιράζω και τα χρέη}{Α.7.6, «Διαίρεση ρητών αριθμών»}

\begin{anagke}
Στο τέλος της εκδρομής ο Νίκος, που κρατάει το ταμείο, ανακοινώνει με σκυμμένο κεφάλι: «Το ταμείο έμεινε στο $-24$ ευρώ. Χρωστάμε $24$ ευρώ». Η Μυρτώ ρωτάει: «Αν το μοιράσουμε εξίσου στα έξι παιδιά της παρέας, πόσο ήταν το καθένα;». Η Πολυξένη το γράφει: $(-24):6$. Ο Βασίλης λέει «τέσσερα» και ο Στέλιος διαμαρτύρεται: «Πόσο, χρέος ή κέρδος;». Αφού το χρέος μοιράζεται, κάθε παιδί χρωστά $4$ ευρώ, δηλαδή το πηλίκο είναι $-4$.

Ο Ηλίας ρωτάει το αντίστροφο: «Αν κάθε παιδί πρέπει να βάλει $3$ ευρώ ($-3$ στο ταμείο του) και το σύνολο είναι $-24$, πόσα παιδιά είναι;». Ο Νίκος σκέφτεται: «Πόσες φορές χωράει το $-3$ στο $-24$; Οκτώ φορές». Άρα $(-24):(-3)=+8$. Η Αγγελική τα βάζει σε εξίσωση: $(-3)\cdot x=-24$ και $x=8$. Και η Πολυξένη θυμάται ότι η διαίρεση είναι ο αντίστροφος δρόμος του πολλαπλασιασμού, όπως η αφαίρεση της πρόσθεσης.

Από τον πολλαπλασιασμό ήδη ξέρουν τους κανόνες των προσήμων. Τους ίδιους κανόνες θα βρουν και για τη διαίρεση, και θα δουν γιατί, πάλι, δεν μπορούν να διαιρέσουν με το $0$.
\end{anagke}

\begin{lexi}{διαίρεση, πηλίκο, αντίστροφος}{διαίρεση, πηλίκο, λόγος}
Η \lexiword{διαίρεση} είναι το «δι-αιρώ», δηλαδή «χωρίζω σε μέρη». Ο \lexiword{διαιρετέος} είναι αυτός που μοιράζεται, ο \lexiword{διαιρέτης} αυτός που μοιράζει και το \lexiword{πηλίκο} (από το «πηλίκος», «πόσος») το αποτέλεσμα. Ο \lexiword{λόγος} δύο ρητών είναι το πηλίκο τους, γραμμένο ως κλάσμα.
\end{lexi}

\begin{metafrasi}
Η διαίρεση $\alpha:\beta$ είναι η λύση της εξίσωσης. & $\beta\cdot x=\alpha$ (με $\beta\neq0$)\\
Μοιράζω χρέος $24$ € σε $6$ παιδιά. & $(-24):6=-4$\\
Πόσες φορές χωράει το $-3$ στο $-24$; & $(-24):(-3)=+8$\\
Διαίρεση σημαίνει πολλαπλασιασμός με τον αντίστροφο. & $\alpha:\beta=\alpha\cdot\dfrac{1}{\beta}$\\
Κλάσμα με αρνητικό. & $\dfrac{-6}{3}=\dfrac{6}{-3}=-2$\\
\end{metafrasi}

\tmima{Η θεωρία}

\begin{orismos}{πηλίκο}
\textbf{Πηλίκο} $\alpha:\beta$ δύο ρητών με $\beta\neq0$ ορίζεται ως η μοναδική λύση της εξίσωσης $\beta\cdot x=\alpha$.
\end{orismos}

Για παράδειγμα το $12:(-3)$ είναι η λύση της $(-3)\cdot x=12$, δηλαδή $x=-4$. Από αυτό προκύπτει ότι τα πρόσημα είναι τα ίδια με του πολλαπλασιασμού:

\begin{kanonas}{διαίρεση ρητών}
Το πηλίκο δύο ρητών έχει απόλυτη τιμή το πηλίκο των απόλυτων τιμών και \textbf{πρόσημο} θετικό αν είναι ομόσημοι και αρνητικό αν είναι ετερόσημοι. Δηλαδή $\alpha:\beta$ και $\alpha\cdot\beta$ είναι \textbf{ομόσημοι}. Το $0:\beta=0$. Δεν διαιρούμε με το $0$.
\end{kanonas}

Έτσι $(-24):6=-4$, $(-24):(-3)=+8$, $18:(-6)=-3$, $0:(-5)=0$.

\tmimasmall{Διαίρεση και αντίστροφος}

Ο \textbf{αντίστροφος} ενός ρητού $\alpha\neq0$ είναι ο $\frac{1}{\alpha}$ (γινόμενο $1$). Ο αντίστροφος του $-4$ είναι ο $-\frac{1}{4}$, του $-\frac{2}{3}$ ο $-\frac{3}{2}$. Το $0$ δεν έχει αντίστροφο. Κάθε διαίρεση γίνεται πολλαπλασιασμός: $\alpha:\beta=\alpha\cdot\frac{1}{\beta}$, όπως και στα κλάσματα. Έτσι $\left(-\frac{3}{4}\right):\frac{1}{2}=\left(-\frac{3}{4}\right)\cdot2=-\frac{3}{2}$.

\tmimasmall{Το πηλίκο ως λόγος}

Το πηλίκο δύο ρητών γράφεται ως κλάσμα (\emph{λόγος}): $\dfrac{-6}{3}=-2$, $\dfrac{6}{-3}=-2$, $\dfrac{-6}{-3}=+2$. Ισχύει $\dfrac{-\alpha}{\beta}=\dfrac{\alpha}{-\beta}=-\dfrac{\alpha}{\beta}$. Τα σύνθετα κλάσματα λύνονται όπως πριν: $\dfrac{-\frac{6}{5}}{\frac{3}{10}}=-\dfrac{6}{5}\cdot\dfrac{10}{3}=-4$.

\begin{prosoxi}
\begin{itemize}[leftmargin=1.2em,itemsep=0.3ex,topsep=2pt]
\item Δεν υπάρχει $\alpha:0$: η $0\cdot x=\alpha$ δεν έχει λύση όταν $\alpha\neq0$ και έχει άπειρες όταν $\alpha=0$.
\item Το $(-12):4$ είναι $-3$, όχι $3$: ετερόσημοι δίνουν αρνητικό.
\item Η διαίρεση δεν είναι αντιμεταθετική: $(-6):3=-2$ αλλά $3:(-6)=-\frac{1}{2}$.
\item Αντίστροφος δεν είναι αντίθετος: ο αντίθετος του $4$ είναι $-4$ και ο αντίστροφος $\frac{1}{4}$.
\end{itemize}
\end{prosoxi}

\tmima{Λυμένα παραδείγματα}

\paradeigma{Το ταμείο μοιράζεται}
\tixero Το ταμείο είναι $-24$ € και μοιράζεται σε $6$ παιδιά.\par
\tithelo Το μερίδιο του καθενός και πόσα παιδιά βάζουν από $-3$.\par
\possindeo $(-24):6=-4$ (ετερόσημοι, αρνητικό). $(-24):(-3)=+8$ (ομόσημοι, θετικό).\par
\apantisi $-4$ € και $8$ παιδιά.\par
\elegxos $6\cdot(-4)=-24$ και $8\cdot(-3)=-24$.

\paradeigma{Με τον ορισμό}
\tixero Το πηλίκο $12:(-3)$.\par
\tithelo Να το βρω ως λύση εξίσωσης.\par
\possindeo Ψάχνω $x$ με $(-3)\cdot x=12$. Αφού $(-3)\cdot(-4)=12$, είναι $x=-4$.\par
\apantisi $-4$.\par
\elegxos Ο κανόνας: ετερόσημοι, αρνητικό, $12:3=4$.

\paradeigma{Πηλίκα ρητών}
\tixero Τα $(-7{,}5):(-2{,}5)$, $\left(-\frac{3}{4}\right):\frac{1}{2}$, $(+0{,}6):(-1{,}5)$.\par
\tithelo Να τα υπολογίσω.\par
\possindeo Ομόσημοι: $7{,}5:2{,}5=3$, άρα $+3$. Ετερόσημοι: $-\frac{3}{4}\cdot2=-\frac{3}{2}$. Ετερόσημοι: $0{,}6:1{,}5=0{,}4$, άρα $-0{,}4$.\par
\apantisi $3$, $-\frac{3}{2}$, $-0{,}4$.\par
\elegxos $(-2{,}5)\cdot3=-7{,}5$ και $(-1{,}5)\cdot(-0{,}4)=0{,}6$.

\paradeigma{Αντίστροφοι}
\tixero Οι αριθμοί $-4$, $-\frac{2}{3}$, $5$, $0$.\par
\tithelo Οι αντίστροφοί τους.\par
\possindeo Ο αντίστροφος του $\alpha$ είναι ο $\frac{1}{\alpha}$ με το ίδιο πρόσημο. Του $-4$: $-\frac{1}{4}$. Του $-\frac{2}{3}$: $-\frac{3}{2}$. Του $5$: $\frac{1}{5}$. Το $0$ δεν έχει.\par
\apantisi $-\frac{1}{4}$, $-\frac{3}{2}$, $\frac{1}{5}$ και το $0$ δεν έχει αντίστροφο.\par
\elegxos $(-4)\cdot\left(-\frac{1}{4}\right)=1$.

\paradeigma{Παραστάσεις και σύνθετο κλάσμα}
\tixero Οι παραστάσεις $(-12):3+(-5)\cdot(-2)$ και $\dfrac{-\frac{6}{5}}{\frac{3}{10}}$.\par
\tithelo Να τις υπολογίσω.\par
\possindeo $(-12):3=-4$ και $(-5)(-2)=10$, άρα $-4+10=6$. Το σύνθετο: $-\frac{6}{5}\cdot\frac{10}{3}=-\frac{60}{15}=-4$.\par
\apantisi $6$ και $-4$.\par
\elegxos $-4\cdot\frac{3}{10}=-\frac{6}{5}$.

\tmima{Ασκήσεις}

\zoni{Το κατάλαβα}{πηλίκα, εξισώσεις, αντίστροφοι}

\ex Υπολόγισε: α) $(-24):6$ \quad β) $(-24):(-3)$ \quad γ) $18:(-6)$ \quad δ) $0:(-5)$ \quad ε) $(-7{,}5):(-2{,}5)$ \quad στ) $(+36):(+9)$.

\ex Λύσε: α) $(-3)x=12$ \quad β) $5x=-20$ \quad γ) $x:(-2)=6$ \quad δ) $(-4)x=-4$ \quad ε) $x:3=-5$ \quad στ) $12:x=-4$.

\ex Υπολόγισε: α) $\left(-\frac{3}{4}\right):\frac{1}{2}$ \quad β) $\left(-\frac{2}{3}\right):\left(-\frac{4}{9}\right)$ \quad γ) $5:\left(-\frac{2}{5}\right)$ \quad δ) $(-1{,}2):0{,}4$ \quad ε) $(+0{,}6):(-1{,}5)$ \quad στ) $\left(-\frac{1}{2}\right):\left(-\frac{1}{2}\right)$.

\ex Βρες τους αντίστροφους των $-4$, $-\frac{2}{3}$, $5$, $-0{,}5$, $\frac{1}{3}$.

\ex Απλοποίησε τα σύνθετα κλάσματα: α) $\dfrac{-\frac{6}{5}}{\frac{3}{10}}$ \quad β) $\dfrac{\frac{2}{3}}{-\frac{4}{9}}$ \quad γ) $\dfrac{-5}{-\frac{1}{2}}$.

\ex Υπολόγισε: α) $(-12):3+(-5)\cdot(-2)$ \quad β) $20:(-4)-(-3)\cdot2$ \quad γ) $[(-18):(-6)]\cdot(-4)$.

\ex Πες το πρόσημο χωρίς να υπολογίσεις: α) $(-45):9$ \quad β) $(-45):(-9)$ \quad γ) $45:(-9)$ \quad δ) $(-1\,000):(-1)$.

\ex Σωστό ή λάθος; α) $(-12):(-4)=-3$. \quad β) $0:(-5)=0$. \quad γ) $(-5):0=0$. \quad δ) $\alpha:\beta=\alpha\cdot\frac{1}{\beta}$. \quad ε) Το πηλίκο ετεροσήμων είναι αρνητικό. \quad στ) $\frac{-6}{3}=\frac{6}{-3}$.

\zoni{Το χρησιμοποιώ}{χρέη, ρυθμοί, μέσοι όροι}

\ex \textbf{Το ταμείο.} Το ταμείο είναι $-24$ € και μοιράζεται εξίσου στα $6$ παιδιά. Πόσο αντιστοιχεί στο καθένα; Αν κάθε παιδί βάζει $-3$ €, πόσα παιδιά έχουν συνεισφέρει συνολικά $-24$ €;

\ex \textbf{Η θερμοκρασία.} Η θερμοκρασία έπεσε $18$ βαθμούς σε $6$ ώρες. Ποιος είναι ο ρυθμός μεταβολής ανά ώρα; Και αν έπεσε $12$ βαθμούς σε $4$ ώρες;

\ex \textbf{Ο δύτης.} Ο δύτης κατέβηκε από το $0$ στο $-36$ μ. σε $9$ λεπτά. Πόσα μέτρα το λεπτό;

\ex \textbf{Σε πόση ώρα;} Η θερμοκρασία πέφτει $2{,}5$ βαθμούς την ώρα. Σε πόσες ώρες θα πέσει $15$ βαθμούς;

\ex \textbf{Ο μέσος όρος.} Οι θερμοκρασίες μιας εβδομάδας ήταν $-3$, $-5$, $2$, $0$, $-4$, $1$, $-1$. Βρες το άθροισμα και τον μέσο όρο στο δέκατο.

\ex \textbf{Οι πόντοι.} Οι διαφορές πόντων της Φανής σε $4$ αγώνες ήταν $-8$, $+4$, $-10$, $+6$. Ποιος είναι ο μέσος όρος;

\ex \textbf{Τα παιδιά.} Το χρέος των $-45$ € μοιράστηκε εξίσου και κάθε παιδί χρωστά $-5$ €. Πόσα είναι τα παιδιά;

\zoni{Το σκέφτομαι}{γιατί όχι με το μηδέν}

\ex Γιατί δεν διαιρούμε με το $0$; Εξέτασε με τον ορισμό τις $0\cdot x=5$ και $0\cdot x=0$.

\ex Ο Στέλιος γράφει $(-12):4=3$. Βρες το λάθος και έλεγξε το σωστό με πολλαπλασιασμό.

\ex Σύγκρινε τα πρόσημα των $(-6):3$, $6:(-3)$, $(-6):(-3)$.

\ex Αν $\alpha:\beta=2$ και $\alpha=-10$, βρες το $\beta$.

\ex Δείξε με $\alpha=6$ και $\beta=3$ ότι $\dfrac{-\alpha}{\beta}=\dfrac{\alpha}{-\beta}=-\dfrac{\alpha}{\beta}$.

\zoni{Για το σπίτι}{μοίρασε και σκέψου}

\ex Φτιάξε ένα πρόβλημα όπου ένα χρέος μοιράζεται ισόποσα και λύσε το.

\paragrafos{Δεκαδική μορφή ρητών: όταν τα ψηφία δεν τελειώνουν}{Α.7.7, «Δεκαδική μορφή ρητών αριθμών»}

\begin{anagke}
Ο Στέλιος παίζει με το κομπιουτεράκι του κινητού και πατάει $1:3$. Βλέπει $0{,}33333333$ και συνεχίζει με το $1:7$: $0{,}14285714$. «Το $1:7$ έχει ένα σωρό ψηφία!». Η Πολυξένη το αντιγράφει και προσπαθεί να δει αν υπάρχει κάποιο μοτίβο: «$142857\ 14\ldots$ Το $142857$ επαναλαμβάνεται!». Δοκιμάζει με το $2:7$ και βρίσκει $0{,}285714\ 285\ldots$: ίδια ψηφία, αλλά ξεκινούν από άλλο σημείο.

Ο Νίκος δοκιμάζει $3:8$ και βρίσκει $0{,}375$ και τελειώνει. Η Αγγελική ρωτάει: «Πότε τελειώνει η διαίρεση και πότε όχι;». Η Μυρτώ κάνει την ευκλείδεια διαίρεση στο χέρι: «Στο $1:7$ τα υπόλοιπα είναι πάντα $1,2,\dots,6$, δεν έχουν άλλες επιλογές. Μετά από έξι βήματα θα ξαναβρούμε κάποιο υπόλοιπο και τα ψηφία θα επαναληφθούν».

Και ο Βασίλης προτείνει το δύσκολο: «Αν πάρω το $0{,}3333\ldots$ και το πολλαπλασιάσω επί τρία, βγαίνει $0{,}9999\ldots$. Αλλά το $\frac{1}{3}\cdot3=1$. Άρα το $0{,}9999\ldots$ ισούται με το $1$;». Σε αυτή την παράγραφο η παρέα βρίσκει τον κανόνα του πότε τελειώνει ένα δεκαδικό ανάπτυγμα, μαθαίνει να γράφει τους άπειρους δεκαδικούς και πώς να γυρίζει πίσω στο κλάσμα.
\end{anagke}

\begin{lexi}{περιοδικός, περίοδος}{περιοδικός, περίοδος}
\lexiword{Περίοδος} από το «περί-οδος», δηλαδή «γύρω-δρόμος»: ένας δρόμος που τον ξανακάνουμε. Ο \lexiword{περιοδικός δεκαδικός} έχει ψηφία που επαναλαμβάνονται σε κύκλο (την περίοδο) για πάντα. Ο \lexiword{ακριβής δεκαδικός} τελειώνει, χωρίς να επαναλαμβάνεται τίποτα.
\end{lexi}

\begin{metafrasi}
Το $\frac{3}{8}$ τελειώνει. & $0{,}375$\\
Το $\frac{1}{3}$ επαναλαμβάνει το $3$. & $0{,}\overline{3}=0{,}333\ldots$\\
Το $\frac{1}{7}$ επαναλαμβάνει το $142857$. & $0{,}\overline{142857}$\\
Το $\frac{5}{6}$: το $8$ μία φορά και μετά το $3$. & $0{,}8\overline{3}=0{,}8333\ldots$\\
Από περιοδικό σε κλάσμα. & $0{,}\overline{27}=\dfrac{27}{99}=\dfrac{3}{11}$\\
\end{metafrasi}

\tmima{Η θεωρία}

Κάθε ρητός $\frac{\kappa}{\lambda}$ γίνεται δεκαδικός αν κάνουμε τη διαίρεση $\kappa:\lambda$. Υπάρχουν δύο ενδεχόμενα:

\begin{itemize}[leftmargin=1.2em,itemsep=0.3ex,topsep=2pt]
\item Η διαίρεση \textbf{τελειώνει} (υπόλοιπο $0$): έχουμε \textbf{ακριβή δεκαδικό} (πεπερασμένο πλήθος δεκαδικών ψηφίων). Π.χ. $\frac{3}{8}=0{,}375$, $\frac{9}{20}=0{,}45$.
\item Η διαίρεση \textbf{δεν τελειώνει}: τα υπόλοιπα είναι πάντα μικρότερα από τον διαιρέτη $\lambda$, άρα κάποτε επαναλαμβάνονται και μαζί τους επαναλαμβάνονται τα ψηφία του πηλίκου. Έχουμε \textbf{περιοδικό δεκαδικό}. Γράφουμε $\frac{5}{6}=0{,}8\overline{3}$ και $\frac{4}{11}=0{,}\overline{36}$.
\end{itemize}

Το τμήμα που επαναλαμβάνεται λέγεται \textbf{περίοδος}. Αν η περίοδος αρχίζει αμέσως μετά την υποδιαστολή, ο δεκαδικός λέγεται \emph{απλός περιοδικός} ($0{,}\overline{36}$), αλλιώς \emph{μικτός} ($0{,}8\overline{3}$).

\begin{kanonas}{πότε τελειώνει}
Ένα ανάγωγο κλάσμα γράφεται ως ακριβής δεκαδικός όταν ο παρονομαστής του έχει \textbf{μόνο παράγοντες $2$ και $5$} (γιατί τότε γίνεται δεκαδικό κλάσμα με παρονομαστή δύναμη του $10$). Αλλιώς είναι περιοδικός.
\end{kanonas}

Έτσι το $\frac{7}{20}$ ($20=2^2\cdot5$) τελειώνει: $0{,}35$. Το $\frac{5}{12}$ ($12=2^2\cdot3$) δεν τελειώνει: $0{,}41\overline{6}$.

\tmimasmall{Από περιοδικό σε κλάσμα}

Αν $x=0{,}\overline{3}$, τότε $10x=3{,}\overline{3}$ και αφαιρώντας $10x-x=3$, άρα $9x=3$ και $x=\frac{1}{3}$. Με την ίδια ιδέα:
\begin{itemize}[leftmargin=1.2em,itemsep=0.3ex,topsep=2pt]
\item \textbf{Απλός περιοδικός:} αριθμητής η περίοδος και παρονομαστής τόσα εννιάρια όσα τα ψηφία της περιόδου: $0{,}\overline{27}=\frac{27}{99}=\frac{3}{11}$, $0{,}\overline{7}=\frac{7}{9}$, $1{,}\overline{3}=1+\frac{3}{9}=\frac{4}{3}$.
\item \textbf{Μικτός:} $0{,}2\overline{3}$: με $x=0{,}2\overline{3}$, $10x=2{,}\overline{3}=2+\frac{1}{3}=\frac{7}{3}$, άρα $x=\frac{7}{30}$.
\end{itemize}
Κάθε περιοδικός δεκαδικός είναι ρητός, και κάθε ρητός είναι ακριβής ή περιοδικός δεκαδικός. Ο $0{,}\overline{9}=1$: με $x=0{,}\overline{9}$, $10x=9{,}\overline{9}$, $9x=9$, $x=1$.

\begin{prosoxi}
\begin{itemize}[leftmargin=1.2em,itemsep=0.3ex,topsep=2pt]
\item Το $\frac{1}{3}$ δεν είναι ίσο με $0{,}33$ αλλά με $0{,}\overline{3}$: το $0{,}33$ είναι προσέγγιση.
\item Η περίοδος είναι μόνο το κομμάτι που επαναλαμβάνεται: στο $0{,}8\overline{3}$ το $8$ δεν επαναλαμβάνεται.
\item Η περίοδος ενός κλάσματος με παρονομαστή $\lambda$ έχει το πολύ $\lambda-1$ ψηφία (λόγω των υπολοίπων).
\item Τα κομπιουτεράκια κόβουν ή στρογγυλοποιούν τα ψηφία: $1:7$ δεν σταματά στο $0{,}14285714$.
\end{itemize}
\end{prosoxi}

\tmima{Λυμένα παραδείγματα}

\paradeigma{Ακριβής ή περιοδικός;}
\tixero Τα κλάσματα $\frac{7}{20}$, $\frac{5}{12}$, $\frac{9}{14}$, $\frac{11}{40}$.\par
\tithelo Ποια έχουν ακριβή δεκαδικό.\par
\possindeo Αναλύω τους παρονομαστές: $20=2^2\cdot5$, $12=2^2\cdot3$, $14=2\cdot7$, $40=2^3\cdot5$. Μόνο παράγοντες $2$ και $5$ έχουν το $20$ και το $40$.\par
\apantisi Ακριβείς: $\frac{7}{20}=0{,}35$ και $\frac{11}{40}=0{,}275$. Περιοδικοί: $\frac{5}{12}=0{,}41\overline{6}$ και $\frac{9}{14}=0{,}6\overline{428571}$.\par
\elegxos Το $0{,}35=\frac{35}{100}=\frac{7}{20}$.

\paradeigma{Η διαίρεση του 1:7}
\tixero Το $\frac{1}{7}$.\par
\tithelo Να βρω τα ψηφία και την περίοδο.\par
\possindeo $1:7$: $10:7=1$ υπόλοιπο $3$, $30:7=4$ υπόλοιπο $2$, $20:7=2$ υπόλοιπο $6$, $60:7=8$ υπόλοιπο $4$, $40:7=5$ υπόλοιπο $5$, $50:7=7$ υπόλοιπο $1$. Το υπόλοιπο $1$ ξαναεμφανίζεται, άρα η περίοδος είναι $142857$.\par
\apantisi $\frac{1}{7}=0{,}\overline{142857}$.\par
\elegxos Η περίοδος έχει $6$ ψηφία, όσα τα δυνατά υπόλοιπα $1,\dots,6$.

\paradeigma{Από περιοδικό σε κλάσμα}
\tixero Οι $0{,}\overline{27}$, $2{,}\overline{6}$ και $0{,}1\overline{6}$.\par
\tithelo Να τους γράψω ως κλάσματα.\par
\possindeo $0{,}\overline{27}=\frac{27}{99}=\frac{3}{11}$. $2{,}\overline{6}=2+\frac{6}{9}=2+\frac{2}{3}=\frac{8}{3}$. Για τον τρίτο $x=0{,}1\overline{6}$, $10x=1{,}\overline{6}=1+\frac{2}{3}=\frac{5}{3}$, άρα $x=\frac{5}{30}=\frac{1}{6}$.\par
\apantisi $\frac{3}{11}$, $\frac{8}{3}$, $\frac{1}{6}$.\par
\elegxos Το $\frac{1}{6}=0{,}1666\ldots$ και το $3:11=0{,}2727\ldots$.

\paradeigma{Στρογγυλοποίηση περιοδικών}
\tixero Οι $\frac{2}{3}$, $\frac{5}{6}$, $\frac{1}{7}$.\par
\tithelo Στρογγυλοποίησε στα εκατοστά.\par
\possindeo $\frac{2}{3}=0{,}666\ldots\approx0{,}67$. $\frac{5}{6}=0{,}8333\ldots\approx0{,}83$. $\frac{1}{7}=0{,}1428\ldots\approx0{,}14$.\par
\apantisi $0{,}67$, $0{,}83$, $0{,}14$.\par
\elegxos Το σφάλμα στο πρώτο: $0{,}67-0{,}666\ldots=0{,}0033\ldots<0{,}005$.

\paradeigma{Ισούται το $0{,}\overline{9}$ με το $1$;}
\tixero Ο αριθμός $0{,}\overline{9}$.\par
\tithelo Να τον συγκρίνω με το $1$.\par
\possindeo Αν $x=0{,}\overline{9}$, τότε $10x=9{,}\overline{9}$. Αφαιρώ: $10x-x=9$, άρα $9x=9$ και $x=1$. Ή: $\frac{1}{3}=0{,}\overline{3}$, άρα $3\cdot\frac{1}{3}=3\cdot0{,}\overline{3}=0{,}\overline{9}$, δηλαδή $1=0{,}\overline{9}$.\par
\apantisi $0{,}\overline{9}=1$.\par
\elegxos Δεν υπάρχει αριθμός ανάμεσα στο $0{,}\overline{9}$ και στο $1$.

\tmima{Ασκήσεις}

\zoni{Το κατάλαβα}{ακριβείς και περιοδικοί}

\ex Γράψε σε δεκαδική μορφή (ακριβή ή περιοδική): α) $\frac{3}{8}$ \quad β) $\frac{5}{6}$ \quad γ) $\frac{4}{11}$ \quad δ) $\frac{7}{9}$ \quad ε) $\frac{9}{20}$ \quad στ) $\frac{2}{15}$ \quad ζ) $\frac{1}{12}$.

\ex Ποια από τα $\frac{7}{20}$, $\frac{5}{12}$, $\frac{9}{14}$, $\frac{11}{40}$, $\frac{3}{16}$, $\frac{13}{25}$ γράφονται ως ακριβείς δεκαδικοί; Χωρίς να κάνεις τη διαίρεση.

\ex Γράψε τα πρώτα οκτώ δεκαδικά ψηφία των $0{,}\overline{36}$, $2{,}1\overline{45}$ και $0{,}\overline{142857}$.

\ex Γράψε ως κλάσματα: α) $0{,}\overline{7}$ \quad β) $0{,}\overline{15}$ \quad γ) $1{,}\overline{3}$ \quad δ) $0{,}\overline{81}$ \quad ε) $0{,}2\overline{3}$.

\ex Σύγκρινε: α) $0{,}\overline{3}$ και $0{,}33$ \quad β) $0{,}\overline{6}$ και $\frac{2}{3}$ \quad γ) $0{,}7\overline{2}$ και $0{,}72$.

\ex Στρογγυλοποίησε στα εκατοστά: α) $\frac{2}{3}$ \quad β) $\frac{5}{6}$ \quad γ) $\frac{1}{7}$ \quad δ) $\frac{4}{9}$.

\ex Σωστό ή λάθος; α) Κάθε ρητός είναι ακριβής ή περιοδικός δεκαδικός. \quad β) $\frac{1}{3}=0{,}33$. \quad γ) $0{,}\overline{9}=1$. \quad δ) Το $\frac{1}{7}$ έχει περίοδο $142857$. \quad ε) Τα κλάσματα με παρονομαστή $8$ είναι ακριβείς δεκαδικοί. \quad στ) $\frac{1}{6}=0{,}16$.

\zoni{Το χρησιμοποιώ}{κομπιουτεράκι, μοιρασιές, μέσοι όροι}

\ex \textbf{Το κομπιουτεράκι.} Ο Στέλιος γράφει $1:7$ και βλέπει $0{,}14285714$. Ποια είναι η περίοδος; Βρες με το κομπιουτεράκι το $2:7$ και το $3:7$ και πες τι παρατηρείς.

\ex \textbf{Η μοιρασιά.} Τα $10$ € μοιράζονται σε $3$ παιδιά. Τι δείχνει το $10:3$ και πόσο παίρνει στην πράξη το καθένα;

\ex \textbf{Η ταχύτητα.} Ο Νίκος τρέχει $100$ μ. σε $12$ s. Πόση είναι η ταχύτητά του σε m/s ως περιοδικός και ως στρογγυλοποιημένος στα εκατοστά;

\ex \textbf{Τα ποσοστά.} Γράψε το $\frac{1}{3}$ και το $\frac{5}{6}$ ως ποσοστά με περιοδικούς και στρογγυλοποίησε στο δέκατο.

\ex \textbf{Ο μέσος όρος.} Οι βαθμοί $7$, $8$, $8$ έχουν μέσο όρο $\frac{23}{3}$. Γράψε τον ως περιοδικό και στρογγυλοποίησέ τον στα εκατοστά.

\zoni{Το σκέφτομαι}{γιατί επαναλαμβάνονται}

\ex Εξήγησε γιατί η διαίρεση $\kappa:7$ έχει περίοδο το πολύ $6$ ψηφίων.

\ex Δείξε ότι $0{,}\overline{9}=1$ με τον τρόπο $x=0{,}\overline{9}$, $10x=9{,}\overline{9}$.

\ex Βρες τη δεκαδική μορφή του $\frac{1}{11}$ και του $\frac{1}{13}$ (δοκίμασε με κομπιουτεράκι και με διαίρεση).

\ex Βρες έναν ρητό ανάμεσα στους $0{,}\overline{3}$ και $0{,}34$.

\ex Γιατί το $\frac{1}{3}$ ισούται με $0{,}\overline{3}$ και όχι με $0{,}333$; Τι μας λέει το $3\cdot0{,}\overline{3}$;

\zoni{Για το σπίτι}{ανακάλυψε τα μοτίβα}

\ex Υπολόγισε με κομπιουτεράκι το $1:\nu$ για $\nu=2,\dots,12$ και βρες ποια έχουν ακριβές ανάπτυγμα. Τι κοινό έχουν οι παρονομαστές τους;

\paragrafos{Δυνάμεις με εκθέτη φυσικό: όταν οι πολλαπλασιασμοί συνεχίζονται}{Α.7.8, «Δυνάμεις ρητών αριθμών με εκθέτη φυσικό»}

\begin{anagke}
Στο διάλειμμα η Πολυξένη λέει στην Αγγελική ένα μυστικό και της ζητάει να το πει μέσα σε μια ώρα σε δύο άλλα παιδιά. Κάθε παιδί που το μαθαίνει το λέει, μέσα στην επόμενη ώρα, σε δύο άλλα. «Μέσα σε δέκα ώρες θα το ξέρει όλο το σχολείο;» ρωτάει ο Ηλίας. Ο Στέλιος γράφει: $1,\ 2,\ 4,\ 8,\ 16,\ \ldots$ και κάθε φορά το πλήθος διπλασιάζεται. «Μετά από $10$ ώρες θα είναι $2\cdot2\cdot2\cdot\ldots$ δέκα φορές».

Ο Βασίλης λέει ότι $2$ δέκα φορές επί $2$ είναι $2\cdot10=20$. Ο Νίκος τον διορθώνει: «Όχι, δέκα φορές πολλαπλασιασμός, όχι δέκα φορές πρόσθεση. Θα βγει $1\,024$!». Η Ρήνη προτείνει μια σύντομη γραφή: $2\cdot2\cdot2\cdot2\cdot2=2^5$, το «δύο στην πέμπτη», που το ξέρουν από την Ε΄ και ΣΤ΄ Δημοτικού για τους φυσικούς. Η κυρία Αλεξάνδρα το βάζει στο νέο επίπεδο: «Μπορούμε να υψώσουμε στη δύναμη και έναν ρητό. Τι νομίζετε ότι είναι το $(-2)^3$; Και το $\left(-\frac{1}{2}\right)^3$;».

Σε αυτή την παράγραφο η παρέα ξαναβρίσκει τις δυνάμεις, μαθαίνει το πρόσημό τους και τις ιδιότητες που τις κάνουν γρήγορες, και προσέχει τα σημεία που μπερδεύουν τους περισσότερους: το $(-3)^2$ και το $-3^2$.
\end{anagke}

\begin{lexi}{δύναμη, βάση, εκθέτης}{δύναμη, βάση, εκθέτης}
Η \lexiword{δύναμη} είναι «ύψωση», δηλαδή το γινόμενο ίσων παραγόντων· στα αρχαία ελληνικά «δύναμις» ήταν και το τετράγωνο. Η \lexiword{βάση} είναι ο αριθμός που πολλαπλασιάζεται με τον εαυτό του και ο \lexiword{εκθέτης} (από το «εκ-τίθημι», «βάζω έξω») αυτός που γράφουμε πάνω δεξιά, δείχνοντας πόσους παράγοντες έχουμε.
\end{lexi}

\begin{metafrasi}
Πέντε φορές ο $2$ πολλαπλασιασμένος. & $2^5=2\cdot2\cdot2\cdot2\cdot2=32$\\
Ο $-2$ τρεις φορές. & $(-2)^3=(-2)(-2)(-2)=-8$\\
Μετά από ν γύρους το μυστικό το ξέρουν. & $2^\nu$ παιδιά\\
Πολλαπλασιάζω δυνάμεις με ίδια βάση. & $\alpha^\mu\cdot\alpha^\nu=\alpha^{\mu+\nu}$\\
Η δύναμη μιας δύναμης. & $(\alpha^\mu)^\nu=\alpha^{\mu\nu}$\\
\end{metafrasi}

\tmima{Η θεωρία}

\begin{orismos}{δύναμη με εκθέτη φυσικό}
Για κάθε ρητό $\alpha$ και φυσικό $\nu\ge2$ η \textbf{δύναμη} $\alpha^\nu$ είναι το γινόμενο $\nu$ παραγόντων ίσων με $\alpha$: $\alpha^\nu=\alpha\cdot\alpha\cdots\alpha$. Ορίζουμε ακόμη $\alpha^1=\alpha$.
\end{orismos}

Για παράδειγμα $(0{,}1)^3=0{,}1\cdot0{,}1\cdot0{,}1=0{,}001$ και $\left(\frac{2}{3}\right)^3=\frac{2}{3}\cdot\frac{2}{3}\cdot\frac{2}{3}=\frac{8}{27}$.

\tmimasmall{Το πρόσημο της δύναμης}

Από τον κανόνα των προσήμων του πολλαπλασιασμού:
\begin{itemize}[leftmargin=1.2em,itemsep=0.3ex,topsep=2pt]
\item Μια δύναμη με θετική βάση είναι πάντα θετική.
\item Μια δύναμη με αρνητική βάση είναι \textbf{θετική} αν ο εκθέτης είναι \textbf{άρτιος} και \textbf{αρνητική} αν είναι \textbf{περιττός}: $(-2)^4=16$, $(-2)^3=-8$, $(-1)^{100}=1$, $(-1)^{101}=-1$.
\end{itemize}

Η βάση είναι ό,τι βρίσκεται μέσα στην παρένθεση: $(-3)^2=9$ ενώ $-3^2=-(3\cdot3)=-9$.

\begin{idiotita}{ιδιότητες των δυνάμεων (εκθέτης φυσικός)}
Για ρητούς $\alpha,\,\beta$ και φυσικούς $\mu,\,\nu$ ισχύουν:
\begin{itemize}[leftmargin=1.2em,itemsep=0.3ex,topsep=2pt]
\item $\alpha^\mu\cdot\alpha^\nu=\alpha^{\mu+\nu}$ \qquad $\alpha^\mu:\alpha^\nu=\alpha^{\mu-\nu}$ (με $\alpha\neq0$, $\mu>\nu$)
\item $(\alpha\cdot\beta)^\nu=\alpha^\nu\cdot\beta^\nu$ \qquad $\left(\dfrac{\alpha}{\beta}\right)^\nu=\dfrac{\alpha^\nu}{\beta^\nu}$ (με $\beta\neq0$)
\item $(\alpha^\mu)^\nu=\alpha^{\mu\cdot\nu}$
\end{itemize}
\end{idiotita}

Γιατί; Το $2^3\cdot2^4=(2\cdot2\cdot2)\cdot(2\cdot2\cdot2\cdot2)=2^7$ (τρεις παράγοντες και τέσσερις: επτά). Το $(3^2)^3=3^2\cdot3^2\cdot3^2=3^6$. Το $(2\cdot5)^3=2\cdot5\cdot2\cdot5\cdot2\cdot5=2^3\cdot5^3=8\cdot125=1\,000$.

\tmimasmall{Προτεραιότητα πράξεων}

Η σειρά είναι: παρενθέσεις, μετά δυνάμεις, μετά πολλαπλασιασμοί και διαιρέσεις, τέλος προσθέσεις και αφαιρέσεις. Για παράδειγμα $2\cdot3^2-(-2)^3+10:(-5)=18-(-8)+(-2)=18+8-2=24$.

\begin{prosoxi}
\begin{itemize}[leftmargin=1.2em,itemsep=0.3ex,topsep=2pt]
\item $2^3=8$, όχι $6$: η δύναμη είναι επανάληψη πολλαπλασιασμού, όχι γινόμενο $2\cdot3$.
\item $(-3)^2=+9$, αλλά $-3^2=-9$ (ο εκθέτης ενεργεί μόνο στο $3$).
\item Το $(\alpha+\beta)^2$ δεν είναι $\alpha^2+\beta^2$: $(3+4)^2=49$ ενώ $3^2+4^2=25$.
\item Οι ιδιότητες $\alpha^\mu\cdot\alpha^\nu=\alpha^{\mu+\nu}$ ισχύουν μόνο για \emph{ίδια βάση}: το $2^3\cdot3^2$ δεν γίνεται $6^5$.
\item Στο $\alpha^\mu\cdot\alpha^\nu$ προσθέτω εκθέτες, δεν πολλαπλασιάζω: $2^3\cdot2^4=2^7$ και όχι $2^{12}$.
\end{itemize}
\end{prosoxi}

\tmima{Λυμένα παραδείγματα}

\paradeigma{Δυνάμεις και πρόσημα}
\tixero Οι δυνάμεις $2^5$, $(-3)^2$, $-3^2$, $(-3)^3$, $\left(-\frac{1}{2}\right)^3$, $(0{,}1)^3$.\par
\tithelo Να τις υπολογίσω.\par
\possindeo $2^5=32$. $(-3)^2=9$. $-3^2=-(3\cdot3)=-9$. $(-3)^3=-27$. $\left(-\frac{1}{2}\right)^3=-\frac{1}{8}$. $(0{,}1)^3=0{,}001$.\par
\apantisi $32$, $9$, $-9$, $-27$, $-\frac{1}{8}$, $0{,}001$.\par
\elegxos Αρνητική βάση με περιττό εκθέτη δίνει αρνητικό: το $(-3)^3$ και το $\left(-\frac{1}{2}\right)^3$.

\paradeigma{Με τις ιδιότητες}
\tixero Οι παραστάσεις $2^3\cdot2^4$, $5^8:5^5$, $(3^2)^3$, $(2\cdot5)^3$, $\left(\frac{2}{3}\right)^3$.\par
\tithelo Να τις υπολογίσω.\par
\possindeo $2^{3+4}=2^7=128$. $5^{8-5}=5^3=125$. $3^{2\cdot3}=3^6=729$. $2^3\cdot5^3=8\cdot125=1\,000$. $\frac{2^3}{3^3}=\frac{8}{27}$.\par
\apantisi $128$, $125$, $729$, $1\,000$, $\frac{8}{27}$.\par
\elegxos $(2\cdot5)^3=10^3=1\,000$.

\paradeigma{Το μυστικό του σχολείου}
\tixero Κάθε παιδί λέει το μυστικό σε $2$ νέα παιδιά κάθε ώρα. Στην αρχή το ξέρει $1$ παιδί.\par
\tithelo Πόσα παιδιά το ξέρουν μετά από $10$ ώρες.\par
\possindeo Μετά από $1$ ώρα: $2$. Μετά από $2$: $4$. Μετά από $3$: $8$. Μετά από $\nu$ ώρες το πλήθος διπλασιάζεται $\nu$ φορές: $2^\nu$. Για $\nu=10$: $2^{10}=1\,024$.\par
\apantisi $1\,024$ παιδιά.\par
\elegxos Από την $9η$ ώρα ($512$) στη $10η$ διπλασιάζεται: $1\,024$.

\paradeigma{Η προτεραιότητα}
\tixero Η παράσταση $2\cdot3^2-(-2)^3+10:(-5)$.\par
\tithelo Να την υπολογίσω.\par
\possindeo Πρώτα οι δυνάμεις: $3^2=9$, $(-2)^3=-8$. Μετά πολλαπλασιασμοί και διαιρέσεις: $2\cdot9=18$, $10:(-5)=-2$. Τέλος $18-(-8)+(-2)=18+8-2=24$.\par
\apantisi $24$.\par
\elegxos Αν έκανα πρώτα τον πολλαπλασιασμό $2\cdot3=6$, θα έβγαινε $6^2=36$: λάθος σειρά.

\tmima{Ασκήσεις}

\zoni{Το κατάλαβα}{δυνάμεις, πρόσημα, ιδιότητες}

\ex Υπολόγισε: α) $2^5$ \quad β) $3^4$ \quad γ) $(-2)^3$ \quad δ) $(-2)^4$ \quad ε) $(-1)^7$ \quad στ) $0^5$.

\ex Υπολόγισε: α) $(-3)^2$ \quad β) $-3^2$ \quad γ) $\left(-\frac{1}{2}\right)^3$ \quad δ) $(0{,}1)^3$ \quad ε) $(-0{,}2)^2$ \quad στ) $\left(\frac{2}{3}\right)^3$.

\ex Πες το πρόσημο χωρίς να υπολογίσεις: α) $(-5)^6$ \quad β) $(-5)^7$ \quad γ) $-(-3)^4$ \quad δ) $(-2)^3\cdot(-3)^2$.

\ex Γράψε ως μία δύναμη: α) $2^3\cdot2^4$ \quad β) $5^8:5^5$ \quad γ) $(3^2)^3$ \quad δ) $7^2\cdot7\cdot7^3$ \quad ε) $(-2)^5\cdot(-2)^2$ \quad στ) $(-4)^6:(-4)^4$.

\ex Υπολόγισε με τις ιδιότητες: α) $(2\cdot5)^3$ \quad β) $(-2\cdot3)^2$ \quad γ) $\left(\frac{2}{3}\right)^3$ \quad δ) $\left(-\frac{3}{4}\right)^2$.

\ex Υπολόγισε: α) $2\cdot3^2-(-2)^3+10:(-5)$ \quad β) $(-2)^2\cdot3-4^2:(-2)$ \quad γ) $5-3\cdot(-2)^2$ \quad δ) $\big((-1)^5+(-1)^6\big)\cdot10$.

\ex Σωστό ή λάθος; α) $2^3=6$. \quad β) $(-2)^4=16$. \quad γ) $-2^4=16$. \quad δ) $(\alpha\beta)^2=\alpha^2\beta^2$. \quad ε) $\alpha^2\cdot\alpha^3=\alpha^6$. \quad στ) $(\alpha^2)^3=\alpha^6$. \quad ζ) $(\alpha+\beta)^2=\alpha^2+\beta^2$.

\zoni{Το χρησιμοποιώ}{φήμες, κύβοι, σκάκι}

\ex \textbf{Το μυστικό.} Κάθε παιδί λέει το μυστικό σε $2$ νέα κάθε ώρα. Φτιάξε πίνακα με το πλήθος των παιδιών μετά από $0,\,1,\,\dots,\,10$ ώρες και πες πότε ξεπερνά τα $1\,000$.

\ex \textbf{Ο κύβος.} Ένας κύβος έχει ακμή $12$ cm. Βρες τον όγκο ($12^3$) και το εμβαδόν της επιφάνειάς του ($6\cdot12^2$).

\ex \textbf{Το σκάκι.} Στο πρώτο τετράγωνο ενός σκακιστικού πίνακα βάζουμε $1$ κόκκο, στο δεύτερο $2$, στο τρίτο $4$ και συνεχίζουμε διπλασιάζοντας. Πόσοι κόκκοι είναι στο $8$ο τετράγωνο και πόσοι στο $11$ο;

\ex \textbf{Τα βακτήρια.} Ένα βακτήριο τριπλασιάζεται κάθε ώρα. Πόσα είναι μετά από $5$ ώρες;

\ex \textbf{Οι τετράγωνοι αριθμοί.} Οι αριθμοί $1,\,4,\,9,\,16,\dots$ είναι τα τετράγωνα. Ποιος είναι ο $20$ός και ο $\nu$-οστός;

\ex \textbf{Η τιμή διπλασιάζεται.} Μια τιμή των $3$ € διπλασιάζεται $4$ φορές. Πόσο γίνεται;

\ex \textbf{Η δύναμη του $10$.} Υπολόγισε $10^3\cdot10^4$ και $10^8:10^5$.

\zoni{Το σκέφτομαι}{όταν το πρόσημο και ο εκθέτης παίζουν}

\ex Ο Βασίλης λέει ότι $2^3=6$. Βρες το λάθος και εξήγησε τη διαφορά ανάμεσα στο $2\cdot3$ και στο $2^3$.

\ex Δείξε με $\alpha=3$ και $\beta=4$ ότι $(\alpha+\beta)^2\neq\alpha^2+\beta^2$.

\ex Πότε ισχύει $(-\alpha)^\nu=\alpha^\nu$ και πότε $(-\alpha)^\nu=-\alpha^\nu$;

\ex Πόσο είναι τα $1^\nu$, $0^\nu$ (για $\nu\ge1$) και $(-1)^\nu$;

\ex Εξήγησε γιατί $2^3\cdot2^4=2^7$ γράφοντας όλους τους παράγοντες.

\ex Ποιο είναι μεγαλύτερο: το $3^4$ ή το $4^3$;

\zoni{Για το σπίτι}{δίπλωσε και μέτρα}

\ex Δίπλωσε ένα χαρτί στη μέση όσες φορές μπορείς. Πόσα στρώματα έχεις μετά από $1,\,2,\,3,\,4,\,5$ διπλώματα;

\paragrafos{Δυνάμεις με εκθέτη ακέραιο: προς τα αριστερά του μηδενός}{Α.7.9, «Δυνάμεις ρητών αριθμών με εκθέτη ακέραιο»}

\begin{anagke}
Η Πολυξένη γράφει στον πίνακα τις δυνάμεις του $2$: $2^4=16$, $2^3=8$, $2^2=4$, $2^1=2$. Ο Στέλιος κοιτάζει και λέει: «Κάθε φορά που ο εκθέτης μειώνεται κατά $1$, ο αριθμός γίνεται το μισό: $16$, $8$, $4$, $2$». «Και αν συνεχίσουμε;» ρωτάει η Αγγελική. Ο επόμενος αριθμός πρέπει να είναι το μισό του $2$, δηλαδή $1$: $2^0=1$. Και μετά το μισό του $1$: $\frac{1}{2}$, δηλαδή $2^{-1}=\frac{1}{2}$. Και $2^{-2}=\frac{1}{4}$, $2^{-3}=\frac{1}{8}$.

«Αρνητικός εκθέτης;» αναρωτιέται ο Βασίλης. «Δεν μπορώ να πολλαπλασιάσω τον $2$ μείον τρεις φορές!». Η κυρία Αλεξάνδρα εξηγεί: «Δεν χρειάζεται να το κάνετε. Αφού ο εκθέτης που μικραίνει διαιρεί με το $2$, μπορούμε να ορίσουμε ότι συνεχίζεται το ίδιο μοτίβο και στους αρνητικούς, ώστε να ισχύουν οι ιδιότητες των δυνάμεων που ξέρετε». Ο Νίκος το συνδέει με τις μονάδες: «Ένα χιλιοστό είναι το $\frac{1}{1\,000}$ του μέτρου, δηλαδή $10^{-3}$ μέτρα. Αυτό θα γράφουμε στο εργαστήριο».

Στην παράγραφο αυτή η παρέα αποκτά ένα εργαλείο για τους πολύ μεγάλους και τους πολύ μικρούς αριθμούς, και δίνει νόημα στο $\alpha^0$ και στο $\alpha^{-\nu}$.
\end{anagke}

\begin{lexi}{αρνητικός εκθέτης}{αντίστροφος, δύναμη}
Στα μαθηματικά οι νέες έννοιες ορίζονται ώστε να «συνεχίζουν το μοτίβο» και να ισχύουν οι παλιοί κανόνες. Η δύναμη με αρνητικό εκθέτη δεν σημαίνει «πολλαπλασιάζω αρνητικές φορές» αλλά «αντίστροφος της δύναμης»: το $2^{-3}$ είναι ο \lexiword{αντίστροφος} του $2^3$.
\end{lexi}

\begin{metafrasi}
Το μισό του $2$. & $2^0=1$\\
Το μισό του $1$. & $2^{-1}=\dfrac{1}{2}$\\
Ένα χιλιοστό του μέτρου. & $10^{-3}$ μ.\\
Η δύναμη με αρνητικό εκθέτη. & $\alpha^{-\nu}=\dfrac{1}{\alpha^{\nu}}$\\
Κλάσμα με αρνητικό εκθέτη. & $\left(\dfrac{\alpha}{\beta}\right)^{-\nu}=\left(\dfrac{\beta}{\alpha}\right)^{\nu}$\\
\end{metafrasi}

\tmima{Η θεωρία}

Ας γράψουμε τις δυνάμεις του $2$ και ας συνεχίσουμε προς τα δεξιά:
\begin{center}
\renewcommand{\arraystretch}{1.3}
\begin{tabular}{c|ccccccc}
$\nu$ & $3$ & $2$ & $1$ & $0$ & $-1$ & $-2$ & $-3$\\\hline
$2^\nu$ & $8$ & $4$ & $2$ & $1$ & $\frac{1}{2}$ & $\frac{1}{4}$ & $\frac{1}{8}$
\end{tabular}
\end{center}
Κάθε φορά που ο εκθέτης μειώνεται κατά $1$, η τιμή διαιρείται με $2$. Έτσι ορίζουμε:

\begin{orismos}{δύναμη με εκθέτη ακέραιο}
Για κάθε ρητό $\alpha\neq0$ και φυσικό $\nu$:
$$\alpha^0=1,\qquad \alpha^{-\nu}=\dfrac{1}{\alpha^{\nu}}.$$
\end{orismos}

Για παράδειγμα $5^0=1$, $(-7)^0=1$, $3^{-2}=\frac{1}{9}$, $(-2)^{-3}=\frac{1}{(-2)^3}=-\frac{1}{8}$, $10^{-3}=\frac{1}{1\,000}=0{,}001$. Το $0^0$ και τα $0^{-\nu}$ δεν ορίζονται.

\tmimasmall{Κλάσμα στη δύναμη}

Αν η βάση είναι κλάσμα, ο αρνητικός εκθέτης αναποδογυρίζει το κλάσμα:
$$\left(\dfrac{\alpha}{\beta}\right)^{-\nu}=\left(\dfrac{\beta}{\alpha}\right)^{\nu}\quad(\alpha,\beta\neq0).$$
Για παράδειγμα $\left(\frac{1}{2}\right)^{-3}=2^3=8$ και $\left(\frac{2}{3}\right)^{-2}=\left(\frac{3}{2}\right)^2=\frac{9}{4}$.

\begin{idiotita}{ιδιότητες με ακέραιο εκθέτη}
Οι ιδιότητες των δυνάμεων ισχύουν και για ακέραιους εκθέτες ($\alpha,\beta\neq0$): $\alpha^\mu\cdot\alpha^\nu=\alpha^{\mu+\nu}$, $\alpha^\mu:\alpha^\nu=\alpha^{\mu-\nu}$, $(\alpha\beta)^\nu=\alpha^\nu\beta^\nu$, $\left(\frac{\alpha}{\beta}\right)^\nu=\frac{\alpha^\nu}{\beta^\nu}$, $(\alpha^\mu)^\nu=\alpha^{\mu\nu}$.
\end{idiotita}

Με το νέο ορισμό η διαίρεση δεν χρειάζεται περιορισμό $\mu>\nu$: $5^3:5^5=5^{-2}=\frac{1}{25}$. Και βγαίνει τώρα φυσιολογικά ότι $\alpha^0=1$: $\alpha^3:\alpha^3=\alpha^{3-3}=\alpha^0$ και, αφού $\alpha^3:\alpha^3=1$, είναι $\alpha^0=1$.

\begin{prosoxi}
\begin{itemize}[leftmargin=1.2em,itemsep=0.3ex,topsep=2pt]
\item Το $2^{-3}$ δεν είναι $-8$ ούτε $-6$: είναι $\frac{1}{8}$, θετικό. Το αρνητικό σημείο στον εκθέτη δεν αλλάζει το πρόσημο.
\item $\alpha^0=1$ και όχι $0$ (για $\alpha\neq0$).
\item $\left(\frac{1}{2}\right)^{-3}=8$, όχι $\frac{1}{8}$: αναποδογυρίζω το κλάσμα και κάνω τον εκθέτη θετικό.
\item Το πρόσημο της βάσης παίζει ρόλο: $(-2)^{-3}=-\frac{1}{8}$ και $(-2)^{-2}=\frac{1}{4}$.
\item Το $0^{-1}$ δεν ορίζεται (θα σήμαινε $\frac{1}{0}$).
\end{itemize}
\end{prosoxi}

\tmima{Λυμένα παραδείγματα}

\paradeigma{Συνεχίζω τον πίνακα}
\tixero Οι δυνάμεις του $3$: $3^3=27$, $3^2=9$, $3^1=3$.\par
\tithelo Να βρω τα $3^0$, $3^{-1}$, $3^{-2}$, $3^{-3}$.\par
\possindeo Κάθε φορά ο εκθέτης μειώνεται και η τιμή διαιρείται με $3$: $3:3=1$, $1:3=\frac{1}{3}$, $\frac{1}{3}:3=\frac{1}{9}$, $\frac{1}{9}:3=\frac{1}{27}$.\par
\apantisi $3^0=1$, $3^{-1}=\frac{1}{3}$, $3^{-2}=\frac{1}{9}$, $3^{-3}=\frac{1}{27}$.\par
\elegxos Σύμφωνα με τον ορισμό $3^{-2}=\frac{1}{3^2}=\frac{1}{9}$.

\paradeigma{Υπολογισμοί}
\tixero Οι $3^{-2}$, $(-2)^{-3}$, $\left(\frac{1}{2}\right)^{-3}$, $\left(\frac{2}{3}\right)^{-2}$, $5^0$.\par
\tithelo Να τους υπολογίσω.\par
\possindeo $3^{-2}=\frac{1}{9}$. $(-2)^{-3}=\frac{1}{(-2)^3}=-\frac{1}{8}$. $\left(\frac{1}{2}\right)^{-3}=2^3=8$. $\left(\frac{2}{3}\right)^{-2}=\left(\frac{3}{2}\right)^2=\frac{9}{4}$. $5^0=1$.\par
\apantisi $\frac{1}{9}$, $-\frac{1}{8}$, $8$, $\frac{9}{4}$, $1$.\par
\elegxos $\left(\frac{1}{2}\right)^{-3}\cdot\left(\frac{1}{2}\right)^{3}=\left(\frac{1}{2}\right)^0=1$ και $8\cdot\frac{1}{8}=1$.

\paradeigma{Με τις ιδιότητες}
\tixero Οι $2^{-3}\cdot2^5$, $10^4:10^{-2}$, $(3^{-1})^2$, $(2\cdot5)^{-2}$.\par
\tithelo Να τους υπολογίσω.\par
\possindeo $2^{-3+5}=2^2=4$. $10^{4-(-2)}=10^6$. $3^{-2}=\frac{1}{9}$. $2^{-2}\cdot5^{-2}=\frac{1}{4}\cdot\frac{1}{25}=\frac{1}{100}$.\par
\apantisi $4$, $10^6=1\,000\,000$, $\frac{1}{9}$, $\frac{1}{100}$.\par
\elegxos $(2\cdot5)^{-2}=10^{-2}=\frac{1}{100}$.

\paradeigma{Παράσταση}
\tixero Η παράσταση $\left(\frac{1}{2}\right)^{-2}+3^0-2^{-1}$.\par
\tithelo Να την υπολογίσω.\par
\possindeo $\left(\frac{1}{2}\right)^{-2}=2^2=4$, $3^0=1$, $2^{-1}=\frac{1}{2}$. Άρα $4+1-\frac{1}{2}=\frac{9}{2}$.\par
\apantisi $\frac{9}{2}$.\par
\elegxos $4{,}5=4+1-0{,}5$.

\paradeigma{Εξίσωση με δύναμη}
\tixero Η εξίσωση $2^x=\frac{1}{16}$.\par
\tithelo Να βρω το $x$.\par
\possindeo Γράφω $\frac{1}{16}=\frac{1}{2^4}=2^{-4}$. Άρα $x=-4$.\par
\apantisi $x=-4$.\par
\elegxos $2^{-4}=\frac{1}{16}$.

\tmima{Ασκήσεις}

\zoni{Το κατάλαβα}{αρνητικοί εκθέτες και μηδέν}

\ex Υπολόγισε: α) $2^{-1}$ \quad β) $2^{-3}$ \quad γ) $10^{-2}$ \quad δ) $5^{-2}$ \quad ε) $4^0$ \quad στ) $(-3)^0$.

\ex Γράψε ως δεκαδικούς: α) $10^{-1}$ \quad β) $10^{-3}$ \quad γ) $10^{-5}$.

\ex Υπολόγισε: α) $(-2)^{-3}$ \quad β) $(-2)^{-2}$ \quad γ) $(-1)^{-5}$ \quad δ) $\left(\frac{1}{2}\right)^{-3}$ \quad ε) $\left(\frac{2}{3}\right)^{-2}$ \quad στ) $\left(-\frac{3}{4}\right)^{-1}$.

\ex Γράψε ως δύναμη με αρνητικό εκθέτη: α) $\frac{1}{9}$ (με βάση $3$) \quad β) $\frac{1}{32}$ (βάση $2$) \quad γ) $0{,}0001$ (βάση $10$) \quad δ) $\frac{1}{(-5)^3}$.

\ex Υπολόγισε με τις ιδιότητες: α) $2^{-3}\cdot2^5$ \quad β) $10^4:10^{-2}$ \quad γ) $(3^{-1})^2$ \quad δ) $(2\cdot5)^{-2}$ \quad ε) $7^{-4}\cdot7^4$ \quad στ) $(-2)^{-3}\cdot(-2)^5$.

\ex Υπολόγισε: α) $\left(\frac{1}{2}\right)^{-2}+3^0-2^{-1}$ \quad β) $2^{-2}+2^{-1}$ \quad γ) $10^2\cdot10^{-5}\cdot10^3$ \quad δ) $(3^2)^{-1}$ \quad ε) $5^0\cdot(-2)^{-2}$.

\ex Βρες το $x$: α) $2^x=\frac{1}{16}$ \quad β) $3^x=\frac{1}{27}$ \quad γ) $10^x=0{,}001$ \quad δ) $\left(\frac{1}{2}\right)^x=8$ \quad ε) $5^x=1$.

\ex Σωστό ή λάθος; α) $2^{-1}=-2$. \quad β) $\alpha^0=1$ για $\alpha\neq0$. \quad γ) $\left(\frac{1}{2}\right)^{-1}=2$. \quad δ) $10^{-3}=0{,}003$. \quad ε) $\alpha^{-\nu}=\frac{1}{\alpha^\nu}$. \quad στ) $2^{-2}\cdot2^2=1$. \quad ζ) $(-2)^{-2}=-\frac{1}{4}$.

\zoni{Το χρησιμοποιώ}{μονάδες, μισά, αξίες}

\ex \textbf{Οι μονάδες.} Γράψε με δυνάμεις του $10$ σε μέτρα: $1$ χιλιόμετρο, $1$ εκατοστό, $1$ χιλιοστό, $1$ μικρόμετρο ($0{,}000001$ μ.), $1$ νανόμετρο ($0{,}000000001$ μ.).

\ex \textbf{Η πίτα.} Κόβουμε κάθε φορά το μισό από ό,τι έμεινε. Τι μέρος της πίτας έχει μείνει μετά από $5$ κοψίματα; Γράψε το ως δύναμη του $2$.

\ex \textbf{Το ραδιενεργό υλικό.} Ένα υλικό 64 g χάνει το μισό της μάζας του κάθε ημέρα. Πόσο θα είναι μετά από $3$ ημέρες; Πόσο ήταν πριν από $2$ ημέρες;

\ex \textbf{Ο πίνακας.} Συμπλήρωσε τις δυνάμεις του $3$ από $3^4$ έως $3^{-3}$.

\ex \textbf{Το κινητό.} Η αξία ενός κινητού των $400$ € μειώνεται στο μισό κάθε χρόνο. Πόση αξία έχει μετά από $3$ έτη; Γράψε το με δύναμη.

\zoni{Το σκέφτομαι}{γιατί το μηδέν δίνει ένα}

\ex Γιατί $\alpha^0=1$ και όχι $0$; Χρησιμοποίησε το μοτίβο των δυνάμεων του $2$.

\ex Δείξε ότι $\left(\frac{2}{3}\right)^{-2}=\left(\frac{3}{2}\right)^2$.

\ex Σε τι διαφέρουν τα $2^{-3}$ και $-2^3$;

\ex Αν ισχύει $\alpha^\mu\cdot\alpha^\nu=\alpha^{\mu+\nu}$ και $\mu=-3$, $\nu=3$, τι σημαίνει $\alpha^0$;

\ex Γιατί δεν ορίζεται το $0^{-1}$;

\ex Υπολόγισε τα $(-1)^{-1}$, $(-1)^{-2}$, $(-1)^{-3}$ και βρες τον κανόνα.

\zoni{Για το σπίτι}{πόσο λεπτό είναι ένα φύλλο;}

\ex Ένα φύλλο χαρτιού έχει πάχος περίπου $0{,}1$ mm. Γράψε το πάχος σε μέτρα με δύναμη του $10$.

\paragrafos{Τυποποιημένη μορφή μεγάλων και μικρών αριθμών}{Α.7.10, «Τυποποιημένη μορφή μεγάλων και μικρών αριθμών»}

\begin{anagke}
Η Α3 επισκέπτεται το εργαστήριο βιολογίας του σχολείου και η κυρία Αλεξάνδρα δείχνει ένα μικροσκόπιο. Κάτω από τον φακό υπάρχουν μια τρίχα, ένα ερυθρό αιμοσφαίριο και ένα βακτήριο. Η καθηγήτρια γράφει στον πίνακα τα μεγέθη: η τρίχα έχει πάχος $0{,}000\,07$ μ., το αιμοσφαίριο $0{,}000\,007$ μ., το βακτήριο $0{,}000\,002$ μ. και ένας ιός $0{,}000\,000\,1$ μ. Ο Στέλιος προσπαθεί να τα συγκρίνει και χάνεται στα μηδενικά: «Ποιο είναι μικρότερο, το $0{,}000\,007$ ή το $0{,}000\,002$;».

Η Πολυξένη θυμάται ότι για τους πολύ μεγάλους αριθμούς η κυρία Αλεξάνδρα τους είχε δείξει να τους γράφουν με δυνάμεις του $10$: $1{,}496\cdot10^{8}$ χιλιόμετρα η απόσταση Γης -- Ήλιου. «Μπορούμε να κάνουμε το ίδιο και για τους πολύ μικρούς;» ρωτάει. «Τώρα που ξέρουμε τους αρνητικούς εκθέτες, ναι», απαντάει η καθηγήτρια. Το $0{,}000\,07$ γράφεται $7\cdot10^{-5}$, και ο ιός $1\cdot10^{-7}$.

Η Μυρτώ προσθέτει μια ιστορική σημείωση: «Αυτά τα μεγέθη μπορούν να τα μετρήσουν χάρη στα επιστημονικά όργανα και στη γλώσσα των αριθμών. Η Υπατία, η μεγάλη μαθηματικός και αστρονόμος της Αλεξάνδρειας (περίπου 370--415 μ.Χ.), είχε δουλέψει σε όργανα μέτρησης του ουρανού, τον αστρολάβο». Η τυποποιημένη μορφή φέρνει τους πολύ μεγάλους και τους πολύ μικρούς αριθμούς στο ίδιο τετράδιο, με την ίδια γλώσσα.
\end{anagke}

\begin{lexi}{τυποποιημένη μορφή, τάξη μεγέθους}{τυποποιημένη μορφή, τάξη μεγέθους}
Όπως είδαμε, \lexiword{τυποποιημένη} είναι η μορφή που έχει το ίδιο «τύπο» για όλους. Η \lexiword{τάξη μεγέθους} ενός αριθμού είναι ο εκθέτης του $10$ στην τυποποιημένη του μορφή: μας λέει «πόσο μεγάλος ή πόσο μικρός» είναι, με μια ματιά.
\end{lexi}

\begin{metafrasi}
Το πάχος της τρίχας. & $0{,}000\,07=7\cdot10^{-5}$ μ.\\
Το αιμοσφαίριο. & $0{,}000\,007=7\cdot10^{-6}$ μ.\\
Ο ιός. & $0{,}000\,000\,1=1\cdot10^{-7}$ μ.\\
Η ταχύτητα του φωτός. & $300\,000\,000=3\cdot10^{8}$ μ./s\\
Η οθόνη του κομπιουτεράκι. & $4.5\mathrm{E}{-}4=4{,}5\cdot10^{-4}$\\
\end{metafrasi}

\tmima{Η θεωρία}

\begin{orismos}{τυποποιημένη μορφή}
Κάθε θετικός ρητός μπορεί να γραφεί στη μορφή $\alpha\cdot10^{\nu}$ με $1\le\alpha<10$ και $\nu$ \textbf{ακέραιο}. Αυτή λέγεται \textbf{τυποποιημένη μορφή}. Για τους μεγάλους αριθμούς ο $\nu$ είναι θετικός, για τους μικρούς ($<1$) αρνητικός.
\end{orismos}

\begin{kanonas}{τυποποίηση}
\begin{itemize}[leftmargin=1.2em,itemsep=0.3ex,topsep=2pt]
\item Μεταφέρω την υποδιαστολή ώστε να μείνει ένα μη μηδενικό ψηφίο πριν από αυτήν και βρίσκω τον $\alpha$.
\item Ο εκθέτης $\nu$ είναι το πλήθος των θέσεων που μετακίνησα, με πρόσημο $+$ αν μετακίνησα την υποδιαστολή αριστερά (αριθμός μεγάλος) και $-$ αν τη μετακίνησα δεξιά (αριθμός μικρός).
\end{itemize}
\end{kanonas}

Έτσι $0{,}00045$: μετακινώ $4$ θέσεις δεξιά και παίρνω $4{,}5$: $0{,}00045=4{,}5\cdot10^{-4}$. Και $3\,420\,000=3{,}42\cdot10^{6}$. Για να γυρίσω πίσω, πολλαπλασιάζω τον $\alpha$ με το $10^\nu$: $3{,}2\cdot10^{-5}=0{,}000032$.

\tmimasmall{Σύγκριση και πράξεις}

Συγκρίνω πρώτα τους εκθέτες: μεγαλύτερος εκθέτης σημαίνει μεγαλύτερο αριθμό ($4\cdot10^{-5}>7\cdot10^{-6}$ αφού $-5>-6$). Αν οι εκθέτες είναι ίσοι συγκρίνω τους $\alpha$. Για πράξεις: στον πολλαπλασιασμό πολλαπλασιάζω τους $\alpha$ και προσθέτω τους εκθέτες, στη διαίρεση διαιρώ τους $\alpha$ και αφαιρώ τους εκθέτες, και μετά ξαναγράφω σε τυποποιημένη μορφή αν χρειάζεται:
$$(3\cdot10^{-3})(2\cdot10^{5})=6\cdot10^{2},\qquad(6\cdot10^{8}):(2\cdot10^{-2})=3\cdot10^{10},\qquad(5\cdot10^{4})(4\cdot10^{-2})=20\cdot10^{2}=2\cdot10^{3}.$$

\begin{prosoxi}
\begin{itemize}[leftmargin=1.2em,itemsep=0.3ex,topsep=2pt]
\item Ο $\alpha$ είναι πάντα από $1$ έως λιγότερο από $10$. Το $12\cdot10^{-4}$ δεν είναι τυποποιημένο: είναι $1{,}2\cdot10^{-3}$.
\item Στους μικρούς αριθμούς ο εκθέτης είναι \emph{αρνητικός}: το $0{,}00045$ δεν είναι $4{,}5\cdot10^{4}$ αλλά $4{,}5\cdot10^{-4}$.
\item Ο εκθέτης $-4$ δεν σημαίνει «τέσσερα μηδενικά» αλλά «η υποδιαστολή πήγε τέσσερις θέσεις δεξιά»: μετά την υποδιαστολή υπάρχουν $3$ μηδενικά και μετά το $4$.
\item Στη σύγκριση: $10^{-5}$ είναι μεγαλύτερο από $10^{-6}$. Όσο πιο αρνητικός ο εκθέτης, τόσο μικρότερος ο αριθμός.
\end{itemize}
\end{prosoxi}

\tmima{Λυμένα παραδείγματα}

\paradeigma{Μικροί αριθμοί}
\tixero Οι $0{,}00045$, $0{,}000\,000\,7$, $0{,}0302$.\par
\tithelo Να τους γράψω σε τυποποιημένη μορφή.\par
\possindeo Μετακινώ την υποδιαστολή δεξιά μέχρι ένα μη μηδενικό ψηφίο πριν από αυτήν. $0{,}00045$: $4$ θέσεις, $4{,}5\cdot10^{-4}$. $0{,}000\,000\,7$: $7$ θέσεις, $7\cdot10^{-7}$. $0{,}0302$: $2$ θέσεις, $3{,}02\cdot10^{-2}$.\par
\apantisi $4{,}5\cdot10^{-4}$, $7\cdot10^{-7}$, $3{,}02\cdot10^{-2}$.\par
\elegxos $4{,}5\cdot10^{-4}=\frac{4{,}5}{10\,000}=0{,}00045$.

\paradeigma{Από τυποποιημένη σε δεκαδική}
\tixero Οι $3{,}2\cdot10^{-5}$, $9{,}1\cdot10^{-3}$, $2{,}5\cdot10^{6}$.\par
\tithelo Να τους γράψω κανονικά.\par
\possindeo Αρνητικός εκθέτης: υποδιαστολή αριστερά. $3{,}2\cdot10^{-5}=0{,}000032$. $9{,}1\cdot10^{-3}=0{,}0091$. Θετικός εκθέτης: υποδιαστολή δεξιά. $2{,}5\cdot10^6=2\,500\,000$.\par
\apantisi $0{,}000032$, $0{,}0091$, $2\,500\,000$.\par
\elegxos $3{,}2\cdot10^{-5}$: $5$ θέσεις αριστερά από το $3$: $0{,}000032$.

\paradeigma{Σύγκριση μεγεθών}
\tixero Η τρίχα $7\cdot10^{-5}$ μ., το αιμοσφαίριο $7\cdot10^{-6}$ μ., το βακτήριο $2\cdot10^{-6}$ μ., ο ιός $1\cdot10^{-7}$ μ.\par
\tithelo Να τα βάλω σε σειρά από το μικρότερο και να βρω πόσες φορές η τρίχα είναι χοντρότερη από το αιμοσφαίριο.\par
\possindeo Συγκρίνω εκθέτες: $-7<-6<-5$. Ο ιός ο μικρότερος. Μετά τα δύο με εκθέτη $-6$: $2<7$, δηλαδή το βακτήριο και μετά το αιμοσφαίριο. Τελευταία η τρίχα. Ο λόγος: $\frac{7\cdot10^{-5}}{7\cdot10^{-6}}=10^{-5-(-6)}=10$.\par
\apantisi $1\cdot10^{-7}<2\cdot10^{-6}<7\cdot10^{-6}<7\cdot10^{-5}$. Η τρίχα είναι $10$ φορές χοντρότερη.\par
\elegxos $0{,}00007=10\cdot0{,}000007$.

\paradeigma{Πράξεις}
\tixero Οι $(3\cdot10^{-3})\cdot(2\cdot10^{5})$, $(6\cdot10^{8}):(2\cdot10^{-2})$, $(5\cdot10^{4})\cdot(4\cdot10^{-2})$.\par
\tithelo Να τις υπολογίσω σε τυποποιημένη μορφή.\par
\possindeo $3\cdot2=6$ και $-3+5=2$: $6\cdot10^2$. $6:2=3$ και $8-(-2)=10$: $3\cdot10^{10}$. $5\cdot4=20$ και $4-2=2$: $20\cdot10^2=2\cdot10^3$.\par
\apantisi $6\cdot10^{2}$, $3\cdot10^{10}$, $2\cdot10^{3}$.\par
\elegxos Το τρίτο: $50\,000\cdot0{,}04=2\,000$.

\paradeigma{Το φως του Ήλιου}
\tixero Το φως ταξιδεύει με $3\cdot10^{8}$ μ./s και ο Ήλιος απέχει $1{,}5\cdot10^{11}$ μ.\par
\tithelo Σε πόσο χρόνο φτάνει το φως στη Γη.\par
\possindeo $\frac{1{,}5\cdot10^{11}}{3\cdot10^{8}}=0{,}5\cdot10^{3}=5\cdot10^{2}$ s.\par
\apantisi $500$ s, περίπου $8{,}3$ λεπτά.\par
\elegxos $500\cdot3\cdot10^8=1{,}5\cdot10^{11}$.

\tmima{Ασκήσεις}

\zoni{Το κατάλαβα}{μικροί και μεγάλοι αριθμοί}

\ex Γράψε σε τυποποιημένη μορφή: α) $0{,}0005$ \quad β) $0{,}000\,072$ \quad γ) $0{,}3$ \quad δ) $45\,000$ \quad ε) $0{,}000\,000\,8$ \quad στ) $6\,300\,000$.

\ex Γράψε σε δεκαδική μορφή: α) $3{,}2\cdot10^{-5}$ \quad β) $9{,}1\cdot10^{-3}$ \quad γ) $4\cdot10^{-1}$ \quad δ) $2{,}5\cdot10^{6}$ \quad ε) $1{,}01\cdot10^{-4}$ \quad στ) $7\cdot10^{0}$.

\ex Γράψε σωστά σε τυποποιημένη μορφή: α) $0{,}3\cdot10^{-2}$ \quad β) $12\cdot10^{-4}$ \quad γ) $5{,}5\cdot10^{-3}$ \quad δ) $10\cdot10^{2}$.

\ex Σύγκρινε: α) $4\cdot10^{-5}$ και $7\cdot10^{-6}$ \quad β) $2{,}5\cdot10^{-3}$ και $2{,}5\cdot10^{-2}$ \quad γ) $6\cdot10^{4}$ και $9\cdot10^{-4}$ \quad δ) $3{,}1\cdot10^{-8}$ και $3{,}09\cdot10^{-8}$.

\ex Γράψε σε αύξουσα σειρά: $3\cdot10^{-3}$, $5\cdot10^{-5}$, $2\cdot10^{2}$, $9\cdot10^{-1}$, $1{,}2\cdot10^{-3}$.

\ex Υπολόγισε σε τυποποιημένη μορφή: α) $(3\cdot10^{-3})(2\cdot10^{5})$ \quad β) $(6\cdot10^{8}):(2\cdot10^{-2})$ \quad γ) $(5\cdot10^{4})(4\cdot10^{-2})$ \quad δ) $(8\cdot10^{-6}):(4\cdot10^{-3})$.

\ex Ο υπολογιστής δείχνει $4.5\mathrm{E}{-}4$, $2.5\mathrm{E}7$ και $7\mathrm{E}{-}9$. Ποιοι αριθμοί είναι;

\ex Σωστό ή λάθος; α) $0{,}004=4\cdot10^{-3}$. \quad β) Το $12\cdot10^{-4}$ είναι τυποποιημένο. \quad γ) Το $10^{-5}$ είναι μικρότερο από το $10^{-6}$. \quad δ) $2\cdot10^{-3}=0{,}002$. \quad ε) Όσο μεγαλύτερος ο εκθέτης, τόσο μεγαλύτερος ο αριθμός. \quad στ) $5{,}5\cdot10^{-3}>5{,}5\cdot10^{-2}$.

\zoni{Το χρησιμοποιώ}{ο μικρόκοσμος και ο μακρόκοσμος}

\ex \textbf{Ο μικρόκοσμος.} Τρίχα $7\cdot10^{-5}$ μ., αιμοσφαίριο $7\cdot10^{-6}$ μ., βακτήριο $2\cdot10^{-6}$ μ., ιός $10^{-7}$ μ. Βάλε τα σε αύξουσα σειρά και βρες πόσες φορές η τρίχα είναι χοντρότερη από το αιμοσφαίριο.

\ex \textbf{Το ηλεκτρόνιο.} Η μάζα του ηλεκτρονίου είναι $9{,}1\cdot10^{-31}$ kg. Πόσα μηδενικά υπάρχουν μετά την υποδιαστολή πριν από το $9$ αν τη γράψεις δεκαδικά;

\ex \textbf{Το φως.} Το φως ταξιδεύει με $3\cdot10^{8}$ μ./s και ο Ήλιος απέχει από τη Γη $1{,}5\cdot10^{11}$ μ. Σε πόσο χρόνο φτάνει το φως στη Γη;

\ex \textbf{Το DNA.} Το πλάτος του DNA είναι $2$ nm και $1\ \mathrm{nm}=10^{-9}$ μ. Γράψε το σε μέτρα σε τυποποιημένη μορφή.

\ex \textbf{Το χαρτί.} Ένα φύλλο χαρτιού έχει πάχος $10^{-4}$ μ. Πόσο είναι το πάχος $500$ φύλλων σε μέτρα και σε εκατοστά;

\ex \textbf{Η άμμος.} Ένας κόκκος άμμου έχει μάζα $10^{-6}$ kg. Πόσοι κόκκοι ζυγίζουν $1$ kg;

\zoni{Το σκέφτομαι}{ο εκθέτης που μικραίνει}

\ex Γιατί ο εκθέτης είναι αρνητικός για αριθμούς μικρότερους από το $1$;

\ex Ποια είναι η τυποποιημένη μορφή του $1$; Και του $10$;

\ex Σύγκρινε τα $2\cdot10^{-3}$ και $5\cdot10^{-4}$ με δύο τρόπους (δεκαδικά και εκθέτες).

\ex Αν πολλαπλασιάσεις το $3\cdot10^{-4}$ με το $10^{6}$ τι παίρνεις;

\ex Γιατί το $12\cdot10^{-4}$ δεν είναι τυποποιημένο; Γράψε το σωστά.

\zoni{Για το σπίτι}{ψάξε τα μεγέθη}

\ex Βρες το μέγεθος ενός ατόμου, ενός ιού και ενός βακτηρίου από αξιόπιστη πηγή και γράψε τα σε τυποποιημένη μορφή και σε αύξουσα σειρά.

\epanalipsi{7}

\tmima{Λεξικό του κεφαλαίου}

\renewcommand{\arraystretch}{1.5}
\begin{tabular*}{\linewidth}{@{\extracolsep{\fill}} >{\raggedright\arraybackslash}p{0.28\linewidth} >{\raggedright\arraybackslash}p{0.62\linewidth}}
\toprule
\textbf{Λέξη ή τύπος} & \textbf{Με μία πρόταση}\\
\midrule
ρητοί & Φυσικοί, κλάσματα και δεκαδικοί μαζί με τους αντίστοιχους αρνητικούς.\\
$|\alpha|$ & Η απόσταση του $\alpha$ από το $0$ στον άξονα: ποτέ αρνητική.\\
αντίθετοι & Ίδια απόλυτη τιμή, ετερόσημοι: $\alpha$ και $-\alpha$, με άθροισμα $0$.\\
σύγκριση & Μεγαλύτερος είναι ο δεξιότερος στον άξονα.\\
$\alpha-\beta$ & $=\alpha+(-\beta)$: προσθέτω τον αντίθετο.\\
γινόμενο, πηλίκο & Ομόσημοι: θετικό. Ετερόσημοι: αρνητικό. Όχι διαίρεση με το $0$.\\
περιοδικός & Δεκαδικός με ψηφία που επαναλαμβάνονται, π.χ. $\frac{1}{3}=0{,}\overline{3}$.\\
$\alpha^{\nu}$, $\alpha^0$, $\alpha^{-\nu}$ & $\alpha\cdots\alpha$, $1$ και $\frac{1}{\alpha^\nu}$.\\
$\alpha\cdot10^{\nu}$ & Τυποποιημένη μορφή: $1\le\alpha<10$, $\nu$ ακέραιος.\\
\bottomrule
\end{tabular*}

\tmima{Ασκήσεις επανάληψης}

\ex Εξέφρασε με ρητό αριθμό: α) $8$ βαθμοί κάτω από το μηδέν \quad β) χρέος $25$ € \quad γ) βάθος $40$ μ. κάτω από την επιφάνεια της θάλασσας.

\ex Ποιοι από τους $-3$, $\frac{4}{2}$, $0{,}5$, $-\frac{7}{7}$, $0$ είναι ακέραιοι;

\ex Βρες: α) $|-7{,}5|$ \quad β) $\left|\frac{3}{4}\right|$ \quad γ) τον αντίθετο του $-\frac{2}{3}$ \quad δ) τους $\alpha$ με $|\alpha|=4$.

\ex Γράψε σε αύξουσα σειρά: $-2{,}5$, $0$, $-\frac{1}{2}$, $3$, $-3$.

\ex Υπολόγισε: α) $(-7)+(+12)$ \quad β) $(-3{,}5)+(-1{,}5)$ \quad γ) $\left(-\frac{1}{2}\right)+\left(+\frac{1}{3}\right)$ \quad δ) $(+8)+(-8)$.

\ex Υπολόγισε: α) $5-(-3)$ \quad β) $(-4)-(+6)$ \quad γ) $(-2{,}5)-(-2{,}5)$ \quad δ) $\left(-\frac{1}{2}\right)-\frac{1}{4}$.

\ex Υπολόγισε το αλγεβρικό άθροισμα $-12+7-(-5)-(3-9)$.

\ex Υπολόγισε: α) $(-6)\cdot(-7)$ \quad β) $(+4)\cdot(-2{,}5)$ \quad γ) $\left(-\frac{1}{2}\right)(-4)(-3)$ \quad δ) $(-2)\cdot(5-8)$.

\ex Υπολόγισε: α) $(-48):(-6)$ \quad β) $15:(-5)$ \quad γ) $\left(-\frac{3}{4}\right):\left(-\frac{9}{2}\right)$ \quad δ) $0:(-7)$.

\ex Γράψε σε δεκαδική μορφή τα $\frac{5}{6}$ και $\frac{7}{20}$ και ως κλάσμα τον $0{,}\overline{18}$.

\ex Υπολόγισε: α) $(-2)^5$ \quad β) $-2^4$ \quad γ) $\left(\frac{1}{3}\right)^{-2}$ \quad δ) $5^0+2^{-3}$.

\ex Υπολόγισε με ιδιότητες: α) $2^5\cdot2^{-3}$ \quad β) $(3^{-1})^2$ \quad γ) $10^4:10^6$.

\ex Γράψε σε τυποποιημένη μορφή τα $0{,}00032$ και $45\,000\,000$, και υπολόγισε $(4\cdot10^{3})(5\cdot10^{-5})$.

\ex \textbf{Η πτώση.} Η θερμοκρασία πέφτει από $+4$° σε $-5$° σε $3$ ώρες. Ποιος είναι ο ρυθμός μεταβολής ανά ώρα;

\ex \textbf{Ερωτήσεις αυτοαξιολόγησης (πράξεις).} Σωστό ή λάθος; α) $|-5|=-5$. \quad β) $(-3)+(-4)=-7$. \quad γ) $5-(-2)=3$. \quad δ) $(-2)(-3)(-4)=-24$. \quad ε) $12:(-4)=3$. \quad στ) Ο αντίθετος του $0$ είναι το $0$.

\ex \textbf{Ερωτήσεις αυτοαξιολόγησης (δυνάμεις, δεκαδικοί).} Σωστό ή λάθος; α) $(-2)^0=1$. \quad β) $2^{-1}=-2$. \quad γ) Το $\frac{1}{7}$ είναι περιοδικός δεκαδικός. \quad δ) $10^{-3}>10^{-4}$. \quad ε) $(\alpha+\beta)^2=\alpha^2+\beta^2$. \quad στ) $3{,}2\cdot10^{-2}=0{,}032$.

\tmima{Ο χάρτης της επανάληψης}

Μετά τις λύσεις, δες σε ποιες ασκήσεις δυσκολεύτηκες. Ο πίνακας σου λέει ποια παράγραφο να ξαναδιαβάσεις.

\renewcommand{\arraystretch}{1.5}
\begin{tabular*}{\linewidth}{@{\extracolsep{\fill}} >{\raggedright\arraybackslash}p{0.38\linewidth} >{\raggedright\arraybackslash}p{0.52\linewidth}}
\toprule
\textbf{Αν δυσκολεύτηκες στις} & \textbf{ξαναδιάβασε τις παραγράφους}\\
\midrule
Ε7.1, Ε7.2 & 7.1 (ρητοί, άξονας)\\
Ε7.3, Ε7.4 & 7.2 (απόλυτη τιμή, αντίθετοι, σύγκριση)\\
Ε7.5 & 7.3 (πρόσθεση)\\
Ε7.6, Ε7.7 & 7.4 (αφαίρεση, παρενθέσεις)\\
Ε7.8 & 7.5 (πολλαπλασιασμός)\\
Ε7.9, Ε7.14 & 7.6 (διαίρεση)\\
Ε7.10 & 7.7 (δεκαδική μορφή)\\
Ε7.11, Ε7.12 & 7.8 και 7.9 (δυνάμεις)\\
Ε7.13 & 7.10 (τυποποιημένη μορφή)\\
Ε7.15, Ε7.16 & 7.1 έως 7.10\\
\bottomrule
\end{tabular*}


\meros{Β}{ΓΕΩΜΕΤΡΙΑ}{Β}
\kefalaio{Βασικές γεωμετρικές έννοιες}{Πριν χτίσουμε, πρέπει να ονομάσουμε τα υλικά}{%
\begin{tabular}{@{}l@{\hspace{1em}}p{0.82\linewidth}@{}}
1.1 & Σημείο, ευθύγραμμο τμήμα, ευθεία, ημιευθεία, επίπεδο, ημιεπίπεδο\\[2pt]
1.2 & Γωνία, γραμμή, επίπεδα και ευθύγραμμα σχήματα, ίσα σχήματα\\[2pt]
1.3 & Μέτρηση, σύγκριση και ισότητα τμημάτων, απόσταση, μέσο\\[2pt]
1.4 & Πρόσθεση και αφαίρεση ευθυγράμμων τμημάτων\\[2pt]
1.5 & Μέτρηση, σύγκριση και ισότητα γωνιών, διχοτόμος\\[2pt]
1.6 & Είδη γωνιών, κάθετες ευθείες\\[2pt]
1.7 & Εφεξής και διαδοχικές γωνίες, άθροισμα γωνιών\\[2pt]
1.8 & Παραπληρωματικές, συμπληρωματικές και κατακορυφήν γωνίες\\[2pt]
1.9 & Θέσεις ευθειών στο επίπεδο\\[2pt]
1.10 & Απόσταση σημείου από ευθεία, απόσταση παραλλήλων\\[2pt]
1.11 & Κύκλος και στοιχεία του κύκλου\\[2pt]
1.12 & Επίκεντρη γωνία, τόξο, μέτρηση τόξου\\[2pt]
1.13 & Θέσεις ευθείας και κύκλου\\[2pt]
 & Επανάληψη του κεφαλαίου\\
\end{tabular}}
\paragrafos{Βασικές έννοιες: σημείο, ευθεία, ημιευθεία, ευθύγραμμο τμήμα, γωνία, κύκλος}{Β.1.1, «Βασικές έννοιες: σημείο, ευθεία, ημιευθεία, ευθύγραμμο τμήμα, γωνία, κύκλος»}

\begin{anagke}
Η κυρία Αλεξάνδρα ανακοινώνει ένα μεγάλο εγχείρημα για τη Γεωμετρία της χρονιάς: «Ο κήπος του σχολείου θα ανανεωθεί και η Α3 θα τον σχεδιάσει. Θα φέρουμε τα δέντρα, τα μονοπάτια, τον φράχτη, ένα στρογγυλό παρτέρι. Αλλά πρώτα πρέπει να μάθουμε να μιλάμε για αυτά που θα σχεδιάσουμε». Η Αγγελική, που βλέπει σχήματα παντού, χαίρεται: «Επιτέλους ένα μάθημα για αυτό που κάνω όλη μέρα!».

Η Αγγελική παίρνει ένα χαρτί και κάνει μια κάτοψη του κήπου. Κάθε δέντρο το σημειώνει με μια τελεία. Το μονοπάτι ανάμεσα σε δύο δέντρα το ενώνει με ένα τεντωμένο νήμα. Ο παππούς Γιάννης, που έρχεται να βοηθήσει, λέει: «Για να φυτέψετε σε ευθεία γραμμή, τεντώστε σχοινί και προχωρήστε όσο θέλετε προς τις δύο μεριές». Ο Στέλιος ρωτάει: «Μέχρι πού; Ο φράχτης τελειώνει στον τοίχο!». Η Ρήνη δείχνει τον φακό της: η δέσμη του φωτός ξεκινά από ένα σημείο και συνεχίζει προς μία μεριά χωρίς να τελειώνει. Για το στρογγυλό παρτέρι ο παππούς καρφώνει ένα παλούκι, δένει ένα σχοινί και γυρίζει τεντωμένο γύρω του. «Όλα τα σημεία της γραμμής απέχουν το ίδιο από το παλούκι», παρατηρεί η Φανή και το γράφει στο τετράδιό της.

Κάθε μία από αυτές τις εικόνες είναι μια γεωμετρική έννοια: σημείο, ευθύγραμμο τμήμα, ευθεία, ημιευθεία, γωνία, κύκλος. Σε αυτή την παράγραφο τις ονομάζουμε προσεκτικά και μαθαίνουμε να τις γράφουμε με σύμβολα.
\end{anagke}

\begin{lexi}{σημείο, ευθεία, επίπεδο, κύκλος}{σημείο, ευθεία, επίπεδο, κύκλος}
Το \lexiword{σημείο} είναι «αυτό που σημαδεύεται»: ένα σημάδι χωρίς διαστάσεις. Η \lexiword{ευθεία} πηγαίνει «ευθύ», χωρίς να στρίβει. Το \lexiword{επίπεδο} είναι ό,τι είναι «επί του πέδου», ίσιο σαν το δάπεδο. Το \lexiword{ευθύγραμμο τμήμα} είναι ένα «κομμένο» κομμάτι ευθείας. Το \lexiword{ημι-} σημαίνει «μισό»: ημιευθεία, ημιεπίπεδο. Ο \lexiword{κύκλος} είναι ό,τι «κυλάει» και γυρίζει. Η \lexiword{ακτίνα} είναι η «αχτίδα» που ξεκινά από το κέντρο.
\end{lexi}

\begin{metafrasi}
Το ίχνος της μύτης του μολυβιού. & σημείο $\mathrm{Α}$\\
Τεντωμένη κλωστή με άκρα $\mathrm{Α}$ και $\mathrm{Β}$. & ευθύγραμμο τμήμα $\mathrm{ΑΒ}$\\
Το ίδιο τμήμα προεκτεινόμενο απεριόριστα και προς τις δύο μεριές. & ευθεία $\mathrm{ΑΒ}$ ή $\varepsilon$ ή $x'x$\\
Φως φακού που ξεκινά από ένα σημείο. & ημιευθεία $\mathrm{Αx}$\\
Δύο ημιευθείες με κοινή αρχή $\mathrm{Ο}$. & γωνία $x\hat{\mathrm{Ο}}y$\\
Τα σημεία που απέχουν $\rho$ από το $\mathrm{Ο}$. & κύκλος $(\mathrm{Ο},\rho)$\\
\end{metafrasi}

\tmima{Η θεωρία}

\tmimasmall{Σημείο και ευθεία}

Ένα \textbf{σημείο} παριστάνεται με μία τελεία και δεν έχει διαστάσεις (ούτε μήκος ούτε πλάτος). Το ίχνος της μύτης του μολυβιού μας δίνει την έννοια του σημείου. Το συμβολίζουμε με ένα κεφαλαίο γράμμα: $\mathrm{Α}$, $\mathrm{Β}$, $\Gamma$.

\textbf{Ευθεία} είναι η γραμμή που εκτείνεται απεριόριστα, χωρίς αρχή και τέλος, και είναι απολύτως ίσια. Τη συμβολίζουμε με ένα μικρό γράμμα, π.χ. $\varepsilon$, ή με δύο γράμματα, π.χ. $x'x$.

\begin{idiotita}{ευθεία και σημεία}
Από ένα σημείο διέρχονται \textbf{άπειρες} ευθείες. Από δύο διαφορετικά σημεία διέρχεται \textbf{μία και μοναδική} ευθεία.
\end{idiotita}

\begin{center}
\begin{tikzpicture}[scale=0.9]
\begin{scope}
\foreach \a in {0,35,70,105,140,-35} {\draw[gblue] ({-1.5*cos(\a)},{-1.5*sin(\a)})--({1.5*cos(\a)},{1.5*sin(\a)});}
\fill (0,0) circle (2pt) node[below right] {$\mathrm{Α}$};
\node at (0,-1.9) {\small άπειρες ευθείες};
\end{scope}
\begin{scope}[xshift=5.5cm]
\draw[gblue,{Stealth}-{Stealth}] (-1.7,0)--(1.7,0);
\fill (-0.8,0) circle (2pt) node[above] {$\mathrm{Α}$}; \fill (0.8,0) circle (2pt) node[above] {$\mathrm{Β}$};
\node at (0,-0.7) {\small μία ευθεία};
\end{scope}
\end{tikzpicture}
\end{center}

\tmimasmall{Ευθύγραμμο τμήμα και ημιευθεία}

Αν σε μια ευθεία πάρουμε δύο σημεία $\mathrm{Α}$ και $\mathrm{Β}$, το σχήμα που αποτελείται από τα $\mathrm{Α}$, $\mathrm{Β}$ και όλα τα σημεία ανάμεσά τους λέγεται \textbf{ευθύγραμμο τμήμα} $\mathrm{ΑΒ}$ (ή $\mathrm{ΒΑ}$). Τα $\mathrm{Α}$ και $\mathrm{Β}$ είναι τα \textbf{άκρα} του. Το τμήμα έχει αρχή και τέλος, ενώ η ευθεία όχι.

Αν προεκτείνουμε το τμήμα $\mathrm{ΑΒ}$ και προς τα δύο άκρα του, σχηματίζεται ευθεία. Αν το προεκτείνουμε \emph{μόνο προς το ένα} άκρο, το σχήμα έχει αρχή αλλά δεν έχει τέλος και λέγεται \textbf{ημιευθεία}. Τη συμβολίζουμε $\mathrm{Αx}$, με πρώτο το γράμμα της αρχής.

\begin{center}
\begin{tikzpicture}[scale=0.95]
\draw[thick] (0,0)--(3,0); \fill (0,0) circle (2pt) node[below] {$\mathrm{Α}$}; \fill (3,0) circle (2pt) node[below] {$\mathrm{Β}$};
\node at (1.5,-0.9) {τμήμα $\mathrm{ΑΒ}$};
\begin{scope}[xshift=4.8cm]
\draw[thick,{Stealth}-{Stealth}] (-0.5,0)--(3.5,0); \fill (0.5,0) circle (2pt) node[below] {$\mathrm{Α}$}; \fill (2.7,0) circle (2pt) node[below] {$\mathrm{Β}$}; \node[above] at (3.2,0) {$\varepsilon$};
\node at (1.5,-0.9) {ευθεία $\varepsilon$};
\end{scope}
\begin{scope}[xshift=9.8cm]
\draw[thick,-{Stealth}] (0,0)--(3.2,0) node[above] {$x$}; \fill (0,0) circle (2pt) node[below] {$\mathrm{Α}$};
\node at (1.5,-0.9) {ημιευθεία $\mathrm{Αx}$};
\end{scope}
\end{tikzpicture}
\end{center}

Ένα σημείο $\mathrm{Α}$ μιας ευθείας $x'x$ τη χωρίζει σε δύο ημιευθείες, τις $\mathrm{Αx}$ και $\mathrm{Αx'}$, που έχουν κοινή αρχή το $\mathrm{Α}$. Οι ημιευθείες αυτές λέγονται \textbf{αντικείμενες}.

\begin{center}
\begin{tikzpicture}
\draw[thick,{Stealth}-{Stealth}] (-3,0)--(3,0); \fill (0,0) circle (2pt) node[above] {$\mathrm{Α}$};
\node[above] at (-2.8,0) {$x'$}; \node[above] at (2.8,0) {$x$};
\node[below,gblue] at (-1.5,0) {\small $\mathrm{Αx'}$}; \node[below,gblue] at (1.5,0) {\small $\mathrm{Αx}$};
\end{tikzpicture}
\end{center}

\tmimasmall{Γωνία}

Σχεδιάζουμε δύο ημιευθείες $\mathrm{Οx}$ και $\mathrm{Οy}$ με κοινή αρχή το σημείο $\mathrm{Ο}$. Οι ημιευθείες χωρίζουν το επίπεδο σε δύο περιοχές. Καθεμία από αυτές τις περιοχές, μαζί με τις δύο ημιευθείες, ονομάζεται \textbf{γωνία}. Το $\mathrm{Ο}$ είναι η \textbf{κορυφή} και οι ημιευθείες οι \textbf{πλευρές} της γωνίας.
\begin{itemize}[leftmargin=1.2em,itemsep=0.3ex,topsep=2pt]
\item Η μικρότερη περιοχή δίνει την \textbf{κυρτή} γωνία.
\item Η μεγαλύτερη περιοχή δίνει τη \textbf{μη κυρτή} γωνία.
\end{itemize}

\begin{center}
\begin{tikzpicture}[scale=1.1]
\fill[gblueL] (0,0)--(2.5,0) arc(0:38:2.5)--cycle;
\draw[thick] (2.9,0)--(0,0)--(2.3,1.8); \fill (0,0) circle (2pt) node[below left] {$\mathrm{Ο}$};
\node[right] at (2.9,0) {$x$}; \node[above right] at (2.3,1.8) {$y$};
\draw[gblue] (0.8,0) arc(0:38:0.8); \node[gblue] at (1.15,0.33) {\small κυρτή};
\draw[ggray] (-0.6,0) arc(180:322:0.6); \node[ggray] at (-1.15,-0.6) {\small μη κυρτή};
\end{tikzpicture}
\end{center}

Συμβολίζουμε τη γωνία με τρία γράμματα, γράφοντας στη μέση το γράμμα της κορυφής: $x\hat{\mathrm{Ο}}y$ ή $y\hat{\mathrm{Ο}}x$. Μπορούμε επίσης να τη συμβολίσουμε $\hat{\mathrm{Ο}}$ (με το γράμμα της κορυφής) ή με ένα μικρό ελληνικό γράμμα, π.χ. $\hat{\omega}$.

\textbf{Σημείωση.} Με τον γνώμονα κατασκευάζουμε ευθείες που λέγονται \textbf{κάθετες}. Η γωνία που σχηματίζεται λέγεται \textbf{ορθή}.

\tmimasmall{Επίπεδο και ημιεπίπεδο}

\textbf{Επίπεδο} είναι η επιφάνεια πάνω στην οποία εφαρμόζει παντού η ευθεία γραμμή. Το επίπεδο επεκτείνεται απεριόριστα και το ονομάζουμε με ένα κεφαλαίο γράμμα, π.χ. $\Pi$. Κάθε ευθεία $\varepsilon$ ενός επιπέδου το χωρίζει σε δύο \textbf{ημιεπίπεδα}.

\begin{center}
\begin{tikzpicture}[scale=0.95]
\fill[gblueL] (0,0)--(4.4,0)--(5.4,1.8)--(1,1.8)--cycle;
\fill[gwarm] (0,0)--(4.4,0)--(3.9,-0.9)--(-0.5,-0.9)--cycle;
\draw (-0.5,-0.9)--(3.9,-0.9)--(5.4,1.8)--(1,1.8)--cycle;
\draw[thick,gblue] (-0.25,-0.45)--(5.15,1.35) coordinate (e);
\node[gblue] at (5.5,1.5) {$\varepsilon$};
\node at (1.4,1.1) {$\Pi_1$}; \node at (3.5,-0.45) {$\Pi_2$};
\end{tikzpicture}
\end{center}

\tmimasmall{Κύκλος}

\begin{orismos}{κύκλος}
\textbf{Κύκλος} λέγεται το σύνολο όλων των σημείων του επιπέδου που απέχουν την ίδια απόσταση από ένα σταθερό σημείο $\mathrm{Ο}$.
\end{orismos}

Το σημείο $\mathrm{Ο}$ λέγεται \textbf{κέντρο} του κύκλου. Η σταθερή απόσταση λέγεται \textbf{ακτίνα} του κύκλου και συμβολίζεται με $\rho$. Έναν κύκλο με κέντρο $\mathrm{Ο}$ και ακτίνα $\rho$ τον γράφουμε $(\mathrm{Ο},\rho)$. Τον σχεδιάζουμε με τον \textbf{διαβήτη}. Οι κύκλοι που έχουν το ίδιο κέντρο λέγονται \textbf{ομόκεντροι}.

\begin{center}
\begin{tikzpicture}
\draw[thick] (0,0) circle (1.2); \draw[thick] (0,0) circle (1.9);
\fill (0,0) circle (2pt) node[below left] {$\mathrm{Ο}$};
\draw[gblue,thick] (0,0)--(0.85,0.85) node[midway,above left] {$\rho$};
\fill (0.85,0.85) circle (2pt) node[above right] {$\mathrm{Α}$};
\node at (0,-2.4) {\small δύο ομόκεντροι κύκλοι};
\end{tikzpicture}
\end{center}

\begin{prosoxi}
\begin{itemize}[leftmargin=1.2em,itemsep=0.3ex,topsep=2pt]
\item Το τμήμα $\mathrm{ΑΒ}$ και το $\mathrm{ΒΑ}$ είναι το ίδιο τμήμα. Αλλά οι ημιευθείες $\mathrm{Αx}$ και $\mathrm{Βx}$ είναι διαφορετικές, γιατί έχουν διαφορετική αρχή.
\item Το τμήμα έχει άκρα και μήκος. Η ευθεία και η ημιευθεία δεν έχουν δύο άκρα και δεν έχουν μήκος.
\item Στο όνομα της γωνίας το γράμμα της κορυφής γράφεται πάντα στη μέση: $x\hat{\mathrm{Ο}}y$, όχι $\mathrm{Ο}xy$.
\item Ο κύκλος είναι μόνο η γραμμή, όχι το εσωτερικό του. Όλες οι ακτίνες του ίδιου κύκλου είναι ίσες.
\end{itemize}
\end{prosoxi}

\tmima{Λυμένα παραδείγματα}

\paradeigma{Πόσες ευθείες;}
\tixero Τέσσερα δέντρα $\mathrm{Α},\,\mathrm{Β},\,\Gamma,\,\Delta$ του κήπου, από τα οποία τρία οποιαδήποτε δεν βρίσκονται στην ίδια ευθεία.\par
\tithelo Πόσες ευθείες ορίζουν.\par
\possindeo Κάθε ζεύγος σημείων ορίζει μία μοναδική ευθεία. Τα ζεύγη είναι $\mathrm{ΑΒ},\,\mathrm{ΑΓ},\,\mathrm{ΑΔ},\,\mathrm{ΒΓ},\,\mathrm{ΒΔ},\,\Gamma\Delta$. Με τον τρόπο των χειραψιών: $\dfrac{4\cdot3}{2}=6$.\par
\apantisi $6$ ευθείες.\par
\elegxos Αν το $\mathrm{Δ}$ ήταν πάνω στην ευθεία $\mathrm{ΑΒ}$, οι ευθείες θα ήταν λιγότερες.

\paradeigma{Ονόματα ημιευθειών}
\tixero Το σημείο $\mathrm{Ο}$ βρίσκεται πάνω στην ευθεία $x'x$. Τα $\mathrm{Α}$ και $\mathrm{Β}$ είναι δεξιά και αριστερά του $\mathrm{Ο}$ αντίστοιχα.\par
\tithelo Ποιες είναι οι δύο ημιευθείες που ορίζει το $\mathrm{Ο}$ και ποια περιέχει το $\mathrm{Α}$.\par
\possindeo Το $\mathrm{Ο}$ χωρίζει την ευθεία σε δύο ημιευθείες με αρχή το $\mathrm{Ο}$. Η μία πηγαίνει προς το $x$ και η άλλη προς το $x'$.\par
\apantisi Οι $\mathrm{Οx}$ (περιέχει το $\mathrm{Α}$) και $\mathrm{Οx'}$ (περιέχει το $\mathrm{Β}$), αντικείμενες.\par
\elegxos Αν τις ενώσω παίρνω ξανά ολόκληρη την ευθεία.

\paradeigma{Το όνομα μιας γωνίας}
\tixero Οι ημιευθείες $\mathrm{Οx}$ και $\mathrm{Οy}$ με κοινή αρχή $\mathrm{Ο}$ σχηματίζουν τη μικρότερη από τις δύο περιοχές του επιπέδου.\par
\tithelo Να ονομάσω τη γωνία με δύο τρόπους και να πω αν είναι κυρτή.\par
\possindeo Η κορυφή είναι το $\mathrm{Ο}$ και μπαίνει στη μέση: $x\hat{\mathrm{Ο}}y$ ή $y\hat{\mathrm{Ο}}x$. Επίσης μπορώ να γράψω απλώς $\hat{\mathrm{Ο}}$.\par
\apantisi Η γωνία $x\hat{\mathrm{Ο}}y$ είναι κυρτή, γιατί αντιστοιχεί στη μικρότερη περιοχή.\par
\elegxos Η μεγαλύτερη περιοχή θα έδινε τη μη κυρτή γωνία.

\paradeigma{Το μονοπάτι χωρίζει τον κήπο}
\tixero Ένα ευθύγραμμο μονοπάτι $\varepsilon$ χωρίζει τον κήπο σε δύο ημιεπίπεδα. Τα δέντρα $\mathrm{Κ}$ και $\Lambda$ είναι στο ένα και το $\mathrm{Μ}$ στο άλλο.\par
\tithelo Ποιες ευθύγραμμες διαδρομές ανάμεσα σε δύο δέντρα διασχίζουν το μονοπάτι.\par
\possindeo Για να πάω από το ένα ημιεπίπεδο στο άλλο πρέπει να περάσω από την ευθεία $\varepsilon$.\par
\apantisi Τα τμήματα $\mathrm{ΚΜ}$ και $\Lambda\mathrm{Μ}$ διασχίζουν το μονοπάτι, το $\mathrm{ΚΛ}$ όχι.\par
\elegxos Τα $\mathrm{Κ}$ και $\Lambda$ είναι στην ίδια πλευρά.

\paradeigma{Το στρογγυλό παρτέρι}
\tixero Το σχοινί του παππού είναι $1{,}5$ μ. και ο παππούς το κρατά τεντωμένο γύρω από το παλούκι $\mathrm{Ο}$. Ένα δέντρο $\mathrm{Δ}$ βρίσκεται $2$ μ. από το παλούκι.\par
\tithelo Ποιος κύκλος χαράσσεται και αν το δέντρο είναι πάνω στη γραμμή του.\par
\possindeo Το παλούκι είναι το κέντρο και το σχοινί η ακτίνα: κύκλος $(\mathrm{Ο},1{,}5)$ μ. Το δέντρο θα ήταν πάνω στον κύκλο αν $\mathrm{ΟΔ}=1{,}5$ μ.\par
\apantisi Ο κύκλος είναι ο $(\mathrm{Ο},1{,}5)$. Το δέντρο έχει $\mathrm{ΟΔ}=2\neq1{,}5$, άρα δεν είναι πάνω στη γραμμή.\par
\elegxos Το δέντρο απέχει περισσότερο από το σχοινί, άρα είναι έξω από το παρτέρι.

\tmima{Ασκήσεις}

\zoni{Το κατάλαβα}{σημεία, ευθείες, ημιευθείες, γωνίες}

\ex Συμπλήρωσε τα κενά. α) Το ίχνος της μύτης του μολυβιού μας δίνει την έννοια του \dots. β) Αν προεκτείνουμε ένα τμήμα και προς τα δύο άκρα του, προκύπτει \dots. γ) Από ένα σημείο διέρχονται \dots ευθείες, ενώ από δύο σημεία διέρχεται \dots ευθεία. δ) Αν προεκτείνουμε ένα τμήμα μόνο προς το ένα άκρο του, προκύπτει \dots. ε) Οι ημιευθείες που έχουν κοινή αρχή και σχηματίζουν ευθεία λέγονται \dots.

\ex Σχεδίασε τρία σημεία $\mathrm{Α},\,\mathrm{Β},\,\Gamma$ που δεν βρίσκονται στην ίδια ευθεία. Φέρε και ονόμασε τις ευθείες που ορίζουν και γράψε όλα τα ευθύγραμμα τμήματα που βλέπεις. Πόσες ευθείες και πόσα τμήματα είναι;

\ex Γράψε με σύμβολα: α) το τμήμα με άκρα $\mathrm{Κ}$ και $\Lambda$ \quad β) η ημιευθεία με αρχή το $\mathrm{Κ}$ που περνά από το $\Lambda$ \quad γ) η ευθεία που ορίζουν τα $\mathrm{Κ}$ και $\Lambda$.

\ex Σχεδίασε ένα σημείο $\mathrm{Ο}$ και δύο διαφορετικές ευθείες που να διέρχονται από αυτό. Ονόμασε τις αντικείμενες ημιευθείες με αρχή το $\mathrm{Ο}$. Πόσες ημιευθείες βλέπεις συνολικά;

\ex Σχεδίασε μια γωνία $x\hat{\mathrm{Ο}}y$ και γράψε το όνομά της με δύο διαφορετικούς τρόπους. Σημείωσε με χρώμα την κυρτή και με άλλο χρώμα τη μη κυρτή γωνία.

\ex Σχεδίασε με διαβήτη τον κύκλο $(\mathrm{Κ},2{,}5)$ cm και τρεις διαφορετικές ακτίνες του. Μέτρησέ τις. Τι παρατηρείς;

\ex Σωστό ή λάθος; α) Η ευθεία έχει δύο άκρα. \quad β) Τα τμήματα $\mathrm{ΑΒ}$ και $\mathrm{ΒΑ}$ είναι ίδια. \quad γ) Οι ημιευθείες $\mathrm{Αx}$ και $\mathrm{Βx}$ είναι ίδιες. \quad δ) Από δύο σημεία περνά μόνο μία ευθεία. \quad ε) Η γωνία έχει δύο ημιευθείες με κοινή αρχή. \quad στ) Οι ομόκεντροι κύκλοι έχουν την ίδια ακτίνα.

\zoni{Το χρησιμοποιώ}{ο κήπος, ο φακός, το στρογγυλό παρτέρι}

\ex \textbf{Η σειρά των δέντρων.} Ο παππούς Γιάννης θέλει να φυτέψει δέντρα σε ευθεία γραμμή. Πόσα παλούκια αρκεί να βάλει για να ορίσει την ευθεία και πώς θα την ελέγξει με ένα σχοινί;

\ex \textbf{Οι ακτίνες.} Ένας φακός φωτίζει με δέσμη που ξεκινά από το σημείο $\mathrm{Α}$ και περνά από το $\mathrm{Β}$. Ποια γεωμετρική έννοια περιγράφει η δέσμη; Αν δύο φακοί στέκονται αντικριστά και φωτίζουν ο ένας τον άλλον, τι σχέση έχουν οι δύο ημιευθείες;

\ex \textbf{Τα σχοινιά.} Πέντε παιδιά στέκονται σε σημεία του κήπου, από τα οποία τρία οποιαδήποτε δεν είναι στην ίδια ευθεία, και ενώνουν κάθε ζευγάρι με σχοινί. Πόσα σχοινιά χρειάζονται;

\ex \textbf{Το μονοπάτι.} Ένα ευθύγραμμο μονοπάτι $\varepsilon$ χωρίζει τον κήπο σε δύο ημιεπίπεδα. Τα δέντρα $\mathrm{Κ}$, $\Lambda$ είναι στην ίδια πλευρά και το $\mathrm{Μ}$ στην άλλη. Ποια από τα τμήματα $\mathrm{ΚΛ}$, $\mathrm{ΚΜ}$, $\Lambda\mathrm{Μ}$ τέμνουν το μονοπάτι;

\ex \textbf{Το παρτέρι.} Η Ρήνη χαράσσει με σχοινί $1{,}2$ μ. ένα στρογγυλό παρτέρι με κέντρο το παλούκι $\mathrm{Ο}$. Τρία φυτά απέχουν από το παλούκι $0{,}9$ μ., $1{,}2$ μ. και $1{,}6$ μ. Ποιο είναι πάνω στη γραμμή του κύκλου και ποιο εκτός του κύκλου;

\zoni{Το σκέφτομαι}{ευθεία ή τμήμα;}

\ex Γιατί η ευθεία δεν έχει μήκος αλλά το τμήμα έχει; Μπορείς να μετρήσεις το μήκος μιας ημιευθείας;

\ex Τρία σημεία $\mathrm{Α},\,\mathrm{Β},\,\Gamma$ βρίσκονται στην ίδια ευθεία με αυτή τη σειρά. Πόσες διαφορετικές ημιευθείες με αρχή ένα από αυτά τα σημεία ορίζονται πάνω στην ευθεία; (Μέτρα και τις δύο κατευθύνσεις.)

\ex Μπορούν δύο διαφορετικές ευθείες να περνούν και οι δύο από τα ίδια δύο σημεία $\mathrm{Α}$ και $\mathrm{Β}$; Εξήγησε με την ιδιότητα της ευθείας.

\ex Δύο ομόκεντροι κύκλοι έχουν ακτίνες $3$ cm και $5$ cm. Μπορεί να υπάρχει σημείο που ανήκει και στους δύο; Γιατί;

\zoni{Για το σπίτι}{γεωμετρία στο δωμάτιο}

\ex Βρες στο δωμάτιό σου πέντε παραδείγματα σημείου, ευθύγραμμου τμήματος και επιπέδου. Πού μοιάζει με ημιευθεία ένα αντικείμενο της καθημερινότητας;

\ex Με νήμα και δύο καρφίτσες σε ένα χαρτόνι δείξε ότι από δύο σημεία περνά μία ευθεία. Μετά σχεδίασε με διαβήτη έναν κύκλο που να περνά από ένα σημείο της επιλογής σου.

\paragrafos{Μήκος ευθύγραμμου τμήματος}{Β.1.2, «Μήκος ευθύγραμμου τμήματος»}

\begin{anagke}
Για να σχεδιάσουν τον κήπο, η παρέα μετράει πρώτα το οικόπεδο. Ο Στέλιος μετράει το μήκος του μονοπατιού με βήματα και βρίσκει $14$. Ο Ηλίας το μετράει με τα δικά του βήματα και βρίσκει $16$. Ποιος έχει δίκιο; «Και οι δύο», λέει η Πολυξένη: «το βήμα του Στέλιου είναι μεγαλύτερο. Έχουμε διαφορετικές μονάδες». Η Αγγελική παίρνει τη μετροταινία του παππού: $9{,}8$ μέτρα. «Αυτό είναι σίγουρο: όλοι θα βρούμε το ίδιο νούμερο».

Μετά η Ρήνη ρωτάει: «Θέλουμε το συντριβάνι να είναι ακριβώς στη μέση του μονοπατιού. Πώς βρίσκουμε τη μέση;». Ο Νίκος πιάνει ένα σχοινί με το μήκος του μονοπατιού και το διπλώνει στα δύο: «Εκεί που διπλώνει είναι η μέση!». Ο Βασίλης προσθέτει: «Και για τον φράχτη γύρω από το παρτέρι πρέπει να προσθέσουμε όλες τις πλευρές του». Η Φανή τα σημειώνει όλα στο τετράδιό της και ρωτάει: «Μα πόσα μέτρα φράχτη είναι αυτά, αν κάποιες πλευρές τις μετρήσαμε σε εκατοστά;».

Σε αυτή την παράγραφο μαθαίνουμε να μετράμε το μήκος ενός τμήματος με σωστές μονάδες, να μετατρέπουμε μονάδες, να βρίσκουμε το μέσο ενός τμήματος και να προσθέτουμε και να αφαιρούμε τμήματα. Έτσι θα υπολογίζουμε και την περίμετρο ενός σχήματος.
\end{anagke}

\begin{lexi}{απόσταση, μέσο, περίμετρος}{απόσταση, μέσο, περίμετρος}
Η \lexiword{απόσταση} δύο σημείων είναι το πόσο «απέχουν». Το \lexiword{μέσο} ενός τμήματος είναι το σημείο που βρίσκεται «στη μέση». Η \lexiword{περίμετρος} είναι το «γύρω μέτρημα» (περί + μέτρο): το συνολικό μήκος του περιγράμματος ενός σχήματος. Το \lexiword{μέτρο} στα αρχαία ελληνικά είναι το «μέτρον», η ποσότητα με την οποία μετράμε.
\end{lexi}

\begin{metafrasi}
Το μήκος του τμήματος $\mathrm{ΑΒ}$ είναι $6{,}4$ εκατοστά. & $\mathrm{ΑΒ}=6{,}4\,\text{cm}$\\
Τα δύο τμήματα έχουν το ίδιο μήκος. & $\mathrm{ΑΒ}=\Gamma\Delta$\\
Η απόσταση των $\mathrm{Α}$ και $\mathrm{Β}$. & το μήκος του $\mathrm{ΑΒ}$\\
Το $\mathrm{Μ}$ είναι το μέσο του $\mathrm{ΑΒ}$. & $\mathrm{ΑΜ}=\mathrm{ΜΒ}$\\
Το $\mathrm{Β}$ είναι ανάμεσα στα $\mathrm{Α}$ και $\Gamma$. & $\mathrm{ΑΓ}=\mathrm{ΑΒ}+\mathrm{ΒΓ}$\\
Το άθροισμα των πλευρών. & περίμετρος\\
\end{metafrasi}

\tmima{Η θεωρία}

\tmimasmall{Μέτρηση μήκους}

Για να συγκρίνουμε δύο μεγέθη, δηλαδή για να δούμε αν το ένα είναι μεγαλύτερο, μικρότερο ή ίσο με το άλλο, τα συγκρίνουμε με ένα σταθερό μέγεθος που λέγεται \textbf{μονάδα μέτρησης}.

\begin{kanonas}{μήκος ευθύγραμμου τμήματος}
Βρίσκουμε το \textbf{μήκος} ενός ευθύγραμμου τμήματος συγκρίνοντάς το με ένα άλλο ευθύγραμμο τμήμα που το λέμε \textbf{μονάδα μήκους}. Η βασική μονάδα μήκους που χρησιμοποιούμε είναι το \textbf{μέτρο} (m).
\end{kanonas}

Για μικρά μήκη χρησιμοποιούμε συνήθως το \textbf{υποδεκάμετρο}, που είναι το ένα δέκατο του μέτρου. Αν το τμήμα είναι πολύ μεγαλύτερο ή πολύ μικρότερο από το μέτρο, διαλέγουμε ένα πολλαπλάσιο ή μια υποδιαίρεση του μέτρου.

\tmimasmall{Μονάδες μέτρησης μήκους}

\renewcommand{\arraystretch}{1.25}
\begin{center}
\begin{tabular}{@{}lll@{}}
\toprule
 & \textbf{Ονομασία} & \textbf{Σύμβολο}\\
\midrule
Πολλαπλάσιο του μέτρου & Χιλιόμετρο & km\\
Βασική μονάδα & Μέτρο & m\\
Υποδιαιρέσεις του μέτρου & Δεκατόμετρο & dm\\
 & Εκατοστόμετρο & cm\\
 & Χιλιοστόμετρο & mm\\
\bottomrule
\end{tabular}
\end{center}

\begin{center}
\begin{tikzpicture}[x=1cm,y=1cm]
\foreach \x/\n in {0/km,2.6/m,5.2/dm,7.8/cm,10.4/mm}{\node[draw,rounded corners=3pt,minimum width=1.5cm,fill=gblueL] (\n) at (\x,0) {\n};}
\draw[-{Stealth},gblue] (km.north) .. controls +(0.4,0.6) and +(-0.4,0.6) .. node[above] {\small $\times1000$} (m.north);
\draw[-{Stealth},gblue] (m.north) .. controls +(0.4,0.6) and +(-0.4,0.6) .. node[above] {\small $\times10$} (dm.north);
\draw[-{Stealth},gblue] (dm.north) .. controls +(0.4,0.6) and +(-0.4,0.6) .. node[above] {\small $\times10$} (cm.north);
\draw[-{Stealth},gblue] (cm.north) .. controls +(0.4,0.6) and +(-0.4,0.6) .. node[above] {\small $\times10$} (mm.north);
\draw[-{Stealth},gwarmD] (mm.south) .. controls +(-0.4,-0.6) and +(0.4,-0.6) .. node[below] {\small $:10$} (cm.south);
\draw[-{Stealth},gwarmD] (cm.south) .. controls +(-0.4,-0.6) and +(0.4,-0.6) .. node[below] {\small $:10$} (dm.south);
\draw[-{Stealth},gwarmD] (dm.south) .. controls +(-0.4,-0.6) and +(0.4,-0.6) .. node[below] {\small $:10$} (m.south);
\draw[-{Stealth},gwarmD] (m.south) .. controls +(-0.4,-0.6) and +(0.4,-0.6) .. node[below] {\small $:1000$} (km.south);
\end{tikzpicture}
\end{center}

Από τη μεγαλύτερη προς τη μικρότερη μονάδα \textbf{πολλαπλασιάζουμε}: $1\,\text{m}=10\,\text{dm}=100\,\text{cm}=1000\,\text{mm}$ και $1\,\text{km}=1000\,\text{m}$. Από τη μικρότερη προς τη μεγαλύτερη \textbf{διαιρούμε}: $1\,\text{mm}=0{,}1\,\text{cm}=0{,}001\,\text{m}$ και $1\,\text{m}=0{,}001\,\text{km}$.

Στην Αγγλία, την Αμερική και σε μερικές ακόμη χώρες η βασική μονάδα μήκους είναι η \textbf{υάρδα} ή γιάρδα (yd). Η $1$ yd διαιρείται σε $3$ πόδια (ft) και ισούται περίπου με $0{,}9144$ m. Το $1$ πόδι διαιρείται σε $12$ ίντσες (in) και ισούται περίπου με $0{,}3048$ m.

\tmimasmall{Ίσα τμήματα, απόσταση, μέσο}

Δύο τμήματα είναι \textbf{ίσα} όταν έχουν το ίδιο μήκος, όταν δηλαδή το ένα εφαρμόζει ακριβώς στο άλλο (π.χ. με τον διαβήτη). Για να συγκρίνουμε δύο μήκη πρέπει να τα εκφράσουμε στην ίδια μονάδα.

\begin{orismos}{απόσταση δύο σημείων}
\textbf{Απόσταση} δύο σημείων $\mathrm{Α}$ και $\mathrm{Β}$ λέγεται το μήκος του ευθύγραμμου τμήματος $\mathrm{ΑΒ}$ που τα ενώνει.
\end{orismos}

\begin{orismos}{μέσο ευθύγραμμου τμήματος}
\textbf{Μέσο} ενός ευθύγραμμου τμήματος $\mathrm{ΑΒ}$ ονομάζουμε το σημείο $\mathrm{Μ}$ του τμήματος που απέχει εξίσου από τα άκρα του: $\mathrm{ΑΜ}=\mathrm{ΜΒ}$.
\end{orismos}

\begin{center}
\begin{tikzpicture}
\draw[thick] (0,0)--(6,0);
\foreach \x/\l in {0/Α,3/Μ,6/Β}{\fill (\x,0) circle (2pt); \node[below] at (\x,-0.1) {$\mathrm{\l}$};}
\node[gblue,above] at (1.5,0.05) {\small $2$ cm}; \node[gblue,above] at (4.5,0.05) {\small $2$ cm};
\end{tikzpicture}
\end{center}

Κάθε τμήμα έχει ένα μοναδικό μέσο, και $\mathrm{ΑΜ}=\mathrm{ΜΒ}=\dfrac{\mathrm{ΑΒ}}{2}$.

\tmimasmall{Άθροισμα και διαφορά τμημάτων}

Αν πάρουμε σε μια ευθεία με τη σειρά τα σημεία $\mathrm{Α}$, $\mathrm{Β}$ και $\Gamma$, τότε:
\begin{itemize}[leftmargin=1.2em,itemsep=0.3ex,topsep=2pt]
\item το τμήμα $\mathrm{ΑΓ}$ είναι το \textbf{άθροισμα} των $\mathrm{ΑΒ}$ και $\mathrm{ΒΓ}$: $\mathrm{ΑΓ}=\mathrm{ΑΒ}+\mathrm{ΒΓ}$.
\item το τμήμα $\mathrm{ΑΒ}$ είναι η \textbf{διαφορά} των $\mathrm{ΑΓ}$ και $\mathrm{ΒΓ}$: $\mathrm{ΑΒ}=\mathrm{ΑΓ}-\mathrm{ΒΓ}$.
\end{itemize}

\begin{center}
\begin{tikzpicture}
\draw[thick] (0,0)--(7,0);
\foreach \x/\l in {0/Α,2.5/Β,7/Γ}{\fill (\x,0) circle (2pt); \node[below] at (\x,-0.1) {$\mathrm{\l}$};}
\draw[gblue,thick] (0,0.4)--(2.5,0.4); \draw[gwarmD,thick] (2.5,0.4)--(7,0.4);
\node[gblue,above] at (1.25,0.45) {$\mathrm{ΑΒ}$}; \node[gwarmD,above] at (4.75,0.45) {$\mathrm{ΒΓ}$};
\draw[thick,ggray] (0,-0.9)--(7,-0.9); \node[ggray,below] at (3.5,-0.95) {$\mathrm{ΑΓ}=\mathrm{ΑΒ}+\mathrm{ΒΓ}$};
\end{tikzpicture}
\end{center}

Για να κατασκευάσουμε με τον διαβήτη ένα τμήμα με μήκος $\mathrm{ΑΒ}+\Gamma\Delta$, σχεδιάζουμε μια ευθεία, παίρνουμε πάνω της ένα τμήμα ίσο με το $\mathrm{ΑΒ}$ και στη συνέχεια, από το άκρο του, ένα τμήμα ίσο με το $\Gamma\Delta$. Για τη διαφορά $\mathrm{ΑΒ}-\Gamma\Delta$ (όταν $\mathrm{ΑΒ}>\Gamma\Delta$) παίρνουμε πάνω στο $\mathrm{ΑΒ}$, από το ένα άκρο του, τμήμα ίσο με το $\Gamma\Delta$ και το κομμάτι που μένει είναι η διαφορά.

\tmimasmall{Περίμετρος ευθύγραμμου σχήματος}

\textbf{Κυρτό} ονομάζεται ένα ευθύγραμμο σχήμα όταν η προέκταση κάθε πλευράς του αφήνει το σχήμα στο ίδιο ημιεπίπεδο. Διαφορετικά ονομάζεται \textbf{μη κυρτό}.

\begin{center}
\begin{tikzpicture}[scale=0.9]
\coordinate (A) at (0,0); \coordinate (B) at (2.5,0); \coordinate (C) at (3.1,1.8); \coordinate (D) at (1.2,2.8); \coordinate (E) at (-0.5,1.5);
\draw[thick,fill=gblueL] (A)--(B)--(C)--(D)--(E)--cycle;
\draw[dashed,ggray] ($(A)!-0.5!(B)$)--($(A)!1.7!(B)$);
\node at (1.3,-0.8) {\small κυρτό};
\begin{scope}[xshift=6cm]
\coordinate (P) at (0,0); \coordinate (Q) at (3,0); \coordinate (R) at (3,2.2); \coordinate (S) at (1.6,0.9); \coordinate (T) at (0,2.2);
\draw[thick,fill=gwarm] (P)--(Q)--(R)--(S)--(T)--cycle;
\draw[dashed,ggray] ($(S)!-0.7!(R)$)--($(S)!2.2!(R)$);
\node at (1.5,-0.8) {\small μη κυρτό};
\end{scope}
\end{tikzpicture}
\end{center}

\begin{orismos}{περίμετρος}
Το άθροισμα των πλευρών ενός ευθύγραμμου σχήματος λέγεται \textbf{περίμετρος} του σχήματος.
\end{orismos}

Για παράδειγμα, αν ένα τετράπλευρο έχει πλευρές $4$ cm, $2{,}6$ cm, $3{,}1$ cm και $3$ cm, η περίμετρός του είναι $4+2{,}6+3{,}1+3=12{,}7$ cm.

\begin{prosoxi}
\begin{itemize}[leftmargin=1.2em,itemsep=0.3ex,topsep=2pt]
\item Για να συγκρίνω ή να προσθέσω μήκη πρέπει πρώτα να τα φέρω στην ίδια μονάδα: τα $0{,}8$ m δεν είναι μικρότερα από τα $75$ cm, γιατί $0{,}8$ m $=80$ cm.
\item Στις μετατροπές μεγάλη $\to$ μικρή μονάδα πολλαπλασιάζω, και μικρή $\to$ μεγάλη διαιρώ. Αν ο αριθμός που βγάζω είναι μεγαλύτερος για τη μεγαλύτερη μονάδα, κάτι πήγε στραβά.
\item Το μέσο \emph{ανήκει} στο τμήμα: δεν αρκεί να απέχει εξίσου από τα άκρα.
\item Η σχέση $\mathrm{ΑΓ}=\mathrm{ΑΒ}+\mathrm{ΒΓ}$ ισχύει όταν το $\mathrm{Β}$ βρίσκεται \emph{ανάμεσα} στα $\mathrm{Α}$ και $\Gamma$ πάνω στην ίδια ευθεία.
\item Η περίμετρος μετράει μόνο το περίγραμμα, όχι την επιφάνεια του σχήματος.
\end{itemize}
\end{prosoxi}

\tmima{Λυμένα παραδείγματα}

\paradeigma{Μετατροπές μονάδων}
\tixero $2{,}4$ km, $35$ cm και $7$ dm.\par
\tithelo Να τα μετατρέψω σε m, m και mm αντίστοιχα.\par
\possindeo Από km σε m πολλαπλασιάζω επί $1000$. Από cm σε m διαιρώ με $100$. Από dm σε mm πολλαπλασιάζω επί $100$.\par
\apantisi $2{,}4$ km $=2400$ m, $35$ cm $=0{,}35$ m, $7$ dm $=700$ mm.\par
\elegxos Το $0{,}35$ m είναι λιγότερο από ένα μέτρο, όπως τα $35$ cm.

\paradeigma{Ποιο είναι μεγαλύτερο;}
\tixero Δύο ράφια έχουν μήκη $0{,}8$ m και $75$ cm.\par
\tithelo Ποιο είναι μεγαλύτερο;\par
\possindeo Τα φέρνω στην ίδια μονάδα: $0{,}8$ m $=80$ cm.\par
\apantisi Μεγαλύτερο είναι το πρώτο ($80$ cm $>75$ cm).\par
\elegxos Σε mm: $800>750$.

\paradeigma{Το συντριβάνι}
\tixero Το μονοπάτι είναι $9{,}8$ μ. και το συντριβάνι θα μπει στο μέσο του.\par
\tithelo Πόσο απέχει το συντριβάνι από κάθε άκρο.\par
\possindeo Το μέσο $\mathrm{Μ}$ έχει $\mathrm{ΑΜ}=\dfrac{\mathrm{ΑΒ}}{2}=\dfrac{9{,}8}{2}$.\par
\apantisi $4{,}9$ μ. από κάθε άκρο.\par
\elegxos $4{,}9+4{,}9=9{,}8$.

\paradeigma{Τα δέντρα στη σειρά}
\tixero Τα δέντρα $\mathrm{Α},\,\mathrm{Β},\,\Gamma,\,\Delta$ είναι με αυτή τη σειρά σε ευθεία και $\mathrm{ΑΒ}=3$ μ., $\mathrm{ΒΓ}=4$ μ., $\Gamma\Delta=2{,}5$ μ.\par
\tithelo Τα $\mathrm{ΑΓ}$, $\mathrm{ΑΔ}$ και $\mathrm{ΒΔ}$.\par
\possindeo $\mathrm{ΑΓ}=\mathrm{ΑΒ}+\mathrm{ΒΓ}=3+4=7$. $\mathrm{ΑΔ}=\mathrm{ΑΓ}+\Gamma\Delta=7+2{,}5=9{,}5$. $\mathrm{ΒΔ}=\mathrm{ΑΔ}-\mathrm{ΑΒ}=9{,}5-3=6{,}5$.\par
\apantisi $\mathrm{ΑΓ}=7$ μ., $\mathrm{ΑΔ}=9{,}5$ μ. και $\mathrm{ΒΔ}=6{,}5$ μ.\par
\elegxos $\mathrm{ΒΔ}=\mathrm{ΒΓ}+\Gamma\Delta=4+2{,}5=6{,}5$.

\paradeigma{Περίμετρος με μικτές μονάδες}
\tixero Ένα τρίγωνο έχει πλευρές $4$ cm, $35$ mm και $0{,}06$ m.\par
\tithelo Η περίμετρός του σε cm.\par
\possindeo Τα φέρνω όλα σε cm: $35$ mm $=3{,}5$ cm και $0{,}06$ m $=6$ cm. Προσθέτω: $4+3{,}5+6$.\par
\apantisi $13{,}5$ cm.\par
\elegxos Αν είχα προσθέσει $4+35+0{,}06$ χωρίς μετατροπή, θα έβγαινε παράλογος αριθμός.

\paradeigma{Η περίμετρος του παρτεριού}
\tixero Ορθογώνιο παρτέρι με πλευρές $4{,}5$ μ. και $3{,}2$ μ.\par
\tithelo Πόσος φράχτης χρειάζεται για να το περιβάλει, και τι γίνεται αν διπλασιάσουμε κάθε πλευρά.\par
\possindeo Η περίμετρος είναι το άθροισμα των πλευρών: $4{,}5+3{,}2+4{,}5+3{,}2=15{,}4$. Με διπλάσιες πλευρές: $9+6{,}4+9+6{,}4=30{,}8$.\par
\apantisi $15{,}4$ μ. και, με διπλάσιες πλευρές, $30{,}8$ μ., δηλαδή διπλάσια περίμετρος.\par
\elegxos $2\cdot15{,}4=30{,}8$.

\tmima{Ασκήσεις}

\zoni{Το κατάλαβα}{μονάδες, μέσο, άθροισμα, περίμετρος}

\ex Ποια μονάδα (km, m, dm, cm, mm) είναι προτιμότερη για τις μετρήσεις; α) απόσταση δύο πόλεων \quad β) μήκος κρεβατιού \quad γ) πλάτος παραθύρου \quad δ) πάχος τετραδίου \quad ε) μήκος ποδηλατόδρομου.

\ex Μετάτρεψε σε m: α) $1{,}78$ km \quad β) $0{,}12$ dm \quad γ) $21{,}3$ cm \quad δ) $0{,}03$ km \quad ε) $2$ dm \quad στ) $10$ mm.

\ex Μετάτρεψε σε cm: α) $0{,}7$ m \quad β) $45$ mm \quad γ) $1{,}2$ dm \quad δ) $0{,}03$ km.

\ex Σύγκρινε και βάλε σε αύξουσα σειρά τα μήκη: $0{,}35$ m, $3{,}4$ dm, $36$ cm, $345$ mm.

\ex Σχεδίασε ένα ευθύγραμμο τμήμα $\mathrm{ΑΒ}=6{,}5$ cm και βρες το μέσο του $\mathrm{Μ}$. Μέτρησε τα $\mathrm{ΑΜ}$ και $\mathrm{ΜΒ}$ και σύγκρινέ τα.

\ex Τα $\mathrm{Α},\,\mathrm{Β},\,\Gamma$ είναι με αυτή τη σειρά σε μια ευθεία. α) Αν $\mathrm{ΑΒ}=4{,}2$ cm και $\mathrm{ΒΓ}=3{,}5$ cm, βρες το $\mathrm{ΑΓ}$. β) Αν $\mathrm{ΑΓ}=10$ cm και $\mathrm{ΑΒ}=6{,}5$ cm, βρες το $\mathrm{ΒΓ}$.

\ex Σχεδίασε τμήματα $\mathrm{ΑΒ}=5$ cm, $\Gamma\Delta=3$ cm και $\mathrm{ΕΖ}=2$ cm. Κατασκεύασε με τον διαβήτη τμήμα με μήκος: α) $\mathrm{ΑΒ}+\Gamma\Delta$ \quad β) $\mathrm{ΑΒ}-\mathrm{ΕΖ}$ \quad γ) $\mathrm{ΑΒ}+\Gamma\Delta-\mathrm{ΕΖ}$. Έλεγξε με τον χάρακα.

\ex Βρες την περίμετρο: α) ορθογωνίου με πλευρές $6$ και $2{,}5$ \quad β) τριγώνου με πλευρές $5$, $6$, $7{,}5$ \quad γ) τετραγώνου με πλευρά $2{,}8$ \quad δ) πενταγώνου με όλες τις πλευρές ίσες με $4$.

\zoni{Το χρησιμοποιώ}{βήματα, σχοινιά, φράχτες, δρόμοι}

\ex \textbf{Τα βήματα.} Ο Στέλιος μετράει το μονοπάτι με $14$ βήματα και ο Ηλίας με $16$. Γιατί βγαίνουν διαφορετικοί αριθμοί; Αν το μονοπάτι είναι $9{,}8$ μ., πόσο είναι το βήμα του καθενός;

\ex \textbf{Το συντριβάνι.} Ένα άλλο μονοπάτι είναι $12{,}6$ μ. και το συντριβάνι μπαίνει στο μέσο του. Πόσο απέχει από κάθε άκρο; Ένα παγκάκι βρίσκεται πάνω στο μονοπάτι, $1{,}5$ μ. από το συντριβάνι προς το ένα άκρο. Πόσο απέχει το παγκάκι από αυτό το άκρο;

\ex \textbf{Το κορδόνι.} Η θεία Μαρία πουλά κορδόνι για κήπο προς $1{,}80$ € το μέτρο. Η Ρήνη αγοράζει $3$ μ. και $45$ cm. Πόσα μέτρα αγόρασε και πόσα ευρώ πλήρωσε;

\ex \textbf{Ο φράχτης.} Το ορθογώνιο παρτέρι έχει πλευρές $4{,}5$ μ. και $3{,}2$ μ. Ο φράχτης κοστίζει $3$ € το μέτρο και ένα μέτρο μένει ανοιχτό για την είσοδο. Πόσο θα κοστίσει;

\ex \textbf{Η διαδρομή.} Ο Ηλίας πηγαίνει από το σπίτι στο σχολείο με στάσεις: $250$ μ., $180$ μ. και $320$ μ. Πόσα μέτρα περπάτησε και πόσα χιλιόμετρα είναι;

\ex \textbf{Αμερικανικά μέτρα.} Ένα ράφι σε αμερικανικό κατάστημα έχει μήκος $4$ ft. Πόσες ίντσες είναι; Πόσα μέτρα είναι περίπου (με $1$ ft $\approx0{,}3048$ m);

\zoni{Το σκέφτομαι}{ίδια μονάδα, διπλασιάζω, αληθές ή ψευδές}

\ex Ο Βασίλης λέει ότι $0{,}8$ m $<75$ cm, επειδή το $0{,}8$ είναι μικρότερο από το $75$. Έχει δίκιο; Πώς θα του το εξηγούσες;

\ex Το $\mathrm{Μ}$ είναι το μέσο του $\mathrm{ΑΒ}=12$ cm και το $\mathrm{Ν}$ το μέσο του $\mathrm{ΑΜ}$. Πόσο είναι τα $\mathrm{ΑΝ}$ και $\mathrm{ΝΒ}$;

\ex Σωστό ή λάθος; α) Το χιλιοστόμετρο είναι μεγαλύτερη μονάδα από το εκατοστόμετρο. \quad β) Αν $\mathrm{Μ}$ μέσο του $\mathrm{ΑΒ}$ τότε $\mathrm{ΑΜ}=\mathrm{ΜΒ}$. \quad γ) Αν $\mathrm{ΜΑ}=\mathrm{ΜΒ}$ τότε το $\mathrm{Μ}$ είναι πάντα το μέσο του $\mathrm{ΑΒ}$. \quad δ) Η περίμετρος ενός σχήματος αλλάζει αν μετρήσω τις πλευρές με άλλη σειρά. \quad ε) $1$ m $=0{,}001$ km.

\ex Αν διπλασιάσουμε κάθε πλευρά ενός ορθογωνίου, πόσο γίνεται η περίμετρος; Θα ίσχυε το ίδιο και για ένα τρίγωνο; Δοκίμασε με παράδειγμα.

\zoni{Για το σπίτι}{μέτρησε και πρόσθεσε}

\ex Μέτρησε με τα βήματά σου και με μετροταινία το ίδιο μήκος στο σπίτι (π.χ. το διάδρομο) και υπολόγισε το μήκος του βήματός σου.

\ex Μέτρησε τις πλευρές του δωματίου σου και υπολόγισε την περίμετρο της κάτοψής του. Βρες και το μέσο ενός σχοινιού διπλώνοντάς το, και έλεγξε με μετροταινία.

\paragrafos{Μέτρο γωνίας}{Β.1.3, «Μέτρο γωνίας»}

\begin{anagke}
Ο παππούς Γιάννης, που ταξίδευε παλιά με καράβια, φέρνει στην τάξη μια πυξίδα και τη δείχνει στην παρέα. «Οι ναυτικοί βρίσκουν την πορεία τους με την πυξίδα», λέει. «Όταν το καράβι αλλάζει πορεία, ο καπετάνιος δεν λέει “στρίψε λίγο”. Λέει πόσο ακριβώς». Η κυρία Αλεξάνδρα ρωτάει: «Αν το πλοίο στρίψει προς τα βορειοανατολικά ή προς τα βορειοδυτικά, πώς ξέρουμε σε ποια περίπτωση η στροφή είναι μεγαλύτερη;».

Στο μεταξύ στον κήπο η Αγγελική σχεδιάζει μια γωνία φράχτη με πολύ μακριές πλευρές, και ο Νίκος μια ίδια γωνία με κοντές πλευρές. «Η δική μου είναι μεγαλύτερη, έχει μεγαλύτερες πλευρές!», λέει η Αγγελική. «Όχι», απαντά η Ρήνη. «Εγώ θα τις μετρήσω». Παίρνει ένα μοιρογνωμόνιο και μετράει και τις δύο: βγαίνει το ίδιο νούμερο.

Για να συγκρίνουμε, να προσθέτουμε και να αφαιρούμε γωνίες χρειαζόμαστε, όπως και στα μήκη, μια μονάδα και ένα όργανο μέτρησης. Σε αυτή την παράγραφο γνωρίζουμε τη μοίρα και το μοιρογνωμόνιο και μαθαίνουμε να υπολογίζουμε με μέτρα γωνιών.
\end{anagke}

\begin{lexi}{μοίρα, μοιρογνωμόνιο}{μοίρα, μοιρογνωμόνιο}
Η \lexiword{μοίρα} είναι το «μέρος», το μερίδιο που παίρνει ο καθένας. Μία μοίρα είναι ένα πολύ μικρό μέρος του πλήρους κύκλου. Το \lexiword{μοιρογνωμόνιο} είναι το όργανο που «γνωρίζει», δηλαδή μετράει, τις μοίρες.
\end{lexi}

\begin{metafrasi}
Μονάδα μέτρησης των γωνιών. & μοίρα, $1^\circ$\\
Η γωνία $\hat{\mathrm{Α}}$ ανοίγει $40$ μοίρες. & $\hat{\mathrm{Α}}=40^\circ$\\
Υποδιαίρεση της μοίρας. & $1^\circ=60'$ (πρώτα λεπτά)\\
Υποδιαίρεση του λεπτού. & $1'=60''$ (δεύτερα λεπτά)\\
Δύο γωνίες με το ίδιο μέτρο. & $\hat{\mathrm{Κ}}=\hat{\Lambda}$\\
\end{metafrasi}

\tmima{Η θεωρία}

\tmimasmall{Η μοίρα και το μοιρογνωμόνιο}

\begin{orismos}{μοίρα}
\textbf{Μονάδα μέτρησης των γωνιών} είναι η \textbf{μοίρα}, που γράφεται $1^\circ$.
\end{orismos}

Η μέτρηση των γωνιών γίνεται με το \textbf{μοιρογνωμόνιο}. Ο αριθμός που βρίσκουμε ονομάζεται \textbf{μέτρο} της γωνίας. Τοποθετούμε το κέντρο του μοιρογνωμονίου στην κορυφή της γωνίας και τη μία πλευρά πάνω στο μηδέν. Διαβάζουμε το μέτρο στο σημείο όπου η άλλη πλευρά διασχίζει την κλίμακα, και προσέχουμε να διαβάζουμε τη σωστή από τις δύο κλίμακες.

\begin{center}
\begin{tikzpicture}[scale=1.15]
\draw[ggray] (-2.2,0) arc(180:0:2.2) -- cycle;
\foreach \a in {0,10,...,180}{\draw[ggray] ({1.95*cos(\a)},{1.95*sin(\a)})--({2.2*cos(\a)},{2.2*sin(\a)});}
\foreach \a in {0,45,90,135,180}{\node[ggray,font=\scriptsize] at ({1.65*cos(\a)},{1.65*sin(\a)}) {$\a^\circ$};}
\draw[thick] (-2.6,0)--(2.6,0); \draw[thick,gblue] (0,0)--({2.6*cos(40)},{2.6*sin(40)});
\fill (0,0) circle (2pt) node[below] {$\mathrm{Α}$};
\draw[gblue,thick] (0.7,0) arc(0:40:0.7); \node[gblue] at (0.95,0.3) {\small $40^\circ$};
\end{tikzpicture}
\end{center}

Παραδείγματα: το μέτρο της γωνίας $\hat{\mathrm{Α}}$ είναι $40^\circ$ και το μέτρο της $\hat{\mathrm{Β}}$ είναι $125^\circ$. Μια γωνία που ανοίγει όσο η ορθή έχει μέτρο $90^\circ$. Μια γωνία με πλευρές αντικείμενες ημιευθείες έχει μέτρο $180^\circ$ και μια μη κυρτή γωνία έχει μέτρο μεγαλύτερο από $180^\circ$.

Η μοίρα έχει υποδιαίρεση, το \textbf{λεπτό}: $1^\circ=60'$ (πρώτα λεπτά) και $1'=60''$ (δεύτερα λεπτά). Για παράδειγμα, $25^\circ30'=25{,}5^\circ$, γιατί $30'$ είναι η μισή μοίρα. Και $48{,}25^\circ=48^\circ15'$, γιατί $0{,}25\cdot60=15$.

\begin{prosoxi}[Προσοχή: δεν είναι δεκαδικά]
Το $1^\circ=60'$ και όχι $100'$. Γι' αυτό το $25^\circ30'$ δεν είναι $25{,}30^\circ$ αλλά $25{,}5^\circ$.
\end{prosoxi}

\begin{kanonas}{το μέτρο εξαρτάται μόνο από το «άνοιγμα»}
Το μέτρο μιας γωνίας είναι μοναδικό και εξαρτάται μόνο από το \textbf{άνοιγμα} των πλευρών της, όχι από το μήκος τους. Οι πλευρές είναι ημιευθείες και δεν έχουν μήκος.
\end{kanonas}

\tmimasmall{Ίσες γωνίες και σύγκριση γωνιών}

Δύο γωνίες είναι \textbf{ίσες} αν έχουν το ίδιο μέτρο. Τότε η μία εφαρμόζει ακριβώς πάνω στην άλλη.

Μπορούμε να συγκρίνουμε δύο γωνίες και με διαφανές χαρτί, χωρίς μοιρογνωμόνιο:
\begin{itemize}[leftmargin=1.2em,itemsep=0.3ex,topsep=2pt]
\item Αντιγράφουμε τη μία από τις δύο γωνίες σε διαφανές χαρτί.
\item Τοποθετούμε το διαφανές χαρτί πάνω στην άλλη γωνία ώστε να έχουν κοινή κορυφή και μία κοινή πλευρά.
\end{itemize}
Η γωνία που μένει μέσα στην άλλη είναι η μικρότερη.

\begin{center}
\begin{tikzpicture}[scale=1.1]
\begin{scope}
\draw[thick] (3,0)--(0,0)--(2.6,1.5); \draw[gblue] (0.7,0) arc(0:30:0.7); \fill (0,0) circle (2pt);
\node[below left] at (0,0) {$\mathrm{Κ}$}; \node[right] at (1.6,0.3) {\small $30^\circ$};
\end{scope}
\begin{scope}[xshift=4.5cm]
\draw[thick] (1.5,0)--(0,0)--(1.3,0.75); \draw[gblue] (0.7,0) arc(0:30:0.7); \fill (0,0) circle (2pt);
\node[below left] at (0,0) {$\Lambda$};
\end{scope}
\begin{scope}[xshift=8cm]
\draw[thick] (3,0)--(0,0)--(2.2,2); \fill[gblueL] (0,0)--(2.2,0) arc(0:42:2.2)--cycle;
\fill[gwarm] (0,0)--(1.7,0) arc(0:28:1.7)--cycle; \fill (0,0) circle (2pt);
\node[below left] at (0,0) {$\mathrm{Ο}$};
\end{scope}
\end{tikzpicture}
\end{center}

Στα δύο πρώτα σχήματα οι γωνίες έχουν το ίδιο μέτρο ($30^\circ$) παρά τα διαφορετικά μήκη των πλευρών τους, άρα είναι ίσες. Στο τρίτο σχήμα η μικρότερη γωνία (στο εσωτερικό της άλλης) έχει μικρότερο μέτρο.

\tmimasmall{Πρόσθεση και αφαίρεση γωνιών}

Δύο γωνίες μπορούν να προστεθούν με \textbf{επίθεση}: τις τοποθετούμε δίπλα δίπλα, με κοινή κορυφή και μία κοινή πλευρά, χωρίς να επικαλύπτονται. Η γωνία που σχηματίζεται από τις δύο εξωτερικές πλευρές έχει μέτρο το \textbf{άθροισμα} των μέτρων τους. Αντίστροφα, αν από μια γωνία αφαιρέσουμε ένα μέρος της, το υπόλοιπο έχει μέτρο τη \textbf{διαφορά}.

\begin{center}
\begin{tikzpicture}[scale=1.1]
\draw[thick] (3,0)--(0,0); \draw[thick] (0,0)--(2.4,1.7); \draw[thick] (0,0)--(-0.9,2.6);
\fill (0,0) circle (2pt) node[below left] {$\mathrm{Ο}$};
\node[right] at (3,0) {$x$}; \node[right] at (2.4,1.7) {$y$}; \node[above] at (-0.9,2.6) {$z$};
\draw[gblue] (0.8,0) arc(0:35:0.8); \node[gblue] at (1.2,0.18) {\small $\alpha$};
\draw[gwarmD] (0.55,0.38) arc(35:110:0.55) ; \node[gwarmD] at (0.12,0.85) {\small $\beta$};
\end{tikzpicture}
\end{center}

Αν οι γωνίες $x\hat{\mathrm{Ο}}y$ και $y\hat{\mathrm{Ο}}z$ τοποθετηθούν δίπλα δίπλα με αυτό τον τρόπο, τότε $x\hat{\mathrm{Ο}}z=x\hat{\mathrm{Ο}}y+y\hat{\mathrm{Ο}}z$ και $y\hat{\mathrm{Ο}}z=x\hat{\mathrm{Ο}}z-x\hat{\mathrm{Ο}}y$.

\begin{prosoxi}
\begin{itemize}[leftmargin=1.2em,itemsep=0.3ex,topsep=2pt]
\item Μια γωνία με μακρύτερες πλευρές δεν είναι μεγαλύτερη: το μέτρο εξαρτάται μόνο από το άνοιγμα.
\item Στο μοιρογνωμόνιο διαβάζω την κλίμακα που ξεκινά από το μηδέν στην πλευρά που έβαλα πάνω στο $0$. Αν η γωνία δείχνει μικρότερη από την ορθή, δεν γράφω αριθμό μεγαλύτερο από $90$.
\item Για να προσθέσω δύο γωνίες πρέπει να έχουν κοινή κορυφή και κοινή πλευρά και να μην επικαλύπτονται.
\item Το $25^\circ30'$ δεν ισούται με $25{,}30^\circ$.
\end{itemize}
\end{prosoxi}

\tmima{Λυμένα παραδείγματα}

\paradeigma{Η σωστή κλίμακα}
\tixero Το μοιρογνωμόνιο έχει δύο κλίμακες. Η μία πλευρά μιας γωνίας δείχνει στο $0$ και η άλλη δείχνει $40$ στη μία κλίμακα και $140$ στην άλλη. Η γωνία φαίνεται μικρότερη από την ορθή.\par
\tithelo Το μέτρο της γωνίας.\par
\possindeo Αφού η γωνία είναι μικρότερη από την ορθή, το μέτρο της είναι μικρότερο από $90^\circ$.\par
\apantisi $40^\circ$.\par
\elegxos $40^\circ<90^\circ$ και $40^\circ+140^\circ=180^\circ$, όπως οι δύο κλίμακες.

\paradeigma{Μοίρες και λεπτά}
\tixero Η γωνία $\hat{\mathrm{Α}}=36^\circ30'$ και η γωνία $\hat{\mathrm{Β}}=48{,}25^\circ$.\par
\tithelo Να γράψω την $\hat{\mathrm{Α}}$ με δεκαδικό και την $\hat{\mathrm{Β}}$ με λεπτά.\par
\possindeo Επειδή $1^\circ=60'$, τα $30'$ είναι $30:60=0{,}5^\circ$. Και $0{,}25^\circ=0{,}25\cdot60'=15'$.\par
\apantisi $\hat{\mathrm{Α}}=36{,}5^\circ$ και $\hat{\mathrm{Β}}=48^\circ15'$.\par
\elegxos $0{,}5^\circ$ είναι η μισή μοίρα, δηλαδή $30'$.

\paradeigma{Ίσες ή όχι;}
\tixero Η Αγγελική έχει γωνία με πλευρές $12$ cm και ο Νίκος γωνία με πλευρές $4$ cm. Και οι δύο μετρούν $\hat{\omega}=62^\circ$.\par
\tithelo Ποια γωνία είναι μεγαλύτερη;\par
\possindeo Το μέτρο μιας γωνίας εξαρτάται από το άνοιγμα και όχι από τα μήκη των πλευρών.\par
\apantisi Είναι ίσες.\par
\elegxos Αν αντιγράψω τη μία σε διαφανές χαρτί, εφαρμόζει πάνω στην άλλη.

\paradeigma{Πρόσθεση και αφαίρεση}
\tixero Η γωνία του φράχτη είναι $x\hat{\mathrm{Ο}}z=84^\circ$ και ένα μονοπάτι $\mathrm{Οy}$ την χωρίζει ώστε $x\hat{\mathrm{Ο}}y=50^\circ$.\par
\tithelo Η γωνία $y\hat{\mathrm{Ο}}z$.\par
\possindeo Η $x\hat{\mathrm{Ο}}z$ είναι το άθροισμα των $x\hat{\mathrm{Ο}}y$ και $y\hat{\mathrm{Ο}}z$. Άρα $y\hat{\mathrm{Ο}}z=84^\circ-50^\circ$.\par
\apantisi $34^\circ$.\par
\elegxos $50^\circ+34^\circ=84^\circ$.

\paradeigma{Οι δείκτες του ρολογιού}
\tixero Ο κύκλος γύρω από το κέντρο του ρολογιού έχει $360^\circ$ και το ρολόι έχει $12$ ώρες.\par
\tithelo Η γωνία ανάμεσα σε δύο διαδοχικές ώρες και η γωνία των δεικτών στις $4{:}00$.\par
\possindeo Ανάμεσα σε δύο διαδοχικές ώρες: $360^\circ:12=30^\circ$. Στις $4{:}00$ ο λεπτοδείκτης δείχνει το $12$ και ο ωροδείκτης το $4$, δηλαδή $4$ διαστήματα ωρών.\par
\apantisi $30^\circ$ και $4\cdot30^\circ=120^\circ$.\par
\elegxos Στις $3{:}00$ η γωνία είναι $3\cdot30^\circ=90^\circ$, δηλαδή ορθή.

\tmima{Ασκήσεις}

\zoni{Το κατάλαβα}{μέτρο, σύγκριση, λεπτά, πρόσθεση}

\ex Σχεδίασε με το μοιρογνωμόνιο τις γωνίες: $\hat{\alpha}=30^\circ$, $\hat{\beta}=58^\circ$, $\hat{\gamma}=90^\circ$, $\hat{\delta}=135^\circ$, $\hat{\varepsilon}=180^\circ$.

\ex Σχεδίασε δύο γωνίες με μέτρο $40^\circ$: τη μία με πλευρές μήκους $3$ cm και την άλλη με πλευρές μήκους $8$ cm. Ποια είναι μεγαλύτερη; Εξήγησε.

\ex Μετάτρεψε: α) $12^\circ30'$ σε μοίρες με δεκαδικό \quad β) $40^\circ15'$ σε μοίρες με δεκαδικό \quad γ) $18{,}75^\circ$ σε μοίρες και λεπτά \quad δ) $1^\circ$ σε δεύτερα λεπτά.

\ex Βάλε σε αύξουσα σειρά τις γωνίες: $45^\circ$, $40^\circ30'$, $40{,}25^\circ$, $45^\circ10'$.

\ex Οι γωνίες $x\hat{\mathrm{Ο}}y$ και $y\hat{\mathrm{Ο}}z$ έχουν κοινή κορυφή, κοινή πλευρά $\mathrm{Οy}$ και δεν επικαλύπτονται. α) Αν $x\hat{\mathrm{Ο}}y=35^\circ$ και $y\hat{\mathrm{Ο}}z=48^\circ$, βρες την $x\hat{\mathrm{Ο}}z$. β) Αν $x\hat{\mathrm{Ο}}z=120^\circ$ και $x\hat{\mathrm{Ο}}y=75^\circ$, βρες την $y\hat{\mathrm{Ο}}z$.

\ex Σχεδίασε ένα τρίγωνο, μέτρησε τις τρεις γωνίες του με το μοιρογνωμόνιο και βρες το άθροισμά τους. Επανάλαβε με άλλο τρίγωνο. Τι παρατηρείς;

\ex Σωστό ή λάθος; α) Μονάδα μέτρησης των γωνιών είναι η μοίρα. \quad β) Όσο μεγαλύτερο μήκος έχουν οι πλευρές μιας γωνίας, τόσο μεγαλύτερη είναι. \quad γ) Δύο γωνίες είναι ίσες όταν έχουν το ίδιο μέτρο. \quad δ) $1^\circ=100'$. \quad ε) Το μέτρο της ορθής γωνίας είναι $90^\circ$. \quad στ) Το μοιρογνωμόνιο μετράει γωνίες.

\zoni{Το χρησιμοποιώ}{πυξίδα, ρολόι, φράχτης, πόρτα}

\ex \textbf{Το πλοίο.} Ένα πλοίο στρίβει πρώτα $45^\circ$ και μετά ακόμη $30^\circ$ προς την ίδια μεριά. Πόση είναι η συνολική στροφή; Αν μετά στρίψει $20^\circ$ προς την αντίθετη μεριά, πόση στροφή απομένει από την αρχική πορεία;

\ex \textbf{Το ρολόι.} Πόσες μοίρες είναι η γωνία ανάμεσα στους δείκτες στις $2{:}00$, στις $5{:}00$ και στις $6{:}00$;

\ex \textbf{Το παρτέρι.} Η γωνία του φράχτη στον κήπο είναι $84^\circ$. Η Αγγελική σχεδιάζει ένα μονοπάτι που σχηματίζει $50^\circ$ με τη μία πλευρά και χωρίζει το παρτέρι σε δύο μέρη. Πόση γωνία σχηματίζει το μονοπάτι με την άλλη πλευρά;

\ex \textbf{Η πόρτα.} Η πόρτα της τάξης ανοίγει πρώτα $35^\circ$ και μετά ανοίγει ακόμη $50^\circ$. Στη συνέχεια η Φανή τη σπρώχνει και κλείνει $20^\circ$. Πόσο ανοιχτή είναι τώρα;

\ex \textbf{Η κλίμακα.} Η Φανή μετράει μια γωνία που φαίνεται μικρότερη από την ορθή και το μοιρογνωμόνιο δείχνει $130$ στη μία κλίμακα και $50$ στην άλλη. Ποιο είναι το σωστό μέτρο και γιατί;

\zoni{Το σκέφτομαι}{τι εξαρτάται από τι}

\ex Γιατί το μέτρο μιας γωνίας δεν εξαρτάται από το μήκος των πλευρών της; Σκέψου τι είναι οι πλευρές μιας γωνίας.

\ex Οι γωνίες $x\hat{\mathrm{Ο}}y=70^\circ$, $y\hat{\mathrm{Ο}}z=50^\circ$ και $z\hat{\mathrm{Ο}}t=40^\circ$ είναι διαδοχικές, χωρίς να επικαλύπτονται. Πόση είναι η $x\hat{\mathrm{Ο}}t$;

\ex Πώς μπορείς να συγκρίνεις δύο γωνίες όταν δεν έχεις μοιρογνωμόνιο; Περιέγραψε τα βήματα.

\ex Σχεδίασε μια γωνία $x\hat{\mathrm{Ο}}y=30^\circ$. Με κοινή κορυφή και κοινή πλευρά $\mathrm{Οx}$ σχεδίασε γωνία $x\hat{\mathrm{Ο}}t$ α) διπλάσιου μέτρου \quad β) μισού μέτρου.

\zoni{Για το σπίτι}{γωνίες γύρω μας}

\ex Βρες στο σπίτι πέντε γωνίες (ανοιχτό βιβλίο, ψαλίδι, ανοιχτή πόρτα, γωνία τραπεζιού, ανοιχτός φορητός υπολογιστής) και μέτρησε το άνοιγμά τους με το μοιρογνωμόνιο.

\ex Άνοιξε έναν διαβήτη σε τρία διαφορετικά ανοίγματα και σχεδίασε στο τετράδιό σου τη γωνία που σχηματίζουν τα σκέλη του κάθε φορά. Μέτρησέ τες.

\paragrafos{Άθροισμα και διαφορά τμημάτων: από την ευθεία στην τεθλασμένη}{Β.1.4, «Πρόσθεση και αφαίρεση ευθυγράμμων τμημάτων»}

\begin{anagke}
Ο παππούς Γιάννης φέρνει τα δέντρα που θα φυτευτούν σε μία σειρά, και ο Ηλίας μετράει τις αποστάσεις: το δεύτερο δέντρο απέχει από το πρώτο $3$ μέτρα, το τρίτο $4$ μέτρα από το δεύτερο και το τέταρτο $2{,}5$ από το τρίτο. «Πόσο απέχει το τέταρτο από το πρώτο;» ρωτάει η Πολυξένη. Ο Ηλίας το βρίσκει εύκολα: $3+4+2{,}5=9{,}5$ μέτρα. «Και το τρίτο από το πρώτο; $3+4=7$. Και το τέταρτο από το δεύτερο;». Εδώ ο Στέλιος προσπαθεί να προσθέσει, αλλά η Αγγελική λέει: «Αφαίρεσε: από την απόσταση του τέταρτου από το πρώτο, αφαιρώ την απόσταση του δεύτερου από το πρώτο: $9{,}5-3=6{,}5$».

Η παρέα ανακαλύπτει ότι οι αποστάσεις προσθέτονται και αφαιρούνται όταν τα σημεία βρίσκονται στην ίδια ευθεία με τη σειρά. Μετά ο Νίκος γυρίζει στο μονοπάτι που στρίβει: «Αν πάω από το πρώτο στο τρίτο δέντρο περνώντας από ένα άλλο σημείο που δεν είναι στην ευθεία, κάνω περισσότερο δρόμο;». Η Ρήνη το δοκιμάζει με σχοινί και δείχνει ότι το ευθύγραμμο τμήμα είναι πάντα το συντομότερο μονοπάτι ανάμεσα σε δύο σημεία. Και η κυρία Αλεξάνδρα προσθέτει: «Το μήκος του περιγράμματος ενός παρτεριού, για το φράχτη του, είναι η περίμετρός του».
\end{anagke}

\begin{lexi}{άθροισμα τμημάτων, περίμετρος}{άθροισμα τμημάτων, περίμετρος}
Η \lexiword{περίμετρος} είναι το «γύρω μέτρημα» (περί + μέτρο): το συνολικό μήκος του περιγράμματος ενός σχήματος. Η \lexiword{τεθλασμένη} γραμμή έχει μήκος το άθροισμα των μηκών των ευθυγράμμων τμημάτων της.
\end{lexi}

\begin{metafrasi}
Τα $\mathrm{Α},\,\mathrm{Β},\,\Gamma$ στη σειρά σε ευθεία. & $\mathrm{ΑΓ}=\mathrm{ΑΒ}+\mathrm{ΒΓ}$\\
Ό,τι περισσεύει μετά το $\mathrm{ΑΒ}$. & $\mathrm{ΒΓ}=\mathrm{ΑΓ}-\mathrm{ΑΒ}$\\
Τεθλασμένη $\mathrm{ΑΒΓΔ}$. & μήκος $=\mathrm{ΑΒ}+\mathrm{ΒΓ}+\Gamma\Delta$\\
Ο συντομότερος δρόμος. & $\mathrm{ΑΒ}<$ μήκος κάθε τεθλασμένης με άκρα $\mathrm{Α},\,\mathrm{Β}$\\
Το άθροισμα των πλευρών. & περίμετρος\\
\end{metafrasi}

\tmima{Η θεωρία}

Αν τα σημεία $\mathrm{Α},\,\mathrm{Β},\,\Gamma$ βρίσκονται με αυτή τη σειρά σε μια ευθεία, το τμήμα $\mathrm{ΑΓ}$ λέγεται \textbf{άθροισμα} των τμημάτων $\mathrm{ΑΒ}$ και $\mathrm{ΒΓ}$:
$$\mathrm{ΑΓ}=\mathrm{ΑΒ}+\mathrm{ΒΓ}.$$
Και το $\mathrm{ΒΓ}$ είναι η \textbf{διαφορά} των $\mathrm{ΑΓ}$ και $\mathrm{ΑΒ}$: $\mathrm{ΒΓ}=\mathrm{ΑΓ}-\mathrm{ΑΒ}$.

\begin{center}
\begin{tikzpicture}
\draw[thick] (0,0)--(7,0);
\foreach \x/\l in {0/Α,2.5/Β,7/Γ}{\fill (\x,0) circle (2pt); \node[below] at (\x,-0.1) {$\mathrm{\l}$};}
\draw[gblue,thick] (0,0.4)--(2.5,0.4); \draw[gblue,thick,dashed] (2.5,0.4)--(7,0.4);
\node[gblue,above] at (1.25,0.45) {$\mathrm{ΑΒ}$}; \node[gblue,above] at (4.75,0.45) {$\mathrm{ΒΓ}$};
\end{tikzpicture}
\end{center}

\tmimasmall{Τεθλασμένη και περίμετρος}

Το μήκος μιας \textbf{τεθλασμένης γραμμής} είναι το άθροισμα των μηκών των ευθυγράμμων τμημάτων της. Το μήκος του ευθυγράμμου τμήματος $\mathrm{ΑΒ}$ είναι μικρότερο από το μήκος κάθε τεθλασμένης γραμμής με τα ίδια άκρα $\mathrm{Α}$, $\mathrm{Β}$ (εκτός αν η τεθλασμένη είναι το ίδιο το τμήμα). Δηλαδή το ευθύγραμμο τμήμα είναι ο συντομότερος δρόμος μεταξύ δύο σημείων.

\begin{center}
\begin{tikzpicture}
\draw[thick] (0,0)--(6,0); \draw[thick,gblue] (0,0)--(2,1.6)--(4,1)--(6,0);
\foreach \p/\l in {(0,0)/Α,(6,0)/Β,(2,1.6)/Γ,(4,1)/Δ}{\fill \p circle (2pt);}
\node[below left] at (0,0) {$\mathrm{Α}$}; \node[below right] at (6,0) {$\mathrm{Β}$}; \node[above] at (2,1.6) {$\Gamma$}; \node[above right] at (4,1) {$\Delta$};
\end{tikzpicture}
\end{center}

Το άθροισμα των πλευρών ενός ευθύγραμμου σχήματος (πολυγώνου) λέγεται \textbf{περίμετρος} του σχήματος.

\begin{prosoxi}
\begin{itemize}[leftmargin=1.2em,itemsep=0.3ex,topsep=2pt]
\item Η σχέση $\mathrm{ΑΓ}=\mathrm{ΑΒ}+\mathrm{ΒΓ}$ ισχύει μόνο όταν το $\mathrm{Β}$ βρίσκεται \emph{ανάμεσα} στα $\mathrm{Α}$ και $\mathrm{Γ}$ πάνω σε ευθεία. Αν δεν είναι συνευθειακά, το $\mathrm{ΑΓ}$ είναι μικρότερο από το $\mathrm{ΑΒ}+\mathrm{ΒΓ}$.
\item Η περίμετρος μετράει μόνο το μήκος του περιγράμματος, όχι την επιφάνεια.
\item Για την αφαίρεση προσέχω τη σειρά: $\mathrm{ΒΓ}=\mathrm{ΑΓ}-\mathrm{ΑΒ}$ και όχι $\mathrm{ΑΒ}-\mathrm{ΑΓ}$ (αρνητικό).
\item Για να συγκρίνω δύο διαδρομές πρέπει να έχουν την ίδια μονάδα μέτρησης.
\end{itemize}
\end{prosoxi}

\tmima{Λυμένα παραδείγματα}

\paradeigma{Τα δέντρα στη σειρά}
\tixero Τα δέντρα $\mathrm{Α},\,\mathrm{Β},\,\Gamma,\,\Delta$ είναι στη σειρά και $\mathrm{ΑΒ}=3$, $\mathrm{ΒΓ}=4$, $\Gamma\Delta=2{,}5$ μ.\par
\tithelo Τα $\mathrm{ΑΓ}$, $\mathrm{ΑΔ}$ και $\mathrm{ΒΔ}$.\par
\possindeo $\mathrm{ΑΓ}=\mathrm{ΑΒ}+\mathrm{ΒΓ}=3+4=7$. $\mathrm{ΑΔ}=\mathrm{ΑΓ}+\Gamma\Delta=7+2{,}5=9{,}5$. $\mathrm{ΒΔ}=\mathrm{ΑΔ}-\mathrm{ΑΒ}=9{,}5-3=6{,}5$.\par
\apantisi $7$, $9{,}5$ και $6{,}5$ μ.\par
\elegxos $\mathrm{ΒΔ}=\mathrm{ΒΓ}+\Gamma\Delta=4+2{,}5=6{,}5$.

\paradeigma{Με αφαίρεση}
\tixero Το $\mathrm{Β}$ είναι ανάμεσα στα $\mathrm{Α}$ και $\mathrm{Γ}$, $\mathrm{ΑΓ}=9$ και $\mathrm{ΑΒ}=4{,}5$.\par
\tithelo Το $\mathrm{ΒΓ}$.\par
\possindeo $\mathrm{ΒΓ}=\mathrm{ΑΓ}-\mathrm{ΑΒ}=9-4{,}5=4{,}5$.\par
\apantisi $4{,}5$.\par
\elegxos Άρα το $\mathrm{Β}$ είναι το μέσο του $\mathrm{ΑΓ}$.

\paradeigma{Ο κοντινός δρόμος}
\tixero Μονοπάτι $\mathrm{ΑΒΓ}$ με $\mathrm{ΑΒ}=6$ μ. και $\mathrm{ΒΓ}=8$ μ. (στρίβει στο $\mathrm{Β}$).\par
\tithelo Πόσο είναι το μήκος του και τι ξέρουμε για το $\mathrm{ΑΓ}$;\par
\possindeo Το μήκος της τεθλασμένης είναι $6+8=14$ μ. Το ευθύγραμμο τμήμα $\mathrm{ΑΓ}$ είναι συντομότερο.\par
\apantisi Τεθλασμένη $14$ μ. και $\mathrm{ΑΓ}<14$ μ.\par
\elegxos Αν το $\mathrm{Β}$ ήταν πάνω στο $\mathrm{ΑΓ}$, θα ήταν ίσα με $14$.

\paradeigma{Η περίμετρος}
\tixero Ένα ορθογώνιο παρτέρι με πλευρές $4{,}5$ μ. και $3{,}2$ μ.\par
\tithelo Το μήκος του φράχτη που το περιβάλλει.\par
\possindeo Η περίμετρος είναι το άθροισμα των πλευρών: $4{,}5+3{,}2+4{,}5+3{,}2=15{,}4$ μ.\par
\apantisi $15{,}4$ μ.\par
\elegxos $2\cdot(4{,}5+3{,}2)=2\cdot7{,}7=15{,}4$.

\paradeigma{Το μέσο και το άθροισμα}
\tixero Το $\mathrm{Γ}$ βρίσκεται πάνω στο $\mathrm{ΑΒ}$ με $\mathrm{ΑΓ}=3$ και $\mathrm{ΓΒ}=7$ και το $\mathrm{Μ}$ είναι το μέσο του $\mathrm{ΑΒ}$.\par
\tithelo Το $\mathrm{ΓΜ}$.\par
\possindeo $\mathrm{ΑΒ}=3+7=10$, άρα $\mathrm{ΑΜ}=5$. Το $\mathrm{Μ}$ είναι πέρα από το $\mathrm{Γ}$ και $\mathrm{ΓΜ}=\mathrm{ΑΜ}-\mathrm{ΑΓ}=5-3=2$.\par
\apantisi $\mathrm{ΓΜ}=2$.\par
\elegxos $\mathrm{ΓΜ}+\mathrm{ΜΒ}=2+5=7=\mathrm{ΓΒ}$.

\tmima{Ασκήσεις}

\zoni{Το κατάλαβα}{άθροισμα, διαφορά, περίμετρος}

\ex Τα $\mathrm{Α},\,\mathrm{Β},\,\Gamma,\,\Delta$ είναι στη σειρά σε μια ευθεία με $\mathrm{ΑΒ}=3$, $\mathrm{ΒΓ}=4$, $\Gamma\Delta=2{,}5$. Βρες τα $\mathrm{ΑΓ}$, $\mathrm{ΑΔ}$, $\mathrm{ΒΔ}$.

\ex Το $\mathrm{Β}$ είναι ανάμεσα στα $\mathrm{Α}$ και $\mathrm{Γ}$. Αν $\mathrm{ΑΓ}=9$ και $\mathrm{ΑΒ}=4{,}5$, βρες το $\mathrm{ΒΓ}$.

\ex Βρες το μήκος της τεθλασμένης $\mathrm{ΑΒΓΔ}$ με $\mathrm{ΑΒ}=5$, $\mathrm{ΒΓ}=3{,}5$, $\Gamma\Delta=6$.

\ex Βρες την περίμετρο: α) ορθογωνίου $4{,}5$ και $3{,}2$ \quad β) τριγώνου με πλευρές $5$, $6$, $7{,}5$ \quad γ) τετραγώνου πλευράς $2{,}8$ \quad δ) κανονικού εξαγώνου πλευράς $4$.

\ex Στην ευθεία $\mathrm{Α},\,\mathrm{Β},\,\Gamma$ με $\mathrm{ΑΒ}=2\mathrm{ΒΓ}$ και $\mathrm{ΑΓ}=12$ βρες τα $\mathrm{ΑΒ}$ και $\mathrm{ΒΓ}$.

\ex Σωστό ή λάθος; α) Αν το $\mathrm{Β}$ είναι ανάμεσα στα $\mathrm{Α}$ και $\Gamma$ τότε $\mathrm{ΑΓ}=\mathrm{ΑΒ}+\mathrm{ΒΓ}$. \quad β) Η τεθλασμένη μεταξύ δύο σημείων είναι πάντα μεγαλύτερη από το τμήμα που τα ενώνει (ή ίση). \quad γ) Η περίμετρος ενός τετραγώνου πλευράς $\alpha$ είναι $\alpha^2$. \quad δ) $\mathrm{ΒΓ}=\mathrm{ΑΒ}-\mathrm{ΑΓ}$ όταν το $\mathrm{Β}$ είναι ανάμεσα στα $\mathrm{Α}$, $\mathrm{Γ}$.

\zoni{Το χρησιμοποιώ}{φράχτες, δρόμοι, κήποι}

\ex \textbf{Ο φράχτης.} Ένα ορθογώνιο παρτέρι $4{,}5$ μ. επί $3{,}2$ μ. θα περιφραχθεί με δίχτυ που κοστίζει $3$ € το μέτρο. Πόσο θα κοστίσει; Μην ξεχάσεις ότι ένα μέτρο το χρειάζεται για την είσοδο που μένει ανοιχτή.

\ex \textbf{Ο κοντινός δρόμος.} Το μονοπάτι $\mathrm{ΑΒΓ}$ έχει $\mathrm{ΑΒ}=6$ μ., $\mathrm{ΒΓ}=8$ μ. Ο Νίκος κόβει δρόμο πηγαίνοντας ευθεία από το $\mathrm{Α}$ στο $\mathrm{Γ}$. Έκανε περισσότερο ή λιγότερο δρόμο από $14$ μ.; Γιατί;

\ex \textbf{Η απόσταση.} Τρία σπίτια $\mathrm{Α}$, $\mathrm{Β}$, $\Gamma$ βρίσκονται στον ίδιο δρόμο, με το $\mathrm{Β}$ στη μέση. Το $\mathrm{Α}$ απέχει από το $\mathrm{Β}$ $350$ μ. και το $\mathrm{Β}$ από το $\Gamma$ $480$ μ. Πόσο απέχουν τα $\mathrm{Α}$ και $\Gamma$;

\ex \textbf{Η τεθλασμένη διαδρομή.} Η Πολυξένη πηγαίνει από το σχολείο στο σπίτι με στάσεις: $250$ μ., $180$ μ. και $320$ μ. Αν πήγαινε ευθεία θα έκανε $600$ μ. Πόσο περισσότερο περπάτησε;

\ex \textbf{Ο φράχτης του κήπου.} Ο κήπος έχει σχήμα πενταγώνου με πλευρές $12$, $9$, $14$, $10$ και $8$ μ. Πόσος φράχτης χρειάζεται;

\zoni{Το σκέφτομαι}{ο ευθύς είναι ο συντομότερος}

\ex Εξήγησε γιατί το ευθύγραμμο τμήμα είναι ο συντομότερος δρόμος ανάμεσα σε δύο σημεία (σκέψου ένα τεντωμένο νήμα).

\ex Αν $\mathrm{ΑΒ}=5$ και $\mathrm{ΒΓ}=7$ και τα $\mathrm{Α},\,\mathrm{Β},\,\Gamma$ δεν είναι συνευθειακά, ανάμεσα σε ποιες τιμές βρίσκεται το $\mathrm{ΑΓ}$;

\ex Ένα σχοινί $\mathrm{ΑΒ}$ μήκους $10$ μ. διπλώνεται ώστε το $\mathrm{Α}$ να πέσει πάνω στο $\mathrm{Β}$. Πού βρίσκεται η πτυχή;

\ex Δείξε με παράδειγμα ότι η περίμετρος δύο σχημάτων μπορεί να είναι ίση χωρίς τα σχήματα να είναι ίσα.

\zoni{Για το σπίτι}{μέτρησε και πρόσθεσε}

\ex Μέτρησε τις πλευρές του δωματίου σου και υπολόγισε την περίμετρο της κάτοψής του.

\ex Φτιάξε μια τεθλασμένη με τρία τμήματα από σπάγκο και σύγκρινε το μήκος της με την απόσταση των άκρων.

\paragrafos{Μέτρηση γωνιών, σύγκριση και διχοτόμος}{Β.1.5, «Μέτρηση, σύγκριση και ισότητα γωνιών – Διχοτόμος γωνίας»}

\begin{anagke}
Στον κήπο υπάρχει μια γωνία φράχτη, και η Α3 θέλει να τη χωρίσει σε δύο ίσα παρτέρια με ένα μονοπάτι. «Πού βάζουμε το μονοπάτι;» ρωτάει ο Βασίλης. Η Αγγελική παίρνει το μοιρογνωμόνιο: «Πρώτα πρέπει να μετρήσουμε πόσο ανοίγει η γωνία». Βρίσκει $84$ μοίρες. «Το μισό είναι $42$», υπολογίζει η Πολυξένη. «Το μονοπάτι θα ξεκινήσει από την κορυφή και θα σχηματίζει $42$ μοίρες με κάθε πλευρά».

Ο Ηλίας, που φοράει ρολόι με δείκτες, παρατηρεί ότι οι δείκτες σχηματίζουν γωνία: «Στις $3{:}00$ οι δείκτες είναι κάθετοι, στις $6{:}00$ ευθεία». Η Ρήνη ρωτάει: «Μπορούμε να ξέρουμε πόσες μοίρες είναι η γωνία των δεικτών κάθε ώρα;». Ο Στέλιος δουλεύει με τον κύκλο: «Ολόκληρος ο κύκλος έχει $360$ μοίρες. Το ρολόι έχει $12$ ώρες. Κάθε ώρα είναι $360:12=30$ μοίρες».

Σε αυτή την παράγραφο μαθαίνουμε να μετράμε γωνίες με μοιρογνωμόνιο, να τις συγκρίνουμε και να τις χωρίζουμε σε δύο ίσα μέρη με τη διχοτόμο.
\end{anagke}

\begin{lexi}{μοίρα, μοιρογνωμόνιο, διχοτόμος}{μοίρα, μοιρογνωμόνιο, διχοτόμος}
Η \lexiword{μοίρα} είναι το «μέρος» (από το «μοίρα», μερίδιο): η μοίρα ($1^\circ$) είναι το $\frac{1}{360}$ του πλήρους κύκλου. Το \lexiword{μοιρογνωμόνιο} «γνωρίζει» (μετράει) τις μοίρες. Η \lexiword{διχοτόμος} «χωρίζει στα δύο» (δίχα + τέμνω) μια γωνία σε δύο ίσες γωνίες.
\end{lexi}

\begin{metafrasi}
Το $\frac{1}{360}$ του κύκλου. & $1^\circ$\\
Η υποδιαίρεση της μοίρας. & $1^\circ=60'$ και $1'=60''$\\
Η γωνία ανοίγει $84$ μοίρες. & $x\mathrm{Ο}y=84^\circ$\\
Η $\mathrm{Οδ}$ είναι διχοτόμος. & $x\mathrm{Ο}\delta=\delta\mathrm{Ο}y=\dfrac{x\mathrm{Ο}y}{2}$\\
Κάθε ώρα στο ρολόι. & $360^\circ:12=30^\circ$\\
\end{metafrasi}

\tmima{Η θεωρία}

\tmimasmall{Μέτρηση γωνιών}

Για να \textbf{μετρήσουμε} μια γωνία τη συγκρίνουμε με μια μονάδα, τη \textbf{μοίρα}. Μία μοίρα ($1^\circ$) είναι το $\frac{1}{360}$ του πλήρους κύκλου. Υποδιαιρείται σε \textbf{πρώτα λεπτά} και \textbf{δεύτερα λεπτά}: $1^\circ=60'$ και $1'=60''$. Για παράδειγμα $25^\circ30'=25{,}5^\circ$ και $48{,}25^\circ=48^\circ15'$, γιατί $0{,}25\cdot60=15$. Μετράμε με το \textbf{μοιρογνωμόνιο}: τοποθετούμε το κέντρο του στην κορυφή και τη μία πλευρά στο μηδέν, και διαβάζουμε πού περνά η άλλη πλευρά. Προσέχουμε σε ποια από τις δύο κλίμακες ξεκινά το μηδέν.

\tmimasmall{Σύγκριση και ισότητα}

Δύο γωνίες είναι \textbf{ίσες} όταν εφαρμόζουν η μία στην άλλη (έχουν το ίδιο μέτρο). Η μία είναι \textbf{μεγαλύτερη} από την άλλη όταν το μέτρο της είναι μεγαλύτερο, ή όταν η μικρή χωράει στο εσωτερικό της μεγάλης με κοινή κορυφή και μία κοινή πλευρά. Το μέτρο μιας γωνίας δεν εξαρτάται από το μήκος των πλευρών της (οι πλευρές είναι ημιευθείες).

\tmimasmall{Διχοτόμος}

\begin{orismos}{διχοτόμος γωνίας}
\textbf{Διχοτόμος} της γωνίας $x\mathrm{Ο}y$ είναι η ημιευθεία $\mathrm{Οδ}$ που βρίσκεται στο εσωτερικό της και τη χωρίζει σε δύο ίσες γωνίες: $x\mathrm{Ο}\delta=\delta\mathrm{Ο}y$. Τότε κάθε μία είναι το μισό της $x\mathrm{Ο}y$.
\end{orismos}

\begin{center}
\begin{tikzpicture}[scale=1.1]
\draw[thick] (3.5,0)--(0,0)--(2.6,2.2); \draw[gblue,thick] (0,0)--(3.2,1.1); \fill (0,0) circle (2pt);
\node[right] at (3.5,0) {$x$}; \node[above] at (2.6,2.2) {$y$}; \node[right,gblue] at (3.2,1.1) {$\delta$}; \node[below left] at (0,0) {$\mathrm{Ο}$};
\draw[gblue] (0.9,0) arc(0:19:0.9); \draw[gblue] (1.1,0.32) arc(19:40:1.1);
\draw (0.6,0.3) -- (0.6,0.3);
\end{tikzpicture}
\end{center}

Κάθε γωνία έχει μία μοναδική διχοτόμο. Την κατασκευάζουμε με μοιρογνωμόνιο (μετράμε, διαιρούμε με το $2$ και σημειώνουμε), ή με διαφανές χαρτί διπλώνοντας ώστε οι δύο πλευρές να συμπέσουν: η πτυχή είναι η διχοτόμος.

\begin{prosoxi}
\begin{itemize}[leftmargin=1.2em,itemsep=0.3ex,topsep=2pt]
\item Το μέτρο της γωνίας δεν εξαρτάται από το μήκος των πλευρών της: μια γωνία με μεγαλύτερο σχέδιο δεν είναι μεγαλύτερη.
\item Διαβάζω τη σωστή κλίμακα του μοιρογνωμονίου: αν η γωνία είναι οξεία το μέτρο της είναι μικρότερο από $90^\circ$.
\item Το $1^\circ=60'$ και όχι $100'$: $25^\circ30'=25{,}5^\circ$, όχι $25{,}30^\circ$.
\item Η διχοτόμος είναι ημιευθεία μέσα στη γωνία, όχι τμήμα που ενώνει δύο σημεία των πλευρών.
\end{itemize}
\end{prosoxi}

\tmima{Λυμένα παραδείγματα}

\paradeigma{Το μονοπάτι του κήπου}
\tixero Η γωνία του φράχτη είναι $84^\circ$.\par
\tithelo Η γωνία που σχηματίζει το μονοπάτι με κάθε πλευρά, αν ήθελαν να τη χωρίσουν σε δύο ίσα παρτέρια.\par
\possindeo Το μονοπάτι θα είναι η διχοτόμος. Κάθε γωνία είναι $84^\circ:2=42^\circ$.\par
\apantisi $42^\circ$.\par
\elegxos $42^\circ+42^\circ=84^\circ$.

\paradeigma{Οι δείκτες του ρολογιού}
\tixero Ο ωροδείκτης και ο λεπτοδείκτης στις $4{:}00$ και στις $5{:}00$.\par
\tithelo Η γωνία που σχηματίζουν.\par
\possindeo Κάθε ώρα αντιστοιχεί σε $360^\circ:12=30^\circ$. Στις $4{:}00$ ο λεπτοδείκτης δείχνει το $12$ και ο ωροδείκτης το $4$: $4\cdot30^\circ=120^\circ$. Στις $5{:}00$: $5\cdot30^\circ=150^\circ$.\par
\apantisi $120^\circ$ και $150^\circ$.\par
\elegxos Στις $6{:}00$ θα είναι $180^\circ$ (ευθεία).

\paradeigma{Μοίρες και λεπτά}
\tixero Οι γωνίες $25^\circ30'$ και $48{,}25^\circ$.\par
\tithelo Να τις γράψω στην άλλη μορφή.\par
\possindeo $30'=\frac{30}{60}^\circ=0{,}5^\circ$, άρα $25^\circ30'=25{,}5^\circ$. Και $0{,}25^\circ=0{,}25\cdot60'=15'$, άρα $48{,}25^\circ=48^\circ15'$.\par
\apantisi $25{,}5^\circ$ και $48^\circ15'$.\par
\elegxos $15'=\frac{15}{60}=0{,}25^\circ$.

\paradeigma{Σύγκριση γωνιών}
\tixero Οι γωνίες $42^\circ$, $4{,}5^\circ$, $38^\circ30'$, $38^\circ$.\par
\tithelo Η σειρά.\par
\possindeo Σε μοίρες: $42$, $4{,}5$, $38{,}5$, $38$.\par
\apantisi $4{,}5^\circ<38^\circ<38^\circ30'<42^\circ$.\par
\elegxos $30'=0{,}5^\circ$.

\paradeigma{Η διχοτόμος και το μέτρο}
\tixero Η $\mathrm{Οδ}$ είναι διχοτόμος της $x\mathrm{Ο}y$ και $x\mathrm{Ο}\delta=28^\circ$.\par
\tithelo Η $x\mathrm{Ο}y$.\par
\possindeo $\delta\mathrm{Ο}y=x\mathrm{Ο}\delta=28^\circ$, άρα $x\mathrm{Ο}y=28^\circ+28^\circ=56^\circ$.\par
\apantisi $56^\circ$.\par
\elegxos Η διχοτόμος της $56^\circ$ δίνει $28^\circ$.

\tmima{Ασκήσεις}

\zoni{Το κατάλαβα}{μοίρες, ρολόι, διχοτόμος}

\ex Σχεδίασε με μοιρογνωμόνιο γωνίες $40^\circ$, $75^\circ$ και $130^\circ$ και έλεγξε με το μέτρημα της διχοτόμου.

\ex Γράψε στην άλλη μορφή: α) $25^\circ30'$ σε μοίρες \quad β) $48{,}25^\circ$ σε μοίρες και λεπτά \quad γ) $3{,}5^\circ$ σε μοίρες και λεπτά \quad δ) $90^\circ$ σε λεπτά.

\ex Βάλε σε αύξουσα σειρά: $42^\circ$, $38^\circ30'$, $38^\circ$, $4{,}5^\circ$, $0{,}7^\circ$.

\ex Βρες τη διχοτόμο: α) $x\mathrm{Ο}y=70^\circ$ \quad β) $x\mathrm{Ο}y=124^\circ$ \quad γ) $x\mathrm{Ο}y=8^\circ30'$. Ποιο είναι το μέτρο κάθε μισής γωνίας;

\ex Η $\mathrm{Οδ}$ είναι διχοτόμος της $x\mathrm{Ο}y$ και $x\mathrm{Ο}\delta=28^\circ$. Βρες την $x\mathrm{Ο}y$.

\ex Βρες τη γωνία των δεικτών του ρολογιού στις: α) $3{:}00$ \quad β) $6{:}00$ \quad γ) $9{:}00$ \quad δ) $2{:}00$ \quad ε) $5{:}00$ \quad στ) $12{:}00$. (Πάρε τη μικρότερη γωνία.)

\ex Σωστό ή λάθος; α) $1^\circ=60'$. \quad β) Η διχοτόμος χωρίζει τη γωνία σε δύο άνισα μέρη. \quad γ) Δύο γωνίες είναι ίσες αν εφαρμόζουν. \quad δ) Η διχοτόμος είναι ημιευθεία. \quad ε) $1'=60''$.

\zoni{Το χρησιμοποιώ}{κήπος, ρολόι, πίτσα}

\ex \textbf{Το μονοπάτι.} Η γωνία του φράχτη είναι $84^\circ$. Με ποια γωνία ξεκινά το μονοπάτι που τη χωρίζει σε δύο ίσα παρτέρια;

\ex \textbf{Το ρολόι.} Πότε (ποια ώρα, με τον λεπτοδείκτη στο $12$) οι δείκτες σχηματίζουν γωνία $150^\circ$; Υπάρχουν δύο ώρες;

\ex \textbf{Η πίτσα.} Μια πίτσα κόβεται σε $8$ ίσες φέτες. Πόση γωνία έχει η κορυφή κάθε φέτας; Αν κόβεται σε $6$ φέτες;

\ex \textbf{Ο ανεμόμυλος.} Ένας ανεμόμυλος έχει $3$ πτερύγια με ίσες γωνίες μεταξύ τους. Πόση γωνία σχηματίζουν δύο διαδοχικά;

\ex \textbf{Το μοιρογνωμόνιο.} Ο Στέλιος διάβασε $140^\circ$ σε μια γωνία που φαίνεται οξεία. Τι πιθανότατα έκανε; Ποιο είναι το σωστό;

\zoni{Το σκέφτομαι}{γιατί $360$;}

\ex Οι Βαβυλώνιοι χωρίζουν τον κύκλο σε $360$ μοίρες. Γιατί ο $360$ είναι βολικός αριθμός (σκέψου τους διαιρέτες του);

\ex Αν διπλασιάσεις το μήκος των πλευρών μιας γωνίας, αλλάζει το μέτρο της; Εξήγησε.

\ex Γιατί κάθε γωνία έχει μία και μοναδική διχοτόμο;

\ex Η διχοτόμος μιας ορθής γωνίας ($90^\circ$) πόσες μοίρες σχηματίζει με κάθε πλευρά;

\zoni{Για το σπίτι}{δίπλωσε το χαρτί}

\ex Κόψε μια γωνία από χαρτί και βρες τη διχοτόμο διπλώνοντάς την. Μέτρησε με μοιρογνωμόνιο.

\ex Μέτρησε με μοιρογνωμόνιο τρεις γωνίες στο σπίτι και γράψε τα μέτρα τους.

\paragrafos{Είδη γωνιών και κάθετες ευθείες}{Β.1.6, «Είδη γωνιών – Κάθετες ευθείες»}

\begin{anagke}
Η Α3 πρέπει να στήσει τον φράχτη ώστε ο κήπος να έχει «σωστές γωνίες» στις γωνιές του. «Σωστή γωνία σημαίνει τετράγωνη, όπως στο θρανίο», λέει ο Νίκος. Η Αγγελική το λέει επίσημα: ορθή γωνία. Αλλά πώς φτιάχνουν ορθή γωνία στο έδαφος, χωρίς μοιρογνωμόνιο; Ο παππούς δείχνει ένα κόλπο που ξέρουν οι χτίστες από την αρχαιότητα: ένα σχοινί με $12$ ίσα διαστήματα, με κόμπους. Το διπλώνει σε τρίγωνο με πλευρές $3$, $4$ και $5$ διαστημάτων: «Η γωνία ανάμεσα στο $3$ και στο $4$ είναι ορθή».

Η Πολυξένη το καταγράφει και ρωτάει για όλες τις άλλες γωνίες που βλέπει: οξείες, αμβλείες, όπως στο φτερό του αεροπλάνου και στη στέγη. Ο Βασίλης θέλει να ταξινομήσει τις γωνίες του κήπου: «Αυτή που ανοίγει λίγο», «αυτή που ανοίγει πολύ», «η ευθεία». Η κυρία Αλεξάνδρα του δίνει ονόματα. Και όταν δύο μονοπάτια τέμνονται σε ορθή γωνία, τα λέμε κάθετα.

Σε αυτή την παράγραφο ταξινομούμε τις γωνίες ανάλογα με το μέτρο τους και ορίζουμε την καθετότητα, μια από τις πιο χρήσιμες έννοιες στη Γεωμετρία.
\end{anagke}

\begin{lexi}{ορθή, οξεία, αμβλεία, κάθετος}{ορθή, οξεία, αμβλεία}
\lexiword{Ορθή} σημαίνει «σωστή, όρθια» (από το «ορθός»): η γωνία ενός όρθιου τοίχου με το δάπεδο. \lexiword{Οξεία} είναι η «αιχμηρή», σαν την οξεία μύτη. \lexiword{Αμβλεία} είναι η «αμβλυμένη», η στρογγυλεμένη, που ανοίγει περισσότερο από την ορθή. \lexiword{Κάθετη} είναι η ευθεία που πέφτει «κάθετα», σαν το νήμα της στάθμης.
\end{lexi}

\begin{metafrasi}
Γωνία μικρότερη από την ορθή. & οξεία: $0^\circ<\omega<90^\circ$\\
Γωνία $90^\circ$. & ορθή\\
Γωνία ανάμεσα στην ορθή και την ευθεία. & αμβλεία: $90^\circ<\omega<180^\circ$\\
Γωνία $180^\circ$. & ευθεία γωνία\\
Δύο ευθείες που τέμνονται σε ορθή γωνία. & $\varepsilon_1\perp\varepsilon_2$\\
\end{metafrasi}

\tmima{Η θεωρία}

\begin{orismos}{είδη γωνιών}
Μια γωνία $\omega$ λέγεται:
\begin{itemize}[leftmargin=1.2em,itemsep=0.3ex,topsep=2pt]
\item \textbf{οξεία} αν $0^\circ<\omega<90^\circ$,
\item \textbf{ορθή} αν $\omega=90^\circ$,
\item \textbf{αμβλεία} αν $90^\circ<\omega<180^\circ$,
\item \textbf{ευθεία} αν $\omega=180^\circ$ (οι πλευρές της είναι αντικείμενες ημιευθείες),
\item \textbf{μη κυρτή} αν $180^\circ<\omega<360^\circ$ και \textbf{πλήρης} αν $\omega=360^\circ$,
\item \textbf{μηδενική} αν $\omega=0^\circ$ (οι πλευρές συμπίπτουν).
\end{itemize}
\end{orismos}

\begin{center}
\begin{tikzpicture}[scale=0.8]
\foreach \a/\n/\x in {40/οξεία/0,90/ορθή/3.6,130/αμβλεία/7.2}{
\begin{scope}[xshift=\x cm]
\draw[thick] (2.2,0)--(0,0)--({2.2*cos(\a)},{2.2*sin(\a)});
\draw[gblue] (0.6,0) arc(0:\a:0.6);
\node at (1.1,-0.7) {\n};
\end{scope}}
\draw[gblue] (3.6+0.25,0)--(3.6+0.25,0.25)--(3.6,0.25);
\end{tikzpicture}
\end{center}

\tmimasmall{Κάθετες ευθείες}

Δύο ευθείες που τέμνονται σχηματίζουν τέσσερις γωνίες. Αν μία από αυτές είναι ορθή, τότε είναι ορθές και οι τέσσερις. Τότε οι ευθείες λέγονται \textbf{κάθετες} και γράφουμε $\varepsilon_1\perp\varepsilon_2$.

\begin{idiotita}{ύπαρξη και μοναδικότητα κάθετης}
Από ένα σημείο του επιπέδου διέρχεται \textbf{μία και μοναδική} ευθεία κάθετη προς μια δοσμένη ευθεία $(\varepsilon)$, είτε το σημείο ανήκει στην ευθεία είτε όχι.
\end{idiotita}

Τις κάθετες τις σχεδιάζουμε με \textbf{γνώμονα} (τρίγωνο με ορθή γωνία) και χάρακα.

\begin{prosoxi}
\begin{itemize}[leftmargin=1.2em,itemsep=0.3ex,topsep=2pt]
\item Η ορθή γωνία είναι ακριβώς $90^\circ$. Η $91^\circ$ είναι αμβλεία και η $89^\circ$ οξεία.
\item Οι κάθετες δεν χρειάζεται να είναι «οριζόντιες και κατακόρυφες»: δύο ευθείες που βρίσκονται σε οποιαδήποτε θέση είναι κάθετες αν σχηματίζουν ορθή γωνία.
\item Η ευθεία γωνία δεν είναι «ευθεία γραμμή» αλλά γωνία $180^\circ$ με αντικείμενες πλευρές.
\item Από σημείο εκτός ευθείας διέρχεται \emph{μία} κάθετη, αλλά άπειρες πλάγιες.
\end{itemize}
\end{prosoxi}

\tmima{Λυμένα παραδείγματα}

\paradeigma{Ταξινομώ γωνίες}
\tixero Οι γωνίες $37^\circ$, $90^\circ$, $123^\circ$, $180^\circ$, $250^\circ$, $360^\circ$.\par
\tithelo Να πω το είδος κάθε μίας.\par
\possindeo Συγκρίνω με τα $90^\circ$, $180^\circ$, $360^\circ$.\par
\apantisi $37^\circ$ οξεία, $90^\circ$ ορθή, $123^\circ$ αμβλεία, $180^\circ$ ευθεία, $250^\circ$ μη κυρτή, $360^\circ$ πλήρης.\par
\elegxos Η $250^\circ$ είναι ανάμεσα στο $180^\circ$ και στο $360^\circ$.

\paradeigma{Το ρολόι}
\tixero Οι δείκτες στις $2{:}00$, $3{:}00$, $4{:}00$.\par
\tithelo Το είδος της γωνίας.\par
\possindeo Κάθε ώρα είναι $30^\circ$. Στις $2{:}00$: $60^\circ$ οξεία. Στις $3{:}00$: $90^\circ$ ορθή. Στις $4{:}00$: $120^\circ$ αμβλεία.\par
\apantisi Οξεία, ορθή, αμβλεία.\par
\elegxos Στις $6{:}00$ θα ήταν ευθεία γωνία.

\paradeigma{Κάθετες ευθείες}
\tixero Δύο ευθείες που τέμνονται σχηματίζοντας γωνία $90^\circ$.\par
\tithelo Οι υπόλοιπες γωνίες.\par
\possindeo Οι διαδοχικές γωνίες ανήκουν σε ευθεία γωνία $180^\circ$: $180^\circ-90^\circ=90^\circ$. Έτσι και οι τέσσερις είναι ορθές.\par
\apantisi $90^\circ$ η καθεμία.\par
\elegxos Το άθροισμα και των τεσσάρων είναι $360^\circ$.

\paradeigma{Η ορθή γωνία του παππού}
\tixero Σχοινί με $12$ ίσα διαστήματα, διπλωμένο σε τρίγωνο με πλευρές $3$, $4$, $5$.\par
\tithelo Γιατί σχηματίζεται ορθή γωνία.\par
\possindeo Το τρίγωνο με πλευρές $3$, $4$, $5$ έχει ορθή γωνία ανάμεσα στις πλευρές $3$ και $4$ (το αποδεικνύουμε αργότερα). Οι αρχαίοι αρπεδονάπτες το χρησιμοποιούσαν για τις γωνίες των κτιρίων και των αγρών.\par
\apantisi Γιατί $3^2+4^2=5^2$ ($9+16=25$).\par
\elegxos Δοκίμασε με γνώμονα.

\tmima{Ασκήσεις}

\zoni{Το κατάλαβα}{είδη γωνιών και κάθετες}

\ex Ταξινόμησε τις γωνίες $37^\circ$, $90^\circ$, $123^\circ$, $180^\circ$, $250^\circ$, $360^\circ$, $0^\circ$, $89^\circ$.

\ex Σχεδίασε με μοιρογνωμόνιο μια οξεία, μια ορθή και μια αμβλεία γωνία και μέτρησε τις διχοτόμους τους.

\ex Βρες το είδος της γωνίας των δεικτών του ρολογιού στις $2{:}00$, $3{:}00$, $4{:}00$, $6{:}00$, $7{:}00$ (πάρε τη μικρότερη).

\ex Δύο ευθείες τέμνονται και μια γωνία είναι $90^\circ$. Πόσες είναι οι άλλες γωνίες;

\ex Με γνώμονα σχεδίασε μια ευθεία $(\varepsilon)$, ένα σημείο $\mathrm{Α}$ πάνω της και ένα σημείο $\mathrm{Β}$ εκτός. Φέρε από το $\mathrm{Α}$ και από το $\mathrm{Β}$ τις κάθετες προς την $(\varepsilon)$.

\ex Πόσες κάθετες διέρχονται από ένα σημείο προς μια ευθεία; Και πόσες πλάγιες;

\ex Σωστό ή λάθος; α) Μια γωνία $91^\circ$ είναι αμβλεία. \quad β) Η ευθεία γωνία είναι $180^\circ$. \quad γ) Δύο κάθετες ευθείες σχηματίζουν γωνίες που δεν είναι όλες ίσες. \quad δ) Η $0^\circ$ είναι οξεία. \quad ε) Η $250^\circ$ είναι κυρτή.

\zoni{Το χρησιμοποιώ}{γωνιές, στέγες, αναβάτες}

\ex \textbf{Ο φράχτης.} Στις γωνιές του κήπου ο παππούς θέλει ορθές γωνίες. Πώς θα το ελέγξεις με ένα φύλλο χαρτί $\mathrm{Α4}$;

\ex \textbf{Το σχοινί του παππού.} Ένα σχοινί με $12$ ίσα διαστήματα διπλώνεται σε τρίγωνο με πλευρές $3$, $4$, $5$. Τι γωνία σχηματίζεται και πού;

\ex \textbf{Τα δωμάτια.} Βρες στην τάξη τρεις ορθές, δύο οξείες και μία αμβλεία γωνία.

\ex \textbf{Η στέγη.} Μια στέγη έχει γωνία $110^\circ$ στην κορυφή. Τι είδους είναι; Αν η στέγη ήταν ορθή γωνία πόσες μοίρες θα ήταν;

\ex \textbf{Ο ανεμοδείκτης.} Ένας ανεμοδείκτης στρέφεται από τη διεύθυνση Βορράς προς Ανατολή. Τι γωνία διέγραψε; Και προς Νότο;

\zoni{Το σκέφτομαι}{όταν μία ορθή φέρνει τρεις}

\ex Δείξε ότι αν μία από τις γωνίες που σχηματίζουν δύο τεμνόμενες ευθείες είναι ορθή, είναι ορθές και οι άλλες τρεις.

\ex Είναι δυνατόν να έχει ένα τρίγωνο δύο ορθές γωνίες; Δοκίμασε να σχεδιάσεις.

\ex Γιατί ο γνώμονας φτιάχνει κάθετες; Πώς θα έλεγχες ότι είναι σωστός (σκέψου την αντιστροφή του);

\ex Αν η γωνία $\omega$ είναι οξεία, τι είδους είναι η $180^\circ-\omega$;

\zoni{Για το σπίτι}{ψάξε γωνίες}

\ex Βρες στο σπίτι πέντε αντικείμενα με ορθές γωνίες και τρία με αμβλείες.

\ex Φτιάξε με σπάγκο και κόμπους το τρίγωνο $3,\,4,\,5$ και έλεγξε με ορθογώνιο χαρτί αν η γωνία είναι ορθή.

\paragrafos{Εφεξής και διαδοχικές γωνίες: οι γωνίες που προστίθενται}{Β.1.7, «Εφεξής και διαδοχικές γωνίες – Άθροισμα γωνιών»}

\begin{anagke}
Στο κέντρο του κήπου θα υπάρχει το συντριβάνι και από αυτό θα ξεκινούν πέντε μονοπάτια προς διαφορετικές κατευθύνσεις. Η Αγγελική σχεδιάζει το σχήμα και μετράει τις γωνίες ανάμεσα στα διαδοχικά μονοπάτια: $60^\circ$, $75^\circ$, $50^\circ$, $90^\circ$. «Η πέμπτη γωνία δεν χωρούσε στο μοιρογνωμόνιο», λέει, «γιατί είναι μη κυρτή;». «Όχι», απαντάει η Πολυξένη, «απλώς πρέπει να βρούμε πόσο μένει μέχρι τον κύκλο. Ολόκληρος ο κύκλος έχει $360^\circ$».

Ο Βασίλης προσθέτει: $60+75+50+90=275$, άρα η πέμπτη είναι $360-275=85$. «Μπορεί να μην βρίσκομαι εκεί, αλλά είμαι σίγουρος!». Ο Στέλιος ρωτάει αν οι γωνίες που έχουν κοινή πλευρά και «χωρίζουν» το μέρος που έχουν κοινή κορυφή έχουν κάποιο όνομα. Η κυρία Αλεξάνδρα απαντάει: «Λέγονται εφεξής, και όταν είναι πολλές στη σειρά, διαδοχικές. Τα μέτρα τους προστίθενται».

Στην παράγραφο αυτή μαθαίνουμε πώς προσθέτουμε και αφαιρούμε γωνίες που έχουν κοινή κορυφή και κοινή πλευρά, και πώς χρησιμοποιούμε το γεγονός ότι γύρω από ένα σημείο υπάρχουν $360^\circ$ και πάνω σε μια ευθεία $180^\circ$.
\end{anagke}

\begin{lexi}{εφεξής, διαδοχικές}{εφεξής, διαδοχικές}
\lexiword{Εφεξής} σημαίνει «η μία μετά την άλλη, στη σειρά» (εφ' εξής). \lexiword{Διαδοχικές} είναι οι γωνίες που «διαδέχονται» η μία την άλλη, έτσι ώστε η κάθε μία να είναι εφεξής με την επόμενη.
\end{lexi}

\begin{metafrasi}
Δύο γωνίες με κοινή κορυφή και κοινή πλευρά, και χωριστά εσωτερικά. & εφεξής γωνίες $x\mathrm{Ο}y$, $y\mathrm{Ο}z$\\
Το άθροισμά τους. & $x\mathrm{Ο}z=x\mathrm{Ο}y+y\mathrm{Ο}z$\\
Όλες οι διαδοχικές γύρω από ένα σημείο. & άθροισμα $360^\circ$\\
Οι διαδοχικές πάνω σε ευθεία. & άθροισμα $180^\circ$\\
\end{metafrasi}

\tmima{Η θεωρία}

\begin{orismos}{εφεξής και διαδοχικές γωνίες}
Δύο γωνίες $x\mathrm{Ο}y$ και $y\mathrm{Ο}z$ λέγονται \textbf{εφεξής} όταν έχουν κοινή κορυφή $\mathrm{Ο}$, κοινή πλευρά $\mathrm{Οy}$ και τα εσωτερικά τους δεν έχουν κοινά σημεία (η μία είναι «πριν» και η άλλη «μετά» την $\mathrm{Οy}$). Πολλές γωνίες λέγονται \textbf{διαδοχικές} όταν καθεμία είναι εφεξής με την επόμενή της.
\end{orismos}

\begin{center}
\begin{tikzpicture}[scale=1.1]
\draw[thick] (3.4,0)--(0,0)--(2.4,2.4); \draw[thick,gblue] (0,0)--(0.6,3.1); \fill (0,0) circle (2pt);
\node[right] at (3.4,0) {$x$}; \node[right] at (2.4,2.4) {$y$}; \node[above] at (0.6,3.1) {$z$}; \node[below left] at (0,0) {$\mathrm{Ο}$};
\draw (0.8,0) arc(0:45:0.8); \draw[gblue] (0.55,0.55) arc(45:79:0.55);
\end{tikzpicture}
\end{center}

\tmimasmall{Άθροισμα γωνιών}

Αν οι $x\mathrm{Ο}y$ και $y\mathrm{Ο}z$ είναι εφεξής, η γωνία $x\mathrm{Ο}z$ λέγεται \textbf{άθροισμά} τους:
$$x\mathrm{Ο}z=x\mathrm{Ο}y+y\mathrm{Ο}z.$$
Και αντίστροφα η διαφορά: $y\mathrm{Ο}z=x\mathrm{Ο}z-x\mathrm{Ο}y$. Το άθροισμα των μέτρων μπορεί να ξεπεράσει τα $180^\circ$ (τότε η γωνία είναι μη κυρτή).

\begin{idiotita}{άθροισμα διαδοχικών γωνιών}
Το άθροισμα των διαδοχικών γωνιών γύρω από ένα σημείο, που καλύπτουν όλο το επίπεδο, είναι $360^\circ$. Το άθροισμα διαδοχικών γωνιών που βρίσκονται από την ίδια μεριά μιας ευθείας (πλευρές αντικείμενες ημιευθείες) είναι $180^\circ$.
\end{idiotita}

\begin{prosoxi}
\begin{itemize}[leftmargin=1.2em,itemsep=0.3ex,topsep=2pt]
\item Δύο γωνίες δεν είναι εφεξής αν δεν έχουν κοινή κορυφή ή αν τα εσωτερικά τους επικαλύπτονται.
\item Το άθροισμα γωνιών αθροίζει τα \emph{μέτρα}: $40^\circ+65^\circ=105^\circ$.
\item Στον κύκλο, όλο το άθροισμα είναι $360^\circ$, όχι $180^\circ$ (αυτό είναι για την ευθεία).
\item Για εφεξής γωνίες το άθροισμα μπορεί να ξεπεράσει τις $180^\circ$ αλλά η γωνία $x\mathrm{Ο}z$ είναι τότε μη κυρτή.
\end{itemize}
\end{prosoxi}

\tmima{Λυμένα παραδείγματα}

\paradeigma{Εφεξής γωνίες}
\tixero $x\mathrm{Ο}y=35^\circ$ και $y\mathrm{Ο}z=50^\circ$ εφεξής.\par
\tithelo Η $x\mathrm{Ο}z$.\par
\possindeo $x\mathrm{Ο}z=35^\circ+50^\circ=85^\circ$.\par
\apantisi $85^\circ$.\par
\elegxos $85^\circ-35^\circ=50^\circ$.

\paradeigma{Με αφαίρεση}
\tixero Οι $\mathrm{ΑΟΒ}$, $\mathrm{ΒΟΓ}$ εφεξής με $\mathrm{ΑΟΒ}=70^\circ$ και $\mathrm{ΑΟΓ}=130^\circ$.\par
\tithelo Η $\mathrm{ΒΟΓ}$.\par
\possindeo $\mathrm{ΒΟΓ}=\mathrm{ΑΟΓ}-\mathrm{ΑΟΒ}=130^\circ-70^\circ=60^\circ$.\par
\apantisi $60^\circ$.\par
\elegxos $70^\circ+60^\circ=130^\circ$.

\paradeigma{Τα πέντε μονοπάτια}
\tixero Τέσσερις διαδοχικές γωνίες γύρω από το συντριβάνι: $60^\circ$, $75^\circ$, $50^\circ$, $90^\circ$.\par
\tithelo Η πέμπτη γωνία.\par
\possindeo Το άθροισμα γύρω από το σημείο είναι $360^\circ$. Τέσσερις γωνίες: $60+75+50+90=275$. Άρα η πέμπτη είναι $360-275=85^\circ$.\par
\apantisi $85^\circ$.\par
\elegxos $275^\circ+85^\circ=360^\circ$.

\paradeigma{Γωνίες πάνω σε ευθεία}
\tixero Διαδοχικές γωνίες $40^\circ$, $65^\circ$ και $\omega$ πάνω από μια ευθεία (η τρίτη κλείνει στην ευθεία).\par
\tithelo Η $\omega$.\par
\possindeo Το άθροισμα είναι $180^\circ$: $40+65+\omega=180$, άρα $\omega=180-105=75^\circ$.\par
\apantisi $75^\circ$.\par
\elegxos $40+65+75=180$.

\paradeigma{Διχοτόμος και άθροισμα}
\tixero Η $\mathrm{Οδ}$ είναι διχοτόμος της $x\mathrm{Ο}y$ και $\delta\mathrm{Ο}y=37^\circ$. Η $y\mathrm{Ο}z=48^\circ$ είναι εφεξής στην $x\mathrm{Ο}y$.\par
\tithelo Η $x\mathrm{Ο}z$.\par
\possindeo $x\mathrm{Ο}y=2\cdot37^\circ=74^\circ$. $x\mathrm{Ο}z=74^\circ+48^\circ=122^\circ$.\par
\apantisi $122^\circ$.\par
\elegxos Η $\mathrm{Οδ}$ πλευρά της $x\mathrm{Ο}z$: $37+37+48=122$.

\tmima{Ασκήσεις}

\zoni{Το κατάλαβα}{άθροισμα και διαφορά γωνιών}

\ex Οι $x\mathrm{Ο}y=35^\circ$ και $y\mathrm{Ο}z=50^\circ$ είναι εφεξής. Βρες την $x\mathrm{Ο}z$.

\ex Οι $\mathrm{ΑΟΒ}$ και $\mathrm{ΒΟΓ}$ είναι εφεξής, $\mathrm{ΑΟΓ}=130^\circ$ και $\mathrm{ΑΟΒ}=70^\circ$. Βρες την $\mathrm{ΒΟΓ}$.

\ex Τέσσερις διαδοχικές γωνίες γύρω από σημείο είναι $60^\circ$, $75^\circ$, $50^\circ$, $90^\circ$. Βρες την πέμπτη.

\ex Οι $40^\circ$, $65^\circ$ και $\omega$ είναι διαδοχικές πάνω από ευθεία. Βρες την $\omega$.

\ex Τρεις διαδοχικές γωνίες γύρω από σημείο είναι ίσες. Πόσες μοίρες είναι η καθεμία;

\ex Οι γωνίες $x\mathrm{Ο}y=24^\circ30'$ και $y\mathrm{Ο}z=37^\circ45'$ είναι εφεξής. Βρες την $x\mathrm{Ο}z$.

\ex Σωστό ή λάθος; α) Δύο γωνίες με κοινή κορυφή είναι πάντα εφεξής. \quad β) Οι εφεξής γωνίες έχουν κοινή πλευρά. \quad γ) Διαδοχικές γωνίες γύρω από σημείο έχουν άθροισμα $360^\circ$. \quad δ) Διαδοχικές γωνίες πάνω από ευθεία έχουν άθροισμα $180^\circ$.

\zoni{Το χρησιμοποιώ}{κήπος, πίτσα, ρολόι}

\ex \textbf{Τα πέντε μονοπάτια.} Στο συντριβάνι καταλήγουν πέντε μονοπάτια και οι γωνίες είναι $60^\circ$, $75^\circ$, $50^\circ$, $90^\circ$ και $\omega$. Βρες την $\omega$.

\ex \textbf{Η πίτσα.} Μια πίτσα κόβεται σε $6$ φέτες: δύο από $75^\circ$ και τρεις από $50^\circ$. Πόσο είναι η έκτη;

\ex \textbf{Ο ήλιος.} Ο Ήλιος σηκώνεται από τον ορίζοντα $12^\circ$ κάθε $48$ λεπτά. Σε πόση ώρα θα φτάσει να ανεβεί $60^\circ$;

\ex \textbf{Το ρολόι.} Στις $4{:}30$ ο ωροδείκτης βρίσκεται στις $135^\circ$ από το $12$ και ο λεπτοδείκτης στο $6$ ($180^\circ$ από το $12$). Πόση γωνία σχηματίζουν;

\ex \textbf{Οι στίβοι.} Τρεις στίβοι ξεκινούν από το ίδιο σημείο σχηματίζοντας ίσες γωνίες. Πόση είναι η γωνία ανάμεσα σε δύο διαδοχικούς;

\zoni{Το σκέφτομαι}{όταν το σύνολο είναι σταθερό}

\ex Γιατί το άθροισμα των διαδοχικών γωνιών γύρω από σημείο είναι $360^\circ$ ανεξάρτητα από τον αριθμό των γωνιών;

\ex Είναι δυνατόν δύο εφεξής γωνίες να είναι και οι δύο αμβλείες; Εξήγησε.

\ex Τρεις διαδοχικές γωνίες γύρω από σημείο είναι $x$, $2x$, $3x$. Βρες το $x$.

\ex Η $\mathrm{Οδ}$ είναι διχοτόμος της $x\mathrm{Ο}y=60^\circ$ και η $y\mathrm{Ο}z=20^\circ$ είναι εφεξής προς αυτήν. Βρες τις $x\mathrm{Ο}z$ και $\delta\mathrm{Ο}z$.

\zoni{Για το σπίτι}{άθροισμα στο χαρτί}

\ex Σχεδίασε τέσσερις διαδοχικές γωνίες γύρω από σημείο και μέτρησέ τις. Πόσο είναι το άθροισμά τους;

\ex Βρες στο σπίτι ένα αντικείμενο με διαδοχικές γωνίες γύρω από σημείο (π.χ. ένα ρολόι ή μια ανεμογεννήτρια).

\paragrafos{Παραπληρωματικές, συμπληρωματικές και κατακορυφήν γωνίες}{Β.1.8, «Παραπληρωματικές και Συμπληρωματικές γωνίες – Κατακορυφήν γωνίες»}

\begin{anagke}
Δύο μονοπάτια του κήπου διασταυρώνονται σαν ένα μεγάλο «Χ». Η Αγγελική μετράει τη μία γωνία και βρίσκει $48^\circ$. «Πόσες είναι οι άλλες τρεις;» ρωτάει. Ο Νίκος μετράει όλες: $48^\circ$, $132^\circ$, $48^\circ$, $132^\circ$. «Οι απέναντι είναι ίσες! Και οι διπλανές κάνουν μαζί $180^\circ$». «Είναι τυχαίο;» ρωτάει η Πολυξένη. «Μήπως αν κάνω το μονοπάτι πιο λοξό, θα αλλάξει;». Δοκιμάζουν με άλλη γωνία, $70^\circ$, και βρίσκουν $70^\circ$, $110^\circ$, $70^\circ$, $110^\circ$.

Η Ρήνη ρωτάει για μια γωνία που τη συμπληρώνουμε μέχρι την ορθή: ένα παρτέρι σχήματος γωνίας $35^\circ$ και το υπόλοιπο μέχρι τη γωνία του φράχτη των $90^\circ$ θα φυτευτεί με τριαντάφυλλα: «Πόσες μοίρες μένουν;». Ο Στέλιος υπολογίζει: «Εδώ το άθροισμα είναι $90^\circ$ και όχι $180^\circ$».

Έτσι η παρέα συναντά τρεις έννοιες: τις παραπληρωματικές γωνίες, που αθροίζουν $180^\circ$, τις συμπληρωματικές, που αθροίζουν $90^\circ$, και τις κατακορυφήν, που βρίσκονται απέναντι όταν δύο ευθείες τέμνονται. Θα βρουν και γιατί οι κατακορυφήν γωνίες είναι ίσες, χωρίς να χρειάζονται μοιρογνωμόνιο.
\end{anagke}

\begin{lexi}{παραπληρωματική, συμπληρωματική, κατακορυφήν}{παραπληρωματική, συμπληρωματική, κατακορυφήν}
Η \lexiword{συμπληρωματική} γωνία «συμπληρώνει» μια γωνία μέχρι την ορθή. Η \lexiword{παραπληρωματική} «παρα-πληρώνει», δηλαδή «συμπληρώνει ως τη μισή στροφή»: ως τις $180^\circ$. Οι \lexiword{κατακορυφήν} γωνίες είναι αυτές που βρίσκονται «κατά κορυφήν», δηλαδή απέναντι, με κοινή κορυφή.
\end{lexi}

\begin{metafrasi}
Δύο γωνίες που αθροίζουν $90^\circ$. & συμπληρωματικές: $\omega+\varphi=90^\circ$\\
Δύο γωνίες που αθροίζουν $180^\circ$. & παραπληρωματικές: $\omega+\varphi=180^\circ$\\
Οι απέναντι γωνίες δύο τεμνόμενων ευθειών. & κατακορυφήν: ίσες\\
Η συμπληρωματική της $35^\circ$. & $90^\circ-35^\circ=55^\circ$\\
Η παραπληρωματική της $130^\circ$. & $180^\circ-130^\circ=50^\circ$\\
\end{metafrasi}

\tmima{Η θεωρία}

\begin{orismos}{συμπληρωματικές, παραπληρωματικές}
Δύο γωνίες λέγονται \textbf{συμπληρωματικές} αν το άθροισμά τους είναι $90^\circ$ και \textbf{παραπληρωματικές} αν το άθροισμά τους είναι $180^\circ$. Η καθεμία λέγεται συμπλήρωμα (ή παραπλήρωμα) της άλλης.
\end{orismos}

Για παράδειγμα η συμπληρωματική της $35^\circ$ είναι η $55^\circ$ και η παραπληρωματική της $130^\circ$ η $50^\circ$. Δύο εφεξής παραπληρωματικές γωνίες έχουν τις μη κοινές πλευρές τους αντικείμενες ημιευθείες. Δύο συμπληρωματικές γωνίες είναι πάντα οξείες, ενώ από δύο παραπληρωματικές η μία είναι αμβλεία και η άλλη οξεία (ή και οι δύο ορθές).

\tmimasmall{Κατακορυφήν γωνίες}

Όταν δύο ευθείες τέμνονται σχηματίζονται τέσσερις γωνίες. Δύο από αυτές που δεν είναι διπλανές λέγονται \textbf{κατακορυφήν}.

\begin{center}
\begin{tikzpicture}[scale=1.1]
\draw[thick,{Stealth}-{Stealth}] (-2.5,-1)--(2.5,1); \draw[thick,{Stealth}-{Stealth}] (-1.2,-1.8)--(1.2,1.8);
\fill (0,0) circle (2pt); \node[below right] at (0.05,0) {$\mathrm{Ο}$};
\node[gblue] at (1.1,0.18) {$\omega$}; \node[gblue] at (-1.1,-0.18) {$\omega$}; \node at (0.18,0.9) {$\varphi$}; \node at (-0.18,-0.9) {$\varphi$};
\end{tikzpicture}
\end{center}

\begin{idiotita}{κατακορυφήν γωνίες}
Οι κατακορυφήν γωνίες είναι \textbf{ίσες}.
\end{idiotita}

Γιατί: η $\omega$ και η διπλανή της $\varphi$ είναι παραπληρωματικές ($\omega+\varphi=180^\circ$). Η κατακορυφήν της $\omega$, έστω $\omega'$, είναι κι αυτή παραπληρωματική της $\varphi$ ($\omega'+\varphi=180^\circ$). Άρα $\omega=180^\circ-\varphi=\omega'$.

\begin{prosoxi}
\begin{itemize}[leftmargin=1.2em,itemsep=0.3ex,topsep=2pt]
\item Συμπληρωματικές είναι $90^\circ$, παραπληρωματικές $180^\circ$. Μην τις μπερδεύεις: η «παραπληρωματική» είναι η μεγαλύτερη λέξη και το μεγαλύτερο άθροισμα ($180^\circ$).
\item Δεν υπάρχει συμπληρωματική γωνία για γωνία μεγαλύτερη ή ίση με $90^\circ$ (θα έβγαινε μηδενική ή αρνητική).
\item Οι διπλανές γωνίες δύο τεμνόμενων ευθειών δεν είναι ίσες (εκτός αν είναι ορθές): είναι παραπληρωματικές.
\item Δύο γωνίες που αθροίζουν $180^\circ$ δεν είναι αναγκαστικά εφεξής: αρκεί το μέτρο.
\end{itemize}
\end{prosoxi}

\tmima{Λυμένα παραδείγματα}

\paradeigma{Συμπλήρωμα και παραπλήρωμα}
\tixero Η γωνία $35^\circ$ και η γωνία $130^\circ$.\par
\tithelo Η συμπληρωματική της πρώτης και η παραπληρωματική της δεύτερης.\par
\possindeo $90^\circ-35^\circ=55^\circ$. $180^\circ-130^\circ=50^\circ$.\par
\apantisi $55^\circ$ και $50^\circ$.\par
\elegxos $35+55=90$ και $130+50=180$.

\paradeigma{Δύο τεμνόμενες ευθείες}
\tixero Δύο ευθείες τέμνονται και μία γωνία είναι $48^\circ$.\par
\tithelo Οι άλλες τρεις γωνίες.\par
\possindeo Η κατακορυφήν της είναι $48^\circ$. Οι διπλανές είναι παραπληρωματικές της: $180^\circ-48^\circ=132^\circ$, και η απέναντί τους επίσης $132^\circ$.\par
\apantisi $48^\circ$, $132^\circ$, $132^\circ$.\par
\elegxos Το άθροισμα και των τεσσάρων: $48+132+48+132=360$.

\paradeigma{Παραπληρωματικές με λόγο}
\tixero Δύο παραπληρωματικές γωνίες έχουν λόγο $2:3$.\par
\tithelo Να τις βρω.\par
\possindeo Αν η μία είναι $2x$ και η άλλη $3x$, $2x+3x=180$, $5x=180$, $x=36$. Οι γωνίες είναι $72^\circ$ και $108^\circ$.\par
\apantisi $72^\circ$ και $108^\circ$.\par
\elegxos $72+108=180$ και $\frac{72}{108}=\frac{2}{3}$.

\paradeigma{Εξίσωση με κατακορυφήν}
\tixero Οι κατακορυφήν γωνίες είναι $2x+10$ και $3x-20$ μοίρες.\par
\tithelo Το $x$ και οι γωνίες.\par
\possindeo Ίσες: $2x+10=3x-20$, άρα $x=30$ και οι γωνίες είναι $2\cdot30+10=70^\circ$.\par
\apantisi $x=30$ και γωνίες $70^\circ$.\par
\elegxos $3\cdot30-20=70$.

\tmima{Ασκήσεις}

\zoni{Το κατάλαβα}{συμπληρωματικές, παραπληρωματικές, κατακορυφήν}

\ex Βρες το συμπλήρωμα των γωνιών: $35^\circ$, $62^\circ$, $8^\circ$, $44^\circ30'$.

\ex Βρες το παραπλήρωμα των γωνιών: $130^\circ$, $90^\circ$, $5^\circ$, $73^\circ15'$.

\ex Δύο ευθείες τέμνονται και μία γωνία είναι $48^\circ$. Βρες τις άλλες τρεις.

\ex Δύο ευθείες τέμνονται και μία γωνία είναι ορθή. Βρες τις άλλες.

\ex Βρες γωνία που είναι ίση με το συμπλήρωμά της. Βρες γωνία ίση με το παραπλήρωμά της.

\ex Δύο παραπληρωματικές γωνίες έχουν λόγο $2:3$. Βρες τες.

\ex Σωστό ή λάθος; α) Οι κατακορυφήν γωνίες είναι ίσες. \quad β) Οι διπλανές γωνίες δύο τεμνόμενων ευθειών είναι συμπληρωματικές. \quad γ) Υπάρχει συμπληρωματική της $100^\circ$. \quad δ) Δύο συμπληρωματικές γωνίες είναι οξείες. \quad ε) Οι παραπληρωματικές γωνίες έχουν άθροισμα $180^\circ$.

\zoni{Το χρησιμοποιώ}{διασταυρώσεις, παρτέρια, καθρέφτες}

\ex \textbf{Η διασταύρωση.} Δύο μονοπάτια του κήπου τέμνονται και μια γωνία είναι $70^\circ$. Πόσες μοίρες έχουν οι άλλες γωνίες; Πόσες διαφορετικές τιμές υπάρχουν;

\ex \textbf{Τα τριαντάφυλλα.} Στη γωνία $90^\circ$ του φράχτη φυτεύεται ένα παρτέρι με γωνία $35^\circ$. Πόση γωνία μένει για τα τριαντάφυλλα;

\ex \textbf{Ο καθρέφτης.} Ένα χέρι του κήπου σχηματίζει γωνία $\omega=2x+10$ και η κατακορυφήν της είναι $3x-20$. Βρες το $x$ και την $\omega$.

\ex \textbf{Η σκάλα.} Μια σκάλα ακουμπά σε έναν τοίχο σχηματίζοντας γωνία $72^\circ$ με το έδαφος. Πόση γωνία σχηματίζει με τον τοίχο; (Ο τοίχος είναι κάθετος στο έδαφος.)

\ex \textbf{Το ψαλίδι.} Ένα ψαλίδι ανοίγει ώστε οι λεπίδες να σχηματίζουν γωνία $40^\circ$. Πόση είναι η γωνία των λαβών (κατακορυφήν);

\zoni{Το σκέφτομαι}{γιατί ίσες}

\ex Εξήγησε γιατί οι κατακορυφήν γωνίες είναι ίσες, χρησιμοποιώντας τα παραπληρώματα.

\ex Δύο γωνίες είναι ίσες. Είναι απαραίτητα κατακορυφήν; Δώσε αντιπαράδειγμα.

\ex Αν η $\omega$ είναι οξεία, τι είδους είναι η παραπληρωματική της; Και η συμπληρωματική της;

\ex Οι διχοτόμοι δύο εφεξής παραπληρωματικών γωνιών: πόση γωνία σχηματίζουν μεταξύ τους; (δοκίμασε με $\omega=50^\circ$).

\zoni{Για το σπίτι}{ψάξε Χ}

\ex Σχεδίασε δύο τεμνόμενες ευθείες, μέτρησε και τις τέσσερις γωνίες και έλεγξε ότι οι κατακορυφήν είναι ίσες.

\ex Βρες στο σπίτι ένα ζευγάρι τεμνόμενων ευθειών (π.χ. ψαλίδι, παραθύρων) και μέτρησε τις γωνίες.

\paragrafos{Θέσεις ευθειών στο επίπεδο: τεμνόμενες και παράλληλες}{Β.1.9, «Θέσεις ευθειών στο επίπεδο»}

\begin{anagke}
Ο παππούς δείχνει πώς θα φυτευτούν οι σειρές με τα καρότα στον κήπο: «Κάθε σειρά θα είναι ευθεία και όλες θα έχουν την ίδια κατεύθυνση. Καμιά δεν θα ακουμπά την άλλη όσο κι αν τις συνεχίσετε». Η Αγγελική ξαφνικά σκέφτεται τις γραμμές στο σχολικό γήπεδο, τις γραμμές στο τετράδιο και τις ράγες του τρένου: «Δεν συναντιούνται ποτέ».

Ο Ηλίας δείχνει δύο μονοπάτια του κήπου που διασταυρώνονται: «Αυτά ακουμπούν σε ένα σημείο». «Μα αν συνεχίσεις ένα μονοπάτι που δεν διασταυρώνεται, μπορεί να διασταυρωθεί αργότερα!» διαμαρτύρεται ο Στέλιος. Η Πολυξένη προτείνει ένα τεστ: «Σκέψου τις ευθείες απεριόριστες. Δύο ευθείες του ίδιου επιπέδου είτε έχουν ένα κοινό σημείο είτε δεν έχουν κανένα».

Στην παράγραφο αυτή ονομάζουμε τις δύο δυνατές θέσεις: τεμνόμενες και παράλληλες. Και ο Ευκλείδης θα μας δώσει ένα μοναδικό αξίωμα για το πώς φέρνουμε παράλληλη από ένα σημείο.
\end{anagke}

\begin{lexi}{παράλληλες, τεμνόμενες}{παράλληλες, τεμνόμενες}
\lexiword{Παράλληλες} είναι οι ευθείες που πηγαίνουν «δίπλα-δίπλα» (παρ' άλληλες), «η μία πλάι στην άλλη», χωρίς ποτέ να συναντιούνται. \lexiword{Τεμνόμενες} είναι οι ευθείες που «τέμνονται», δηλαδή σταυρώνουν.
\end{lexi}

\begin{metafrasi}
Δύο ευθείες χωρίς κοινό σημείο. & παράλληλες: $\varepsilon_1\parallel\varepsilon_2$\\
Δύο ευθείες με ένα κοινό σημείο. & τεμνόμενες\\
Το κοινό τους σημείο. & σημείο τομής\\
Τμήματα πάνω σε παράλληλες ευθείες. & $\mathrm{ΑΒ}\parallel\Gamma\Delta$\\
\end{metafrasi}

\tmima{Η θεωρία}

\begin{orismos}{παράλληλες και τεμνόμενες ευθείες}
Δύο ευθείες του ίδιου επιπέδου λέγονται \textbf{παράλληλες} αν δεν έχουν κανένα κοινό σημείο, όσο κι αν προεκταθούν. Γράφουμε $\varepsilon_1\parallel\varepsilon_2$. Λέγονται \textbf{τεμνόμενες} αν έχουν ένα (και μόνο ένα) κοινό σημείο, που λέγεται \textbf{σημείο τομής}.
\end{orismos}

Δύο ευθείες του ίδιου επιπέδου ή είναι παράλληλες ή τέμνονται (αν έχουν δύο κοινά σημεία ταυτίζονται, αφού από δύο σημεία περνά μία ευθεία). Δύο τμήματα λέγονται \textbf{παράλληλα} όταν βρίσκονται σε παράλληλες ευθείες.

\begin{center}
\begin{tikzpicture}[scale=0.9]
\draw[thick,{Stealth}-{Stealth}] (0,0)--(4,0) node[right] {$\varepsilon_1$}; \draw[thick,{Stealth}-{Stealth}] (0,1.3)--(4,1.3) node[right] {$\varepsilon_2$};
\node at (2,-0.6) {$\varepsilon_1\parallel\varepsilon_2$};
\begin{scope}[xshift=7cm]
\draw[thick,{Stealth}-{Stealth}] (0,0)--(4,1.6); \draw[thick,{Stealth}-{Stealth}] (0,1.5)--(4,-0.1); \fill (2,0.77) circle (2pt);
\node at (2,-0.6) {τεμνόμενες};
\end{scope}
\end{tikzpicture}
\end{center}

\begin{idiotita}{αξίωμα παραλλήλων (Ευκλείδη)}
Από ένα σημείο $\mathrm{Α}$ εκτός ευθείας $(\varepsilon)$ διέρχεται \textbf{μία και μοναδική} ευθεία παράλληλη προς την $(\varepsilon)$.
\end{idiotita}

Το αξίωμα αυτό ανήκει στο έργο του Ευκλείδη, τα «Στοιχεία», που γράφτηκε στην Αλεξάνδρεια περίπου το $300$ π.Χ. Για πολλούς αιώνες οι μαθηματικοί προσπαθούσαν να το αποδείξουν από τα άλλα, και τελικά κατάλαβαν ότι είναι αναγκαίο: υπάρχουν και γεωμετρίες όπου δεν ισχύει.

Μπορούμε να φέρουμε παράλληλη με \textbf{γνώμονα και χάρακα}: ακουμπάμε τον χάρακα στη δοσμένη ευθεία $(\varepsilon)$ και τον γνώμονα πάνω στον χάρακα. Ολισθαίνουμε τον γνώμονα κατά μήκος του χάρακα μέχρι να περάσει από το σημείο $\mathrm{Α}$ και χαράσσουμε. Η ευθεία που προκύπτει είναι παράλληλη.

\begin{prosoxi}
\begin{itemize}[leftmargin=1.2em,itemsep=0.3ex,topsep=2pt]
\item Για να είναι παράλληλες δύο ευθείες δεν αρκεί να μην συναντιούνται στο σχέδιο: πρέπει να μην έχουν κοινό σημείο όσο κι αν προεκταθούν.
\item Δύο παράλληλα τμήματα δεν χρειάζεται να έχουν το ίδιο μήκος.
\item Στο επίπεδο δεν υπάρχουν δύο ευθείες που να μην είναι ούτε τεμνόμενες ούτε παράλληλες (στον χώρο υπάρχουν, οι ασύμβατες).
\item Αν μια ευθεία είναι παράλληλη σε δύο άλλες, αυτές είναι παράλληλες μεταξύ τους.
\end{itemize}
\end{prosoxi}

\tmima{Λυμένα παραδείγματα}

\paradeigma{Παράλληλες ή τεμνόμενες;}
\tixero Οι ευθείες: τα όρια του γηπέδου, οι ράγες, δύο μονοπάτια που διασταυρώνονται.\par
\tithelo Να πω ποιες είναι παράλληλες.\par
\possindeo Παράλληλες: οι ράγες και τα απέναντι όρια του γηπέδου. Τεμνόμενες: τα μονοπάτια.\par
\apantisi Ράγες και απέναντι πλευρές: παράλληλες. Μονοπάτια: τεμνόμενες.\par
\elegxos Οι παράλληλες δεν έχουν κοινό σημείο.

\paradeigma{Φέρνω παράλληλη}
\tixero Μια ευθεία $(\varepsilon)$ και ένα σημείο $\mathrm{Α}$ εκτός της.\par
\tithelo Να φέρω την παράλληλη.\par
\possindeo (1) Τοποθετώ τον χάρακα πάνω στην $(\varepsilon)$. (2) Τοποθετώ τον γνώμονα πάνω στον χάρακα, με μια πλευρά κατά μήκος του. (3) Σύρω τον γνώμονα μέχρι η άλλη του πλευρά να περάσει από το $\mathrm{Α}$. (4) Χαράσσω.\par
\apantisi Η ευθεία από το $\mathrm{Α}$ είναι η μοναδική παράλληλη προς την $(\varepsilon)$.\par
\elegxos Μετρώ δύο φορές την απόσταση των ευθειών: βγαίνει ίδια.

\paradeigma{Πόσα σημεία τομής;}
\tixero Τρεις διαφορετικές ευθείες του ίδιου επιπέδου.\par
\tithelo Πόσα το πολύ και πόσα το λιγότερο σημεία τομής μπορούν να έχουν.\par
\possindeo Αν είναι όλες παράλληλες, $0$. Αν δύο είναι παράλληλες και η τρίτη τέμνει και τις δύο: $2$. Αν περνούν και οι τρεις από το ίδιο σημείο: $1$. Αν ανά δύο τέμνονται σε διαφορετικά σημεία: $3$.\par
\apantisi $0$, $1$, $2$ ή $3$ σημεία.\par
\elegxos Δεν μπορεί να υπάρχει περισσότερα από $\frac{3\cdot2}{2}=3$ σημεία (ένα ανά ζευγάρι).

\paradeigma{Παράλληλα τμήματα}
\tixero Οι ευθείες $\mathrm{ΑΒ}$ και $\Gamma\Delta$ είναι παράλληλες, με $\mathrm{ΑΒ}=3$ και $\Gamma\Delta=7$.\par
\tithelo Τι λέμε για τα τμήματα.\par
\possindeo Δύο τμήματα σε παράλληλες ευθείες είναι παράλληλα.\par
\apantisi $\mathrm{ΑΒ}\parallel\Gamma\Delta$, παρότι έχουν διαφορετικά μήκη.\par
\elegxos Η παραλληλία εξαρτάται από τις ευθείες και όχι από τα μήκη.

\tmima{Ασκήσεις}

\zoni{Το κατάλαβα}{παράλληλες και τεμνόμενες}

\ex Ποιες από τις παρακάτω ευθείες είναι παράλληλες και ποιες τεμνόμενες: α) οι ράγες του τρένου \quad β) τα δύο μονοπάτια του κήπου που διασταυρώνονται \quad γ) οι αράδες του τετραδίου \quad δ) δύο ευθείες που έχουν ένα κοινό σημείο.

\ex Σχεδίασε μια ευθεία $(\varepsilon)$ και ένα σημείο $\mathrm{Α}$ εκτός της. Φέρε από το $\mathrm{Α}$ την παράλληλη με γνώμονα και χάρακα.

\ex Πόσες παράλληλες προς μια ευθεία μπορούν να φέρουμε από ένα σημείο εκτός της; Από ένα σημείο πάνω της;

\ex Πόσα σημεία τομής μπορεί να έχουν $3$ διαφορετικές ευθείες; Σχεδίασε κάθε περίπτωση.

\ex Πόσα το πολύ σημεία τομής έχουν $4$ ευθείες; Και $\nu$ ευθείες;

\ex Αν $\varepsilon_1\parallel\varepsilon_2$ και $\varepsilon_2\parallel\varepsilon_3$, τι ξέρεις για τις $\varepsilon_1$, $\varepsilon_3$;

\ex Σωστό ή λάθος; α) Δύο ευθείες του επιπέδου που δεν τέμνονται είναι παράλληλες. \quad β) Δύο παράλληλα τμήματα έχουν ίσα μήκη. \quad γ) Αν δύο ευθείες έχουν δύο κοινά σημεία, ταυτίζονται. \quad δ) Από σημείο εκτός ευθείας περνούν δύο παράλληλες προς αυτήν.

\zoni{Το χρησιμοποιώ}{κήπος, ράγες, γήπεδο}

\ex \textbf{Οι σειρές.} Οι σειρές των λαχανικών είναι παράλληλες ευθείες και απέχουν $60$ cm. Πώς θα τις σημάνεις στο έδαφος με σχοινί και παλούκια;

\ex \textbf{Το γήπεδο.} Πόσα ζεύγη παράλληλων πλευρών έχει ένα ορθογώνιο γήπεδο; Πόσα ζεύγη κάθετων πλευρών;

\ex \textbf{Τα μονοπάτια.} Πέντε μονοπάτια του κήπου είναι ευθείες και δεν είναι δύο παράλληλα. Πόσα σημεία τομής έχουν το πολύ;

\ex \textbf{Οι ράγες.} Γιατί οι ράγες του τρένου πρέπει να είναι παράλληλες;

\ex \textbf{Το μοτίβο.} Σχεδίασε τρεις παράλληλες ευθείες και μια ευθεία που τις τέμνει. Πόσα σημεία τομής έχεις;

\zoni{Το σκέφτομαι}{ο Ευκλείδης και η παραλληλία}

\ex Γιατί από δύο σημεία περνά μία ευθεία ενώ από σημείο εκτός ευθείας περνά μία παράλληλη; Σε τι διαφέρουν οι δύο προτάσεις;

\ex Αν δύο ευθείες δεν έχουν κοινό σημείο στο χαρτί σου, μπορεί να τέμνονται σε σημείο πέρα από το χαρτί; Εξήγησε.

\ex Δύο διαφορετικές ευθείες είναι κάθετες σε μια τρίτη. Τι σχέση έχουν μεταξύ τους;

\ex Μπορείς να φέρεις δύο διαφορετικές παράλληλες από ένα σημείο προς μια ευθεία; Τι θα συνέβαινε;

\zoni{Για το σπίτι}{ψάξε παράλληλες}

\ex Βρες στο σπίτι πέντε ζεύγη παράλληλων ευθειών και τρία ζεύγη τεμνόμενων.

\ex Σχεδίασε με γνώμονα και χάρακα τρεις παράλληλες ευθείες και μέτρησε την απόσταση ανάμεσά τους.

\paragrafos{Απόσταση σημείου από ευθεία και απόσταση παραλλήλων}{Β.1.10, «Απόσταση σημείου από ευθεία – Απόσταση παραλλήλων»}

\begin{anagke}
Στον κήπο υπάρχει ένα ευθύγραμμο μονοπάτι και ένα δέντρο που βρίσκεται λίγο πιο πέρα, εκτός μονοπατιού. Ο Νίκος θέλει να φτάσει από το δέντρο στο μονοπάτι όσο πιο γρήγορα γίνεται. Δοκιμάζει τρεις δρόμους με ένα σχοινί: έναν λοξό προς τα αριστερά, έναν λοξό προς τα δεξιά και έναν κατευθείαν «πάνω». Ο ευθύς είναι ο μικρότερος. «Ο δρόμος που πέφτει κάθετα στο μονοπάτι είναι ο συντομότερος!» ανακαλύπτει.

Η Αγγελική προσθέτει ένα δεύτερο σχέδιο: δύο παράλληλες σειρές με δέντρα. «Πόσο απέχουν οι δύο σειρές;» Μετράει από ένα δέντρο της μιας στην άλλη, κάθετα, και βρίσκει $2{,}4$ μέτρα. Μετράει από ένα άλλο δέντρο και βρίσκει πάλι $2{,}4$. «Σε όλα τα σημεία η απόσταση είναι ίδια!». Η Πολυξένη το συμπληρώνει: «Γι' αυτό οι ράγες του τρένου έχουν σταθερή απόσταση: αλλιώς το τρένο δεν χωράει».

Σε αυτή την παράγραφο ορίζουμε την απόσταση ενός σημείου από μια ευθεία και την απόσταση δύο παραλλήλων, δύο έννοιες που θα τις χρειαζόμαστε συνεχώς.
\end{anagke}

\begin{lexi}{απόσταση, κάθετο, πλάγιο}{απόσταση, κάθετο, πλάγιο}
Η \lexiword{απόσταση} σημείου από ευθεία είναι ο «συντομότερος δρόμος» από το σημείο προς την ευθεία. Το \lexiword{κάθετο} τμήμα πέφτει με ορθή γωνία πάνω στην ευθεία, ενώ τα υπόλοιπα, τα \lexiword{πλάγια}, είναι «λοξά» και μακρύτερα.
\end{lexi}

\begin{metafrasi}
Το κάθετο τμήμα από το $\mathrm{Α}$ στην $(\varepsilon)$. & $\mathrm{ΑΜ}\perp(\varepsilon)$\\
Η απόσταση του $\mathrm{Α}$ από την $(\varepsilon)$. & $\mathrm{ΑΜ}$\\
Η απόσταση δύο παραλλήλων. & η απόσταση οποιουδήποτε σημείου της μιας από την άλλη\\
Κάθε πλάγιο τμήμα. & μεγαλύτερο από το κάθετο\\
\end{metafrasi}

\tmima{Η θεωρία}

\begin{orismos}{απόσταση σημείου από ευθεία}
\textbf{Απόσταση} ενός σημείου $\mathrm{Α}$ από μια ευθεία $(\varepsilon)$ είναι το μήκος του \textbf{κάθετου} τμήματος $\mathrm{ΑΜ}$ από το $\mathrm{Α}$ προς την $(\varepsilon)$, όπου $\mathrm{Μ}$ το ίχνος της καθέτου.
\end{orismos}

Το κάθετο τμήμα είναι το \textbf{συντομότερο} από όλα τα τμήματα που ενώνουν το $\mathrm{Α}$ με σημεία της $(\varepsilon)$. Κάθε \textbf{πλάγιο} τμήμα είναι μεγαλύτερο από το κάθετο. Αν το σημείο ανήκει στην ευθεία, η απόστασή του από αυτήν είναι $0$.

\begin{center}
\begin{tikzpicture}[scale=1.1]
\draw[thick,{Stealth}-{Stealth}] (-2.5,0)--(3.5,0) node[right] {$(\varepsilon)$};
\fill (0.5,2) circle (2pt) node[above] {$\mathrm{Α}$}; \fill (0.5,0) circle (2pt) node[below] {$\mathrm{Μ}$};
\draw[gblue,thick] (0.5,2)--(0.5,0); \draw (0.5,0.25)--(0.75,0.25)--(0.75,0);
\draw (0.5,2)--(-1.5,0); \draw (0.5,2)--(2.5,0);
\node[gblue,right] at (0.5,1) {$\mathrm{ΑΜ}$};
\end{tikzpicture}
\end{center}

\tmimasmall{Απόσταση παραλλήλων}

Αν $\varepsilon_1\parallel\varepsilon_2$, όλα τα σημεία της $\varepsilon_1$ απέχουν την ίδια απόσταση από την $\varepsilon_2$. Η απόσταση αυτή λέγεται \textbf{απόσταση των παραλλήλων}. Δύο παράλληλες ευθείες έχουν σταθερή απόσταση, δηλαδή δεν πλησιάζουν ούτε απομακρύνονται.

Το σύνολο των σημείων που απέχουν απόσταση $\delta$ από μια ευθεία $(\varepsilon)$ είναι δύο παράλληλες ευθείες προς την $(\varepsilon)$, μία σε κάθε ημιεπίπεδο, που απέχουν $\delta$ από αυτήν.

\begin{prosoxi}
\begin{itemize}[leftmargin=1.2em,itemsep=0.3ex,topsep=2pt]
\item Η απόσταση σημείου από ευθεία μετριέται πάντα κάθετα, όχι λοξά.
\item Το πλάγιο τμήμα δεν δείχνει την απόσταση: είναι μεγαλύτερο.
\item Η απόσταση παραλλήλων είναι η ίδια όπου κι αν τη μετρήσεις, αλλά πρέπει να τη μετράς κάθετα.
\item Η απόσταση ενός σημείου της ευθείας από την ευθεία είναι $0$.
\end{itemize}
\end{prosoxi}

\tmima{Λυμένα παραδείγματα}

\paradeigma{Ο συντομότερος δρόμος}
\tixero Ένα δέντρο $\mathrm{Α}$ εκτός ευθύγραμμου μονοπατιού $(\varepsilon)$.\par
\tithelo Ποιος δρόμος από το δέντρο προς το μονοπάτι είναι ο συντομότερος.\par
\possindeo Η απόσταση του $\mathrm{Α}$ από την $(\varepsilon)$ είναι το κάθετο τμήμα $\mathrm{ΑΜ}$. Κάθε πλάγιο $\mathrm{ΑΒ}$ είναι μεγαλύτερο.\par
\apantisi Ο κάθετος δρόμος $\mathrm{ΑΜ}$.\par
\elegxos Με σχοινί: ο λοξός δρόμος βγαίνει μακρύτερος.

\paradeigma{Η απόσταση παραλλήλων}
\tixero Δύο σειρές δέντρων παράλληλες και $\mathrm{ΑΜ}\perp\varepsilon_2$ με $\mathrm{ΑΜ}=2{,}4$ μ.\par
\tithelo Η απόσταση ενός άλλου δέντρου $\mathrm{Β}$ της $\varepsilon_1$ από την $\varepsilon_2$.\par
\possindeo Οι παράλληλες έχουν σταθερή απόσταση.\par
\apantisi $2{,}4$ μ.\par
\elegxos Κάθε κάθετο τμήμα ανάμεσα στις παράλληλες είναι $2{,}4$ μ.

\paradeigma{Τρεις σειρές}
\tixero Τρεις παράλληλες σειρές δέντρων σε ίσες αποστάσεις $2{,}4$ μ.\par
\tithelo Η απόσταση της πρώτης από την τρίτη.\par
\possindeo Η απόσταση από την πρώτη στη δεύτερη είναι $2{,}4$ και από τη δεύτερη στην τρίτη $2{,}4$, αφού οι παράλληλες είναι διαδοχικές. $2{,}4+2{,}4=4{,}8$.\par
\apantisi $4{,}8$ μ.\par
\elegxos Η μεσαία είναι στη μέση των δύο άλλων.

\paradeigma{Σημεία με δοσμένη απόσταση}
\tixero Μια ευθεία $(\varepsilon)$.\par
\tithelo Ποια σημεία απέχουν $3$ cm από αυτήν.\par
\possindeo Σε κάθε ημιεπίπεδο υπάρχει μία παράλληλη προς την $(\varepsilon)$ σε απόσταση $3$ cm.\par
\apantisi Δύο παράλληλες ευθείες, μία σε κάθε ημιεπίπεδο.\par
\elegxos Η απόσταση ανάμεσά τους είναι $6$ cm.

\tmima{Ασκήσεις}

\zoni{Το κατάλαβα}{κάθετα και πλάγια}

\ex Σχεδίασε μια ευθεία $(\varepsilon)$ και ένα σημείο $\mathrm{Α}$ εκτός της. Φέρε την κάθετη από το $\mathrm{Α}$ και δύο πλάγια. Μέτρησε και σύγκρινε τα μήκη τους.

\ex Μέτρησε την απόσταση του $\mathrm{Α}$ από την $(\varepsilon)$. Πόση είναι η απόσταση ενός σημείου που βρίσκεται πάνω στην $(\varepsilon)$;

\ex Σχεδίασε δύο παράλληλες ευθείες και μέτρησε την απόστασή τους σε τρία διαφορετικά σημεία. Τι παρατηρείς;

\ex Τρεις παράλληλες ευθείες απέχουν διαδοχικά $2{,}4$ cm. Πόση είναι η απόσταση της πρώτης από την τρίτη;

\ex Σχεδίασε όλα τα σημεία που απέχουν $2$ cm από μια ευθεία $(\varepsilon)$.

\ex Σωστό ή λάθος; α) Το κάθετο τμήμα είναι το μικρότερο από όλα όσα ενώνουν ένα σημείο με την ευθεία. \quad β) Η απόσταση παραλλήλων δεν είναι σταθερή. \quad γ) Η απόσταση σημείου της ευθείας από την ευθεία είναι $0$. \quad δ) Το πλάγιο είναι μικρότερο από το κάθετο.

\zoni{Το χρησιμοποιώ}{κήπος, δρόμοι, σχοινιά}

\ex \textbf{Το κοντινότερο.} Ένα δέντρο απέχει $3$ μ. κάθετα από ένα μονοπάτι. Ο Νίκος δοκιμάζει να φτάσει λοξά και βρίσκει $3{,}6$ μ. Μπορεί να έχει υπολογίσει τον κάθετο δρόμο λάθος; Εξήγησε.

\ex \textbf{Οι σειρές.} Δύο παράλληλες σειρές λαχανικών απέχουν $60$ cm. Μια τρίτη σειρά, παράλληλη προς αυτές και προς την άλλη πλευρά της δεύτερης, απέχει $60$ cm από τη δεύτερη. Πόσο απέχει από την πρώτη;

\ex \textbf{Το πλάγιο.} Ένα σημείο $\mathrm{Α}$ απέχει $4$ cm από μια ευθεία $(\varepsilon)$. Ένα σημείο $\mathrm{Β}$ της $(\varepsilon)$ μπορεί να απέχει από το $\mathrm{Α}$ $3$ cm; $4$ cm; $6$ cm;

\ex \textbf{Το ράφι.} Ένα ράφι είναι παράλληλο προς το πάτωμα και απέχει $1{,}2$ μ. Πόσο απέχει κάθε σημείο του από το πάτωμα;

\ex \textbf{Το χωράφι.} Ένα ορθογώνιο χωράφι έχει διαστάσεις $40$ και $25$ μ. Πόση είναι η απόσταση ανάμεσα στις μεγάλες απέναντι πλευρές του; Και ανάμεσα στις μικρές;

\zoni{Το σκέφτομαι}{ο συντομότερος δρόμος}

\ex Εξήγησε γιατί το κάθετο τμήμα είναι το συντομότερο (σκέψου ένα τεντωμένο νήμα και μια ράβδο που στρέφεται).

\ex Σε ποια περίπτωση η απόσταση δύο παραλλήλων είναι $0$;

\ex Βρες τα σημεία που απέχουν $3$ cm από την $(\varepsilon)$ και $5$ cm από μια παράλληλη προς αυτήν σε απόσταση $8$ cm.

\ex Δύο παράλληλες ευθείες απέχουν $6$ cm. Πόσες ευθείες παράλληλες προς αυτές απέχουν $3$ cm και από τις δύο;

\zoni{Για το σπίτι}{μέτρησε κάθετα}

\ex Μέτρησε την απόσταση ανάμεσα σε δύο παράλληλες ακμές του τραπεζιού σου σε δύο διαφορετικά σημεία.

\ex Βρες τη μικρότερη απόσταση ενός σημείου του δωματίου από έναν τοίχο και εξήγησε τι μέτρησες.

\paragrafos{Ο κύκλος και τα στοιχεία του}{Β.1.11, «Κύκλος και στοιχεία του κύκλου»}

\begin{anagke}
Το πιο όμορφο σημείο του κήπου θα είναι ένα στρογγυλό παρτέρι με ένα συντριβάνι στο κέντρο. Ο παππούς Γιάννης δείχνει πώς θα το χαράξουν στο χώμα: καρφώνει ένα παλούκι στο κέντρο, δένει ένα σχοινί και, κρατώντας το τεντωμένο, γυρίζει και χαράσσει τη γραμμή. «Όλα τα σημεία της γραμμής έχουν την ίδια απόσταση από το παλούκι, όσο είναι το σχοινί», λέει.

Η Αγγελική το αναγνωρίζει: «Αυτό είναι ο κύκλος! Και το σχοινί είναι η ακτίνα!». Ο Νίκος ρωτάει πώς λέγεται το πιο μακρύ τμήμα που μπορούν να τεντώσουν από ένα άκρο του κύκλου ως το άλλο: «Η διάμετρος», απαντά η Πολυξένη, «περνάει από το κέντρο και είναι διπλάσια από την ακτίνα». Ο Βασίλης στέκεται κάπου και ρωτάει αν είναι «μέσα» ή «έξω» από τον κύκλο: «Αν η απόστασή σου από το παλούκι είναι μικρότερη από το σχοινί, είσαι μέσα».

Ο κύκλος είναι το πιο συμμετρικό σχήμα και όλα τα στοιχεία του προκύπτουν από τη μέτρηση αποστάσεων από το κέντρο.
\end{anagke}

\begin{lexi}{κύκλος, ακτίνα, διάμετρος, χορδή}{κύκλος, ακτίνα, διάμετρος, χορδή}
Ο \lexiword{κύκλος} είναι ό,τι «κυλάει», γυρίζει. Η \lexiword{ακτίνα} είναι η «αχτίδα», η ακτίνα του φωτός που ξεκινά από το κέντρο. Η \lexiword{διάμετρος} είναι το «δια-μέτρο», αυτό που μετράει «δια μέσου» του κέντρου. Η \lexiword{χορδή} είναι η «χορδή» του μουσικού οργάνου: ένα τεντωμένο νήμα ανάμεσα σε δύο σημεία. Το \lexiword{ημικύκλιο} είναι το μισό του κύκλου.
\end{lexi}

\begin{metafrasi}
Ο κύκλος με κέντρο $\mathrm{Ο}$ και ακτίνα $\rho$. & $(\mathrm{Ο},\rho)$\\
Τα σημεία πάνω στον κύκλο. & $\mathrm{ΟΜ}=\rho$\\
Τα σημεία μέσα στον κύκλο. & $\mathrm{ΟΜ}<\rho$\\
Τα σημεία έξω από τον κύκλο. & $\mathrm{ΟΜ}>\rho$\\
Η διάμετρος. & $\delta=2\rho$\\
\end{metafrasi}

\tmima{Η θεωρία}

\begin{orismos}{κύκλος}
\textbf{Κύκλος} με κέντρο $\mathrm{Ο}$ και ακτίνα $\rho$ είναι το σύνολο όλων των σημείων του επιπέδου που απέχουν από το $\mathrm{Ο}$ απόσταση ίση με $\rho$. Συμβολίζεται $(\mathrm{Ο},\rho)$. Η απόσταση $\rho$ λέγεται \textbf{ακτίνα} και το $\mathrm{Ο}$ \textbf{κέντρο}.
\end{orismos}

Για να σχεδιάσουμε κύκλο χρησιμοποιούμε \textbf{διαβήτη}. Δύο κύκλοι με ίσες ακτίνες είναι ίσοι.

\begin{center}
\begin{tikzpicture}[scale=1.15]
\draw[thick] (0,0) circle (1.6);
\fill (0,0) circle (2pt) node[below left] {$\mathrm{Ο}$};
\draw[gblue,thick] (0,0)--(1.13,1.13) node[midway,above left,gblue] {$\rho$}; \fill (1.13,1.13) circle (2pt) node[above right] {$\mathrm{Α}$};
\draw (-1.6,0)--(1.6,0); \fill (-1.6,0) circle (2pt) node[left] {$\mathrm{Β}$}; \fill (1.6,0) circle (2pt) node[right] {$\mathrm{Γ}$};
\node[below] at (0,-0.1) {\small διάμετρος $\mathrm{ΒΓ}$};
\draw[dashed] (-1.2,1.06)--(0.8,-1.38); \fill (-1.2,1.06) circle (2pt); \fill (0.8,-1.38) circle (2pt);
\node at (-1.2,-0.9) {\small χορδή};
\end{tikzpicture}
\end{center}

\begin{itemize}[leftmargin=1.2em,itemsep=0.3ex,topsep=2pt]
\item \textbf{Χορδή} είναι το ευθύγραμμο τμήμα που ενώνει δύο σημεία του κύκλου.
\item \textbf{Διάμετρος} είναι η χορδή που περνά από το κέντρο. Είναι η μεγαλύτερη χορδή, ίση με $\delta=2\rho$, και χωρίζει τον κύκλο σε δύο ίσα μέρη, τα \textbf{ημικύκλια}.
\item Ένα σημείο $\mathrm{Μ}$ είναι \textbf{πάνω} στον κύκλο αν $\mathrm{ΟΜ}=\rho$, \textbf{μέσα} σε αυτόν αν $\mathrm{ΟΜ}<\rho$ και \textbf{έξω} αν $\mathrm{ΟΜ}>\rho$.
\item Ο κύκλος μαζί με το εσωτερικό του λέγεται \textbf{κυκλικός δίσκος}.
\end{itemize}

\begin{prosoxi}
\begin{itemize}[leftmargin=1.2em,itemsep=0.3ex,topsep=2pt]
\item Ο κύκλος είναι μόνο η γραμμή, όχι το εσωτερικό: το εσωτερικό μαζί με τη γραμμή είναι ο κυκλικός δίσκος.
\item Η ακτίνα είναι τμήμα (ή μήκος), όχι ολόκληρη ευθεία.
\item Η διάμετρος είναι πάντα $2\rho$, και κάθε διάμετρος είναι χορδή, αλλά όχι κάθε χορδή είναι διάμετρος.
\item Σε έναν κύκλο όλες οι ακτίνες είναι ίσες, όλες οι διάμετροι είναι ίσες.
\end{itemize}
\end{prosoxi}

\tmima{Λυμένα παραδείγματα}

\paradeigma{Το κυκλικό παρτέρι}
\tixero Το σχοινί του παππού είναι $2{,}5$ μ.\par
\tithelo Η ακτίνα και η διάμετρος του παρτεριού.\par
\possindeo Το σχοινί είναι η ακτίνα: $\rho=2{,}5$ μ. Η διάμετρος $\delta=2\cdot2{,}5=5$ μ.\par
\apantisi $\rho=2{,}5$ μ. και $\delta=5$ μ.\par
\elegxos Το παρτέρι έχει πλάτος $5$ μ.

\paradeigma{Μέσα, πάνω ή έξω;}
\tixero Κύκλος $(\mathrm{Ο},4)$ και σημεία με $\mathrm{ΟΑ}=3$, $\mathrm{ΟΒ}=4$, $\mathrm{ΟΓ}=5$.\par
\tithelo Η θέση κάθε σημείου.\par
\possindeo Συγκρίνω την απόσταση με την ακτίνα $4$. $\mathrm{ΟΑ}<4$: μέσα. $\mathrm{ΟΒ}=4$: πάνω. $\mathrm{ΟΓ}>4$: έξω.\par
\apantisi Το $\mathrm{Α}$ μέσα, το $\mathrm{Β}$ πάνω, το $\mathrm{Γ}$ έξω.\par
\elegxos Η απόσταση από το κέντρο καθορίζει τη θέση.

\paradeigma{Μία χορδή}
\tixero Ο κύκλος $(\mathrm{Ο},5)$ και μια χορδή $\mathrm{ΑΒ}$.\par
\tithelo Πόσο μεγάλη μπορεί να είναι η χορδή.\par
\possindeo Η μεγαλύτερη χορδή είναι η διάμετρος, $2\cdot5=10$.\par
\apantisi $\mathrm{ΑΒ}\le10$, ίσο με $10$ μόνο όταν περνά από το κέντρο.\par
\elegxos Αν $\mathrm{ΑΒ}=10$ τότε το κέντρο είναι το μέσο του $\mathrm{ΑΒ}$.

\paradeigma{Ίσοι κύκλοι}
\tixero Δύο κύκλοι: $(\mathrm{Ο},3)$ και $(\mathrm{Κ},3)$.\par
\tithelo Είναι ίσοι;\par
\possindeo Έχουν ίσες ακτίνες, άρα μπορούν να εφαρμοστούν ο ένας πάνω στον άλλο μετακινώντας το κέντρο.\par
\apantisi Ναι.\par
\elegxos Ίσες ακτίνες δίνουν ίσους κύκλους.

\tmima{Ασκήσεις}

\zoni{Το κατάλαβα}{ακτίνα, διάμετρος, χορδή}

\ex Σχεδίασε με διαβήτη κύκλο $(\mathrm{Ο},3\ \mathrm{cm})$. Φέρε μια ακτίνα, μια διάμετρο και μια χορδή που δεν περνά από το κέντρο.

\ex Συμπλήρωσε: α) $\rho=4$, $\delta=\ldots$ \quad β) $\delta=9$, $\rho=\ldots$ \quad γ) $\rho=2{,}5$, $\delta=\ldots$ \quad δ) $\delta=7{,}2$, $\rho=\ldots$.

\ex Κύκλος $(\mathrm{Ο},4)$. Πες αν είναι μέσα, πάνω ή έξω τα σημεία με $\mathrm{ΟΑ}=3$, $\mathrm{ΟΒ}=4$, $\mathrm{ΟΓ}=5$, $\mathrm{ΟΔ}=0$.

\ex Στον κύκλο $(\mathrm{Ο},5)$ πόσο μεγάλη είναι η μεγαλύτερη χορδή; Πόσες διάμετροι υπάρχουν;

\ex Σχεδίασε δύο κύκλους με την ίδια ακτίνα και δείξε ότι είναι ίσοι με διαφανές χαρτί.

\ex Σωστό ή λάθος; α) Η διάμετρος είναι χορδή. \quad β) Η μεγαλύτερη χορδή είναι η ακτίνα. \quad γ) $\delta=2\rho$. \quad δ) Ένα σημείο με $\mathrm{ΟΜ}<\rho$ είναι πάνω στον κύκλο. \quad ε) Όλες οι ακτίνες ενός κύκλου είναι ίσες.

\zoni{Το χρησιμοποιώ}{παρτέρια, ρολόγια, ρόδες}

\ex \textbf{Το παρτέρι.} Το κυκλικό παρτέρι χαράσσεται με σχοινί $2{,}5$ μ. Πόσο είναι το πλάτος του;

\ex \textbf{Το ποδήλατο.} Ο τροχός ενός ποδηλάτου έχει διάμετρο $66$ cm. Πόση είναι η ακτίνα του; Σε ύψος πόσο βρίσκεται ο άξονάς του από το έδαφος;

\ex \textbf{Το ρολόι.} Ο λεπτοδείκτης έχει μήκος $8$ cm. Σε ποιο κύκλο κινείται η άκρη του και πόση διάμετρο έχει;

\ex \textbf{Η σέντρα.} Ο κύκλος της σέντρας σε ένα γήπεδο ποδοσφαίρου έχει ακτίνα $9{,}15$ μ. Πόση είναι η διάμετρός του;

\ex \textbf{Η στρογγυλή πίτα.} Μια πίτα έχει διάμετρο $26$ cm. Θα χωρέσει σε ταψί διαστάσεων $25\times25$ cm; Εξήγησε.

\zoni{Το σκέφτομαι}{όλα ίσα από το κέντρο}

\ex Γιατί η διάμετρος είναι η μεγαλύτερη χορδή; (Σκέψου το τμήμα ανάμεσα σε δύο σημεία και την απόστασή τους από το κέντρο.)

\ex Αν σε έναν κύκλο οι ακτίνες είναι $5$ cm, πόσα σημεία του κύκλου απέχουν $5$ cm από ένα δοσμένο σημείο του κύκλου; Πρόβλεψε και σχεδίασε.

\ex Πόσοι κύκλοι μπορούν να περάσουν από δύο δοσμένα σημεία $\mathrm{Α}$ και $\mathrm{Β}$; Πού βρίσκονται τα κέντρα τους;

\ex Μπορούν δύο διαφορετικοί κύκλοι με το ίδιο κέντρο να έχουν κοινό σημείο; Εξήγησε.

\zoni{Για το σπίτι}{στρογγυλά αντικείμενα}

\ex Βρες πέντε κυκλικά αντικείμενα στο σπίτι και μέτρησε διάμετρο και ακτίνα.

\ex Σχεδίασε έναν κύκλο με νήμα και καρφίτσα στο χαρτόνι.

\paragrafos{Επίκεντρη γωνία, τόξο και μέτρηση τόξου}{Β.1.12, «Επίκεντρη γωνία – Σχέση επίκεντρης γωνίας και του αντίστοιχου τόξου – Μέτρηση τόξου»}

\begin{anagke}
Στο πάρκο δίπλα στο σχολείο υπάρχει μια ρόδα λούνα παρκ με $8$ ίσες καμπίνες. Η παρέα την κοιτάζει από ψηλά: οι καμπίνες είναι στερεωμένες σε ίσες αποστάσεις πάνω σε ένα κύκλο. «Πόσες μοίρες γυρίζει ο τροχός για να πάει η μία καμπίνα στη θέση της επόμενης;» ρωτάει ο Ηλίας. Η Πολυξένη σκέφτεται: «Ο πλήρης κύκλος είναι $360^\circ$ και υπάρχουν $8$ ίσα διαστήματα: $360:8=45^\circ$».

Η Αγγελική σχεδιάζει δύο ακτίνες προς δύο διαδοχικές καμπίνες και παρατηρεί τη γωνία ανάμεσά τους με κορυφή το κέντρο: «Είναι $45^\circ$, και το κομμάτι του κύκλου ανάμεσα στις δύο καμπίνες το μετράμε και αυτό $45^\circ$». Ο Στέλιος ρωτάει: «Το μήκος του κύκλου και οι μοίρες είναι το ίδιο;». Η κυρία Αλεξάνδρα απαντάει: «Όχι. Άλλο το μήκος σε μέτρα και άλλο το άνοιγμα σε μοίρες. Σήμερα μιλάμε για το άνοιγμα, το λεγόμενο \emph{μέτρο του τόξου}».

Στην παράγραφο αυτή ορίζουμε την επίκεντρη γωνία και το αντίστοιχο τόξο, και μετράμε τα τόξα με μοίρες.
\end{anagke}

\begin{lexi}{επίκεντρη γωνία, τόξο}{επίκεντρη γωνία, τόξο}
Η \lexiword{επίκεντρη} γωνία είναι η γωνία που έχει την κορυφή της «επί του κέντρου» του κύκλου. Το \lexiword{τόξο} είναι το κομμάτι του κύκλου, σαν το τόξο του τοξοβόλου, καμπυλωμένο. Το \lexiword{ημικύκλιο} είναι το μισό του κύκλου.
\end{lexi}

\begin{metafrasi}
Γωνία με κορυφή το κέντρο. & επίκεντρη γωνία $\widehat{\mathrm{ΑΟΒ}}$\\
Το κομμάτι του κύκλου μέσα στη γωνία. & αντίστοιχο τόξο $\overset{\frown}{\mathrm{ΑΒ}}$\\
Το μέτρο του τόξου. & ίσο με το μέτρο της επίκεντρης γωνίας\\
Ο πλήρης κύκλος. & $360^\circ$\\
Το ημικύκλιο. & $180^\circ$\\
\end{metafrasi}

\tmima{Η θεωρία}

\begin{orismos}{επίκεντρη γωνία και τόξο}
\textbf{Επίκεντρη} γωνία ενός κύκλου $(\mathrm{Ο},\rho)$ είναι κάθε γωνία $x\mathrm{Ο}y$ με κορυφή το κέντρο $\mathrm{Ο}$. Αν οι πλευρές της τέμνουν τον κύκλο στα $\mathrm{Α}$ και $\mathrm{Β}$, το τμήμα του κύκλου που βρίσκεται στο εσωτερικό της γωνίας λέγεται \textbf{αντίστοιχο τόξο} της και γράφεται $\overset{\frown}{\mathrm{ΑΒ}}$.
\end{orismos}

\begin{center}
\begin{tikzpicture}[scale=1.2]
\draw (0,0) circle (1.6);
\fill (0,0) circle (2pt) node[below left] {$\mathrm{Ο}$};
\draw[thick] (0,0)--(1.6,0) (0,0)--({1.6*cos(50)},{1.6*sin(50)});
\fill (1.6,0) circle (2pt) node[right] {$\mathrm{Α}$}; \fill ({1.6*cos(50)},{1.6*sin(50)}) circle (2pt) node[above right] {$\mathrm{Β}$};
\draw[gblue,very thick] (1.6,0) arc(0:50:1.6);
\draw[gblue] (0.5,0) arc(0:50:0.5); \node[gblue] at (0.8,0.2) {\small$\omega$};
\end{tikzpicture}
\end{center}

\begin{kanonas}{μέτρο τόξου}
Το \textbf{μέτρο ενός τόξου} σε μοίρες είναι ίσο με το μέτρο της αντίστοιχης επίκεντρης γωνίας. Ολόκληρος ο κύκλος είναι τόξο $360^\circ$ και το ημικύκλιο $180^\circ$. Ίσες επίκεντρες γωνίες αντιστοιχούν σε ίσα τόξα, και αντίστροφα.
\end{kanonas}

Έτσι το τόξο $\overset{\frown}{\mathrm{ΑΒ}}$ του σχήματος είναι $50^\circ$. Όταν η γωνία είναι μη κυρτή, το τόξο που αντιστοιχεί σε αυτήν είναι το μεγάλο: $360^\circ-50^\circ=310^\circ$.

Το μέτρο τόξου (σε μοίρες) είναι διαφορετικό από το \emph{μήκος} του τόξου (σε μέτρα): δύο ομόκεντροι κύκλοι έχουν τόξα με ίδιο μέτρο αλλά διαφορετικό μήκος.

\begin{prosoxi}
\begin{itemize}[leftmargin=1.2em,itemsep=0.3ex,topsep=2pt]
\item Το μέτρο του τόξου σε μοίρες δεν είναι το μήκος του σε μέτρα.
\item Κάθε χορδή $\mathrm{ΑΒ}$ ορίζει δύο τόξα (ένα μικρό και ένα μεγάλο), εκτός αν είναι διάμετρος (δύο ημικύκλια).
\item Η επίκεντρη γωνία έχει την κορυφή της στο κέντρο: αν η κορυφή είναι αλλού, δεν είναι επίκεντρη.
\item Σε διαφορετικούς κύκλους, ίσα τόξα (σε μοίρες) δεν σημαίνει ίσα μήκη.
\end{itemize}
\end{prosoxi}

\tmima{Λυμένα παραδείγματα}

\paradeigma{Η ρόδα με τις $8$ καμπίνες}
\tixero Οκτώ ίσες καμπίνες πάνω σε κύκλο.\par
\tithelo Το μέτρο του τόξου ανάμεσα σε δύο διαδοχικές.\par
\possindeo Ο κύκλος είναι $360^\circ$ και τα $8$ τόξα είναι ίσα: $360^\circ:8=45^\circ$.\par
\apantisi $45^\circ$.\par
\elegxos $8\cdot45=360$.

\paradeigma{Τόξο και επίκεντρη}
\tixero Η επίκεντρη γωνία $\widehat{\mathrm{ΑΟΒ}}=50^\circ$.\par
\tithelo Το μικρό και το μεγάλο τόξο $\overset{\frown}{\mathrm{ΑΒ}}$.\par
\possindeo Το μικρό τόξο είναι $50^\circ$. Το μεγάλο είναι $360^\circ-50^\circ=310^\circ$.\par
\apantisi $50^\circ$ και $310^\circ$.\par
\elegxos $50+310=360$.

\paradeigma{Διάμετρος και τόξα}
\tixero Η $\mathrm{ΑΒ}$ είναι διάμετρος του κύκλου.\par
\tithelo Το μέτρο των τόξων που ορίζει.\par
\possindeo Η διάμετρος χωρίζει τον κύκλο σε δύο ημικύκλια. Κάθε ημικύκλιο είναι $180^\circ$.\par
\apantisi $180^\circ$ το καθένα.\par
\elegxos Η επίκεντρη γωνία $\widehat{\mathrm{ΑΟΒ}}$ είναι ευθεία.

\paradeigma{Ένα ρολόι}
\tixero Ο λεπτοδείκτης κινείται από το $12$ στο $4$.\par
\tithelo Το τόξο που διέγραψε η άκρη του.\par
\possindeo Κάθε λεπτό ($1/60$ του κύκλου) είναι $6^\circ$, ή κάθε πεντάλεπτο $30^\circ$. Από το $12$ στο $4$ είναι $4$ πεντάλεπτα, $4\cdot30^\circ=120^\circ$.\par
\apantisi $120^\circ$.\par
\elegxos Στα $20$ λεπτά, $20\cdot6^\circ=120^\circ$.

\tmima{Ασκήσεις}

\zoni{Το κατάλαβα}{τόξα και μοίρες}

\ex Σχεδίασε κύκλο και επίκεντρη γωνία $70^\circ$. Ποιο είναι το μέτρο του μικρού και του μεγάλου τόξου που αντιστοιχούν;

\ex Ένας κύκλος χωρίζεται σε $5$ ίσα τόξα. Πόσες μοίρες είναι το καθένα; Και σε $9$; Σε $12$;

\ex Πόσες μοίρες είναι το $\frac{1}{4}$, το $\frac{1}{8}$, το $\frac{3}{4}$ και το $\frac{5}{6}$ του κύκλου;

\ex Ένα τόξο έχει μέτρο $60^\circ$. Τι μέρος του κύκλου είναι;

\ex Μια διάμετρος $\mathrm{ΑΒ}$ ορίζει δύο τόξα. Πόσες μοίρες είναι το καθένα;

\ex Σχεδίασε κύκλο και τρεις επίκεντρες διαδοχικές γωνίες $100^\circ$, $120^\circ$, $140^\circ$. Πόσο είναι τα τόξα και πόσο το άθροισμά τους;

\ex Σωστό ή λάθος; α) Το μέτρο του τόξου είναι ίσο με της επίκεντρης γωνίας. \quad β) Το ημικύκλιο είναι $180^\circ$. \quad γ) Ίσα τόξα έχουν πάντα ίσα μήκη σε οποιουσδήποτε κύκλους. \quad δ) Ο κύκλος είναι τόξο $360^\circ$.

\zoni{Το χρησιμοποιώ}{ρόδες, ρολόγια, πίτες}

\ex \textbf{Η ρόδα.} Μια ρόδα λούνα παρκ έχει $12$ ίσες καμπίνες. Πόσες μοίρες γυρίζει για να έρθει η επόμενη καμπίνα στη θέση της;

\ex \textbf{Το ρολόι.} Πόσες μοίρες διαγράφει η άκρη του λεπτοδείκτη σε $20$ λεπτά; Σε $45$ λεπτά; Σε μία ώρα;

\ex \textbf{Η πίτα.} Μια πίτα κόβεται σε φέτες με επίκεντρες γωνίες $45^\circ$. Πόσες φέτες βγαίνουν; Τι μέρος της πίτας είναι $3$ φέτες;

\ex \textbf{Το ραντάρ.} Ένα ραντάρ περιστρέφεται και σαρώνει ένα τόξο $75^\circ$. Τι μέρος του κύκλου καλύπτει;

\ex \textbf{Ο ανεμόμυλος.} Ένας ανεμόμυλος με $4$ πτερύγια γυρίζει $90^\circ$. Πόσα τέτοια τόξα χρειάζονται για μία πλήρη περιστροφή;

\zoni{Το σκέφτομαι}{μέτρο ή μήκος;}

\ex Δύο κύκλοι έχουν ίδιο κέντρο και διαφορετικές ακτίνες. Οι ίδιες δύο ακτίνες ορίζουν στον καθένα ένα τόξο. Είναι ίσα τα μέτρα τους; Είναι ίσα τα μήκη τους;

\ex Γιατί ένας πλήρης κύκλος είναι $360^\circ$ και όχι $100^\circ$; Θα άλλαζε τίποτα αν το διαλέγαμε $100$;

\ex Μια χορδή $\mathrm{ΑΒ}$ ορίζει τόξο $120^\circ$. Πόσο είναι το άλλο τόξο της;

\ex Αν η επίκεντρη γωνία διπλασιαστεί, τι γίνεται το τόξο;

\zoni{Για το σπίτι}{ο κύκλος και οι μοίρες}

\ex Σχεδίασε έναν κύκλο και χώρισέ τον σε $6$ ίσα τόξα με διαβήτη και μοιρογνωμόνιο. Πόσες μοίρες είναι το καθένα;

\ex Βρες κάτι στο σπίτι (ρολόι, πίτα, ρόδα) που χωρίζεται σε ίσα τόξα και υπολόγισε τις μοίρες κάθε τόξου.

\paragrafos{Θέσεις ευθείας και κύκλου: εξωτερική, εφαπτομένη, τέμνουσα}{Β.1.13, «Θέσεις ευθείας και κύκλου»}

\begin{anagke}
Γύρω από το στρογγυλό παρτέρι θα περάσει ένα μονοπάτι, και η Α3 συζητάει πού θα το τοποθετήσει. Ο Νίκος προτείνει να περνά ακριβώς πάνω από το παρτέρι. «Θα πατάμε τα λουλούδια!», διαμαρτύρεται η Ρήνη. Ο Βασίλης προτείνει ένα μονοπάτι που περνά μακριά από το παρτέρι, χωρίς να το αγγίζει. «Τότε δεν θα το βλέπουμε από κοντά», λέει η Αγγελική. Η Πολυξένη προτείνει το τρίτο: ένα μονοπάτι που ακουμπά το παρτέρι σε ένα μόνο σημείο, σαν τον δρόμο που ακουμπά μια ρόδα.

Η κυρία Αλεξάνδρα το ονομάζει: τρεις θέσεις μιας ευθείας ως προς έναν κύκλο. Δείχνει πώς καταλαβαίνουμε ποια είναι από την απόσταση του κέντρου από την ευθεία: «Μετρήστε την απόσταση του κέντρου από το μονοπάτι και συγκρίνετέ την με την ακτίνα».

Στο τέλος του κεφαλαίου δένουμε την απόσταση σημείου από ευθεία (1.10) με τον κύκλο (1.11): τρεις θέσεις ευθείας και κύκλου, και τι ισχύει όταν η ευθεία εφάπτεται.
\end{anagke}

\begin{lexi}{εφαπτομένη, τέμνουσα}{εφαπτομένη, τέμνουσα, εξωτερική}
Η \lexiword{εφαπτομένη} είναι αυτή που «εφάπτεται», δηλαδή «ακουμπά», μόνο σε ένα σημείο. Η \lexiword{τέμνουσα} είναι αυτή που «τέμνει», περνά μέσα από τον κύκλο και τον τέμνει σε δύο σημεία. Η \lexiword{εξωτερική} ευθεία βρίσκεται «έξω» από τον κύκλο και δεν τον ακουμπά.
\end{lexi}

\begin{metafrasi}
Η απόσταση του κέντρου από την ευθεία. & $\mathrm{ΟΜ}=d$\\
Η ευθεία είναι μακριά από τον κύκλο. & $d>\rho$: εξωτερική\\
Η ευθεία ακουμπά τον κύκλο σε ένα σημείο. & $d=\rho$: εφαπτομένη\\
Η ευθεία περνά μέσα στον κύκλο. & $d<\rho$: τέμνουσα\\
\end{metafrasi}

\tmima{Η θεωρία}

Έστω κύκλος $(\mathrm{Ο},\rho)$ και ευθεία $(\varepsilon)$ και $d=\mathrm{ΟΜ}$ η απόσταση του κέντρου από την ευθεία (με $\mathrm{ΟΜ}\perp\varepsilon$). Υπάρχουν τρεις περιπτώσεις:

\begin{center}
\begin{tikzpicture}[scale=0.85]
\begin{scope}
\draw (0,0) circle (1.2); \fill (0,0) circle (1.5pt) node[above left,font=\small] {$\mathrm{Ο}$};
\draw[thick,{Stealth}-{Stealth}] (-1.8,1.9)--(1.8,1.9); \draw[gblue] (0,0)--(0,1.9);
\node at (0,-1.7) {\small $d>\rho$: εξωτερική};
\end{scope}
\begin{scope}[xshift=4.5cm]
\draw (0,0) circle (1.2); \fill (0,0) circle (1.5pt);
\draw[thick,{Stealth}-{Stealth}] (-1.8,1.2)--(1.8,1.2); \draw[gblue] (0,0)--(0,1.2); \fill (0,1.2) circle (2pt) node[above right,font=\small] {$\mathrm{Μ}$};
\node at (0,-1.7) {\small $d=\rho$: εφαπτομένη};
\end{scope}
\begin{scope}[xshift=9cm]
\draw (0,0) circle (1.2); \fill (0,0) circle (1.5pt);
\draw[thick,{Stealth}-{Stealth}] (-1.8,0.6)--(1.8,0.6); \draw[gblue] (0,0)--(0,0.6);
\fill (-1.04,0.6) circle (2pt) node[above left,font=\small] {$\mathrm{Α}$}; \fill (1.04,0.6) circle (2pt) node[above right,font=\small] {$\mathrm{Β}$};
\node at (0,-1.7) {\small $d<\rho$: τέμνουσα};
\end{scope}
\end{tikzpicture}
\end{center}

\begin{itemize}[leftmargin=1.2em,itemsep=0.3ex,topsep=2pt]
\item Αν $d>\rho$, η ευθεία δεν έχει κανένα κοινό σημείο με τον κύκλο: είναι \textbf{εξωτερική} του κύκλου.
\item Αν $d=\rho$, η ευθεία έχει \textbf{ένα} κοινό σημείο $\mathrm{Μ}$ με τον κύκλο και λέγεται \textbf{εφαπτομένη} του κύκλου στο $\mathrm{Μ}$. Το $\mathrm{Μ}$ λέγεται σημείο επαφής και η εφαπτομένη είναι \textbf{κάθετη} στην ακτίνα $\mathrm{ΟΜ}$.
\item Αν $d<\rho$, η ευθεία έχει \textbf{δύο} κοινά σημεία $\mathrm{Α}$, $\mathrm{Β}$ με τον κύκλο και λέγεται \textbf{τέμνουσα}.
\end{itemize}

Αντίστροφα, από το πλήθος των κοινών σημείων προκύπτει η θέση: $0$ κοινά σημεία σημαίνει εξωτερική, $1$ εφαπτομένη, $2$ τέμνουσα. Μια ευθεία δεν μπορεί να έχει με κύκλο περισσότερα από δύο κοινά σημεία.

\begin{prosoxi}
\begin{itemize}[leftmargin=1.2em,itemsep=0.3ex,topsep=2pt]
\item Η απόσταση του κέντρου από την ευθεία είναι η κάθετη απόσταση, όχι οποιοδήποτε πλάγιο τμήμα.
\item Η εφαπτομένη έχει ένα μόνο κοινό σημείο με τον κύκλο. Αν έχει δύο, είναι τέμνουσα.
\item Η εφαπτομένη είναι κάθετη στην ακτίνα του σημείου επαφής, όχι σε οποιαδήποτε ακτίνα.
\item Σε ένα σημείο του κύκλου υπάρχει μία και μοναδική εφαπτομένη.
\end{itemize}
\end{prosoxi}

\tmima{Λυμένα παραδείγματα}

\paradeigma{Η θέση του μονοπατιού}
\tixero Κύκλος $(\mathrm{Ο},3)$ και ευθείες με αποστάσεις $d$ από το κέντρο: $5$, $3$, $1{,}5$.\par
\tithelo Η θέση καθεμίας.\par
\possindeo Συγκρίνω το $d$ με το $\rho=3$. $5>3$: εξωτερική. $3=3$: εφαπτομένη. $1{,}5<3$: τέμνουσα.\par
\apantisi Εξωτερική, εφαπτομένη, τέμνουσα.\par
\elegxos Στην τέμνουσα υπάρχουν δύο κοινά σημεία.

\paradeigma{Κοινά σημεία}
\tixero Μια ευθεία με απόσταση $d=4$ από το κέντρο κύκλου ακτίνας $\rho=4$.\par
\tithelo Πόσα κοινά σημεία έχει.\par
\possindeo $d=\rho$, εφαπτομένη, άρα ένα κοινό σημείο.\par
\apantisi Ένα.\par
\elegxos Το σημείο επαφής είναι το ίχνος της καθέτου από το κέντρο.

\paradeigma{Η εφαπτομένη σε σημείο}
\tixero Ένα σημείο $\mathrm{Μ}$ του κύκλου $(\mathrm{Ο},\rho)$.\par
\tithelo Πώς σχεδιάζω την εφαπτομένη στο $\mathrm{Μ}$;\par
\possindeo Φέρνω την ακτίνα $\mathrm{ΟΜ}$ και την κάθετη στην $\mathrm{ΟΜ}$ στο σημείο $\mathrm{Μ}$.\par
\apantisi Η εφαπτομένη είναι η κάθετη στην ακτίνα $\mathrm{ΟΜ}$ στο $\mathrm{Μ}$.\par
\elegxos Η απόσταση του $\mathrm{Ο}$ από την ευθεία είναι $\mathrm{ΟΜ}=\rho$.

\paradeigma{Πόση απόσταση;}
\tixero Κύκλος $(\mathrm{Ο},5)$ και ευθεία που τον τέμνει σε δύο σημεία.\par
\tithelo Τι ξέρουμε για την απόσταση $d$ του κέντρου από την ευθεία.\par
\possindeo Αφού είναι τέμνουσα, $d<5$.\par
\apantisi $0\le d<5$.\par
\elegxos Αν $d=0$ η ευθεία περνά από το κέντρο και είναι διάμετρος.

\tmima{Ασκήσεις}

\zoni{Το κατάλαβα}{τρεις θέσεις}

\ex Κύκλος $(\mathrm{Ο},3)$. Πες τη θέση της ευθείας όταν η απόσταση του κέντρου από αυτήν είναι: α) $5$ \quad β) $3$ \quad γ) $1{,}5$ \quad δ) $0$ \quad ε) $3{,}01$.

\ex Πόσα κοινά σημεία έχει ευθεία με κύκλο όταν είναι εξωτερική, εφαπτομένη, τέμνουσα;

\ex Σχεδίασε κύκλο με ακτίνα $3$ cm και φέρε μια εξωτερική, μια εφαπτομένη και μια τέμνουσα. Μέτρησε την απόσταση του κέντρου από την καθεμία.

\ex Σχεδίασε την εφαπτομένη ενός κύκλου σε ένα σημείο του $\mathrm{Μ}$.

\ex Κύκλος $(\mathrm{Ο},5)$ και ευθεία που τον τέμνει σε δύο σημεία. Μεταξύ ποιων τιμών βρίσκεται η απόσταση $d$ του κέντρου από την ευθεία;

\ex Σωστό ή λάθος; α) Η εφαπτομένη έχει ένα κοινό σημείο με τον κύκλο. \quad β) Η τέμνουσα έχει τρία κοινά σημεία. \quad γ) Η εφαπτομένη είναι κάθετη στην ακτίνα του σημείου επαφής. \quad δ) Αν $d>\rho$ η ευθεία τέμνει τον κύκλο.

\zoni{Το χρησιμοποιώ}{μονοπάτια και ρόδες}

\ex \textbf{Το μονοπάτι.} Το παρτέρι έχει ακτίνα $2{,}5$ μ. Ένα ευθύγραμμο μονοπάτι περνά σε απόσταση $2{,}5$ μ. από το κέντρο. Πατάμε τα λουλούδια;

\ex \textbf{Τρία μονοπάτια.} Τρία ευθύγραμμα μονοπάτια περνούν σε αποστάσεις $1{,}8$ μ., $2{,}5$ μ. και $3{,}4$ μ. από το κέντρο του παρτεριού ακτίνας $2{,}5$ μ. Ποιο τέμνει το παρτέρι, ποιο το εφάπτεται και ποιο δεν το ακουμπά;

\ex \textbf{Ο τροχός.} Ένας τροχός ακτίνας $33$ cm κυλά σε ευθύγραμμο δρόμο. Πόσο απέχει το κέντρο του από το δρόμο; Πόσα κοινά σημεία έχει ο τροχός με τον δρόμο;

\ex \textbf{Ο φακός.} Ένας φακός φωτίζει ευθεία γραμμή που εφάπτεται σε ένα κυκλικό δωμάτιο. Πού ακουμπά το φως στον τοίχο;

\ex \textbf{Η ρόδα.} Γιατί μια ρόδα κυλά ομαλά πάνω σε επίπεδο δρόμο; (σκέψου ότι ο δρόμος εφάπτεται στη ρόδα).

\zoni{Το σκέφτομαι}{όταν η ευθεία ακουμπά}

\ex Γιατί η εφαπτομένη είναι κάθετη στην ακτίνα του σημείου επαφής; (Σκέψου ότι είναι η μικρότερη απόσταση του κέντρου από την ευθεία.)

\ex Ευθεία περνά από δύο διαφορετικά σημεία του κύκλου. Μπορεί να έχει και τρίτο κοινό σημείο; Εξήγησε.

\ex Πόσες εφαπτομένες διέρχονται από σημείο του κύκλου; Και από σημείο εκτός κύκλου; (Σχεδίασε.)

\ex Δύο παράλληλες ευθείες εφάπτονται σε κύκλο ακτίνας $\rho$. Πόση είναι η απόσταση ανάμεσά τους;

\zoni{Για το σπίτι}{εφαπτόμενα σχήματα}

\ex Βρες στο σπίτι πέντε περιπτώσεις ευθείας και κύκλου με κοινό ένα σημείο (π.χ. τροχός πάνω στο πάτωμα).

\ex Σχεδίασε έναν κύκλο και τρεις εφαπτομένες του σε διαφορετικά σημεία.

\epanalipsi{Β1}

\tmima{Λεξικό του κεφαλαίου}

\renewcommand{\arraystretch}{1.25}
\begin{tabular*}{\linewidth}{@{\extracolsep{\fill}} >{\raggedright\arraybackslash}p{0.28\linewidth} >{\raggedright\arraybackslash}p{0.62\linewidth}}
\toprule
\textbf{Λέξη ή σύμβολο} & \textbf{Με μία πρόταση}\\
\midrule
ευθύγραμμο τμήμα, ευθεία, ημιευθεία & Με δύο άκρα, χωρίς άκρα, με μία αρχή. Από δύο σημεία περνά μία ευθεία.\\
μέσο & Το σημείο $\mathrm{Μ}$ του $\mathrm{ΑΒ}$ με $\mathrm{ΑΜ}=\mathrm{ΜΒ}$.\\
περίμετρος & Το άθροισμα των πλευρών ενός πολυγώνου.\\
γωνία, διχοτόμος & Δύο ημιευθείες με κοινή αρχή· η διχοτόμος τη χωρίζει σε δύο ίσες.\\
οξεία, ορθή, αμβλεία & $<90^\circ$, $=90^\circ$, ανάμεσα σε $90^\circ$ και $180^\circ$.\\
συμπληρωματικές, παραπληρωματικές & Άθροισμα $90^\circ$ και $180^\circ$ αντίστοιχα.\\
κατακορυφήν & Οι απέναντι γωνίες τεμνόμενων ευθειών: ίσες.\\
$\varepsilon_1\parallel\varepsilon_2$, $\varepsilon_1\perp\varepsilon_2$ & Παράλληλες: χωρίς κοινό σημείο. Κάθετες: σχηματίζουν ορθή.\\
κύκλος $(\mathrm{Ο},\rho)$ & Τα σημεία που απέχουν $\rho$ από το $\mathrm{Ο}$. Διάμετρος $\delta=2\rho$.\\
τόξο & Το μέτρο του (σε μοίρες) ισούται με της επίκεντρης γωνίας.\\
εφαπτομένη & Ευθεία με ένα κοινό σημείο με τον κύκλο, κάθετη στην ακτίνα: $d=\rho$.\\
\bottomrule
\end{tabular*}

\tmima{Ασκήσεις επανάληψης}

\ex Πόσες ευθείες ορίζουν τρία μη συνευθειακά σημεία και πόσες τέσσερα σημεία, αν κανένα τριάδα δεν είναι συνευθειακή;

\ex Πόσες διαγώνιους έχουν το πεντάγωνο και το εξάγωνο;

\ex Ένα τμήμα $\mathrm{ΑΒ}=7{,}2$ cm. Βρες το μέσο του και σύγκρινέ το με τμήμα $70$ mm.

\ex Τα $\mathrm{Α},\,\mathrm{Β},\,\Gamma,\,\Delta$ είναι στη σειρά με $\mathrm{ΑΒ}=2{,}5$, $\mathrm{ΒΓ}=3$, $\Gamma\Delta=4$. Βρες τα $\mathrm{ΑΔ}$ και $\mathrm{ΒΔ}$.

\ex Βρες την περίμετρο ενός ορθογωνίου με πλευρές $6{,}5$ και $4$.

\ex Για τη γωνία $124^\circ$ βρες τη διχοτόμο (μέτρο μισής γωνίας) και το παραπλήρωμά της. Βρες το συμπλήρωμα της γωνίας $36^\circ$.

\ex Δύο ευθείες τέμνονται και μία γωνία είναι $65^\circ$. Βρες τις άλλες τρεις.

\ex Τρεις διαδοχικές γωνίες γύρω από σημείο είναι $90^\circ$, $80^\circ$, $110^\circ$. Βρες την τέταρτη.

\ex Τι είδους γωνία σχηματίζουν οι δείκτες του ρολογιού στις $5{:}00$;

\ex Πόσα σημεία τομής έχουν το πολύ $4$ ευθείες του επιπέδου; Από σημείο εκτός ευθείας πόσες παράλληλες διέρχονται;

\ex Τρεις παράλληλες ευθείες απέχουν διαδοχικά $1{,}8$ cm. Πόση απόσταση έχει η πρώτη από την τρίτη;

\ex Στον κύκλο $(\mathrm{Ο},6)$: πόση είναι η διάμετρος; Πες αν τα σημεία με $\mathrm{ΟΑ}=4$ και $\mathrm{ΟΒ}=7$ είναι μέσα ή έξω.

\ex Ένας κύκλος χωρίζεται σε $9$ ίσα τόξα. Πόσες μοίρες είναι το καθένα; Πόσο είναι ένα τόξο που αντιστοιχεί σε επίκεντρη γωνία $150^\circ$;

\ex Κύκλος $(\mathrm{Ο},5)$. Πες τη θέση ευθείας με απόσταση από το κέντρο: α) $5{,}3$ \quad β) $5$ \quad γ) $2$. Πόσα κοινά σημεία έχει καθεμία;

\ex \textbf{Ερωτήσεις αυτοαξιολόγησης (τμήματα, ευθείες).} Σωστό ή λάθος; α) Από δύο σημεία περνά μία ευθεία. \quad β) Η ημιευθεία έχει δύο άκρα. \quad γ) Το μέσο ενός τμήματος είναι μοναδικό. \quad δ) Η περίμετρος ισούται με το άθροισμα των πλευρών. \quad ε) Το ευθύγραμμο τμήμα είναι συντομότερο από κάθε τεθλασμένη με τα ίδια άκρα.

\ex \textbf{Ερωτήσεις αυτοαξιολόγησης (γωνίες, κύκλος).} Σωστό ή λάθος; α) Οι κατακορυφήν γωνίες είναι ίσες. \quad β) Οι συμπληρωματικές γωνίες έχουν άθροισμα $180^\circ$. \quad γ) Δύο ευθείες του επιπέδου ή τέμνονται ή είναι παράλληλες. \quad δ) Η εφαπτομένη είναι κάθετη στην ακτίνα του σημείου επαφής. \quad ε) Η διάμετρος είναι η μεγαλύτερη χορδή. \quad στ) Το κάθετο τμήμα από σημείο σε ευθεία είναι το μεγαλύτερο.

\tmima{Ο χάρτης της επανάληψης}

Μετά τις λύσεις, δες σε ποιες ασκήσεις δυσκολεύτηκες. Ο πίνακας σου λέει ποια παράγραφο να ξαναδιαβάσεις.

\renewcommand{\arraystretch}{1.5}
\begin{tabular*}{\linewidth}{@{\extracolsep{\fill}} >{\raggedright\arraybackslash}p{0.38\linewidth} >{\raggedright\arraybackslash}p{0.52\linewidth}}
\toprule
\textbf{Αν δυσκολεύτηκες στις} & \textbf{ξαναδιάβασε τις παραγράφους}\\
\midrule
ΕΒ1.1 & 1.1 (σημείο, ευθεία, ημιευθεία)\\
ΕΒ1.2 & 1.2 (γωνία, πολύγωνα)\\
ΕΒ1.3, ΕΒ1.4, ΕΒ1.5 & 1.3 και 1.4 (μήκη, μέσο, περίμετρος)\\
ΕΒ1.6 & 1.5 (μοίρες, διχοτόμος) και 1.8\\
ΕΒ1.7 & 1.8 (κατακορυφήν γωνίες)\\
ΕΒ1.8, ΕΒ1.9 & 1.6 και 1.7 (είδη γωνιών, άθροισμα)\\
ΕΒ1.10, ΕΒ1.11 & 1.9 και 1.10 (παράλληλες, αποστάσεις)\\
ΕΒ1.12 & 1.11 (κύκλος)\\
ΕΒ1.13 & 1.12 (τόξα)\\
ΕΒ1.14 & 1.13 (ευθεία και κύκλος)\\
ΕΒ1.15, ΕΒ1.16 & 1.1 έως 1.13\\
\bottomrule
\end{tabular*}


\kefalaio{Συμμετρία}{Ο καθρέφτης δείχνει ίσα, αλλά αντίστροφα}{%
\begin{tabular}{@{}l@{\hspace{1em}}p{0.82\linewidth}@{}}
2.1 & Συμμετρία ως προς άξονα\\[2pt]
2.2 & Άξονας συμμετρίας σχήματος\\[2pt]
2.3 & Μεσοκάθετος ευθυγράμμου τμήματος\\[2pt]
2.4 & Συμμετρία ως προς σημείο\\[2pt]
2.5 & Κέντρο συμμετρίας\\[2pt]
2.6 & Παράλληλες ευθείες που τέμνονται από μια άλλη ευθεία\\[2pt]
 & Επανάληψη του κεφαλαίου\\
\end{tabular}}
\paragrafos{Συμμετρία ως προς άξονα: ο καθρέφτης του κήπου}{Β.2.1, «Συμμετρία ως προς άξονα»}

\begin{anagke}
Η Αγγελική σχεδιάζει πια το δεύτερο σχέδιο του κήπου: έναν κήπο με κεντρικό ευθύγραμμο μονοπάτι, και δύο όμοια παρτέρια, ένα δεξιά και ένα αριστερά του. «Θέλω να είναι το ένα η εικόνα του άλλου στον καθρέφτη», λέει. Ο Στέλιος ρωτάει: «Ποιο σημείο του αριστερού παρτεριού αντιστοιχεί στο σημείο του δεξιού;». Η Αγγελική δοκιμάζει ένα κόλπο: σχεδιάζει το αριστερό παρτέρι στο φύλλο, διπλώνει το φύλλο κατά μήκος του μονοπατιού και με μια σταθερή πίεση αντιγράφει το σχήμα στην άλλη πλευρά.

«Κάθε σημείο πέφτει ακριβώς πάνω σε ένα σημείο της άλλης πλευράς!», λέει ο Νίκος. Η Μυρτώ το ελέγχει με τον χάρακα: το δέντρο που απέχει $1{,}2$ μέτρα από το μονοπάτι έχει «ζευγάρι» ένα δέντρο που απέχει επίσης $1{,}2$ μέτρα, από την άλλη πλευρά και «ακριβώς απέναντι», πάνω στην κάθετη. Η Πολυξένη γράφει στον πίνακα: «Συμμετρία».

Σε αυτή την παράγραφο ορίζουμε ποιο είναι το συμμετρικό ενός σημείου και ενός σχήματος ως προς μια ευθεία, και ανακαλύπτουμε ότι η συμμετρία διατηρεί τα μήκη και τις γωνίες.
\end{anagke}

\begin{lexi}{συμμετρία, άξονας}{συμμετρία, άξονας}
Η λέξη \lexiword{συμμετρία} (σύν + μέτρο) σημαίνει «κοινό μέτρο»: ό,τι έχει ίσα μέτρα από δύο μεριές. Ο \lexiword{άξονας} είναι η ευθεία γύρω από την οποία «στρέφεται» ο καθρέφτης, αυτός που χωρίζει τα δύο συμμετρικά μέρη.
\end{lexi}

\begin{metafrasi}
Το σημείο που πέφτει πάνω στο $\mathrm{Β}$ όταν διπλώνω. & το συμμετρικό $\mathrm{Β}'$ του $\mathrm{Β}$ ως προς την $\varepsilon$\\
Ίση απόσταση, κάθετα. & $\mathrm{ΒΒ}'\perp\varepsilon$ και η $\varepsilon$ περνά από το μέσο του $\mathrm{ΒΒ}'$\\
Σημείο πάνω στον άξονα. & είναι συμμετρικό του εαυτού του\\
Ο καθρέφτης δεν αλλάζει μήκη. & $\mathrm{ΑΒ}=\mathrm{Α}'\mathrm{Β}'$\\
\end{metafrasi}

\tmima{Η θεωρία}

\begin{orismos}{συμμετρικό σημείου ως προς ευθεία}
\textbf{Συμμετρικό} του σημείου $\mathrm{Β}$ ως προς την ευθεία $\varepsilon$ είναι το σημείο $\mathrm{Β}'$ με το οποίο συμπίπτει το $\mathrm{Β}$ όταν διπλώσουμε το επίπεδο κατά μήκος της $\varepsilon$. Η $\varepsilon$ λέγεται \textbf{άξονας συμμετρίας}. Κάθε σημείο της $\varepsilon$ είναι συμμετρικό του εαυτού του.
\end{orismos}

Από τη δίπλωση προκύπτει ότι το τμήμα $\mathrm{ΒΒ}'$ είναι \textbf{κάθετο} στην $\varepsilon$ και η $\varepsilon$ περνά από το \textbf{μέσο} του: τα $\mathrm{Β}$ και $\mathrm{Β}'$ απέχουν το ίδιο από την $\varepsilon$, από δύο αντίθετες μεριές. Έτσι κατασκευάζουμε το συμμετρικό με \emph{γνώμονα}: φέρνουμε την κάθετη από το $\mathrm{Β}$ προς την $\varepsilon$, και στην προέκτασή της παίρνουμε σημείο $\mathrm{Β}'$ με την ίδια απόσταση από την $\varepsilon$.

\begin{center}
\begin{tikzpicture}[scale=1.1]
\draw[thick,{Stealth}-{Stealth}] (0,-0.2)--(0,3.4) node[above] {$\varepsilon$};
\fill (-1.6,2.5) circle (2pt) node[left] {$\mathrm{Β}$}; \fill (1.6,2.5) circle (2pt) node[right] {$\mathrm{Β}'$};
\draw[dashed,gblue] (-1.6,2.5)--(1.6,2.5);
\fill (-1.2,0.6) circle (2pt) node[left] {$\mathrm{Α}$}; \fill (1.2,0.6) circle (2pt) node[right] {$\mathrm{Α}'$};
\draw[dashed,gblue] (-1.2,0.6)--(1.2,0.6);
\draw (-1.6,2.5)--(-1.2,0.6); \draw (1.6,2.5)--(1.2,0.6);
\draw (0,2.5)++(0.15,0)--++(0,-0.15)--++(-0.15,0);
\end{tikzpicture}
\end{center}

\begin{kanonas}{συμμετρικά σχήματα}
Δύο σχήματα $(\Sigma_1)$ και $(\Sigma_2)$ λέγονται \textbf{συμμετρικά ως προς την $\varepsilon$} όταν το καθένα αποτελείται από τα συμμετρικά των σημείων του άλλου. Τα συμμετρικά ως προς ευθεία σχήματα είναι \textbf{ίσα}: έχουν ίσα μήκη και ίσες γωνίες (με δίπλωση συμπίπτουν).
\end{kanonas}

Έτσι το συμμετρικό ενός ευθυγράμμου τμήματος είναι ευθύγραμμο τμήμα με το ίδιο μήκος και το συμμετρικό μιας γωνίας είναι γωνία με το ίδιο μέτρο. Για να σχεδιάσω το συμμετρικό ενός πολυγώνου, βρίσκω τα συμμετρικά των κορυφών και τα ενώνω με την ίδια σειρά.

\begin{prosoxi}
\begin{itemize}[leftmargin=1.2em,itemsep=0.3ex,topsep=2pt]
\item Το τμήμα $\mathrm{ΒΒ}'$ είναι \emph{κάθετο} στον άξονα (όχι παράλληλο).
\item Η απόσταση του $\mathrm{Β}$ από τον άξονα είναι ίση με της $\mathrm{Β}'$, όχι ίση με την απόσταση $\mathrm{ΒΒ}'$: το $\mathrm{ΒΒ}'$ είναι διπλάσιο.
\item Το συμμετρικό ενός σχήματος είναι ίσο με αυτό, αλλά «ανάποδο» (σαν καθρέφτης): τη φορά δεν την κρατάει.
\item Το συμμετρικό του $\mathrm{Β}'$ είναι πάλι το $\mathrm{Β}$.
\end{itemize}
\end{prosoxi}

\tmima{Λυμένα παραδείγματα}

\paradeigma{Το συμμετρικό σημείου}
\tixero Το δέντρο $\mathrm{Β}$ απέχει $1{,}2$ μ. από την ευθεία $\varepsilon$.\par
\tithelo Πού βρίσκεται το συμμετρικό του και πόσο απέχουν τα δύο.\par
\possindeo Το $\mathrm{Β}'$ είναι στην κάθετη από το $\mathrm{Β}$ προς την $\varepsilon$, από την άλλη μεριά, με απόσταση $1{,}2$ μ. Η απόσταση $\mathrm{ΒΒ}'=2{,}4$ μ.\par
\apantisi $\mathrm{Β}'$ απέχει $1{,}2$ μ. από την $\varepsilon$ και $\mathrm{ΒΒ}'=2{,}4$ μ.\par
\elegxos Η $\varepsilon$ είναι στο μέσο του $\mathrm{ΒΒ}'$: $1{,}2+1{,}2=2{,}4$.

\paradeigma{Σημείο πάνω στον άξονα}
\tixero Το σημείο $\mathrm{Μ}$ βρίσκεται πάνω στην $\varepsilon$.\par
\tithelo Το συμμετρικό του.\par
\possindeo Αφού διπλώνω κατά μήκος της $\varepsilon$, το $\mathrm{Μ}$ δεν κινείται.\par
\apantisi Το συμμετρικό του είναι το ίδιο το $\mathrm{Μ}$.\par
\elegxos Η απόσταση του $\mathrm{Μ}$ από την $\varepsilon$ είναι $0$.

\paradeigma{Συμμετρικό τριγώνου}
\tixero Τρίγωνο $\mathrm{ΑΒΓ}$ με $\mathrm{ΑΒ}=5$ cm, $\widehat{\mathrm{Β}}=40^\circ$ και το συμμετρικό του $\mathrm{Α}'\mathrm{Β}'\Gamma'$.\par
\tithelo Το $\mathrm{Α}'\mathrm{Β}'$ και η $\widehat{\mathrm{Β}'}$.\par
\possindeo Η συμμετρία διατηρεί μήκη και γωνίες.\par
\apantisi $\mathrm{Α}'\mathrm{Β}'=5$ cm και $\widehat{\mathrm{Β}'}=40^\circ$.\par
\elegxos Τα συμμετρικά σχήματα είναι ίσα.

\paradeigma{Κατασκευή συμμετρικού}
\tixero Η ευθεία $\varepsilon$ και το σημείο $\mathrm{Α}$ εκτός αυτής.\par
\tithelo Να κατασκευάσω το συμμετρικό του.\par
\possindeo (1) Φέρνω από το $\mathrm{Α}$ την κάθετη προς την $\varepsilon$ με γνώμονα και χάρακα. (2) Μετράω την απόσταση $\mathrm{ΑΜ}$ ως το ίχνος $\mathrm{Μ}$. (3) Στην προέκταση βρίσκω το $\mathrm{Α}'$ με $\mathrm{ΜΑ}'=\mathrm{ΜΑ}$.\par
\apantisi Το σημείο $\mathrm{Α}'$.\par
\elegxos Διπλώνω το χαρτί κατά μήκος της $\varepsilon$: το $\mathrm{Α}$ πέφτει πάνω στο $\mathrm{Α}'$.

\tmima{Ασκήσεις}

\zoni{Το κατάλαβα}{συμμετρικό σημείου και σχήματος}

\ex Σχεδίασε ευθεία $\varepsilon$ και σημείο $\mathrm{Α}$ εκτός της. Βρες το συμμετρικό του με γνώμονα και με δίπλωση.

\ex Ποιο είναι το συμμετρικό ενός σημείου που ανήκει στην $\varepsilon$;

\ex Το $\mathrm{Α}$ απέχει $3{,}5$ cm από την $\varepsilon$. Πόσο απέχει το συμμετρικό του από την $\varepsilon$ και πόσο από το $\mathrm{Α}$;

\ex Σχεδίασε ένα τρίγωνο $\mathrm{ΑΒΓ}$ και το συμμετρικό του ως προς μια ευθεία $\varepsilon$ που δεν το τέμνει.

\ex Αν στο τρίγωνο $\mathrm{ΑΒΓ}$ η γωνία $\widehat{\mathrm{Β}}$ είναι $40^\circ$ και $\mathrm{ΑΒ}=5$ cm, τι ξέρεις για το συμμετρικό τρίγωνο;

\ex Σχεδίασε το συμμετρικό ενός ευθυγράμμου τμήματος $\mathrm{ΑΒ}=4{,}2$ cm ως προς ευθεία $\varepsilon$. Πόσο είναι το μήκος του;

\ex Σωστό ή λάθος; α) Κάθε σημείο της $\varepsilon$ είναι συμμετρικό του εαυτού του. \quad β) Το $\mathrm{ΑΑ}'$ είναι παράλληλο προς την $\varepsilon$. \quad γ) Συμμετρικά σχήματα είναι ίσα. \quad δ) Το συμμετρικό του $\mathrm{Α}'$ είναι το $\mathrm{Α}$. \quad ε) Η συμμετρία αλλάζει τα μήκη.

\zoni{Το χρησιμοποιώ}{κήπος, καθρέφτες, γράμματα}

\ex \textbf{Ο καθρέφτης του κήπου.} Ένα δέντρο απέχει $1{,}2$ μ. από το κεντρικό μονοπάτι. Πού πρέπει να φυτευτεί το συμμετρικό του; Πόσο απέχουν τα δύο δέντρα;

\ex \textbf{Η λίμνη.} Η αντανάκλαση ενός δέντρου στη λίμνη είναι το συμμετρικό του ως προς την ακτή. Αν το δέντρο είναι $4$ μ. ψηλό, πόσο «βαθιά» φαίνεται η αντανάκλαση;

\ex \textbf{Το ρολόι.} Στον καθρέφτη ένα ρολόι δείχνει $3{:}00$. Τι ώρα είναι πραγματικά; (Σκέψου πού είναι ο ωροδείκτης.)

\ex \textbf{Τα γράμματα.} Ποια από τα κεφαλαία γράμματα $\mathrm{Α}$, $\mathrm{Β}$, $\mathrm{Η}$, $\mathrm{Μ}$, $\mathrm{Ο}$, $\mathrm{Τ}$, $\mathrm{Ζ}$ μένουν ίδια όταν τα καθρεφτίσεις ως προς κατακόρυφη ευθεία;

\ex \textbf{Η πεταλούδα.} Σχεδίασε στο τετράδιο μια ευθεία και το μισό μιας πεταλούδας στη μία πλευρά. Συμπλήρωσε το άλλο μισό ως συμμετρικό.

\zoni{Το σκέφτομαι}{γιατί κάθετα;}

\ex Γιατί το τμήμα $\mathrm{ΑΑ}'$ είναι κάθετο στον άξονα; (Σκέψου τη δίπλωση.)

\ex Ποια σημεία του επιπέδου μένουν στη θέση τους όταν καθρεφτίζουμε ως προς την $\varepsilon$;

\ex Αν το $\mathrm{Α}'$ είναι συμμετρικό του $\mathrm{Α}$ ως προς την $\varepsilon$, και το $\mathrm{Α}''$ συμμετρικό του $\mathrm{Α}'$, πού βρίσκεται το $\mathrm{Α}''$;

\ex Το τμήμα $\mathrm{ΑΒ}$ τέμνει την $\varepsilon$ στο $\mathrm{Μ}$. Εξήγησε γιατί το συμμετρικό $\mathrm{Α}'\mathrm{Β}'$ περνά επίσης από το $\mathrm{Μ}$.

\zoni{Για το σπίτι}{πείραμα με μελάνι}

\ex Βάλε λίγο μελάνι στη μέση ενός χαρτιού και δίπλωσέ το. Άνοιξέ το και περιέγραψε τι σχήμα βγαίνει και πού είναι ο άξονας συμμετρίας.

\ex Βρες στο σπίτι πέντε αντικείμενα που φαίνονται συμμετρικά.

\paragrafos{Άξονας συμμετρίας σχήματος}{Β.2.2, «Άξονας συμμετρίας»}

\begin{anagke}
Η Α3 διαλέγει τα σχήματα για τα παρτέρια του κήπου: τετράγωνο, ορθογώνιο, κύκλο, ισόπλευρο τρίγωνο και ένα παραλληλόγραμμο. Ο Βασίλης προτείνει να τα κόψουν από χαρτόνι και να τα διπλώσουν. «Αν διπλώσω ένα τετράγωνο στη μέση, συμπίπτουν τα δύο μισά», λέει. «Αλλά μπορώ να το διπλώσω και με άλλον τρόπο: από γωνία σε γωνία, και πάλι συμπίπτουν!». Η Αγγελική μετράει: το τετράγωνο διπλώνεται με τέσσερις διαφορετικούς τρόπους. Το ορθογώνιο με δύο. Το παραλληλόγραμμο καθόλου: όσο και αν προσπαθήσουν, τα δύο μισά δεν συμπίπτουν.

«Και ο κύκλος;» ρωτάει ο Νίκος. «Μπορώ να τον διπλώσω όπως θέλω, αρκεί η δίπλωση να περνά από το κέντρο!». Η Πολυξένη το γράφει: «Κάθε διπλωμένη ευθεία είναι ένας άξονας συμμετρίας του σχήματος». Η Ρήνη βρίσκει ότι ένα ισοσκελές τρίγωνο έχει έναν άξονα και ένα ισόπλευρο τρεις.

Στην παράγραφο αυτή ορίζουμε τον άξονα συμμετρίας ενός σχήματος και μετράμε πόσους άξονες έχουν τα γνωστά μας σχήματα.
\end{anagke}

\begin{lexi}{άξονας συμμετρίας σχήματος}{άξονας συμμετρίας σχήματος}
Ένα σχήμα έχει \lexiword{άξονα συμμετρίας} όταν υπάρχει ευθεία που το χωρίζει σε δύο μέρη που εφαρμόζουν τέλεια όταν διπλώσουμε κατά μήκος της: το σχήμα είναι δηλαδή το δικό του καθρέφτισμα.
\end{lexi}

\begin{metafrasi}
Η ευθεία που διπλώνει το σχήμα σε δύο ίσα μισά. & άξονας συμμετρίας\\
Το συμμετρικό του σχήματος ως προς τον άξονά του. & το ίδιο το σχήμα\\
Τετράγωνο. & $4$ άξονες\\
Ορθογώνιο (όχι τετράγωνο). & $2$ άξονες\\
Κύκλος. & άπειροι άξονες (κάθε διάμετρος)\\
\end{metafrasi}

\tmima{Η θεωρία}

\begin{orismos}{άξονας συμμετρίας}
\textbf{Άξονας συμμετρίας} ενός σχήματος είναι η ευθεία που το χωρίζει σε δύο μέρη τα οποία συμπίπτουν όταν διπλώσουμε το σχήμα κατά μήκος της. Τότε το συμμετρικό του σχήματος ως προς αυτήν είναι το ίδιο το σχήμα.
\end{orismos}

\begin{center}
\begin{tikzpicture}[scale=0.8]
\begin{scope}
\draw[thick] (0,0) rectangle (2,2); \draw[gblue,dashed] (1,-0.3)--(1,2.3) (-0.3,1)--(2.3,1) (-0.3,-0.3)--(2.3,2.3) (-0.3,2.3)--(2.3,-0.3);
\node at (1,-0.8) {\small τετράγωνο: $4$};
\end{scope}
\begin{scope}[xshift=4cm]
\draw[thick] (0,0) rectangle (3,1.8); \draw[gblue,dashed] (1.5,-0.3)--(1.5,2.1) (-0.3,0.9)--(3.3,0.9);
\node at (1.5,-0.8) {\small ορθογώνιο: $2$};
\end{scope}
\begin{scope}[xshift=9cm]
\draw[thick] (0,0)--(2.4,0)--(1.2,2.08)--cycle; \draw[gblue,dashed] (1.2,-0.3)--(1.2,2.4) (0,0)--(1.8,1.56) (2.4,0)--(0.6,1.56);
\node at (1.2,-0.8) {\small ισόπλευρο: $3$};
\end{scope}
\end{tikzpicture}
\end{center}

Μερικά παραδείγματα:
\begin{itemize}[leftmargin=1.2em,itemsep=0.3ex,topsep=2pt]
\item \textbf{Κύκλος}: κάθε ευθεία που περνά από το κέντρο είναι άξονας συμμετρίας (άπειροι).
\item \textbf{Τετράγωνο} $4$ (οι δύο μεσοκάθετοι των πλευρών και οι δύο διαγώνιοι). \textbf{Ορθογώνιο} $2$ (οι μεσοκάθετοι των πλευρών). \textbf{Ρόμβος} $2$ (οι διαγώνιοι).
\item \textbf{Παραλληλόγραμμο} (που δεν είναι ορθογώνιο ή ρόμβος): κανένας.
\item \textbf{Ισοσκελές τρίγωνο}: $1$, η ευθεία από την κορυφή που διχοτομεί τη γωνία της. \textbf{Ισόπλευρο}: $3$. \textbf{Σκαληνό}: κανένας.
\item \textbf{Κανονικό $\nu$-γωνο}: $\nu$ άξονες.
\item \textbf{Ευθύγραμμο τμήμα}: $2$ (η ευθεία του και η μεσοκάθετός του). \textbf{Γωνία}: ο φορέας της διχοτόμου της.
\end{itemize}

\begin{prosoxi}
\begin{itemize}[leftmargin=1.2em,itemsep=0.3ex,topsep=2pt]
\item Για να είναι άξονας συμμετρίας δεν αρκεί να χωρίζει το σχήμα σε δύο ίσα εμβαδά: πρέπει τα δύο μέρη να συμπίπτουν. Η διαγώνιος ενός (μη τετράγωνου) ορθογωνίου το χωρίζει σε δύο ίσα εμβαδά αλλά δεν είναι άξονας.
\item Το παραλληλόγραμμο δεν έχει άξονα συμμετρίας (παρότι έχει κέντρο συμμετρίας, όπως θα δούμε).
\item Ο άξονας συμμετρίας είναι ολόκληρη \emph{ευθεία}, όχι μόνο το κομμάτι μέσα στο σχήμα.
\item Ένα σχήμα μπορεί να έχει πολλούς άξονες συμμετρίας, και όλοι διέρχονται από το ίδιο σημείο όταν είναι πάνω από έναν.
\end{itemize}
\end{prosoxi}

\tmima{Λυμένα παραδείγματα}

\paradeigma{Τετράγωνο και ορθογώνιο}
\tixero Ένα τετράγωνο και ένα ορθογώνιο $4\times6$.\par
\tithelo Πόσους άξονες έχει το καθένα.\par
\possindeo Το τετράγωνο διπλώνεται: στη μέση οριζόντια, στη μέση κατακόρυφα και κατά μήκος κάθε διαγωνίου: $4$. Το ορθογώνιο μόνο οριζόντια και κατακόρυφα στη μέση: $2$ (οι διαγώνιοι δεν δουλεύουν).\par
\apantisi $4$ και $2$.\par
\elegxos Αν διπλώσω το ορθογώνιο στη διαγώνιο, μια γωνία «βγαίνει έξω».

\paradeigma{Ο κύκλος}
\tixero Ο κύκλος $(\mathrm{Ο},\rho)$.\par
\tithelo Οι άξονες συμμετρίας.\par
\possindeo Κάθε διάμετρος $\delta$ χωρίζει τον κύκλο σε δύο ημικύκλια που εφαρμόζουν. Η ευθεία της διαμέτρου είναι άξονας.\par
\apantisi Άπειροι: οι ευθείες που περνούν από το $\mathrm{Ο}$.\par
\elegxos Δεν υπάρχει άξονας που να μην περνά από το κέντρο.

\paradeigma{Τρίγωνα}
\tixero Ισόπλευρο, ισοσκελές, σκαληνό τρίγωνο.\par
\tithelo Οι άξονες συμμετρίας.\par
\possindeo Το ισόπλευρο έχει από κάθε κορυφή έναν άξονα (στη μέση της απέναντι πλευράς): $3$. Το ισοσκελές έχει μόνο τον άξονα της κορυφής: $1$. Το σκαληνό: κανένα.\par
\apantisi $3$, $1$, $0$.\par
\elegxos Στο ισοσκελές, οι δύο ίσες πλευρές είναι συμμετρικές.

\paradeigma{Τα γράμματα}
\tixero Τα γράμματα $\mathrm{Α}$, $\mathrm{Β}$, $\mathrm{Η}$, $\mathrm{Ο}$, $\mathrm{Ζ}$.\par
\tithelo Ποια έχουν άξονα συμμετρίας.\par
\possindeo Το $\mathrm{Α}$ έχει κατακόρυφο άξονα. Το $\mathrm{Β}$ έχει οριζόντιο. Το $\mathrm{Η}$ έχει δύο (οριζόντιο και κατακόρυφο). Το $\mathrm{Ο}$ έχει δύο (αν είναι έλλειψη) ή άπειρους (αν κύκλος). Το $\mathrm{Ζ}$ κανέναν.\par
\apantisi Όλα εκτός από το $\mathrm{Ζ}$.\par
\elegxos Το $\mathrm{Ζ}$ έχει όμως κέντρο συμμετρίας.

\tmima{Ασκήσεις}

\zoni{Το κατάλαβα}{πόσοι άξονες}

\ex Κόψε από χαρτί ένα τετράγωνο και βρες όλους τους άξονες συμμετρίας του με δίπλωση.

\ex Πόσους άξονες συμμετρίας έχει: α) το ορθογώνιο ($4\times6$) \quad β) ο ρόμβος \quad γ) το ισόπλευρο τρίγωνο \quad δ) το ισοσκελές \quad ε) το σκαληνό \quad στ) το παραλληλόγραμμο.

\ex Πόσους άξονες συμμετρίας έχει το κανονικό πεντάγωνο, το κανονικό εξάγωνο και ο κύκλος;

\ex Σχεδίασε ένα ευθύγραμμο τμήμα $\mathrm{ΑΒ}$. Βρες τους άξονες συμμετρίας του.

\ex Σχεδίασε μια γωνία $x\mathrm{Ο}y$. Βρες τον άξονα συμμετρίας της.

\ex Ποια από τα κεφαλαία γράμματα του ελληνικού αλφαβήτου $\mathrm{Α}$, $\mathrm{Β}$, $\Gamma$, $\Delta$, $\mathrm{Ε}$, $\mathrm{Η}$, $\mathrm{Θ}$, $\mathrm{Ι}$, $\mathrm{Κ}$, $\Lambda$, $\mathrm{Μ}$, $\mathrm{Ν}$, $\mathrm{Ο}$, $\mathrm{Π}$, $\mathrm{Τ}$, $\mathrm{Φ}$, $\mathrm{Χ}$, $\Psi$ έχουν άξονα συμμετρίας;

\ex Σωστό ή λάθος; α) Κάθε σχήμα έχει άξονα συμμετρίας. \quad β) Ο κύκλος έχει άπειρους άξονες. \quad γ) Η διαγώνιος ορθογωνίου είναι άξονας συμμετρίας του. \quad δ) Το τετράγωνο έχει $4$ άξονες. \quad ε) Το ισοσκελές τρίγωνο έχει τουλάχιστον έναν άξονα.

\zoni{Το χρησιμοποιώ}{κήπος, σχέδια, σημαίες}

\ex \textbf{Το παρτέρι.} Ένα παρτέρι είναι κανονικό εξάγωνο. Πόσοι είναι οι άξονες συμμετρίας του; Πόσοι περνούν από δύο κορυφές;

\ex \textbf{Η πόρτα.} Η πόρτα του σχολείου έχει σχήμα ορθογωνίου διαστάσεων $1{,}2\times2{,}1$ μ. Πόσους άξονες συμμετρίας έχει το σχήμα της;

\ex \textbf{Το συντριβάνι.} Σε ένα ορθογώνιο παρτέρι $4\times6$ μ. θέλουμε να βάλουμε ένα συντριβάνι. Σε ποια σημεία θα είναι «συμμετρικό» ως προς όλους τους άξονες;

\ex \textbf{Το λογότυπο.} Σχεδίασε ένα λογότυπο για την Α3 που να έχει ακριβώς δύο άξονες συμμετρίας.

\ex \textbf{Ο χαρταετός.} Ένας χαρταετός έχει σχήμα ρόμβου. Πόσους άξονες συμμετρίας έχει; Πού τέμνονται;

\zoni{Το σκέφτομαι}{όταν η δίπλωση δεν δουλεύει}

\ex Γιατί η διαγώνιος ενός ορθογωνίου (που δεν είναι τετράγωνο) δεν είναι άξονας συμμετρίας, παρότι το χωρίζει σε δύο ίσα τρίγωνα;

\ex Μπορεί ένα σχήμα να έχει ακριβώς δύο άξονες συμμετρίας; Ακριβώς τρεις; Δώσε παραδείγματα.

\ex Αν ένα σχήμα έχει δύο άξονες που τέμνονται, πού βρίσκεται ένα σημείο που «ανήκει» και στους δύο;

\ex Γιατί ο άξονας συμμετρίας ενός ισοσκελούς τριγώνου διχοτομεί τη γωνία της κορυφής;

\zoni{Για το σπίτι}{συμμετρίες γύρω σου}

\ex Βρες πέντε σχήματα στο σπίτι με έναν άξονα, τρία με δύο άξονες και ένα με τρεις ή περισσότερους.

\ex Δίπλωσε ένα χαρτί και κόψε ένα σχήμα. Άνοιξέ το και βρες τον άξονα συμμετρίας του.

\paragrafos{Μεσοκάθετος ευθυγράμμου τμήματος}{Β.2.3, «Μεσοκάθετος ευθυγράμμου τμήματος»}

\begin{anagke}
Στον κήπο υπάρχουν δύο μεγάλα δέντρα, το $\mathrm{Α}$ και το $\mathrm{Β}$, και η παρέα θέλει να τοποθετήσει ένα παγκάκι ώστε να είναι σε ακριβώς ίση απόσταση από τα δύο, για να έχουν όλοι τον ίδιο ίσκιο. Ο Βασίλης προτείνει το μέσο του τμήματος $\mathrm{ΑΒ}$. «Αυτό ισαπέχει», λέει η Πολυξένη, «αλλά είναι μόνο ένα σημείο. Υπάρχουν και άλλα;».

Η Αγγελική παίρνει σχοινί μήκους $5$ μέτρων, καρφώνει το ένα άκρο στο $\mathrm{Α}$ και το άλλο στο $\mathrm{Β}$ και τεντώνει από το μέσο ένα δεύτερο σχοινί του ίδιου μήκους προς κάθε μεριά. Όπου ενώνονται, βρίσκουν σημεία που ισαπέχουν από τα δύο δέντρα. Τα σημεία ενώνονται σε μια ευθεία κάθετη στο $\mathrm{ΑΒ}$, που περνά από το μέσο του. «Αυτή είναι η μεσοκάθετος», λέει η κυρία Αλεξάνδρα.

Η παρέα ανακαλύπτει ότι όλα τα σημεία της μεσοκαθέτου ισαπέχουν από τα άκρα του τμήματος, και ότι κάθε σημείο που ισαπέχει είναι πάνω της. Θα δουν και πώς κατασκευάζεται με τους δύο κλασικούς κανόνες του Ευκλείδη: τον χάρακα χωρίς υποδιαιρέσεις και τον διαβήτη.
\end{anagke}

\begin{lexi}{μεσοκάθετος, ισαπέχω}{μεσοκάθετος, ισαπέχω}
Η \lexiword{μεσοκάθετος} είναι η «κάθετη στο μέσο»: κάθετη προς το τμήμα και μάλιστα στο μέσο του. \lexiword{Ισαπέχω} σημαίνει «απέχω ίσα» (ίση απόσταση) από δύο ή περισσότερα σημεία.
\end{lexi}

\begin{metafrasi}
Η κάθετη στο μέσο του $\mathrm{ΑΒ}$. & η μεσοκάθετος του $\mathrm{ΑΒ}$\\
Ίση απόσταση από τα $\mathrm{Α}$ και $\mathrm{Β}$. & $\mathrm{ΜΑ}=\mathrm{ΜΒ}$\\
Ο άξονας συμμετρίας του τμήματος. & η μεσοκάθετος\\
Η κατασκευή με διαβήτη. & δύο ίσοι κύκλοι, κέντρα $\mathrm{Α}$ και $\mathrm{Β}$\\
\end{metafrasi}

\tmima{Η θεωρία}

\begin{orismos}{μεσοκάθετος}
\textbf{Μεσοκάθετος} ενός ευθυγράμμου τμήματος $\mathrm{ΑΒ}$ είναι η ευθεία που είναι \textbf{κάθετη} στο $\mathrm{ΑΒ}$ και διέρχεται από το \textbf{μέσο} του.
\end{orismos}

\begin{idiotita}{μεσοκάθετος και ισαπόσταση}
(α) Κάθε σημείο της μεσοκαθέτου του $\mathrm{ΑΒ}$ ισαπέχει από τα $\mathrm{Α}$ και $\mathrm{Β}$.\\
(β) Κάθε σημείο που ισαπέχει από τα $\mathrm{Α}$ και $\mathrm{Β}$ ανήκει στη μεσοκάθετο του $\mathrm{ΑΒ}$.\\
(γ) Η μεσοκάθετος του $\mathrm{ΑΒ}$ είναι άξονας συμμετρίας του τμήματος.
\end{idiotita}

Η (γ) λέει ότι το $\mathrm{Α}$ και το $\mathrm{Β}$ είναι συμμετρικά ως προς τη μεσοκάθετο (το $\mathrm{ΑΒ}$ είναι κάθετο στον άξονα και ο άξονας περνά από το μέσο του). Η (α) προκύπτει έτσι: αν $\mathrm{Μ}$ είναι σημείο του άξονα, το συμμετρικό του είναι το ίδιο, άρα η συμμετρία φέρνει το τμήμα $\mathrm{ΜΑ}$ πάνω στο $\mathrm{ΜΒ}$ και $\mathrm{ΜΑ}=\mathrm{ΜΒ}$.

\tmimasmall{Κατασκευή με κανόνα και διαβήτη}

Ο Ευκλείδης στήριζε τις κατασκευές του σε τρεις κανόνες («αιτήματα»): (1) από δύο σημεία διέρχεται μία ευθεία, (2) ένα τμήμα προεκτείνεται απεριόριστα, (3) ο κύκλος ορίζεται από κέντρο και ακτίνα. Έτσι οι αρχαίοι Έλληνες κατασκεύαζαν σχήματα με ένα χάρακα χωρίς υποδιαιρέσεις («κανόνα») και ένα διαβήτη. Για να φέρουμε τη μεσοκάθετο:

\begin{enumerate}[leftmargin=1.6em,itemsep=0.3ex,topsep=2pt]
\item Με κέντρο το $\mathrm{Α}$ και ακτίνα μεγαλύτερη από το μισό του $\mathrm{ΑΒ}$ γράφω κύκλο.
\item Με κέντρο το $\mathrm{Β}$ και την ίδια ακτίνα γράφω δεύτερο κύκλο. Οι δύο κύκλοι τέμνονται σε δύο σημεία $\mathrm{Κ}$ και $\Lambda$.
\item Η ευθεία $\mathrm{ΚΛ}$ είναι η μεσοκάθετος του $\mathrm{ΑΒ}$.
\end{enumerate}

\begin{center}
\begin{tikzpicture}[scale=0.9]
\draw (0,0) circle (2.4); \draw (3.2,0) circle (2.4);
\fill (0,0) circle (2pt) node[below left] {$\mathrm{Α}$}; \fill (3.2,0) circle (2pt) node[below right] {$\mathrm{Β}$};
\draw[thick] (0,0)--(3.2,0);
\coordinate (K) at (1.6,1.789); \coordinate (L) at (1.6,-1.789);
\fill (K) circle (2pt) node[above] {$\mathrm{Κ}$}; \fill (L) circle (2pt) node[below] {$\Lambda$};
\draw[gblue,thick,{Stealth}-{Stealth}] (1.6,2.3)--(1.6,-2.3);
\fill (1.6,0) circle (2pt) node[above right] {$\mathrm{Μ}$};
\end{tikzpicture}
\end{center}

Το σημείο $\mathrm{Μ}$ όπου η $\mathrm{ΚΛ}$ τέμνει το $\mathrm{ΑΒ}$ είναι το μέσο του.

\begin{prosoxi}
\begin{itemize}[leftmargin=1.2em,itemsep=0.3ex,topsep=2pt]
\item Η μεσοκάθετος είναι ευθεία, όχι τμήμα: συνεχίζεται και πέρα από το μέσο.
\item Δεν αρκεί να περνά από το μέσο: πρέπει να είναι και κάθετη. Μια λοξή ευθεία από το μέσο δεν είναι μεσοκάθετος.
\item Στην κατασκευή, και οι δύο κύκλοι πρέπει να έχουν την ίδια ακτίνα, μεγαλύτερη από το μισό του $\mathrm{ΑΒ}$ (αλλιώς δεν τέμνονται).
\item Το μέσο βρίσκεται με $\mathrm{ΑΜ}=\mathrm{ΜΒ}$, αλλά η μεσοκάθετος είναι η κάθετη από αυτό στο $\mathrm{ΑΒ}$.
\end{itemize}
\end{prosoxi}

\tmima{Λυμένα παραδείγματα}

\paradeigma{Ισαπέχω από τα δύο δέντρα}
\tixero Το παγκάκι $\mathrm{Μ}$ βρίσκεται πάνω στη μεσοκάθετο του $\mathrm{ΑΒ}$ και $\mathrm{ΜΑ}=5$ μ.\par
\tithelo Η απόσταση $\mathrm{ΜΒ}$.\par
\possindeo Κάθε σημείο της μεσοκαθέτου ισαπέχει από τα άκρα.\par
\apantisi $\mathrm{ΜΒ}=5$ μ.\par
\elegxos Αν ήταν $\mathrm{ΜΒ}\neq5$, το $\mathrm{Μ}$ δεν θα ήταν στη μεσοκάθετο.

\paradeigma{Κατασκευή μεσοκαθέτου}
\tixero Ένα τμήμα $\mathrm{ΑΒ}=6$ cm.\par
\tithelo Να φέρω τη μεσοκάθετο και το μέσο του.\par
\possindeo Γράφω δύο κύκλους ακτίνας $4$ cm (μεγαλύτερης από το $3$) με κέντρα $\mathrm{Α}$ και $\mathrm{Β}$. Οι δύο κύκλοι τέμνονται στα $\mathrm{Κ}$ και $\Lambda$. Η ευθεία $\mathrm{ΚΛ}$ τέμνει το $\mathrm{ΑΒ}$ στο μέσο $\mathrm{Μ}$.\par
\apantisi $\mathrm{ΑΜ}=\mathrm{ΜΒ}=3$ cm και $\mathrm{ΚΛ}\perp\mathrm{ΑΒ}$.\par
\elegxos Μετρώ: $\mathrm{ΑΜ}=3$ cm και η γωνία στο $\mathrm{Μ}$ είναι $90^\circ$.

\paradeigma{Σημείο που ισαπέχει}
\tixero Σημείο $\mathrm{Π}$ με $\mathrm{ΠΑ}=\mathrm{ΠΒ}$.\par
\tithelo Πού βρίσκεται το $\mathrm{Π}$.\par
\possindeo Αφού ισαπέχει από τα άκρα, βρίσκεται πάνω στη μεσοκάθετο του $\mathrm{ΑΒ}$.\par
\apantisi Πάνω στη μεσοκάθετο του $\mathrm{ΑΒ}$.\par
\elegxos Το μέσο του $\mathrm{ΑΒ}$ είναι το ένα τέτοιο σημείο, αλλά υπάρχουν άπειρα.

\paradeigma{Τρία δέντρα}
\tixero Τρία δέντρα $\mathrm{Α}$, $\mathrm{Β}$, $\Gamma$ που δεν είναι συνευθειακά.\par
\tithelo Πώς βρίσκω σημείο που ισαπέχει από τα τρία.\par
\possindeo Η μεσοκάθετος του $\mathrm{ΑΒ}$ περιέχει τα σημεία που ισαπέχουν από τα $\mathrm{Α}$, $\mathrm{Β}$. Η μεσοκάθετος του $\mathrm{ΒΓ}$ τα σημεία που ισαπέχουν από τα $\mathrm{Β}$, $\Gamma$. Το σημείο τομής τους $\mathrm{Ο}$ ισαπέχει και από τα τρία.\par
\apantisi Το σημείο τομής δύο μεσοκαθέτων.\par
\elegxos Τότε $\mathrm{ΟΑ}=\mathrm{ΟΒ}=\mathrm{ΟΓ}$ και ο κύκλος με κέντρο $\mathrm{Ο}$ περνά από τα τρία δέντρα.

\tmima{Ασκήσεις}

\zoni{Το κατάλαβα}{μεσοκάθετος και ισαπόσταση}

\ex Σχεδίασε τμήμα $\mathrm{ΑΒ}=6$ cm και φέρε τη μεσοκάθετο με κανόνα και διαβήτη. Μέτρησε το $\mathrm{ΑΜ}$ και την $\widehat{\mathrm{ΑΜΚ}}$.

\ex Πάνω στη μεσοκάθετο του $\mathrm{ΑΒ}$ πάρε τρία διαφορετικά σημεία και μέτρησε τις αποστάσεις τους από τα $\mathrm{Α}$ και $\mathrm{Β}$.

\ex Αν το $\mathrm{Μ}$ ανήκει στη μεσοκάθετο του $\mathrm{ΑΒ}$ και $\mathrm{ΜΑ}=5$, πόσο είναι το $\mathrm{ΜΒ}$;

\ex Αν $\mathrm{ΠΑ}=\mathrm{ΠΒ}$, πού βρίσκεται το $\mathrm{Π}$;

\ex Σχεδίασε τρίγωνο και φέρε τις μεσοκαθέτους των δύο πλευρών του. Σημείωσε το σημείο τομής τους και μέτρησε τις αποστάσεις του από τις τρεις κορυφές.

\ex Σωστό ή λάθος; α) Η μεσοκάθετος περνά από το μέσο του τμήματος. \quad β) Κάθε ευθεία που περνά από το μέσο είναι μεσοκάθετος. \quad γ) Η μεσοκάθετος είναι άξονας συμμετρίας του τμήματος. \quad δ) Κάθε σημείο εκτός της μεσοκαθέτου ισαπέχει από τα $\mathrm{Α}$, $\mathrm{Β}$.

\zoni{Το χρησιμοποιώ}{παγκάκια, πηγάδια, δρόμοι}

\ex \textbf{Το παγκάκι.} Τα δέντρα $\mathrm{Α}$ και $\mathrm{Β}$ απέχουν $8$ μ. Βρες δύο θέσεις για το παγκάκι ώστε να απέχει εξίσου από τα δύο δέντρα και ακριβώς $5$ μ. από το καθένα.

\ex \textbf{Το πηγάδι.} Δύο χωριά $\mathrm{Α}$ και $\mathrm{Β}$ θέλουν ένα πηγάδι που να ισαπέχει από αυτά και να είναι πάνω στο ευθύγραμμο δρόμο που περνά από το μέσο τους (κάθετος στο $\mathrm{ΑΒ}$). Πού θα το βάλουν;

\ex \textbf{Οι τρεις φίλοι.} Τρεις φίλοι μένουν σε διαφορετικά σημεία και θέλουν να συναντιούνται σε σημείο που να ισαπέχει από τα σπίτια τους. Πώς θα το βρουν;

\ex \textbf{Ο κύκλος.} Βρες το κέντρο ενός κύκλου με μία χορδή του $\mathrm{ΑΒ}$ και μία δεύτερη $\mathrm{ΒΓ}$ φέρνοντας μεσοκαθέτους.

\ex \textbf{Το φανάρι.} Ένα φανάρι πρέπει να τοποθετηθεί ώστε να ισαπέχει από δύο σχολεία $\mathrm{Α}$ και $\mathrm{Β}$ και να είναι στον ευθύγραμμο δρόμο $\varepsilon$ που τέμνει το $\mathrm{ΑΒ}$. Πώς βρίσκεται η θέση;

\zoni{Το σκέφτομαι}{η αρχαία κατασκευή}

\ex Γιατί στην κατασκευή της μεσοκαθέτου οι δύο κύκλοι πρέπει να έχουν ακτίνα μεγαλύτερη από το μισό του $\mathrm{ΑΒ}$;

\ex Γιατί η μεσοκάθετος είναι άξονας συμμετρίας του τμήματος; Τι είναι το συμμετρικό του $\mathrm{Α}$;

\ex Αν δύο σημεία ισαπέχουν από τα $\mathrm{Α}$ και $\mathrm{Β}$, πώς είναι η ευθεία που τα ενώνει ως προς το $\mathrm{ΑΒ}$;

\ex Πόσα σημεία έχουν απόσταση $5$ cm και από τα δύο άκρα ενός τμήματος $\mathrm{ΑΒ}=6$ cm;

\zoni{Για το σπίτι}{κατασκευές}

\ex Φέρε με κανόνα και διαβήτη τη μεσοκάθετο ενός τμήματος και έλεγξέ την δίπλωνοντας το χαρτί ώστε το $\mathrm{Α}$ να πέσει πάνω στο $\mathrm{Β}$.

\ex Βρες το κέντρο ενός κυκλικού αντικειμένου, π.χ. ενός δίσκου, χρησιμοποιώντας δύο χορδές.

\paragrafos{Συμμετρία ως προς σημείο: μισή στροφή}{Β.2.4, «Συμμετρία ως προς σημείο»}

\begin{anagke}
Στο πάρκο δίπλα στο σχολείο υπάρχει μια ρόδα λούνα παρκ. Η Αγγελική κοιτάζει δύο καμπίνες που βρίσκονται η μία απέναντι από την άλλη, και παρατηρεί: «Όταν η ρόδα κάνει μισή στροφή, η πρώτη καμπίνα πηγαίνει ακριβώς στη θέση της δεύτερης!». Ο Νίκος ρωτάει: «Και η δεύτερη στη θέση της πρώτης». Το κέντρο της ρόδας είναι ακριβώς στη μέση ανάμεσα στις δύο καμπίνες.

Η Πολυξένη τα αντιγράφει: «Αν ένα σημείο $\mathrm{Α}$ περιστραφεί γύρω από το κέντρο $\mathrm{Ο}$ κατά $180^\circ$, πηγαίνει στο σημείο $\mathrm{Α}'$ με το $\mathrm{Ο}$ στη μέση του $\mathrm{ΑΑ}'$». Η Μυρτώ το δοκιμάζει με ένα διαφανές χαρτί και μια καρφίτσα: καρφώνει στο $\mathrm{Ο}$, σχεδιάζει ένα τρίγωνο, γυρίζει το χαρτί μισή στροφή και σχεδιάζει το νέο. «Βγήκε ανάποδο, αλλά ίσο!».

Σε αυτή την παράγραφο ορίζουμε τη συμμετρία ως προς σημείο και βλέπουμε πώς ταιριάζει με τη συμμετρία ως προς άξονα. Θα δούμε και ότι οι αντίθετοι αριθμοί στον άξονα είναι συμμετρικά σημεία ως προς το $0$.
\end{anagke}

\begin{lexi}{συμμετρικό ως προς σημείο}{κέντρο συμμετρίας, συμμετρικό}
Το \lexiword{κέντρο συμμετρίας} είναι το σημείο γύρω από το οποίο περιστρέφουμε. Το \lexiword{συμμετρικό} του $\mathrm{Α}$ ως προς το $\mathrm{Ο}$ είναι το σημείο στην «απέναντι» θέση, σε ίση απόσταση από το $\mathrm{Ο}$ και στην ίδια ευθεία.
\end{lexi}

\begin{metafrasi}
Μισή στροφή γύρω από το $\mathrm{Ο}$. & στροφή $180^\circ$\\
Το $\mathrm{Ο}$ είναι το μέσο του $\mathrm{ΑΑ}'$. & $\mathrm{ΟΑ}=\mathrm{ΟΑ}'$ και $\mathrm{Α},\mathrm{Ο},\mathrm{Α}'$ συνευθειακά\\
Το συμμετρικό του $\mathrm{Ο}$ ως προς το $\mathrm{Ο}$. & το $\mathrm{Ο}$\\
Οι αντίθετοι αριθμοί. & συμμετρικά ως προς το $0$\\
\end{metafrasi}

\tmima{Η θεωρία}

\begin{orismos}{συμμετρικό σημείου ως προς σημείο}
\textbf{Συμμετρικό} του σημείου $\mathrm{Α}$ ως προς σημείο $\mathrm{Ο}$ είναι το σημείο $\mathrm{Α}'$ για το οποίο το $\mathrm{Ο}$ είναι το \textbf{μέσο} του τμήματος $\mathrm{ΑΑ}'$. Το $\mathrm{Ο}$ λέγεται \textbf{κέντρο συμμετρίας}. Το συμμετρικό του $\mathrm{Ο}$ είναι το ίδιο το $\mathrm{Ο}$.
\end{orismos}

Με άλλα λόγια, το $\mathrm{Α}'$ προκύπτει αν περιστρέψουμε το $\mathrm{Α}$ κατά $180^\circ$ γύρω από το $\mathrm{Ο}$. Για να το κατασκευάσω: ενώνω το $\mathrm{Α}$ με το $\mathrm{Ο}$, προεκτείνω την $\mathrm{ΑΟ}$ και με τον διαβήτη παίρνω $\mathrm{ΟΑ}'=\mathrm{ΟΑ}$.

\begin{center}
\begin{tikzpicture}[scale=1.1]
\fill (0,0) circle (2pt) node[above right] {$\mathrm{Ο}$};
\fill (-1.8,1.1) circle (2pt) node[above left] {$\mathrm{Α}$}; \fill (1.8,-1.1) circle (2pt) node[below right] {$\mathrm{Α}'$};
\draw[thick] (-1.8,1.1)--(1.8,-1.1);
\fill (-2.2,-0.4) circle (2pt) node[left] {$\mathrm{Β}$}; \fill (2.2,0.4) circle (2pt) node[right] {$\mathrm{Β}'$};
\draw[thick] (-2.2,-0.4)--(2.2,0.4);
\draw[gblue,thick] (-1.8,1.1)--(-2.2,-0.4); \draw[gblue,thick] (1.8,-1.1)--(2.2,0.4);
\end{tikzpicture}
\end{center}

\begin{kanonas}{συμμετρία ως προς σημείο}
Δύο σχήματα είναι \textbf{συμμετρικά ως προς σημείο} $\mathrm{Ο}$ αν το καθένα αποτελείται από τα συμμετρικά των σημείων του άλλου. Τα συμμετρικά ως προς σημείο σχήματα είναι \textbf{ίσα}. Το συμμετρικό τμήματος $\mathrm{ΑΒ}$ είναι τμήμα $\mathrm{Α}'\mathrm{Β}'$ ίσο και \textbf{παράλληλο} προς το $\mathrm{ΑΒ}$.
\end{kanonas}

Στον άξονα των ρητών, το συμμετρικό του αριθμού $\alpha$ ως προς το $0$ είναι ο $-\alpha$, ο αντίθετός του. Ως προς το $1$: το συμμετρικό του $3{,}5$ είναι το $x$ με $\frac{3{,}5+x}{2}=1$, δηλαδή $x=-1{,}5$.

\begin{prosoxi}
\begin{itemize}[leftmargin=1.2em,itemsep=0.3ex,topsep=2pt]
\item Στη συμμετρία ως προς σημείο, τα $\mathrm{Α}$, $\mathrm{Ο}$, $\mathrm{Α}'$ είναι πάνω στην ίδια ευθεία (στη συμμετρία ως προς άξονα το $\mathrm{ΑΑ}'$ είναι κάθετο στον άξονα).
\item Το $\mathrm{Ο}$ είναι το μέσο του $\mathrm{ΑΑ}'$: $\mathrm{ΑΑ}'=2\mathrm{ΟΑ}$.
\item Δεν είναι ίδιο με τη συμμετρία ως προς άξονα: μια στροφή $180^\circ$ δεν «αναποδογυρίζει» το σχήμα σαν καθρέφτης.
\item Το συμμετρικό τμήματος είναι \emph{παράλληλο} στο αρχικό (ή στην ίδια ευθεία, αν η ευθεία περνά από το $\mathrm{Ο}$).
\end{itemize}
\end{prosoxi}

\tmima{Λυμένα παραδείγματα}

\paradeigma{Η ρόδα}
\tixero Η καμπίνα $\mathrm{Α}$ απέχει $12$ μ. από το κέντρο $\mathrm{Ο}$ της ρόδας.\par
\tithelo Πού είναι η συμμετρική καμπίνα και πόσο απέχουν οι δύο.\par
\possindeo Το $\mathrm{Α}'$ είναι στην προέκταση της $\mathrm{ΑΟ}$, με $\mathrm{ΟΑ}'=12$ μ. Το $\mathrm{ΑΑ}'=2\cdot12=24$ μ. (διάμετρος της ρόδας).\par
\apantisi $\mathrm{ΟΑ}'=12$ μ. και $\mathrm{ΑΑ}'=24$ μ.\par
\elegxos Το $\mathrm{Ο}$ είναι το μέσο του $\mathrm{ΑΑ}'$.

\paradeigma{Στον άξονα των αριθμών}
\tixero Το σημείο $\mathrm{Α}(3{,}5)$ του άξονα.\par
\tithelo Το συμμετρικό του ως προς το $0$ και ως προς το $1$.\par
\possindeo Ως προς το $0$: ο αντίθετος $-3{,}5$. Ως προς το $1$: το $1$ είναι το μέσο, $\frac{3{,}5+x}{2}=1$, $x=-1{,}5$.\par
\apantisi $-3{,}5$ και $-1{,}5$.\par
\elegxos Η απόσταση του $3{,}5$ από το $1$ είναι $2{,}5$ και του $-1{,}5$ από το $1$ επίσης $2{,}5$.

\paradeigma{Συμμετρικό τμήματος}
\tixero Τμήμα $\mathrm{ΑΒ}=4$ cm και το συμμετρικό του $\mathrm{Α}'\mathrm{Β}'$ ως προς $\mathrm{Ο}$.\par
\tithelo Το μήκος του και η θέση του.\par
\possindeo Τα συμμετρικά σχήματα είναι ίσα, άρα $\mathrm{Α}'\mathrm{Β}'=4$ cm, και το $\mathrm{Α}'\mathrm{Β}'$ είναι παράλληλο προς το $\mathrm{ΑΒ}$.\par
\apantisi $\mathrm{Α}'\mathrm{Β}'=4$ cm και $\mathrm{Α}'\mathrm{Β}'\parallel\mathrm{ΑΒ}$.\par
\elegxos Το $\mathrm{ΑΒΒ}'\mathrm{Α}'$ έχει τα $\mathrm{Α}\mathrm{Α}'$ και $\mathrm{ΒΒ}'$ να περνούν από το $\mathrm{Ο}$.

\paradeigma{Κατασκευή}
\tixero Το τρίγωνο $\mathrm{ΑΒΓ}$ και το σημείο $\mathrm{Ο}$ εκτός αυτού.\par
\tithelo Να κατασκευάσω το συμμετρικό του.\par
\possindeo Βρίσκω τα $\mathrm{Α}'$, $\mathrm{Β}'$, $\Gamma'$ (το καθένα με $\mathrm{ΟΑ}'=\mathrm{ΟΑ}$ κ.λπ. στην προέκταση) και ενώνω.\par
\apantisi Το τρίγωνο $\mathrm{Α}'\mathrm{Β}'\Gamma'$ είναι ίσο με το $\mathrm{ΑΒΓ}$.\par
\elegxos Με διαφανές χαρτί και στροφή $180^\circ$ γύρω από το $\mathrm{Ο}$ το $\mathrm{ΑΒΓ}$ πέφτει πάνω στο $\mathrm{Α}'\mathrm{Β}'\Gamma'$.

\tmima{Ασκήσεις}

\zoni{Το κατάλαβα}{σημεία, τμήματα, σχήματα}

\ex Σχεδίασε σημείο $\mathrm{Ο}$ και σημείο $\mathrm{Α}$ και βρες το συμμετρικό του ως προς $\mathrm{Ο}$ με διαβήτη.

\ex Αν $\mathrm{ΟΑ}=3{,}2$ cm, πόσο είναι το $\mathrm{ΟΑ}'$ και πόσο το $\mathrm{ΑΑ}'$;

\ex Σχεδίασε τμήμα $\mathrm{ΑΒ}$ και το συμμετρικό του ως προς σημείο $\mathrm{Ο}$ εκτός του. Τι σχέση έχουν;

\ex Σχεδίασε τρίγωνο $\mathrm{ΑΒΓ}$ και το συμμετρικό του ως προς ένα σημείο $\mathrm{Ο}$.

\ex Βρες το συμμετρικό των σημείων του άξονα $\mathrm{Α}(3{,}5)$, $\mathrm{Β}(-2)$, $\Gamma(0)$ ως προς την αρχή $0$.

\ex Βρες το συμμετρικό του $3{,}5$ ως προς το $1$ και του $-2$ ως προς το $3$.

\ex Σωστό ή λάθος; α) Στη συμμετρία ως προς σημείο, τα $\mathrm{Α}$, $\mathrm{Ο}$, $\mathrm{Α}'$ είναι συνευθειακά. \quad β) Το $\mathrm{Ο}$ είναι το μέσο του $\mathrm{ΑΑ}'$. \quad γ) Τα συμμετρικά σχήματα ως προς σημείο είναι ίσα. \quad δ) Το συμμετρικό τμήματος είναι κάθετο σε αυτό.

\zoni{Το χρησιμοποιώ}{ρόδες, ρολόγια, παρτέρια}

\ex \textbf{Η ρόδα.} Μια καμπίνα απέχει $12$ μ. από το κέντρο της ρόδας. Πού είναι η καμπίνα μετά από μισή στροφή; Πόσο απέχουν οι δύο θέσεις;

\ex \textbf{Τα δύο παρτέρια.} Δύο ίδια παρτέρια του κήπου είναι συμμετρικά ως προς το συντριβάνι $\mathrm{Ο}$. Ένα φυτό στο πρώτο απέχει $2{,}5$ μ. από το $\mathrm{Ο}$. Πού είναι το αντίστοιχο στο δεύτερο;

\ex \textbf{Το ρολόι.} Ο λεπτοδείκτης στις $12{:}00$ δείχνει προς τα πάνω. Πού δείχνει μετά από $30$ λεπτά (μισή στροφή); Τι σχέση έχουν οι δύο θέσεις της άκρης του;

\ex \textbf{Ο ανεμόμυλος.} Ένας ανεμόμυλος με δύο αντίθετα πτερύγια περιστρέφεται. Πού βρίσκεται το ένα πτερύγιο όταν το άλλο έχει μετακινηθεί στη θέση του αντιθέτου του;

\ex \textbf{Το χαρτί.} Κάρφωσε με καρφίτσα ένα διαφανές χαρτί, σχεδίασε μια λέξη και γύρνα το $180^\circ$. Πώς φαίνεται η λέξη;

\zoni{Το σκέφτομαι}{ίδιο ή διαφορετικό από τον καθρέφτη;}

\ex Σε τι διαφέρει η συμμετρία ως προς σημείο από τη συμμετρία ως προς άξονα στο συμμετρικό ενός τριγώνου;

\ex Αν το $\mathrm{Α}'$ είναι συμμετρικό του $\mathrm{Α}$ ως προς $\mathrm{Ο}$, ποιο είναι το συμμετρικό του $\mathrm{Α}'$ ως προς το ίδιο $\mathrm{Ο}$;

\ex Εξήγησε γιατί το συμμετρικό ενός τμήματος ως προς σημείο είναι παράλληλο προς αυτό. (Βοήθεια: ένωσε τα $\mathrm{Α}$, $\mathrm{Β}'$ και δες το $\mathrm{ΑΒΑ}'\mathrm{Β}'$.)

\ex Ποια είναι η ευθεία που είναι συμμετρική του εαυτού της ως προς σημείο $\mathrm{Ο}$;

\zoni{Για το σπίτι}{στροφή μισή}

\ex Σχεδίασε ένα σχήμα σε διαφανές χαρτί, κάρφωσε ένα σημείο και γύρνα μισή στροφή. Σχεδίασε το συμμετρικό.

\ex Βρες στο σπίτι δύο αντικείμενα που ταιριάζουν αν περιστραφούν $180^\circ$.

\paragrafos{Κέντρο συμμετρίας σχήματος}{Β.2.5, «Κέντρο συμμετρίας»}

\begin{anagke}
Η Ρήνη βρίσκει στο συρτάρι μια τράπουλα. «Κοίτα το φύλλο με το καρό: αν το γυρίσω ανάποδα, φαίνεται ακριβώς το ίδιο!». Ο Βασίλης δοκιμάζει με το γράμμα Ν της λέξης «Νίκος»: το γυρίζει μισή στροφή και πάλι Ν. Με το γράμμα Γ όμως δεν συμβαίνει το ίδιο.

Η Αγγελική ρωτάει: «Υπάρχει ένα σημείο, γύρω από το οποίο όταν γυρίζουμε το σχήμα μισή στροφή, αυτό πέφτει πάνω στον εαυτό του;». Η Πολυξένη απαντά: «Όταν υπάρχει, το λέμε κέντρο συμμετρίας του σχήματος».

Σε αυτή την παράγραφο ψάχνουμε ποια σχήματα έχουν κέντρο συμμετρίας και πού βρίσκεται.
\end{anagke}

\begin{lexi}{κέντρο συμμετρίας σχήματος}{κέντρο, συμμετρία}
Ένα σχήμα έχει \lexiword{κέντρο συμμετρίας} αν υπάρχει σημείο $\mathrm{Ο}$ τέτοιο ώστε, όταν το περιστρέψουμε κατά $180^\circ$ γύρω από αυτό, το σχήμα να συμπέσει με τον εαυτό του.
\end{lexi}

\begin{metafrasi}
Μισή στροφή και το σχήμα φαίνεται το ίδιο. & το σχήμα έχει κέντρο συμμετρίας\\
Το συμμετρικό του σχήματος ως προς $\mathrm{Ο}$ είναι το ίδιο το σχήμα. & το $\mathrm{Ο}$ είναι κέντρο συμμετρίας\\
Κάθε σημείο έχει «ταίρι» απέναντι. & για κάθε $\mathrm{Α}$ υπάρχει $\mathrm{Α}'$ στο σχήμα με μέσο το $\mathrm{Ο}$\\
\end{metafrasi}

\tmima{Η θεωρία}

\begin{orismos}{κέντρο συμμετρίας σχήματος}
\textbf{Κέντρο συμμετρίας} σχήματος είναι ένα σημείο $\mathrm{Ο}$ γύρω από το οποίο, αν περιστραφεί το σχήμα κατά $180^\circ$, συμπίπτει με το αρχικό. Τότε το συμμετρικό του σχήματος ως προς το $\mathrm{Ο}$ είναι το ίδιο το σχήμα.
\end{orismos}

Παραδείγματα: το \textbf{παραλληλόγραμμο} έχει κέντρο συμμετρίας το σημείο τομής των διαγωνίων του, ο \textbf{κύκλος} το κέντρο του, το \textbf{ευθύγραμμο τμήμα} το μέσο του. Το \textbf{ισόπλευρο τρίγωνο} και το \textbf{ισοσκελές} δεν έχουν κέντρο συμμετρίας.

\begin{kanonas}{συμμετρική ευθεία}
Το συμμετρικό μιας ευθείας $\varepsilon$ ως προς σημείο $\mathrm{Ο}$ είναι ευθεία $\varepsilon'$ \textbf{παράλληλη} προς την $\varepsilon$. Αν η $\varepsilon$ περνά από το $\mathrm{Ο}$, τότε $\varepsilon'=\varepsilon$.
\end{kanonas}

\begin{center}\small
\begin{tabular}{lcc}
\toprule
Σχήμα & Άξονες συμμετρίας & Κέντρο συμμετρίας\\
\midrule
Ευθύγραμμο τμήμα & $2$ & ναι (μέσο)\\
Ισοσκελές τρίγωνο & $1$ & όχι\\
Ισόπλευρο τρίγωνο & $3$ & όχι\\
Παραλληλόγραμμο & κανένας & ναι\\
Ορθογώνιο & $2$ & ναι\\
Ρόμβος & $2$ & ναι\\
Τετράγωνο & $4$ & ναι\\
Κύκλος & άπειροι & ναι (κέντρο)\\
\bottomrule
\end{tabular}
\end{center}

\begin{prosoxi}
\begin{itemize}[leftmargin=1.2em,itemsep=0.3ex,topsep=2pt]
\item Άξονας συμμετρίας και κέντρο συμμετρίας είναι \emph{ανεξάρτητα}: υπάρχουν σχήματα με το ένα και χωρίς το άλλο (το ισοσκελές τρίγωνο έχει άξονα αλλά όχι κέντρο, το παραλληλόγραμμο κέντρο αλλά όχι άξονα).
\item Ένα σχήμα έχει το πολύ ένα κέντρο συμμετρίας (εκτός από τα άπειρα σχήματα, όπως η ευθεία).
\item Το κέντρο συμμετρίας δεν χρειάζεται να ανήκει στο σχήμα (π.χ. στο δαχτυλίδι).
\end{itemize}
\end{prosoxi}

\tmima{Λυμένα παραδείγματα}

\paradeigma{Το γράμμα}
\tixero Τα γράμματα Ο, Ζ, Μ, Η, Ε.\par
\tithelo Ποια έχουν κέντρο συμμετρίας.\par
\possindeo Κάνω στροφή $180^\circ$ και συγκρίνω: το Ο και το Η συμπίπτουν· το Ζ γυρισμένο φαίνεται Ζ (ναι)· το Μ φαίνεται ανάποδο, ως W (όχι)· το Ε γίνεται «καθρεφτισμένο» (όχι).\par
\apantisi Κέντρο συμμετρίας έχουν τα Ο, Η, Ζ.\par
\elegxos Κάθε σημείο της μισής μορφής έχει ταίρι στην άλλη μισή.

\paradeigma{Το ορθογώνιο}
\tixero Ορθογώνιο $\mathrm{ΑΒΓΔ}$ με διαστάσεις $8$ cm και $6$ cm.\par
\tithelo Το κέντρο συμμετρίας και η απόστασή του από κάθε κορυφή.\par
\possindeo Το κέντρο είναι το σημείο τομής των διαγωνίων. Η διαγώνιος είναι $\sqrt{8^2+6^2}=10$ cm και το κέντρο είναι το μέσο της, άρα απέχει $5$ cm από κάθε κορυφή.\par
\apantisi $\mathrm{ΟΑ}=\mathrm{ΟΒ}=\mathrm{ΟΓ}=\mathrm{ΟΔ}=5$ cm.\par
\elegxos Τα συμμετρικά ως προς $\mathrm{Ο}$ ζεύγη κορυφών ($\mathrm{Α},\mathrm{Γ}$) και ($\mathrm{Β},\mathrm{Δ}$) έχουν το $\mathrm{Ο}$ μέσο.

\paradeigma{Η ευθεία}
\tixero Η ευθεία $\varepsilon$ και σημείο $\mathrm{Ο}$ εκτός αυτής.\par
\tithelo Η συμμετρική ευθεία.\par
\possindeo Παίρνω δύο σημεία $\mathrm{Α},\mathrm{Β}$ της $\varepsilon$, βρίσκω τα $\mathrm{Α}',\mathrm{Β}'$ και φέρνω την ευθεία από αυτά.\par
\apantisi Η $\varepsilon'$ είναι παράλληλη προς την $\varepsilon$.\par
\elegxos Αν το $\mathrm{Ο}$ ήταν πάνω στην $\varepsilon$, η $\varepsilon'$ θα ταυτιζόταν με την $\varepsilon$.

\tmima{Ασκήσεις}

\zoni{Το κατάλαβα}{ποιος έχει κέντρο}

\ex Ποια από τα κεφαλαία γράμματα Α, Η, Ι, Λ, Ν, Ο, Σ, Χ έχουν κέντρο συμμετρίας;

\ex Έχει κέντρο συμμετρίας το ισοσκελές τρίγωνο; Και το ορθογώνιο;

\ex Πόσους άξονες και πόσα κέντρα συμμετρίας έχει ο κύκλος;

\ex Βρες το κέντρο συμμετρίας ενός ρόμβου με διαγώνιους $6$ cm και $8$ cm και πόσο απέχει από κάθε κορυφή του.

\ex Σχεδίασε παραλληλόγραμμο και βρες το κέντρο συμμετρίας του.

\ex Σχεδίασε ευθεία $\varepsilon$ και σημείο $\mathrm{Ο}$ εκτός της. Σχεδίασε τη συμμετρική της.

\ex Σωστό ή λάθος; α) Κάθε σχήμα με άξονα συμμετρίας έχει και κέντρο συμμετρίας. \quad β) Το τετράγωνο έχει κέντρο συμμετρίας. \quad γ) Ο κύκλος έχει άπειρους άξονες συμμετρίας.

\zoni{Το χρησιμοποιώ}{αντικείμενα γύρω μας}

\ex \textbf{Η τράπουλα.} Ο ρόμβος που είναι τυπωμένος στο καρό φύλλο, αν γυρίσεις το φύλλο ανάποδα, μοιάζει ίδιος. Τι σημαίνει αυτό για το κέντρο του ρόμβου;

\ex \textbf{Το πλακάκι.} Ένα ορθογώνιο πλακάκι $30\times20$ cm έχει κέντρο συμμετρίας. Πού βρίσκεται και πόσο απέχει από τις πλευρές του;

\ex \textbf{Η ρόδα.} Μια ρόδα με $8$ ίσα ακτινωτά κομμάτια έχει κέντρο συμμετρίας το κέντρο της. Πόσα ζεύγη «ταίρι» κομμάτια υπάρχουν;

\ex \textbf{Το σήμα κυκλοφορίας.} Το σήμα «STOP» και το «Απαγορεύεται η είσοδος» έχουν κέντρο συμμετρίας; Εξήγησε.

\ex \textbf{Το ταβλί.} Ένα ταμπλό ταβλιού έχει δύο όμοια μισά. Σχεδίασε το κέντρο συμμετρίας.

\zoni{Το σκέφτομαι}{συνδυασμοί συμμετριών}

\ex Βρες σχήμα με κέντρο αλλά χωρίς άξονα συμμετρίας.

\ex Βρες σχήμα με άξονα αλλά χωρίς κέντρο συμμετρίας.

\ex Αν ένα σχήμα έχει δύο κάθετους άξονες συμμετρίας, έχει κέντρο συμμετρίας; Πού;

\ex Εξήγησε γιατί το συμμετρικό ευθείας ως προς σημείο εκτός της είναι παράλληλο. (Βοήθεια: κάθετο τμήμα $\mathrm{ΟΑ}$ από το $\mathrm{Ο}$ προς την $\varepsilon$.)

\zoni{Για το σπίτι}{κέντρα παντού}

\ex Βρες πέντε αντικείμενα στο σπίτι με κέντρο συμμετρίας.

\ex Φτιάξε ένα σχήμα με κέντρο συμμετρίας από διπλωμένο χαρτί.

\paragrafos{Παράλληλες ευθείες και τέμνουσα}{Β.2.6, «Παράλληλες ευθείες που τέμνονται από μια άλλη ευθεία»}

\begin{anagke}
Ο δρόμος που περνά μπροστά από το σχολείο έχει δύο παράλληλα πεζοδρόμια. Ένα σοκάκι τον διασχίζει λοξά. Ο Νίκος παρατηρεί ότι στο ένα πεζοδρόμιο η γωνία μεταξύ σοκακιού και δρόμου είναι πολύ μυτερή, και ρωτάει: «Στο απέναντι πεζοδρόμιο θα είναι ίδια;». Ο Βασίλης αποφασίζει να τη μετρήσει: $50^\circ$ και στα δύο.

Η Αγγελική σχεδιάζει δύο παράλληλες ευθείες και μία τρίτη που τις τέμνει, και βρίσκει $8$ γωνίες. Η Πολυξένη λέει: «Αν ξέρω μία, ξέρω όλες!». Αλλά για να το πει με ακρίβεια, χρειαζόμαστε ονόματα για τις γωνίες: «εντός», «εκτός», «εναλλάξ», «επί τα αυτά».
\end{anagke}

\begin{lexi}{εντός, εκτός, εναλλάξ, επί τα αυτά}{θέσεις γωνιών}
Εντός σημαίνει «μέσα» στη λωρίδα των παραλλήλων, εκτός «έξω». Οι γωνίες \lexiword{εναλλάξ} είναι από αντίθετες πλευρές της τέμνουσας, ενώ οι γωνίες \lexiword{επί τα αυτά} είναι από την ίδια πλευρά της (το «επί τα αυτά μέρη» = προς τα ίδια μέρη).
\end{lexi}

\begin{metafrasi}
Στη λωρίδα ανάμεσα στις παράλληλες. & εντός\\
Έξω από τη λωρίδα. & εκτός\\
Από αντίθετες πλευρές της τέμνουσας. & εναλλάξ\\
Από την ίδια πλευρά της τέμνουσας. & επί τα αυτά\\
\end{metafrasi}

\tmima{Η θεωρία}

Έστω ευθείες $\varepsilon_1\parallel\varepsilon_2$ και μια \textbf{τέμνουσα} $\delta$ που τις τέμνει στα $\mathrm{Α}$ και $\mathrm{Β}$. Σχηματίζονται $8$ γωνίες· ονομάζουμε $\hat{\mathrm{Α}}_1,\hat{\mathrm{Α}}_2,\hat{\mathrm{Α}}_3,\hat{\mathrm{Α}}_4$ γύρω από το $\mathrm{Α}$ και $\hat{\mathrm{Β}}_1,\ldots,\hat{\mathrm{Β}}_4$ γύρω από το $\mathrm{Β}$.

\begin{center}
\begin{tikzpicture}[scale=1.0]
\draw[thick] (-2.6,1.2)--(2.6,1.2) node[right] {$\varepsilon_1$};
\draw[thick] (-2.6,-1.2)--(2.6,-1.2) node[right] {$\varepsilon_2$};
\draw[thick,gblue] (-1.2,2.4)--(1.2,-2.4) node[below right] {$\delta$};
\fill (-0.6,1.2) circle (2pt) node[above left] {$\mathrm{Α}$};
\fill (0.6,-1.2) circle (2pt) node[below right] {$\mathrm{Β}$};
\node at (-1.0,1.55) {\small $1$}; \node at (-0.2,1.55) {\small $2$};
\node at (-0.2,0.85) {\small $3$}; \node at (-1.0,0.85) {\small $4$};
\node at (0.2,-0.85) {\small $1$}; \node at (1.0,-0.85) {\small $2$};
\node at (1.0,-1.55) {\small $3$}; \node at (0.2,-1.55) {\small $4$};
\end{tikzpicture}
\end{center}

\begin{orismos}{ονομασίες γωνιών}
Οι γωνίες $\hat{\mathrm{Α}}_3,\hat{\mathrm{Α}}_4,\hat{\mathrm{Β}}_1,\hat{\mathrm{Β}}_2$ είναι \textbf{εντός}, οι υπόλοιπες \textbf{εκτός}. Δύο γωνίες είναι \textbf{εναλλάξ} αν βρίσκονται σε διαφορετικά ημιεπίπεδα της $\delta$ και \textbf{επί τα αυτά} αν βρίσκονται στο ίδιο. Προκύπτουν $6$ ονομασίες για τα $16$ ζεύγη γωνιών: εντός εναλλάξ, εκτός εναλλάξ, εντός και επί τα αυτά, εκτός και επί τα αυτά, εντός--εκτός εναλλάξ, εντός--εκτός επί τα αυτά.
\end{orismos}

\begin{kanonas}{παράλληλες και τέμνουσα}
Αν $\varepsilon_1\parallel\varepsilon_2$, τότε:
\begin{itemize}[leftmargin=1.2em,itemsep=0.3ex,topsep=2pt]
\item οι \textbf{εντός εναλλάξ} είναι ίσες,
\item οι \textbf{εκτός εναλλάξ} είναι ίσες,
\item οι \textbf{εντός--εκτός επί τα αυτά} είναι ίσες,
\item οι \textbf{εντός επί τα αυτά} είναι παραπληρωματικές ($\text{άθροισμα}=180^\circ$),
\item οι \textbf{εκτός επί τα αυτά} είναι παραπληρωματικές.
\end{itemize}
Δηλαδή όλες οι οξείες γωνίες του σχήματος είναι ίσες ($\omega$), όλες οι αμβλείες ίσες ($\varphi$) και $\omega+\varphi=180^\circ$.
\end{kanonas}

Ισχύει και το \textbf{αντίστροφο}: αν δύο ευθείες που τέμνονται από τρίτη σχηματίζουν εντός εναλλάξ γωνίες ίσες (ή εντός επί τα αυτά παραπληρωματικές, ή εντός--εκτός επί τα αυτά ίσες), τότε είναι παράλληλες. Έτσι ελέγχουμε αν δύο ευθείες είναι παράλληλες.

\begin{prosoxi}
\begin{itemize}[leftmargin=1.2em,itemsep=0.3ex,topsep=2pt]
\item Οι ιδιότητες ισχύουν \emph{μόνο όταν} οι ευθείες είναι παράλληλες.
\item Οι εντός επί τα αυτά είναι παραπληρωματικές (όχι ίσες).
\item Οι κατακορυφήν γωνίες είναι ίσες και οι διαδοχικές στο ίδιο σημείο παραπληρωματικές: $\hat{\mathrm{Α}}_1+\hat{\mathrm{Α}}_2=180^\circ$.
\end{itemize}
\end{prosoxi}

\tmima{Λυμένα παραδείγματα}

\paradeigma{Όλες οι γωνίες}
\tixero $\varepsilon_1\parallel\varepsilon_2$ και $\hat{\mathrm{Α}}_1=40^\circ$.\par
\tithelo Όλες οι άλλες γωνίες.\par
\possindeo Οι $\hat{\mathrm{Α}}_1$ και $\hat{\mathrm{Α}}_3$ είναι κατακορυφήν, άρα $\hat{\mathrm{Α}}_3=40^\circ$. Η $\hat{\mathrm{Α}}_2=180^\circ-40^\circ=140^\circ$ και η $\hat{\mathrm{Α}}_4=140^\circ$. Στο $\mathrm{Β}$ οι γωνίες είναι ίσες με τις αντίστοιχες του $\mathrm{Α}$ (εντός--εκτός επί τα αυτά).\par
\apantisi Σε κάθε σημείο: δύο γωνίες $40^\circ$ και δύο $140^\circ$.\par
\elegxos $40^\circ+140^\circ=180^\circ$ και το άθροισμα των γωνιών γύρω από σημείο είναι $2\cdot40^\circ+2\cdot140^\circ=360^\circ$.

\paradeigma{Εντός επί τα αυτά}
\tixero $\varepsilon_1\parallel\varepsilon_2$ και μία εντός επί τα αυτά γωνία είναι $65^\circ$.\par
\tithelo Η άλλη.\par
\possindeo Οι εντός επί τα αυτά είναι παραπληρωματικές: $x=180^\circ-65^\circ=115^\circ$.\par
\apantisi $115^\circ$.\par
\elegxos $65^\circ+115^\circ=180^\circ$.

\paradeigma{Παράλληλες ή όχι}
\tixero Δύο ευθείες που τέμνονται από τρίτη σχηματίζουν εντός εναλλάξ γωνίες $72^\circ$ και $75^\circ$.\par
\tithelo Είναι παράλληλες;\par
\possindeo Για να είναι παράλληλες θα έπρεπε οι εντός εναλλάξ να είναι ίσες.\par
\apantisi Όχι, γιατί $72^\circ\neq75^\circ$.\par
\elegxos Αν ήταν $72^\circ$ και $72^\circ$, θα ήταν παράλληλες.

\tmima{Ασκήσεις}

\zoni{Το κατάλαβα}{ονόματα και ισότητες}

\ex Αν $\varepsilon_1\parallel\varepsilon_2$ και μια οξεία γωνία είναι $55^\circ$, πόσες από τις $8$ γωνίες είναι $55^\circ$ και πόσες $125^\circ$;

\ex Αν δύο εντός εναλλάξ γωνίες είναι $x$ και $2x-30^\circ$ και οι ευθείες είναι παράλληλες, βρες το $x$.

\ex Δύο εντός επί τα αυτά γωνίες παραλλήλων είναι $x$ και $x+40^\circ$. Βρες τες.

\ex Δύο εκτός εναλλάξ γωνίες είναι $3x$ και $x+50^\circ$ και οι ευθείες παράλληλες. Βρες το $x$.

\ex Για κάθε ζεύγος γωνιών, πες πώς ονομάζεται: α) $\hat{\mathrm{Α}}_3,\hat{\mathrm{Β}}_1$ β) $\hat{\mathrm{Α}}_3,\hat{\mathrm{Β}}_2$ γ) $\hat{\mathrm{Α}}_1,\hat{\mathrm{Β}}_3$ δ) $\hat{\mathrm{Α}}_1,\hat{\mathrm{Β}}_1$.

\ex Σχεδίασε δύο παράλληλες που απέχουν $4$ cm και μια τέμνουσα που σχηματίζει $72^\circ$ με τη μία. Υπολόγισε τις υπόλοιπες γωνίες.

\ex Σωστό ή λάθος; α) Οι εντός εναλλάξ γωνίες είναι πάντα ίσες. \quad β) Αν δύο ευθείες έχουν εντός εναλλάξ γωνίες ίσες, είναι παράλληλες. \quad γ) Οι εντός επί τα αυτά παραλλήλων είναι ίσες.

\zoni{Το χρησιμοποιώ}{δρόμοι και σκάλες}

\ex \textbf{Το σοκάκι.} Στο πεζοδρόμιο η γωνία με το σοκάκι είναι $50^\circ$. Πόση είναι στο απέναντι πεζοδρόμιο (εντός εναλλάξ) και πόση η γωνία που είναι εντός επί τα αυτά με αυτήν;

\ex \textbf{Η σκάλα.} Μια σκάλα ακουμπά σε δύο παράλληλα ράφια. Η γωνία της με το πρώτο ράφι είναι $70^\circ$. Πόση είναι με το δεύτερο, από την ίδια πλευρά;

\ex \textbf{Οι ράγες.} Δύο παράλληλες σιδηροτροχιές τέμνονται από μια διάβαση. Αν μια γωνία είναι $110^\circ$, βρες όλες τις άλλες.

\ex \textbf{Ο φράκτης.} Δύο πασσάλοι κεκλιμένοι προς το έδαφος σχηματίζουν γωνίες $64^\circ$ και $64^\circ$ (εντός εναλλάξ) με το σχοινί. Είναι παράλληλοι; Γιατί;

\ex \textbf{Το ράφι.} Δύο ράφια τέμνονται από κατακόρυφο στύλο και σχηματίζουν εντός επί τα αυτά γωνίες $90^\circ$ και $90^\circ$. Τι ισχύει για τα ράφια;

\zoni{Το σκέφτομαι}{αποδείξεις}

\ex Εξήγησε γιατί αν οι εντός εναλλάξ είναι ίσες τότε και οι εντός επί τα αυτά είναι παραπληρωματικές.

\ex Στο τρίγωνο $\mathrm{ΑΒΓ}$ φέρε από το $\mathrm{Α}$ παράλληλη προς τη $\mathrm{ΒΓ}$. Τι ισχύει για τις γωνίες που σχηματίζονται;

\ex Αν $\varepsilon_1\parallel\varepsilon_2$ και $\varepsilon_3\parallel\varepsilon_4$, και μια γωνία τους είναι $70^\circ$, βρες όλες τις γωνίες του «παραλληλογράμμου» που σχηματίζεται.

\ex Αν η $\mathrm{Βδ}$ είναι διχοτόμος της γωνίας $\hat{\mathrm{Β}}=56^\circ$ και $\varepsilon_1\parallel\varepsilon_2$, βρες τη γωνία που σχηματίζει η διχοτόμος με την $\varepsilon_1$.

\zoni{Για το σπίτι}{ρίξε μια ματιά}

\ex Βρες στο σπίτι δύο παράλληλες ευθείες τεμνόμενες από τρίτη και μέτρησε τις γωνίες.

\ex Σχεδίασε τετράπλευρο $\mathrm{ΑΒΓΔ}$ με $\mathrm{ΑΒ}\parallel\mathrm{ΔΓ}$, $\mathrm{ΑΔ}\parallel\mathrm{ΒΓ}$ και $\hat{\mathrm{Α}}=63^\circ$. Βρες τις άλλες γωνίες.

\epanalipsi{Β2}

\tmima{Λεξικό του κεφαλαίου}

\renewcommand{\arraystretch}{1.5}
\begin{tabular*}{\linewidth}{@{\extracolsep{\fill}} >{\raggedright\arraybackslash}p{0.28\linewidth} >{\raggedright\arraybackslash}p{0.62\linewidth}}
\toprule
\textbf{Λέξη ή σύμβολο} & \textbf{Με μία πρόταση}\\
\midrule
συμμετρικό ως προς άξονα $\xi$ & Το $\mathrm{Α}'$ με τον $\xi$ μεσοκάθετο του $\mathrm{ΑΑ}'$. Τα συμμετρικά σχήματα είναι ίσα.\\
άξονας συμμετρίας & Ευθεία που χωρίζει το σχήμα σε δύο σχήματα που καθρεφτίζονται.\\
μεσοκάθετος & Η κάθετη στο μέσο ενός τμήματος· κάθε σημείο της ισαπέχει από τα άκρα.\\
συμμετρικό ως προς σημείο $\mathrm{Ο}$ & Το $\mathrm{Α}'$ με το $\mathrm{Ο}$ μέσο του $\mathrm{ΑΑ}'$ (στροφή $180^\circ$).\\
κέντρο συμμετρίας & Σημείο $\mathrm{Ο}$ γύρω από το οποίο το σχήμα, με στροφή $180^\circ$, πέφτει πάνω στον εαυτό του.\\
τέμνουσα & Ευθεία που τέμνει δύο άλλες ευθείες.\\
εντός εναλλάξ & Ίσες γωνίες, όταν οι ευθείες είναι παράλληλες.\\
εντός επί τα αυτά & Παραπληρωματικές γωνίες, όταν οι ευθείες είναι παράλληλες.\\
\bottomrule
\end{tabular*}

\tmima{Ασκήσεις επανάληψης}

\ex Βρες το συμμετρικό του σημείου $\mathrm{Α}$, που απέχει $3{,}4$ cm από τον άξονα $\xi$. Πόσο απέχουν τα $\mathrm{Α}$ και $\mathrm{Α}'$;

\ex Πόσους άξονες συμμετρίας έχουν: ισόπλευρο τρίγωνο, ορθογώνιο μη τετράγωνο, ρόμβος, κύκλος;

\ex Σημείο $\mathrm{Μ}$ της μεσοκαθέτου του $\mathrm{ΑΒ}$ απέχει $5$ cm από το $\mathrm{Α}$. Πόσο απέχει από το $\mathrm{Β}$;

\ex Ένα τμήμα $\mathrm{ΑΒ}=10$ cm. Πόσο απέχει η μεσοκάθετός του από το $\mathrm{Α}$;

\ex Το τρίγωνο $\mathrm{ΑΒΓ}$ έχει $\mathrm{ΑΒ}=\mathrm{ΑΓ}$. Ποιος είναι ένας άξονας συμμετρίας του;

\ex Βρες το συμμετρικό του $\mathrm{Α}(-2)$ στον άξονα των αριθμών ως προς το $0$ και ως προς το $1$.

\ex Σημείο $\mathrm{Α}$ απέχει $7$ cm από το $\mathrm{Ο}$. Πόσο απέχουν τα $\mathrm{Α}$ και $\mathrm{Α}'$ (συμμετρικό ως προς $\mathrm{Ο}$);

\ex Ποια κεφαλαία γράμματα από τα Α, Β, Η, Ν, Σ, Τ, Φ έχουν άξονα και ποια κέντρο συμμετρίας;

\ex Δύο παράλληλες ευθείες τέμνονται από τρίτη και μία γωνία είναι $80^\circ$. Πόσες γωνίες είναι $80^\circ$ και πόσες $100^\circ$;

\ex Οι εντός επί τα αυτά γωνίες δύο παραλλήλων είναι $2x$ και $x+60^\circ$. Βρες το $x$.

\ex Δύο εντός εναλλάξ γωνίες είναι $64^\circ$ και $67^\circ$. Είναι παράλληλες οι ευθείες; Γιατί;

\ex Η μεσοκάθετος και η διαγώνιος: Σε ρόμβο $\mathrm{ΑΒΓΔ}$ η $\mathrm{ΒΔ}$ είναι ποιας διαγωνίου μεσοκάθετος;

\ex Ένα ορθογώνιο $12\times5$ έχει διαγώνιο $13$. Πού βρίσκεται το κέντρο συμμετρίας και πόσο απέχει από κάθε κορυφή;

\ex Πόσο είναι το μήκος του συμμετρικού τμήματος $\mathrm{Α}'\mathrm{Β}'$ ενός τμήματος $\mathrm{ΑΒ}=8{,}5$ cm; Τι σχέση έχει με το $\mathrm{ΑΒ}$ αν η συμμετρία είναι ως προς σημείο εκτός της ευθείας $\mathrm{ΑΒ}$;

\ex \textbf{Ερωτήσεις αυτοαξιολόγησης (συμμετρίες).} Σωστό ή λάθος; α) Τα συμμετρικά σχήματα ως προς άξονα είναι ίσα. \quad β) Η μεσοκάθετος περνά από το μέσο του τμήματος. \quad γ) Το ισοσκελές τρίγωνο έχει κέντρο συμμετρίας. \quad δ) Ο κύκλος έχει άπειρους άξονες.

\ex \textbf{Ερωτήσεις αυτοαξιολόγησης (παράλληλες).} Σωστό ή λάθος; α) Οι εντός εναλλάξ γωνίες δύο παραλλήλων είναι ίσες. \quad β) Οι εντός επί τα αυτά είναι ίσες. \quad γ) Αν οι εντός εναλλάξ είναι ίσες, οι ευθείες είναι παράλληλες. \quad δ) Οι κατακορυφήν γωνίες είναι ίσες.

\tmima{Ο χάρτης της επανάληψης}

Μετά τις λύσεις, δες σε ποιες ασκήσεις δυσκολεύτηκες. Ο πίνακας σου λέει ποια παράγραφο να ξαναδιαβάσεις.

\renewcommand{\arraystretch}{1.5}
\begin{tabular*}{\linewidth}{@{\extracolsep{\fill}} >{\raggedright\arraybackslash}p{0.38\linewidth} >{\raggedright\arraybackslash}p{0.52\linewidth}}
\toprule
\textbf{Αν δυσκολεύτηκες στις} & \textbf{ξαναδιάβασε τις παραγράφους}\\
\midrule
ΕΒ2.1, ΕΒ2.14 & 2.1 (συμμετρία ως προς άξονα)\\
ΕΒ2.2, ΕΒ2.5, ΕΒ2.8 & 2.2 (άξονες συμμετρίας)\\
ΕΒ2.3, ΕΒ2.4, ΕΒ2.12 & 2.3 (μεσοκάθετος)\\
ΕΒ2.6, ΕΒ2.7 & 2.4 (συμμετρία ως προς σημείο)\\
ΕΒ2.13 & 2.5 (κέντρο συμμετρίας)\\
ΕΒ2.9, ΕΒ2.10, ΕΒ2.11 & 2.6 (παράλληλες και τέμνουσα)\\
ΕΒ2.15, ΕΒ2.16 & 2.1 έως 2.6\\
\bottomrule
\end{tabular*}


\kefalaio{Γεωμετρία του επιπέδου}{Από τη γωνία στον κύκλο: όλα τα σχήματα του επιπέδου με ένα κανόνα και έναν διαβήτη}{%
\begin{tabular}{@{}l@{\hspace{1em}}p{0.82\linewidth}@{}}
3.1 & Είδη γωνιών\\[2pt]
3.2 & Μεσοκάθετος ευθύγραμμου τμήματος\\[2pt]
3.3 & Διχοτόμος γωνίας\\[2pt]
3.4 & Είδη τριγώνων και στοιχεία τους\\[2pt]
3.5 & Ιδιότητες ισοσκελούς και ισόπλευρου τριγώνου\\[2pt]
3.6 & Κατακορυφήν, συμπληρωματικές και παραπληρωματικές γωνίες\\[2pt]
3.7 & Σχετικές θέσεις ευθειών στο επίπεδο\\[2pt]
3.8 & Γωνίες που σχηματίζονται από δύο παράλληλες ευθείες και μια τέμνουσα\\[2pt]
3.9 & Άθροισμα γωνιών τριγώνου\\[2pt]
3.10 & Τετράπλευρα, παραλληλόγραμμα, τραπέζια\\[2pt]
3.11 & Είδη παραλληλογράμμων και ιδιότητες: ορθογώνιο, ρόμβος, τετράγωνο\\[2pt]
3.12 & Κύκλος και στοιχεία του\\[2pt]
3.13 & Σχετικές θέσεις ευθείας και κύκλου\\[2pt]
 & Επανάληψη του κεφαλαίου\\
\end{tabular}}
\paragrafos{Στοιχεία και είδη τριγώνων}{Β.3.1, «Στοιχεία τριγώνου – Είδη τριγώνων»}

\begin{anagke}
Στο εργαστήριο τεχνολογίας του σχολείου η Αγγελική και ο Νίκος φτιάχνουν μια μικρή γέφυρα από ξυλάκια. Ο Νίκος πρώτα φτιάχνει ένα τετράγωνο, αλλά αυτό «κουνιέται» και στραβώνει. Η κυρία Ελένη λέει: «Βάλε ένα ξυλάκι διαγώνια και θα γίνουν δύο τρίγωνα. Το τρίγωνο δεν παραμορφώνεται!».

Η Πολυξένη μετράει τα τρίγωνα της γέφυρας και βλέπει ότι δεν είναι όλα ίδια: άλλα έχουν όλες τις πλευρές ίσες, άλλα δύο ίσες, άλλα καμία. Ο Βασίλης προσθέτει ότι άλλα έχουν ορθή γωνία και άλλα όχι. Για να συνεννοούνται, χρειάζονται ονόματα: ποια τρίγωνα λέμε ισοσκελή, ποια ορθογώνια, ποια οξυγώνια.

Σε αυτή την παράγραφο μαθαίνουμε τα στοιχεία του τριγώνου, ταξινομούμε τα τρίγωνα με δύο κριτήρια (πλευρές και γωνίες) και γνωρίζουμε τη διάμεσο, το ύψος και τη διχοτόμο.
\end{anagke}

\begin{lexi}{τρίγωνο, ισοσκελές, σκαληνό, διάμεσος}{ονόματα τριγώνων}
Το \lexiword{τρίγωνο} έχει τρεις γωνίες. Το \lexiword{ισοσκελές} έχει «ίσα σκέλη» (πλευρές), το \lexiword{ισόπλευρο} όλες τις πλευρές ίσες. Το \lexiword{σκαληνό} (από το «σκαληνός» = λοξός, ανώμαλος) δεν έχει δύο ίσες πλευρές. Η \lexiword{διάμεσος} χωρίζει την πλευρά στη «μέση».
\end{lexi}

\begin{metafrasi}
Τρεις κορυφές, τρεις πλευρές, τρεις γωνίες. & $\mathrm{Α},\mathrm{Β},\Gamma$, $\mathrm{ΑΒ},\mathrm{ΒΓ},\Gamma\mathrm{Α}$, $\hat{\mathrm{Α}},\hat{\mathrm{Β}},\hat{\Gamma}$\\
Ένα τρίγωνο με ορθή γωνία. & ορθογώνιο ($\hat{\mathrm{Α}}=90^\circ$)\\
Ένα τρίγωνο με όλες τις πλευρές ίσες. & ισόπλευρο\\
Ένα τμήμα από κορυφή στο μέσο της απέναντι πλευράς. & διάμεσος\\
Ένα τμήμα από κορυφή κάθετο στην απέναντι πλευρά. & ύψος\\
\end{metafrasi}

\tmima{Η θεωρία}

\begin{orismos}{στοιχεία τριγώνου}
Κάθε τρίγωνο $\mathrm{ΑΒΓ}$ έχει τρεις \textbf{κορυφές} $\mathrm{Α},\mathrm{Β},\Gamma$, τρεις \textbf{πλευρές} $\mathrm{ΑΒ},\mathrm{ΒΓ},\Gamma\mathrm{Α}$ και τρεις \textbf{γωνίες} $\hat{\mathrm{Α}},\hat{\mathrm{Β}},\hat{\Gamma}$. Τα $\mathrm{ΑΒ},\mathrm{ΒΓ},\Gamma\mathrm{Α}$ δηλώνουν και τα μήκη των πλευρών. Περίμετρος είναι το άθροισμα $\mathrm{ΑΒ}+\mathrm{ΒΓ}+\Gamma\mathrm{Α}$.
\end{orismos}

\begin{kanonas}{ταξινόμηση τριγώνων}
\textbf{Ως προς τις γωνίες:} \emph{ορθογώνιο} (μία ορθή γωνία), \emph{αμβλυγώνιο} (μία αμβλεία γωνία), \emph{οξυγώνιο} (όλες οι γωνίες οξείες).\\
\textbf{Ως προς τις πλευρές:} \emph{ισόπλευρο} (τρεις ίσες), \emph{ισοσκελές} (δύο ίσες), \emph{σκαληνό} (όλες άνισες).
\end{kanonas}

\begin{center}
\begin{tikzpicture}[scale=0.62]
\draw[thick] (0,0)--(0,2.2)--(2.8,0)--cycle; \node at (1,-0.6) {\small ορθογώνιο};
\begin{scope}[xshift=4.6cm]
\draw[thick] (0,0)--(3,0)--(-0.8,1.6)--cycle; \node at (1.1,-0.6) {\small αμβλυγώνιο};
\end{scope}
\begin{scope}[xshift=9.6cm]
\draw[thick] (0,0)--(3,0)--(1.4,2.4)--cycle; \node at (1.5,-0.6) {\small οξυγώνιο};
\end{scope}
\end{tikzpicture}
\end{center}

\begin{center}
\begin{tikzpicture}[scale=0.62]
\draw[thick] (0,0)--(3,0)--(1.5,2.6)--cycle; \node at (1.5,-0.6) {\small ισόπλευρο};
\begin{scope}[xshift=5cm]
\draw[thick] (0,0)--(3,0)--(1.5,2.6)--cycle; \node at (1.5,-0.6) {\small ισοσκελές};
\end{scope}
\begin{scope}[xshift=10cm]
\draw[thick] (0,0)--(3.3,0)--(0.9,2.2)--cycle; \node at (1.5,-0.6) {\small σκαληνό};
\end{scope}
\end{tikzpicture}
\end{center}

\begin{orismos}{δευτερεύοντα στοιχεία}
\textbf{Διάμεσος} τριγώνου: τμήμα που ενώνει μια κορυφή με το μέσο της απέναντι πλευράς.\\
\textbf{Ύψος}: τμήμα από μια κορυφή, κάθετο στην ευθεία της απέναντι πλευράς.\\
\textbf{Διχοτόμος}: το τμήμα της διχοτόμου μιας γωνίας του τριγώνου από την κορυφή ως την απέναντι πλευρά.
\end{orismos}

Κάθε τρίγωνο έχει τρεις διαμέσους, τρία ύψη και τρεις διχοτόμους. Οι τρεις διάμεσοι τέμνονται σε ένα σημείο (βαρύκεντρο), τα τρία ύψη επίσης (ορθόκεντρο). Στο αμβλυγώνιο τρίγωνο τα ύψη φέρονται εκτός τριγώνου, πάνω στις προεκτάσεις των πλευρών. Στο ορθογώνιο τρίγωνο οι δύο κάθετες πλευρές είναι και ύψη.

\begin{prosoxi}
\begin{itemize}[leftmargin=1.2em,itemsep=0.3ex,topsep=2pt]
\item Το ισόπλευρο είναι και ισοσκελές (έχει δύο ίσες πλευρές), αλλά όχι ανάποδα.
\item Ένα τρίγωνο έχει το πολύ μία ορθή ή μία αμβλεία γωνία.
\item Το ύψος είναι \emph{κάθετο} τμήμα, η διάμεσος πάει στο \emph{μέσο}. Δεν είναι πάντα το ίδιο τμήμα.
\item Το ύψος μπορεί να βρίσκεται έξω από το τρίγωνο.
\end{itemize}
\end{prosoxi}

\tmima{Λυμένα παραδείγματα}

\paradeigma{Η περίμετρος και το είδος}
\tixero Τρίγωνο με πλευρές $5$, $5$ και $8$ cm.\par
\tithelo Το είδος του ως προς τις πλευρές και την περίμετρο.\par
\possindeo Δύο πλευρές είναι ίσες, άρα ισοσκελές. Περίμετρος $5+5+8=18$ cm.\par
\apantisi Ισοσκελές, περίμετρος $18$ cm.\par
\elegxos Το άθροισμα των δύο μικρότερων ($5+5=10$) είναι μεγαλύτερο από την τρίτη ($8$), άρα το τρίγωνο υπάρχει.

\paradeigma{Το ύψος}
\tixero Τρίγωνο $\mathrm{ΑΒΓ}$ με $\mathrm{ΒΓ}=10$ cm και ύψος $\mathrm{ΑΔ}=4$ cm.\par
\tithelo Σχεδίασε τα ύψη όταν το τρίγωνο είναι οξυγώνιο, ορθογώνιο ή αμβλυγώνιο.\par
\possindeo Από κάθε κορυφή φέρνω κάθετο στην απέναντι πλευρά (ή την προέκτασή της, στο αμβλυγώνιο).\par
\apantisi Τα τρία ύψη περνούν από το ίδιο σημείο (ορθόκεντρο).\par
\elegxos Στο ορθογώνιο τρίγωνο το ορθόκεντρο είναι η κορυφή της ορθής γωνίας.

\paradeigma{Η διάμεσος}
\tixero Στο τρίγωνο $\mathrm{ΑΒΓ}$ η πλευρά $\mathrm{ΒΓ}=7{,}4$ cm και η $\mathrm{ΑΜ}$ είναι διάμεσος.\par
\tithelo Το μήκος $\mathrm{ΒΜ}$.\par
\possindeo Το $\mathrm{Μ}$ είναι το μέσο της $\mathrm{ΒΓ}$, άρα $\mathrm{ΒΜ}=7{,}4:2=3{,}7$ cm.\par
\apantisi $3{,}7$ cm.\par
\elegxos $\mathrm{ΒΜ}+\mathrm{ΜΓ}=3{,}7+3{,}7=7{,}4$.

\tmima{Ασκήσεις}

\zoni{Το κατάλαβα}{ονόματα και είδη}

\ex Ονόμασε τα τρίγωνα με πλευρές: α) $4,4,4$ β) $5,5,3$ γ) $3,4,5$ ως προς τις πλευρές.

\ex Βρες την περίμετρο τριγώνου με πλευρές $6{,}5$, $7$ και $8{,}5$ cm.

\ex Σχεδίασε οξυγώνιο, ορθογώνιο και αμβλυγώνιο τρίγωνο και φέρε τα ύψη τους.

\ex Στο τρίγωνο $\mathrm{ΑΒΓ}$ με $\mathrm{ΒΓ}=4{,}4$ cm φέρε τη διάμεσο $\mathrm{ΑΜ}$. Πόσο είναι το $\mathrm{ΒΜ}$;

\ex Σχεδίασε τρίγωνο $\mathrm{ΑΒΓ}$, βρες τα μέσα $\mathrm{Δ}$, $\mathrm{Ε}$, $\mathrm{Ζ}$ των πλευρών και σχεδίασε τις τρεις διαμέσους. Τι παρατηρείς;

\ex Σχεδίασε τρίγωνο και τη διχοτόμο μιας γωνίας του. Μέτρησε τις δύο γωνίες που σχηματίζονται.

\ex Σωστό ή λάθος; α) Το ισόπλευρο είναι ισοσκελές. \quad β) Το αμβλυγώνιο έχει δύο αμβλείες. \quad γ) Το ορθογώνιο τρίγωνο μπορεί να είναι ισοσκελές. \quad δ) Το ισόπλευρο είναι πάντα οξυγώνιο. \quad ε) Ένα σκαληνό δεν μπορεί να είναι ορθογώνιο.

\zoni{Το χρησιμοποιώ}{ξύλινες κατασκευές}

\ex \textbf{Η γέφυρα.} Ένα ξύλινο τρίγωνο της γέφυρας έχει πλευρές $30$, $30$ και $40$ cm. Τι είδους είναι και πόση είναι η περίμετρος;

\ex \textbf{Το μαντήλι.} Ένα τριγωνικό μαντήλι έχει πλευρές $60$, $60$ και $60$ cm. Πόσα εκατοστά κορδόνι χρειάζεται για να το γυρίσει κανείς γύρω γύρω;

\ex \textbf{Το σχοινί.} Ένα σχοινί $24$ μ. σχηματίζει ισόπλευρο τρίγωνο. Πόση είναι η πλευρά του;

\ex \textbf{Το τρίγωνο της κορνίζας.} Ισοσκελές τρίγωνο έχει περίμετρο $50$ cm και βάση $20$ cm. Πόσο μακριές είναι οι δύο ίσες πλευρές;

\ex \textbf{Η κατασκευή.} Μπορείς να φτιάξεις τρίγωνο με ξυλάκια $3$, $4$ και $9$ cm; Και με $3$, $4$ και $5$ cm; Εξήγησε.

\zoni{Το σκέφτομαι}{σύγκριση ειδών}

\ex Σε ποια είδη τριγώνων ως προς τις γωνίες μπορεί να ανήκει ένα ισοσκελές τρίγωνο;

\ex Μπορεί ένα ορθογώνιο τρίγωνο να είναι ισόπλευρο; Εξήγησε.

\ex Πόσα το πολύ ύψη του τριγώνου βρίσκονται εκτός αυτού; (Σκέψου το αμβλυγώνιο.)

\ex Συμπλήρωσε πίνακα: ποιοι συνδυασμοί (ορθογώνιο, αμβλυγώνιο, οξυγώνιο) × (σκαληνό, ισοσκελές, ισόπλευρο) υπάρχουν;

\zoni{Για το σπίτι}{τρίγωνα στην πόλη}

\ex Βρες πέντε τριγωνικά αντικείμενα στο σπίτι (σήμα, φέτα πίτσας, ράφι, στέγη κ.λπ.) και πες το είδος του καθενός.

\ex Κόψε από χαρτί τρίγωνο και δίπλωσε το ώστε να βρεις το μέσο μιας πλευράς και τη διάμεσο.

\paragrafos{Άθροισμα γωνιών τριγώνου}{Β.3.2, «Άθροισμα γωνιών τριγώνου – Ιδιότητες ισοσκελούς τριγώνου»}

\begin{anagke}
Ο Βασίλης κόβει από χαρτί ένα τρίγωνο και σκίζει τις τρεις γωνίες του. Τις βάζει δίπλα δίπλα, με τις κορυφές στο ίδιο σημείο, και οι τρεις μυτούλες σχηματίζουν μια ευθεία γραμμή. «Ευθεία γωνία!», φωνάζει. Η Αγγελική κάνει το ίδιο με ένα άλλο τρίγωνο, πιο λεπτό και μακρύ: πάλι ευθεία. Ο Νίκος το κάνει με αμβλυγώνιο: ξανά ευθεία.

Η Πολυξένη ρωτάει: «Ισχύει πάντα; Και πώς το ξέρουμε χωρίς να κόψουμε όλα τα τρίγωνα του κόσμου;». Η απάντηση είναι ότι το αποδεικνύουμε με όσα ξέρουμε για τις παράλληλες ευθείες. Η Ρήνη προσθέτει: «Και τα ισοσκελή; Αν διπλώσω ένα ισοσκελές στη μέση, οι δύο μισές πέφτουν η μία πάνω στην άλλη!».

Σε αυτή την παράγραφο μαθαίνουμε ότι το άθροισμα των γωνιών κάθε τριγώνου είναι $180^\circ$, τι ισχύει για την εξωτερική γωνία, και τις ιδιότητες του ισοσκελούς και του ισόπλευρου τριγώνου.
\end{anagke}

\begin{lexi}{εξωτερική γωνία, βάση}{ονόματα}
Η \lexiword{εξωτερική γωνία} είναι αυτή που σχηματίζεται από μια πλευρά και την προέκταση της διπλανής. Η \lexiword{βάση} του ισοσκελούς είναι η πλευρά που δεν είναι από τις δύο ίσες, και οι \lexiword{προσκείμενες γωνίες} της είναι οι δύο γωνίες στα άκρα της.
\end{lexi}

\begin{metafrasi}
Οι τρεις γωνίες φτιάχνουν ευθεία γωνία. & $\hat{\mathrm{Α}}+\hat{\mathrm{Β}}+\hat{\Gamma}=180^\circ$\\
Οι δύο οξείες του ορθογωνίου τριγώνου. & συμπληρωματικές: $\hat{\mathrm{Β}}+\hat{\Gamma}=90^\circ$\\
Η εξωτερική γωνία στο $\Gamma$. & $\hat{\varphi}=\hat{\mathrm{Α}}+\hat{\mathrm{Β}}$\\
Οι γωνίες στη βάση του ισοσκελούς. & $\hat{\mathrm{Β}}=\hat{\Gamma}$\\
Όλες οι γωνίες του ισόπλευρου. & $60^\circ$\\
\end{metafrasi}

\tmima{Η θεωρία}

\begin{kanonas}{άθροισμα γωνιών τριγώνου}
Σε κάθε τρίγωνο $\mathrm{ΑΒΓ}$: $\hat{\mathrm{Α}}+\hat{\mathrm{Β}}+\hat{\Gamma}=180^\circ$.
\end{kanonas}

\textbf{Γιατί ισχύει.} Φέρνω από το $\mathrm{Α}$ την ευθεία $x\mathrm{Α}y$ παράλληλη προς τη $\mathrm{ΒΓ}$. Η $x\hat{\mathrm{Α}}\mathrm{Β}$ είναι ίση με την $\hat{\mathrm{Β}}$ (εντός εναλλάξ) και η $y\hat{\mathrm{Α}}\Gamma$ ίση με την $\hat{\Gamma}$ (εντός εναλλάξ). Οι τρεις γωνίες γύρω από το $\mathrm{Α}$ σχηματίζουν ευθεία γωνία, άρα $\hat{\mathrm{Β}}+\hat{\mathrm{Α}}+\hat{\Gamma}=180^\circ$.

\begin{center}
\begin{tikzpicture}[scale=0.9]
\draw[thick] (-0.4,0)--(3.8,0) ;
\node[below left] at (0,0) {$\mathrm{Β}$}; \node[below right] at (3.4,0) {$\Gamma$};
\draw[thick] (0,0)--(1.4,1.8)--(3.4,0)--cycle;
\draw[thick,gblue] (-0.6,1.8)--(3.6,1.8);
\node[above] at (1.4,1.8) {$\mathrm{Α}$};
\node at (0.55,0.2) {\small $\omega$}; \node at (2.9,0.2) {\small $\theta$};
\node at (0.7,1.6) {\small $\omega$}; \node at (2.1,1.6) {\small $\theta$};
\node at (1.4,1.45) {\small $\hat{\mathrm{Α}}$};
\end{tikzpicture}
\end{center}

\begin{kanonas}{συνέπειες}
\begin{itemize}[leftmargin=1.2em,itemsep=0.3ex,topsep=2pt]
\item Στο \textbf{ορθογώνιο} τρίγωνο οι δύο οξείες γωνίες είναι \textbf{συμπληρωματικές}: $\hat{\mathrm{Β}}+\hat{\Gamma}=90^\circ$.
\item Η \textbf{εξωτερική γωνία} τριγώνου ισούται με το άθροισμα των δύο \emph{απέναντι} εσωτερικών: $\hat{\varphi}=\hat{\mathrm{Α}}+\hat{\mathrm{Β}}$ (γιατί $\hat{\varphi}=180^\circ-\hat{\Gamma}$).
\item Στο \textbf{ισόπλευρο} τρίγωνο κάθε γωνία είναι $180^\circ:3=60^\circ$.
\end{itemize}
\end{kanonas}

\begin{kanonas}{ισοσκελές τρίγωνο}
Σε ισοσκελές τρίγωνο ($\mathrm{ΑΒ}=\mathrm{ΑΓ}$, κορυφή $\mathrm{Α}$, βάση $\mathrm{ΒΓ}$):
\begin{itemize}[leftmargin=1.2em,itemsep=0.3ex,topsep=2pt]
\item η ευθεία της διαμέσου προς τη βάση είναι \textbf{άξονας συμμετρίας},
\item η διάμεσος προς τη βάση είναι και \textbf{ύψος} και \textbf{διχοτόμος},
\item οι \textbf{προσκείμενες στη βάση γωνίες είναι ίσες}: $\hat{\mathrm{Β}}=\hat{\Gamma}$.
\end{itemize}
Στο \textbf{ισόπλευρο}: οι ευθείες των διαμέσων είναι άξονες συμμετρίας, κάθε διάμεσος είναι ύψος και διχοτόμος, όλες οι πλευρές και γωνίες είναι ίσες.
\end{kanonas}

\begin{prosoxi}
\begin{itemize}[leftmargin=1.2em,itemsep=0.3ex,topsep=2pt]
\item Στο ισοσκελές, η ιδιότητα «διάμεσος = ύψος = διχοτόμος» ισχύει για το τμήμα προς τη \emph{βάση}, όχι για τα άλλα.
\item Στο ισοσκελές ορθογώνιο τρίγωνο οι δύο οξείες είναι $45^\circ$ η καθεμία.
\item Η εξωτερική γωνία είναι ίση με το άθροισμα των \emph{δύο άλλων} γωνιών (όχι της διπλανής της).
\end{itemize}
\end{prosoxi}

\tmima{Λυμένα παραδείγματα}

\paradeigma{Η τρίτη γωνία}
\tixero Τρίγωνο με γωνίες $\hat{\mathrm{Α}}=47^\circ$ και $\hat{\mathrm{Β}}=68^\circ$.\par
\tithelo Η $\hat{\Gamma}$.\par
\possindeo $\hat{\Gamma}=180^\circ-(47^\circ+68^\circ)=180^\circ-115^\circ=65^\circ$.\par
\apantisi $65^\circ$ (οξυγώνιο τρίγωνο).\par
\elegxos $47+68+65=180$.

\paradeigma{Η εξωτερική γωνία}
\tixero Τρίγωνο με $\hat{\mathrm{Α}}=50^\circ$, $\hat{\mathrm{Β}}=60^\circ$.\par
\tithelo Η εξωτερική γωνία στο $\Gamma$.\par
\possindeo $\hat{\varphi}=\hat{\mathrm{Α}}+\hat{\mathrm{Β}}=110^\circ$. Ή: $\hat{\Gamma}=70^\circ$ και $\hat{\varphi}=180^\circ-70^\circ=110^\circ$.\par
\apantisi $110^\circ$.\par
\elegxos Οι δύο τρόποι δίνουν το ίδιο.

\paradeigma{Ισοσκελές με γωνία κορυφής}
\tixero Ισοσκελές τρίγωνο με γωνία κορυφής $\hat{\mathrm{Α}}=40^\circ$.\par
\tithelo Οι γωνίες της βάσης.\par
\possindeo $\hat{\mathrm{Β}}=\hat{\Gamma}$ και $\hat{\mathrm{Β}}+\hat{\Gamma}=180^\circ-40^\circ=140^\circ$. Άρα $\hat{\mathrm{Β}}=\hat{\Gamma}=70^\circ$.\par
\apantisi $70^\circ$ και $70^\circ$.\par
\elegxos $40+70+70=180$.

\paradeigma{Ορθογώνιο και ισοσκελές}
\tixero Ένα τρίγωνο είναι ορθογώνιο και ισοσκελές.\par
\tithelo Οι γωνίες του.\par
\possindeo Η ορθή είναι $90^\circ$. Οι άλλες δύο είναι ίσες και με άθροισμα $90^\circ$, άρα $45^\circ$ η καθεμία.\par
\apantisi $90^\circ,45^\circ,45^\circ$.\par
\elegxos $90+45+45=180$.

\tmima{Ασκήσεις}

\zoni{Το κατάλαβα}{γωνίες τριγώνων}

\ex Βρες την τρίτη γωνία τριγώνου με γωνίες: α) $35^\circ,\,80^\circ$ β) $90^\circ,\,28^\circ$ γ) $112^\circ,\,41^\circ$.

\ex Ισοσκελές τρίγωνο έχει γωνία κορυφής $50^\circ$. Βρες τις γωνίες της βάσης.

\ex Ισοσκελές τρίγωνο έχει γωνία στη βάση $72^\circ$. Βρες τη γωνία κορυφής.

\ex Πόσο είναι κάθε γωνία ισοπλεύρου τριγώνου; Πόσο η εξωτερική του γωνία;

\ex Οι οξείες ενός ορθογωνίου τριγώνου είναι $x$ και $2x$. Βρες τες.

\ex Σε τρίγωνο η εξωτερική γωνία στο $\Gamma$ είναι $125^\circ$ και $\hat{\mathrm{Α}}=48^\circ$. Βρες τη $\hat{\mathrm{Β}}$ και τη $\hat{\Gamma}$.

\ex Σωστό ή λάθος; α) Τρίγωνο μπορεί να έχει δύο ορθές γωνίες. \quad β) Οι γωνίες βάσης ισοσκελούς είναι ίσες. \quad γ) Σε ισόπλευρο τρίγωνο όλες οι γωνίες είναι $60^\circ$. \quad δ) Οι οξείες γωνίες ορθογωνίου είναι παραπληρωματικές.

\zoni{Το χρησιμοποιώ}{στέγες και σήματα}

\ex \textbf{Η στέγη.} Μια στέγη έχει σχήμα ισοσκελούς τριγώνου με γωνία κορυφής $100^\circ$. Πόση είναι η γωνία που σχηματίζει κάθε πλευρά με το οριζόντιο δοκάρι;

\ex \textbf{Το σήμα.} Ένα τριγωνικό σήμα έχει ισόπλευρο σχήμα. Πόσο είναι το άθροισμα των τριών γωνιών του;

\ex \textbf{Η σκάλα.} Μια σκάλα ακουμπά στον τοίχο και σχηματίζει με αυτόν γωνία $25^\circ$. Ο τοίχος είναι κάθετος στο έδαφος. Πόση γωνία σχηματίζει η σκάλα με το έδαφος;

\ex \textbf{Ο χαρταετός.} Ένα τριγωνικό κομμάτι ενός χαρταετού έχει γωνίες $55^\circ$ και $75^\circ$. Πόση είναι η τρίτη γωνία;

\ex \textbf{Η γωνιά.} Ένας ξύλινος γνώμονας έχει σχήμα ισοσκελούς ορθογωνίου τριγώνου. Πόσες μοίρες είναι οι δύο ίσες γωνίες του;

\zoni{Το σκέφτομαι}{γιατί συμβαίνει}

\ex Εξήγησε με παράλληλη από την κορυφή γιατί το άθροισμα των γωνιών ενός τριγώνου είναι $180^\circ$.

\ex Εξήγησε γιατί τρίγωνο δεν μπορεί να έχει δύο αμβλείες γωνίες.

\ex Αν σε ισοσκελές τρίγωνο μια γωνία είναι $100^\circ$, ποια θέση έχει (κορυφής ή βάσης); Βρες τις άλλες δύο.

\ex Στο ισοσκελές τρίγωνο $\mathrm{ΑΒΓ}$ ($\mathrm{ΑΒ}=\mathrm{ΑΓ}$) η $\mathrm{ΑΔ}$ είναι διάμεσος προς τη $\mathrm{ΒΓ}$ και $\mathrm{ΒΓ}=12$ cm. Πόσο είναι το $\mathrm{ΒΔ}$ και πόση η γωνία $\mathrm{ΑΔΒ}$;

\zoni{Για το σπίτι}{γωνίες με το χέρι}

\ex Κόψε ένα τρίγωνο από χαρτί, σκίσε τις γωνίες του και τοποθέτησέ τες δίπλα δίπλα. Τι σχηματίζουν;

\ex Δίπλωσε ένα ισοσκελές τρίγωνο στη μέση και δες ποιες γωνίες πέφτουν η μία πάνω στην άλλη.

\paragrafos{Διχοτόμος γωνίας}{Β.3.3, «Διχοτόμος γωνίας»}

\begin{anagke}
Στο διάλειμμα ο Νίκος και η Αγγελική μοιράζονται ένα κομμάτι πίτσα σε σχήμα «φέτας» με μυτερή γωνία. «Θέλω ακριβώς τη μισή», λέει η Αγγελική, «και όχι με το μάτι». Ο Νίκος κοιτάζει τη φέτα: «Αν κόψω από την μύτη προς την άκρη, πώς θα ξέρω ότι τα δύο κομμάτια είναι ίσα;». Η Φανή το σημειώνει: ζητούμενο είναι μια ευθεία από την κορυφή που χωρίζει τη γωνία σε δύο ίσες γωνίες.

Η Ρήνη βρίσκει και ένα δεύτερο πρόβλημα στο διπλανό γήπεδο: «Το καφενείο θέλει να βάλει ένα τραπέζι σε μια γωνία, ώστε να απέχει ίσα από τους δύο τοίχους». Ο Ηλίας υποψιάζεται ότι οι δύο προβλήματα είναι ένα: η ευθεία που κόβει τη γωνία στη μέση είναι και η ευθεία των σημείων που ισαπέχουν από τις πλευρές της.

Σε αυτή την παράγραφο μαθαίνουμε τι είναι η διχοτόμος μιας γωνίας, ποια ιδιότητα έχουν τα σημεία της και πώς τη σχεδιάζουμε με μοιρογνωμόνιο, με δίπλωση και με διαβήτη.
\end{anagke}

\begin{lexi}{διχοτόμος, απόσταση}{διχοτόμος}
Η \lexiword{διχοτόμος} είναι αυτή που «κόβει δίχα», δηλαδή στα δύο (δίχα + τέμνω): χωρίζει τη γωνία σε δύο ίσα μέρη. Η \lexiword{απόσταση} σημείου από ευθεία είναι το μήκος του κάθετου τμήματος που φέρνουμε από το σημείο προς την ευθεία.
\end{lexi}

\begin{metafrasi}
Η ημιευθεία που κόβει τη γωνία στη μέση. & η διχοτόμος $\mathrm{Οδ}$ της $x\hat{\mathrm{Ο}}y$\\
Οι δύο μισές γωνίες είναι ίσες. & $x\hat{\mathrm{Ο}}\delta=\delta\hat{\mathrm{Ο}}y$\\
Κάθε μισή είναι το μισό της γωνίας. & $x\hat{\mathrm{Ο}}\delta=\dfrac{x\hat{\mathrm{Ο}}y}{2}$\\
Το σημείο $\mathrm{Α}$ έχει ίσες αποστάσεις από τις δύο πλευρές. & $\mathrm{ΑΒ}=\mathrm{ΑΓ}$\\
\end{metafrasi}

\tmima{Η θεωρία}

\begin{orismos}{διχοτόμος γωνίας}
\textbf{Διχοτόμος} μιας γωνίας $x\hat{\mathrm{Ο}}y$ λέγεται η ημιευθεία $\mathrm{Οδ}$ που έχει αρχή την κορυφή της γωνίας και τη χωρίζει σε δύο ίσες γωνίες: $x\hat{\mathrm{Ο}}\delta=\delta\hat{\mathrm{Ο}}y$.
\end{orismos}

Κάθε γωνία έχει μία και μόνο μία διχοτόμο. Κάθε μία από τις δύο μισές γωνίες έχει μέτρο ίσο με το μισό της αρχικής.

\begin{center}
\begin{tikzpicture}[scale=0.95]
\draw[thick] (5,0)--(0,0)--({5*cos(60)},{5*sin(60)});
\draw[gblue,thick] (0,0)--({5*cos(30)},{5*sin(30)});
\fill (0,0) circle (2pt) node[below left] {$\mathrm{Ο}$};
\node[right] at (5,0) {$x$}; \node[above] at ({5*cos(60)},{5*sin(60)}) {$y$}; \node[right,gblue] at ({5*cos(30)},{5*sin(30)}) {$\delta$};
\draw[gblue] (1.1,0) arc(0:30:1.1); \draw[gblue] (1.4,0) arc(0:30:1.4); \draw[gwarmD] (1.3,0.75) arc(30:60:1.5); \draw[gwarmD] (1.6,0.92) arc(30:60:1.85);
\end{tikzpicture}
\end{center}

Θυμόμαστε ότι \textbf{απόσταση} ενός σημείου $\mathrm{Α}$ από μια ευθεία $\varepsilon$ είναι το μήκος του κάθετου τμήματος $\mathrm{ΑΒ}$ που φέρνουμε από το $\mathrm{Α}$ προς την $\varepsilon$.

\begin{idiotita}{σημεία της διχοτόμου}
Κάθε σημείο της διχοτόμου μιας γωνίας \textbf{ισαπέχει} από τις πλευρές της.
\end{idiotita}

Ισχύει και το αντίστροφο: κάθε σημείο στο εσωτερικό της γωνίας που ισαπέχει από τις πλευρές της ανήκει στη διχοτόμο της. Έτσι η διχοτόμος είναι ακριβώς το σύνολο των σημείων που ισαπέχουν από τις πλευρές της γωνίας, όπως η μεσοκάθετος είναι το σύνολο των σημείων που ισαπέχουν από δύο σημεία.

\begin{center}
\begin{tikzpicture}[scale=0.95]
\draw[thick] (5,0)--(0,0)--({5*cos(60)},{5*sin(60)});
\draw[gblue,thick] (0,0)--({4.5*cos(30)},{4.5*sin(30)});
\fill (0,0) circle (2pt) node[below left] {$\mathrm{Ο}$};
\coordinate (P) at (2.944,1.7); \coordinate (B) at (2.944,0); \coordinate (G) at (1.472,2.55);
\fill (P) circle (2pt) node[above right] {$\mathrm{Α}$};
\draw[gwarmD,thick] (P)--(B) node[midway,right] {$\delta$}; \draw[gwarmD,thick] (P)--(G) node[midway,above right] {$\delta$};
\node[below] at (B) {$\mathrm{Β}$}; \node[left] at (G) {$\Gamma$};
\draw (2.744,0)--(2.744,0.2)--(2.944,0.2);
\draw (1.372,2.377)--(1.545,2.277)--(1.645,2.45);
\end{tikzpicture}
\end{center}

\tmimasmall{Πώς σχεδιάζουμε τη διχοτόμο}

\textbf{Με μοιρογνωμόνιο.} Μετράμε τη γωνία, διαιρούμε το μέτρο με το $2$ και σημειώνουμε αυτή τη γωνία πάνω στη μία πλευρά.

\textbf{Με δίπλωση.} Διπλώνουμε διαφανές χαρτί ώστε οι δύο πλευρές της γωνίας να συμπέσουν. Η πτυχή είναι η διχοτόμος.

\textbf{Με κανόνα και διαβήτη.} Με κέντρο την κορυφή $\mathrm{Ο}$ γράφουμε τόξο που τέμνει τις πλευρές στα $\mathrm{Κ}$ και $\Lambda$. Με κέντρα τα $\mathrm{Κ}$ και $\Lambda$ και την ίδια ακτίνα γράφουμε δύο τόξα που τέμνονται στο $\mathrm{Π}$. Η $\mathrm{ΟΠ}$ είναι η διχοτόμος.

\begin{center}
\begin{tikzpicture}[scale=0.95]
\draw[thick] (5,0)--(0,0)--({5*cos(70)},{5*sin(70)});
\fill (0,0) circle (2pt) node[below left] {$\mathrm{Ο}$};
\draw[gblue] (2,0) arc(0:70:2);
\coordinate (K) at (2,0); \coordinate (L) at ({2*cos(70)},{2*sin(70)});
\fill (K) circle (2pt) node[below] {$\mathrm{Κ}$}; \fill (L) circle (2pt) node[left] {$\Lambda$};
\draw[gwarmD] (K)++(55:1.8) arc(55:95:1.8);
\draw[gwarmD] (L)++(-25:1.8) arc(-25:15:1.8);
\fill (2.478,1.735) circle (2pt) node[above right] {$\mathrm{Π}$};
\draw[gblue,thick] (0,0)--(3.9,2.73);
\end{tikzpicture}
\end{center}

\begin{prosoxi}
\begin{itemize}[leftmargin=1.2em,itemsep=0.3ex,topsep=2pt]
\item Η διχοτόμος είναι ημιευθεία μέσα στη γωνία, με αρχή την κορυφή. Δεν είναι τμήμα που ενώνει δύο σημεία των πλευρών.
\item Η απόσταση σημείου από ευθεία μετριέται πάντα κατά το \emph{κάθετο} τμήμα, όχι κατά ένα λοξό.
\item Σημείο στο εσωτερικό της γωνίας με άνισες αποστάσεις από τις πλευρές δεν ανήκει στη διχοτόμο.
\item Οι δύο μισές γωνίες είναι ίσες σε \emph{μέτρο}, ακόμη κι αν στο σχέδιο οι πλευρές τους έχουν διαφορετικό μήκος.
\end{itemize}
\end{prosoxi}

\tmima{Λυμένα παραδείγματα}

\paradeigma{Η φέτα της πίτσας}
\tixero Η γωνία της κορυφής της φέτας είναι $60^\circ$.\par
\tithelo Με ποια γωνία ξεκινά η κοψιά που τη χωρίζει σε δύο ίσα μέρη;\par
\possindeo Η κοψιά είναι η διχοτόμος της γωνίας, άρα χωρίζει τη γωνία σε δύο ίσες γωνίες.\par
\apantisi $60^\circ:2=30^\circ$ με κάθε πλευρά.\par
\elegxos $30^\circ+30^\circ=60^\circ$.

\paradeigma{Η διχοτόμος και το μέτρο}
\tixero Η $\mathrm{Οδ}$ είναι διχοτόμος της $x\hat{\mathrm{Ο}}y$ και $x\hat{\mathrm{Ο}}\delta=28^\circ$.\par
\tithelo Το μέτρο της $x\hat{\mathrm{Ο}}y$.\par
\possindeo Οι δύο μισές γωνίες είναι ίσες: $\delta\hat{\mathrm{Ο}}y=x\hat{\mathrm{Ο}}\delta=28^\circ$.\par
\apantisi $x\hat{\mathrm{Ο}}y=28^\circ+28^\circ=56^\circ$.\par
\elegxos Η διχοτόμος των $56^\circ$ δίνει δύο γωνίες των $28^\circ$.

\paradeigma{Το ρολόι ως διχοτόμος}
\tixero Στις $4{:}00$ ο λεπτοδείκτης δείχνει το $12$ και ο ωροδείκτης το $4$. Κάθε ώρα είναι $30^\circ$.\par
\tithelo Ποιον αριθμό του ρολογιού δείχνει η διχοτόμος της γωνίας των δεικτών;\par
\possindeo Η γωνία των δεικτών είναι $4\cdot30^\circ=120^\circ$ και η μισή της είναι $60^\circ$, που αντιστοιχεί σε $60^\circ:30^\circ=2$ ώρες από το $12$.\par
\apantisi Η διχοτόμος δείχνει το $2$.\par
\elegxos Από το $12$ ως το $2$ είναι $60^\circ$ και από το $2$ ως το $4$ άλλες $60^\circ$.

\paradeigma{Το σημείο που ισαπέχει}
\tixero Μια γωνία $x\hat{\mathrm{Ο}}y=80^\circ$.\par
\tithelo Να βρω σημείο $\mathrm{Α}$ στο εσωτερικό της που να ισαπέχει από τις πλευρές της και να απέχει $5$ cm από την κορυφή.\par
\possindeo Τα σημεία που ισαπέχουν από τις πλευρές είναι τα σημεία της διχοτόμου. Φέρνω τη διχοτόμο (γωνία $40^\circ$ με κάθε πλευρά) και παίρνω πάνω της σημείο $\mathrm{Α}$ με $\mathrm{ΟΑ}=5$ cm.\par
\apantisi Το $\mathrm{Α}$ είναι πάνω στη διχοτόμο, σε απόσταση $5$ cm από το $\mathrm{Ο}$.\par
\elegxos Μετράω τις κάθετες από το $\mathrm{Α}$ προς τις πλευρές: βγαίνουν ίσες.

\tmima{Ασκήσεις}

\zoni{Το κατάλαβα}{διχοτόμος και ισαπόσταση}

\ex Βρες το μέτρο κάθε μισής γωνίας όταν η γωνία είναι: α) $70^\circ$ \quad β) $124^\circ$ \quad γ) $36^\circ$ \quad δ) $90^\circ$ \quad ε) $180^\circ$.

\ex Σχεδίασε μια γωνία $x\hat{\mathrm{Ο}}y=64^\circ$ και φέρε τη διχοτόμο της: α) με μοιρογνωμόνιο, β) με δίπλωση διαφανούς χαρτιού. Συμπίπτουν;

\ex Η $\mathrm{Οδ}$ είναι διχοτόμος της $x\hat{\mathrm{Ο}}y$ και $x\hat{\mathrm{Ο}}\delta=37^\circ$. Βρες το μέτρο της $x\hat{\mathrm{Ο}}y$.

\ex Σχεδίασε ένα τρίγωνο $\mathrm{ΑΒΓ}$ και φέρε τις διχοτόμους των γωνιών του. Τι παρατηρείς;

\ex Σωστό ή λάθος; α) Η διχοτόμος είναι ημιευθεία. \quad β) Η διχοτόμος χωρίζει τη γωνία σε δύο άνισα μέρη. \quad γ) Κάθε σημείο της διχοτόμου ισαπέχει από τις πλευρές της γωνίας. \quad δ) Η διχοτόμος γωνίας $80^\circ$ σχηματίζει με κάθε πλευρά γωνία $40^\circ$. \quad ε) Η διχοτόμος της ευθείας γωνίας είναι κάθετη στις πλευρές της.

\ex Σχεδίασε μια ευθεία $\varepsilon$ και ένα σημείο $\mathrm{Α}$ εκτός της. Φέρε με τον γνώμονα την κάθετη από το $\mathrm{Α}$ προς την $\varepsilon$ και μέτρησε την απόσταση του $\mathrm{Α}$ από την $\varepsilon$.

\ex Σχεδίασε γωνία $x\hat{\mathrm{Ο}}y=80^\circ$ και πάρε πάνω στη διχοτόμο της ένα σημείο $\mathrm{Α}$ με $\mathrm{ΟΑ}=5$ cm. Μέτρησε τις αποστάσεις του $\mathrm{Α}$ από τις $\mathrm{Ο}x$ και $\mathrm{Ο}y$.

\zoni{Το χρησιμοποιώ}{πίτσες, ρολόγια, δρόμοι}

\ex \textbf{Η πίτσα.} Μια φέτα πίτσα έχει γωνία κορυφής $50^\circ$. Ο Νίκος τη χωρίζει στα δύο με την διχοτόμο. Πόση γωνία έχει καθένα από τα δύο κομμάτια;

\ex \textbf{Το μονοπάτι.} Η γωνία της γωνιάς του φράχτη στον κήπο της θείας Μαρίας είναι $84^\circ$. Ένα μονοπάτι τη χωρίζει σε δύο ίσα παρτέρια. Με ποια γωνία ξεκινά το μονοπάτι από κάθε πλευρά του φράχτη;

\ex \textbf{Το ρολόι.} Στις $8{:}00$, ποιον αριθμό του ρολογιού δείχνει η διχοτόμος της (μικρότερης) γωνίας των δεικτών; Και στις $2{:}00$;

\ex \textbf{Ο ανεμιστήρας.} Ένας ανεμιστήρας της Αγγελικής έχει τρία πτερύγια που σχηματίζουν μεταξύ τους ίσες γωνίες. Πόση γωνία έχει η διχοτόμος δύο διαδοχικών πτερυγίων με καθένα από αυτά;

\ex \textbf{Το περίπτερο.} Δύο ευθείες οδοί τέμνονται σε γωνία $60^\circ$ και ένα περίπτερο πρέπει να ισαπέχει από τις δύο οδούς και να είναι $6$ cm από τη διασταύρωση στο σχέδιο. Σχεδίασε πού μπαίνει.

\zoni{Το σκέφτομαι}{γιατί και πώς}

\ex Γιατί κάθε γωνία έχει μία και μοναδική διχοτόμο; Τι γίνεται με τη διχοτόμο της μηδενικής γωνίας;

\ex Ένα σημείο στο εσωτερικό μιας γωνίας απέχει $3$ cm από τη μία πλευρά και $5$ cm από την άλλη. Ανήκει στη διχοτόμο; Πού να το μετακινήσεις ώστε να ανήκει;

\ex Δύο γωνίες είναι διαδοχικές με άθροισμα $180^\circ$. Η πρώτη είναι $50^\circ$. Σχεδίασε τις διχοτόμους και βρες τη γωνία που σχηματίζουν μεταξύ τους.

\ex Δες τον τρίτο τρόπο, του διαβήτη, για τη διχοτόμο. Σχεδίασε γωνία $70^\circ$ και κατασκεύασε τη διχοτόμο της χωρίς μοιρογνωμόνιο. Έλεγξέ την με μοιρογνωμόνιο.

\zoni{Για το σπίτι}{με χαρτί}

\ex Κόψε από χαρτί μια οξεία και μια αμβλεία γωνία και βρες τη διχοτόμο τους με δίπλωση. Μέτρησε με μοιρογνωμόνιο ότι οι δύο μισές είναι ίσες.

\ex Στο σπίτι σου βρες τρία πράγματα με γωνίες (π.χ. ένα βιβλίο ανοιχτό, ένα ψαλίδι) και πες πού θα περνούσε η διχοτόμος τους.

\paragrafos{Ιδιότητες των τετραπλεύρων}{Β.3.4, «Ιδιότητες παραλληλογράμμου – ορθογωνίου – ρόμβου – τετραγώνου – τραπεζίου – ισοσκελούς τραπεζίου»}

\begin{anagke}
Ο Νίκος έχει ένα μεταλλικό πλαίσιο παραθύρου σε σχήμα ορθογωνίου και θέλει να βεβαιωθεί ότι είναι σωστά στερεωμένο, χωρίς γωνιόμετρο. Η Αγγελική του δείχνει ένα κόλπο: «Μέτρα τις δύο διαγώνιες. Αν είναι ίσες, το πλαίσιο είναι σωστό ορθογώνιο!». Ο Βασίλης δοκιμάζει σε ρόμβο από χαρτόνι: οι διαγώνιες του δεν είναι ίσες, αλλά τέμνονται κάθετα και διχοτομούνται.

Η Μυρτώ διπλώνει ένα χάρτινο παραλληλόγραμμο. Δεν βρίσκει άξονα συμμετρίας, αλλά αν το περιστρέψει μισή στροφή γύρω από το σημείο τομής των διαγωνίων, πέφτει ακριβώς πάνω στον εαυτό του. Η Πολυξένη: «Το κέντρο συμμετρίας του είναι εκεί!».

Σε αυτή την παράγραφο μαθαίνουμε τις ιδιότητες κάθε τετραπλεύρου: πλευρές, γωνίες, διαγώνιες, άξονες και κέντρο συμμετρίας.
\end{anagke}

\begin{lexi}{διαγώνιος, διχοτομούνται}{ονόματα}
Η \lexiword{διαγώνιος} ενώνει δύο μη διαδοχικές κορυφές (διά γωνίας). Δύο τμήματα \lexiword{διχοτομούνται} όταν το καθένα περνά από το μέσο του άλλου.
\end{lexi}

\begin{metafrasi}
Οι διαγώνιες περνούν από το μέσο η μία της άλλης. & διχοτομούνται\\
Οι απέναντι πλευρές είναι ίσες. & $\mathrm{ΑΒ}=\mathrm{ΔΓ}$, $\mathrm{ΑΔ}=\mathrm{ΒΓ}$\\
Οι απέναντι γωνίες είναι ίσες. & $\hat{\mathrm{Α}}=\hat{\Gamma}$, $\hat{\mathrm{Β}}=\hat{\mathrm{Δ}}$\\
Οι διαγώνιες σχηματίζουν ορθή γωνία. & κάθετες: $\mathrm{ΑΓ}\perp\mathrm{ΒΔ}$\\
\end{metafrasi}

\tmima{Η θεωρία}

\begin{kanonas}{παραλληλόγραμμο}
Σε κάθε παραλληλόγραμμο:
\begin{itemize}[leftmargin=1.2em,itemsep=0.3ex,topsep=2pt]
\item το σημείο τομής των διαγωνίων είναι \textbf{κέντρο συμμετρίας},
\item οι διαγώνιες \textbf{διχοτομούνται},
\item οι απέναντι πλευρές είναι \textbf{ίσες},
\item οι απέναντι γωνίες είναι \textbf{ίσες}.
\end{itemize}
\end{kanonas}

\begin{kanonas}{ορθογώνιο}
Εκτός από τις ιδιότητες του παραλληλογράμμου: οι \textbf{διαγώνιες είναι ίσες}, και οι μεσοκάθετοι των πλευρών του είναι \textbf{άξονες συμμετρίας} (δύο άξονες).
\end{kanonas}

\begin{kanonas}{ρόμβος}
Εκτός από τις ιδιότητες του παραλληλογράμμου: οι ευθείες των διαγωνίων είναι \textbf{άξονες συμμετρίας}, οι διαγώνιες είναι \textbf{κάθετες} και \textbf{διχοτόμοι των γωνιών} του.
\end{kanonas}

\begin{kanonas}{τετράγωνο}
Έχει όλες τις ιδιότητες του ορθογωνίου και του ρόμβου: οι διαγώνιες είναι ίσες, κάθετες, διχοτομούνται, διχοτομούν τις γωνίες ($45^\circ$ η καθεμία). Έχει \textbf{$4$ άξονες συμμετρίας} (δύο διαγώνιες και δύο μεσοκάθετοι των πλευρών) και κέντρο συμμετρίας.
\end{kanonas}

\begin{kanonas}{ισοσκελές τραπέζιο}
Η ευθεία που περνά από τα μέσα των βάσεων είναι \textbf{άξονας συμμετρίας} και μεσοκάθετος των βάσεων. Οι \textbf{προσκείμενες στη βάση γωνίες} είναι ίσες.
\end{kanonas}

\begin{center}\small
\begin{tabular}{lcccc}
\toprule
Ιδιότητα & Παραλληλόγραμμο & Ορθογώνιο & Ρόμβος & Τετράγωνο\\
\midrule
Απέναντι πλευρές ίσες & ναι & ναι & ναι & ναι\\
Απέναντι γωνίες ίσες & ναι & ναι & ναι & ναι\\
Διαγώνιες διχοτομούνται & ναι & ναι & ναι & ναι\\
Διαγώνιες ίσες & όχι & ναι & όχι & ναι\\
Διαγώνιες κάθετες & όχι & όχι & ναι & ναι\\
Άξονες συμμετρίας & $0$ & $2$ & $2$ & $4$\\
Κέντρο συμμετρίας & ναι & ναι & ναι & ναι\\
\bottomrule
\end{tabular}
\end{center}

\begin{prosoxi}
\begin{itemize}[leftmargin=1.2em,itemsep=0.3ex,topsep=2pt]
\item Σε ένα (πλάγιο) παραλληλόγραμμο οι διαγώνιες δεν είναι ούτε ίσες ούτε κάθετες, και δεν υπάρχει άξονας συμμετρίας.
\item Ό,τι ισχύει για το παραλληλόγραμμο ισχύει και για το ορθογώνιο, τον ρόμβο και το τετράγωνο, γιατί είναι παραλληλόγραμμα.
\item Οι διαγώνιες του ισοσκελούς τραπεζίου είναι ίσες, αλλά δεν διχοτομούνται.
\end{itemize}
\end{prosoxi}

\tmima{Λυμένα παραδείγματα}

\paradeigma{Το ορθογώνιο παράθυρο}
\tixero Ορθογώνιο παράθυρο με διαστάσεις $90$ cm και $120$ cm.\par
\tithelo Η διαγώνιος και πόσο απέχει το κέντρο συμμετρίας από μια κορυφή.\par
\possindeo Από το πυθαγόρειο: $\mathrm{ΑΓ}=\sqrt{90^2+120^2}=\sqrt{8100+14400}=150$ cm. Οι διαγώνιες διχοτομούνται και είναι ίσες, άρα το κέντρο απέχει $150:2=75$ cm από κάθε κορυφή.\par
\apantisi $150$ cm και $75$ cm.\par
\elegxos Οι τέσσερις αποστάσεις από το κέντρο είναι ίσες.

\paradeigma{Οι διαγώνιες του ρόμβου}
\tixero Ρόμβος με διαγώνιες $10$ και $24$ cm.\par
\tithelo Η πλευρά του.\par
\possindeo Οι διαγώνιες διχοτομούνται και είναι κάθετες: σχηματίζουν ορθογώνιο τρίγωνο με κάθετες $5$ και $12$. Η πλευρά είναι $\sqrt{5^2+12^2}=13$ cm.\par
\apantisi $13$ cm (περίμετρος $52$ cm).\par
\elegxos $5^2+12^2=25+144=169=13^2$.

\paradeigma{Οι γωνίες του παραλληλογράμμου}
\tixero Παραλληλόγραμμο με $\hat{\mathrm{Α}}=65^\circ$.\par
\tithelo Οι άλλες γωνίες.\par
\possindeo Οι απέναντι είναι ίσες: $\hat{\Gamma}=65^\circ$. Οι διαδοχικές είναι παραπληρωματικές (εντός επί τα αυτά): $\hat{\mathrm{Β}}=180^\circ-65^\circ=115^\circ=\hat{\mathrm{Δ}}$.\par
\apantisi $65^\circ,115^\circ,65^\circ,115^\circ$.\par
\elegxos Το άθροισμα είναι $360^\circ$.

\paradeigma{Το ισοσκελές τραπέζιο}
\tixero Ισοσκελές τραπέζιο με μια γωνία στη μεγάλη βάση $70^\circ$.\par
\tithelo Οι άλλες γωνίες.\par
\possindeo Οι προσκείμενες στη βάση γωνίες είναι ίσες: δύο γωνίες $70^\circ$. Οι γωνίες στη μικρή βάση είναι παραπληρωματικές των πλαγίων: $180^\circ-70^\circ=110^\circ$ η καθεμία.\par
\apantisi $70^\circ,70^\circ,110^\circ,110^\circ$.\par
\elegxos $70+70+110+110=360$.

\tmima{Ασκήσεις}

\zoni{Το κατάλαβα}{πλευρές, γωνίες, διαγώνιοι}

\ex Παραλληλόγραμμο έχει πλευρές $8$ cm και $5$ cm. Βρες τις άλλες δύο και την περίμετρο.

\ex Παραλληλόγραμμο έχει γωνία $50^\circ$. Βρες τις άλλες τρεις γωνίες.

\ex Οι διαγώνιες ορθογωνίου είναι $13$ cm. Πόσο απέχει το κέντρο από κάθε κορυφή;

\ex Ρόμβος έχει διαγώνιες $12$ και $16$ cm. Βρες την πλευρά και την περίμετρο.

\ex Πόσους άξονες συμμετρίας έχει το τετράγωνο; Σχεδίασέ τους.

\ex Σχεδίασε ορθογώνιο ΑΒΓΔ με διάμετρο μια διαγώνιο και δες ότι ο κύκλος περνά από όλες τις κορυφές.

\ex Σωστό ή λάθος; α) Οι διαγώνιες του ορθογωνίου είναι ίσες. \quad β) Οι διαγώνιες του ρόμβου είναι ίσες. \quad γ) Οι διαγώνιες του παραλληλογράμμου διχοτομούνται. \quad δ) Το ορθογώνιο έχει $4$ άξονες συμμετρίας.

\zoni{Το χρησιμοποιώ}{κατασκευές και μετρήσεις}

\ex \textbf{Το πλαίσιο.} Ορθογώνιο πλαίσιο είναι $60$ cm $\times$ $80$ cm. Πόση πρέπει να είναι η διαγώνιος για να είναι σωστό ορθογώνιο;

\ex \textbf{Ο χαρταετός.} Ένας χαρταετός σε σχήμα ρόμβου έχει διαγώνιες $60$ και $80$ cm. Πόση είναι η πλευρά του και πόση η περίμετρός του;

\ex \textbf{Το τραπεζάκι.} Ισοσκελές τραπέζιο έχει γωνία $65^\circ$ στη μεγάλη βάση. Βρες τις άλλες γωνίες του.

\ex \textbf{Το οικόπεδο.} Ένα παραλληλόγραμμο οικόπεδο έχει γωνία $72^\circ$. Βρες τις υπόλοιπες γωνίες.

\ex \textbf{Ο γνώμονας.} Ένας τεχνίτης ελέγχει ότι ένα τραπέζι είναι ορθογώνιο μετρώντας τις διαγώνιες. Αρκεί αυτό αν το τραπέζι είναι σίγουρα παραλληλόγραμμο;

\zoni{Το σκέφτομαι}{γιατί συμβαίνει}

\ex Εξήγησε γιατί οι απέναντι πλευρές παραλληλογράμμου είναι ίσες. (Βοήθεια: δες το κέντρο συμμετρίας.)

\ex Σε ποια παραλληλόγραμμα οι διαγώνιες είναι διχοτόμοι των γωνιών; Εξήγησε.

\ex Βρες τους άξονες συμμετρίας του ισοσκελούς τραπεζίου. Εξήγησε γιατί οι προσκείμενες γωνίες βάσης είναι ίσες.

\ex Ποιο σχήμα προκύπτει αν ενώσουμε διαδοχικά τα μέσα των πλευρών ενός ορθογωνίου; (Ένα πείραμα με χαρτί.)

\zoni{Για το σπίτι}{πείραμα με χαρτί}

\ex Κόψε από χαρτί ένα παραλληλόγραμμο και τσάκισέ το στις διαγώνιες. Τι διαπιστώνεις;

\ex Μέτρησε τις διαγώνιες μιας πόρτας ή ενός παραθύρου και έλεγξε αν είναι ίσες.

\paragrafos{Ιδιότητες ισοσκελούς και ισόπλευρου τριγώνου}{Β.3.5, «Ιδιότητες ισοσκελούς και ισόπλευρου τριγώνου»}

\begin{anagke}
Η Αγγελική σχεδιάζει στο τετράδιό της ένα έλατο με τρίγωνα: κάθε τρίγωνο έχει δύο ίσες πλευρές. Τα κόβει με το ψαλίδι και τα διπλώνει στη μέση, και παρατηρεί ότι οι δύο μισές πέφτουν ακριβώς η μία πάνω στην άλλη. «Και οι γωνίες στη βάση είναι ίσες;», ρωτά η Ζωή. «Γιατί να είναι;» Η Πολυξένη το γράφει στο τετράδιό της με κωδικούς, γιατί της αρέσουν τα μοτίβα: ισοσκελές $\to$ ίσες γωνίες βάσης. Ο Ηλίας το μετρά σε τρία διαφορετικά ισοσκελή: τρεις φορές βγαίνουν ίσες.

Η κυρία Αλεξάνδρα ρωτά: «Η πτυχή του χαρτιού είναι και ύψος του τριγώνου; Και διχοτόμος; Και διάμεσος;». Η παρέα δοκιμάζει με ισόπλευρο τρίγωνο και βρίσκει ότι δεν υπάρχει μία πτυχή, αλλά τρεις.

Σε αυτή την παράγραφο μαθαίνουμε γιατί το ισοσκελές τρίγωνο έχει άξονα συμμετρίας, ποιες γωνίες του είναι ίσες, και ότι σε αυτό η διάμεσος, το ύψος και η διχοτόμος προς τη βάση είναι το ίδιο τμήμα. Θα δούμε και τις αντίστοιχες ιδιότητες του ισόπλευρου τριγώνου.
\end{anagke}

\begin{lexi}{βάση, προσκείμενες γωνίες, ταυτίζονται}{ονόματα και ιδιότητες}
Η \lexiword{βάση} του ισοσκελούς είναι η πλευρά που δεν είναι από τις δύο ίσες· στηρίζεται πάνω της σαν στο έδαφος. \lexiword{Προσκείμενες} (προς + κείμαι) λέμε τις γωνίες που «κείτονται πάνω» στην ίδια πλευρά, στα άκρα της. \lexiword{Ταυτίζονται} σημαίνει ότι δύο πράγματα είναι ένα και το ίδιο: δύο τμήματα που ταυτίζονται δεν είναι δύο, είναι ένα.
\end{lexi}

\begin{metafrasi}
Το ισοσκελές τρίγωνο με ίσες πλευρές $\mathrm{ΑΒ}$ και $\mathrm{ΑΓ}$. & $\mathrm{ΑΒ}=\mathrm{ΑΓ}$, βάση η $\mathrm{ΒΓ}$\\
Οι γωνίες στη βάση είναι ίσες. & $\hat{\mathrm{Β}}=\hat{\Gamma}$\\
Η διάμεσος $\mathrm{ΑΜ}$ προς τη βάση είναι και ύψος. & $\mathrm{ΒΜ}=\mathrm{ΜΓ}$ και $\mathrm{ΑΜ}\perp\mathrm{ΒΓ}$\\
Η διάμεσος $\mathrm{ΑΜ}$ είναι και διχοτόμος. & $\hat{\mathrm{Α}}_1=\hat{\mathrm{Α}}_2$\\
Το ισόπλευρο τρίγωνο. & $\mathrm{ΑΒ}=\mathrm{ΒΓ}=\mathrm{ΓΑ}$\\
\end{metafrasi}

\tmima{Η θεωρία}

\tmimasmall{Ισοσκελές τρίγωνο}

Ένα ισοσκελές τρίγωνο $\mathrm{ΑΒΓ}$ έχει $\mathrm{ΑΒ}=\mathrm{ΑΓ}$. Το $\mathrm{Α}$ είναι η \textbf{κορυφή} και η $\mathrm{ΒΓ}$ η \textbf{βάση} του. Αν πάρουμε το μέσο $\mathrm{Μ}$ της βάσης και φέρουμε την $\mathrm{ΑΜ}$, διπλώνοντας το σχήμα κατά την $\mathrm{ΑΜ}$ το $\mathrm{Β}$ πέφτει πάνω στο $\Gamma$.

\begin{center}
\begin{tikzpicture}[scale=0.95]
\draw[thick] (0,0)--(4,0)--(2,3)--cycle;
\draw[gblue,thick,dashed] (2,3.4)--(2,-0.4);
\node[below left] at (0,0) {$\mathrm{Β}$}; \node[below right] at (4,0) {$\Gamma$}; \node[above right] at (2,3) {$\mathrm{Α}$}; \node[below right] at (2,0) {$\mathrm{Μ}$};
\draw (2,0.3)--(2.3,0.3)--(2.3,0);
\draw (0.9,1.5)--(1.1,1.4); \draw (2.9,1.4)--(3.1,1.5); \draw (1,-0.1)--(1,0.1); \draw (3,-0.1)--(3,0.1);
\draw (0.6,0) arc(0:56:0.6); \draw (3.4,0) arc(180:124:0.6);
\end{tikzpicture}
\end{center}

\begin{kanonas}{ιδιότητες ισοσκελούς τριγώνου}
Σε ισοσκελές τρίγωνο $\mathrm{ΑΒΓ}$ με $\mathrm{ΑΒ}=\mathrm{ΑΓ}$:
\begin{itemize}[leftmargin=1.2em,itemsep=0.3ex,topsep=2pt]
\item Η ευθεία της διαμέσου $\mathrm{ΑΜ}$ προς τη βάση είναι \textbf{άξονας συμμετρίας} του τριγώνου.
\item Οι προσκείμενες στη βάση γωνίες είναι \textbf{ίσες}: $\hat{\mathrm{Β}}=\hat{\Gamma}$.
\item Η διάμεσος, το ύψος και η διχοτόμος προς τη βάση \textbf{ταυτίζονται}: είναι το ίδιο τμήμα $\mathrm{ΑΜ}$.
\end{itemize}
\end{kanonas}

Γιατί ισχύουν. Στη συμμετρία ως προς τον άξονα $\mathrm{ΑΜ}$ το $\mathrm{Β}$ γίνεται $\Gamma$, και το $\mathrm{Α}$ μένει στη θέση του. Έτσι η γωνία $\hat{\mathrm{Β}}$ πέφτει πάνω στη $\hat{\Gamma}$ και είναι ίσες. Η $\hat{\mathrm{Α}}_1$ πέφτει πάνω στην $\hat{\mathrm{Α}}_2$, άρα η $\mathrm{ΑΜ}$ είναι διχοτόμος. Και η γωνία $\mathrm{ΑΜ}\hat{\mathrm{Β}}$ πέφτει πάνω στην $\mathrm{ΑΜ}\hat{\Gamma}$, άρα είναι ίσες και μαζί κάνουν ευθεία γωνία: καθεμία είναι $90^\circ$, δηλαδή η $\mathrm{ΑΜ}$ είναι και ύψος.

Η ιδιότητα αυτή συνδέεται με τη μεσοκάθετο: αφού $\mathrm{ΑΒ}=\mathrm{ΑΓ}$, το $\mathrm{Α}$ ισαπέχει από τα $\mathrm{Β}$ και $\Gamma$, άρα βρίσκεται στη μεσοκάθετο της $\mathrm{ΒΓ}$. Και η $\mathrm{ΑΜ}$ είναι ακριβώς αυτή η μεσοκάθετος.

\tmimasmall{Ισόπλευρο τρίγωνο}

Το ισόπλευρο τρίγωνο έχει και τις τρεις πλευρές του ίσες. Αφού κάθε πλευρά του μπορεί να παίξει τον ρόλο της βάσης, όσα ισχύουν για το ισοσκελές ισχύουν και για καθεμία από τις τρεις πλευρές.

\begin{center}
\begin{tikzpicture}[scale=0.95]
\draw[thick] (0,0)--(4,0)--(2,3.464)--cycle;
\draw[gblue,dashed] (2,3.9)--(2,-0.4); \draw[gblue,dashed] (4,0)--(0.5,2.598)--(-0.3,3.17); \draw[gblue,dashed] (0,0)--(3.5,2.598)--(4.3,3.17);
\node[below left] at (0,0) {$\mathrm{Β}$}; \node[below right] at (4,0) {$\Gamma$}; \node[above] at (2,3.9) {$\mathrm{Α}$};
\end{tikzpicture}
\end{center}

\begin{kanonas}{ιδιότητες ισόπλευρου τριγώνου}
\begin{itemize}[leftmargin=1.2em,itemsep=0.3ex,topsep=2pt]
\item Οι ευθείες των τριών διαμέσων είναι \textbf{άξονες συμμετρίας} του τριγώνου: έχει τρεις άξονες συμμετρίας.
\item Όλες οι γωνίες του είναι \textbf{ίσες}. Θα αποδείξουμε στην 3.9 ότι κάθε μία είναι $60^\circ$.
\item Η διάμεσος, το ύψος και η διχοτόμος που αντιστοιχούν σε οποιαδήποτε πλευρά \textbf{ταυτίζονται}.
\end{itemize}
\end{kanonas}

\begin{prosoxi}
\begin{itemize}[leftmargin=1.2em,itemsep=0.3ex,topsep=2pt]
\item Οι ίσες γωνίες του ισοσκελούς είναι οι δύο γωνίες της \emph{βάσης}. Η γωνία της κορυφής είναι γενικά διαφορετική από αυτές.
\item Στο ισοσκελές η διάμεσος προς τη βάση ταυτίζεται με το ύψος και τη διχοτόμο. Η διάμεσος προς μία από τις ίσες πλευρές δεν είναι κατά κανόνα ύψος ούτε διχοτόμος.
\item Το ισοσκελές τρίγωνο που δεν είναι ισόπλευρο έχει έναν μόνο άξονα συμμετρίας, ενώ το ισόπλευρο έχει τρεις.
\item Η ιδιότητα $\hat{\mathrm{Β}}=\hat{\Gamma}$ ισχύει όταν $\mathrm{ΑΒ}=\mathrm{ΑΓ}$. Σε ένα σκαληνό τρίγωνο οι γωνίες δεν είναι ίσες, όσο κι αν στο σχέδιο μοιάζουν.
\end{itemize}
\end{prosoxi}

\tmima{Λυμένα παραδείγματα}

\paradeigma{Οι γωνίες της βάσης}
\tixero Το τρίγωνο $\mathrm{ΑΒΓ}$ είναι ισοσκελές με $\mathrm{ΑΒ}=\mathrm{ΑΓ}$ και $\hat{\mathrm{Β}}=70^\circ$.\par
\tithelo Η γωνία $\hat{\Gamma}$.\par
\possindeo Οι προσκείμενες στη βάση $\mathrm{ΒΓ}$ γωνίες του ισοσκελούς είναι ίσες.\par
\apantisi $\hat{\Gamma}=\hat{\mathrm{Β}}=70^\circ$.\par
\elegxos Σε σχέδιο με μοιρογνωμόνιο μετράμε και τις δύο γωνίες της βάσης: βγαίνουν $70^\circ$.

\paradeigma{Η διάμεσος του ισοσκελούς}
\tixero Το τρίγωνο $\mathrm{ΑΒΓ}$ είναι ισοσκελές με $\mathrm{ΑΒ}=\mathrm{ΑΓ}$. Η $\mathrm{ΑΜ}$ είναι διάμεσος, $\mathrm{ΒΓ}=10$ cm και $\hat{\mathrm{Α}}=50^\circ$.\par
\tithelo Τα $\mathrm{ΒΜ}$, $\mathrm{ΑΜ}\hat{\mathrm{Β}}$ και $\mathrm{ΒΑΜ}$.\par
\possindeo Η διάμεσος προς τη βάση είναι και ύψος και διχοτόμος. Το $\mathrm{Μ}$ είναι μέσο της $\mathrm{ΒΓ}$.\par
\apantisi $\mathrm{ΒΜ}=10:2=5$ cm. Είναι ύψος: $\mathrm{ΑΜ}\hat{\mathrm{Β}}=90^\circ$. Είναι διχοτόμος: $\mathrm{ΒΑΜ}=50^\circ:2=25^\circ$.\par
\elegxos Οι δύο μισές γωνίες $25^\circ+25^\circ=50^\circ$.

\paradeigma{Οι άξονες}
\tixero Ένα ισόπλευρο τρίγωνο πλευράς $6$ cm.\par
\tithelo Πόσους άξονες συμμετρίας έχει και ποιοι είναι;\par
\possindeo Κάθε διάμεσος είναι άξονας συμμετρίας και έχει τρεις διαμέσους.\par
\apantisi Τρεις άξονες: οι ευθείες των τριών διαμέσων, που περνούν από το ίδιο σημείο.\par
\elegxos Αν διπλώσω το τρίγωνο κατά οποιονδήποτε από τους τρεις άξονες, οι δύο μισές συμπίπτουν.

\paradeigma{Ποιο τρίγωνο διπλώνεται;}
\tixero Η Ρήνη διπλώνει ένα χάρτινο τρίγωνο ώστε το $\mathrm{Β}$ να πέσει πάνω στο $\Gamma$, και οι δύο μισές συμπίπτουν ακριβώς.\par
\tithelo Τι είδους τρίγωνο είναι;\par
\possindeo Η πτυχή είναι άξονας συμμετρίας στον οποίο το $\mathrm{Β}$ γίνεται $\Gamma$ και το $\mathrm{Α}$ μένει στη θέση του, άρα $\mathrm{ΑΒ}=\mathrm{ΑΓ}$.\par
\apantisi Ισοσκελές, με κορυφή το $\mathrm{Α}$.\par
\elegxos Η πτυχή περνά από το $\mathrm{Α}$ και από το μέσο της $\mathrm{ΒΓ}$.

\tmima{Ασκήσεις}

\zoni{Το κατάλαβα}{ισοσκελές και ισόπλευρο}

\ex Σωστό ή λάθος; α) Οι προσκείμενες στη βάση γωνίες ισοσκελούς τριγώνου είναι ίσες. \quad β) Σε κάθε ισοσκελές η διάμεσος προς τη βάση είναι και διχοτόμος. \quad γ) Σε κάθε ισόπλευρο τρίγωνο οι ευθείες των υψών είναι και μεσοκάθετοι των πλευρών του. \quad δ) Το ισοσκελές τρίγωνο έχει πάντα τρεις άξονες συμμετρίας. \quad ε) Σε κάθε ισοσκελές η διάμεσος προς τη βάση είναι άξονας συμμετρίας.

\ex Το τρίγωνο $\mathrm{ΑΒΓ}$ είναι ισοσκελές με $\mathrm{ΑΒ}=\mathrm{ΑΓ}$ και η $\mathrm{ΑΜ}$ είναι ύψος. Αν $\mathrm{ΒΓ}=8$ cm και $\hat{\mathrm{Α}}=50^\circ$, βρες τα $\mathrm{ΒΜ}$, $\mathrm{ΑΜ}\hat{\mathrm{Β}}$ και $\mathrm{ΒΑΜ}$.

\ex Στο ισοσκελές τρίγωνο $\mathrm{ΑΒΓ}$ με κορυφή το $\mathrm{Α}$ δίνονται: α) $\hat{\mathrm{Β}}=70^\circ$, β) $\hat{\Gamma}=48^\circ$. Βρες σε κάθε περίπτωση την άλλη γωνία της βάσης και δικαιολόγησε.

\ex Στο ισόπλευρο τρίγωνο $\mathrm{ΑΒΓ}$ η $\mathrm{ΑΔ}$ είναι διάμεσος και $\mathrm{ΑΒ}=6$ cm. Βρες το $\mathrm{ΒΓ}$, το $\mathrm{ΒΔ}$ και τη γωνία $\mathrm{ΑΔ}\hat{\mathrm{Β}}$. Πόσο είναι η γωνία $\mathrm{ΒΑΔ}$ αν ξέρουμε ότι $\hat{\mathrm{Α}}=60^\circ$;

\ex Κατασκεύασε με κανόνα και διαβήτη ένα ισοσκελές τρίγωνο με $\mathrm{ΑΒ}=\mathrm{ΑΓ}=5$ cm και $\mathrm{ΒΓ}=6$ cm. Μέτρησε τις γωνίες της βάσης και βρες τον άξονα συμμετρίας με δίπλωση.

\ex Κατασκεύασε με κανόνα και διαβήτη ένα ισόπλευρο τρίγωνο με πλευρά $5$ cm. Φέρε τους τρεις άξονες συμμετρίας και μέτρησε τις γωνίες του.

\ex Πόσους άξονες συμμετρίας έχει: α) ένα ισοσκελές τρίγωνο που δεν είναι ισόπλευρο, β) ένα ισόπλευρο, γ) ένα σκαληνό;

\zoni{Το χρησιμοποιώ}{στέγες, πανιά, πινακίδες}

\ex \textbf{Η στέγη.} Η στέγη του σπιτιού της γιαγιάς Μυρτώς έχει σχήμα ισοσκελούς τριγώνου με γωνία βάσης $35^\circ$. Πόση είναι η άλλη γωνία της βάσης;

\ex \textbf{Το πανί.} Ένα ισοσκελές τριγωνικό πανί έχει βάση $12$ cm. Η διάμεσος από την κορυφή προς τη βάση έχει μήκος $8$ cm. Πόσο απέχει το πόδι της από κάθε άκρο της βάσης; Τι γωνία σχηματίζει με τη βάση;

\ex \textbf{Η σημαία.} Στο ισοσκελές τρίγωνο $\mathrm{ΑΒΓ}$ με $\mathrm{ΑΒ}=\mathrm{ΑΓ}$ η $\mathrm{ΑΔ}$ είναι η διχοτόμος της γωνίας $\hat{\mathrm{Α}}=64^\circ$. Βρες τη γωνία $\mathrm{ΒΑΔ}$ και τη γωνία $\mathrm{ΑΔ}\hat{\mathrm{Β}}$.

\ex \textbf{Η πινακίδα.} Η πινακίδα «Δώσε προτεραιότητα» έχει σχήμα ισόπλευρου τριγώνου πλευράς $40$ cm. Βρες την περίμετρο και πόσους άξονες συμμετρίας έχει.

\ex \textbf{Η δίπλωση του Στέλιου.} Ο Στέλιος διπλώνει ένα τρίγωνο και οι δύο μισές συμπίπτουν. Λέει: «Άρα το τρίγωνο είναι ισόπλευρο!». Έχει δίκιο; Τι ακριβώς είναι;

\zoni{Το σκέφτομαι}{γιατί συμβαίνει}

\ex Εξήγησε γιατί οι γωνίες στη βάση ενός ισοσκελούς είναι ίσες, με βάση τη δίπλωση.

\ex Γιατί η ευθεία $\mathrm{ΑΜ}$ στο ισοσκελές τρίγωνο είναι η μεσοκάθετος της βάσης $\mathrm{ΒΓ}$; (Σκέψου το $\mathrm{Α}$ και τις αποστάσεις του από τα $\mathrm{Β}$, $\Gamma$.)

\ex Σε ισοσκελές τρίγωνο $\mathrm{ΑΒΓ}$ ($\mathrm{ΑΒ}=\mathrm{ΑΓ}$) φέρε τη διάμεσο προς την πλευρά $\mathrm{ΑΓ}$. Είναι και ύψος; Είναι και διχοτόμος; Έλεγξε με σχέδιο.

\ex Σχεδίασε ένα ορθογώνιο ισοσκελές τρίγωνο με κάθετες πλευρές $4$ cm. Μέτρησε τις δύο άλλες γωνίες. Τι παρατηρείς;

\zoni{Για το σπίτι}{με χαρτί}

\ex Κόψε ένα ισοσκελές τρίγωνο από χαρτί, δίπλωσέ το στη μέση και μέτρησε τις γωνίες της βάσης.

\ex Κόψε ένα ισόπλευρο τρίγωνο από χαρτί και βρες με δίπλωση τους τρεις άξονες συμμετρίας του. Πού τέμνονται;

\paragrafos{Κατακορυφήν, συμπληρωματικές και παραπληρωματικές γωνίες}{Β.3.6, «Κατακορυφήν, συμπληρωματικές και παραπληρωματικές γωνίες»}

\begin{anagke}
Δύο μονοπάτια στον κήπο της θείας Μαρίας διασταυρώνονται σαν ένα μεγάλο «Χ». Η Αγγελική μετράει με το μοιρογνωμόνιο τη μία γωνία και βρίσκει $48^\circ$. «Πόσες είναι οι άλλες τρεις;», ρωτά. Ο Νίκος μετράει όλες: $48^\circ$, $132^\circ$, $48^\circ$, $132^\circ$. «Οι απέναντι είναι ίσες! Και οι διπλανές κάνουν μαζί $180^\circ$». Η Πολυξένη ρωτά: «Είναι τυχαίο; Αν στρέψω τη μία ευθεία, θα αλλάξει;». Δοκιμάζουν με άλλη γωνία, $70^\circ$, και βρίσκουν $70^\circ$, $110^\circ$, $70^\circ$, $110^\circ$.

Η Ρήνη θέλει να φυτέψει τριαντάφυλλα σε μια γωνιά του φράχτη που είναι ορθή. Ένα παρτέρι με λουλούδια παίρνει γωνία $35^\circ$ και ρωτά πόση γωνία μένει για τα τριαντάφυλλα. Ο Στέλιος υπολογίζει: «Εδώ το άθροισμα είναι $90^\circ$, όχι $180^\circ$».

Έτσι η παρέα συναντά τρεις έννοιες: τις γωνίες που αθροίζουν $90^\circ$ (συμπληρωματικές), εκείνες που αθροίζουν $180^\circ$ (παραπληρωματικές) και αυτές που βρίσκονται απέναντι όταν δύο ευθείες τέμνονται (κατακορυφήν). Σε αυτή την παράγραφο θα δούμε και γιατί οι κατακορυφήν είναι ίσες, χωρίς μοιρογνωμόνιο.
\end{anagke}

\begin{lexi}{εφεξής, συμπληρωματική, παραπληρωματική, κατακορυφήν}{γωνίες που σχετίζονται}
Δύο γωνίες είναι \lexiword{εφεξής} όταν η μία έρχεται «από κοντά» στην άλλη. Η \lexiword{συμπληρωματική} γωνία «συμπληρώνει» μια γωνία μέχρι την ορθή. Η \lexiword{παραπληρωματική} «παρα-πληρώνει», δηλαδή συμπληρώνει ως τη μισή στροφή, τις $180^\circ$. Οι \lexiword{κατακορυφήν} γωνίες είναι αυτές που βρίσκονται «κατά κορυφήν», δηλαδή απέναντι, με κοινή κορυφή.
\end{lexi}

\begin{metafrasi}
Δύο γωνίες με κοινή κορυφή και κοινή πλευρά, χωρίς άλλο κοινό σημείο. & εφεξής\\
Δύο γωνίες με άθροισμα $90^\circ$. & συμπληρωματικές: $\hat{\alpha}+\hat{\beta}=90^\circ$\\
Δύο γωνίες με άθροισμα $180^\circ$. & παραπληρωματικές: $\hat{\alpha}+\hat{\beta}=180^\circ$\\
Οι απέναντι γωνίες δύο τεμνόμενων ευθειών. & κατακορυφήν: $\hat{\alpha}=\hat{\gamma}$\\
Η συμπληρωματική της $35^\circ$. & $90^\circ-35^\circ=55^\circ$\\
Η παραπληρωματική της $130^\circ$. & $180^\circ-130^\circ=50^\circ$\\
\end{metafrasi}

\tmima{Η θεωρία}

\tmimasmall{Εφεξής και διαδοχικές γωνίες}

\begin{orismos}{εφεξής γωνίες}
Δύο γωνίες λέγονται \textbf{εφεξής} όταν έχουν κοινή κορυφή, μία κοινή πλευρά και κανένα άλλο κοινό σημείο.
\end{orismos}

Μια σειρά από γωνίες που η καθεμία είναι εφεξής με την επόμενη λέγονται \textbf{διαδοχικές}. Για παράδειγμα στο σχήμα οι γωνίες $\mathrm{ΑΟΒ}$, $\mathrm{ΒΟΓ}$ και $\mathrm{ΓΟΔ}$ είναι διαδοχικές. Το άθροισμά τους είναι η γωνία $\mathrm{ΑΟΔ}$.

\begin{center}
\begin{tikzpicture}[scale=0.9]
\draw[thick] (0,0)--(4.5,0) node[right] {$\mathrm{Α}$};
\draw[thick] (0,0)--({4.2*cos(35)},{4.2*sin(35)}) node[above right] {$\mathrm{Β}$};
\draw[thick] (0,0)--({4.2*cos(80)},{4.2*sin(80)}) node[above] {$\Gamma$};
\draw[thick] (0,0)--({4.2*cos(130)},{4.2*sin(130)}) node[above left] {$\Delta$};
\fill (0,0) circle (2pt) node[below left] {$\mathrm{Ο}$};
\draw[gblue] (1.0,0) arc(0:35:1.0); \draw[gblue] (1.2,0.55) arc(25:80:1.2); \draw[gblue] (0.7,0.5) arc(35:130:0.7);
\end{tikzpicture}
\end{center}

\tmimasmall{Συμπληρωματικές και παραπληρωματικές}

\begin{orismos}{συμπληρωματικές και παραπληρωματικές}
Δύο γωνίες λέγονται \textbf{συμπληρωματικές} όταν το άθροισμά τους είναι $90^\circ$ και \textbf{παραπληρωματικές} όταν το άθροισμά τους είναι $180^\circ$. Κάθε μία από τις δύο λέγεται συμπληρωματική (ή παραπληρωματική) της άλλης.
\end{orismos}

Για παράδειγμα η συμπληρωματική της $35^\circ$ είναι η $55^\circ$ και η παραπληρωματική της $130^\circ$ είναι η $50^\circ$. Δύο γωνίες δεν χρειάζεται να είναι εφεξής για να είναι συμπληρωματικές ή παραπληρωματικές: αρκεί το άθροισμα των μέτρων τους.

\begin{center}
\begin{tikzpicture}[scale=0.8]
\draw[thick] (-0.3,0)--(3.5,0); \draw[thick] (0,-0.3)--(0,3.5); \draw[thick] (0,0)--({3.3*cos(35)},{3.3*sin(35)});
\draw[gblue] (0.8,0) arc(0:35:0.8); \draw[gwarmD] (0.6,0.4) arc(35:90:0.6); \draw (0.3,0)--(0.3,0.3)--(0,0.3);
\node[gblue] at (1.2,0.3) {\small $\hat{\alpha}$}; \node[gwarmD] at (0.4,1.1) {\small $\hat{\beta}$}; \node at (1.7,-0.7) {\small $\hat{\alpha}+\hat{\beta}=90^\circ$};
\begin{scope}[xshift=6.2cm]
\draw[thick] (-2.2,0)--(2.6,0); \draw[thick] (0,0)--({2.3*cos(60)},{2.3*sin(60)});
\draw[gblue] (0.7,0) arc(0:60:0.7); \draw[gwarmD] (-0.55,0) arc(180:60:0.55);
\node[gblue] at (1.05,0.3) {\small $\hat{\alpha}$}; \node[gwarmD] at (-0.3,0.75) {\small $\hat{\beta}$}; \node at (0.3,-0.7) {\small $\hat{\alpha}+\hat{\beta}=180^\circ$};
\end{scope}
\end{tikzpicture}
\end{center}

\tmimasmall{Κατακορυφήν γωνίες}

\begin{orismos}{κατακορυφήν γωνίες}
Δύο γωνίες λέγονται \textbf{κατακορυφήν} όταν έχουν κοινή κορυφή και οι πλευρές τους είναι αντικείμενες ημιευθείες.
\end{orismos}

Όταν δύο ευθείες τέμνονται σχηματίζονται τέσσερις γωνίες, και οι δύο ζεύγη απέναντι γωνιών είναι κατακορυφήν.

\begin{center}
\begin{tikzpicture}[scale=1.1]
\draw[thick,{Stealth}-{Stealth}] (-2.5,-1)--(2.5,1); \draw[thick,{Stealth}-{Stealth}] (-1.2,-1.8)--(1.2,1.8);
\fill (0,0) circle (2pt); \node[below right] at (0.05,0) {$\mathrm{Ο}$};
\node[gblue] at (1.2,0.15) {$\hat{\alpha}$}; \node[gblue] at (-1.2,-0.15) {$\hat{\gamma}$}; \node at (0.2,0.9) {$\hat{\beta}$}; \node at (-0.2,-0.9) {$\hat{\delta}$};
\end{tikzpicture}
\end{center}

\begin{idiotita}{κατακορυφήν γωνίες}
Οι κατακορυφήν γωνίες είναι \textbf{ίσες}.
\end{idiotita}

Γιατί: Η $\hat{\alpha}$ και η $\hat{\beta}$ είναι παραπληρωματικές ($\hat{\alpha}+\hat{\beta}=180^\circ$), και η $\hat{\gamma}$ και η $\hat{\beta}$ επίσης ($\hat{\gamma}+\hat{\beta}=180^\circ$). Άρα $\hat{\alpha}=180^\circ-\hat{\beta}=\hat{\gamma}$.

\begin{prosoxi}
\begin{itemize}[leftmargin=1.2em,itemsep=0.3ex,topsep=2pt]
\item Συμπληρωματικές: $90^\circ$. Παραπληρωματικές: $180^\circ$. Η «παραπληρωματική» είναι η μεγαλύτερη λέξη και το μεγαλύτερο άθροισμα.
\item Δεν υπάρχει συμπληρωματική γωνία για γωνία ίση ή μεγαλύτερη από $90^\circ$ (θα βγαίνε μηδέν ή αρνητικό).
\item Οι διπλανές γωνίες δύο τεμνόμενων ευθειών δεν είναι ίσες (εκτός αν είναι ορθές): είναι παραπληρωματικές.
\item Δύο ίσες γωνίες δεν είναι πάντα κατακορυφήν: αρκεί να έχουν ίδιο μέτρο.
\end{itemize}
\end{prosoxi}

\tmima{Λυμένα παραδείγματα}

\paradeigma{Συμπλήρωμα και παραπλήρωμα}
\tixero Οι γωνίες $35^\circ$ και $130^\circ$.\par
\tithelo Η συμπληρωματική της πρώτης και η παραπληρωματική της δεύτερης.\par
\possindeo Αφαιρώ από το $90^\circ$ και από το $180^\circ$ αντίστοιχα.\par
\apantisi $90^\circ-35^\circ=55^\circ$ και $180^\circ-130^\circ=50^\circ$.\par
\elegxos $35^\circ+55^\circ=90^\circ$ και $130^\circ+50^\circ=180^\circ$.

\paradeigma{Δύο τεμνόμενες ευθείες}
\tixero Δύο ευθείες τέμνονται και μία από τις γωνίες είναι $50^\circ$.\par
\tithelo Οι άλλες τρεις γωνίες.\par
\possindeo Η απέναντί της είναι κατακορυφήν, άρα $50^\circ$. Οι διπλανές της είναι παραπληρωματικές της, δηλαδή $180^\circ-50^\circ=130^\circ$, και η τέταρτη είναι κατακορυφήν της μίας από αυτές, άρα επίσης $130^\circ$.\par
\apantisi $50^\circ$, $130^\circ$ και $130^\circ$.\par
\elegxos Το άθροισμα και των τεσσάρων: $50^\circ+130^\circ+50^\circ+130^\circ=360^\circ$.

\paradeigma{Η συμπληρωματική με το διπλάσιο}
\tixero Μια γωνία είναι συμπληρωματική με μια άλλη που είναι διπλάσια από αυτήν.\par
\tithelo Οι δύο γωνίες.\par
\possindeo Αν η μία είναι $\omega$, η άλλη είναι $2\omega$ και το άθροισμά τους είναι $90^\circ$: $3\omega=90^\circ$, άρα $\omega=30^\circ$.\par
\apantisi $30^\circ$ και $60^\circ$.\par
\elegxos $30^\circ+60^\circ=90^\circ$ και η δεύτερη είναι διπλάσια της πρώτης.

\paradeigma{Όλες οι κατακορυφήν}
\tixero Δύο ευθείες τέμνονται και μία γωνία είναι $\hat{\alpha}=2x+10^\circ$ ... \par
\tithelo ...\par
\possindeo ...\par
\apantisi ...\par
\elegxos ...

\tmima{Ασκήσεις}

\zoni{Το κατάλαβα}{συμπληρωματικές, παραπληρωματικές, κατακορυφήν}

\ex Σωστό ή λάθος; α) Δύο γωνίες με άθροισμα $90^\circ$ λέγονται παραπληρωματικές. \quad β) Δύο κατακορυφήν γωνίες είναι πάντα ίσες. \quad γ) Η συμπληρωματική της μηδενικής γωνίας είναι η ορθή. \quad δ) Η παραπληρωματική μιας αμβλείας είναι πάντα οξεία. \quad ε) Δύο ίσες γωνίες είναι πάντα κατακορυφήν. \quad στ) Δύο εφεξής γωνίες έχουν κοινή κορυφή και μία κοινή πλευρά.

\ex Βρες τη συμπληρωματική των γωνιών: $35^\circ$, $62^\circ$, $8^\circ$, $89^\circ$.

\ex Βρες την παραπληρωματική των γωνιών: $130^\circ$, $90^\circ$, $5^\circ$, $111^\circ$.

\ex Δύο ευθείες τέμνονται και μία γωνία είναι $48^\circ$. Βρες τις άλλες τρεις γωνίες.

\ex Δύο ευθείες $x\mathrm{Ο}x'$ και $y\mathrm{Ο}y'$ τέμνονται στο $\mathrm{Ο}$. Γράψε όλα τα ζεύγη κατακορυφήν γωνιών και όλα τα ζεύγη παραπληρωματικών γωνιών που σχηματίζονται.

\ex Σχεδίασε μια γωνία $40^\circ$ και μετά τη συμπληρωματική της και την παραπληρωματική της.

\ex Τι είδους γωνία είναι: α) η παραπληρωματική μιας οξείας, β) η παραπληρωματική μιας ορθής, γ) η παραπληρωματική μιας αμβλείας, δ) η συμπληρωματική μιας οξείας;

\ex Οι γωνίες $\mathrm{ΑΟΒ}=30^\circ$, $\mathrm{ΒΟΓ}=45^\circ$ και $\mathrm{ΓΟΔ}=50^\circ$ είναι διαδοχικές. Πόση είναι η γωνία $\mathrm{ΑΟΔ}$; Είναι οξεία, ορθή ή αμβλεία;

\zoni{Το χρησιμοποιώ}{διασταυρώσεις, σκάλες, ψαλίδια}

\ex \textbf{Η διασταύρωση.} Δύο μονοπάτια του κήπου τέμνονται και μία γωνία είναι $70^\circ$. Πόσες μοίρες έχουν οι άλλες γωνίες; Πόσες διαφορετικές τιμές εμφανίζονται;

\ex \textbf{Τα τριαντάφυλλα.} Στη γωνιά του φράχτη, που είναι ορθή, φυτεύεται ένα παρτέρι με γωνία $35^\circ$ και το υπόλοιπο φυτεύεται με τριαντάφυλλα. Πόση γωνία μένει για τα τριαντάφυλλα;

\ex \textbf{Η σκάλα.} Μια σκάλα ακουμπά σε έναν τοίχο και σχηματίζει με το έδαφος γωνία $72^\circ$. Πόση γωνία σχηματίζει με τον τοίχο; (Ο τοίχος είναι κάθετος στο έδαφος.)

\ex \textbf{Το ψαλίδι.} Ο παππούς Γιάννης ανοίγει ένα ψαλίδι ώστε οι δύο λεπίδες να σχηματίζουν γωνία $40^\circ$. Πόση γωνία σχηματίζουν οι λαβές του;

\ex \textbf{Η πίτσα.} Μια γωνία είναι ίση με τα $\frac{2}{5}$ της ορθής. Πόσες μοίρες είναι; Βρες τη συμπληρωματική και την παραπληρωματική της.

\zoni{Το σκέφτομαι}{γιατί και αν}

\ex Εξήγησε γιατί οι κατακορυφήν γωνίες είναι ίσες, χρησιμοποιώντας τις παραπληρωματικές γωνίες.

\ex Δύο ίσες γωνίες είναι απαραίτητα κατακορυφήν; Δώσε ένα αντιπαράδειγμα.

\ex Βρες το μέτρο της γωνίας για την οποία ισχύει ότι: α) η συμπληρωματική της είναι κατά $20^\circ$ μεγαλύτερη από αυτήν, β) η παραπληρωματική της είναι κατά $40^\circ$ μικρότερη από αυτήν, γ) η συμπληρωματική της είναι διπλάσια από αυτήν, δ) η παραπληρωματική της είναι τριπλάσια από αυτήν.

\ex Υπάρχει γωνία ίση με τη συμπληρωματική της; Και γωνία ίση με την παραπληρωματική της; Βρες τες.

\zoni{Για το σπίτι}{ψάξε Χ}

\ex Σχεδίασε δύο τεμνόμενες ευθείες, μέτρησε και τις τέσσερις γωνίες και έλεγξε ότι οι κατακορυφήν είναι ίσες.

\ex Βρες στο σπίτι ένα ζευγάρι τεμνόμενων ευθειών (ψαλίδι, παράθυρο, αναδιπλούμενη καρέκλα) και μέτρησε τις γωνίες που σχηματίζουν.

\paragrafos{Σχετικές θέσεις ευθειών στο επίπεδο}{Β.3.7, «Σχετικές θέσεις ευθειών στο επίπεδο»}

\begin{anagke}
Ο παππούς Γιάννης δείχνει πώς θα φυτευτούν οι σειρές με τα καρότα στον κήπο: «Κάθε σειρά θα είναι ευθεία, όλες με την ίδια κατεύθυνση. Καμία δεν θα ακουμπά την άλλη, όσο κι αν τις συνεχίσεις». Η Αγγελική σκέφτεται τις γραμμές στην εθνική οδό και τις ράγες του τρένου: «Δεν συναντιούνται ποτέ, ούτε αν προχωρήσουμε χιλιόμετρα». Ο Ηλίας δείχνει δύο μονοπάτια του κήπου που διασταυρώνονται: «Αυτά ακουμπούν σε ένα σημείο».

«Μα αν συνεχίσεις ένα μονοπάτι που φαινομενικά δεν διασταυρώνεται, μπορεί να διασταυρωθεί αργότερα!», διαμαρτύρεται ο Στέλιος. Η Πολυξένη προτείνει ένα τεστ: «Ας φανταστούμε ότι οι ευθείες δεν τελειώνουν ποτέ, ούτε στο χαρτί. Τότε δύο ευθείες ή έχουν ένα κοινό σημείο ή δεν έχουν κανένα». Η κυρία Αλεξάνδρα προσθέτει μια ερώτηση: «Και πόσες παράλληλες στην ευθεία σας περνούν από ένα σημείο που βρίσκεται δίπλα της;».

Σε αυτή την παράγραφο ονομάζουμε τις δύο δυνατές σχετικές θέσεις δύο ευθειών του επιπέδου, τις τεμνόμενες και τις παράλληλες. Θα μάθουμε και το αίτημα του Ευκλείδη για τις παράλληλες, καθώς και πώς φέρνουμε παράλληλη από ένα σημείο.
\end{anagke}

\begin{lexi}{παράλληλες, τεμνόμενες}{παράλληλες, τεμνόμενες}
Οι \lexiword{παράλληλες} (παρ' άλληλες) πηγαίνουν η μία «δίπλα στην άλλη» χωρίς να συναντιούνται ποτέ. Οι \lexiword{τεμνόμενες} είναι εκείνες που «τέμνονται», δηλαδή «σταυρώνουν». Το κοινό σημείο τους λέγεται \lexiword{σημείο τομής}.
\end{lexi}

\begin{metafrasi}
Δύο ευθείες χωρίς κανένα κοινό σημείο. & παράλληλες: $\varepsilon_1\parallel\varepsilon_2$\\
Δύο ευθείες με ένα κοινό σημείο. & τεμνόμενες\\
Το κοινό τους σημείο. & σημείο τομής $\mathrm{Α}$\\
Τμήματα πάνω σε παράλληλες ευθείες. & $\mathrm{ΑΒ}\parallel\Gamma\Delta$\\
Δύο ευθείες που τέμνονται σε ορθές γωνίες. & κάθετες: $\varepsilon_1\perp\varepsilon_2$\\
\end{metafrasi}

\tmima{Η θεωρία}

\begin{orismos}{παράλληλες και τεμνόμενες ευθείες}
Δύο ευθείες του επιπέδου λέγονται \textbf{παράλληλες} όταν δεν έχουν κανένα κοινό σημείο, όσο κι αν προεκταθούν. Γράφουμε $\varepsilon_1\parallel\varepsilon_2$.\\
Δύο ευθείες του επιπέδου λέγονται \textbf{τεμνόμενες} όταν έχουν ένα μόνο κοινό σημείο. Το σημείο αυτό λέγεται \textbf{σημείο τομής} τους.
\end{orismos}

\begin{center}
\begin{tikzpicture}[scale=0.9]
\draw[thick,{Stealth}-{Stealth}] (0,0)--(4,0); \node[right] at (4,0) {$\varepsilon_2$}; \draw[thick,{Stealth}-{Stealth}] (0,1.3)--(4,1.3); \node[right] at (4,1.3) {$\varepsilon_1$};
\node at (2,-0.6) {$\varepsilon_1\parallel\varepsilon_2$};
\begin{scope}[xshift=7cm]
\draw[thick,{Stealth}-{Stealth}] (0,0)--(4,1.6); \draw[thick,{Stealth}-{Stealth}] (0,1.5)--(4,-0.1); \fill (2,0.77) circle (2pt) node[below] {$\mathrm{Α}$};
\node at (2,-0.6) {τεμνόμενες};
\end{scope}
\end{tikzpicture}
\end{center}

Δύο διαφορετικές ευθείες του ίδιου επιπέδου έχουν δύο μόνο δυνατότητες: ή είναι \emph{τεμνόμενες} ή είναι \emph{παράλληλες}. (Αν είχαν δύο κοινά σημεία, θα ταυτίζονταν, γιατί από δύο σημεία περνά μία μόνο ευθεία.) Δύο τμήματα λέγονται παράλληλα όταν βρίσκονται σε παράλληλες ευθείες.

Οι κάθετες ευθείες είναι μια ειδική περίπτωση τεμνόμενων ευθειών: τέμνονται και σχηματίζουν ορθές γωνίες.

\begin{kanonas}{δύο κάθετες στην ίδια ευθεία}
Δύο ευθείες που είναι κάθετες στην ίδια ευθεία είναι μεταξύ τους \textbf{παράλληλες}.
\end{kanonas}

\begin{center}
\begin{tikzpicture}[scale=0.9]
\draw[thick] (0,0.2)--(5,0.2) node[right] {$\varepsilon$};
\draw[thick] (1,-0.8)--(1,3) node[above] {$\varepsilon_1$}; \draw[thick] (3.6,-0.8)--(3.6,3) node[above] {$\varepsilon_2$};
\draw (1,0.5)--(1.3,0.5)--(1.3,0.2); \draw (3.6,0.5)--(3.9,0.5)--(3.9,0.2);
\node at (2.3,-1.2) {$\varepsilon_1\perp\varepsilon$, $\varepsilon_2\perp\varepsilon$, άρα $\varepsilon_1\parallel\varepsilon_2$};
\end{tikzpicture}
\end{center}

Έτσι φέρνουμε παράλληλη προς μια ευθεία $\zeta$ από ένα σημείο $\mathrm{Α}$ εκτός αυτής: φέρνουμε από το $\mathrm{Α}$ την κάθετη $\mathrm{ΑΒ}$ προς τη $\zeta$ και μετά, στο $\mathrm{Α}$, την κάθετη προς την $\mathrm{ΑΒ}$. Στην πράξη το κάνουμε με γνώμονα και χάρακα: ακουμπάμε τον χάρακα στη δοσμένη ευθεία, τον γνώμονα πάνω στον χάρακα και τον σύρουμε μέχρι να περάσει από το σημείο.

\begin{idiotita}{αίτημα του Ευκλείδη}
Από ένα σημείο εκτός ευθείας διέρχεται \textbf{μία και μόνο μία} ευθεία παράλληλη προς αυτή.
\end{idiotita}

Το αίτημα αυτό το δεχόμαστε χωρίς απόδειξη. Ανήκει στα «Στοιχεία» του Ευκλείδη, που γράφτηκαν στην Αλεξάνδρεια γύρω στο $300$ π.Χ. Από αυτό προκύπτει ότι αν δύο ευθείες είναι παράλληλες προς μια τρίτη, τότε είναι παράλληλες μεταξύ τους: αλλιώς θα είχαν σημείο τομής και από αυτό θα περνούσαν δύο παράλληλες προς την τρίτη.

\tmimasmall{Απόσταση δύο παράλληλων ευθειών}

Αν από ένα σημείο της μίας από δύο παράλληλες ευθείες φέρουμε το κάθετο τμήμα προς την άλλη, το μήκος του είναι η \textbf{απόσταση των παραλλήλων}. Όπου κι αν πάρουμε το σημείο, το κάθετο τμήμα έχει το ίδιο μήκος: οι παράλληλες είναι «ισαπέχουσες».

\begin{center}
\begin{tikzpicture}[scale=0.9]
\draw[thick] (0,0)--(5,0) node[right] {$\varepsilon_2$}; \draw[thick] (0,1.5)--(5,1.5) node[right] {$\varepsilon_1$};
\draw[gblue,thick] (1,0)--(1,1.5) node[midway,left] {\small $\delta$}; \draw[gblue,thick] (3.2,0)--(3.2,1.5) node[midway,right] {\small $\delta$};
\draw (1,0.25)--(1.25,0.25)--(1.25,0); \draw (3.2,0.25)--(3.45,0.25)--(3.45,0);
\end{tikzpicture}
\end{center}

\begin{prosoxi}
\begin{itemize}[leftmargin=1.2em,itemsep=0.3ex,topsep=2pt]
\item Για να είναι παράλληλες δύο ευθείες δεν αρκεί να μην συναντιούνται στο σχέδιο: πρέπει να μην έχουν κοινό σημείο όσο κι αν προεκταθούν. Δύο τμήματα που δεν ακουμπούν μπορεί να βρίσκονται σε τεμνόμενες ευθείες.
\item Δύο παράλληλα τμήματα δεν χρειάζεται να έχουν το ίδιο μήκος.
\item Οι κάθετες είναι τεμνόμενες, αλλά οι τεμνόμενες δεν είναι πάντα κάθετες.
\end{itemize}
\end{prosoxi}

\tmima{Λυμένα παραδείγματα}

\paradeigma{Παράλληλες ή τεμνόμενες;}
\tixero Οι ευθείες: οι ράγες του τρένου, δύο μονοπάτια που διασταυρώνονται, οι γραμμές του τετραδίου.\par
\tithelo Να πω ποιες ευθείες είναι παράλληλες και ποιες τεμνόμενες.\par
\possindeo Παράλληλες είναι όσες δεν έχουν κοινό σημείο, τεμνόμενες όσες έχουν ένα.\par
\apantisi Ράγες και γραμμές του τετραδίου: παράλληλες. Μονοπάτια: τεμνόμενες.\par
\elegxos Οι ράγες δεν θα συναντηθούν ούτε αν προεκταθούν.

\paradeigma{Φέρνω παράλληλη}
\tixero Μια ευθεία $\zeta$ και ένα σημείο $\mathrm{Α}$ εκτός της.\par
\tithelo Η ευθεία $\varepsilon$ από το $\mathrm{Α}$ παράλληλη προς τη $\zeta$.\par
\possindeo Φέρνω από το $\mathrm{Α}$ το κάθετο τμήμα $\mathrm{ΑΒ}$ προς τη $\zeta$. Στο $\mathrm{Α}$ φέρνω την κάθετη στην $\mathrm{ΑΒ}$, δηλαδή την $\varepsilon$. Οι $\zeta$ και $\varepsilon$ είναι κάθετες στην ίδια ευθεία.\par
\apantisi Η $\varepsilon$ είναι παράλληλη προς τη $\zeta$ και είναι η μοναδική που περνά από το $\mathrm{Α}$.\par
\elegxos Η απόσταση των $\varepsilon$ και $\zeta$ είναι $\mathrm{ΑΒ}$, και μετρώντας σε δεύτερο σημείο βρίσκουμε το ίδιο μήκος.

\paradeigma{Δύο παράλληλες σε απόσταση}
\tixero Μια ευθεία $\varepsilon$.\par
\tithelo Ευθείες παράλληλες προς την $\varepsilon$ που απέχουν $3$ cm από αυτήν.\par
\possindeo Παίρνω ένα σημείο $\mathrm{Α}$ της $\varepsilon$ και φέρνω το κάθετο τμήμα $\mathrm{ΑΒ}=3$ cm προς τη μία μεριά και $\mathrm{ΑΒ}'=3$ cm προς την άλλη. Από τα $\mathrm{Β}$, $\mathrm{Β}'$ φέρνω παράλληλες προς την $\varepsilon$.\par
\apantisi Υπάρχουν δύο τέτοιες ευθείες, μία σε κάθε μεριά της $\varepsilon$.\par
\elegxos Και οι δύο έχουν απόσταση $3$ cm από την $\varepsilon$.

\paradeigma{Ευθεία ή τμήμα;}
\tixero Δύο τμήματα $\mathrm{ΑΒ}$ και $\Gamma\Delta$ που δεν ακουμπούν στο χαρτί.\par
\tithelo Είναι παράλληλα;\par
\possindeo Προεκτείνω και τα δύο: αν οι ευθείες τους τέμνονται κάπου εκτός χαρτιού, τα τμήματα δεν είναι παράλληλα.\par
\apantisi Μπορεί να μην είναι: αποφασίζω από τις ευθείες, όχι από τα τμήματα.\par
\elegxos Αν οι προεκτάσεις δεν συναντιούνται, είναι παράλληλα.

\tmima{Ασκήσεις}

\zoni{Το κατάλαβα}{παράλληλες και τεμνόμενες}

\ex Σωστό ή λάθος; α) Δύο ευθείες του επιπέδου που δεν έχουν κανένα κοινό σημείο είναι παράλληλες. \quad β) Δύο τμήματα που δεν έχουν κανένα κοινό σημείο είναι πάντα παράλληλα. \quad γ) Αν δύο ευθείες είναι κάθετες μεταξύ τους, τότε είναι τεμνόμενες. \quad δ) Από ένα σημείο εκτός ευθείας περνά μία μόνο παράλληλη προς αυτήν. \quad ε) Δύο ευθείες κάθετες στην ίδια ευθεία είναι παράλληλες.

\ex Πες αν οι ευθείες είναι παράλληλες ή τεμνόμενες: α) οι δύο γραμμές διαχωρισμού μιας λωρίδας του δρόμου, β) οι ράγες του τρένου, γ) δύο δρόμοι που διασταυρώνονται, δ) οι απέναντι πλευρές ενός θρανίου, ε) τα δύο χέρια του ρολογιού στις $3{:}00$.

\ex Σχεδίασε μια ευθεία $\varepsilon$ και δύο σημεία $\mathrm{Α}$, $\mathrm{Β}$ εκτός της. Από τα $\mathrm{Α}$ και $\mathrm{Β}$ φέρε παράλληλες προς την $\varepsilon$. Τι παρατηρείς; Πότε ταυτίζονται οι δύο ευθείες;

\ex Σχεδίασε μια ευθεία $\varepsilon$ και ένα σημείο $\mathrm{Α}$ εκτός της. Φέρε από το $\mathrm{Α}$ την κάθετη προς την $\varepsilon$ και μετά, στο $\mathrm{Α}$, την κάθετη προς αυτήν. Τι σχέση έχει η δεύτερη κάθετη με την $\varepsilon$;

\ex Σχεδίασε μια γωνία $x\hat{\mathrm{Ο}}y=50^\circ$. Από ένα σημείο $\mathrm{Α}$ της $\mathrm{Ο}x$ φέρε παράλληλη προς την $\mathrm{Ο}y$ και από ένα σημείο $\mathrm{Β}$ της $\mathrm{Ο}y$ παράλληλη προς την $\mathrm{Ο}x$. Αν $\mathrm{Κ}$ είναι το σημείο τομής τους, μέτρησε τη γωνία $\mathrm{ΑΚΒ}$.

\ex Σχεδίασε μια ευθεία $\varepsilon$ και δύο ευθείες που απέχουν $4$ cm από αυτήν και είναι παράλληλες προς αυτήν. Πόσες τέτοιες ευθείες υπάρχουν;

\ex Σχεδίασε δύο τμήματα που δεν τέμνονται πάνω στο χαρτί, αλλά οι ευθείες τους τέμνονται εκτός χαρτιού. Είναι παράλληλα τα τμήματα;

\zoni{Το χρησιμοποιώ}{κήποι, ράγες, πλακάκια}

\ex \textbf{Οι σειρές των καρότων.} Ο παππούς Γιάννης φυτεύει πέντε παράλληλες σειρές καρότα και κάθε δύο διαδοχικές απέχουν $40$ cm. Πόση είναι η απόσταση ανάμεσα στην πρώτη και στην τελευταία σειρά;

\ex \textbf{Οι ράγες.} Οι ράγες του τρένου απέχουν $1{,}5$ μ. σε ένα σημείο. Πόσο απέχουν $2$ χιλιόμετρα πιο κάτω; Εξήγησε.

\ex \textbf{Οι δρόμοι.} Οι οδοί $\varepsilon_1$ και $\varepsilon_2$ είναι κάθετες στην οδό $\varepsilon_3$. Πώς είναι οι $\varepsilon_1$ και $\varepsilon_2$ μεταξύ τους;

\ex \textbf{Τα πλακάκια.} Η θεία Μαρία στρώνει πλακάκια στην ταράτσα. Οι αρμοί προς τη μία κατεύθυνση είναι κάθετοι στους αρμούς της άλλης κατεύθυνσης. Πώς είναι οι αρμοί της ίδιας κατεύθυνσης μεταξύ τους; Εξήγησε.

\ex \textbf{Το τετράδιο.} Πώς θα σχεδιάσεις με χάρακα και γνώμονα μια ευθεία παράλληλη προς μια δοσμένη, που να περνά από ένα δοσμένο σημείο; Γράψε τα βήματα.

\zoni{Το σκέφτομαι}{θέσεις και αιτήματα}

\ex Γιατί δύο διαφορετικές ευθείες του επιπέδου δεν μπορούν να έχουν δύο κοινά σημεία;

\ex Αν $\varepsilon_1\parallel\varepsilon_2$ και $\varepsilon_2\parallel\varepsilon_3$, τι σχέση έχουν οι $\varepsilon_1$ και $\varepsilon_3$; Εξήγησε με το αίτημα του Ευκλείδη.

\ex Από ένα σημείο εκτός ευθείας $\varepsilon$ πόσες ευθείες περνούν που δεν τέμνουν την $\varepsilon$; Και πόσες που την τέμνουν;

\ex Είναι κάθετες οι τεμνόμενες ευθείες; Είναι τεμνόμενες οι κάθετες ευθείες; Δώσε παραδείγματα.

\zoni{Για το σπίτι}{ψάξε παράλληλες}

\ex Βρες στο σπίτι σου πέντε ζεύγη παράλληλων ευθειών (ράφια, πλευρές παραθύρου κ.λπ.) και δύο ζεύγη κάθετων.

\ex Δίπλωσε ένα φύλλο στη μέση. Δίπλωσέ το ξανά ώστε η πτυχή να πέσει πάνω στον εαυτό της. Έτσι έχεις φτιάξει κάθετες πτυχές. Φτιάξε με τον ίδιο τρόπο δεύτερη κάθετη στην ίδια πτυχή. Πώς είναι οι δύο μεταξύ τους;

\paragrafos{Γωνίες σε παράλληλες ευθείες με τέμνουσα}{Β.3.8, «Γωνίες σε παράλληλες ευθείες με τέμνουσα»}

\begin{anagke}
Η θεία Μαρία φυτεύει ντομάτες σε δύο ευθείες σειρές, που είναι παράλληλες, και ένα λοξό μονοπάτι τις διασχίζει και τις δύο. Η Ρήνη, που βοηθά στον κήπο, μετρά με τον μοιρογνωμόνιο τη γωνία ανάμεσα στο μονοπάτι και στην πρώτη σειρά: $62^\circ$. «Να μετρήσω και την άλλη σειρά;», ρωτά. «Δεν χρειάζεται», λέει η κυρία Αλεξάνδρα, που περνούσε τυχαία, «αν ξέρεις τι σημαίνει να είναι παράλληλες, τη βρίσκεις με το μυαλό». Ο Ηλίας μετρά, για να βεβαιωθεί: $62^\circ$. Και στις υπόλοιπες θέσεις; Βγαίνουν συνολικά οκτώ γωνίες και, όπως φαίνεται, μόνο δύο διαφορετικά μέτρα.

Για να μιλήσουμε με ακρίβεια για το ποια γωνία είναι ίση με ποια, χρειαζόμαστε ονόματα: «εντός», «εκτός», «εναλλάξ», «επί τα αυτά».

Σε αυτή την παράγραφο μαθαίνουμε να ονομάζουμε τις οκτώ γωνίες που σχηματίζονται όταν δύο παράλληλες ευθείες τέμνονται από μια τρίτη, και να βρίσκουμε όλες τις γωνίες από μία μόνο.
\end{anagke}

\begin{lexi}{τέμνουσα, εντός, εκτός, εναλλάξ, επί τα αυτά}{θέσεις γωνιών}
Η \lexiword{τέμνουσα} είναι η ευθεία που «κόβει» τις άλλες δύο. Το \lexiword{εντός} σημαίνει «μέσα» και το \lexiword{εκτός} «έξω». Το \lexiword{εναλλάξ} είναι οι θέσεις που «αλλάζουν» πλευρά, δηλαδή αντίθετες. Το \lexiword{επί τα αυτά} σημαίνει «προς τα ίδια μέρη», δηλαδή προς την ίδια πλευρά.
\end{lexi}

\begin{metafrasi}
Η ευθεία που κόβει δύο άλλες ευθείες. & τέμνουσα\\
Μέσα στη λωρίδα ανάμεσα στις παράλληλες. & εντός\\
Έξω από τη λωρίδα των παραλλήλων. & εκτός\\
Στην ίδια πλευρά της τέμνουσας. & επί τα αυτά\\
Σε αντίθετες πλευρές της τέμνουσας. & εναλλάξ\\
\end{metafrasi}

\tmima{Η θεωρία}

Δύο ευθείες $\varepsilon_1$ και $\varepsilon_2$ του ίδιου επιπέδου είναι \textbf{παράλληλες} ($\varepsilon_1\parallel\varepsilon_2$) όταν δεν έχουν κανένα κοινό σημείο, όσο κι αν προεκταθούν. Αν μια τρίτη ευθεία $\varepsilon_3$ τις τέμνει και τις δύο, λέγεται \textbf{τέμνουσά} τους. Τότε σχηματίζονται οκτώ γωνίες, τις οποίες ονομάζουμε $\alpha,\beta,\gamma,\delta$ γύρω από το πρώτο σημείο τομής και $\varepsilon,\zeta,\eta,\theta$ γύρω από το δεύτερο.

\begin{center}
\begin{tikzpicture}[scale=1.05]
\draw[thick] (-3,1.2)--(3,1.2) node[right] {$\varepsilon_1$};
\draw[thick] (-3,-1.2)--(3,-1.2) node[right] {$\varepsilon_2$};
\draw[thick,gblue] (1.4,2.8)--(-1.4,-2.8) node[below left] {$\varepsilon_3$};
\fill (0.6,1.2) circle (1.6pt); \fill (-0.6,-1.2) circle (1.6pt);
\node at (1.05,1.47) {\small $\alpha$}; \node at (0.33,1.65) {\small $\beta$};
\node at (0.15,0.75) {\small $\gamma$}; \node at (1.0,0.75) {\small $\delta$};
\node at (-0.15,-0.93) {\small $\varepsilon$}; \node at (-0.87,-0.75) {\small $\zeta$};
\node at (-1.05,-1.47) {\small $\eta$}; \node at (-0.15,-1.65) {\small $\theta$};
\end{tikzpicture}
\end{center}

\textbf{Ως προς τις παράλληλες:} οι γωνίες $\gamma,\delta,\varepsilon,\zeta$ βρίσκονται ανάμεσα στις $\varepsilon_1,\varepsilon_2$ και λέγονται \textbf{εντός}. Οι $\alpha,\beta,\eta,\theta$ βρίσκονται έξω από τη λωρίδα και λέγονται \textbf{εκτός}.

\textbf{Ως προς την τέμνουσα:} δύο γωνίες που βρίσκονται προς το ίδιο μέρος της $\varepsilon_3$ λέγονται \textbf{επί τα αυτά}. Αν βρίσκονται σε διαφορετικά μέρη, λέγονται \textbf{εναλλάξ}.

Συνδυάζοντας τα δύο, παίρνουμε τα ζεύγη του πίνακα.

\begin{center}
\renewcommand{\arraystretch}{1.25}
\begin{tabular}{@{}p{0.37\linewidth}p{0.29\linewidth}p{0.28\linewidth}@{}}
\toprule
Ζεύγος & Ονομασία & Σχέση\\
\midrule
$\gamma,\varepsilon$ και $\delta,\zeta$ & εντός εναλλάξ & ίσες\\
$\alpha,\eta$ και $\beta,\theta$ & εκτός εναλλάξ & ίσες\\
$\alpha,\varepsilon$ και $\beta,\zeta$ και $\gamma,\eta$ και $\delta,\theta$ & εντός--εκτός και επί τα αυτά & ίσες\\
$\delta,\varepsilon$ και $\gamma,\zeta$ & εντός και επί τα αυτά & παραπληρωματικές\\
$\alpha,\theta$ και $\beta,\eta$ & εκτός και επί τα αυτά & παραπληρωματικές\\
\bottomrule
\end{tabular}
\end{center}

\begin{kanonas}{παράλληλες ευθείες και τέμνουσα}
Αν $\varepsilon_1\parallel\varepsilon_2$, τότε:
\begin{itemize}[leftmargin=1.2em,itemsep=0.3ex,topsep=2pt]
\item οι \textbf{εντός εναλλάξ} γωνίες είναι ίσες,
\item οι \textbf{εκτός εναλλάξ} γωνίες είναι ίσες,
\item οι \textbf{εντός και εκτός επί τα αυτά} γωνίες είναι ίσες,
\item οι \textbf{εντός και επί τα αυτά} γωνίες έχουν άθροισμα $180^\circ$ (είναι παραπληρωματικές), και το ίδιο οι \textbf{εκτός και επί τα αυτά}.
\end{itemize}
Με απλά λόγια: όλες οι \textbf{οξείες} γωνίες του σχήματος είναι ίσες μεταξύ τους, όλες οι \textbf{αμβλείες} είναι ίσες μεταξύ τους, και μια οξεία με μια αμβλεία έχουν άθροισμα $180^\circ$.
\end{kanonas}

Γιατί συμβαίνει αυτό; Όταν δύο ευθείες είναι παράλληλες, η τέμνουσα «γέρνει» προς αυτές με την ίδια κλίση και στις δύο, άρα οι αντίστοιχες γωνίες είναι ίδιες. Οι υπόλοιπες προκύπτουν από τις κατακορυφήν γωνίες (ίσες) και τις παραπληρωματικές (άθροισμα $180^\circ$).

Ισχύει και το \textbf{αντίστροφο}, που μας βοηθά να ελέγχουμε αν δύο ευθείες είναι παράλληλες: αν δύο ευθείες τεμνόμενες από τρίτη σχηματίζουν \emph{εντός εναλλάξ} γωνίες ίσες (ή εντός και επί τα αυτά γωνίες με άθροισμα $180^\circ$), τότε οι ευθείες είναι παράλληλες.

\begin{prosoxi}
\begin{itemize}[leftmargin=1.2em,itemsep=0.3ex,topsep=2pt]
\item Οι ισότητες ισχύουν \emph{μόνο όταν οι ευθείες είναι παράλληλες}. Αν δεν είναι, οι εντός εναλλάξ γωνίες δεν είναι ίσες.
\item Οι εντός και επί τα αυτά γωνίες είναι \emph{παραπληρωματικές}, όχι ίσες (εκτός αν είναι και οι δύο $90^\circ$).
\item Μην μπερδεύεις το «εναλλάξ» με το «επί τα αυτά»: το πρώτο σημαίνει διαφορετικές πλευρές της τέμνουσας, το δεύτερο την ίδια.
\end{itemize}
\end{prosoxi}

\tmima{Λυμένα παραδείγματα}

\paradeigma{Όλες οι γωνίες από μία}
\tixero $\varepsilon_1\parallel\varepsilon_2$ και στο σχήμα της θεωρίας $\gamma=52^\circ$.\par
\tithelo Οι υπόλοιπες επτά γωνίες.\par
\possindeo Η $\alpha$ είναι κατακορυφήν με τη $\gamma$, άρα $\alpha=52^\circ$. Η $\beta$ είναι παραπληρωματική της $\gamma$ (σχηματίζουν ευθεία), άρα $\beta=180^\circ-52^\circ=128^\circ$, και η $\delta=\beta=128^\circ$ (κατακορυφήν). Η $\varepsilon$ είναι εντός εναλλάξ με τη $\gamma$, άρα $\varepsilon=52^\circ$. Η $\zeta$ είναι εντός και επί τα αυτά με τη $\gamma$, άρα $\zeta=180^\circ-52^\circ=128^\circ$. Τέλος $\eta=\varepsilon=52^\circ$ και $\theta=\zeta=128^\circ$ (κατακορυφήν).\par
\apantisi $\alpha=\gamma=\varepsilon=\eta=52^\circ$ και $\beta=\delta=\zeta=\theta=128^\circ$.\par
\elegxos Τέσσερις οξείες ίσες, τέσσερις αμβλείες ίσες, και $52^\circ+128^\circ=180^\circ$.

\paradeigma{Εντός εναλλάξ με άγνωστο}
\tixero $\varepsilon_1\parallel\varepsilon_2$ και δύο εντός εναλλάξ γωνίες είναι $x+20^\circ$ και $2x-40^\circ$.\par
\tithelo Το $x$ και οι δύο γωνίες.\par
\possindeo Οι εντός εναλλάξ γωνίες είναι ίσες, άρα $x+20^\circ=2x-40^\circ$. Προσθέτω $40^\circ$ και στα δύο μέλη: $x+60^\circ=2x$, άρα $x=60^\circ$.\par
\apantisi $x=60^\circ$ και κάθε γωνία είναι $80^\circ$.\par
\elegxos $60^\circ+20^\circ=80^\circ$ και $2\cdot60^\circ-40^\circ=80^\circ$.

\paradeigma{Εντός και επί τα αυτά με άγνωστο}
\tixero $\varepsilon_1\parallel\varepsilon_2$ και δύο εντός και επί τα αυτά γωνίες είναι $4x$ και $x+20^\circ$.\par
\tithelo Οι δύο γωνίες.\par
\possindeo Είναι παραπληρωματικές: $4x+x+20^\circ=180^\circ$, άρα $5x=160^\circ$ και $x=32^\circ$.\par
\apantisi $4x=128^\circ$ και $x+20^\circ=52^\circ$.\par
\elegxos $128^\circ+52^\circ=180^\circ$.

\paradeigma{Παράλληλες ή όχι;}
\tixero Δύο ευθείες τέμνονται από τρίτη και οι εντός εναλλάξ γωνίες είναι $72^\circ$ και $75^\circ$.\par
\tithelo Είναι παράλληλες;\par
\possindeo Αν ήταν παράλληλες, οι εντός εναλλάξ γωνίες θα ήταν ίσες.\par
\apantisi Όχι, γιατί $72^\circ\neq75^\circ$.\par
\elegxos Αν και οι δύο ήταν $72^\circ$, θα ήταν παράλληλες.

\tmima{Ασκήσεις}

\zoni{Το κατάλαβα}{ονόματα και ισότητες}

\ex Στο σχήμα της θεωρίας, πες πώς ονομάζεται κάθε ζεύγος γωνιών: α) $\alpha,\varepsilon$ β) $\gamma,\varepsilon$ γ) $\beta,\theta$ δ) $\delta,\varepsilon$ ε) $\gamma,\zeta$ στ) $\beta,\zeta$.

\ex Οι $\varepsilon_1,\varepsilon_2$ είναι παράλληλες και $\alpha=70^\circ$ στο σχήμα της θεωρίας. Βρες και τις επτά άλλες γωνίες.

\ex Αν $\varepsilon_1\parallel\varepsilon_2$ και $\delta=118^\circ$, βρες τις $\beta,\zeta,\theta,\varepsilon,\alpha$ και εξήγησε με ποια ιδιότητα βρίσκεις την καθεμία.

\ex Δύο εντός εναλλάξ γωνίες παραλλήλων είναι $3x-10^\circ$ και $x+50^\circ$. Βρες το $x$ και τις γωνίες.

\ex Δύο εντός και επί τα αυτά γωνίες παραλλήλων είναι $x$ και $4x$. Βρες τες.

\ex Σωστό ή λάθος; α) Οι εντός εναλλάξ γωνίες είναι πάντα ίσες. β) Αν οι ευθείες είναι παράλληλες, οι εκτός εναλλάξ γωνίες είναι ίσες. γ) Οι εντός και επί τα αυτά γωνίες παραλλήλων είναι ίσες. δ) Δύο κατακορυφήν γωνίες είναι πάντα ίσες, είτε οι ευθείες είναι παράλληλες είτε όχι. ε) Αν οι εντός εναλλάξ γωνίες είναι ίσες, οι ευθείες είναι παράλληλες.

\ex Σχεδίασε δύο παράλληλες ευθείες που απέχουν $3$ cm και μια τέμνουσα που σχηματίζει $55^\circ$ με τη μία. Υπολόγισε όλες τις γωνίες και μέτρησε τες με τον μοιρογνωμόνιο για να ελέγξεις.

\zoni{Το χρησιμοποιώ}{κήπος, δρόμος, σχολείο}

\ex \textbf{Το μονοπάτι.} Στον κήπο της θείας Μαρίας οι δύο σειρές ντομάτας είναι παράλληλες. Μία από τις γωνίες ανάμεσα στο μονοπάτι και μια σειρά είναι $62^\circ$. Πόσες από τις οκτώ γωνίες είναι $62^\circ$ και πόσες είναι $118^\circ$;

\ex \textbf{Ο ποδηλατόδρομος.} Δύο παράλληλες λωρίδες ποδηλατόδρομου τέμνονται από μια διάβαση πεζών που σχηματίζει με την πρώτη λωρίδα μια γωνία $48^\circ$, προς το εσωτερικό του δρόμου. Πόση είναι η εντός εναλλάξ γωνία της στη δεύτερη λωρίδα; Και πόση η εντός και επί τα αυτά;

\ex \textbf{Το κιγκλίδωμα.} Σε ένα κιγκλίδωμα δύο οριζόντιες ράβδοι είναι παράλληλες και μια λοξή ράβδος τις διασχίζει. Μια γωνία της λοξής με την πάνω ράβδο, προς το εσωτερικό, είναι $125^\circ$. Πόση είναι η γωνία της με την κάτω ράβδο, επί τα αυτά μέρη και εντός;

\ex \textbf{Ο φράχτης του Βασίλη.} Ο Βασίλης λέει ότι τα δύο δοκάρια του φράχτη του σχολείου είναι παράλληλα. Μια λοξή σανίδα σχηματίζει με αυτά εντός εναλλάξ γωνίες $74^\circ$ και $76^\circ$. Έχει δίκιο;

\ex \textbf{Ο πίνακας της Φανής.} Η Φανή ελέγχει αν οι δύο απέναντι πλευρές ενός πίνακα είναι παράλληλες. Με μια λοξή ράβδο μετρά εντός και επί τα αυτά γωνίες $98^\circ$ και $82^\circ$. Είναι παράλληλες; Δικαιολόγησε.

\zoni{Το σκέφτομαι}{γιατί συμβαίνει}

\ex Εξήγησε γιατί, αν οι εντός εναλλάξ γωνίες δύο παραλλήλων είναι ίσες, τότε οι εντός και επί τα αυτά είναι παραπληρωματικές.

\ex Στο σχήμα, $\mathrm{ΒΕ}$ και $\mathrm{ΕΓ}$ είναι δύο τμήματα που ξεκινούν από το $\mathrm{Ε}$, με $\mathrm{Β}$ στην $\varepsilon_1$ και $\mathrm{Γ}$ στην $\varepsilon_2$, και $\varepsilon_1\parallel\varepsilon_2$. Η $\mathrm{ΕΒ}$ σχηματίζει με την $\varepsilon_1$ γωνία $40^\circ$ και η $\mathrm{ΕΓ}$ με την $\varepsilon_2$ γωνία $30^\circ$. Βρες τη γωνία $\widehat{\mathrm{ΒΕΓ}}$. (Βοήθεια: φέρε από το $\mathrm{Ε}$ παράλληλη προς τις δύο ευθείες.)

\begin{center}
\begin{tikzpicture}[scale=0.9]
\draw[thick] (-0.8,1.4)--(5.2,1.4) node[right] {$\varepsilon_1$};
\draw[thick] (-0.8,-1.4)--(5.2,-1.4) node[right] {$\varepsilon_2$};
\draw[thick] (3.4,1.4)--(1.73,0)--(4.15,-1.4);
\fill (3.4,1.4) circle (1.6pt) node[above] {$\mathrm{Β}$};
\fill (1.73,0) circle (1.6pt) node[left] {$\mathrm{Ε}$};
\fill (4.15,-1.4) circle (1.6pt) node[below] {$\mathrm{Γ}$};
\node at (2.5,1.15) {\small $40^\circ$};
\node at (3.2,-1.15) {\small $30^\circ$};
\end{tikzpicture}
\end{center}

\ex Στο παραλληλόγραμμο που σχηματίζουν οι παράλληλες $\varepsilon_1\parallel\varepsilon_2$ και $\varepsilon_3\parallel\varepsilon_4$, μια γωνία είναι $70^\circ$. Βρες τις τέσσερις γωνίες του τετραπλεύρου που περικλείεται.

\ex Μια οξεία και μια αμβλεία γωνία που σχηματίζονται από παράλληλες και τέμνουσα είναι εντός και επί τα αυτά, και η αμβλεία είναι τριπλάσια της οξείας. Βρες και τις δύο.

\zoni{Για το σπίτι}{ρίξε μια ματιά}

\ex Βρες στο σπίτι ή στον δρόμο δύο παράλληλες ευθείες (σιδηροτροχιά, κάγκελα, ράφια) που τέμνονται από μια τρίτη. Μέτρησε με τον μοιρογνωμόνιο δύο εντός εναλλάξ γωνίες. Είναι ίσες;

\ex Σχεδίασε με χάρακα και γνώμονα δύο παράλληλες ευθείες και μια τέμνουσα. Βάψε με το ίδιο χρώμα τις γωνίες που είναι ίσες και με άλλο χρώμα τις παραπληρωματικές τους.

\paragrafos{Άθροισμα γωνιών τριγώνου}{Β.3.9, «Άθροισμα γωνιών τριγώνου»}

\begin{anagke}
Η κυρία Αλεξάνδρα ζωγραφίζει στον πίνακα ένα τρίγωνο και ζητά από τον Στέλιο να μετρήσει τις γωνίες του. Εκείνος μετρά δύο γωνίες, $62^\circ$ και $48^\circ$, αλλά το μοιρογνωμόνιο γλιστράει και η τρίτη μέτρηση χάνεται. «Δεν πειράζει», λέει η Φανή, «δεν χρειάζεται να τη μετρήσουμε ξανά». Ο Στέλιος δεν πείθεται: «Και πώς θα τη βρούμε χωρίς να τη μετρήσουμε;». Η Ζωή ρωτά: «Μήπως οι τρεις γωνίες ενός τριγώνου έχουν πάντα το ίδιο άθροισμα, όποιο τρίγωνο κι αν σχεδιάσεις;». Η κυρία Αλεξάνδρα χαμογελά: «Κόψτε ένα τρίγωνο από χαρτί, σκίστε τις τρεις γωνίες και βάλτε τες δίπλα δίπλα με τις κορυφές στο ίδιο σημείο. Τι σχηματίζεται;».

Σε αυτή την παράγραφο μαθαίνουμε ότι το άθροισμα των γωνιών κάθε τριγώνου είναι πάντα το ίδιο, και το χρησιμοποιούμε για να βρίσκουμε γωνίες που δεν μπορούμε ή δεν θέλουμε να μετρήσουμε.
\end{anagke}

\begin{lexi}{άθροισμα, προσκείμενες γωνίες, συμπληρωματικές}{άθροισμα γωνιών}
Το \lexiword{άθροισμα} είναι το «ύψος» που φτάνει ένας σωρός όταν βάζεις τα πράγματα ένα πάνω στο άλλο, δηλαδή το αποτέλεσμα της πρόσθεσης. Οι \lexiword{προσκείμενες γωνίες} μιας πλευράς είναι όσες «κάθονται» πάνω της (έχουν την πλευρά αυτή κοινή). Οι \lexiword{συμπληρωματικές} γωνίες «συμπληρώνουν» η μία την άλλη μέχρι την ορθή γωνία.
\end{lexi}

\begin{metafrasi}
Οι τρεις γωνίες του τριγώνου μαζί. & $\hat{\mathrm{Α}}+\hat{\mathrm{Β}}+\hat{\mathrm{Γ}}$\\
Μια ευθεία γωνία έχει μέτρο. & $180^\circ$\\
Οι γωνίες που κάθονται στη βάση του ισοσκελούς. & προσκείμενες στη βάση, $\hat{\mathrm{Β}}=\hat{\mathrm{Γ}}$\\
Δύο γωνίες που μαζί κάνουν ορθή. & συμπληρωματικές ($90^\circ$)\\
\end{metafrasi}

\tmima{Η θεωρία}

Το πείραμα με το χαρτί δείχνει ότι οι τρεις γωνίες ενός τριγώνου σχηματίζουν μαζί μια ευθεία γωνία. Μπορούμε να το αποδείξουμε με όσα ξέρουμε για τις παράλληλες ευθείες.

\begin{kanonas}{άθροισμα γωνιών τριγώνου}
Το άθροισμα των γωνιών κάθε τριγώνου είναι $180^\circ$:
\[\hat{\mathrm{Α}}+\hat{\mathrm{Β}}+\hat{\mathrm{Γ}}=180^\circ.\]
\end{kanonas}

\begin{center}
\begin{tikzpicture}[scale=0.95]
\draw[thick] (0,0)--(5,0)--(1.6,2.4)--cycle;
\draw[gblue,thick] (-0.8,2.4)--(5.8,2.4) node[right] {$\varepsilon$};
\node[below left] at (0,0) {$\mathrm{Β}$}; \node[below right] at (5,0) {$\mathrm{Γ}$}; \node[above] at (1.6,2.4) {$\mathrm{Α}$};
\fill (1.6,2.4) circle (1.6pt);
\node at (0.8,0.2) {\small $\hat{\mathrm{Β}}$};
\node at (4.2,0.2) {\small $\hat{\mathrm{Γ}}$};
\node at (1.5,1.9) {\small $\hat{\mathrm{Α}}$};
\node at (0.85,2.15) {\small $\hat{\mathrm{Α}}_1$};
\node at (2.9,2.15) {\small $\hat{\mathrm{Α}}_2$};
\end{tikzpicture}
\end{center}

\textbf{Απόδειξη.} Από την κορυφή $\mathrm{Α}$ φέρνουμε την ευθεία $\varepsilon$ παράλληλη προς τη $\mathrm{ΒΓ}$. Τότε $\hat{\mathrm{Α}}_1=\hat{\mathrm{Β}}$ (εντός εναλλάξ, με τέμνουσα την $\mathrm{ΑΒ}$) και $\hat{\mathrm{Α}}_2=\hat{\mathrm{Γ}}$ (εντός εναλλάξ, με τέμνουσα την $\mathrm{ΑΓ}$). Οι γωνίες $\hat{\mathrm{Α}}_1$, $\hat{\mathrm{Α}}$, $\hat{\mathrm{Α}}_2$ σχηματίζουν μαζί ευθεία γωνία, άρα $\hat{\mathrm{Α}}_1+\hat{\mathrm{Α}}+\hat{\mathrm{Α}}_2=180^\circ$. Αντικαθιστώντας, παίρνουμε $\hat{\mathrm{Β}}+\hat{\mathrm{Α}}+\hat{\mathrm{Γ}}=180^\circ$.

Από τον κανόνα προκύπτουν χρήσιμα συμπεράσματα.

\begin{idiotita}{συνέπειες}
\begin{itemize}[leftmargin=1.2em,itemsep=0.3ex,topsep=2pt]
\item Σε \textbf{ισόπλευρο} τρίγωνο όλες οι γωνίες είναι ίσες, άρα κάθε γωνία είναι $180^\circ:3=60^\circ$.
\item Σε \textbf{ορθογώνιο} τρίγωνο οι δύο οξείες γωνίες είναι \textbf{συμπληρωματικές}: $\hat{\mathrm{Β}}+\hat{\mathrm{Γ}}=90^\circ$.
\item Σε \textbf{ορθογώνιο και ισοσκελές} τρίγωνο οι οξείες γωνίες είναι $45^\circ$ και $45^\circ$.
\item Σε \textbf{ισοσκελές} τρίγωνο οι γωνίες της βάσης είναι ίσες, άρα
\[\text{κορυφή}=180^\circ-2\cdot(\text{βάση}),\]
\[\text{βάση}=\bigl(180^\circ-\text{κορυφή}\bigr):2.\]
\item Ένα τρίγωνο έχει \textbf{το πολύ μία} ορθή ή αμβλεία γωνία, και \textbf{τουλάχιστον δύο} οξείες.
\end{itemize}
\end{idiotita}

\begin{prosoxi}
\begin{itemize}[leftmargin=1.2em,itemsep=0.3ex,topsep=2pt]
\item Το άθροισμα είναι πάντα $180^\circ$, όσο μεγάλο ή μικρό κι αν είναι το τρίγωνο. Το μέγεθος δεν αλλάζει το άθροισμα.
\item Αν σε ισοσκελές τρίγωνο σου δίνουν μία γωνία, πρέπει να σκεφτείς αν είναι η γωνία της κορυφής ή της βάσης. Συχνά υπάρχουν \emph{δύο} δυνατές περιπτώσεις.
\item Το ισοσκελές τρίγωνο έχει ίσες μόνο τις γωνίες της βάσης. Και οι τρεις γωνίες είναι ίσες μόνο στο ισόπλευρο.
\end{itemize}
\end{prosoxi}

\tmima{Λυμένα παραδείγματα}

\paradeigma{Η τρίτη γωνία του Στέλιου}
\tixero Σε τρίγωνο δύο γωνίες είναι $62^\circ$ και $48^\circ$.\par
\tithelo Η τρίτη γωνία και το είδος του τριγώνου.\par
\possindeo Το άθροισμα των τριών γωνιών είναι $180^\circ$. Οι δύο γωνίες μαζί κάνουν $62^\circ+48^\circ=110^\circ$, άρα η τρίτη είναι $180^\circ-110^\circ=70^\circ$.\par
\apantisi $70^\circ$. Όλες οι γωνίες είναι οξείες, άρα το τρίγωνο είναι οξυγώνιο.\par
\elegxos $62^\circ+48^\circ+70^\circ=180^\circ$.

\paradeigma{Ισοσκελές με γνωστή κορυφή}
\tixero Ισοσκελές τρίγωνο $\mathrm{ΑΒΓ}$ ($\mathrm{ΑΒ}=\mathrm{ΑΓ}$) με $\hat{\mathrm{Α}}=50^\circ$.\par
\tithelo Οι γωνίες $\hat{\mathrm{Β}}$ και $\hat{\mathrm{Γ}}$.\par
\possindeo Οι γωνίες της βάσης είναι ίσες, $\hat{\mathrm{Β}}=\hat{\mathrm{Γ}}=x$. Τότε $50^\circ+x+x=180^\circ$, άρα $2x=130^\circ$ και $x=65^\circ$.\par
\apantisi $\hat{\mathrm{Β}}=\hat{\mathrm{Γ}}=65^\circ$.\par
\elegxos $50^\circ+65^\circ+65^\circ=180^\circ$.

\paradeigma{Ισοσκελές με μία γνωστή γωνία}
\tixero Ένα ισοσκελές τρίγωνο έχει μια γωνία $80^\circ$.\par
\tithelo Οι άλλες δύο γωνίες.\par
\possindeo Η γωνία των $80^\circ$ μπορεί να είναι η γωνία της κορυφής ή μία γωνία της βάσης.\par
\textbf{1η περίπτωση: κορυφή.} Οι γωνίες της βάσης είναι ίσες: $(180^\circ-80^\circ):2=50^\circ$ η καθεμία.\par
\textbf{2η περίπτωση: βάση.} Η άλλη γωνία της βάσης είναι επίσης $80^\circ$, και η κορυφή είναι $180^\circ-2\cdot80^\circ=20^\circ$.\par
\apantisi $(80^\circ,50^\circ,50^\circ)$ ή $(80^\circ,80^\circ,20^\circ)$.\par
\elegxos Και στις δύο περιπτώσεις το άθροισμα είναι $180^\circ$.

\paradeigma{Γωνίες με άγνωστο}
\tixero Οι γωνίες ενός τριγώνου είναι $x$, $x+20^\circ$ και $2x-20^\circ$.\par
\tithelo Το $x$ και οι γωνίες του τριγώνου.\par
\possindeo Το άθροισμα είναι $180^\circ$: $x+(x+20^\circ)+(2x-20^\circ)=180^\circ$. Στο άθροισμα, τα $+20^\circ$ και $-20^\circ$ αλληλοαναιρούνται και μένει $4x=180^\circ$, άρα $x=45^\circ$.\par
\apantisi Οι γωνίες είναι $45^\circ$, $65^\circ$ και $70^\circ$.\par
\elegxos $45^\circ+65^\circ+70^\circ=180^\circ$. Το τρίγωνο είναι οξυγώνιο.

\paradeigma{Παράλληλες και τρίγωνα}
\tixero Οι $\mathrm{ΑΒ}$ και $\mathrm{ΓΔ}$ είναι παράλληλες και τα τμήματα $\mathrm{ΑΔ}$ και $\mathrm{ΒΓ}$ τέμνονται στο $\mathrm{Ο}$. Είναι $\widehat{\mathrm{ΟΑΒ}}=40^\circ$ και $\widehat{\mathrm{ΟΓΔ}}=35^\circ$.\par
\tithelo Η γωνία $\widehat{\mathrm{ΑΟΒ}}$.\par
\possindeo Η $\widehat{\mathrm{ΟΒΑ}}$ είναι εντός εναλλάξ με τη $\widehat{\mathrm{ΟΓΔ}}$ (παράλληλες $\mathrm{ΑΒ}$, $\mathrm{ΓΔ}$ και τέμνουσα $\mathrm{ΒΓ}$), άρα $\widehat{\mathrm{ΟΒΑ}}=35^\circ$. Στο τρίγωνο $\mathrm{ΟΑΒ}$ είναι $40^\circ+35^\circ+\widehat{\mathrm{ΑΟΒ}}=180^\circ$, άρα $\widehat{\mathrm{ΑΟΒ}}=105^\circ$.\par
\apantisi $\widehat{\mathrm{ΑΟΒ}}=105^\circ$.\par
\elegxos Η $\widehat{\mathrm{ΓΟΔ}}$ είναι κατακορυφήν με την $\widehat{\mathrm{ΑΟΒ}}$, άρα $105^\circ$, και στο τρίγωνο $\mathrm{ΟΓΔ}$ έχουμε $35^\circ+40^\circ+105^\circ=180^\circ$.

\begin{center}
\begin{tikzpicture}[scale=0.85]
\draw[thick] (-1.6,1.34)--(1.92,1.34);
\draw[thick] (-1.92,-1.34)--(1.6,-1.34);
\draw[thick,gblue] (-1.6,1.34)--(1.6,-1.34);
\draw[thick,gblue] (1.92,1.34)--(-1.92,-1.34);
\fill (-1.6,1.34) circle (1.6pt) node[above left] {$\mathrm{Α}$};
\fill (1.92,1.34) circle (1.6pt) node[above right] {$\mathrm{Β}$};
\fill (-1.92,-1.34) circle (1.6pt) node[below left] {$\mathrm{Γ}$};
\fill (1.6,-1.34) circle (1.6pt) node[below right] {$\mathrm{Δ}$};
\fill (0,0) circle (1.6pt) node[above] {\small $\mathrm{Ο}$};
\node at (-0.65,1.1) {\small $40^\circ$};
\node at (-1.1,-1.1) {\small $35^\circ$};
\end{tikzpicture}
\end{center}

\tmima{Ασκήσεις}

\zoni{Το κατάλαβα}{άθροισμα $180^\circ$}

\ex Σωστό ή λάθος; α) Υπάρχει τρίγωνο με δύο ορθές γωνίες. β) Υπάρχει τρίγωνο με γωνίες $50^\circ$, $60^\circ$ και $70^\circ$. γ) Οι οξείες γωνίες ενός ορθογωνίου τριγώνου είναι συμπληρωματικές. δ) Υπάρχει ισόπλευρο τρίγωνο με γωνία $70^\circ$. ε) Αν η γωνία της κορυφής ισοσκελούς τριγώνου είναι $30^\circ$, κάθε γωνία της βάσης είναι $75^\circ$. στ) Ένα τρίγωνο μπορεί να έχει δύο αμβλείες γωνίες.

\ex Βρες την τρίτη γωνία τριγώνου αν οι δύο άλλες είναι: α) $45^\circ$ και $65^\circ$ β) $90^\circ$ και $34^\circ$ γ) $108^\circ$ και $25^\circ$ δ) $59^\circ$ και $61^\circ$.

\ex Σε ισοσκελές τρίγωνο η γωνία της κορυφής είναι: α) $40^\circ$ β) $100^\circ$ γ) $90^\circ$ δ) $20^\circ$. Βρες τις γωνίες της βάσης.

\ex Σε ισοσκελές τρίγωνο η γωνία της βάσης είναι: α) $65^\circ$ β) $30^\circ$ γ) $45^\circ$ δ) $80^\circ$. Βρες τη γωνία της κορυφής.

\ex Οι γωνίες ενός τριγώνου είναι $x$, $x+10^\circ$ και $x+20^\circ$. Βρες το $x$ και τις τρεις γωνίες.

\ex Ποιες από τις παρακάτω τριάδες μπορεί να είναι οι γωνίες ενός τριγώνου; α) $70^\circ,60^\circ,50^\circ$ β) $90^\circ,50^\circ,50^\circ$ γ) $100^\circ,40^\circ,40^\circ$ δ) $120^\circ,30^\circ,40^\circ$ ε) $60^\circ,60^\circ,60^\circ$ στ) $35^\circ,45^\circ,90^\circ$.

\ex Πες αν το τρίγωνο με τις παρακάτω γωνίες είναι οξυγώνιο, ορθογώνιο ή αμβλυγώνιο: α) $50^\circ,60^\circ,70^\circ$ β) $90^\circ,35^\circ,55^\circ$ γ) $105^\circ,40^\circ,35^\circ$ δ) $80^\circ,50^\circ,50^\circ$.

\zoni{Το χρησιμοποιώ}{στέγες, σκάλες, χωράφια}

\ex \textbf{Η στέγη.} Το αέτωμα του εξοχικού της θείας Μαρίας έχει σχήμα ισοσκελούς τριγώνου με γωνία στην κορυφή $110^\circ$. Πόση γωνία σχηματίζει κάθε πλευρά της στέγης με το οριζόντιο δοκάρι της βάσης;

\ex \textbf{Το σημαιάκι.} Η Ρήνη κόβει ένα ισοσκελές τρίγωνο για σημαιάκι και μία γωνία του είναι $50^\circ$. Βρες όλες τις δυνατές τριάδες γωνιών που μπορεί να έχει.

\ex \textbf{Η σκάλα.} Ο Ηλίας ακουμπά μια σκάλα σε κατακόρυφο τοίχο. Η σκάλα σχηματίζει με το οριζόντιο δάπεδο γωνία $72^\circ$. Πόση γωνία σχηματίζει με τον τοίχο;

\ex \textbf{Το χωράφι.} Ένα τριγωνικό χωράφι του παππού Γιάννη έχει γωνίες $2x$, $3x$ και $4x$. Βρες τις γωνίες του και το είδος του τριγώνου ως προς τις γωνίες.

\ex \textbf{Το σχέδιο της Αγγελικής.} Η Αγγελική ζωγραφίζει ισοσκελές τρίγωνο όπου η γωνία της κορυφής είναι τριπλάσια από τη γωνία της βάσης. Βρες τις γωνίες του.

\zoni{Το σκέφτομαι}{γιατί και πώς}

\ex Εξήγησε γιατί κάθε τρίγωνο έχει τουλάχιστον δύο οξείες γωνίες.

\ex Το τρίγωνο $\mathrm{ΑΒΓ}$ είναι ισοσκελές με $\mathrm{ΑΒ}=\mathrm{ΑΓ}$ και $\hat{\mathrm{Α}}=44^\circ$. Η ευθεία $\mathrm{ΔΕ}$ είναι παράλληλη προς τη $\mathrm{ΒΓ}$, με το $\mathrm{Δ}$ στην $\mathrm{ΑΒ}$ και το $\mathrm{Ε}$ στην $\mathrm{ΑΓ}$. Βρες τις γωνίες $\widehat{\mathrm{ΑΔΕ}}$, $\widehat{\mathrm{ΑΕΔ}}$ και $\widehat{\mathrm{ΔΕΓ}}$.

\begin{center}
\begin{tikzpicture}[scale=0.62]
\draw[thick] (0,0)--(4,0)--(2,4.95)--cycle;
\draw[thick,gblue] (0.8,1.98)--(3.2,1.98);
\node[below left] at (0,0) {$\mathrm{Β}$}; \node[below right] at (4,0) {$\mathrm{Γ}$}; \node[above] at (2,4.95) {$\mathrm{Α}$};
\node[left] at (0.8,1.98) {$\mathrm{Δ}$}; \node[right] at (3.2,1.98) {$\mathrm{Ε}$};
\node at (2,4.2) {\small $44^\circ$};
\end{tikzpicture}
\end{center}

\ex Οι $\varepsilon_1$ και $\varepsilon_2$ είναι παράλληλες και μια τέμνουσα τις τέμνει στα $\mathrm{Α}$ και $\mathrm{Β}$. Οι εντός και επί τα αυτά γωνίες είναι $70^\circ$ στο $\mathrm{Α}$ και $110^\circ$ στο $\mathrm{Β}$. Οι διχοτόμοι τους τέμνονται στο $\mathrm{Κ}$. Βρες τη γωνία $\widehat{\mathrm{ΑΚΒ}}$. Θα έβγαινε το ίδιο με άλλες τιμές γωνιών;

\ex Στο τρίγωνο $\mathrm{ΑΒΓ}$ είναι $\hat{\mathrm{Β}}=30^\circ$ και η γωνία $\hat{\mathrm{Α}}$ είναι τετραπλάσια της $\hat{\mathrm{Γ}}$. Βρες τις γωνίες του τριγώνου και το είδος του ως προς τις γωνίες και τις πλευρές του.

\zoni{Για το σπίτι}{ρίξε μια ματιά}

\ex Κόψε από χαρτί ένα τρίγωνο της αρεσκείας σου, χρωμάτισε τις τρεις γωνίες του με διαφορετικά χρώματα και σκίσε τες. Βάλε τες δίπλα δίπλα με τις κορυφές στο ίδιο σημείο. Τι σχηματίζουν; Δοκίμασε και με ένα αμβλυγώνιο τρίγωνο.

\ex Σχεδίασε τρίγωνο $\mathrm{ΑΒΓ}$ με $\mathrm{ΒΓ}=6$ cm, $\hat{\mathrm{Β}}=50^\circ$ και $\hat{\mathrm{Γ}}=60^\circ$. Υπολόγισε τη γωνία $\hat{\mathrm{Α}}$ και μετά μέτρησέ τη. Συμφωνούν;

\paragrafos{Τετράπλευρα, παραλληλόγραμμα, τραπέζια}{Β.3.10, «Τετράπλευρα, παραλληλόγραμμα, τραπέζια»}

\begin{anagke}
Η θεία Μαρία στρώνει πλακάκια στην αυλή και έχει μπροστά της τρεις στοίβες. Τα πλακάκια της πρώτης έχουν τα απέναντι ζευγάρια πλευρών παράλληλα και εφαρμόζουν τέλεια το ένα δίπλα στο άλλο. Της δεύτερης έχουν μόνο ένα ζευγάρι παράλληλων πλευρών και μοιάζουν με μικρά τραπεζάκια. Της τρίτης δεν έχουν καθόλου παράλληλες πλευρές και δύσκολα «κάθονται» με τα άλλα. «Και τα τρία έχουν τέσσερις πλευρές», λέει η Αγγελική, «αλλά δεν είναι το ίδιο σχήμα». Ο Στέλιος προσπαθεί να τα βάλει όλα στην ίδια κατηγορία, η Ρήνη ρωτά ποια ιδιότητα κάνει τα πλακάκια της πρώτης στοίβας τόσο βολικά, και η Ζωή θέλει να ξέρει γιατί.

Σε αυτή την παράγραφο ονομάζουμε τα τετράπλευρα, μελετάμε το παραλληλόγραμμο και το τραπέζιο και μαθαίνουμε τις ιδιότητές τους.
\end{anagke}

\begin{lexi}{τετράπλευρο, παραλληλόγραμμο, τραπέζιο, διαγώνιος}{ονόματα τετραπλεύρων}
Το \lexiword{τετράπλευρο} έχει τέσσερις πλευρές. Το \lexiword{παραλληλόγραμμο} είναι το σχήμα που ορίζεται από «παράλληλες γραμμές». Το \lexiword{τραπέζιο} έρχεται από το «τράπεζα», το τραπέζι, γιατί μοιάζει με τραπεζάκι. Η \lexiword{διαγώνιος} είναι το τμήμα που περνά «δια γωνίας», από γωνία σε γωνία. Ισοσκελές σημαίνει «με ίσα σκέλη».
\end{lexi}

\begin{metafrasi}
Σχήμα με τέσσερις ευθύγραμμες πλευρές. & τετράπλευρο\\
Τα απέναντι ζεύγη πλευρών είναι παράλληλα. & παραλληλόγραμμο ($\mathrm{ΑΒ}\parallel\mathrm{ΓΔ}$, $\mathrm{ΑΔ}\parallel\mathrm{ΒΓ}$)\\
Μόνο ένα ζεύγος παράλληλων πλευρών. & τραπέζιο\\
Τραπέζιο με ίσες τις μη παράλληλες πλευρές. & ισοσκελές τραπέζιο\\
Η απόσταση ανάμεσα στις παράλληλες πλευρές. & ύψος\\
\end{metafrasi}

\tmima{Η θεωρία}

\begin{orismos}{τετράπλευρο}
\textbf{Τετράπλευρο} λέγεται το ευθύγραμμο σχήμα με τέσσερις πλευρές. Έχει τέσσερις κορυφές και τέσσερις γωνίες. Τα τμήματα που ενώνουν απέναντι κορυφές λέγονται \textbf{διαγώνιοι}.
\end{orismos}

\begin{center}
\begin{tikzpicture}[scale=0.8]
\draw[thick] (0.3,2.2)--(3.6,2.6)--(4,0.2)--(0,0)--cycle;
\draw[gblue,dashed] (0.3,2.2)--(4,0.2); \draw[gblue,dashed] (3.6,2.6)--(0,0);
\node[above left] at (0.3,2.2) {$\mathrm{Α}$}; \node[above right] at (3.6,2.6) {$\mathrm{Β}$};
\node[below right] at (4,0.2) {$\mathrm{Γ}$}; \node[below left] at (0,0) {$\mathrm{Δ}$};
\end{tikzpicture}
\end{center}

\begin{kanonas}{άθροισμα γωνιών τετραπλεύρου}
Μια διαγώνιος χωρίζει το τετράπλευρο σε δύο τρίγωνα, καθένα με άθροισμα γωνιών $180^\circ$. Άρα το άθροισμα των γωνιών κάθε τετραπλεύρου είναι $2\cdot180^\circ=360^\circ$.
\end{kanonas}

\begin{orismos}{παραλληλόγραμμο}
\textbf{Παραλληλόγραμμο} λέγεται το τετράπλευρο που έχει τις απέναντι πλευρές του παράλληλες: $\mathrm{ΑΒ}\parallel\mathrm{ΔΓ}$ και $\mathrm{ΑΔ}\parallel\mathrm{ΒΓ}$.
\end{orismos}

Ειδικές περιπτώσεις παραλληλογράμμου είναι το ορθογώνιο, ο ρόμβος και το τετράγωνο, που θα δούμε στην επόμενη παράγραφο.

\begin{center}
\begin{tikzpicture}[scale=0.9]
\draw[thick] (0,0)--(4,0)--(5.2,2)--(1.2,2)--cycle;
\draw[gblue,dashed] (1.2,2)--(4,0); \draw[gblue,dashed] (5.2,2)--(0,0);
\draw[orange!80!black,thick] (1.2,2)--(1.2,0); \draw (1.2,0.2)--(1.4,0.2)--(1.4,0);
\fill (2.6,1) circle (1.6pt) node[above right] {\small $\mathrm{Ο}$};
\node[above left] at (1.2,2) {$\mathrm{Α}$}; \node[above right] at (5.2,2) {$\mathrm{Β}$};
\node[below right] at (4,0) {$\mathrm{Γ}$}; \node[below left] at (0,0) {$\mathrm{Δ}$};
\node[left,orange!80!black] at (1.2,1) {\small $\upsilon$};
\end{tikzpicture}
\end{center}

\begin{idiotita}{παραλληλογράμμου}
Σε κάθε παραλληλόγραμμο $\mathrm{ΑΒΓΔ}$ ισχύουν:
\begin{itemize}[leftmargin=1.2em,itemsep=0.3ex,topsep=2pt]
\item οι \textbf{απέναντι πλευρές} είναι ίσες: $\mathrm{ΑΒ}=\mathrm{ΓΔ}$ και $\mathrm{ΑΔ}=\mathrm{ΒΓ}$,
\item οι \textbf{απέναντι γωνίες} είναι ίσες: $\hat{\mathrm{Α}}=\hat{\mathrm{Γ}}$ και $\hat{\mathrm{Β}}=\hat{\mathrm{Δ}}$,
\item οι \textbf{διαγώνιοι διχοτομούνται}: $\mathrm{ΟΑ}=\mathrm{ΟΓ}$ και $\mathrm{ΟΒ}=\mathrm{ΟΔ}$,
\item οι \textbf{διαδοχικές γωνίες} είναι παραπληρωματικές: $\hat{\mathrm{Α}}+\hat{\mathrm{Δ}}=180^\circ$.
\end{itemize}
\end{idiotita}

Η τελευταία ιδιότητα προκύπτει από όσα ξέρουμε για τις παράλληλες: οι $\hat{\mathrm{Α}}$ και $\hat{\mathrm{Δ}}$ είναι εντός και επί τα αυτά γωνίες των παραλλήλων $\mathrm{ΑΒ}$, $\mathrm{ΔΓ}$ με τέμνουσα την $\mathrm{ΑΔ}$.

Το σημείο τομής $\mathrm{Ο}$ των διαγωνίων λέγεται \textbf{κέντρο} του παραλληλογράμμου. Κάθε πλευρά μπορεί να ονομαστεί \textbf{βάση}, και η απόσταση της βάσης από την απέναντι πλευρά λέγεται \textbf{ύψος} ($\upsilon$ στο σχήμα). Η \textbf{περίμετρος} είναι το άθροισμα των πλευρών, δηλαδή $2\cdot(\mathrm{ΑΒ}+\mathrm{ΒΓ})$.

\begin{orismos}{τραπέζιο}
\textbf{Τραπέζιο} λέγεται το τετράπλευρο που έχει \emph{μόνο δύο} πλευρές παράλληλες. Οι παράλληλες λέγονται \textbf{βάσεις} (η μεγάλη και η μικρή) και η απόστασή τους \textbf{ύψος}. Αν οι άλλες δύο πλευρές είναι ίσες, το τραπέζιο λέγεται \textbf{ισοσκελές}.
\end{orismos}

\begin{center}
\begin{tikzpicture}[scale=0.75]
\draw[thick] (0,0)--(5,0)--(3.6,2)--(1,2)--cycle;
\node[below left] at (0,0) {$\mathrm{Δ}$}; \node[below right] at (5,0) {$\mathrm{Γ}$};
\node[above right] at (3.6,2) {$\mathrm{Β}$}; \node[above left] at (1,2) {$\mathrm{Α}$};
\node at (2.5,-0.7) {\small τραπέζιο ($\mathrm{ΑΒ}\parallel\mathrm{ΔΓ}$)};
\begin{scope}[xshift=7.5cm]
\draw[thick] (0,0)--(5,0)--(4,2)--(1,2)--cycle;
\draw[dashed,gblue] (2.5,2.5)--(2.5,-0.5);
\fill (2.5,2) circle (1.4pt) node[above right] {\small $\mathrm{Μ}_1$}; \fill (2.5,0) circle (1.4pt) node[below right] {\small $\mathrm{Μ}_2$};
\node[below left] at (0,0) {$\mathrm{Δ}$}; \node[below right] at (5,0) {$\mathrm{Γ}$};
\node[above right] at (4,2) {$\mathrm{Β}$}; \node[above left] at (1,2) {$\mathrm{Α}$};
\node at (2.5,-1.05) {\small ισοσκελές τραπέζιο};
\end{scope}
\end{tikzpicture}
\end{center}

\begin{idiotita}{ισοσκελούς τραπεζίου}
Σε ισοσκελές τραπέζιο $\mathrm{ΑΒΓΔ}$ ($\mathrm{ΑΒ}\parallel\mathrm{ΔΓ}$, $\mathrm{ΑΔ}=\mathrm{ΒΓ}$):
\begin{itemize}[leftmargin=1.2em,itemsep=0.3ex,topsep=2pt]
\item οι γωνίες που βρίσκονται στην ίδια βάση είναι ίσες: $\hat{\mathrm{Δ}}=\hat{\mathrm{Γ}}$ και $\hat{\mathrm{Α}}=\hat{\mathrm{Β}}$,
\item η ευθεία των μέσων $\mathrm{Μ}_1$, $\mathrm{Μ}_2$ των βάσεων είναι \textbf{άξονας συμμετρίας} και μεσοκάθετος των δύο βάσεων.
\end{itemize}
\end{idiotita}

Σε κάθε τραπέζιο οι γωνίες που βρίσκονται στην ίδια μη παράλληλη πλευρά είναι παραπληρωματικές (εντός και επί τα αυτά γωνίες των βάσεων): $\hat{\mathrm{Α}}+\hat{\mathrm{Δ}}=180^\circ$ και $\hat{\mathrm{Β}}+\hat{\mathrm{Γ}}=180^\circ$.

\begin{prosoxi}
\begin{itemize}[leftmargin=1.2em,itemsep=0.3ex,topsep=2pt]
\item Το τραπέζιο έχει \emph{ακριβώς} ένα ζεύγος παράλληλων πλευρών. Αν είχε και το δεύτερο, θα ήταν παραλληλόγραμμο.
\item Το ισοσκελές τραπέζιο έχει ίσες τις μη παράλληλες πλευρές, αλλά δεν είναι παραλληλόγραμμο.
\item Οι διαγώνιοι ενός παραλληλογράμμου διχοτομούνται, αλλά δεν είναι γενικά ίσες ούτε κάθετες.
\item Το ύψος είναι \emph{απόσταση}, κάθετη στις βάσεις, και συνήθως δεν είναι πλευρά του σχήματος.
\end{itemize}
\end{prosoxi}

\tmima{Λυμένα παραδείγματα}

\paradeigma{Πλευρές και περίμετρος}
\tixero Παραλληλόγραμμο $\mathrm{ΑΒΓΔ}$ με $\mathrm{ΑΒ}=7$ cm και $\mathrm{ΒΓ}=4$ cm.\par
\tithelo Οι άλλες δύο πλευρές και η περίμετρος.\par
\possindeo Οι απέναντι πλευρές είναι ίσες: $\mathrm{ΓΔ}=\mathrm{ΑΒ}=7$ cm και $\mathrm{ΑΔ}=\mathrm{ΒΓ}=4$ cm. Η περίμετρος είναι $7+4+7+4=22$ cm.\par
\apantisi $\mathrm{ΓΔ}=7$ cm, $\mathrm{ΑΔ}=4$ cm, περίμετρος $22$ cm.\par
\elegxos $2\cdot(7+4)=22$ cm.

\paradeigma{Γωνίες παραλληλογράμμου}
\tixero Παραλληλόγραμμο $\mathrm{ΑΒΓΔ}$ με $\hat{\mathrm{Α}}=65^\circ$.\par
\tithelo Οι άλλες γωνίες.\par
\possindeo Η $\hat{\mathrm{Γ}}$ είναι απέναντι από την $\hat{\mathrm{Α}}$, άρα $\hat{\mathrm{Γ}}=65^\circ$. Οι $\hat{\mathrm{Β}}$ και $\hat{\mathrm{Δ}}$ είναι διαδοχικές της $\hat{\mathrm{Α}}$, άρα $\hat{\mathrm{Β}}=\hat{\mathrm{Δ}}=180^\circ-65^\circ=115^\circ$.\par
\apantisi $65^\circ,115^\circ,65^\circ,115^\circ$.\par
\elegxos $65+115+65+115=360$.

\paradeigma{Οι διαγώνιοι}
\tixero Στο παραλληλόγραμμο $\mathrm{ΑΒΓΔ}$ οι διαγώνιοι τέμνονται στο $\mathrm{Ο}$ και $\mathrm{ΟΑ}=3{,}5$ cm, $\mathrm{ΟΒ}=2{,}4$ cm.\par
\tithelo Τα μήκη των διαγωνίων.\par
\possindeo Οι διαγώνιοι διχοτομούνται, άρα $\mathrm{ΑΓ}=2\cdot3{,}5=7$ cm και $\mathrm{ΒΔ}=2\cdot2{,}4=4{,}8$ cm.\par
\apantisi $\mathrm{ΑΓ}=7$ cm και $\mathrm{ΒΔ}=4{,}8$ cm.\par
\elegxos Το $\mathrm{Ο}$ είναι μέσο και των δύο διαγωνίων.

\paradeigma{Ισοσκελές τραπέζιο}
\tixero Ισοσκελές τραπέζιο $\mathrm{ΑΒΓΔ}$ με $\mathrm{ΑΒ}\parallel\mathrm{ΔΓ}$ και $\hat{\mathrm{Δ}}=70^\circ$.\par
\tithelo Οι άλλες γωνίες.\par
\possindeo Οι γωνίες της βάσης $\mathrm{ΔΓ}$ είναι ίσες: $\hat{\mathrm{Γ}}=70^\circ$. Οι $\hat{\mathrm{Α}}$ και $\hat{\mathrm{Δ}}$ είναι παραπληρωματικές, άρα $\hat{\mathrm{Α}}=180^\circ-70^\circ=110^\circ$, και $\hat{\mathrm{Β}}=\hat{\mathrm{Α}}=110^\circ$.\par
\apantisi $\hat{\mathrm{Α}}=\hat{\mathrm{Β}}=110^\circ$ και $\hat{\mathrm{Γ}}=\hat{\mathrm{Δ}}=70^\circ$.\par
\elegxos $110+110+70+70=360$.

\paradeigma{Τι είδους τετράπλευρο;}
\tixero Στο τετράπλευρο $\mathrm{ΑΒΓΔ}$ είναι $\hat{\mathrm{Α}}=105^\circ$, $\hat{\mathrm{Β}}=100^\circ$ και $\hat{\mathrm{Δ}}=75^\circ$.\par
\tithelo Η $\hat{\mathrm{Γ}}$ και το είδος του τετραπλεύρου.\par
\possindeo Το άθροισμα είναι $360^\circ$, άρα $\hat{\mathrm{Γ}}=360^\circ-(105^\circ+100^\circ+75^\circ)=80^\circ$. Με τέμνουσα την $\mathrm{ΑΔ}$ οι $\hat{\mathrm{Α}}$, $\hat{\mathrm{Δ}}$ έχουν άθροισμα $180^\circ$, άρα $\mathrm{ΑΒ}\parallel\mathrm{ΔΓ}$. Με τέμνουσα την $\mathrm{ΑΒ}$ οι $\hat{\mathrm{Α}}$, $\hat{\mathrm{Β}}$ έχουν άθροισμα $205^\circ\neq180^\circ$, άρα οι $\mathrm{ΑΔ}$ και $\mathrm{ΒΓ}$ δεν είναι παράλληλες.\par
\apantisi $\hat{\mathrm{Γ}}=80^\circ$ και το τετράπλευρο είναι τραπέζιο (όχι παραλληλόγραμμο).\par
\elegxos Στο παραλληλόγραμμο θα ήταν $\hat{\mathrm{Α}}+\hat{\mathrm{Β}}=180^\circ$.

\tmima{Ασκήσεις}

\zoni{Το κατάλαβα}{ονόματα και ιδιότητες}

\ex Σωστό ή λάθος; α) Κάθε παραλληλόγραμμο είναι τετράπλευρο. β) Κάθε τετράπλευρο είναι παραλληλόγραμμο. γ) Το τραπέζιο έχει ακριβώς ένα ζεύγος παράλληλων πλευρών. δ) Οι διαγώνιοι του παραλληλογράμμου διχοτομούνται. ε) Οι διαγώνιοι κάθε παραλληλογράμμου είναι ίσες. στ) Οι απέναντι γωνίες παραλληλογράμμου είναι ίσες.

\ex Σχεδίασε ένα παραλληλόγραμμο $\mathrm{ΑΒΓΔ}$, τις διαγωνίους του και το κέντρο $\mathrm{Ο}$. Μέτρησε τα $\mathrm{ΟΑ}$ και $\mathrm{ΟΓ}$, και φέρε ένα ύψος του.

\ex Στο παραλληλόγραμμο $\mathrm{ΑΒΓΔ}$ είναι $\mathrm{ΑΒ}=9$ cm και $\mathrm{ΒΓ}=5$ cm. Βρες τις $\mathrm{ΓΔ}$, $\mathrm{ΑΔ}$ και την περίμετρο.

\ex Στο παραλληλόγραμμο $\mathrm{ΑΒΓΔ}$ είναι $\hat{\mathrm{Α}}=72^\circ$. Βρες τις άλλες γωνίες.

\ex Οι διαγώνιοι παραλληλογράμμου $\mathrm{ΑΒΓΔ}$ τέμνονται στο $\mathrm{Ο}$ και είναι $\mathrm{ΑΓ}=10$ cm, $\mathrm{ΒΔ}=8$ cm. Βρες τα $\mathrm{ΟΑ}$, $\mathrm{ΟΓ}$, $\mathrm{ΟΒ}$, $\mathrm{ΟΔ}$.

\ex Στο ισοσκελές τραπέζιο $\mathrm{ΑΒΓΔ}$ ($\mathrm{ΑΒ}\parallel\mathrm{ΔΓ}$) είναι $\hat{\mathrm{Δ}}=68^\circ$. Βρες τις άλλες γωνίες.

\ex Στο τραπέζιο $\mathrm{ΑΒΓΔ}$ ($\mathrm{ΑΒ}\parallel\mathrm{ΔΓ}$) είναι $\hat{\mathrm{Α}}=95^\circ$ και $\hat{\mathrm{Β}}=110^\circ$. Βρες τις $\hat{\mathrm{Δ}}$ και $\hat{\mathrm{Γ}}$.

\zoni{Το χρησιμοποιώ}{πλακάκια και ζαρντινιέρες}

\ex \textbf{Το πλακάκι.} Ένα πλακάκι της θείας Μαρίας είναι παραλληλόγραμμο με πλευρές $30$ cm και $20$ cm. Θέλει να κολλήσει ταινία γύρω από την άκρη του. Πόση ταινία χρειάζεται;

\ex \textbf{Η ζαρντινιέρα.} Η πλαϊνή όψη μιας ζαρντινιέρας είναι ισοσκελές τραπέζιο με μικρή βάση $40$ cm, μεγάλη βάση $60$ cm και μη παράλληλες πλευρές $25$ cm. Οι γωνίες στη μεγάλη βάση είναι $70^\circ$. Βρες τις άλλες γωνίες και την περίμετρο.

\ex \textbf{Ο μηχανισμός του Ηλία.} Ο Ηλίας έχει τέσσερις ράβδους που ενώνονται με καρφιά σε παραλληλόγραμμο. Στην αρχή η $\hat{\mathrm{Α}}$ είναι $90^\circ$. Τον πιέζει ώστε να γίνει $\hat{\mathrm{Α}}=60^\circ$. Βρες τις άλλες γωνίες. Αλλάζουν τα μήκη των πλευρών;

\ex \textbf{Η περίμετρος.} Ένα παραλληλόγραμμο οικόπεδο έχει περίμετρο $46$ m και η μια πλευρά του είναι κατά $3$ m μεγαλύτερη από την άλλη. Βρες τις πλευρές του.

\ex \textbf{Το ύψος.} Σχεδίασε παραλληλόγραμμο $\mathrm{ΑΒΓΔ}$ με $\mathrm{ΑΒ}=5$ cm, $\mathrm{ΑΔ}=3$ cm και $\hat{\mathrm{Α}}=60^\circ$. Φέρε το ύψος από το $\mathrm{Δ}$ προς την $\mathrm{ΑΒ}$ και μέτρησέ το.

\zoni{Το σκέφτομαι}{γιατί και πώς}

\ex Εξήγησε γιατί σε κάθε παραλληλόγραμμο $\mathrm{ΑΒΓΔ}$ οι γωνίες $\hat{\mathrm{Α}}$ και $\hat{\mathrm{Δ}}$ έχουν άθροισμα $180^\circ$.

\ex Το τετράπλευρο $\mathrm{ΑΒΓΔ}$ έχει $\hat{\mathrm{Α}}=110^\circ$ και $\hat{\mathrm{Δ}}=65^\circ$. Μπορεί να είναι παραλληλόγραμμο; Δικαιολόγησε.

\ex Στο τετράπλευρο $\mathrm{ΕΖΗΔ}$ είναι $\hat{\mathrm{Δ}}=\hat{\mathrm{Η}}=90^\circ$, $\mathrm{ΔΕ}=4$ cm και $\mathrm{ΗΖ}=7$ cm. Δείξε ότι $\mathrm{ΔΕ}\parallel\mathrm{ΗΖ}$ και ότι το τετράπλευρο είναι τραπέζιο και όχι παραλληλόγραμμο.

\ex Στο ισοσκελές τραπέζιο $\mathrm{ΑΒΓΔ}$ ($\mathrm{ΑΒ}\parallel\mathrm{ΔΓ}$) είναι $\hat{\mathrm{Δ}}=\hat{\mathrm{Γ}}$. Εξήγησε γιατί $\hat{\mathrm{Α}}=\hat{\mathrm{Β}}$.

\zoni{Για το σπίτι}{ρίξε μια ματιά}

\ex Βρες στο σπίτι ή στον δρόμο τρία αντικείμενα που έχουν σχήμα παραλληλογράμμου ή τραπεζίου (παράθυρο, πλακάκι, ράφι, κερκίδα). Σχεδίασέ τα και σημείωσε ποιες πλευρές είναι παράλληλες.

\ex Σχεδίασε παραλληλόγραμμο $\mathrm{ΑΒΓΔ}$ με $\mathrm{ΑΒ}=6$ cm, $\mathrm{ΑΔ}=4$ cm και $\hat{\mathrm{Α}}=70^\circ$. Βρες με υπολογισμό και με μέτρηση τις άλλες γωνίες του.

\paragrafos{Ορθογώνια, ρόμβοι, τετράγωνα}{Β.3.11, «Είδη παραλληλογράμμων και ιδιότητες: ορθογώνια, ρόμβοι, τετράγωνα»}

\begin{anagke}
Η Αγγελική διπλώνει χαρτί σε τρία σχήματα: ένα ορθογώνιο, έναν ρόμβο και ένα τετράγωνο. Για το καθένα ψάχνει με πόσους τρόπους μπορεί να το διπλώσει ώστε τα δύο μισά να κάθονται ακριβώς το ένα πάνω στο άλλο. «Το ορθογώνιο δύο, ο ρόμβος δύο, το τετράγωνο τέσσερις!», λέει. Ο Στέλιος αναφωνεί: «Το τετράγωνο τα έχει όλα, άρα ό,τι ισχύει για τα άλλα ισχύει και για αυτό!». Η Ρήνη όμως ρωτά: «Δηλαδή το τετράγωνο είναι και ορθογώνιο και ρόμβος; Δεν βγάζει νόημα». Η Φανή γράφει στο τετράδιό της και ρωτά τι ακριβώς κάνει κάθε σχήμα «ξεχωριστό».

Σε αυτή την παράγραφο ορίζουμε το ορθογώνιο, τον ρόμβο και το τετράγωνο, βρίσκουμε τις ιδιότητές τους και καταλαβαίνουμε πώς σχετίζονται.
\end{anagke}

\begin{lexi}{ορθογώνιο, ρόμβος, τετράγωνο}{ονόματα παραλληλογράμμων}
Το \lexiword{ορθογώνιο} έχει όλες τις γωνίες \emph{ορθές} (από το «ορθός» και «γωνία»). Ο \lexiword{ρόμβος} ήταν στην αρχαιότητα η σβούρα, που «γυρίζει». Το \lexiword{τετράγωνο} έχει «τέσσερις γωνίες» ίσες, και τέσσερις ίσες πλευρές.
\end{lexi}

\begin{metafrasi}
Παραλληλόγραμμο με όλες τις γωνίες ορθές. & ορθογώνιο\\
Παραλληλόγραμμο με όλες τις πλευρές ίσες. & ρόμβος\\
Ορθογώνιο και ρόμβος μαζί. & τετράγωνο\\
Η ευθεία που διπλώνει το σχήμα σε δύο ίσα μέρη. & άξονας συμμετρίας\\
\end{metafrasi}

\tmima{Η θεωρία}

Τα τρία αυτά σχήματα είναι παραλληλόγραμμα, άρα έχουν όλες τις ιδιότητες του παραλληλογράμμου (ίσες απέναντι πλευρές και γωνίες, διχοτομούμενες διαγώνιες). Έχουν όμως και κάτι παραπάνω.

\begin{orismos}{ορθογώνιο}
\textbf{Ορθογώνιο} (παραλληλόγραμμο) λέγεται το παραλληλόγραμμο που έχει όλες τις γωνίες του ορθές.
\end{orismos}

\begin{idiotita}{ορθογωνίου}
Εκτός από τις ιδιότητες του παραλληλογράμμου, στο ορθογώνιο:
\begin{itemize}[leftmargin=1.2em,itemsep=0.3ex,topsep=2pt]
\item οι \textbf{διαγώνιοι είναι ίσες}: $\mathrm{ΑΓ}=\mathrm{ΒΔ}$,
\item οι \textbf{μεσοκάθετοι των πλευρών} του είναι \textbf{άξονες συμμετρίας} (έχει δύο άξονες).
\end{itemize}
\end{idiotita}

\begin{orismos}{ρόμβος}
\textbf{Ρόμβος} λέγεται το παραλληλόγραμμο που έχει όλες τις πλευρές του ίσες.
\end{orismos}

\begin{idiotita}{ρόμβου}
Εκτός από τις ιδιότητες του παραλληλογράμμου, στον ρόμβο:
\begin{itemize}[leftmargin=1.2em,itemsep=0.3ex,topsep=2pt]
\item οι \textbf{διαγώνιοι τέμνονται κάθετα}: $\mathrm{ΑΓ}\perp\mathrm{ΒΔ}$,
\item οι \textbf{διαγώνιοι διχοτομούν τις γωνίες} του,
\item οι \textbf{ευθείες των διαγωνίων} είναι \textbf{άξονες συμμετρίας} (έχει δύο άξονες).
\end{itemize}
\end{idiotita}

\begin{orismos}{τετράγωνο}
\textbf{Τετράγωνο} λέγεται το παραλληλόγραμμο που έχει όλες τις γωνίες του ορθές και όλες τις πλευρές του ίσες.
\end{orismos}

\begin{idiotita}{τετραγώνου}
Το τετράγωνο είναι ταυτόχρονα \textbf{ορθογώνιο και ρόμβος}, άρα έχει όλες τις ιδιότητές τους. Οι διαγώνιοι του είναι ίσες, κάθετες και διχοτομούν τις γωνίες του (κάθε γωνία της διαγωνίου με την πλευρά είναι $45^\circ$). Έχει \textbf{τέσσερις άξονες συμμετρίας}: τις δύο διαγωνίους και τις δύο μεσοκαθέτους των πλευρών.
\end{idiotita}

\begin{center}
\begin{tikzpicture}[scale=0.62]
\draw[thick] (0,0)--(4,0)--(4,2.6)--(0,2.6)--cycle;
\draw[gblue,dashed] (-0.4,1.3)--(4.4,1.3); \draw[gblue,dashed] (2,-0.4)--(2,3);
\node at (2,-1) {\small ορθογώνιο};
\begin{scope}[xshift=7cm]
\draw[thick] (2,-0.3)--(4,1.3)--(2,2.9)--(0,1.3)--cycle;
\draw[gblue,dashed] (-0.4,1.3)--(4.4,1.3); \draw[gblue,dashed] (2,-0.7)--(2,3.3);
\node at (2,-1) {\small ρόμβος};
\end{scope}
\begin{scope}[xshift=14cm]
\draw[thick] (0.2,0)--(3.8,0)--(3.8,3.6)--(0.2,3.6)--cycle;
\draw[gblue,dashed] (0,0)--(4,3.6); \draw[gblue,dashed] (0,3.6)--(4,0);
\draw[gblue,dashed] (-0.2,1.8)--(4.2,1.8); \draw[gblue,dashed] (2,-0.2)--(2,3.8);
\node at (2,-1) {\small τετράγωνο};
\end{scope}
\end{tikzpicture}
\end{center}

Το διάγραμμα δείχνει πώς συνδέονται όλα τα τετράπλευρα που μάθαμε.

\begin{center}
\begin{tikzpicture}[scale=0.9,node distance=0.8cm,
  box/.style={draw,rounded corners=2pt,fill=gblueL,inner sep=4pt,font=\small,align=center}]
\node[box] (t) at (3,3.6) {τετράπλευρα};
\node[box] (p) at (0.8,2.1) {παραλληλόγραμμα};
\node[box] (tr) at (5.2,2.1) {τραπέζια};
\node[box] (o) at (-0.6,0.6) {ορθογώνια};
\node[box] (r) at (3,0.6) {ρόμβοι};
\node[box] (q) at (1.2,-0.8) {τετράγωνα};
\draw[->,thick] (t)--(p); \draw[->,thick] (t)--(tr);
\draw[->,thick] (p)--(o); \draw[->,thick] (p)--(r);
\draw[->,thick] (o)--(q); \draw[->,thick] (r)--(q);
\end{tikzpicture}
\end{center}

Οι διαγώνιοι χωρίζουν τα σχήματα σε τρίγωνα: στο ορθογώνιο σε τέσσερα ισοσκελή τρίγωνα (γιατί $\mathrm{ΟΑ}=\mathrm{ΟΒ}=\mathrm{ΟΓ}=\mathrm{ΟΔ}$), στον ρόμβο σε τέσσερα ορθογώνια τρίγωνα (γιατί οι διαγώνιοι είναι κάθετες) και στο τετράγωνο σε τέσσερα ορθογώνια και ισοσκελή τρίγωνα. Η περίμετρος του ορθογωνίου είναι $2\cdot(\mathrm{ΑΒ}+\mathrm{ΒΓ})$ και του ρόμβου και του τετραγώνου $4\cdot(\text{πλευρά})$.

Σχετικά με το ύψος: όλες οι πλευρές του ορθογωνίου και του τετραγώνου είναι και ύψη, γιατί είναι κάθετες μεταξύ τους. Στον ρόμβο τα ύψη που φέρνουμε από μια κορυφή είναι ίσα, ενώ σε ένα τυχαίο παραλληλόγραμμο δεν είναι αναγκαστικά ίσα.

\begin{prosoxi}
\begin{itemize}[leftmargin=1.2em,itemsep=0.3ex,topsep=2pt]
\item Κάθε τετράγωνο είναι ρόμβος και ορθογώνιο, αλλά \emph{όχι} κάθε ρόμβος ή ορθογώνιο είναι τετράγωνο.
\item Οι διαγώνιοι του ορθογωνίου είναι ίσες αλλά \emph{όχι κάθετες}, ενώ του ρόμβου είναι κάθετες αλλά \emph{όχι ίσες}.
\item Το ορθογώνιο έχει δύο άξονες συμμετρίας και όχι τέσσερις· οι διαγώνιοί του \emph{δεν} είναι άξονες.
\item Ο ρόμβος δεν έχει ορθές γωνίες, εκτός αν είναι τετράγωνο.
\end{itemize}
\end{prosoxi}

\tmima{Λυμένα παραδείγματα}

\paradeigma{Το ορθογώνιο}
\tixero Ορθογώνιο $\mathrm{ΑΒΓΔ}$ με $\mathrm{ΑΒ}=8$ cm, $\mathrm{ΒΓ}=6$ cm και $\mathrm{ΑΓ}=10$ cm.\par
\tithelo Η περίμετρος, η $\mathrm{ΒΔ}$ και τα τμήματα $\mathrm{ΟΑ}$, $\mathrm{ΟΒ}$ ($\mathrm{Ο}$ το κέντρο).\par
\possindeo Η περίμετρος είναι $2\cdot(8+6)=28$ cm. Οι διαγώνιοι του ορθογωνίου είναι ίσες, άρα $\mathrm{ΒΔ}=\mathrm{ΑΓ}=10$ cm. Διχοτομούνται, άρα $\mathrm{ΟΑ}=\mathrm{ΟΒ}=10:2=5$ cm.\par
\apantisi Περίμετρος $28$ cm, $\mathrm{ΒΔ}=10$ cm, $\mathrm{ΟΑ}=\mathrm{ΟΒ}=5$ cm.\par
\elegxos Το τρίγωνο $\mathrm{ΟΑΒ}$ είναι ισοσκελές ($\mathrm{ΟΑ}=\mathrm{ΟΒ}$).

\paradeigma{Ο ρόμβος και οι γωνίες του}
\tixero Ρόμβος $\mathrm{ΑΒΓΔ}$ με $\hat{\mathrm{Α}}=40^\circ$.\par
\tithelo Οι γωνίες του και η γωνία $\widehat{\mathrm{ΑΒΟ}}$, όπου $\mathrm{Ο}$ το σημείο τομής των διαγωνίων.\par
\possindeo Η $\hat{\mathrm{Γ}}=40^\circ$ (απέναντι) και $\hat{\mathrm{Β}}=\hat{\mathrm{Δ}}=180^\circ-40^\circ=140^\circ$. Η διαγώνιος $\mathrm{ΑΓ}$ διχοτομεί την $\hat{\mathrm{Α}}$, άρα $\widehat{\mathrm{ΟΑΒ}}=20^\circ$. Οι διαγώνιοι είναι κάθετες, άρα $\widehat{\mathrm{ΑΟΒ}}=90^\circ$. Στο τρίγωνο $\mathrm{ΑΟΒ}$: $\widehat{\mathrm{ΑΒΟ}}=180^\circ-20^\circ-90^\circ=70^\circ$.\par
\apantisi $40^\circ,140^\circ,40^\circ,140^\circ$ και $\widehat{\mathrm{ΑΒΟ}}=70^\circ$.\par
\elegxos Η $\widehat{\mathrm{ΑΒΟ}}$ είναι η μισή της $\hat{\mathrm{Β}}$: $140^\circ:2=70^\circ$.

\paradeigma{Ρόμβος με διαγωνίους}
\tixero Ρόμβος με πλευρά $6{,}5$ cm και διαγωνίους $12$ cm και $5$ cm.\par
\tithelo Η περίμετρος του ρόμβου και η περίμετρος του τριγώνου $\mathrm{ΑΟΒ}$.\par
\possindeo Η περίμετρος του ρόμβου είναι $4\cdot6{,}5=26$ cm. Οι διαγώνιοι διχοτομούνται: $\mathrm{ΟΑ}=12:2=6$ cm και $\mathrm{ΟΒ}=5:2=2{,}5$ cm. Το τρίγωνο $\mathrm{ΑΟΒ}$ έχει πλευρές $6$, $2{,}5$ και $6{,}5$.\par
\apantisi $26$ cm και $6+2{,}5+6{,}5=15$ cm.\par
\elegxos Η περίμετρος του ρόμβου είναι $4$ φορές την πλευρά.

\paradeigma{Το τετράγωνο με άγνωστο}
\tixero Στο τετράγωνο $\mathrm{ΑΒΓΔ}$ είναι $\mathrm{ΑΒ}=3x+1$ και $\mathrm{ΒΓ}=2x+6$.\par
\tithelo Το $x$, η πλευρά και η περίμετρος.\par
\possindeo Όλες οι πλευρές του τετραγώνου είναι ίσες, άρα $3x+1=2x+6$, δηλαδή $x=5$. Η πλευρά είναι $3\cdot5+1=16$ και η περίμετρος $4\cdot16=64$.\par
\apantisi $x=5$, πλευρά $16$, περίμετρος $64$.\par
\elegxos $2\cdot5+6=16$.

\paradeigma{Τρίγωνα από τις διαγωνίους}
\tixero Ένα ορθογώνιο, ένας ρόμβος και ένα τετράγωνο με τις διαγωνίους τους.\par
\tithelo Σε τι τρίγωνα χωρίζονται;\par
\possindeo Στο ορθογώνιο οι διαγώνιοι είναι ίσες και διχοτομούνται, άρα τα τέσσερα τρίγωνα έχουν δύο ίσες πλευρές: ισοσκελή. Στον ρόμβο οι διαγώνιοι είναι κάθετες: ορθογώνια τρίγωνα. Στο τετράγωνο και τα δύο: ορθογώνια και ισοσκελή.\par
\apantisi Ορθογώνιο: ισοσκελή. Ρόμβος: ορθογώνια. Τετράγωνο: ορθογώνια και ισοσκελή.\par
\elegxos Το τετράγωνο έχει τις ιδιότητες και των δύο, οπότε βρίσκουμε και τα δύο είδη.

\tmima{Ασκήσεις}

\zoni{Το κατάλαβα}{ποιο σχήμα;}

\ex Σωστό ή λάθος; α) Κάθε τετράγωνο είναι ρόμβος. β) Κάθε ρόμβος είναι τετράγωνο. γ) Οι διαγώνιοι του ρόμβου είναι κάθετες. δ) Οι διαγώνιοι του ορθογωνίου είναι κάθετες. ε) Το ορθογώνιο έχει τέσσερις άξονες συμμετρίας. στ) Το τετράγωνο έχει τέσσερις άξονες συμμετρίας.

\ex Για κάθε ιδιότητα πες σε ποια από τα τρία σχήματα (ορθογώνιο, ρόμβος, τετράγωνο) ισχύει: α) όλες οι γωνίες ορθές β) όλες οι πλευρές ίσες γ) οι διαγώνιοι ίσες δ) οι διαγώνιοι κάθετες ε) οι διαγώνιοι διχοτομούν τις γωνίες στ) τέσσερις άξονες συμμετρίας.

\ex Σχεδίασε: α) ένα ορθογώνιο $5$ cm επί $3$ cm, β) έναν ρόμβο με πλευρά $3$ cm και μία γωνία $60^\circ$, γ) ένα τετράγωνο με πλευρά $4$ cm. Φέρε τις διαγωνίους τους και πες σε τι τρίγωνα χωρίζονται.

\ex Ένας ρόμβος έχει $\hat{\mathrm{Α}}=70^\circ$. Βρες τις άλλες γωνίες του.

\ex Βρες: α) την περίμετρο ορθογωνίου με πλευρές $12$ cm και $7$ cm, β) την πλευρά τετραγώνου με περίμετρο $36$ cm, γ) την πλευρά ρόμβου με περίμετρο $30$ cm.

\ex Στο ορθογώνιο $\mathrm{ΑΒΓΔ}$ η διαγώνιος $\mathrm{ΑΓ}$ είναι $14$ cm. Βρες τη $\mathrm{ΒΔ}$ και τα $\mathrm{ΟΑ}$, $\mathrm{ΟΒ}$, $\mathrm{ΟΓ}$, $\mathrm{ΟΔ}$.

\ex Στον ρόμβο $\mathrm{ΑΒΓΔ}$ είναι $\hat{\mathrm{Α}}=50^\circ$. Βρες τις γωνίες $\widehat{\mathrm{ΟΑΒ}}$ και $\widehat{\mathrm{ΟΒΑ}}$.

\zoni{Το χρησιμοποιώ}{αυλές, παράθυρα, χαρταετοί}

\ex \textbf{Η αυλή.} Η αυλή της θείας Μαρίας είναι τετράγωνη με πλευρά $4{,}5$ m. Πόσα μέτρα περιφράκτη χρειάζεται και πόσους άξονες συμμετρίας έχει το σχήμα της;

\ex \textbf{Το παράθυρο.} Ο παππούς Γιάννης ελέγχει αν ένα παραλληλόγραμμο παράθυρο είναι ορθογώνιο μετρώντας τις διαγωνίους του: $130$ cm και $128$ cm. Είναι ορθογώνιο; Γιατί;

\ex \textbf{Ο χαρταετός.} Ο Νίκος φτιάχνει χαρταετό σε σχήμα ρόμβου με πλευρά $40$ cm. Οι ξύλινες ράβδοι πάνω στις διαγωνίους είναι $64$ cm και $48$ cm. Βρες την περίμετρο του χαρταετού και την περίμετρο ενός από τα τέσσερα τρίγωνα που σχηματίζονται.

\ex \textbf{Τα πλακάκια.} Η Φανή τοποθετεί $5$ τετράγωνα πλακάκια πλευράς $25$ cm σε μία σειρά. Τι σχήμα σχηματίζουν και ποια είναι η περίμετρός του;

\ex \textbf{Ο ισχυρισμός του Στέλιου.} Ο Στέλιος λέει: «Οι διαγώνιοι του ρόμβου είναι κάθετες, άρα και οι διαγώνιοι του ορθογωνίου». Έχει δίκιο; Φέρε ένα παράδειγμα με ορθογώνιο $5$ cm επί $3$ cm.

\zoni{Το σκέφτομαι}{γιατί και πώς}

\ex Εξήγησε με αναδίπλωση γιατί οι διαγώνιοι του ρόμβου τέμνονται κάθετα.

\ex Στο τετράγωνο $\mathrm{ΑΒΓΔ}$ εξήγησε γιατί η διαγώνιος $\mathrm{ΒΔ}$ σχηματίζει γωνία $45^\circ$ με την πλευρά $\mathrm{ΑΒ}$.

\ex Στο ορθογώνιο $\mathrm{ΑΒΓΔ}$ με κέντρο $\mathrm{Ο}$ είναι $\widehat{\mathrm{ΑΒΔ}}=30^\circ$. Βρες τις $\widehat{\mathrm{ΔΒΓ}}$, $\widehat{\mathrm{ΑΟΒ}}$, $\widehat{\mathrm{ΒΟΓ}}$ και πες τι είδους τρίγωνο είναι το $\mathrm{ΟΒΓ}$.

\ex Ένας ρόμβος έχει περίμετρο $52$ cm και μία πλευρά του είναι $3x-2$ cm. Βρες το $x$.

\zoni{Για το σπίτι}{ρίξε μια ματιά}

\ex Κόψε από χαρτί ένα τετράγωνο και ένα ορθογώνιο. Δίπλωσέ τα για να βρεις τους άξονες συμμετρίας τους. Πόσους βρήκες σε κάθε σχήμα;

\ex Βρες στο σπίτι ένα αντικείμενο σχήματος ορθογωνίου, ένα σχήματος τετραγώνου και ένα σχήματος ρόμβου (π.χ. στο σκάκι, σε μοτίβο, σε μπλούζα). Μέτρησε τις διαγωνίους τους και πες ποιες είναι ίσες.

\paragrafos{Κύκλος και στοιχεία του}{Β.3.12, «Κύκλος και στοιχεία του»}

\begin{anagke}
Ο μπαμπάς της Ρήνης τοποθετεί έναν δρομολογητή Wi-Fi στο κέντρο της αυλής του σχολείου, και το σήμα του φτάνει ακριβώς $15$ μέτρα προς κάθε κατεύθυνση. Ο Ηλίας ρωτά: «Πού μπορούμε να καθίσουμε και να έχουμε σήμα;». Ο Νίκος τρέχει $15$ μέτρα από τον δρομολογητή και σταματά. Ο Στέλιος πηγαίνει $20$ μέτρα μακριά, και η Φανή μένει $10$ μέτρα μακριά. «Όλοι όσοι είναι στα $15$ μέτρα σχηματίζουν μια γραμμή», λέει η Αγγελική και με ένα σχοινί την ζωγραφίζει στο χώμα: στρογγυλή. Η Ζωή ρωτά πώς λέγονται τα διάφορα κομμάτια του σχήματος και πόσο «μεγάλο» είναι ένα τμήμα του.

Σε αυτή την παράγραφο μαθαίνουμε τον κύκλο και τα στοιχεία του, μετράμε τόξα με μοίρες και φτιάχνουμε σχήματα με τον διαβήτη.
\end{anagke}

\begin{lexi}{κύκλος, ακτίνα, χορδή, τόξο}{τα στοιχεία του κύκλου}
Ο \lexiword{κύκλος} είναι αυτό που «κυλάει». Η \lexiword{ακτίνα} είναι η «αχτίδα», μια γραμμή που ξεκινά από το κέντρο όπως οι αχτίδες του ήλιου. Το \lexiword{τόξο} είναι το όπλο που τεντώνεται, και η \lexiword{χορδή} είναι το νήμα του τόξου, που ενώνει τα δύο άκρα του. Η \lexiword{διάμετρος} είναι αυτό που μετρά «δια μέσου» του κέντρου.
\end{lexi}

\begin{metafrasi}
Όλα τα σημεία σε ίδια απόσταση από το κέντρο. & κύκλος $(\mathrm{Ο},\rho)$\\
Η απόσταση από το κέντρο. & ακτίνα $\rho$\\
Τμήμα που ενώνει δύο σημεία του κύκλου. & χορδή\\
Χορδή που περνά από το κέντρο. & διάμετρος $\delta=2\rho$\\
Μέρος του κύκλου ανάμεσα σε δύο σημεία. & τόξο $\overset{\frown}{\mathrm{ΑΒ}}$\\
\end{metafrasi}

\tmima{Η θεωρία}

\begin{orismos}{κύκλος}
\textbf{Κύκλος} λέγεται το σύνολο όλων των σημείων του επιπέδου που απέχουν την ίδια απόσταση από ένα σταθερό σημείο $\mathrm{Ο}$. Το $\mathrm{Ο}$ λέγεται \textbf{κέντρο} και η κοινή απόσταση \textbf{ακτίνα} $\rho$. Ο κύκλος με κέντρο $\mathrm{Ο}$ και ακτίνα $\rho$ συμβολίζεται $(\mathrm{Ο},\rho)$ και σχεδιάζεται με τον \textbf{διαβήτη}.
\end{orismos}

Δύο κύκλοι είναι \textbf{ίσοι} όταν έχουν ίσες ακτίνες. Δύο κύκλοι με το ίδιο κέντρο λέγονται \textbf{ομόκεντροι}.

\begin{center}
\begin{tikzpicture}[scale=0.95]
\draw[thick] (0,0) circle (1.8);
\fill (0,0) circle (1.8pt) node[below left] {\small $\mathrm{Ο}$};
\draw[gblue,thick] (0,0)--(1.27,1.27) node[midway,above left] {\small $\rho$};
\fill (1.27,1.27) circle (1.6pt) node[above right] {$\mathrm{Α}$};
\draw (-1.8,0)--(1.8,0); \fill (-1.8,0) circle (1.6pt) node[left] {$\mathrm{Β}$}; \fill (1.8,0) circle (1.6pt) node[right] {$\mathrm{Γ}$};
\node[above] at (-0.7,0.02) {\small διάμετρος};
\draw[orange!80!black,thick] (-1.5,-0.995)--(0.9,-1.56);
\fill (-1.5,-0.995) circle (1.6pt) node[left] {$\mathrm{Δ}$}; \fill (0.9,-1.56) circle (1.6pt) node[below right] {$\mathrm{Ε}$};
\node[orange!80!black] at (-0.2,-1.0) {\small χορδή};
\draw[red!70!black,very thick] (1.8,0) arc (0:45:1.8);
\node[red!70!black] at (2.35,0.75) {\small τόξο};
\end{tikzpicture}
\end{center}

\begin{itemize}[leftmargin=1.2em,itemsep=0.3ex,topsep=2pt]
\item \textbf{Χορδή} είναι το ευθύγραμμο τμήμα που ενώνει δύο σημεία του κύκλου.
\item \textbf{Διάμετρος} είναι η χορδή που περνά από το κέντρο. Είναι η μεγαλύτερη χορδή, ίση με το διπλάσιο της ακτίνας ($\delta=2\rho$), και το κέντρο είναι το μέσο της.
\item Δύο σημεία $\mathrm{Α}$, $\mathrm{Β}$ του κύκλου τον χωρίζουν σε δύο μέρη, καθένα από τα οποία λέγεται \textbf{τόξο} με άκρα τα $\mathrm{Α}$ και $\mathrm{Β}$, $\overset{\frown}{\mathrm{ΑΒ}}$. Ένα τόξο που είναι το μισό του κύκλου λέγεται \textbf{ημικύκλιο}.
\item Ο \textbf{κυκλικός δίσκος} είναι ο κύκλος μαζί με το εσωτερικό του. Ένα σημείο $\mathrm{Μ}$ του δίσκου έχει $\mathrm{ΟΜ}\le\rho$: πάνω στον κύκλο όταν $\mathrm{ΟΜ}=\rho$, μέσα του όταν $\mathrm{ΟΜ}<\rho$, έξω όταν $\mathrm{ΟΜ}>\rho$.
\end{itemize}

\begin{orismos}{επίκεντρη γωνία και μέτρο τόξου}
Μια γωνία λέγεται \textbf{επίκεντρη} όταν η κορυφή της είναι το κέντρο του κύκλου. Οι πλευρές της τέμνουν τον κύκλο στα $\mathrm{Α}$, $\mathrm{Β}$ και το τόξο $\overset{\frown}{\mathrm{ΑΒ}}$ που περιέχεται στο εσωτερικό της λέγεται \textbf{αντίστοιχο τόξο}. Το \textbf{μέτρο ενός τόξου} (σε μοίρες) είναι ίσο με το μέτρο της αντίστοιχης επίκεντρης γωνίας.
\end{orismos}

\begin{center}
\begin{tikzpicture}[scale=0.95]
\draw[thick] (0,0) circle (1.7);
\fill (0,0) circle (1.8pt) node[below] {\small $\mathrm{Ο}$};
\draw[thick] (0,0)--(1.7,0) (0,0)--(0.85,1.47);
\fill (1.7,0) circle (1.6pt) node[right] {$\mathrm{Α}$}; \fill (0.85,1.47) circle (1.6pt) node[above right] {$\mathrm{Β}$};
\draw[red!70!black,very thick] (1.7,0) arc (0:60:1.7);
\draw[gblue] (0.45,0) arc (0:60:0.45); \node[gblue] at (0.72,0.3) {\small $\omega$};
\end{tikzpicture}
\end{center}

\begin{idiotita}{τόξα και επίκεντρες γωνίες}
\begin{itemize}[leftmargin=1.2em,itemsep=0.3ex,topsep=2pt]
\item Ο κύκλος είναι ένα τόξο $360^\circ$ και το ημικύκλιο ένα τόξο $180^\circ$.
\item Δύο κάθετες διάμετροι χωρίζουν τον κύκλο σε τέσσερα ίσα τόξα των $90^\circ$.
\item Στον ίδιο κύκλο (ή σε ίσους κύκλους) δύο ίσες επίκεντρες γωνίες έχουν ίσα αντίστοιχα τόξα και αντίστροφα: δύο ίσα τόξα έχουν ίσες επίκεντρες γωνίες.
\end{itemize}
\end{idiotita}

Ο διαβήτης είναι και εργαλείο κατασκευής. Αν θέλουμε τρίγωνο με πλευρές $a,b,c$, σχεδιάζουμε ένα τμήμα $\mathrm{ΑΒ}=c$ και δύο κύκλους, $(\mathrm{Α},b)$ και $(\mathrm{Β},a)$. Το σημείο τομής τους είναι η τρίτη κορυφή.

\begin{prosoxi}
\begin{itemize}[leftmargin=1.2em,itemsep=0.3ex,topsep=2pt]
\item Ο κύκλος είναι μόνο η γραμμή· το εσωτερικό μαζί με τη γραμμή είναι ο κυκλικός δίσκος.
\item Κάθε διάμετρος είναι χορδή, αλλά όχι κάθε χορδή είναι διάμετρος. Καμία χορδή δεν είναι μεγαλύτερη από τη διάμετρο.
\item Το \emph{μέτρο} ενός τόξου (σε μοίρες) δεν είναι το \emph{μήκος} του. Δύο τόξα με το ίδιο μέτρο σε κύκλους με διαφορετική ακτίνα έχουν διαφορετικό μήκος.
\item Δεν υπάρχει πάντα τρίγωνο με τρία δοσμένα μήκη: οι κύκλοι μπορεί να μη τέμνονται.
\end{itemize}
\end{prosoxi}

\tmima{Λυμένα παραδείγματα}

\paradeigma{Ακτίνα και διάμετρος}
\tixero Ένας κύκλος έχει ακτίνα $3{,}5$ cm. Ένας άλλος έχει διάμετρο $5$ cm.\par
\tithelo Η διάμετρος του πρώτου και η ακτίνα του δεύτερου.\par
\possindeo $\delta=2\rho=2\cdot3{,}5=7$ cm. $\rho=\delta:2=5:2=2{,}5$ cm.\par
\apantisi $7$ cm και $2{,}5$ cm.\par
\elegxos Η διάμετρος είναι διπλάσια της ακτίνας.

\paradeigma{Μέσα, πάνω ή έξω;}
\tixero Κύκλος $(\mathrm{Ο},4)$ cm και σημεία με $\mathrm{ΟΑ}=3$ cm, $\mathrm{ΟΒ}=4$ cm, $\mathrm{ΟΓ}=5$ cm.\par
\tithelo Η θέση κάθε σημείου.\par
\possindeo Συγκρίνουμε με την ακτίνα $4$. Το $\mathrm{Α}$ έχει $3<4$, το $\mathrm{Β}$ έχει $4=4$, το $\mathrm{Γ}$ έχει $5>4$.\par
\apantisi Το $\mathrm{Α}$ είναι μέσα, το $\mathrm{Β}$ πάνω στον κύκλο, το $\mathrm{Γ}$ έξω.\par
\elegxos Μόνο το $\mathrm{Β}$ ανήκει στον κύκλο· τα $\mathrm{Α}$, $\mathrm{Β}$ ανήκουν στον κυκλικό δίσκο.

\paradeigma{Επίκεντρη γωνία και τόξο}
\tixero Σε κύκλο με κέντρο $\mathrm{Ο}$ η επίκεντρη γωνία $\widehat{\mathrm{ΑΟΒ}}$ είναι $70^\circ$.\par
\tithelo Το μέτρο του τόξου $\overset{\frown}{\mathrm{ΑΒ}}$ και του άλλου τόξου με άκρα $\mathrm{Α}$, $\mathrm{Β}$.\par
\possindeo Το τόξο που αντιστοιχεί στη γωνία έχει το ίδιο μέτρο, $70^\circ$. Το υπόλοιπο τόξο συμπληρώνει τον κύκλο των $360^\circ$.\par
\apantisi $70^\circ$ και $360^\circ-70^\circ=290^\circ$.\par
\elegxos $70^\circ+290^\circ=360^\circ$.

\paradeigma{Κατασκευή τριγώνου}
\tixero Θέλουμε τρίγωνο $\mathrm{ΑΒΓ}$ με $\mathrm{ΑΒ}=6$ cm, $\mathrm{ΑΓ}=5$ cm, $\mathrm{ΒΓ}=4$ cm. Υπάρχει τρίγωνο με πλευρές $2$, $3$ και $6$ cm;\par
\tithelo Η κατασκευή και η απάντηση για το δεύτερο τρίγωνο.\par
\possindeo Σχεδιάζουμε $\mathrm{ΑΒ}=6$ cm. Γράφουμε κύκλο $(\mathrm{Α},5)$ και κύκλο $(\mathrm{Β},4)$. Ένα από τα σημεία τομής τους είναι το $\mathrm{Γ}$. Με πλευρές $2$, $3$, $6$ θα χρειαζόμασταν τους κύκλους $(\mathrm{Α},3)$ και $(\mathrm{Β},2)$ με $\mathrm{ΑΒ}=6$. Οι δύο ακτίνες μαζί φτάνουν μόνο $3+2=5$ cm, λιγότερο από τα $6$ cm, άρα οι κύκλοι δεν τέμνονται.\par
\apantisi Το πρώτο τρίγωνο κατασκευάζεται, το δεύτερο δεν υπάρχει.\par
\elegxos Για να τέμνονται οι δύο κύκλοι πρέπει η τρίτη πλευρά να είναι μικρότερη από το άθροισμα των άλλων δύο.

\begin{center}
\begin{tikzpicture}[scale=0.55]
\draw[gray] (0,0) circle (5); \draw[gray] (6,0) circle (4);
\draw[thick] (0,0)--(6,0)--(3.75,3.31)--cycle;
\fill (0,0) circle (2pt) node[below left] {$\mathrm{Α}$};
\fill (6,0) circle (2pt) node[below right] {$\mathrm{Β}$};
\fill (3.75,3.31) circle (2pt) node[above] {$\mathrm{Γ}$};
\node at (3,-0.5) {\small $6$ cm}; \node at (1.5,2.0) {\small $5$ cm}; \node at (5.5,2.0) {\small $4$ cm};
\end{tikzpicture}
\end{center}

\paradeigma{Κύκλος που περνά από τρία σημεία}
\tixero Τρία σημεία $\mathrm{Α}$, $\mathrm{Β}$, $\mathrm{Γ}$ που δεν είναι στην ίδια ευθεία.\par
\tithelo Ο κύκλος που περνά από τα τρία.\par
\possindeo Το κέντρο $\mathrm{Ο}$ απέχει το ίδιο από το $\mathrm{Α}$ και το $\mathrm{Β}$, άρα βρίσκεται στη μεσοκάθετο του $\mathrm{ΑΒ}$. Ομοίως βρίσκεται στη μεσοκάθετο του $\mathrm{ΑΓ}$. Το κέντρο είναι το σημείο τομής των δύο μεσοκαθέτων. Ακτίνα είναι το $\mathrm{ΟΑ}$.\par
\apantisi Κέντρο η τομή των μεσοκαθέτων, ακτίνα $\mathrm{ΟΑ}$.\par
\elegxos Ο κύκλος περνά και από το $\mathrm{Γ}$, γιατί $\mathrm{ΟΓ}=\mathrm{ΟΑ}$.

\begin{center}
\begin{tikzpicture}[scale=0.8]
\draw[gblue,thick] (2,1) circle (2.236);

\draw[dashed] (2,-0.8)--(2,3.2); \draw[dashed] (-0.4,1.8)--(3.2,0.6);
\draw (0,0)--(4,0)--(1,3)--cycle;
\fill (0,0) circle (1.6pt) node[below left] {$\mathrm{Α}$};
\fill (4,0) circle (1.6pt) node[below right] {$\mathrm{Β}$};
\fill (1,3) circle (1.6pt) node[above] {$\mathrm{Γ}$};
\fill (2,1) circle (1.8pt) node[above right] {$\mathrm{Ο}$};
\end{tikzpicture}
\end{center}

\tmima{Ασκήσεις}

\zoni{Το κατάλαβα}{στοιχεία του κύκλου}

\ex Συμπλήρωσε: α) Ο κύκλος είναι ένα τόξο \ldots\ μοιρών. β) Το ημικύκλιο είναι ένα τόξο \ldots\ μοιρών. γ) Δύο κάθετες διάμετροι χωρίζουν τον κύκλο σε $4$ ίσα τόξα των \ldots\ μοιρών. δ) Αν $\rho=7$ cm, τότε $\delta=\ldots$ cm. ε) Αν $\delta=18$ cm, τότε $\rho=\ldots$ cm.

\ex Σωστό ή λάθος; α) Δύο κύκλοι είναι ίσοι όταν έχουν ίσες ακτίνες. β) Κάθε χορδή είναι διάμετρος. γ) Η διάμετρος είναι η μεγαλύτερη χορδή. δ) Όλα τα σημεία του κυκλικού δίσκου απέχουν από το κέντρο το πολύ όσο η ακτίνα. ε) Το μέτρο ενός τόξου είναι ίσο με το μέτρο της αντίστοιχης επίκεντρης γωνίας. στ) Μια χορδή μπορεί να είναι μεγαλύτερη από τη διάμετρο.

\ex Σχεδίασε ένα σημείο $\mathrm{Ο}$. Με τρία χρώματα βάψε: τα σημεία που απέχουν λιγότερο από $2$ cm από το $\mathrm{Ο}$, αυτά που απέχουν ακριβώς $2$ cm και αυτά που απέχουν περισσότερο από $2$ cm.

\ex Σχεδίασε κύκλο $(\mathrm{Ο},3\text{ cm})$ και πάρε πάνω του τρία σημεία $\mathrm{Α}$, $\mathrm{Β}$, $\mathrm{Γ}$. Σχεδίασε τη διάμετρο $\mathrm{ΑΑ'}$, την επίκεντρη γωνία $\widehat{\mathrm{ΑΟΒ}}$ και τη χορδή $\mathrm{ΒΓ}$.

\ex Δύο ομόκεντροι κύκλοι έχουν ακτίνα $2$ cm και διάμετρο $6$ cm αντίστοιχα. Βρες την ακτίνα του δεύτερου και πόσο απέχουν οι δύο κύκλοι.

\ex Ένας κύκλος έχει ακτίνα $6{,}5$ cm. Πόση είναι η διάμετρός του; Ποια είναι η μεγαλύτερη δυνατή χορδή του; Μπορεί να έχει χορδή $14$ cm;

\ex Σε κύκλο με κέντρο $\mathrm{Ο}$ είναι $\widehat{\mathrm{ΑΟΒ}}=125^\circ$. Βρες το μέτρο των δύο τόξων με άκρα $\mathrm{Α}$ και $\mathrm{Β}$.

\zoni{Το χρησιμοποιώ}{σήμα, ρολόγια, πίτσες}

\ex \textbf{Το Wi-Fi.} Ο δρομολογητής της αυλής έχει εμβέλεια $15$ m. Ο Ηλίας κάθεται $12$ m μακριά, η Ρήνη ακριβώς $15$ m και ο Στέλιος $18$ m. Ποιοι έχουν σήμα; Πόση είναι η μέγιστη απόσταση δύο ανθρώπων που έχουν σήμα όταν κάθονται αντίθετα, εκατέρωθεν του δρομολογητή;

\ex \textbf{Το ρολόι.} Η Αγγελική ζωγραφίζει ένα ρολόι. Πόσες μοίρες έχει η επίκεντρη γωνία ανάμεσα σε δύο διαδοχικούς αριθμούς; Πόσες ανάμεσα στο $12$ και στο $3$; Πόσο είναι το τόξο από το $12$ μέχρι το $5$; Πόσες μοίρες διαγράφει ο λεπτοδείκτης σε $10$ λεπτά;

\ex \textbf{Η πίτσα.} Ο Στέλιος κόβει μια στρογγυλή πίτσα σε $8$ ίσα κομμάτια. Πόση είναι η γωνία της κορυφής κάθε κομματιού στο κέντρο; Πόσο είναι το τόξο $3$ κομματιών;

\ex \textbf{Το παρτέρι.} Το κυκλικό παρτέρι του παππού Γιάννη έχει ακτίνα $2{,}4$ m. Πόσο μήκος έχει η μεγαλύτερη ευθεία διαδρομή μέσα στο παρτέρι; Μπορεί να τραβήξει σχοινί $3$ m ανάμεσα σε δύο σημεία του κύκλου; Και $5$ m;

\ex \textbf{Ο στόχος.} Ο Βασίλης παίζει βελάκια σε στόχο με ομόκεντρους κύκλους ακτίνας $5$, $10$ και $15$ cm. Το βελάκι του καρφώνεται $12$ cm από το κέντρο. Ανάμεσα σε ποιους δύο κύκλους βρίσκεται;

\zoni{Το σκέφτομαι}{κατασκευές}

\ex α) Κατασκεύασε τρίγωνο με πλευρές $5$, $4$ και $3$ cm. β) Υπάρχει τρίγωνο με πλευρές $7$, $2$ και $4$ cm; Δικαιολόγησε.

\ex Κατασκεύασε ισόπλευρο τρίγωνο πλευράς $5$ cm μόνο με χάρακα και διαβήτη. Εξήγησε γιατί οι τρεις πλευρές του είναι ίσες.

\ex Η Ζωή βρήκε ένα σπασμένο κυκλικό πιάτο και θέλει να βρει το κέντρο του. Πώς μπορεί, με τρία σημεία στην άκρη του πιάτου και διαβήτη, να το βρει;

\ex Η Πολυξένη σχεδιάζει κύκλο ακτίνας $3$ cm και, με το διαβήτη ανοιχτό $3$ cm, πηγαίνει από σημείο σε σημείο πάνω στον κύκλο. Μετά από πόσα βήματα επιστρέφει στην αρχή; Εξήγησε με τα τρίγωνα που σχηματίζονται.

\zoni{Για το σπίτι}{ρίξε μια ματιά}

\ex Βρες τρία κυκλικά αντικείμενα στο σπίτι. Μέτρησε τη διάμετρό τους και βρες την ακτίνα τους. Πώς θα βρεις το κέντρο ενός χάρτινου κύκλου διπλώνοντάς τον;

\ex Σχεδίασε κύκλο με διάμετρο $7$ cm, χώρισέ τον σε $6$ ίσα τόξα και φτιάξε ένα διακοσμητικό «λουλούδι».

\paragrafos{Ευθεία και κύκλος}{Β.3.13, «Σχετικές θέσεις ευθείας και κύκλου»}

\begin{anagke}
Στο κέντρο του κήπου η θεία Μαρία έχει στρογγυλό ποτιστικό που ρίχνει νερό σε απόσταση $4$ μέτρων προς όλες τις κατευθύνσεις. Κοντά στο ποτιστικό περνά ένα ευθύγραμμο μονοπάτι και η Ζωή ρωτά: «Θα βραχούμε περπατώντας;». Ο Ηλίας μετρά την απόσταση του μονοπατιού από το ποτιστικό: $2{,}5$ μέτρα. Ο Βασίλης λέει ότι θα βραχεί σε δύο σημεία, ο Νίκος ότι σε ένα και η Φανή ότι θα προσέξει μήπως δεν βραχεί καθόλου. «Εξαρτάται», λέει η κυρία Αλεξάνδρα, «από το πόσο μακριά περνά το μονοπάτι». Η Ρήνη βρίσκει ότι υπάρχει μια ειδική απόσταση όπου το μονοπάτι μόλις και μετά βίας ακουμπά το νερό.

Σε αυτή την παράγραφο βρίσκουμε πόσα κοινά σημεία μπορεί να έχουν μια ευθεία και ένας κύκλος και πώς σχετίζεται αυτό με την απόσταση της ευθείας από το κέντρο.
\end{anagke}

\begin{lexi}{τέμνουσα, εφαπτομένη, εξωτερική}{θέσεις ευθείας και κύκλου}
Η \lexiword{τέμνουσα} «κόβει» τον κύκλο. Η \lexiword{εφαπτομένη} είναι αυτή που «εφάπτεται», δηλαδή απλώς αγγίζει (το ρήμα «άπτομαι» σημαίνει «αγγίζω»). Η \lexiword{εξωτερική} ευθεία μένει «έξω» από τον κύκλο.
\end{lexi}

\begin{metafrasi}
Η ευθεία και ο κύκλος έχουν δύο κοινά σημεία. & τέμνουσα, $\mathrm{ΟΜ}<\rho$\\
Η ευθεία και ο κύκλος έχουν ένα κοινό σημείο. & εφαπτομένη, $\mathrm{ΟΜ}=\rho$\\
Η ευθεία και ο κύκλος δεν έχουν κοινό σημείο. & εξωτερική, $\mathrm{ΟΜ}>\rho$\\
\end{metafrasi}

\tmima{Η θεωρία}

Ένας κύκλος $(\mathrm{Ο},\rho)$ και μια ευθεία $\varepsilon$ μπορεί να έχουν \textbf{δύο}, \textbf{ένα} ή \textbf{κανένα} κοινό σημείο. Ονομάζουμε $\mathrm{ΟΜ}$ την \textbf{απόσταση} του κέντρου $\mathrm{Ο}$ από την ευθεία $\varepsilon$, δηλαδή το κάθετο τμήμα από το $\mathrm{Ο}$ προς την $\varepsilon$.

\begin{center}
\begin{tikzpicture}[scale=0.85]
\draw[thick] (0,0) circle (1.3);
\fill (0,0) circle (1.6pt) node[above left] {\small $\mathrm{Ο}$};
\draw[thick,gblue] (-2,-0.7)--(2,-0.7); \draw[dashed] (0,0)--(0,-0.7);
\fill (0,-0.7) circle (1.4pt) node[below right] {\small $\mathrm{Μ}$};
\node at (0,-2.9) {\small τέμνουσα, $\mathrm{ΟΜ}<\rho$};
\begin{scope}[xshift=4.8cm]
\draw[thick] (0,0) circle (1.3);
\fill (0,0) circle (1.6pt) node[above left] {\small $\mathrm{Ο}$};
\draw[thick,gblue] (-2,-1.3)--(2,-1.3); \draw[dashed] (0,0)--(0,-1.3);
\fill (0,-1.3) circle (1.4pt) node[below right] {\small $\mathrm{Μ}$};
\node at (0,-2.9) {\small εφαπτομένη, $\mathrm{ΟΜ}=\rho$};
\end{scope}
\begin{scope}[xshift=9.6cm]
\draw[thick] (0,0) circle (1.3);
\fill (0,0) circle (1.6pt) node[above left] {\small $\mathrm{Ο}$};
\draw[thick,gblue] (-2,-1.9)--(2,-1.9); \draw[dashed] (0,0)--(0,-1.9);
\fill (0,-1.9) circle (1.4pt) node[below right] {\small $\mathrm{Μ}$};
\node at (0,-2.9) {\small εξωτερική, $\mathrm{ΟΜ}>\rho$};
\end{scope}
\end{tikzpicture}
\end{center}

\begin{kanonas}{σχετικές θέσεις ευθείας και κύκλου}
Αν $d=\mathrm{ΟΜ}$ είναι η απόσταση του κέντρου από την ευθεία $\varepsilon$, τότε:
\begin{itemize}[leftmargin=1.2em,itemsep=0.3ex,topsep=2pt]
\item αν $d<\rho$, η $\varepsilon$ είναι \textbf{τέμνουσα} του κύκλου και έχουν \textbf{δύο} κοινά σημεία,
\item αν $d=\rho$, η $\varepsilon$ είναι \textbf{εφαπτομένη} του κύκλου και έχουν \textbf{ένα} κοινό σημείο, το \textbf{σημείο επαφής},
\item αν $d>\rho$, η $\varepsilon$ είναι \textbf{εξωτερική} και \textbf{δεν έχουν} κοινό σημείο.
\end{itemize}
\end{kanonas}

Όταν η $\varepsilon$ εφάπτεται στο σημείο $\mathrm{Μ}$, το $\mathrm{ΟΜ}$ είναι ακτίνα και η απόστασή της από το $\mathrm{Ο}$, άρα:

\begin{idiotita}{εφαπτομένη}
Η εφαπτομένη ενός κύκλου σε ένα σημείο του είναι \textbf{κάθετη στην ακτίνα} που καταλήγει στο σημείο επαφής. Από κάθε σημείο του κύκλου περνά \textbf{μία μόνο} εφαπτομένη.
\end{idiotita}

Αν από ένα σημείο $\mathrm{Μ}$ έξω από τον κύκλο φέρουμε τις δύο εφαπτόμενες που ακουμπούν στα $\mathrm{Α}$ και $\mathrm{Β}$, τα τμήματα $\mathrm{ΜΑ}$ και $\mathrm{ΜΒ}$ λέγονται \textbf{εφαπτόμενα τμήματα}.

\begin{prosoxi}
\begin{itemize}[leftmargin=1.2em,itemsep=0.3ex,topsep=2pt]
\item Η \emph{απόσταση} από το κέντρο είναι το κάθετο τμήμα $\mathrm{ΟΜ}$, όχι οποιοδήποτε τμήμα προς την ευθεία.
\item Όσο μικρότερη είναι η απόσταση, τόσο περισσότερο «μπαίνει» η ευθεία στον κύκλο: μικρή απόσταση σημαίνει τέμνουσα, όχι εξωτερική.
\item Η ευθεία που περνά από το κέντρο είναι τέμνουσα (απόσταση $0$) και όχι εφαπτομένη.
\item Η εφαπτομένη είναι κάθετη στην \emph{ακτίνα} του σημείου επαφής, όχι σε οποιαδήποτε ακτίνα.
\end{itemize}
\end{prosoxi}

\tmima{Λυμένα παραδείγματα}

\paradeigma{Η θέση της ευθείας}
\tixero Κύκλος $(\mathrm{Ο},5)$ cm και τρεις ευθείες με αποστάσεις από το $\mathrm{Ο}$: $3$ cm, $5$ cm και $7$ cm.\par
\tithelo Η θέση κάθε ευθείας ως προς τον κύκλο.\par
\possindeo Συγκρίνουμε την απόσταση με την ακτίνα $\rho=5$.\par
\apantisi Η ευθεία με $3<5$ είναι τέμνουσα (δύο κοινά σημεία). Η ευθεία με $5=5$ είναι εφαπτομένη (ένα κοινό σημείο). Η ευθεία με $7>5$ είναι εξωτερική (κανένα).\par
\elegxos Όσο η ευθεία απομακρύνεται, τα κοινά σημεία μειώνονται: $2$, $1$, $0$.

\paradeigma{Εφαπτομένη σε σημείο του κύκλου}
\tixero Κύκλος $(\mathrm{Ο},\rho)$ και σημείο $\mathrm{Α}$ πάνω του.\par
\tithelo Σχεδίαση της εφαπτομένης στο $\mathrm{Α}$.\par
\possindeo Φέρνουμε την ακτίνα $\mathrm{ΟΑ}$. Στο $\mathrm{Α}$ σχεδιάζουμε την ευθεία $\varepsilon$ που είναι κάθετη στην $\mathrm{ΟΑ}$ (με γνώμονα).\par
\apantisi Η $\varepsilon$ είναι η εφαπτομένη στο $\mathrm{Α}$.\par
\elegxos Η απόσταση του $\mathrm{Ο}$ από την $\varepsilon$ είναι $\mathrm{ΟΑ}=\rho$, άρα η $\varepsilon$ εφάπτεται.

\paradeigma{Κύκλος που εφάπτεται σε ευθεία}
\tixero Ευθεία $\varepsilon$ και σημείο $\mathrm{Α}$ πάνω της.\par
\tithelo Κύκλος που εφάπτεται στην $\varepsilon$ στο $\mathrm{Α}$.\par
\possindeo Φέρνουμε την κάθετη στην $\varepsilon$ στο $\mathrm{Α}$. Παίρνουμε πάνω της ένα οποιοδήποτε σημείο $\mathrm{Κ}$ και γράφουμε τον κύκλο $(\mathrm{Κ},\mathrm{ΚΑ})$.\par
\apantisi Ο κύκλος $(\mathrm{Κ},\mathrm{ΚΑ})$.\par
\elegxos Η απόσταση του $\mathrm{Κ}$ από την $\varepsilon$ είναι $\mathrm{ΚΑ}$, ίση με την ακτίνα. Υπάρχουν άπειροι τέτοιοι κύκλοι, ανάλογα με τη θέση του $\mathrm{Κ}$.

\paradeigma{Εφαπτόμενες στα άκρα χορδής}
\tixero Κύκλος $(\mathrm{Ο},\rho)$ και χορδή $\mathrm{ΑΒ}$.\par
\tithelo Οι εφαπτόμενες στα $\mathrm{Α}$ και $\mathrm{Β}$.\par
\possindeo Σχεδιάζουμε τις ακτίνες $\mathrm{ΟΑ}$ και $\mathrm{ΟΒ}$. Στο $\mathrm{Α}$ φέρνουμε κάθετη στην $\mathrm{ΟΑ}$ και στο $\mathrm{Β}$ κάθετη στην $\mathrm{ΟΒ}$. Αυτές είναι οι εφαπτόμενες και τέμνονται σε ένα σημείο $\mathrm{Μ}$.\par
\apantisi Οι εφαπτόμενες τέμνονται στο $\mathrm{Μ}$ και τα $\mathrm{ΜΑ}$, $\mathrm{ΜΒ}$ είναι τα εφαπτόμενα τμήματα.\par
\elegxos Το $\mathrm{Μ}$ είναι έξω από τον κύκλο.

\begin{center}
\begin{tikzpicture}[scale=0.9]
\draw[thick] (0,0) circle (1.6);
\fill (0,0) circle (1.6pt) node[above] {\small $\mathrm{Ο}$};
\fill (-1.028,-1.226) circle (1.6pt) node[left] {$\mathrm{Α}$};
\fill (1.028,-1.226) circle (1.6pt) node[right] {$\mathrm{Β}$};
\fill (0,-2.09) circle (1.6pt) node[below] {$\mathrm{Μ}$};
\draw (-1.028,-1.226)--(1.028,-1.226);
\draw[gblue,thick] (-1.75,-0.02)--(0,-2.09)--(1.75,-0.02);
\draw[dashed] (0,0)--(-1.028,-1.226) (0,0)--(1.028,-1.226);
\end{tikzpicture}
\end{center}

\tmima{Ασκήσεις}

\zoni{Το κατάλαβα}{τέμνουσα, εφαπτομένη, εξωτερική}

\ex Για κύκλο με ακτίνα $\rho=6$ cm και ευθείες με απόσταση από το κέντρο $d=4$ cm, $d=6$ cm και $d=9$ cm, πες αν η ευθεία είναι τέμνουσα, εφαπτομένη ή εξωτερική και πόσα κοινά σημεία έχει με τον κύκλο.

\ex Βρες πόσα κοινά σημεία έχουν ο κύκλος και η ευθεία όταν: α) $\rho=8$, $d=3$ β) $\rho=4$, $d=4$ γ) $\rho=3$, $d=6$ δ) $\rho=2{,}5$, $d=2{,}5$ ε) $\rho=10$, $d=12$.

\ex Σωστό ή λάθος; α) Η εφαπτομένη έχει ένα μόνο κοινό σημείο με τον κύκλο. β) Η εφαπτομένη είναι κάθετη στην ακτίνα του σημείου επαφής. γ) Αν μια ευθεία είναι εξωτερική, η απόστασή της από το κέντρο είναι μικρότερη από την ακτίνα. δ) Από ένα σημείο του κύκλου περνά μία μόνο εφαπτομένη. ε) Μια ευθεία μπορεί να έχει τρία κοινά σημεία με έναν κύκλο. στ) Κάθε ευθεία που περνά από το κέντρο είναι τέμνουσα.

\ex Σχεδίασε κύκλο $(\mathrm{Ο},4\text{ cm})$ και τρεις ευθείες που απέχουν από το $\mathrm{Ο}$ $3$ cm, $4$ cm και $5$ cm. Πες τη θέση της καθεμιάς.

\ex Σχεδίασε έναν κύκλο, πάρε ένα σημείο $\mathrm{Α}$ πάνω του και σχεδίασε την εφαπτομένη στο $\mathrm{Α}$.

\ex Κύκλος $(\mathrm{Ο},3\text{ cm})$ και ευθεία $\varepsilon$ με απόσταση $3$ cm από το $\mathrm{Ο}$. Πόσα κοινά σημεία έχουν; Πώς πρέπει να αλλάξει η απόσταση για να έχουν δύο κοινά σημεία;

\zoni{Το χρησιμοποιώ}{κήπος, τροχοί, μπάλες}

\ex \textbf{Το ποτιστικό.} Το ποτιστικό της θείας Μαρίας βρέχει σε ακτίνα $4$ m. Τρία ευθύγραμμα μονοπάτια απέχουν από αυτό $2{,}5$ m, $4$ m και $6$ m. Για κάθε μονοπάτι πες πόσα σημεία του βρέχονται.

\ex \textbf{Ο τροχός.} Ο Ηλίας κινεί το ποδήλατό του σε ευθύγραμμο δρόμο. Ο τροχός έχει ακτίνα $35$ cm. Πόσα σημεία του τροχού ακουμπούν στον δρόμο και πόσο ψηλά από το έδαφος βρίσκεται ο άξονας (το κέντρο) του τροχού;

\ex \textbf{Η μπάλα.} Ο Νίκος ακουμπά μια μπάλα του μπάσκετ ακτίνας $12$ cm στη γωνία δαπέδου και τοίχου. Πόσο απέχει το κέντρο της από τον τοίχο και πόσο από το δάπεδο; Τι ισχύει για τον τοίχο και το δάπεδο ως προς τον κύκλο που σχηματίζει η μπάλα στο πλάι;

\ex \textbf{Η λίμνη.} Στο πάρκο υπάρχει κυκλική λίμνη ακτίνας $50$ m. Ένας δρόμος περνά σε απόσταση $38$ m από το κέντρο της. Είναι εφικτός ο δρόμος; Πόση πρέπει να είναι η απόσταση ώστε ο δρόμος να περνά ακριβώς στην άκρη της λίμνης, χωρίς να τη διασχίζει;

\ex \textbf{Το σήμα.} Ο δρομολογητής Wi-Fi της Ρήνης έχει εμβέλεια $25$ m. Ο δρόμος απέναντι απέχει $30$ m από αυτόν. Φτάνει το σήμα στον δρόμο; Αν ο δρόμος ήταν σε απόσταση $25$ m, σε πόσα σημεία θα έφτανε; Και αν ήταν σε απόσταση $20$ m;

\zoni{Το σκέφτομαι}{κατασκευές και συμπεράσματα}

\ex Σε κύκλο $(\mathrm{Ο},\rho)$ φέρε δύο κάθετες διαμέτρους $\mathrm{ΑΒ}$ και $\mathrm{ΓΔ}$ και σχεδίασε τις εφαπτόμενες στα $\mathrm{Α}$, $\mathrm{Β}$, $\mathrm{Γ}$, $\mathrm{Δ}$. Τι σχήμα σχηματίζουν οι τέσσερις εφαπτόμενες και πόση είναι η πλευρά του;

\ex Σχεδίασε κύκλο $(\mathrm{Ο},4{,}5\text{ cm})$ και μια διάμετρο $\mathrm{ΑΒ}$. Σχεδίασε τις εφαπτόμενες του κύκλου που είναι παράλληλες προς την $\mathrm{ΑΒ}$. Πόσες είναι;

\ex Οι ευθείες $\varepsilon_1$ και $\varepsilon_2$ είναι εφαπτόμενες ενός κύκλου στα σημεία $\mathrm{Α}$ και $\mathrm{Β}$. Περίγραψε πώς μπορούμε να βρούμε το κέντρο του κύκλου.

\ex Πόσες εφαπτόμενες του κύκλου περνούν από ένα σημείο: α) πάνω στον κύκλο; β) έξω από τον κύκλο; γ) μέσα στον κύκλο; Δικαιολόγησε.

\zoni{Για το σπίτι}{ρίξε μια ματιά}

\ex Πάρε ένα κυκλικό καπάκι και έναν χάρακα. Τοποθέτησε τον χάρακα ώστε να εφάπτεται στο καπάκι. Μέτρησε την απόσταση του κέντρου από την άκρη του χάρακα. Τι βρίσκεις;

\ex Σχεδίασε ένα ποδήλατο: δύο ίσους κύκλους που εφάπτονται στην ίδια ευθεία (τον δρόμο). Πού βρίσκονται τα κέντρα τους ως προς τον δρόμο;

\epanalipsi{Β3}

\tmima{Λεξικό του κεφαλαίου}

\renewcommand{\arraystretch}{1.5}
\begin{tabular*}{\linewidth}{@{\extracolsep{\fill}} >{\raggedright\arraybackslash}p{0.28\linewidth} >{\raggedright\arraybackslash}p{0.62\linewidth}}
\toprule
\textbf{Λέξη ή σύμβολο} & \textbf{Με μία πρόταση}\\
\midrule
ισόπλευρο, ισοσκελές, σκαληνό & Τρεις, δύο ή καμία ίση πλευρά.\\
οξυγώνιο, ορθογώνιο, αμβλυγώνιο & Όλες οι γωνίες οξείες, μία ορθή, μία αμβλεία.\\
διάμεσος, ύψος, διχοτόμος & Προς το μέσο της απέναντι πλευράς, κάθετο στην απέναντι πλευρά, χωρίζει τη γωνία στη μέση.\\
$\hat{\mathrm{Α}}+\hat{\mathrm{Β}}+\hat{\Gamma}=180^\circ$ & Το άθροισμα των γωνιών τριγώνου.\\
εξωτερική γωνία & Ίση με το άθροισμα των δύο απέναντι εσωτερικών.\\
παραλληλόγραμμο & Απέναντι πλευρές παράλληλες, ίσες· διαγώνιες διχοτομούνται.\\
ορθογώνιο, ρόμβος, τετράγωνο & Ορθές γωνίες· ίσες πλευρές· και τα δύο μαζί.\\
τραπέζιο, ισοσκελές τραπέζιο & Μόνο δύο παράλληλες πλευρές· και ίσες οι άλλες δύο.\\
\bottomrule
\end{tabular*}

\tmima{Ασκήσεις επανάληψης}

\ex Τρίγωνο έχει πλευρές $9$, $9$ και $12$ cm. Ποιο είναι το είδος του και η περίμετρός του;

\ex Βρες την τρίτη γωνία τριγώνου με γωνίες $58^\circ$ και $64^\circ$. Τι είδους είναι;

\ex Ισοσκελές τρίγωνο έχει γωνία κορυφής $36^\circ$. Βρες τις γωνίες βάσης.

\ex Σε τρίγωνο η εξωτερική γωνία στο $\mathrm{Β}$ είναι $130^\circ$ και η $\hat{\mathrm{Α}}=45^\circ$. Βρες τις $\hat{\mathrm{Β}}$ και $\hat{\Gamma}$.

\ex Οι οξείες ενός ορθογωνίου τριγώνου είναι $x$ και $x+20^\circ$. Βρες τες.

\ex Πόσα ύψη έχει ένα τρίγωνο; Πού βρίσκονται σε αμβλυγώνιο;

\ex Βρες την περίμετρο ρόμβου με πλευρά $7{,}5$ cm και παραλληλογράμμου με πλευρές $9$ και $6$ cm.

\ex Ορθογώνιο $9\times12$ cm. Βρες τη διαγώνιο και την απόσταση του κέντρου συμμετρίας από μια κορυφή.

\ex Ρόμβος με διαγώνιες $14$ και $48$ cm. Βρες την πλευρά και την περίμετρό του.

\ex Παραλληλόγραμμο έχει γωνία $75^\circ$. Βρες τις άλλες.

\ex Ισοσκελές τραπέζιο έχει γωνίες βάσης $58^\circ$. Βρες τις άλλες.

\ex Πόσους άξονες συμμετρίας έχουν: ισόπλευρο τρίγωνο, ορθογώνιο, ρόμβος, τετράγωνο, ισοσκελές τραπέζιο;

\ex Η περίμετρος ισοσκελούς τραπεζίου είναι $46$ cm, οι βάσεις $16$ και $10$ cm. Βρες τις ίσες πλευρές.

\ex Ποιο παραλληλόγραμμο έχει ίσες και κάθετες διαγώνιες;

\ex \textbf{Ερωτήσεις αυτοαξιολόγησης (τρίγωνα).} Σωστό ή λάθος; α) Το ισόπλευρο έχει όλες τις γωνίες $60^\circ$. \quad β) Τρίγωνο έχει δύο ορθές γωνίες. \quad γ) Οι γωνίες βάσης ισοσκελούς είναι ίσες. \quad δ) Η εξωτερική γωνία ισούται με το άθροισμα των δύο απέναντι.

\ex \textbf{Ερωτήσεις αυτοαξιολόγησης (τετράπλευρα).} Σωστό ή λάθος; α) Κάθε ρόμβος είναι παραλληλόγραμμο. \quad β) Οι διαγώνιες ρόμβου είναι ίσες. \quad γ) Οι διαγώνιες ορθογωνίου είναι ίσες. \quad δ) Το τραπέζιο έχει τέσσερις παράλληλες πλευρές.

\tmima{Ο χάρτης της επανάληψης}

Μετά τις λύσεις, δες σε ποιες ασκήσεις δυσκολεύτηκες. Ο πίνακας σου λέει ποια παράγραφο να ξαναδιαβάσεις.

\renewcommand{\arraystretch}{1.5}
\begin{tabular*}{\linewidth}{@{\extracolsep{\fill}} >{\raggedright\arraybackslash}p{0.38\linewidth} >{\raggedright\arraybackslash}p{0.52\linewidth}}
\toprule
\textbf{Αν δυσκολεύτηκες στις} & \textbf{ξαναδιάβασε τις παραγράφους}\\
\midrule
ΕΒ3.1, ΕΒ3.6 & 3.1 (στοιχεία και είδη)\\
ΕΒ3.2, ΕΒ3.3, ΕΒ3.4, ΕΒ3.5 & 3.2 (άθροισμα γωνιών)\\
ΕΒ3.7 & 3.3 (είδη τετραπλεύρων)\\
ΕΒ3.8, ΕΒ3.9, ΕΒ3.10, ΕΒ3.11, ΕΒ3.13 & 3.4 (ιδιότητες)\\
ΕΒ3.12, ΕΒ3.14 & 3.4 και Β2 (συμμετρία)\\
ΕΒ3.15, ΕΒ3.16 & 3.1 έως 3.4\\
\bottomrule
\end{tabular*}



\backmatter
\cleardoublepage
\addcontentsline{toc}{part}{Λύσεις}
\markboth{Λύσεις}{Λύσεις}
{\color{gblue}\bfseries\Large Λύσεις}\par\vspace{2mm}
{\color{ggray}\small Κοίτα τες αφού προσπαθήσεις, όχι πριν.\par}
\lyseismeros{Α}{ΑΡΙΘΜΗΤΙΚΗ – ΑΛΓΕΒΡΑ}
\lyseiskef{1}{Φυσικοί αριθμοί}
\lyseis{Α1.0}{Εισαγωγή: οι φυσικοί αριθμοί και οι πράξεις τους}
\sol{1} Ο $0$ δεν έχει προηγούμενο και έχει επόμενο τον $1$. Ο $99$: $98$ και $100$. Ο $1.000$: $999$ και $1.001$. Ο $25.099$: $25.098$ και $25.100$.
\sol{2} Από αριστερά: $4$ εκατομμύρια, $6$ εκατοντάδες χιλιάδων, $3$ δεκάδες χιλιάδων, $5$ χιλιάδες, $1$ εκατοντάδα, $8$ δεκάδες, $2$ μονάδες. Το $6$ είναι στις εκατοντάδες χιλιάδων: αξία $600.000$. Το $3$ είναι στις δεκάδες χιλιάδων: αξία $30.000$.
\sol{3} α) $<$ β) $>$ (ίδιο πλήθος ψηφίων, τα τρία πρώτα ίσα και στη συνέχεια $7>0$) γ) $>$ (ο πρώτος έχει περισσότερα ψηφία) δ) $=$.
\sol{4} α) αντιμεταθετική της πρόσθεσης β) προσεταιριστική της πρόσθεσης γ) το $1$ είναι ουδέτερο στοιχείο του πολλαπλασιασμού δ) το $0$ είναι απορροφητικό στοιχείο ε) επιμεριστική.
\sol{5} α) $7\cdot(100+2)=700+14=714$. β) $8\cdot(100-1)=800-8=792$. γ) $25\cdot(10+2)=250+50=300$.
\sol{6} Πρώτος τρόπος: $6\cdot 25+4\cdot 25=150+100=250$ κιλά. Δεύτερος τρόπος: $(6+4)\cdot 25=10\cdot 25=250$ κιλά. Είναι η επιμεριστική ιδιότητα.
\sol{7} α) $40-30+8=18$. β) $35\cdot 6+8=210+8=218$. γ) $2\cdot 7+10=14+10=24$. δ) $3+8\cdot 5=3+40=43$.
\sol{8} Ο μικρότερος τετραψήφιος είναι ο $1.000$ και ο μεγαλύτερος τριψήφιος ο $999$. Οι διψήφιοι είναι από $10$ ως $99$, δηλαδή $99-10+1=90$. Οι άρτιοι είναι οι $10,12,\dots,98$, δηλαδή $45$ (οι μισοί).
\sol{9} Έχουμε $8-5=3$, ενώ το $5-8$ δεν υπάρχει στους φυσικούς. Άρα η αντιμεταθετική δεν ισχύει. Επίσης $(8-3)-2=5-2=3$, ενώ $8-(3-2)=8-1=7$: διαφέρουν, άρα δεν ισχύει ούτε η προσεταιριστική.
\sol{10} Μεγαλύτερος: $8.520$. Μικρότερος: $2.058$ (το $0$ δεν μπορεί να μπει πρώτο). Διαφορά: $8.520-2.058=6.462$.

\lyseis{Α1.1}{Δυνάμεις φυσικών αριθμών}
\sol{1} α) $5^4$ β) $7^3$ γ) $10^6$ δ) $x^4$ ε) $\alpha^2\cdot\beta^3$.
\sol{2} α) $32$ β) $81$ γ) $216$ δ) $100.000$ ε) $1$ στ) $0$ ζ) $225$.
\sol{3} Στο $6^7$ βάση $6$, εκθέτης $7$. Στο $15^2$ βάση $15$, εκθέτης $2$ (τετράγωνο του $15$). Στο $1^4$ βάση $1$, εκθέτης $4$. Στο $10^9$ βάση $10$, εκθέτης $9$. Κύβος δεν υπάρχει· «τετράγωνο» είναι μόνο η $15^2$.
\sol{4} $121,\ 144,\ 169,\ 196,\ 225$.
\sol{5} α) Λάθος ($16$). β) Λάθος ($8\neq 9$). γ) Σωστό. δ) Λάθος ($1^{20}=1$). ε) Σωστό. στ) Λάθος ($25\neq 10$).
\sol{6} α) $8<9$, άρα $<$. β) $64<81$, άρα $<$. γ) $100<1.024$, άρα $<$. δ) $1=1$, άρα $=$.
\sol{7} α) $20-12=8$. β) $17\cdot 4=68$. γ) $6\cdot 2=12$. δ) $36:12=3$. ε) $4+9=13$. στ) $7^2=49$.
\sol{8} $10\cdot 10\cdot 10\cdot 10=10^4=10.000$ κωδικοί. Με $6$ ψηφία: $10^6=1.000.000$ κωδικοί.
\sol{9} $3\cdot 3\cdot 3\cdot 3=3^4=81$ φυτά.
\sol{10} $12\cdot 12=12^2=144$ πλακάκια.
\sol{11} $4^3=4\cdot 4\cdot 4=64$ κυβικά δεκατόμετρα, δηλαδή $64$ λίτρα.
\sol{12} Τα ρέστα είναι $20-(4\cdot 2+3\cdot 3)=20-(8+9)=20-17=3$ €.
\sol{13} Το $2^5$ είναι $2\cdot 2\cdot 2\cdot 2\cdot 2=32$ και όχι $2\cdot 5=10$. Το $4^3=4\cdot 4\cdot 4=64$ και όχι $4\cdot 3=12$. Ο Νίκος πολλαπλασίασε βάση και εκθέτη, ενώ ο εκθέτης δείχνει πόσες φορές γράφεται η βάση.
\sol{14} α) $(2+3)\cdot 4=20$. β) $(12-4):2=4$. γ) $(2\cdot 3)^2=36$.
\sol{15} $10^{12}=1.000.000.000.000$ (το $1$ και $12$ μηδενικά). $2\cdot 10^5=200.000$ και $35\cdot 10^3=35.000$.
\sol{16} Τα τελευταία ψηφία των $2^1,\dots,2^8$ είναι $2,4,8,6,2,4,8,6$: επαναλαμβάνονται ανά τέσσερα. Ο εκθέτης $20$ είναι πολλαπλάσιο του $4$, όπως οι εκθέτες $4$ και $8$ που δίνουν $6$, άρα το $2^{20}$ τελειώνει σε $6$.
\sol{17} $5.082=5\cdot 10^3+0\cdot 10^2+8\cdot 10+2$ και $40.730=4\cdot 10^4+0\cdot 10^3+7\cdot 10^2+3\cdot 10+0$. Με ψηφία: $6.407$ και $200.080$.
\sol{18} $\alpha=8-(9-7)=8-2=6$. $\beta=4\cdot(9-8)+(4-4)^3=4\cdot 1+0=4$. $\mathrm{Α}=(6+4)^2-2\cdot 6\cdot 4=100-48=52$.

\lyseis{1.2}{Πρόσθεση, αφαίρεση, πολλαπλασιασμός}
\sol{1} Στο $24+15=39$: προσθετέοι και άθροισμα. Στο $50-18=32$: μειωτέος, αφαιρετέος και διαφορά. Στο $6\cdot 9=54$: παράγοντες και γινόμενο.
\sol{2} α) αντιμεταθετική της πρόσθεσης \quad β) προσεταιριστική της πρόσθεσης \quad γ) επιμεριστική του πολλαπλασιασμού ως προς την πρόσθεση \quad δ) το 1 δεν αλλάζει το γινόμενο \quad ε) ο πολλαπλασιασμός με το 0 δίνει 0 \quad στ) το 0 δεν αλλάζει το άθροισμα \quad ζ) αντιμεταθετική του πολλαπλασιασμού \quad η) επιμεριστική ως προς την αφαίρεση.
\sol{3} α) 630 \quad β) 20.500 \quad γ) 38.000 \quad δ) 700.000.
\sol{4} α) $71-26=45$ \quad β) $100-38=62$ \quad γ) $45+19=64$ \quad δ) $90-57=33$. Έλεγχος, π.χ. στο δ: $90-33=57$.
\sol{5} α) $2\cdot 50\cdot 37=100\cdot 37=3.700$ \quad β) $25\cdot 4\cdot 13=100\cdot 13=1.300$ \quad γ) $(16+84)+(39+61)=100+100=200$ \quad δ) $5\cdot 20\cdot 49=100\cdot 49=4.900$.
\sol{6} α) $56\cdot(8+2)=56\cdot 10=560$ \quad β) $(37+63)\cdot 64=100\cdot 64=6.400$ \quad γ) $91\cdot(15-5)=91\cdot 10=910$ \quad δ) $47\cdot(100-1)=4.700-47=4.653$ \quad ε) $24\cdot(100+1)=2.400+24=2.424$.
\sol{7} α) Λάθος: το $5-12$ δεν υπάρχει στους φυσικούς. β) Λάθος: $(30-10)-5=15$, ενώ $30-(10-5)=25$. γ) Λάθος: $4\cdot(6+1)=28$, ενώ $4\cdot 6+1=25$. δ) Λάθος: $0\cdot 9=0$. ε) Σωστό. στ) Σωστό.
\sol{8} α) 300 \quad β) $(117+83)+46=246$ \quad γ) 528 \quad δ) 3.486 (έλεγχος: $2.518+3.486=6.004$) \quad ε) $90-25=65$.
\sol{9} α) $(8+5)+2=15$ και $8+(5+2)=15$. β) $3\cdot(4\cdot 5)=60$ και $(3\cdot 4)\cdot 5=60$. Είναι η προσεταιριστική ιδιότητα, της πρόσθεσης στο α) και του πολλαπλασιασμού στο β).
\sol{10} $98\cdot 5=(100-2)\cdot 5=500-10=490$ τοστ.
\sol{11} Αφετηρία ως γέφυρα: $47$ χλμ. Γέφυρα ως βενζινάδικο: $112-47=65$ χλμ. Βενζινάδικο ως χωριό: $175-112=63$ χλμ. Χωριό ως θάλασσα: $230-175=55$ χλμ. Βενζινάδικο ως θάλασσα: $230-112=118$ χλμ. Έλεγχος: $63+55=118$ και $47+65+63+55=230$.
\sol{12} Μειωτέος είναι τα αυτοκόλλητα που είχε: $Μ=17+38=55$. Έλεγχος: $55-17=38$.
\sol{13} Στην αφαίρεση $82-47$ το 2 στις μονάδες είναι μικρότερο από το 7. Δεν επιτρέπεται να γυρίσουμε τους αριθμούς ανάποδα: δανειζόμαστε μία δεκάδα και κάνουμε $12-7=5$, και μετά $7-4=3$. Σωστό: 35. Έλεγχος με $Μ=Α+\Delta$: $47+35=82$. Με το 45 του Ηλία βγαίνει $47+45=92$, που δεν είναι 82.
\sol{14} $48\cdot 25=12\cdot 4\cdot 25=12\cdot 100=1.200$ πλακάκια.
\sol{15} $13\cdot 8=(10+3)\cdot 8=10\cdot 8+3\cdot 8=80+24=104$ τ.μ. Ένα ορθογώνιο $10\times 8$ και ένα $3\times 8$, κολλημένα.
\sol{16} $3\cdot 99=3\cdot(100-1)=300-3=297$ €.
\sol{17} Τα ζευγάρια είναι 30 και το καθένα κάνει 61 ($1+60$, $2+59$, \dots, $30+31$), άρα $30\cdot 61=1.830$ κούτσουρα. Για το $1+\dots+8$: 4 ζευγάρια και καθένα κάνει 9, άρα $4\cdot 9=36$. Έλεγχος με πρόσθεση: $1+2+3+4+5+6+7+8=36$.
\sol{18} Οι αριθμοί είναι 13 και 7: $13+7=20$ και $13-7=6$.
\sol{19} Πέντε κουτιά με 0 κουλουράκια έχουν συνολικά $0+0+0+0+0=0$. Γενικά $5\cdot 0=0+0+0+0+0=0$.
\sol{20} Το ορθογώνιο με 3 σειρές των 5 έχει 15 τετράγωνα. Αν το γυρίσεις στο πλάι, βλέπεις 5 σειρές των 3, και τα τετράγωνα είναι πάλι 15. Η εικόνα είναι διαφορετική (3 ομάδες των 5 ή 5 ομάδες των 3), το αποτέλεσμα ίδιο.
\sol{21} $(20-8)-4=12-4=8$ και $20-(8-4)=20-4=16$. Οι τιμές διαφέρουν, άρα η αφαίρεση δεν είναι προσεταιριστική: έχει σημασία ποιες δύο αφαιρέσεις γίνονται πρώτες.
\sol{22} $3\cdot 3=9$, $9\cdot 3=27$, $27\cdot 3=81$, $81\cdot 3=243$, $243\cdot 3=729$. Έχει 3 ψηφία. Πιο σύντομα: $3^6$, «το τρία στην έκτη».

\lyseis{Α1.3}{Πρώτοι και σύνθετοι αριθμοί}
\sol{1} α) Σωστό ($24=8\cdot 3$). β) Λάθος (ο $12$ είναι μεγαλύτερος από τον $6$). γ) Σωστό. δ) Λάθος (ο $2$ είναι άρτιος και πρώτος). ε) Σωστό ($121=11\cdot 11$). στ) Λάθος (το $1$ δεν είναι ούτε πρώτος ούτε σύνθετος).
\sol{2} $36$: $1,2,3,4,6,9,12,18,36$. $50$: $1,2,5,10,25,50$ (πρώτοι $2,5$). $83$: μόνο $1$ και $83$, άρα πρώτος.
\sol{3} $2,3,5,7,11,13,17,19,23,29$.
\sol{4} α) πρώτος. β) σύνθετος ($3\cdot 13$). γ) σύνθετος ($3\cdot 17$). δ) σύνθετος ($3\cdot 19$). ε) πρώτος. στ) σύνθετος ($7\cdot 13$).
\sol{5} Ο $120$: $2,3,4,5,10$ (όχι $9$, άθροισμα $3$). Ο $345$: $3,5$ (άθροισμα $12$). Ο $1.008$: $2,3,4,9$ (άθροισμα $9$, $08=8$ διαιρείται με το $4$).
\sol{6} α) $640$, το ψηφίο είναι $0$. β) Το άθροισμα είναι $5+3+8+x=16+x$. Διαιρείται με το $3$ για $x=2,5,8$. Με το $9$ διαιρείται για $x=2$ ($18$). Άρα $x=5$ ή $x=8$. γ) Άθροισμα $1+0+3+4+x=8+x$, άρτιος: $x=4$ ($12$ διαιρείται με το $3$). Άρα $10.344$. δ) Για το $10$ πρέπει $x=0$: $9.700$, και $00$ άρα διαιρείται και με το $4$.
\sol{7} α) Ναι. β) Όχι ($4$ δεν είναι πρώτος). γ) Ναι. δ) Όχι ($9$ δεν είναι πρώτος). ε) Ναι.
\sol{8} $420=2^2\cdot 3\cdot 5\cdot 7$. $792=2^3\cdot 3^2\cdot 11$. $1.350=2\cdot 3^3\cdot 5^2$.
\sol{9} Οι τρόποι αντιστοιχούν στους διαιρέτες: ο $24$ έχει $8$ διαιρέτες ($1,2,3,4,6,8,12,24$), άρα $8$ τρόποι. Ο $13$ είναι πρώτος και έχει $2$ διαιρέτες: μία σειρά των $13$ ή $13$ σειρές του ενός. Δεν υπάρχει άλλη ορθογώνια διάταξη.
\sol{10} Ο $51$ έχει άθροισμα ψηφίων $5+1=6$, που διαιρείται με το $3$, άρα ο $51$ διαιρείται με το $3$: $51=3\cdot 17$. Άρα είναι σύνθετος.
\sol{11} Ο αριθμός είναι $2\cdot2\cdot 3\cdot 5\cdot 7$. Το $6=2\cdot 3$ και το $35=5\cdot 7$ είναι γινόμενα παραγόντων του, άρα τον διαιρούν. Ομοίως $60=2^2\cdot 3\cdot 5$ και $21=3\cdot 7$. Με το $9=3^2$ δεν διαιρείται, γιατί το $3$ εμφανίζεται μόνο μία φορά.
\sol{12} $41,43,47,53,59$.
\sol{13} $(3,5),(5,7),(11,13),(17,19),(29,31),(41,43)$.
\sol{14} α) $8\cdot 3\cdot 7=168$. β) Ο $2\cdot 3\cdot 5\cdot 7=210$.
\sol{15} Π.χ. $33=3\cdot 11$, $63=7\cdot 9$, $93=3\cdot 31$ (και $143=11\cdot 13$).
\sol{16} Αν ένας σύνθετος αριθμός μέχρι το $100$ είχε όλους τους πρώτους παράγοντές του μεγαλύτερους ή ίσους του $11$, θα ήταν τουλάχιστον $11\cdot 11=121>100$. Άρα έχει πρώτο παράγοντα μικρότερο του $11$, δηλαδή έναν από τους $2,3,5,7$. Έτσι τα πολλαπλάσια των $2,3,5,7$ αρκούν.
\sol{17} $6.030$: με $2,3,5,9,10$ (άθροισμα $9$, όχι με το $4$ γιατί $30:4$ δεν είναι τέλεια). $2.800$: με $2,4,5,10$ (άθροισμα $10$, όχι με $3$ ή $9$). $4.617$: με $3,9$ (άθροισμα $18$), όχι με $2,4,5,10$.
\sol{18} $325=5\cdot 65=5\cdot 5\cdot 13=5^2\cdot 13$. $338=2\cdot 169=2\cdot 13^2$.

\lyseis{1.4}{Ευκλείδεια διαίρεση, διαιρετότητα}
\sol{1} α) $1.752:24=73$ \quad β) $2.405:37=65$ \quad γ) $3.136:56=56$ \quad δ) $9.100:25=364$ \quad στ) $7.007:77=91$: όλες τέλειες, με υπόλοιπο 0 (δοκιμή, π.χ. $77\cdot 91=7.007$). ε) $8.514:42$ δίνει πηλίκο 202 και υπόλοιπο 30, γιατί $42\cdot 202=8.484$ και $8.514-8.484=30$. Δοκιμή: $42\cdot 202+30=8.514$.
\sol{2} α) $58=9\cdot 6+4$ \quad β) $200=7\cdot 28+4$ \quad γ) $365=30\cdot 12+5$ \quad δ) $1.000=64\cdot 15+40$ \quad ε) $47=50\cdot 0+47$ (το 50 δεν χωράει στο 47, πηλίκο 0) \quad στ) $0=12\cdot 0+0$.
\sol{3} α) Ναι, με διαιρέτη το 14 (πηλίκο 6): $\upsilon=6<14$. Όχι με διαιρέτη το 6, γιατί το υπόλοιπο 6 δεν είναι μικρότερο από το 6. β) Ναι, με διαιρέτη το 17 (πηλίκο 14) και με διαιρέτη το 14 (πηλίκο 17), γιατί $12<14$. γ) Όχι: το υπόλοιπο 25 δεν είναι μικρότερο από το 25, και είναι μεγαλύτερο από το 15. δ) Όχι: το υπόλοιπο 17 είναι μεγαλύτερο και από το 10 και από το 6.
\sol{4} Τα υπόλοιπα είναι μικρότερα από το 6: 0, 1, 2, 3, 4 ή 5.
\sol{5} $\Delta=11\cdot 48+7=528+7=535$. Δοκιμή: $535:11$ δίνει πηλίκο 48 και υπόλοιπο 7.
\sol{6} α) $\Delta=15\cdot 9+4=139$ \quad β) $250=12\cdot 20+10$, άρα $\pi=20$, $\upsilon=10$ \quad γ) $83=\delta\cdot 5+3$, άρα $5\delta=80$, $\delta=16$. Έλεγχος: $3<16$.
\sol{7} α) Λάθος: το υπόλοιπο είναι μικρότερο από τον διαιρέτη, άρα το πολύ 8. β) Σωστό. γ) Λάθος: δεν διαιρούμε με το 0. δ) Σωστό. ε) Λάθος: $23:1=23$. στ) Σωστό.
\sol{8} Η ισότητα ισχύει: $8\cdot 10+13=93$. Όμως το υπόλοιπο 13 δεν είναι μικρότερο από τον διαιρέτη 8: ο 8 χωράει άλλη μία φορά στο 13. Σωστό: $93=8\cdot 11+5$.
\sol{9} $214=52\cdot 4+6$: χρειάζονται 5 λεωφορεία. Στο πέμπτο υπάρχουν 6 παιδιά και μένουν $52-6=46$ κενές θέσεις. Με 4 λεωφορεία μένουν στο σχολείο 6 παιδιά.
\sol{10} $150=7\cdot 21+3$: Τρίτη και 3 μέρες, δηλαδή Παρασκευή. Πριν από 8 ημέρες: $8=7+1$, άρα μία μέρα πριν από την Τρίτη, δηλαδή Δευτέρα.
\sol{11} $500=24\cdot 20+20$: γεμίζει 20 αυλάκια και περισσεύουν 20 φυτά. Όχι, δεν φτάνουν για άλλο γεμάτο αυλάκι, γιατί $20<24$.
\sol{12} $87=13\cdot 6+9$: αγοράζει 6 μπάλες (78 €) και του μένουν 9 €.
\sol{13} $120:4=30$, $120:5=24$, $120:8=15$, $120:12=10$: τέλειες διαιρέσεις, γίνονται ίσες ομάδες. Το $120:7$ δίνει πηλίκο 17 και υπόλοιπο 1 ($120=7\cdot 17+1$) και το $120:9$ πηλίκο 13 και υπόλοιπο 3 ($120=9\cdot 13+3$): περισσεύουν 1 και 3 παιδιά αντίστοιχα.
\sol{14} $50=24\cdot 2+2$: μετά από 2 ολόκληρα εικοσιτετράωρα το ρολόι δείχνει πάλι 10:00 και μένουν 2 ώρες. Θα είναι 12:00, δύο μέρες αργότερα.
\sol{15} Οι αριθμοί είναι 32, 37, 42, 47, 52, 57 ($32=5\cdot 6+2$, $37=5\cdot 7+2$, \dots). Η διαφορά ανάμεσα σε δύο διαδοχικούς από αυτούς είναι πάντα 5, δηλαδή ο διαιρέτης.
\sol{16} Για $\nu=10$: $10:3$ έχει υπόλοιπο 1. Για 16: $16=3\cdot 5+1$, υπόλοιπο 1. Για 22: $22=3\cdot 7+1$, υπόλοιπο 1. Πάντα 1. Ο λόγος: $\nu=6\cdot\kappa+4$, το $6\kappa$ διαιρείται με το 3, και το 4 αφήνει υπόλοιπο 1.
\sol{17} Ο διαιρέτης πρέπει να είναι μεγαλύτερος από το υπόλοιπο 5, άρα ο μικρότερος είναι ο 6. Τότε $\Delta=6\cdot 9+5=54+5=59$.
\sol{18} Αν $6:0=x$, τότε $6=0\cdot x$, που είναι αδύνατο, γιατί το $0\cdot x$ είναι πάντα 0. Δεν υπάρχει $x$. Στο $0:0=x$ θα έπρεπε $0=0\cdot x$, που ισχύει για κάθε $x$: δεν υπάρχει ένας μοναδικός αριθμός. Και στις δύο περιπτώσεις η πράξη δεν ορίζεται.
\sol{19} $1.000.000=7\cdot 142.857+1$. Το υπόλοιπο είναι 1, άρα μία μέρα μετά την Τρίτη: Τετάρτη.

\lyseis{1.5}{Χαρακτήρες διαιρετότητας, πρώτοι αριθμοί, ΜΚΔ και ΕΚΠ}
\sol{1} Με το 2: 2.480, 6.002, 1.050. Με το 5: 2.480, 3.115, 1.050. Με το 10: 2.480, 1.050.
\sol{2} Άθροισμα ψηφίων: 252 $\to$ 9 (ναι), 3.614 $\to$ 14 (όχι), 8.001 $\to$ 9 (ναι), 55.551 $\to$ 21 (ναι), 1.370 $\to$ 11 (όχι). Έλεγχος, π.χ. $252:3=84$.
\sol{3} 738 $\to$ 18 (ναι), 5.427 $\to$ 18 (ναι), 92.808 $\to$ 27 (ναι), 6.661 $\to$ 19 (όχι), 14.445 $\to$ 18 (ναι).
\sol{4} 524 (24, ναι), 1.018 (18, όχι), 6.300 (00, ναι), 3.772 (72, ναι), 90 (το 90 δεν διαιρείται με το 4, όχι).
\sol{5} 450 (50, ναι), 2.125 (25, ναι), 3.300 (00, ναι), 1.080 (80, όχι), 7.675 (75, ναι).
\sol{6} \begin{tabular}[t]{@{}lccccccc@{}}\toprule & 2 & 3 & 4 & 5 & 9 & 10 & 25\\\midrule 360 & ναι & ναι & ναι & ναι & ναι & ναι & όχι\\ 1.125 & όχι & ναι & όχι & ναι & ναι & όχι & ναι\\ 5.724 & ναι & ναι & ναι & όχι & ναι & όχι & όχι\\ 40.500 & ναι & ναι & ναι & ναι & ναι & ναι & ναι\\\bottomrule\end{tabular}
\sol{7} Το 78 (άρτιος και $7+8=15$), το 3.402 (άρτιος και $3+4+0+2=9$) και το 9.996 (άρτιος και $9+9+9+6=33$). Το 250 δεν διαιρείται με το 3 ($2+5+0=7$), και το 1.113 είναι περιττός.
\sol{8} α) $6+3+8=17$, χρειάζεται 18, άρα $\square=1$: 6.318. \quad β) $4+2=6$ και πρέπει να φτάσει σε πολλαπλάσιο του 3: $\square=0,\,3,\,6,\,9$ (402, 432, 462, 492). \quad γ) Τα δύο τελευταία ψηφία $\square 6$ πρέπει να διαιρούνται με το 4: $16, 36, 56, 76, 96$, άρα $\square=1,\,3,\,5,\,7,\,9$. \quad δ) Πρέπει να τελειώνει σε 0 (διαιρείται με 2 και 5): 290.
\sol{9} α) Λάθος: το 16 τελειώνει σε 6 και δεν διαιρείται με το 6. β) Σωστό: το 9 διαιρείται με το 3. γ) Σωστό. δ) Λάθος: το 24 διαιρείται και με τα δύο, αλλά όχι με το 48. ε) Σωστό. στ) Σωστό.
\sol{10} Το 4.752 έχει άθροισμα ψηφίων 18, άρα διαιρείται με το 3 και το 9. Τελειώνει σε 2, άρα δεν διαιρείται με το 5. Τα δύο τελευταία ψηφία 52 διαιρούνται με το 4. Είναι άρτιος και διαιρείται με το 3, άρα και με το 6. Σε 3 ομάδες: $4.752:3=1.584$ €. Σε 4: $1.188$ €. Σε 6: $792$ €. Σε 9: $528$ €. Σε 5 όχι.
\sol{11} Το 150 τελειώνει σε 0, άρα διαιρείται με το 2, το 5 και το 10. Το άθροισμα ψηφίων είναι 6, άρα διαιρείται με το 3, και επειδή είναι και άρτιος, με το 6. Τα δύο τελευταία ψηφία 50 δεν διαιρούνται με το 4. Γίνεται με 2, 3, 5, 6 και 10 στοίβες.
\sol{12} Ο αριθμός είναι $54\square 0$, με άθροισμα ψηφίων $9+\square$. Πρέπει να διαιρείται με το 9, άρα $\square=0$ ή $9$. Επειδή το τρίτο ψηφίο είναι μεγαλύτερο από το 4, είναι το 9: ο κωδικός είναι 5.490.
\sol{13} Με το 5: $90:5=18$ θέσεις. Με το 10: 9 θέσεις (10, 20, \dots, 90). Και με το 3 και με το 5: οι 15, 30, 45, 60, 75, 90, δηλαδή 6 θέσεις.
\sol{14} Το 2.024 είναι δίσεκτο ($24:4=6$). Το 2.000 είναι δίσεκτο (διαιρείται με το 400). Το 1.900 δεν είναι (διαιρείται με το 100 αλλά όχι με το 400). Το 2.100 δεν είναι. Το 2.030 δεν είναι (30 δεν διαιρείται με το 4).
\sol{15} Αν ένας αριθμός τελειώνει σε 0, είναι πολλαπλάσιο του 10: $10\cdot\kappa$. Επειδή $10=2\cdot 5$, έχουμε $10\kappa=2\cdot(5\kappa)=5\cdot(2\kappa)$. Άρα διαιρείται και με το 2 και με το 5.
\sol{16} Τα κομμάτια $2\cdot 999$, $7\cdot 99$ και $6\cdot 9$ διαιρούνται με το 9 και με το 3. Αρκεί να ελέγξουμε το $2+7+6+3=18$. Το 18 διαιρείται και με το 3 και με το 9, άρα ο 2.763 διαιρείται και με τα δύο ($2.763=9\cdot 307$).
\sol{17} Πρέπει να διαιρείται με το 4, το 5 και το 6. Τα πολλαπλάσια του $\EKP(4,\,5,\,6)=60$ είναι 60, 120, \dots, και το πρώτο τετραψήφιο είναι $60\cdot 17=1.020$.
\sol{18} Το μεσαίο ψηφίο είναι το άθροισμα των δύο άκρων: $1+1=2$, $1+2=3$, $1+3=4$, $1+4=5$. Το 165 ταιριάζει ($1+5=6$) και $165=11\cdot 15$. Το 175 δεν ταιριάζει ($1+5=6\neq 7$) και $175:11$ δίνει υπόλοιπο 10.
\sol{19} Για το 4 δεν αρκεί το τελευταίο ψηφίο: κοιτάμε τα δύο τελευταία ψηφία μαζί. Το 34 δεν διαιρείται με το 4 ($34=4\cdot 8+2$), άρα ο 5.034 δεν διαιρείται με το 4.
\sol{20} 16: 1, 2, 4, 8, 16. \quad 20: 1, 2, 4, 5, 10, 20. \quad 28: 1, 2, 4, 7, 14, 28. \quad 42: 1, 2, 3, 6, 7, 14, 21, 42. \quad 49: 1, 7, 49.
\sol{21} Του 3: 0, 3, 6, 9, 12, 15. Του 8: 0, 8, 16, 24, 32, 40. Του 15: 0, 15, 30, 45, 60, 75.
\sol{22} Πρώτοι: 3, 23, 31, 43, 59, 83. Το 1 δεν είναι ούτε πρώτος ούτε σύνθετος. Σύνθετοι: 15 (3), 27 (3), 39 (3), 49 (7), 51 (3), 63 (3), 77 (7).
\sol{23} 83, 89, 97, 101, 103, 107, 109.
\sol{24} α) $36=2^2\cdot 3^2$ \quad β) $90=2\cdot 3^2\cdot 5$ \quad γ) $126=2\cdot 3^2\cdot 7$ \quad δ) $200=2^3\cdot 5^2$ \quad ε) $360=2^3\cdot 3^2\cdot 5$ \quad στ) $462=2\cdot 3\cdot 7\cdot 11$.
\sol{25} α) $16\cdot 3=48$ \quad β) $2\cdot 25\cdot 7=350$ \quad γ) $9\cdot 5\cdot 7=315$ \quad δ) $8\cdot 3\cdot 11=264$.
\sol{26} α) Σωστό. β) Λάθος. γ) Λάθος: $77=7\cdot 11$. δ) Λάθος: $3+5=8$. ε) Σωστό: έχει διαιρέτες τουλάχιστον τους δύο παράγοντες. στ) Σωστό: $0=\alpha\cdot 0$. ζ) Λάθος: $57=3\cdot 19$.
\sol{27} Ο 323 δεν διαιρείται με το 2, το 3 ($3+2+3=8$), το 5, το 7 ($7\cdot 46=322$), το 11 ($11\cdot 29=319$) ούτε το 13 ($13\cdot 24=312$). Με το 17: $17\cdot 19=323$. Οι πρώτοι είναι 17 και 19.
\sol{28} Οι διαιρέτες του 54 είναι 1, 2, 3, 6, 9, 18, 27, 54: σακουλάκια με όποιον από αυτούς τους αριθμούς καραμελών. Το 53 είναι πρώτος, άρα έχει μόνο τους διαιρέτες 1 και 53: δύο μόνο τρόποι.
\sol{29} Οι διαιρέτες του 72 είναι 1, 2, 3, 4, 6, 8, 9, 12, 18, 24, 36, 72. Αφαιρώ τον 1 και τον 72: μένουν 10, άρα 10 τρόποι (2, 3, 4, 6, 8, 9, 12, 18, 24 ή 36 παιδιά σε κάθε ομάδα). Το 73 είναι πρώτος, άρα κανένας τρόπος.
\sol{30} Οι πρώτοι ανάμεσα στο 30 και στο 50 είναι οι 31, 37, 41, 43 και 47. Άθροισμα ψηφίων: 4, 10, 5, 7, 11. Μόνο ο 43 έχει άθροισμα 7: ο αριθμός είναι ο 43.
\sol{31} Οι πρώτοι μέχρι το 60 είναι 17: 2, 3, 5, 7, 11, 13, 17, 19, 23, 29, 31, 37, 41, 43, 47, 53, 59. Το 2 σβήνει 29 αριθμούς (4, 6, 8, \dots, 60). Το 3 σβήνει ακόμα 9 αριθμούς: 9, 15, 21, 27, 33, 39, 45, 51, 57 (τα άλλα πολλαπλάσιά του είναι ήδη σβησμένα, γιατί είναι άρτια).
\sol{32} Για το 6: $1+2+3=6$. Για το 28: $1+2+4+7+14=28$. Σε κάθε περίπτωση το άθροισμα των διαιρετών (εκτός από τον ίδιο) ισούται με τον αριθμό.
\sol{33} $(3,\,5)$, $(5,\,7)$, $(11,\,13)$, $(17,\,19)$, $(29,\,31)$, $(41,\,43)$.
\sol{34} Οι αριθμοί είναι $1, 3, 9, 2, 6, 18, 4, 12, 36, 8, 24, 72$, δηλαδή 12 αριθμοί. Είναι όλοι οι διαιρέτες του 72.
\sol{35} Αν το 1 ήταν πρώτος, το 6 θα είχε πολλές αναλύσεις: $2\cdot 3$, $1\cdot 2\cdot 3$, $1\cdot 1\cdot 2\cdot 3$, \dots, όσες θέλεις. Η ανάλυση δεν θα ήταν μοναδική.
\sol{36} $101\cdot 103=10.100+303=10.403$. Αν δοκιμάζαμε τους πρώτους με τη σειρά (2, 3, 5, \dots, 97), θα αποτυγχάναμε 25 φορές, και μόνο ο 101 θα πετύχαινε. Με ακόμα μεγαλύτερους αριθμούς η δουλειά γίνεται τεράστια, και σε αυτό στηρίζεται η ασφάλεια των κωδικών.
\sol{37} α) Διαιρέτες του 14: 1, 2, 7, 14. Του 35: 1, 5, 7, 35. $\MKD=7$. \quad β) Του 24: 1, 2, 3, 4, 6, 8, 12, 24. Του 36: 1, 2, 3, 4, 6, 9, 12, 18, 36. $\MKD=12$. \quad γ) Του 16: 1, 2, 4, 8, 16. Του 45: 1, 3, 5, 9, 15, 45. Κοινός μόνο ο 1: $\MKD=1$. \quad δ) Του 28: 1, 2, 4, 7, 14, 28. Του 42: 1, 2, 3, 6, 7, 14, 21, 42. $\MKD=14$.
\sol{38} α) $\EKP(6,\,8)=24$ \quad β) $\EKP(9,\,12)=36$ \quad γ) $\EKP(10,\,15)=30$ \quad δ) $\EKP(4,\,9)=36$.
\sol{39} α) $30=2\cdot 3\cdot 5$, $45=3^2\cdot 5$: $\MKD=3\cdot 5=15$, $\EKP=2\cdot 3^2\cdot 5=90$. \quad β) $72=2^3\cdot 3^2$, $108=2^2\cdot 3^3$: $\MKD=2^2\cdot 3^2=36$, $\EKP=2^3\cdot 3^3=216$. \quad γ) $84=2^2\cdot 3\cdot 7$, $126=2\cdot 3^2\cdot 7$: $\MKD=2\cdot 3\cdot 7=42$, $\EKP=2^2\cdot 3^2\cdot 7=252$. \quad δ) $150=2\cdot 3\cdot 5^2$, $225=3^2\cdot 5^2$: $\MKD=3\cdot 5^2=75$, $\EKP=2\cdot 3^2\cdot 5^2=450$.
\sol{40} $20=2^2\cdot 5$, $30=2\cdot 3\cdot 5$, $50=2\cdot 5^2$. $\MKD=2\cdot 5=10$ (κοινοί παράγοντες 2 και 5, με μικρότερο εκθέτη 1). $\EKP=2^2\cdot 3\cdot 5^2=300$.
\sol{41} Πρώτοι μεταξύ τους είναι τα (α), (β), (δ) και (στ). Τα (γ) 18 και 27 έχουν κοινό διαιρέτη το 9, και τα (ε) 35 και 49 το 7.
\sol{42} α) $\MKD=6$, $\EKP=48$ \quad β) $13$ και $52$ \quad γ) $20$ και $100$. Όταν ο ένας αριθμός διαιρεί τον άλλον, το ΜΚΔ είναι ο μικρότερος και το ΕΚΠ ο μεγαλύτερος.
\sol{43} α) Λάθος: το ΜΚΔ είναι 4. β) Λάθος: το ΕΚΠ είναι 18. γ) Σωστό. δ) Σωστό. ε) Σωστό. στ) Σωστό.
\sol{44} $10=2\cdot 5$ και $25=5^2$, άρα $\EKP=2\cdot 5^2=50$ δευτερόλεπτα. Στα 3 λεπτά (180 δευτ.) συμπίπτουν στα 50, 100 και 150 δευτερόλεπτα: 3 φορές.
\sol{45} $90=2\cdot 3^2\cdot 5$, $150=2\cdot 3\cdot 5^2$, άρα $\MKD=2\cdot 3\cdot 5=30$. Το πολύ 30 καλάθια, με $90:30=3$ μήλα και $150:30=5$ πορτοκάλια το καθένα.
\sol{46} $240=2^4\cdot 3\cdot 5$, $180=2^2\cdot 3^2\cdot 5$, άρα $\MKD=2^2\cdot 3\cdot 5=60$. Μέγιστη πλευρά 60 εκ. Πλακάκια: $240:60=4$ στο μήκος και $180:60=3$ στο πλάτος, δηλαδή $4\cdot 3=12$ πλακάκια.
\sol{47} $12=2^2\cdot 3$ και $16=2^4$, άρα $\EKP=2^4\cdot 3=48$ λεπτά: ξανά μαζί στις 7:48. Μαζί βρίσκονται στις 7:00, 7:48, 8:36 και 9:24 (το επόμενο θα ήταν 10:12, μετά τις 10:00): 4 φορές.
\sol{48} $\EKP(4,\,6,\,9)$: $4=2^2$, $6=2\cdot 3$, $9=3^2$, άρα $2^2\cdot 3^2=36$. Σε 36 μέρες.
\sol{49} $(10,\,25)$: ΜΚΔ 5, ΕΚΠ 50, $5\cdot 50=250=10\cdot 25$. $(16,\,24)$: 8 και 48, $8\cdot 48=384=16\cdot 24$. $(21,\,35)$: 7 και 105, $7\cdot 105=735=21\cdot 35$. $(18,\,30)$: 6 και 90, $6\cdot 90=540=18\cdot 30$. Το γινόμενο ΜΚΔ·ΕΚΠ είναι πάντα ίσο με το γινόμενο των δύο αριθμών.
\sol{50} $180=84\cdot 2+12$, και $\MKD(180,\,84)=12=\MKD(84,\,12)$. Συνεχίζω: $84=12\cdot 7+0$. Ο τελευταίος διαιρέτης πριν το υπόλοιπο 0, το 12, είναι το ΜΚΔ.
\sol{51} $9=3^2$ και $20=2^2\cdot 5$ δεν έχουν κοινό παράγοντα. Το ΕΚΠ πρέπει να περιέχει όλους τους παράγοντες και των δύο, με τον μεγαλύτερο εκθέτη, και επειδή δεν υπάρχει κανένας κοινός, τους παίρνει όλους: $2^2\cdot 3^2\cdot 5=180=9\cdot 20$.
\sol{52} Το γινόμενο των δύο αριθμών πρέπει να είναι $4\cdot 48=192$. Ένα ζευγάρι είναι το 4 και το 48 (ΜΚΔ 4, ΕΚΠ 48). Δεύτερο ζευγάρι: 12 και 16 (ΜΚΔ 4, ΕΚΠ 48, γινόμενο 192).
\sol{53} Όχι. Το ΜΚΔ είναι μικρότερο ή ίσο από τον μικρότερο αριθμό και το ΕΚΠ μεγαλύτερο ή ίσο από τον μεγαλύτερο. Αν οι αριθμοί είναι διαφορετικοί, ο μικρότερος είναι αυστηρά μικρότερος από τον μεγαλύτερο, άρα $\MKD<\EKP$. Ισότητα έχουμε μόνο όταν οι δύο αριθμοί είναι ίσοι.
\sol{54} $\EKP(5,\,6,\,8)$: $5$, $6=2\cdot 3$, $8=2^3$, άρα $2^3\cdot 3\cdot 5=120$. Σε 120 ημέρες.
\sol{55} Το πλήθος είναι πολλαπλάσιο του 4, του 6 και του 9, άρα του $\EKP(4,\,6,\,9)=36$. Το λιγότερο έχει 36 κουλουράκια. Τα πολλαπλάσια του 36 είναι 36, 72, 108, 144, 180, 216, 252, \dots, άρα ανάμεσα στο 150 και στο 250 είναι 180 ή 216.
\sol{56} $\EKP(6,\,8)=24$ ημέρες. Η Αγγελική πάει $24:6=4$ φορές μέχρι τότε, άρα 3 φορές ανάμεσα. Ο Νίκος $24:8=3$ φορές, άρα 2 φορές ανάμεσα.
\sol{57} Οι κοινοί διαιρέτες δύο αριθμών είναι ακριβώς οι διαιρέτες του ΜΚΔ τους. Ο 18 έχει διαιρέτες και τους 2, 3, 6, 9, άρα και οι δύο αριθμοί διαιρούνται με όλους αυτούς.
\sol{58} Σύνθετος. Το τριπλάσιο $3p$ έχει διαιρέτες τουλάχιστον τους $1$, $3$ και $3p$, και ο 3 είναι διαφορετικός από το 1 και από το $3p$ (αφού $3p\geq 6$). Άρα έχει διαιρέτη εκτός από το 1 και τον εαυτό του.
\sol{59} α) 1, 2, 4, 5, 8, 10, 20, 40 \quad β) 1, 2, 4, 7, 8, 14, 28, 56 \quad γ) 1, 2, 3, 6, 11, 22, 33, 66 \quad δ) 1, 7, 13, 91 \quad ε) 1, 2, 3, 5, 6, 10, 15, 25, 30, 50, 75, 150 \quad στ) 1, 13, 17, 221.
\sol{60} α) $66=2\cdot 3\cdot 11$ \quad β) $275=5^2\cdot 11$ \quad γ) $840=2^3\cdot 3\cdot 5\cdot 7$ \quad δ) $1.001=7\cdot 11\cdot 13$.
\sol{61} Με το 6: $42+18=60=6\cdot 10$ και $42-18=24=6\cdot 4$. Γενικά αν $\delta\mid\beta$ και $\delta\mid\gamma$, τότε $\beta=\delta\kappa$ και $\gamma=\delta\lambda$, άρα $\beta+\gamma=\delta(\kappa+\lambda)$. Για το $5\alpha+15$: ο 5 διαιρεί το $5\alpha$ και το 15, άρα και το άθροισμά τους. Ή $5\alpha+15=5(\alpha+3)$.
\sol{62} $24\cdot 101=2.424$ και $58\cdot 101=5.858$: το διψήφιο γράφεται δύο φορές. Εξήγηση: $24\cdot 101=24\cdot(100+1)=2.400+24=2.424$. Το $24\cdot 100$ βάζει το 24 στις εκατοντάδες και το $+24$ στις μονάδες.

\lyseisE{Επανάληψη κεφαλαίου Α1}
\solr{Α1}{1} α) $81$ β) $1.000.000$ γ) $128$ δ) $1$ ε) $0$.
\solr{Α1}{2} $2\cdot 9+6^2:6-5=18+36:6-5=18+6-5=19$.
\solr{Α1}{3} $12\cdot 22=264$ και $275-264=11$, άρα $275=12\cdot 22+11$ με $11<12$.
\solr{Α1}{4} $\Delta=9\cdot 8+7=72+7=79$.
\solr{Α1}{5} $100=6\cdot 16+4$. Τα $16$ κουτιά γεμίζουν και περισσεύουν $4$ μπαλάκια, άρα χρειάζεται $17$ κουτί. Το τελευταίο έχει $4$ μπαλάκια.
\solr{Α1}{6} Άρτιος: με το $2$. Άθροισμα ψηφίων $3+4+5+6=18$: με το $3$ και το $9$. Το $56$ διαιρείται με το $4$ ($56=4\cdot 14$): με το $4$. Δεν τελειώνει σε $0$ ή $5$. Άρα με $2,3,4,9$.
\solr{Α1}{7} Πρώτοι: $29$, $61$, $97$. Σύνθετοι: $51=3\cdot 17$, $57=3\cdot 19$, $87=3\cdot 29$.
\solr{Α1}{8} $180=2^2\cdot 3^2\cdot 5$.
\solr{Α1}{9} $180=2^2\cdot 3^2\cdot 5$, $150=2\cdot 3\cdot 5^2$. $\MKD=2\cdot 3\cdot 5=30$. $\EKP=2^2\cdot 3^2\cdot 5^2=900$.
\solr{Α1}{10} $\EKP(6,8)=24$ ($6=2\cdot 3$, $8=2^3$, $2^3\cdot 3=24$). Μετά από $24$ λεπτά.
\solr{Α1}{11} $\MKD(42,63)$: $42=2\cdot 3\cdot 7$, $63=3^2\cdot 7$, $\MKD=21$. Φτιάχνει $21$ ανθοδέσμες με $42:21=2$ κόκκινα και $63:21=3$ άσπρα μπαλόνια.
\solr{Α1}{12} $(60,84,132)\to(60,24,12)\to(0,0,12)$. Άρα $\MKD=12$.
\solr{Α1}{13} $32+0+8+4+0+1=45$.
\solr{Α1}{14} $37=32+4+1$, άρα $(100101)_2$. (Με διαιρέσεις: υπόλοιπα $1,0,1,0,0,1$ από πάνω προς τα κάτω.)
\solr{Α1}{15} α) Λάθος ($32$). β) Λάθος. γ) Λάθος (πρέπει να είναι μικρότερο από τον διαιρέτη). δ) Λάθος. ε) Σωστό.
\solr{Α1}{16} α) Σωστό ($2+1$). β) Σωστό ($9=3^2$, $16=2^4$, ΜΚΔ $=1$). γ) Σωστό ($15=3\cdot 5$ είναι παράγοντας του γινομένου). δ) Λάθος (είναι μεγαλύτερο ή ίσο). ε) Λάθος (άθροισμα ψηφίων $10$, δεν διαιρείται με το $3$).


\lyseiskef{2}{Τα κλάσματα}
\lyseis{2.1}{Η έννοια του κλάσματος}
\sol{1} α) $\frac{3}{8}$ \quad β) $\frac{5}{6}$ \quad γ) $\frac{1}{4}$ \quad δ) $\frac{7}{10}$.
\sol{2} Χωρίζεις κάθε ορθογώνιο σε ίσα μέρη, με το ίδιο πλάτος: σε 5 για το α) και χρωματίζεις 2, σε 4 για το β) και χρωματίζεις 3, σε 6 για το γ) και χρωματίζεις 5.
\sol{3} α) $\frac{3}{7}$ \quad β) $\frac{5}{4}$ \quad γ) $\frac{1}{10}$ \quad δ) $\frac{7}{9}$. Διαβάζονται: «τέσσερα ένατα», «ένα δωδέκατο», «έντεκα έκτα», «τρία εκατοστά».
\sol{4} $5:8=\frac{5}{8}$, $3:11=\frac{3}{11}$, $10:3=\frac{10}{3}$, $7:7=\frac{7}{7}$ (που ισούται με 1). Και $\frac{9}{4}=9:4$, $\frac{1}{6}=1:6$, $\frac{20}{5}=20:5$ (που ισούται με 4).
\sol{5} $7=\frac{7}{1}=\frac{28}{4}=\frac{70}{10}$ (γιατί $7\cdot 4=28$ και $7\cdot 10=70$). Και $0=\frac{0}{5}$.
\sol{6} Μικρότερα από το 1: $\frac{3}{4}$, $\frac{5}{12}$, $\frac{1}{100}$. Ίσα με το 1: $\frac{7}{7}$, $\frac{20}{20}$. Μεγαλύτερα από το 1: $\frac{10}{9}$, $\frac{18}{17}$.
\sol{7} α) $\frac{25}{60}$ \quad β) $\frac{250}{1000}$ \quad γ) $\frac{2}{7}$ \quad δ) $\frac{3}{12}$.
\sol{8} α) $20:4\cdot 3=15$ \quad β) $35:5\cdot 2=14$ \quad γ) $48:6\cdot 5=40$ \quad δ) $200:10\cdot 7=140$.
\sol{9} α) $14:2\cdot 3=21$ \quad β) $27:3\cdot 5=45$ \quad γ) $36:4\cdot 9=81$ \quad δ) $56:7\cdot 8=64$.
\sol{10} Το ένα πέμπτο των 20 είναι 4, άρα τα τρία πέμπτα είναι $3\cdot 4=12$ βολές. Έχασε $20-12=8$ βολές, δηλαδή τα $\frac{2}{5}$ ($1-\frac{3}{5}$).
\sol{11} Το ένα ένατο των 36 είναι 4, άρα τα $\frac{5}{9}$ είναι $5\cdot 4=20$ χρωματιστά. Τα υπόλοιπα είναι $36-20=16$, δηλαδή τα $\frac{4}{9}$ των μολυβιών.
\sol{12} Το ένα όγδοο των 2.000 γρ. είναι 250 γρ., άρα τα $\frac{3}{8}$ είναι $3\cdot 250=750$ γραμμάρια.
\sol{13} Ο Βασίλης: $27:9\cdot 4=12$ ψήφους. Η Ζωή: $27:3=9$ ψήφους. Η Μυρτώ: $27-12-9=6$ ψήφους, δηλαδή τα $\frac{6}{27}$, που είναι ίσα με τα $\frac{2}{9}$ ($\frac{9}{9}-\frac{4}{9}-\frac{3}{9}$).
\sol{14} Τα $\frac{5}{8}$ είναι 150 χλμ., άρα το ένα όγδοο είναι $150:5=30$ χλμ. και όλη η διαδρομή $8\cdot 30=240$ χλμ. Μένουν $240-150=90$ χλμ.
\sol{15} Τα $\frac{3}{7}$ είναι 18 λίτρα, άρα το ένα έβδομο είναι $18:3=6$ λίτρα και ολόκληρο το δοχείο χωράει $7\cdot 6=42$ λίτρα.
\sol{16} $\frac{15}{60}$, $\frac{20}{60}$ και $\frac{45}{60}$ της ώρας. Η σχολική μέρα είναι τα $\frac{6}{24}$ του εικοσιτετραώρου.
\sol{17} Αν κόψεις και τις δύο πίτες σε 9 κομμάτια, το $\frac{1}{9}$ κάθε πίτας είναι ένα κομμάτι. Παίρνεις ένα από κάθε πίτα, άρα 2 κομμάτια των $\frac{1}{9}$, δηλαδή $\frac{2}{9}$ μιας πίτας. Με υπολογισμό: $2\cdot\frac{1}{9}=\frac{2}{9}$.
\sol{18} Το συμπέρασμα είναι σωστό: $\frac{3}{5}<\frac{3}{4}$. Αν κόψεις μια πίτα σε 4 κομμάτια, κάθε κομμάτι είναι μεγαλύτερο από ό,τι αν την κόψεις σε 5 (περισσότερα κομμάτια σημαίνει μικρότερα κομμάτια). Τρία μεγάλα κομμάτια ($\frac{3}{4}$) είναι περισσότερο φαγητό από τρία μικρά ($\frac{3}{5}$).
\sol{19} Για να είναι το $\frac{7}{\nu}>1$ πρέπει $\nu<7$. Για να είναι $<2$ πρέπει $7<2\nu$, δηλαδή $\nu\geq 4$. Άρα $\nu=4,\,5,\,6$ και τα κλάσματα είναι τα $\frac{7}{4}$, $\frac{7}{5}$, $\frac{7}{6}$: τρία.
\sol{20} Αν $\frac{5}{0}=5:0=x$, τότε $5=0\cdot x$, που είναι αδύνατο, γιατί το $0\cdot x$ είναι πάντα 0. Δεν υπάρχει αριθμός που να είναι το $5:0$. Γι' αυτό ο παρονομαστής δεν μπορεί να είναι 0.
\sol{21} Χωρίζω ολόκληρη την πίτα σε 6 ίσα κομμάτια. Η μισή πίτα είναι 3 κομμάτια. Τα $\frac{2}{3}$ της μισής είναι 2 από αυτά τα 3 κομμάτια, δηλαδή $\frac{2}{6}$ της ολόκληρης πίτας. (Το $\frac{2}{6}$ είναι ίσο με το $\frac{1}{3}$, όπως θα δούμε στην παράγραφο 2.2.)
\sol{22} $\mathrm{ΓΔ}=12:4\cdot 3=9$ εκ. και $\mathrm{ΕΖ}=12:6\cdot 5=10$ εκ.
\sol{23} Ενδεικτικά: μια σοκολάτα με 12 τετράγωνα, από τα οποία έφαγες 3: $\frac{3}{12}$. Ένα ημερολόγιο: ο Μάρτιος είναι $\frac{1}{12}$ του έτους. Μια αυγοθήκη των 6 με 2 αυγά: $\frac{2}{6}$.

\lyseis{2.2}{Ισοδύναμα κλάσματα}
\sol{1} $\frac{2}{3}=\frac{2\cdot 2}{3\cdot 2}=\frac{4}{6}$, $\frac{2}{3}=\frac{6}{9}$ και $\frac{2}{3}=\frac{20}{30}$.
\sol{2} α) $\frac{3}{4}=\frac{15}{20}$ (επί 5) \quad β) $\frac{5}{7}=\frac{35}{49}$ (επί 7) \quad γ) $\frac{18}{24}=\frac{3}{4}$ (διά 6) \quad δ) $\frac{5}{9}=\frac{20}{36}$ (επί 4).
\sol{3} α) Ναι: $4\cdot 15=60=6\cdot 10$. β) Όχι: $5\cdot 32=160$, $8\cdot 15=120$. γ) Ναι: $7\cdot 27=189=9\cdot 21$. δ) Ναι: $12\cdot 21=252=18\cdot 14$. ε) Όχι: $3\cdot 7=21$, $5\cdot 4=20$.
\sol{4} α) $\frac{18}{24}=\frac{3}{4}$ (ΜΚΔ 6) \quad β) $\frac{35}{42}=\frac{5}{6}$ (ΜΚΔ 7) \quad γ) $\frac{48}{72}=\frac{2}{3}$ (ΜΚΔ 24) \quad δ) $\frac{100}{250}=\frac{2}{5}$ (ΜΚΔ 50) \quad ε) $\frac{56}{98}=\frac{4}{7}$ (ΜΚΔ 14).
\sol{5} Ανάγωγα είναι τα $\frac{9}{16}$, $\frac{15}{28}$ και $\frac{25}{36}$. Τα υπόλοιπα: $\frac{12}{18}=\frac{2}{3}$, $\frac{21}{35}=\frac{3}{5}$, $\frac{17}{51}=\frac{1}{3}$.
\sol{6} α) $\frac{75}{100}$ \quad β) $\frac{35}{100}$ \quad γ) $\frac{40}{100}$ \quad δ) $\frac{18}{100}$ \quad ε) $\frac{52}{100}$.
\sol{7} α) $\frac{10}{12}$ \quad β) $\frac{9}{12}$ \quad γ) $\frac{8}{12}$ \quad δ) $\frac{6}{12}$.
\sol{8} α) $\frac{8}{12}$ και $\frac{9}{12}$ \quad β) $\frac{20}{24}$ και $\frac{21}{24}$ \quad γ) $\frac{8}{10}$, $\frac{5}{10}$ και $\frac{7}{10}$ \quad δ) ΕΚΠ(6, 8, 12) = 24, άρα $\frac{4}{24}$, $\frac{9}{24}$ και $\frac{10}{24}$.
\sol{9} α) Σωστό. β) Λάθος: $\frac{3}{5}=\frac{6}{10}$. γ) Λάθος: διαιρούμε και τους δύο όρους με τον ίδιο αριθμό. δ) Σωστό. ε) Λάθος: $\frac{6}{9}=\frac{2}{3}$ (ΜΚΔ 3). στ) Σωστό.
\sol{10} $\frac{6}{8}=\frac{3}{4}$, $\frac{9}{12}=\frac{3}{4}$, και $\frac{4}{6}=\frac{2}{3}$. Με κοινό παρονομαστή 12: Ρήνη $\frac{9}{12}$, Φανή $\frac{9}{12}$, Ηλίας $\frac{9}{12}$, Στέλιος $\frac{8}{12}$. Οι Ρήνη, Φανή και Ηλίας έφαγαν το ίδιο και περισσότερο από τον Στέλιο.
\sol{11} Έφαγε $\frac{18}{24}$, και με ΜΚΔ 6 γίνεται $\frac{3}{4}$ της σοκολάτας.
\sol{12} $\frac{45}{60}=\frac{3}{4}$ \quad $\frac{20}{60}=\frac{1}{3}$ \quad $\frac{40}{60}=\frac{2}{3}$ \quad $\frac{12}{60}=\frac{1}{5}$.
\sol{13} $\frac{250}{1000}=\frac{1}{4}$ \quad $\frac{750}{1000}=\frac{3}{4}$ \quad $\frac{400}{1000}=\frac{2}{5}$ \quad $\frac{125}{1000}=\frac{1}{8}$.
\sol{14} Χωράνε 2 όγδοα στο τέταρτο και 4 δέκατα έκτα: $\frac{1}{4}=\frac{2}{8}=\frac{4}{16}$.
\sol{15} Το $\frac{1}{2}$ είναι η μισή πίτσα, ενώ τα $\frac{4}{5}$ είναι σχεδόν ολόκληρη. Δεν είναι ίσα. Ο Στέλιος πρόσθεσε το 3 και στους δύο όρους, αλλά πρέπει να πολλαπλασιάσει: $\frac{1}{2}=\frac{1\cdot 3}{2\cdot 3}=\frac{3}{6}$.
\sol{16} Πολλαπλασιάζω τους όρους του $\frac{3}{5}$ με 1, 2, 3, \dots: $\frac{3}{5}$, $\frac{6}{10}$, $\frac{9}{15}$, $\frac{12}{20}$, $\frac{15}{25}$, $\frac{18}{30}$, $\frac{21}{35}$, $\frac{24}{40}$, $\frac{27}{45}$. Είναι 9 κλάσματα (το επόμενο, $\frac{30}{50}$, έχει παρονομαστή 50).
\sol{17} Το κλάσμα είναι $\frac{5k}{8k}$ και $5k+8k=13k=91$, άρα $k=7$ και το κλάσμα είναι $\frac{35}{56}$.
\sol{18} Το κλάσμα είναι $\frac{3k}{7k}$ και $7k-3k=4k=20$, άρα $k=5$ και το κλάσμα είναι $\frac{15}{35}$.
\sol{19} Αν τα 2 και 3 είχαν κοινό διαιρέτη $d>1$, τότε το $42\cdot d$ θα ήταν κοινός διαιρέτης των 84 και 126, και μεγαλύτερος από το 42. Αυτό είναι αδύνατο, γιατί το 42 είναι ο \emph{μέγιστος} κοινός διαιρέτης. Άρα $\MKD(2,\,3)=1$ και το κλάσμα είναι ανάγωγο.
\sol{20} Ισοδύναμα με το 0: $\frac{0}{3}$, $\frac{0}{7}$, $\frac{0}{12}$. Ισοδύναμα με το 1: $\frac{4}{4}$, $\frac{9}{9}$, $\frac{15}{15}$.
\sol{21} $\frac{2}{3}$, $\frac{4}{6}$, $\frac{6}{9}$, $\frac{8}{12}$: όλα είναι ίσα με $\frac{2}{3}$ (κάθε φορά πολλαπλασιάζονται οι όροι $2$ και $3$ με το ίδιο νούμερο).
\sol{22} Με ένα δίπλωμα: $\frac{1}{2}$ του φύλλου κάθε μέρος. Με δύο διπλώματα: $\frac{1}{4}$, και τα δύο μέρη μαζί $\frac{2}{4}$. Με τρία: $\frac{1}{8}$, και τέσσερα μέρη μαζί $\frac{4}{8}$. Άρα $\frac{1}{2}=\frac{2}{4}=\frac{4}{8}=\frac{8}{16}$.

\lyseis{2.3}{Σύγκριση κλασμάτων}
\sol{1} α) $\frac{5}{9}<\frac{7}{9}$ \quad β) $\frac{4}{5}>\frac{4}{9}$ \quad γ) $\frac{6}{11}<\frac{6}{7}$ \quad δ) $\frac{13}{20}>\frac{9}{20}$ \quad ε) $\frac{1}{3}>\frac{1}{4}$.
\sol{2} α) $\frac{10}{15}>\frac{9}{15}$, άρα $\frac{2}{3}>\frac{3}{5}$ \quad β) $\frac{15}{18}>\frac{14}{18}$, άρα $\frac{5}{6}>\frac{7}{9}$ \quad γ) $\frac{45}{60}>\frac{44}{60}$, άρα $\frac{3}{4}>\frac{11}{15}$ \quad δ) $\frac{14}{24}<\frac{15}{24}$, άρα $\frac{7}{12}<\frac{5}{8}$ \quad ε) $\frac{81}{90}>\frac{80}{90}$, άρα $\frac{9}{10}>\frac{8}{9}$.
\sol{3} Ίδιος αριθμητής: μεγαλύτερο είναι αυτό με τον μικρότερο παρονομαστή. Άρα $\frac{17}{10}>\frac{17}{11}>\frac{17}{12}>\frac{17}{13}>\frac{17}{14}$.
\sol{4} Μικρότερα από το 1: $\frac{7}{9}$, $\frac{4}{21}$, $\frac{99}{100}$. Ίσο με το 1: $\frac{12}{12}$. Μεγαλύτερα από το 1: $\frac{15}{11}$, $\frac{100}{99}$.
\sol{5} Με κοινό παρονομαστή 12: $\frac{9}{12}$, $\frac{10}{12}$, $\frac{8}{12}$, $\frac{7}{12}$, $\frac{11}{12}$. Σε αύξουσα σειρά: $\frac{7}{12}<\frac{2}{3}<\frac{3}{4}<\frac{5}{6}<\frac{11}{12}$.
\sol{6} $\frac{7}{3}$: $2<\frac{7}{3}<3$ (γιατί $2=\frac{6}{3}$). $\frac{11}{4}$: ανάμεσα στο 2 και 3. $\frac{25}{6}$: ανάμεσα στο 4 και 5 ($4\cdot 6=24$). $\frac{40}{9}$: ανάμεσα στο 4 και 5 ($4\cdot 9=36$, $5\cdot 9=45$). $\frac{100}{7}$: ανάμεσα στο 14 και 15 ($14\cdot 7=98$).
\sol{7} Με μονάδα 8 εκ.: το $\frac{1}{4}$ στα 2 εκ., το $\frac{1}{2}$ στα 4 εκ., το $\frac{5}{4}$ στα 10 εκ., το $\frac{3}{2}$ στα 12 εκ. και το $\frac{7}{4}$ στα 14 εκ. από το 0.
\sol{8} Α: $\frac{5}{6}$, Β: $\frac{7}{6}$ (που είναι $1+\frac{1}{6}$), Γ: $\frac{10}{6}$ (ίσο με $\frac{5}{3}$), Δ: $\frac{11}{6}$.
\sol{9} α) Σωστό. β) Λάθος: μεγαλύτερο είναι αυτό με τον μεγαλύτερο αριθμητή. γ) Σωστό. δ) Λάθος: $\frac{5}{5}=1$. ε) Σωστό. στ) Λάθος: $\frac{2}{3}=\frac{10}{15}>\frac{9}{15}=\frac{3}{5}$.
\sol{10} Κοινός παρονομαστής το $\EKP(5,\,9,\,4)=180$: $\frac{4}{5}=\frac{144}{180}$, $\frac{7}{9}=\frac{140}{180}$, $\frac{3}{4}=\frac{135}{180}$. Το πιο επιτυχημένο είναι αυτό με τα $\frac{4}{5}$ και το λιγότερο αυτό με τα $\frac{3}{4}$.
\sol{11} Κοινός παρονομαστής το $\EKP(5,\,7,\,6)=210$: $\frac{3}{5}=\frac{126}{210}$, $\frac{2}{7}=\frac{60}{210}$, $\frac{1}{6}=\frac{35}{210}$. Έτρεξε περισσότερο την πρώτη μέρα, μετά τη δεύτερη και λιγότερο την τρίτη.
\sol{12} $\frac{3}{4}=\frac{15}{20}$ και $\frac{7}{10}=\frac{14}{20}$. Επειδή $15>14$, χρειάζεται περισσότερο από όσο έχει: το πακέτο δεν φτάνει. Λείπει $\frac{1}{20}$ του κιλού.
\sol{13} Κοινός παρονομαστής 24: Πολυξένη $\frac{15}{24}$, Βασίλης $\frac{16}{24}$, Στέλιος $\frac{14}{24}$. Πιο φορτισμένο είναι του Βασίλη, μετά της Πολυξένης και τελευταίο του Στέλιου.
\sol{14} Κοινός παρονομαστής 9: Βασίλης $\frac{4}{9}$, Ζωή $\frac{1}{3}=\frac{3}{9}$, Μυρτώ $\frac{2}{9}$. Σειρά: Βασίλης, Ζωή, Μυρτώ.
\sol{15} $3\cdot 7=21$ και $5\cdot 4=20$, και $21>20$, άρα $\frac{3}{5}>\frac{4}{7}$. Δουλεύει γιατί με κοινό παρονομαστή το $5\cdot 7=35$ έχουμε $\frac{3}{5}=\frac{3\cdot 7}{35}=\frac{21}{35}$ και $\frac{4}{7}=\frac{4\cdot 5}{35}=\frac{20}{35}$: οι αριθμητές που συγκρίνουμε είναι ακριβώς τα χιαστί γινόμενα.
\sol{16} $\frac{2}{5}=\frac{8}{20}$ και $\frac{3}{5}=\frac{12}{20}$. Ανάμεσά τους: $\frac{9}{20}$, $\frac{10}{20}$, $\frac{11}{20}$.
\sol{17} $\frac{1}{3}=\frac{4}{12}$ και $\frac{5}{6}=\frac{10}{12}$. Ανάμεσα: $\frac{5}{12}$, $\frac{6}{12}$, $\frac{7}{12}$, $\frac{8}{12}$, $\frac{9}{12}$, δηλαδή 5 κλάσματα.
\sol{18} Το $\frac{10}{21}$ έχει μεγαλύτερους όρους από το $\frac{2}{3}$, αλλά είναι μικρότερο: $10\cdot 3=30<42=21\cdot 2$. Άρα οι μεγάλοι αριθμοί δεν σημαίνουν μεγάλο κλάσμα.
\sol{19} Όχι. Όσο κοντά κι αν είναι δύο κλάσματα, ανάμεσά τους χωράει άλλο: αν γράψουμε τα δύο κλάσματα με διπλάσιους (ή μεγαλύτερους) όρους, οι αριθμητές απομακρύνονται και χωράει ένας αριθμός ανάμεσα. Άρα δεν υπάρχει «επόμενο» κλάσμα, αντίθετα με τους φυσικούς.
\sol{20} Ενδεικτικά: ύπνος $\frac{8}{24}$, σχολείο $\frac{6}{24}$, διάβασμα $\frac{3}{24}$, παιχνίδι $\frac{3}{24}$, φαγητό και άλλα $\frac{4}{24}$. Με την ίδια ονομασία, συγκρίνεις απλώς τους αριθμητές: ύπνος, σχολείο, φαγητό, και μετά διάβασμα και παιχνίδι ισοδύναμα.
\sol{21} Με μονάδα 10 εκ.: $\frac{1}{5}$ στα 2 εκ., $\frac{3}{10}$ στα 3 εκ., $\frac{1}{2}$ στα 5 εκ., $\frac{7}{10}$ στα 7 εκ., $\frac{4}{5}$ στα 8 εκ., $\frac{6}{5}$ στα 12 εκ., $\frac{3}{2}$ στα 15 εκ. και $\frac{9}{5}$ στα 18 εκ.

\lyseis{Α2.4}{Πολλαπλασιασμός ακεραίων}
\sol{1} α) $6$ \quad β) $-24$ \quad γ) $42$ \quad δ) $-25$ \quad ε) $0$ \quad στ) $-207$.
\sol{2} α) θετικός. β) αρνητικός. γ) αρνητικών. δ) ετερόσημοι.
\sol{3} α) Πέντε αρνητικοί (περιττό πλήθος): $-1$. β) Τέσσερις αρνητικοί ($-1$, $-2$, $-4$, $-5$), άρτιο πλήθος, και $1\cdot2\cdot3\cdot4\cdot5=120$: $+120$. γ) Τρεις αρνητικοί, και $5\cdot10\cdot5\cdot10\cdot2=5000$: $-5000$.
\sol{4} Γραμμή $-5$: $-15$, $20$, $30$, $0$. Γραμμή $-1$: $-3$, $4$, $6$, $0$. Γραμμή $+2$: $6$, $-8$, $-12$, $0$. Γραμμή $7$: $21$, $-28$, $-42$, $0$.
\sol{5} α) $(+3)\cdot(-4)=-12$ και $-12=(-2)\cdot6$, άρα $6$. β) $(-5)\cdot(-6)=30=3\cdot10$, άρα $10$. γ) $9\cdot8=72=(-6)\cdot(-12)$, άρα $-12$. δ) $(-3)\cdot8=-24=12\cdot(-2)$, άρα $-2$.
\sol{6} α) $2\cdot(-2)+4\cdot5=-4+20=16$. β) $(-9)\cdot(-1)=9$ (γιατί $-7+11-13=-9$ και $-5+4=-1$). γ) $-18+(-4)-40=-62$ (ο όρος $2\cdot(-5)\cdot(-4)=40$ αφαιρείται). δ) $3\cdot[-8-(-10)]=3\cdot2=6$.
\sol{7} α) Σωστό. β) Λάθος: $+12$. γ) Σωστό. δ) Λάθος: οι $100$ αρνητικοί είναι άρτιο πλήθος, άρα $+1$. ε) Σωστό (αντιμεταθετική). στ) Λάθος: $(-2)\cdot(-3)=6$. ζ) Σωστό.
\sol{8} $\alpha\beta=-4$, $\beta\gamma=-20$, $\alpha\beta\gamma=(-4)\cdot(-5)=20$, $\alpha\gamma-\beta=5-4=1$, $\alpha-\beta\gamma=-1-(-20)=19$.
\sol{9} α) $4\cdot(-3)=-12$, δηλαδή $-12$°. β) $(-3)\cdot(-2)=+6$, δηλαδή $+6$° (είχε $6$ βαθμούς περισσότερο).
\sol{10} Ηλίας: $7\cdot5+4\cdot(-2)=35-8=27$ πόντους. Στέλιος: $3\cdot5+9\cdot(-2)=15-18=-3$ πόντους.
\sol{11} $6+5\cdot(-2)=6-10=-4$. Βρίσκεται στον όροφο $-4$ (τέταρτο υπόγειο).
\sol{12} Μετά από $6$ μέρες: $6\cdot(-15)=-90$, δηλαδή $90$ λίτρα λιγότερα. Πριν από $4$ μέρες: $(-15)\cdot(-4)=+60$, δηλαδή υπήρχαν $60$ λίτρα περισσότερα.
\sol{13} $(-6)\cdot(-1)=6$ και $(-6)\cdot(-2)=12$. Κάθε γραμμή είναι κατά $6$ μεγαλύτερη από την προηγούμενη, άρα το γινόμενο δύο αρνητικών είναι θετικό.
\sol{14} Αν $\alpha\beta>0$, οι $\alpha,\,\beta$ είναι ομόσημοι (και οι δύο θετικοί ή και οι δύο αρνητικοί). Αν $\alpha\beta<0$, είναι ετερόσημοι.
\sol{15} Άρτιο πλήθος αρνητικών: $0$, $2$ ή $4$.
\sol{16} Ο κανόνας «μείον επί μείον κάνει συν» ισχύει μόνο στον πολλαπλασιασμό. Στην πρόσθεση δύο αρνητικοί δίνουν αρνητικό: $(-3)+(-4)=-7$, ενώ $(-3)\cdot(-4)=+12$.
\sol{17} Οι αριθμοί είναι ετερόσημοι. Δοκιμάζω: $4\cdot(-3)=-12$ και $4+(-3)=1$. Αυτό είναι το ζεύγος $4$ και $-3$.
\sol{18} Ελεύθερη κατασκευή. Η γραμμή και η στήλη του $0$ έχουν παντού $0$. Στα κελιά με ομόσημους αριθμούς το γινόμενο είναι θετικό και στα κελιά με ετερόσημους αρνητικό.
\sol{19} Ελεύθερη απάντηση. Π.χ. αν μια ταινία που δείχνει κάποιον να κατεβαίνει $2$ ορόφους την ώρα γυρίσει ανάποδα (αρνητικός χρόνος), δείχνει να ανεβαίνει: $(-2)\cdot(-1)=+2$.

\lyseis{2.5}{Πολλαπλασιασμός κλασμάτων}
\sol{1} α) $\frac{10}{21}$ \quad β) $\frac{24}{36}=\frac{2}{3}$ \quad γ) $\frac{60}{150}=\frac{2}{5}$ \quad δ) $\frac{140}{210}=\frac{2}{3}$ \quad ε) $\frac{12}{72}=\frac{1}{6}$ \quad στ) $\frac{63}{168}=\frac{3}{8}$.
\sol{2} α) $\frac{8}{5}$ \quad β) $\frac{3\cdot 14}{7}=6$ \quad γ) $\frac{30}{9}=\frac{10}{3}$ \quad δ) $\frac{36}{8}=\frac{9}{2}$ \quad ε) $\frac{5\cdot 18}{6}=15$ \quad στ) $\frac{700}{10}=70$.
\sol{3} α) $\frac{1}{4}$ \quad β) $\frac{1}{15}$ \quad γ) $\frac{2}{3}\cdot\frac{3}{4}=\frac{1}{2}$ \quad δ) $\frac{3}{4}\cdot\frac{4}{5}=\frac{3}{5}$.
\sol{4} α) $\frac{7}{3}\cdot 3=7$ \quad β) $\frac{3}{2}\cdot\frac{8}{3}=4$ \quad γ) $\frac{13}{4}\cdot 2=\frac{13}{2}=6\frac{1}{2}$ \quad δ) $\frac{5}{2}\cdot\frac{8}{5}=4$.
\sol{5} $\frac{5}{7}$, $\frac{1}{3}$, $8$, $1$, $\frac{4}{11}$, και για το $2\frac{1}{2}=\frac{5}{2}$ ο αντίστροφος είναι $\frac{2}{5}$.
\sol{6} α) $\frac{4}{3}$ \quad β) $\frac{9}{5}$ \quad γ) $\frac{1}{6}$ \quad δ) $\frac{7}{2}$.
\sol{7} α) $\frac{14\cdot 25}{15\cdot 28}=\frac{350}{420}=\frac{5}{6}$ \quad β) $\frac{378}{945}=\frac{2}{5}$ \quad γ) $\frac{120}{144}=\frac{5}{6}$ \quad δ) $\frac{140}{252}=\frac{5}{9}$.
\sol{8} α) $\frac{2}{3}+\frac{3}{8}=\frac{16}{24}+\frac{9}{24}=\frac{25}{24}=1\frac{1}{24}$ \quad β) $\frac{7}{6}\cdot\frac{3}{4}=\frac{21}{24}=\frac{7}{8}$ \quad γ) $\frac{1}{2}\cdot\left(\frac{9}{12}-\frac{2}{12}\right)=\frac{1}{2}\cdot\frac{7}{12}=\frac{7}{24}$.
\sol{9} α) Σωστό. β) Λάθος: $\frac{3}{5}\cdot 4=\frac{12}{5}$. γ) Λάθος: $\frac{1}{2}\cdot\frac{1}{3}=\frac{1}{6}$. δ) Λάθος: το 0 δεν έχει αντίστροφο. ε) Σωστό. στ) Λάθος: $2\frac{1}{2}\cdot 2=\frac{5}{2}\cdot 2=5$.
\sol{10} $\frac{3}{4}\cdot 2\frac{1}{2}=\frac{3}{4}\cdot\frac{5}{2}=\frac{15}{8}=1\frac{7}{8}$ κιλά αλεύρι.
\sol{11} $\frac{2}{5}\cdot 1\frac{1}{2}=\frac{2}{5}\cdot\frac{3}{2}=\frac{3}{5}$ του λίτρου.
\sol{12} $\frac{3}{5}\cdot\frac{2}{3}=\frac{6}{15}=\frac{2}{5}$ του κήπου είναι ντομάτες. Σε $300$ τ.μ.: $\frac{2}{5}\cdot 300=120$ τ.μ.
\sol{13} $3\frac{1}{2}\cdot 2\frac{2}{5}=\frac{7}{2}\cdot\frac{12}{5}=\frac{84}{10}=\frac{42}{5}=8\frac{2}{5}$ τ.μ.
\sol{14} $4\frac{1}{2}\cdot\frac{2}{3}=\frac{9}{2}\cdot\frac{2}{3}=3$ χλμ.
\sol{15} Η έκπτωση είναι $\frac{3}{10}\cdot 80=24$ €. Πληρώνει $80-24=56$ €, που είναι και τα $\frac{7}{10}$ της τιμής: $\frac{7}{10}\cdot 80=56$.
\sol{16} Το μισό από το ένα τρίτο είναι $\frac{1}{2}\cdot\frac{1}{3}=\frac{1}{6}$, δηλαδή το ένα τρίτο χωρισμένο στη μέση. Είναι μικρότερο και από το $\frac{1}{3}$ και από το $\frac{1}{2}$. Ο πολλαπλασιασμός με κλάσμα μικρότερο του 1 δίνει λιγότερο από ό,τι είχες.
\sol{17} $12\cdot\frac{3}{4}=9$ (μικρότερο από 12) και $12\cdot\frac{5}{4}=15$ (μεγαλύτερο από 12). Όταν πολλαπλασιάζεις με κλάσμα μικρότερο του 1 ο αριθμός μικραίνει, και όταν πολλαπλασιάζεις με κλάσμα μεγαλύτερο του 1 μεγαλώνει.
\sol{18} Ο αντίστροφος του $\frac{3}{8}$ είναι ο $\frac{8}{3}$. Το κλάσμα που ψάχνουμε $x$ ικανοποιεί $x\cdot\frac{3}{8}=\frac{1}{2}$. Πολλαπλασιάζοντας και τα δύο μέλη με το $\frac{8}{3}$ βρίσκουμε $x=\frac{1}{2}\cdot\frac{8}{3}=\frac{4}{3}$. Έλεγχος: $\frac{4}{3}\cdot\frac{3}{8}=\frac{12}{24}=\frac{1}{2}$.
\sol{19} Οι αριθμητές και οι παρονομαστές απλοποιούνται διαδοχικά: $\frac{1\cdot 2\cdot 3\cdot 4\cdot 5}{2\cdot 3\cdot 4\cdot 5\cdot 6}=\frac{1}{6}$. Το γινόμενο ως το $\frac{9}{10}$ βγαίνει $\frac{1}{10}$.
\sol{20} Αν ο αριθμός $x$ ήταν αντίστροφος του 0, θα έπρεπε $0\cdot x=1$. Μα $0\cdot x=0$ για κάθε $x$, άρα η εξίσωση δεν έχει λύση. Το 0 δεν έχει αντίστροφο.
\sol{21} Με τις 4 στήλες και τις 3 σειρές το φύλλο έχει 12 ίσα μέρη. Η σκιασμένη περιοχή (3 στήλες) έχει 9 μέρη και από αυτά οι 2 σειρές είναι 6. Άρα τα σκούρα είναι $\frac{6}{12}=\frac{1}{2}$ του φύλλου. Και $\frac{2}{3}\cdot\frac{3}{4}=\frac{6}{12}=\frac{1}{2}$.
\sol{22} Κάθε «μισό του» διαιρεί με το 2: $80:2=40$, $20$, $10$, $5$. Άρα $\left(\frac{1}{2}\right)^4\cdot 80=\frac{80}{16}=5$.

\lyseis{2.6}{Διαίρεση κλασμάτων}
\sol{1} α) $\frac{3}{5}\cdot 10=6$ \quad β) $\frac{5}{6}\cdot\frac{12}{5}=2$ \quad γ) $\frac{7}{8}\cdot\frac{4}{3}=\frac{7}{6}=1\frac{1}{6}$ \quad δ) $\frac{4}{9}\cdot\frac{27}{8}=\frac{3}{2}$ \quad ε) $\frac{2}{3}\cdot\frac{1}{2}=\frac{1}{3}$ \quad στ) $6\cdot\frac{4}{3}=8$.
\sol{2} α) $12\cdot\frac{3}{2}=18$ \quad β) $5\cdot 4=20$ \quad γ) $\frac{3}{7}\cdot\frac{1}{3}=\frac{1}{7}$ \quad δ) $10\cdot\frac{2}{5}=4$ \quad ε) $\frac{1}{2}\cdot 6=3$ \quad στ) $\frac{9}{10}\cdot\frac{5}{3}=\frac{3}{2}$.
\sol{3} α) $1:\frac{1}{4}=4$ φορές \quad β) $\frac{2}{3}:\frac{1}{3}=2$ φορές \quad γ) $\frac{5}{12}:\frac{1}{6}=\frac{5}{12}\cdot 6=\frac{5}{2}$ φορές (δύο φορές και μισή) \quad δ) $2:\frac{2}{5}=5$ φορές.
\sol{4} α) $\frac{5}{2}\cdot\frac{6}{5}=3$ \quad β) $\frac{10}{3}\cdot\frac{1}{2}=\frac{5}{3}=1\frac{2}{3}$ \quad γ) $4\cdot\frac{3}{4}=3$ \quad δ) $\frac{7}{4}\cdot\frac{2}{5}=\frac{14}{20}=\frac{7}{10}$.
\sol{5} α) $\frac{2}{3}\cdot\frac{9}{4}=\frac{18}{12}=\frac{3}{2}$ \quad β) $5\cdot\frac{6}{5}=6$ \quad γ) $\frac{7}{8}\cdot\frac{1}{7}=\frac{1}{8}$ \quad δ) $\frac{3}{4}\cdot\frac{10}{9}=\frac{30}{36}=\frac{5}{6}$.
\sol{6} Ψάχνω $x$ με $x\cdot\frac{3}{5}=\frac{9}{10}$, άρα $x=\frac{9}{10}:\frac{3}{5}=\frac{9}{10}\cdot\frac{5}{3}=\frac{3}{2}$. Έλεγχος: $\frac{3}{2}\cdot\frac{3}{5}=\frac{9}{10}$.
\sol{7} α) Το όλο είναι τρία «τρίτα»: $\frac{7}{9}:\frac{1}{3}=\frac{7}{9}\cdot 3=\frac{7}{3}=2\frac{1}{3}$. β) $\frac{6}{7}:\frac{2}{5}=\frac{6}{7}\cdot\frac{5}{2}=\frac{15}{7}=2\frac{1}{7}$.
\sol{8} α) $\left(\frac{1}{2}:\frac{1}{4}\right)\cdot\frac{1}{3}=2\cdot\frac{1}{3}=\frac{2}{3}$ \quad β) $\frac{1}{2}:\frac{1}{12}=6$ \quad γ) $1:\frac{1}{2}=2$ \quad δ) $\frac{3}{4}+\left(\frac{1}{4}:\frac{1}{2}\right)=\frac{3}{4}+\frac{1}{2}=\frac{5}{4}$.
\sol{9} α) Λάθος: $\frac{3}{4}:\frac{1}{2}=\frac{3}{4}\cdot 2=\frac{3}{2}$. β) Σωστό. γ) Σωστό. δ) Σωστό. ε) Λάθος: $\frac{1}{2}:3=\frac{1}{6}$. στ) Λάθος: αντιστρέφω μόνο τον διαιρέτη.
\sol{10} $\frac{5}{6}:\frac{1}{12}=\frac{5}{6}\cdot 12=10$ μπουκαλάκια. Με μπουκαλάκια των $\frac{1}{5}$: $\frac{5}{6}:\frac{1}{5}=\frac{25}{6}=4\frac{1}{6}$: 4 γεμάτα και το $\frac{1}{6}$ του πέμπτου μπουκαλιού.
\sol{11} $\frac{9}{2}:\frac{3}{4}=\frac{9}{2}\cdot\frac{4}{3}=6$ κομμάτια.
\sol{12} Σε μια ώρα (4 φορές το $\frac{1}{4}$ της ώρας) θα έτρεχε $4\cdot 2\frac{1}{2}=10$ χλμ. Με διαίρεση: $2\frac{1}{2}:\frac{1}{4}=\frac{5}{2}\cdot 4=10$ χλμ. την ώρα.
\sol{13} $3\frac{1}{2}:\frac{7}{4}=\frac{7}{2}\cdot\frac{4}{7}=2$ βάρδιες.
\sol{14} $3:\frac{3}{8}=3\cdot\frac{8}{3}=8$ μερίδες.
\sol{15} $\frac{3}{4}:6=\frac{3}{4}\cdot\frac{1}{6}=\frac{3}{24}=\frac{1}{8}$ του κιλού ανά άτομο.
\sol{16} Ο Στέλιος αντέστρεψε και τα δύο κλάσματα. Αντιστρέφουμε μόνο τον διαιρέτη: $\frac{3}{4}:\frac{1}{8}=\frac{3}{4}\cdot 8=6$.
\sol{17} $\frac{2}{3}:\frac{3}{2}=\frac{2}{3}\cdot\frac{2}{3}=\frac{4}{9}$. $\frac{3}{5}:\frac{5}{3}=\frac{3}{5}\cdot\frac{3}{5}=\frac{9}{25}$. Βγαίνει το τετράγωνο του κλάσματος: $\frac{\alpha}{\beta}:\frac{\beta}{\alpha}=\frac{\alpha^2}{\beta^2}$.
\sol{18} Το $6:\frac{1}{2}$ ρωτάει πόσα μισά χωράνε στο 6. Σε κάθε ολόκληρο χωράνε 2 μισά, άρα στα 6 χωράνε 12. Τα μισά είναι μικρότερα από το 1, γι' αυτό χωράνε περισσότερα.
\sol{19} $\frac{1}{14}+\frac{1}{42}=\frac{3}{42}+\frac{1}{42}=\frac{4}{42}=\frac{2}{21}$ και $\frac{2}{3}\cdot\frac{1}{7}=\frac{2}{21}$. Για το $\frac{2}{3}\cdot\frac{1}{5}=\frac{2}{15}$: $\frac{1}{10}+\frac{1}{30}=\frac{3}{30}+\frac{1}{30}=\frac{4}{30}=\frac{2}{15}$.
\sol{20} $x\cdot\frac{4}{5}=12$ δίνει $x=12:\frac{4}{5}=12\cdot\frac{5}{4}=15$. Και $x:\frac{4}{5}=12$ δίνει $x=12\cdot\frac{4}{5}=\frac{48}{5}=9\frac{3}{5}$.
\sol{21} $1\frac{1}{2}:\frac{1}{4}=\frac{3}{2}\cdot 4=6$ ποτήρια.
\sol{22} Αριθμητής: $\frac{1}{2}+\frac{1}{3}=\frac{5}{6}$. Παρονομαστής: $\frac{1}{2}-\frac{1}{3}=\frac{1}{6}$. Άρα $\frac{5}{6}:\frac{1}{6}=5$.

\lyseisE{Επανάληψη κεφαλαίου Α2}
\solr{Α2}{1} α) $-6$ \quad β) $+40$ \quad γ) $-3$ \quad δ) $-1$.
\solr{Α2}{2} Θετικοί: $12$, $+3$. Αρνητικοί: $-5$, $-1$, $-20$. Φυσικοί: $0$, $12$, $+3$. Ομόσημοι π.χ. $-5$ και $-1$ ή $12$ και $+3$. Ετερόσημοι π.χ. $-5$ και $12$. Το $0$ δεν έχει πρόσημο.
\solr{Α2}{3} $\mathrm{Α}(-5)$, $\mathrm{Β}(-2)$, $\Gamma(1)$, $\mathrm{Δ}(4)$. Ο $-3$ βρίσκεται ανάμεσα στο $\mathrm{Α}$ και στο $\mathrm{Β}$, μία μονάδα αριστερά του $\mathrm{Β}$.
\solr{Α2}{4} α) $9$ \quad β) $4$ \quad γ) $+15$ \quad δ) $\alpha=6$ ή $\alpha=-6$.
\solr{Α2}{5} $-12<-7<-1<0<3<5$.
\solr{Α2}{6} Οι $-3,\,-2,\,-1,\,0,\,1,\,2$ (ο $-4$ δεν μπαίνει, ο $2$ μπαίνει).
\solr{Α2}{7} α) $5$ \quad β) $-15$ \quad γ) $0$ \quad δ) $-12$.
\solr{Α2}{8} α) $5+3=8$ \quad β) $-4+(-6)=-10$ \quad γ) $0$ \quad δ) $0+9=9$.
\solr{Α2}{9} $15+4-9-6+10=14$.
\solr{Α2}{10} α) $42$ \quad β) $-20$ \quad γ) $-24$ (τρεις αρνητικοί παράγοντες) \quad δ) $0$.
\solr{Α2}{11} α) $-32$ \quad β) $-16$ \quad γ) $9$ \quad δ) $1$.
\solr{Α2}{12} $3\cdot4-4\cdot(-1)+(-3)^2=12+4+9=25$.
\solr{Α2}{13} $-4+2\cdot(-3)+9=-4-6+9=-1$. Η τελική θερμοκρασία είναι $-1$°.
\solr{Α2}{14} $6\cdot5+5\cdot(-3)=30-15=15$ πόντους.
\solr{Α2}{15} α) Λάθος: $5$. β) Σωστό. γ) Λάθος: $5-(-2)=7$. δ) Σωστό. ε) Σωστό. στ) Σωστό.
\solr{Α2}{16} α) Σωστό. β) Λάθος: $-16$. γ) Σωστό. δ) Λάθος: $1\neq100$. ε) Λάθος: ο $-3$ είναι ακέραιος αλλά όχι φυσικός. στ) Λάθος: $-6\neq9$.


\lyseiskef{3}{Οι δεκαδικοί αριθμοί}
\lyseis{Α3.1}{Η έννοια του ρητού}
\sol{1} Ο $5$: φυσικός, ακέραιος, ρητός. Ο $-1{,}7$: μόνο ρητός. Ο $\frac{2}{3}$: μόνο ρητός. Ο $-10$: ακέραιος και ρητός. Το $0$: φυσικός, ακέραιος, ρητός. Ο $\frac{12}{4}=3$: φυσικός, ακέραιος, ρητός.
\sol{2} α) Σωστό: $-\frac{6}{2}=-3$. β) Σωστό. γ) Λάθος: ο $\frac{1}{2}$ είναι ρητός αλλά όχι ακέραιος. δ) Σωστό. ε) Λάθος: το $\frac{15}{7}$ είναι το πηλίκο $15:7$. στ) Σωστό: $0=\frac{0}{1}$.
\sol{3} α) $20:4=5$ και $3\cdot5=15$. β) $35:5=7$ και $7\cdot7=49$. γ) $24:6=4$ (και $18:6=3$). δ) $36:4=9$ άρα το $20$ διαιρέθηκε με το $4$: $20:4=5$.
\sol{4} α) $5\cdot9=45=3\cdot15$: ισοδύναμα. β) $4\cdot9=36=6\cdot6$: ισοδύναμα. γ) $16\cdot35=560=7\cdot80$: ισοδύναμα. δ) $21\cdot11=231$ και $13\cdot19=247$: όχι. ε) $14\cdot25=350$ και $15\cdot24=360$: όχι.
\sol{5} Το $\frac{60}{15}=4$ απλοποιείται. Ανάγωγα είναι τα $\frac{24}{35}$ (αφού $24=2^3\cdot3$ και $35=5\cdot7$) και $\frac{19}{6}$. Τα άλλα: $\frac{12}{15}=\frac{4}{5}$ και $\frac{3}{18}=\frac{1}{6}$.
\sol{6} α) $0{,}375$ \quad β) $0{,}35$ \quad γ) $4{,}8$ \quad δ) $-2{,}75$ \quad ε) $0{,}8333\ldots$ \quad στ) $2{,}08$.
\sol{7} α) $\frac{45}{10}=\frac{9}{2}$ \quad β) $\frac{31}{100}$ \quad γ) $\frac{441}{10}$ \quad δ) $-\frac{25}{10}=-\frac{5}{2}$ \quad ε) $\frac{2}{100}=\frac{1}{50}$ \quad στ) $-\frac{1\,125}{1\,000}=-\frac{9}{8}$.
\sol{8} $\mathrm{Α}=-2{,}5=-\frac{5}{2}$, $\mathrm{Β}=-0{,}5=-\frac{1}{2}$, $\Gamma=0{,}5=\frac{1}{2}$, $\Delta=1{,}5=\frac{3}{2}$, $\mathrm{Ε}=2{,}5=\frac{5}{2}$. Κάθε μονάδα έχει διαιρεθεί στη μέση, γι' αυτό όλα τα σημεία είναι «μισά».
\sol{9} Έφαγε $\frac{18}{24}=\frac{18:6}{24:6}=\frac{3}{4}=0{,}75$ της σοκολάτας.
\sol{10} $\frac{45}{60}=\frac{3}{4}$, \ $\frac{20}{60}=\frac{1}{3}$, \ $\frac{40}{60}=\frac{2}{3}$, \ $\frac{12}{60}=\frac{1}{5}$.
\sol{11} Πατάτες $2{,}75=\frac{275}{100}=\frac{11}{4}$ κιλά. Τυρί $0{,}4=\frac{4}{10}=\frac{2}{5}$ του κιλού. Ντομάτες $1{,}25=\frac{125}{100}=\frac{5}{4}$ κιλά.
\sol{12} Κλάσμα $\frac{35}{100}$, δεκαδικός $0{,}35$ και ανάγωγο $\frac{35:5}{100:5}=\frac{7}{20}$.
\sol{13} Φούρνος: $-\frac{1}{2}=-0{,}5$, αριστερά του σπιτιού, $2$ cm αριστερά του $0$. Σχολείο: $\frac{3}{2}=1{,}5$, δεξιά του σπιτιού, $6$ cm δεξιά του $0$. Περίπτερο: $-0{,}25$, αριστερά, $1$ cm αριστερά του $0$. Ο φούρνος και το περίπτερο είναι στην ίδια πλευρά.
\sol{14} Με δεκαδικούς: $\frac{1}{2}=0{,}5$ ενώ $\frac{4}{5}=0{,}8$, άρα δεν είναι ίσα. Στην πίτσα: το μισό δεν είναι ίσο με τα τέσσερα πέμπτα. Σωστά: πολλαπλασιάζω και τους δύο όρους επί $5$: $\frac{1}{2}=\frac{1\cdot5}{2\cdot5}=\frac{5}{10}$.
\sol{15} Αν οι όροι είναι $3k$ και $7k$, η διαφορά τους είναι $4k=20$, άρα $k=5$ και το κλάσμα είναι το $\frac{15}{35}$. Έλεγχος: $15\cdot7=105=35\cdot3$ και $35-15=20$.
\sol{16} $\frac{1}{3}=0{,}333\ldots$ και $\frac{2}{3}=0{,}666\ldots$ Το ψηφίο επαναλαμβάνεται χωρίς τέλος. Δεν γράφεται ως δεκαδικός που τελειώνει, γιατί η διαίρεση $1:3$ δεν δίνει ποτέ υπόλοιπο $0$.
\sol{17} Για παράδειγμα $-\frac{3}{4}=\frac{-3}{4}=\frac{3}{-4}$ και ισοδύναμο $-\frac{6}{8}$ και δεκαδικός $-0{,}75$. Βρίσκεται αριστερά του $0$, στα τρία τέταρτα της διαδρομής προς το $-1$, ανάμεσα στο $-1$ και στο $-0{,}5$.
\sol{18} Ελεύθερη απάντηση. Ο μαθητής πρέπει να γράψει σωστά τα ισοδύναμα, π.χ. $\frac{1}{2}$ φλιτζάνι $=0{,}5=\frac{2}{4}$ και $\frac{3}{4}$ του κιλού $=0{,}75=\frac{6}{8}$.
\sol{19} Ελεύθερη κατασκευή. Οι αρνητικοί και οι αριθμοί με κλασματικό ή δεκαδικό μέρος (όπως $-3$, $0{,}5$, $\frac{1}{3}$) δεν μπαίνουν στους φυσικούς, και όσοι δεν είναι ακέραιοι (όπως $0{,}5$ και $\frac{1}{3}$) δεν μπαίνουν στους ακέραιους. Όλοι μπαίνουν στους ρητούς.

\lyseis{Α3.2}{Απόλυτη τιμή, σύγκριση, διάταξη, αντίθετοι}
\sol{1} α) $-\alpha$ \quad β) μεγαλύτερη \quad γ) μικρότερη \quad δ) αριθμητή \quad ε) μικρότερο.
\sol{2} Στη στήλη του $-3$: αντίθετος $3$, απόλυτη τιμή $3$. Για τον $0{,}35$: $-0{,}35$ και $0{,}35$. Για τον $\frac{1}{9}$: $-\frac{1}{9}$ και $\frac{1}{9}$. Στην τέταρτη στήλη ο αντίθετος είναι ο $-0{,}1$, άρα ο αριθμός είναι ο $0{,}1$ και η απόλυτη τιμή $0{,}1$. Στην πέμπτη ο αντίθετος είναι ο $\frac{7}{4}$, άρα ο αριθμός είναι ο $-\frac{7}{4}$ και η απόλυτη τιμή $\frac{7}{4}$. Για τον $-\frac{2}{5}$: αντίθετος $\frac{2}{5}$ και απόλυτη τιμή $\frac{2}{5}$.
\sol{3} α) $>$ (ομώνυμα, $21>12$). β) $<$ (ίδιος αριθμητής, μεγαλύτερος παρονομαστής, άρα μικρότερο κλάσμα). γ) $\frac{2}{3}=\frac{16}{24}<\frac{21}{24}=\frac{7}{8}$, άρα $<$. δ) $>$ (θετικός και αρνητικός). ε) $-0{,}8>-1$, άρα $>$. στ) $-\frac{8}{3}=-2{,}666\ldots>-3$, άρα $>$. ζ) $\frac{9}{8}=1{,}125<1{,}2$, άρα $<$.
\sol{4} α) $3{,}16<3{,}20$, άρα $<$. β) $0{,}092>0{,}087$ (στα εκατοστά $9>8$), άρα $>$. γ) $-1{,}36<-0{,}01$ (αρνητικοί, ο $-1{,}36$ έχει μεγαλύτερη απόλυτη τιμή), άρα $<$. δ) $<$ (αρνητικός και θετικός). ε) $-1{,}3=-1{,}30>-1{,}33$, άρα $>$. στ) $-1{,}10=-1{,}1$, άρα $=$.
\sol{5} $-3<-0{,}03<0<0{,}2<1<1{,}37$.
\sol{6} α) $-\frac{5}{4}<-\frac{3}{4}<-\frac{1}{4}<\frac{1}{4}<\frac{3}{4}<\frac{5}{4}$. β) Οι αρνητικοί: $-\frac{5}{2}=-2{,}5$, $-\frac{5}{3}=-1{,}666\ldots$, $-\frac{5}{7}=-0{,}714\ldots$ Οι θετικοί, με ίδιο αριθμητή: όσο μεγαλύτερος ο παρονομαστής, τόσο μικρότερο το κλάσμα: $\frac{5}{8}<\frac{5}{4}<\frac{5}{3}$. Άρα $-\frac{5}{2}<-\frac{5}{3}<-\frac{5}{7}<\frac{5}{8}<\frac{5}{4}<\frac{5}{3}$.
\sol{7} α) Μικρότερα του $1$ είναι τα $\frac{3}{16},\ \frac{7}{16},\ \frac{15}{16}$ και $\frac{3}{16}<\frac{7}{16}<\frac{15}{16}$. β) Μεγαλύτερα του $1$ είναι τα $\frac{9}{8}=1{,}125$, $\frac{5}{4}=1{,}25$, $\frac{9}{4}=2{,}25$ και $\frac{9}{8}<\frac{5}{4}<\frac{9}{4}$. γ) $\frac{3}{16}<\frac{7}{16}<\frac{15}{16}<\frac{9}{8}<\frac{5}{4}<\frac{9}{4}$.
\sol{8} α) Σωστό. β) Λάθος: $|-5|=5$. γ) Σωστό. δ) Λάθος: $-9<-2$. ε) Σωστό (ίδιος αριθμητής, μεγαλύτερος παρονομαστής, μικρότερο κλάσμα). στ) Σωστό.
\sol{9} Από την πιο κρύα: Φλώρινα $-10$, Μέτσοβο $-5{,}5$, Ιωάννινα $-3$, Θεσσαλονίκη $-1{,}5$, Αθήνα $7$.
\sol{10} α) Ύπνος $24:3=8$ ώρες, Σχολείο $24:4=6$ ώρες, Διάβασμα $24:8=3$ ώρες, Διασκέδαση $24:6=4$ ώρες. β) Ύπνος ($8$), Σχολείο ($6$), Διασκέδαση ($4$), Διάβασμα ($3$). γ) Μένουν $24-21=3$ ώρες, δηλαδή $\frac{3}{24}=\frac{1}{8}$ του εικοσιτετραώρου. Έλεγχος με κλάσματα: $\frac{8}{24}+\frac{6}{24}+\frac{3}{24}+\frac{4}{24}=\frac{21}{24}$.
\sol{11} Κέρδισε ο Βασίλης: $9{,}75<9{,}80$ και ο μικρότερος χρόνος είναι ο καλύτερος. Το λάθος του: συνέκρινε το $75$ με το $8$ σαν φυσικούς. Τα $9{,}8$ είναι $9{,}80$, δηλαδή $80$ εκατοστά, και $80>75$.
\sol{12} Το σχολείο είναι στο $0$. Το βιβλιοπωλείο ($-2$) απέχει $2$ μονάδες, δηλαδή $200$ μ. Ο φούρνος ($+3$) απέχει $300$ μ. Πιο κοντά είναι το βιβλιοπωλείο. Τα δύο μαγαζιά απέχουν $2+3=5$ μονάδες, δηλαδή $500$ μ.
\sol{13} $-11\,034<-430<0<2\,917$: τάφρος Μαριανών, Νεκρά Θάλασσα, επιφάνεια θάλασσας, Όλυμπος.
\sol{14} $\frac{2}{5}=0{,}4$ και $\frac{3}{5}=0{,}6$. Ανάμεσα βρίσκονται π.χ. οι $0{,}45$, $0{,}5=\frac{1}{2}$ και $0{,}55$. Υπάρχουν άπειροι άλλοι: δεν υπάρχει «επόμενος».
\sol{15} α) Π.χ. $-0{,}75$ και $-0{,}72$ (ο $-0{,}8<-0{,}75<-0{,}7$). β) $\frac{3}{10}=0{,}3$ και $\frac{2}{5}=0{,}4$: π.χ. $0{,}35=\frac{7}{20}$ και $0{,}33$. γ) $\frac{7}{9}=0{,}777\ldots$ και $\frac{5}{6}=0{,}8333\ldots$: π.χ. $0{,}8=\frac{4}{5}$ και $0{,}79$.
\sol{16} $|\alpha|=\alpha$ για κάθε θετικό και για το $0$. $|\alpha|=-\alpha$ για κάθε αρνητικό και για το $0$. Ο $\alpha$ και ο $-\alpha$ είναι συμμετρικοί ως προς το $0$, άρα απέχουν το ίδιο από το $0$ και $|\alpha|=|-\alpha|$.
\sol{17} α) $0$ και $1$. β) $-2$ και $-1$. γ) $\frac{7}{2}=3{,}5$: $3$ και $4$. δ) $-\frac{11}{3}=-3{,}666\ldots$: $-4$ και $-3$. ε) $0$ και $1$ (είναι λίγο μικρότερος του $1$).
\sol{18} Ελεύθερη απάντηση. Ο μαθητής πρέπει να παρατηρήσει ότι ο πιο αριστερός αριθμός στον άξονα είναι ο μικρότερος και να τοποθετήσει σωστά τα πρόσημα.
\sol{19} Ελεύθερη απάντηση. Π.χ. $8$ ώρες ύπνου: $\frac{8}{24}=\frac{1}{3}$ και $9$ ώρες: $\frac{9}{24}=\frac{3}{8}$. Για να συγκριθούν: $\frac{1}{3}=\frac{8}{24}<\frac{9}{24}=\frac{3}{8}$.

\lyseis{3.3}{Υπολογιστής τσέπης}
\sol{1} α) $365{,}852$ \quad β) $912{,}346$ \quad γ) $75{,}03$ \quad δ) $570$.
\sol{2} α) $50\cdot6=300$, άρα $287{,}33$. β) $20:5=4$, άρα $4{,}4$.
\sol{3} α) \key{(}\key{4.5}\kpl\key{2.7}\key{)}\kpo\key{1.25}\kis\ $\to9$. β) \key{36}\kdi\key{(}\key{0.9}\kpl\key{2.1}\key{)}\kis\ $\to12$. (Αν δεν έχει παρενθέσεις: πρώτα $0{,}9+2{,}1=3$ και μετά $36:3$.)
\sol{4} α) $6{,}5\cdot4-3{,}2\cdot5=26-16=10$. β) $(1{,}2+0{,}8):0{,}5=2:0{,}5=4$.
\sol{5} \key{MC} \key{2.5}\kpo\key{8}\kis\key{M+} \key{1.5}\kpo\key{6}\kis\key{M+} \key{4.25}\kpo\key{2}\kis\key{M$-$} \key{MR}. Αποτέλεσμα $20+9-8{,}5=20{,}5$.
\sol{6} α) $10-2{,}5=7{,}5$ \quad β) $12{,}6:2{,}1=6$ \quad γ) $0{,}25\cdot6{,}8+1{,}2=1{,}7+1{,}2=2{,}9$.
\sol{7} α) Λάθος: κάνει σωστά ό,τι του πατήσουμε, και το δεδομένο μπορεί να είναι λάθος. β) Λάθος: δείχνει προσέγγιση. γ) Σωστό. δ) Λάθος: τα απλά εκτελούν με τη σειρά πληκτρολόγησης. ε) Λάθος: η εκτίμηση βρίσκει τα λάθη πληκτρολόγησης.
\sol{8} $57{,}6-18{,}35=39{,}25$ €. Εκτίμηση: $58-18=40$.
\sol{9} Έβαλε $67{,}5$ αντί για $6{,}75$ (η υποδιαστολή έφυγε μία θέση δεξιά): $67{,}5\cdot1{,}2=81$. Το σωστό είναι $6{,}75\cdot1{,}2=8{,}1$. Εκτίμηση: $7\cdot1=7$.
\sol{10} $312:100\cdot7{,}8=3{,}12\cdot7{,}8=24{,}336$ λίτρα, δηλαδή $24{,}3$ λίτρα.
\sol{11} Η οθόνη δείχνει $14{,}285714$ (ή $14{,}285714286$ με περισσότερα ψηφία). Στρογγυλοποίηση: $14{,}29$ €. Συνολικά $7\cdot14{,}29=100{,}03$ €, δηλαδή τρία λεπτά παραπάνω από το ταμείο.
\sol{12} $0{,}06\cdot4{,}5=0{,}27$ kWh την ημέρα. Σε $30$ ημέρες $8{,}1$ kWh. Κόστος $8{,}1\cdot0{,}18=1{,}458\approx1{,}46$ €.
\sol{13} Το $1:3$ στην οθόνη δείχνει $0{,}33333333$ (προσέγγιση). Άρα $3\cdot0{,}33333333=0{,}99999999\neq1$. Η ακριβής τιμή $\frac{1}{3}$ έχει άπειρα ψηφία.
\sol{14} Εκτέλεσε τις πράξεις με τη σειρά που πατήθηκαν: $2+3=5$ και μετά $5\cdot4=20$. Για να βγει το σωστό πατάω πρώτα το γινόμενο: \key{3}\kpo\key{4}\kpl\key{2}\kis{} (δίνει $12+2=14$), ή υπολογίζω το $3\cdot4$ και το φυλάω με \key{M+}.
\sol{15} Ο $0{,}98$ είναι μικρότερος από το $1$, άρα το $487\cdot0{,}98$ είναι λίγο μικρότερο από το $487$, δηλαδή περίπου $480$ (βγαίνει $477{,}26$). Το $4\,773{,}3$ είναι δέκα φορές μεγαλύτερο.
\sol{16} $25$ φορές, γιατί $2{,}5:0{,}1=25$.
\sol{17} Η οθόνη δείχνει $0{,}14285714$ (οκτώ ψηφία). Στα εκατοστά: $0{,}14$. Τα ψηφία $142857$ επαναλαμβάνονται συνεχώς, επομένως $\frac{1}{7}=0{,}142857142857\ldots$.
\sol{18} Ελεύθερη απάντηση. Προσοχή να στοιχίσεις σωστά τα δεκαδικά και να ελέγξεις με εκτίμηση.
\sol{19} Ελεύθερη απάντηση. Το άθροισμα των θερμοκρασιών διαιρείται με το πλήθος τους, π.χ. $\frac{18{,}5+19+21+20{,}5+17+16{,}5+22}{7}=\frac{134{,}5}{7}\approx19{,}2$.

\lyseis{3.4}{Τυποποιημένη μορφή μεγάλων αριθμών}
\sol{1} α) $7{,}2\cdot10^4$ \quad β) $1{,}25\cdot10^6$ \quad γ) $9{,}03\cdot10^8$ \quad δ) $6\cdot10^9$ \quad ε) $4{,}87\cdot10^4$ \quad στ) $1{,}5\cdot10^{13}$.
\sol{2} α) $270\,000$ \quad β) $5\,060$ \quad γ) $19\,000\,000$ \quad δ) $80\,000$ \quad ε) $3\,125\,000$ \quad στ) $4\,008\,000\,000$.
\sol{3} α) Ναι. β) Όχι: ο $52$ δεν είναι μικρότερος από $10$. γ) Όχι: ο $0{,}52$ είναι μικρότερος από $1$. δ) Όχι: ο $10$ δεν είναι μικρότερος από $10$. ε) Ναι.
\sol{4} α) $24\cdot10^5=2{,}4\cdot10^6$ \quad β) $0{,}35\cdot10^7=3{,}5\cdot10^6$ \quad γ) $150\cdot10^3=1{,}5\cdot10^5$ \quad δ) $0{,}008\cdot10^9=8\cdot10^6$.
\sol{5} α) $>$ (εκθέτες $5>4$) \quad β) $>$ (ίδιος εκθέτης, $3{,}5>2{,}9$) \quad γ) $>$ ($6>5{,}99$) \quad δ) $<$ (εκθέτες $3<5$).
\sol{6} $7{,}9\cdot10^5<9\cdot10^5<3{,}2\cdot10^6<4{,}5\cdot10^6<1{,}1\cdot10^7$.
\sol{7} α) $45\,000\,000$ \quad β) $25\,000\,000\,000$. Το γινόμενο είναι $36\,000\,000\,000\,000$ και η οθόνη δείχνει $3.6\mathrm{E}13$.
\sol{8} $3{,}844\cdot10^5$ χλμ. και $1{,}496\cdot10^8$ χλμ. Ο λόγος είναι περίπου $\frac{1{,}496\cdot10^8}{3{,}844\cdot10^5}\approx389$, δηλαδή περίπου $390$ φορές μακρύτερα.
\sol{9} $9{,}461\cdot10^{12}$ χλμ. Τα $4{,}3$ έτη φωτός είναι $4{,}3\cdot9{,}461\cdot10^{12}=40{,}68\cdot10^{12}=4{,}07\cdot10^{13}$ χλμ.
\sol{10} $\frac{1{,}496\cdot10^8}{3\cdot10^5}\approx499$ δευτερόλεπτα, δηλαδή $499:60\approx8{,}3$ λεπτά.
\sol{11} $8\cdot10^9$ και $1{,}04\cdot10^7$. Ο λόγος είναι $\frac{8\cdot10^9}{1{,}04\cdot10^7}\approx769$, περίπου $770$ φορές.
\sol{12} $5\,976\,000\,000\,000\,000\,000\,000\,000$ κιλά. Ο λόγος $\frac{1{,}983\cdot10^{30}}{5{,}976\cdot10^{24}}\approx3{,}3\cdot10^5$, δηλαδή περίπου $330\,000$ φορές.
\sol{13} $\frac{1{,}28\cdot10^{11}}{4\cdot10^6}=0{,}32\cdot10^5=3{,}2\cdot10^4=32\,000$ φωτογραφίες.
\sol{14} Επειδή ο πρώτος παράγοντας $10$ δεν είναι μικρότερος από $10$. Γράφεται $10\cdot10^5=10^6=1\cdot10^6$.
\sol{15} Ο $4{,}56\cdot10^{10}$ έχει $11$ ψηφία ($\nu+1$, όπου $\nu=10$ ο εκθέτης). Ο $1\cdot10^{\nu}$ έχει $\nu+1$ ψηφία: το $1$ και $\nu$ μηδενικά.
\sol{16} $(3\cdot10^4)\cdot(2\cdot10^5)=6\cdot10^9$. $(4\cdot10^3)\cdot(5\cdot10^6)=20\cdot10^9=2\cdot10^{10}$. Κανόνας: πολλαπλασιάζω τους $\alpha$, προσθέτω τους εκθέτες και, αν ο νέος $\alpha$ βγει ίσος ή μεγαλύτερος από $10$, ξαναγράφω σε τυποποιημένη μορφή.
\sol{17} Όχι. Ο $9{,}9\cdot10^3=9\,900$ και ο $1\cdot10^4=10\,000$, άρα $9{,}9\cdot10^3<1\cdot10^4$. Συγκρίνουμε πρώτα τους εκθέτες.
\sol{18} $\frac{10^{15}}{3{,}15\cdot10^7}\approx3{,}17\cdot10^7$ έτη, δηλαδή περίπου $32$ εκατομμύρια χρόνια.
\sol{19} Ελεύθερη απάντηση. Π.χ. μάζα του Ήλιου $1{,}989\cdot10^{30}$ κιλά, ακτίνα της Γης $6{,}371\cdot10^3$ χλμ. (οι τιμές διαφέρουν λίγο ανά πηγή).
\sol{20} Ελεύθερη απάντηση. Έλεγξε ότι ο $\alpha$ είναι τουλάχιστον $1$ και λιγότερο από $10$ και ότι ο εκθέτης είναι ίσος με το πλήθος των θέσεων της υποδιαστολής.

\lyseis{Α3.5}{Διαίρεση ρητών}
\sol{1} α) $-3$ \quad β) $-4$ \quad γ) $-9$ \quad δ) $+70$.
\sol{2} α) $\frac{2}{5}\cdot\frac{15}{4}=\frac{30}{20}=\frac{3}{2}$ \quad β) $\frac{1}{2}\cdot(-3)=-\frac{3}{2}$ \quad γ) $4\cdot2=8$ \quad δ) $\frac{4}{3}\cdot\frac{1}{8}=\frac{4}{24}=\frac{1}{6}$ \quad ε) $\frac{18}{5}\cdot\frac{5}{6}=3$ \quad στ) $\frac{12}{21}\cdot\frac{4}{7}=\frac{48}{147}=\frac{16}{49}$.
\sol{3} α) $\frac{7}{11}\cdot\frac{22}{3}=\frac{14}{3}$ \quad β) $\frac{27}{49}\cdot\frac{14}{9}=\frac{6}{7}$ \quad γ) $4\cdot\frac{9}{56}=\frac{9}{14}$ \quad δ) $\frac{11}{2}\cdot\frac{1}{121}=\frac{1}{22}$ \quad ε) $2\frac{1}{8}=\frac{17}{8}$, άρα $\frac{3}{4}\cdot\frac{8}{17}=\frac{6}{17}$.
\sol{4} α) $450:15=30$ \quad β) $0{,}58$ (γιατί $49\cdot0{,}58=28{,}42$) \quad γ) $357{,}6:15=23{,}84$ \quad δ) $10$ (πολλαπλασιάζω επί $1\,000$: $292\,620:29\,262=10$) \quad ε) $4{,}8:6=0{,}8$.
\sol{5} α) $4{,}57$ \quad β) $0{,}032$ \quad γ) $0{,}0005$ \quad δ) $64$ \quad ε) $35$ \quad στ) $12\,000$.
\sol{6} α) Σωστό. β) Λάθος. γ) Σωστό. δ) Λάθος: $(-12):(-4)=+3$. ε) Λάθος: $3:\frac{1}{3}=3\cdot3=9$. στ) Λάθος: π.χ. $3:\frac{1}{2}=6>3$.
\sol{7} α) $-\frac{3}{4}+\frac{5}{8}+\frac{1}{2}=-\frac{6}{8}+\frac{5}{8}+\frac{4}{8}=\frac{3}{8}$. β) $1-\frac{5}{4}+\frac{1}{8}=\frac{8}{8}-\frac{10}{8}+\frac{1}{8}=-\frac{1}{8}$.
\sol{8} α) $x=12:(-3)=-4$ \quad β) $x=4:\frac{2}{3}=4\cdot\frac{3}{2}=6$ \quad γ) $x=\left(-\frac{1}{2}\right):(-5)=\frac{1}{10}$.
\sol{9} $6:\frac{3}{4}=6\cdot\frac{4}{3}=8$ μπουκάλια.
\sol{10} $4{,}5:0{,}15=450:15=30$ κομμάτια.
\sol{11} $6:0{,}75=600:75=8$. Το κιλό κοστίζει $8$ €.
\sol{12} $(-24{,}60):6=-4{,}10$. Κάθε παιδί έχει μερίδιο $-4{,}10$ €, δηλαδή χρωστάει $4{,}10$ €.
\sol{13} $\frac{3}{5}:\frac{1}{10}=\frac{3}{5}\cdot10=6$. Τρέχει $6$ χιλιόμετρα σε μία ώρα.
\sol{14} Π.χ. $3:\frac{1}{2}=6>3$. Το πηλίκο είναι μεγαλύτερο από τον (θετικό) διαιρετέο όταν ο διαιρέτης είναι θετικός και μικρότερος από το $1$, γιατί τότε ο διαιρέτης «χωράει» περισσότερες φορές.
\sol{15} Αν υπήρχε $x=5:0$, θα έπρεπε $x\cdot0=5$. Όμως κάθε αριθμός επί $0$ δίνει $0$, άρα δεν υπάρχει τέτοιος αριθμός.
\sol{16} $3:\frac{1}{2}=6$ και $3:\frac{1}{3}=9$. Το $\frac{1}{3}$ χωράει περισσότερες φορές στο $3$ (τρία τρίτα σε κάθε ένα), άρα το $3:\frac{1}{3}$ είναι μεγαλύτερο.
\sol{17} Αν $\frac{\alpha}{\beta}>0$, οι $\alpha,\beta$ είναι ομόσημοι (και οι δύο θετικοί ή και οι δύο αρνητικοί). Αν $\frac{\alpha}{\beta}<0$, είναι ετερόσημοι.
\sol{18} Αναμενόμενο: $1{,}2:0{,}15=120:15=8$ κομμάτια. Ο μαθητής ελέγχει κόβοντας.
\sol{19} Ελεύθερη απάντηση. Π.χ. $\dfrac{\frac{1}{2}}{\frac{3}{4}}=\frac{1}{2}\cdot\frac{4}{3}=\frac{2}{3}$.

\lyseisE{Επανάληψη κεφαλαίου 3}
\solr{3}{1} $\frac{7}{25}=\frac{28}{100}=0{,}28$, $\frac{13}{8}=1{,}625$, $\frac{4}{5}=0{,}8$. Και $0{,}075=\frac{75}{1\,000}=\frac{3}{40}$.
\solr{3}{2} α) $0{,}609<0{,}610$ \quad β) $\frac{5}{8}=0{,}625>0{,}6$ \quad γ) $1{,}3=\frac{13}{10}$.
\solr{3}{3} Δέκατα $12{,}3$, εκατοστά $12{,}35$, χιλιοστά $12{,}346$.
\solr{3}{4} $24{,}3+5{,}78=30{,}08$ και $30{,}08-9{,}085=20{,}995$.
\solr{3}{5} α) $0{,}9$ \quad β) $450:18=25$ \quad γ) $0{,}008$ \quad δ) $730-50=680$.
\solr{3}{6} $2{,}4\cdot3{,}25=7{,}8$ και $0{,}9:0{,}3=3$, άρα $10{,}8$. Εκτίμηση: $2\cdot3+3=9$ περίπου.
\solr{3}{7} $3{,}6\cdot10^8$ και $7{,}5\cdot10^3$. Και $2{,}05\cdot10^6=2\,050\,000$.
\solr{3}{8} α) $4\,200$ m \quad β) $0{,}0085$ m$^2$ \quad γ) $2\,500$ cm$^3$ \quad δ) $4\,800$ s.
\solr{3}{9} $1{,}25\cdot0{,}85=1{,}0625$ και $0{,}5\cdot1{,}4=0{,}7$. Σύνολο $1{,}7625\approx1{,}76$ €.
\solr{3}{10} $12:0{,}75=1\,200:75=16$ μπουκάλια.
\solr{3}{11} $24{,}5^2=600{,}25$ m$^2$ $=0{,}60025$ στρέμματα $\approx0{,}6$ στρέμματα.
\solr{3}{12} $3\cdot10^5\cdot60=180\cdot10^5=1{,}8\cdot10^7$ χλμ.
\solr{3}{13} $7$ h $45$ min $+2$ h $50$ min $=9$ h $95$ min $=10$ h $35$ min. Φτάνει στις $10{:}35$.
\solr{3}{14} α) Σωστό. β) Λάθος: $0{,}07<0{,}10$. γ) Λάθος: $2{,}30>2{,}25$. δ) Σωστό. ε) Λάθος: υπάρχουν άπειροι, π.χ. $0{,}15$. στ) Σωστό.
\solr{3}{15} α) Σωστό. β) Λάθος: $5:0{,}1=50$. γ) Σωστό. δ) Λάθος: $(0{,}1)^2=0{,}01$. ε) Σωστό. στ) Σωστό.
\solr{3}{16} α) Λάθος: ο $52$ δεν είναι μικρότερος από $10$. β) Σωστό. γ) Σωστό. δ) Λάθος: $10\,000$. ε) Λάθος: $2$ h $30$ min. στ) Σωστό.


\lyseiskef{4}{Εξισώσεις και προβλήματα}
\lyseis{Α4.1}{Κανονικότητα, έννοια και αναπαραστάσεις}
\sol{1} α) Λ, Μ. Μονάδα επανάληψης: ΚΛΜ. β) $1,\ 1$. Μονάδα επανάληψης: $4,1,1$. γ) Β, Β. Μονάδα επανάληψης: ΑΒΒΓ.
\sol{2} α) $23,\ 28$ (κανόνας $+5$). β) $32,\ 25$ ($-7$). γ) $2,\ 5$ ($+3$). δ) $61,\ 72$ ($+11$).
\sol{3} α) $81,\ 243$ (πολλαπλασιάζω με $3$). β) $80,\ 160$ (με $2$). γ) $64,\ 32$ (κάθε όρος είναι το μισό του προηγούμενου, δηλαδή διαιρώ με $2$). δ) $32,\ -64$ (πολλαπλασιάζω με $-2$, γι' αυτό τα πρόσημα εναλλάσσονται).
\sol{4} α) Οι διαφορές είναι $1,2,3,4$, άρα επόμενες $5$ και $6$: $16$ και $22$. β) Οι διαφορές είναι $1,2,3,4$: επόμενα $17$ και $23$. γ) Οι διαφορές είναι $1,3,5,7$: επόμενες $9$ και $11$, άρα $25$ και $36$ (είναι τα τετράγωνα $0^2,1^2,2^2,\dots$).
\sol{5} α) Το $4$ο σχήμα είναι ένα «Γ» με $4$ τετράγωνα στην οριζόντια και $4$ στην κατακόρυφη πλευρά, με το γωνιακό κοινό: $7$ τετράγωνα (ακολουθία $1,3,5,7$). β) Το $4$ο σχήμα έχει $4$ γεμάτα τετράγωνα και ένα κενό: $5$ τετράγωνα (ακολουθία $2,3,4,5$).
\sol{6} α) Ο κανόνας είναι $+3$, άρα το τελευταίο πρέπει να είναι $13$. β) Ο κανόνας είναι $-10$, άρα $100,90,80,70,60$: το $75$ είναι λάθος. γ) Ο κανόνας είναι $\cdot2$, άρα το $49$ πρέπει να είναι $48$. δ) Το μοτίβο είναι ΚΜ ΚΜ ΚΜ\ldots Ο $5$ος όρος πρέπει να είναι Κ και όχι Μ.
\sol{7} α) Λάθος: υπάρχει πάντα κανόνας. β) Σωστό. γ) Λάθος: η μονάδα επανάληψης είναι η ΑΒ. δ) Σωστό ($+2$). ε) Σωστό ($+4$). στ) Σωστό. ζ) Λάθος: μπορούν να είναι και αρνητικοί αριθμοί ή σχήματα.
\sol{8} α) ΑΑΣΣΣ. β) Κάθε $5$ δευτερόλεπτα. γ) ΑΑΣΣΣ ΑΑΣΣΣ ΑΑΣΣΣ. δ) Τα δευτερόλεπτα $16$ και $17$ είναι τα δύο πρώτα της $4$ης μονάδας, δηλαδή ΑΑ: τα λαμπάκια είναι αναμμένα.
\sol{9} α) Οι χρόνοι είναι $4,5,6,7,8,9$: το $6$ο σκαλί παίρνει $9$ δευτερόλεπτα. β) Το $4$ είναι στη θέση $1$ και το $15$ είναι $11$ μονάδες μετά, άρα στη θέση $12$: το $12$ο σκαλί.
\sol{10} Οι κωδικοί διπλασιάζονται: $3,6,12,24,48,96,192$, άρα ο $7$ος είναι ο $192$. Συνεχίζουμε: $384,\ 768$. Ο πρώτος που ξεπερνά το $500$ είναι ο $768$ (ο $9$ος).
\sol{11} α) Τα πλακάκια είναι $1,3,6,10,15$: το $4$ο έχει $10$ και το $5$ο έχει $15$ (προσθέτω κάθε φορά $2,3,4,5$). β) Η περίμετρος αυξάνεται κατά $4$ κάθε φορά: $16$ για το $4$ο και $20$ για το $5$ο. Κανόνας: $+4$.
\sol{12} α) Κάθε όρος είναι κατά $3$ μεγαλύτερος από τον προηγούμενο ($+3$). β) Τα σημεία είναι τα $(1,2)$, $(2,5)$, $(3,8)$, $(4,11)$: ανεβαίνουν κατά $3$ όταν το $\nu$ αυξάνεται κατά $1$. γ) $5$ος όρος: $14$, $6$ος όρος: $17$.
\sol{13} Και οι δύο μπορεί να έχουν δίκιο: αν ο κανόνας είναι «διπλασιάζω», έχουμε $2,4,8,16$. Αν ο κανόνας είναι «οι διαφορές αυξάνονται κατά $2$» ($2,\,4,\,6,\dots$), έχουμε $2,4,8,14$. Τρεις όροι δεν αρκούν για να ξεχωρίσουμε τον κανόνα.
\sol{14} Μετά από $7$ ημέρες ξανά Δευτέρα. Η $1$η, $8$η, $15$η, \ldots είναι Δευτέρες. Η $99$η ($=1+14\cdot7$) είναι Δευτέρα, άρα η $100$ή είναι Τρίτη.
\sol{15} Μονάδα επανάληψης: ΑΒΑ ($3$ γράμματα). Στα $12$ γράμματα έχουμε $4$ μονάδες και κάθε μονάδα έχει $2$ Α, άρα $8$ Α.
\sol{16} Οι διαφορές είναι $1,2,3,\dots,9$. Οι όροι είναι $1,2,4,7,11,16,22,29,37,46$. Ο $10$ος όρος είναι $46$.
\sol{17} Ελεύθερη απάντηση. Π.χ. οι εποχές (κανόνας: ξαναρχίζουν κάθε $4$), η σειρά των φαναριών (πράσινο, πορτοκαλί, κόκκινο), οι ημέρες της εβδομάδας.
\sol{18} Ελεύθερη κατασκευή. Ο κανόνας πρέπει να ισχύει σε όλους τους όρους και να επιτρέπει να προβλέψεις τους επόμενους.

\lyseis{Α4.2}{Κανονικότητες της μορφής α·ν}
\sol{1} Οι επόμενοι όροι είναι $45$ και $54$. Γενικός όρος ο $9\nu$ (β): ο $1$ος όρος είναι $9\cdot1=9$ και κάθε όρος αυξάνεται κατά $9$. Για τη σειρά $4,8,12,16$: ο τύπος $4\nu$ (β) δίνει $4,8,12,16$. Οι άλλοι δεν κάνουν: $4+\nu$ δίνει $5$ για $\nu=1$, $4+4\nu$ δίνει $8$, και $8\nu-4$ δίνει $4$, $12$.
\sol{2} α) $8,\ 16,\ 24,\ 32,\ 40$. β) $11,\ 22,\ 33,\ 44,\ 55$. γ) $-3,\ -6,\ -9,\ -12,\ -15$. δ) $0{,}5,\ 1,\ 1{,}5,\ 2,\ 2{,}5$.
\sol{3} α) $30,\ 36$ και γενικός όρος $6\nu$. β) $-10,\ -12$ και $-2\nu$. γ) $1{,}25,\ 1{,}5$ και $0{,}25\nu$ (ο πρώτος όρος είναι $0{,}25$ και το βήμα $0{,}25$).
\sol{4} $9\cdot10=90$, $9\cdot25=225$, $9\cdot100=900$.
\sol{5} $6\nu=72$, $\nu=72:6=12$ (ο $12$ος). $6\nu=240$, $\nu=40$ (ο $40$ός). $6\nu=1\,500$, $\nu=250$ (ο $250$ός).
\sol{6} α) Ναι: $91:7=13$, είναι ο $13$ος όρος. β) Όχι: $150:7=21{,}4\ldots$ δεν είναι φυσικός. γ) Ναι: $700:7=100$, ο $100$ός όρος. δ) Όχι: $5:7<1$, δεν είναι φυσικός.
\sol{7} α) Λάθος: είναι $5\nu$ (το $5+\nu$ δίνει $6$ για $\nu=1$). β) Σωστό: $3\cdot4=12$. γ) Σωστό. δ) Σωστό. ε) Λάθος: το $\nu$ είναι θέση, φυσικός αριθμός. στ) Σωστό: $\alpha\cdot1=\alpha$. ζ) Σωστό: $2\cdot9=18$.
\sol{8} α) $3\cdot7=21$ οδοντογλυφίδες. β) $3\nu=120$, $\nu=40$: πλευρά $40$. γ) $3\nu=100$, $\nu=100:3=33{,}3\ldots$, που δεν είναι φυσικός: όχι.
\sol{9} Για $3,4,5,6$ πλευρές οι ακμές είναι $6,8,10,12$. Για $8$-γωνο $16$, για $9$-γωνο $18$ και για $20$-γωνο $40$. Γενικός όρος: $2\nu$ (το διπλάσιο των πλευρών της βάσης).
\sol{10} $\nu$ παγωτά κοστίζουν $1{,}5\nu$ €. Τα $12$ κοστίζουν $1{,}5\cdot12=18$ €. Με $27$ €: $1{,}5\nu=27$, $\nu=27:1{,}5=18$ παγωτά.
\sol{11} Μετά από $12$ δευτερόλεπτα: $\frac{3}{4}\cdot12=9$ μέτρα. Για τα $15$ μέτρα: $\frac{3}{4}\nu=15$, $\nu=15\cdot\frac{4}{3}=20$ δευτερόλεπτα.
\sol{12} Η θερμοκρασία μετά από $\nu$ ώρες είναι $-1{,}5\nu$. Μετά από $8$ ώρες: $-1{,}5\cdot8=-12$°. Για τους $-27$°: $-1{,}5\nu=-27$, $\nu=(-27):(-1{,}5)=18$: μετά από $18$ ώρες.
\sol{13} Για $\nu=1$ ο τύπος του Στέλιου δίνει $6+6\cdot1=12$, ενώ ο $1$ος όρος είναι $6$. Ο σωστός τύπος είναι $6\nu$.
\sol{14} $\alpha\cdot4=10$, άρα $\alpha=10:4=2{,}5$. Οι δύο πρώτοι όροι είναι $2{,}5$ και $5$. Ο $50$ός είναι $2{,}5\cdot50=125$.
\sol{15} Πρέπει $\alpha\nu=30$ με $\nu$ φυσικό, δηλαδή ο $\alpha$ να είναι διαιρέτης του $30$: $\alpha=1,\,2,\,3,\,5,\,6,\,10,\,15,\,30$.
\sol{16} Η σειρά της Ζωής έχει γενικό όρο $3\nu$, γιατί ο πρώτος όρος ($3$) ισούται με το βήμα. Η σειρά του Ηλία έχει βήμα $3$, αλλά ο πρώτος όρος είναι $4$: τα $4,\,7,\,10,\dots$ είναι $3\nu+1$, δεν έχουν τη μορφή $\alpha\nu$ ($3\cdot1=3\neq4$).
\sol{17} Ελεύθερη απάντηση. Π.χ. $12\nu$: αυγά σε $\nu$ δωδεκάδες, μήνες σε $\nu$ έτη. $0{,}5\nu$: λίτρα νερού για $\nu$ άτομα με μισό λίτρο το καθένα, ή ευρώ για $\nu$ καραμέλες των $0{,}5$ €.
\sol{18} Οι όροι είναι $2{,}5,\ 5,\ 7{,}5,\ 10,\ 12{,}5,\ 15,\ 17{,}5,\ 20$. Τα σημεία $(1;2{,}5)$, $(2;5)$, $\dots$ βρίσκονται σε μια ευθεία γραμμή που περνά από την αρχή των αξόνων.

\lyseis{4.3}{Παραδείγματα επίλυσης προβλημάτων}
\sol{1} $3x+3=111$, $3x=108$, $x=36$: οι αριθμοί $36,\,37,\,38$.
\sol{2} $x+2x=51$, $3x=51$, $x=17$: οι αριθμοί $17$ και $34$.
\sol{3} $x+(x+2)+(x+4)+(x+6)=100$, $4x+12=100$, $x=22$: οι $22,\,24,\,26,\,28$.
\sol{4} $\frac{3}{5}x=36$, $x=36:\frac{3}{5}=36\cdot\frac{5}{3}=60$.
\sol{5} $x-\frac{x}{3}-4=10$, $\frac{2x}{3}=14$, $x=21$ €. Έλεγχος: $21-7-4=10$.
\sol{6} $x+(x+5)=29$, $2x=24$, $x=12$ (η Πολυξένη) και ο αδελφός $17$.
\sol{7} $250\cdot8+100x=4\,800$, $100x=2\,800$, $x=28$ ημέρες.
\sol{8} $\frac{6\cdot5}{2}=15$ αγώνες. Με δέκα ομάδες $\frac{10\cdot9}{2}=45$.
\sol{9} $\frac{3}{8}x=1{,}2$, $x=1{,}2\cdot\frac{8}{3}=3{,}2$ λίτρα.
\sol{10} Η αμοιβή είναι $6x$ και το μισό της $270$: $3x=270$, $x=90$ € ανά ημέρα.
\sol{11} $3{,}5+0{,}8x=13{,}1$, $0{,}8x=9{,}6$, $x=12$ χλμ.
\sol{12} Πλησιάζουν με $8+12=20$ χλμ. την ώρα, άρα $x=5:20=0{,}25$ ώρες $=15$ λεπτά.
\sol{13} Μέτρησε κάθε χειραψία δύο φορές. Σωστά: $\frac{12\cdot11}{2}=66$.
\sol{14} $x+\frac{x}{2}=36$, $\frac{3x}{2}=36$, $x=24$. Έλεγχος: $24+12=36$.
\sol{15} Τα ζευγάρια $(0,10),\,(1,9),\dots,(10,0)$: $11$ ζευγάρια. Έχει πολλές λύσεις γιατί έχουμε δύο αγνώστους και μία μόνο ισότητα.
\sol{16} Κάθε ομάδα παίζει δύο φορές με καθεμία από τις άλλες $7$: ο συνολικός αριθμός είναι $8\cdot7=56$ αγώνες (χωρίς διαίρεση με το $2$, γιατί ο αγώνας εντός και ο εκτός είναι διαφορετικοί).
\sol{17} Δέκα ζευγάρια με άθροισμα $21$ το καθένα ($1+20$, $2+19$, \dots). Άθροισμα $10\cdot21=210$.
\sol{18} Ελεύθερη απάντηση. Π.χ. οι αριθμοί τριών γειτονικών σπιτιών έχουν άθροισμα $93$: $x+(x+1)+(x+2)=93$, $x=30$.
\sol{19} Ελεύθερη απάντηση. Ο Διόφαντος ονόμαζε τον άγνωστο «αριθμό».

\lyseisE{Επανάληψη κεφαλαίου 4}
\solr{4}{1} $5x-3$ και $2\nu+1$.
\solr{4}{2} $(7-3)x+(2+5)y=4x+7y$.
\solr{4}{3} α) $6{,}6$ \quad β) $6$ \quad γ) $12{,}5$ \quad δ) $10$ \quad ε) $4{,}5$.
\solr{4}{4} α) $3x=27$, $x=9$ \quad β) $5x=20$, $x=4$ \quad γ) $3x=33$, $x=11$ \quad δ) $\frac{x}{4}=4$, $x=16$.
\solr{4}{5} α) Ταυτότητα. β) Αδύνατη. γ) Μία λύση, $x=0$. δ) Αδύνατη.
\solr{4}{6} $2{,}5\cdot1{,}2\cdot x=9$, $3x=9$, $x=3$ μ.
\solr{4}{7} $5+0{,}03x=9{,}2$, $0{,}03x=4{,}2$, $x=140$ λεπτά.
\solr{4}{8} $\frac{3}{4}x=18$, $x=18\cdot\frac{4}{3}=24$ κομμάτια.
\solr{4}{9} $3x+3=87$, $x=28$: $28,\,29,\,30$.
\solr{4}{10} $x+(x+12)=100$, $x=44$ (η Αγγελική) και $56$ (ο Νίκος).
\solr{4}{11} $\frac{15\cdot14}{2}=105$ χειραψίες.
\solr{4}{12} Η αμοιβή είναι $5x$ και $\frac{3}{5}\cdot5x=3x=240$, $x=80$ € ανά ημέρα.
\solr{4}{13} α) Σωστό. β) Σωστό. γ) Λάθος: $3x+4x=7x$. δ) Σωστό. ε) Λάθος: $x=12-5=7$. στ) Λάθος: κάθε αριθμός την επαληθεύει.
\solr{4}{14} α) Σωστό. β) Λάθος: επαληθεύω και στο πρόβλημα. γ) Λάθος. δ) Λάθος: $x=4$. ε) Σωστό.


\lyseiskef{5}{Ποσοστά}
\lyseis{Α5.1}{Αλγεβρική παράσταση και επιμεριστική ιδιότητα}
\sol{1} α) αριθμητική \quad β) αλγεβρική \quad γ) τρεις όρους ($2x$, $6\alpha$ και $-2$) \quad δ) $12-3^2=12-9=3$.
\sol{2} α) $2x$ \quad β) $3x-2$ \quad γ) $x+1$ \quad δ) $2x$ \quad ε) $2x+1$ \quad στ) $\dfrac{x}{2}+3$ \quad ζ) $x^2+y^2$ \quad η) $(x+y)^2$.
\sol{3} α) Το πενταπλάσιο του $\alpha$. \quad β) Ο $\alpha$ αυξημένος κατά $7$. \quad γ) Το ένα τρίτο του $\alpha$ ελαττωμένο κατά $1$. \quad δ) Το γινόμενο των $\alpha$ και $\beta$ αυξημένο κατά $1$. \quad ε) Το διπλάσιο του αθροίσματος των $\alpha$ και $\beta$.
\sol{4} α) $2\cdot5+6=16$ \quad β) $3\cdot(-2)-4=-10$ \quad γ) $4^2-3\cdot4=16-12=4$ \quad δ) $10-\dfrac{6}{2}=7$ \quad ε) $3\cdot(-2)+3=-6+3=-3$.
\sol{5} α) $(3+5-6)x=2x$ \quad β) $(7-2+1)\alpha=6\alpha$ \quad γ) $6x+4y$ \quad δ) $(9-4-5)\omega=0$ \quad ε) $2\kappa+9$.
\sol{6} α) Όροι $2x^2$, $-x$, $8$. Συντελεστής του $x^2$ είναι το $2$ και του $x$ το $-1$. \quad β) Όροι $5\alpha$, $-3\beta$, $4$. Συντελεστής του $\alpha$ είναι το $5$ και του $\beta$ το $-3$.
\sol{7} α) $3x+12$ \quad β) $10-5x$ \quad γ) $6x+12$ \quad δ) $-2x+6$ \quad ε) $8\alpha+4\beta$ \quad στ) $-3+y$.
\sol{8} α) Λάθος: $2x+6$. \quad β) Λάθος: $x+x=2x$. \quad γ) Σωστό. \quad δ) Σωστό. \quad ε) Σωστό. \quad στ) Λάθος: οι όροι έχουν διαφορετικά γράμματα και δεν ενώνονται.
\sol{9} α) $8\cdot(50+1)=400+8=408$ \quad β) $7\cdot(100-1)=700-7=693$ \quad γ) $25\cdot(100+2)=2500+50=2550$ \quad δ) $9\cdot(50-3)=450-27=423$.
\sol{10} α) Τετράγωνο $4x$, τρίγωνο $3x$, ορθογώνιο $2x+2(x+3)=2x+2x+6=4x+6$. \quad β) Για $x=1$: $4,\ 3,\ 10$. Για $x=2$: $8,\ 6,\ 14$. Για $x=5$: $20,\ 15,\ 26$. Για $x=10$: $40,\ 30,\ 46$.
\sol{11} Λογαριασμός $7+0{,}05x$ ευρώ. Για $40$ λεπτά $7+2=9$ €, για $100$ λεπτά $7+5=12$ €, για $250$ λεπτά $7+12{,}5=19{,}5$ €.
\sol{12} Ως ένα ορθογώνιο $4(x+3)$, ως δύο ορθογώνια $4x+4\cdot3=4x+12$. Για $x=5$: $4\cdot8=32$ και $20+12=32$ m$^2$.
\sol{13} Κόστος $3x+2y$. Για $x=7$ και $y=4{,}5$: $21+9=30$ €. Με έκπτωση $2$ €: $3x+2y-2=28$ €.
\sol{14} α) $2x+6+4x-4=6x+2$ \quad β) $10\alpha-5\beta-3\alpha-3\beta=7\alpha-8\beta$ \quad γ) $3x+6-x+4=2x+10$ \quad δ) $7-2x+6=13-2x$.
\sol{15} α) $A=4x+6x-15-7x+6=3x-9$. \quad β) Για $x=4$: $3\cdot4-9=3$. Στην αρχική: $16+3\cdot3-(28-6)=16+9-22=3$. Συμφωνούν.
\sol{16} α) Αδελφός $x+3$, μητέρα $(x+3)+25=x+28$. \quad β) $x+(x+3)+(x+28)=3x+31$. \quad γ) $3\cdot11+31=64$ έτη (έλεγχος: $11+14+39=64$).
\sol{17} $2(x+7)-4-2x=2x+14-4-2x=10$. Τα $x$ φεύγουν και μένει πάντα $10$, όποιον αριθμό και να σκεφτεί κανείς.
\sol{18} Σωστά: $3(x+4)=3x+12$ (το $3$ πολλαπλασιάζει και το $4$) και $-2(x-5)=-2x+10$ (το $-2$ επί $-5$ δίνει $+10$). Η $4x+3y$ δεν συντομεύεται, μένει όπως είναι, γιατί οι όροι έχουν διαφορετικά γράμματα.
\sol{19} α) $3x+3y=3(x+y)=30$ \quad β) $2(x+y)-5=20-5=15$ \quad γ) $-x-y=-(x+y)=-10$.
\sol{20} α) Το μεγάλο ορθογώνιο έχει $3\cdot6=18$ τετράγωνα. Τα δύο μικρά έχουν $3\cdot4=12$ και $3\cdot2=6$ τετράγωνα και $12+6=18$. \quad β) $3\cdot(x+2)=3x+6$.
\sol{21} Για $x=0,1,2,3$: α) $6,\ 8,\ 10,\ 12$ \quad β) $6,\ 8,\ 10,\ 12$ \quad γ) $3,\ 5,\ 7,\ 9$ \quad δ) $6,\ 8,\ 10,\ 12$. Οι α, β, δ έχουν πάντα την ίδια τιμή: $2(x+3)=2x+6$ (επιμεριστική) και $x+3+x+3=2x+6$ (αναγωγή). Η γ διαφέρει, γιατί είναι η $2x+3$.
\sol{22} Ελεύθερη απάντηση. Π.χ. «Πέντε εισιτήρια των $x$ ευρώ και $12$ € για τη συμμετοχή» ή «Πέντε ίσα κουτιά βιβλίων των $x$ κιλών και ένα δέμα των $12$ κιλών». Για $x=3$ η τιμή είναι $15+12=27$.
\sol{23} Κόστος $0{,}8\alpha+1{,}1\beta$. Για $\alpha=5$ και $\beta=3$: $4+3{,}3=7{,}3$ €.
\sol{24} Ελεύθερη απάντηση. Π.χ. «Σκέψου έναν αριθμό, πρόσθεσε $3$, διπλασίασε, αφαίρεσε $6$: βρίσκεις το διπλάσιο του αριθμού». Απόδειξη: $2(x+3)-6=2x+6-6=2x$.

\lyseis{5.2}{Προβλήματα με ποσοστά}
\sol{1} α) $0{,}8\cdot45=36$ € \quad β) $9{,}6$ € \quad γ) $104$ €.
\sol{2} α) $1{,}15\cdot40=46$ € \quad β) $230$ € \quad γ) $20{,}7$ €.
\sol{3} α) $\frac{10}{50}=20\%$ αύξηση \quad β) $\frac{20}{80}=25\%$ μείωση \quad γ) $\frac{10}{25}=40\%$ αύξηση \quad δ) $\frac{30}{120}=25\%$ μείωση.
\sol{4} α) $1{,}2x=72$, $x=60$ \quad β) $0{,}6x=45$, $x=75$ \quad γ) $0{,}15x=12$, $x=80$.
\sol{5} Για $248$ €: $\frac{24}{124}\cdot248=48$ € ΦΠΑ και αξία $200$ €. Για $62$ €: ΦΠΑ $12$ € και αξία $50$ €.
\sol{6} α) $0{,}04\cdot1\,500=60$ € \quad β) $30$ €.
\sol{7} Με $30\%$ και $20\%$: $0{,}7\cdot0{,}8\cdot90=50{,}40$ €. Με $50\%$: $45$ €. Διαφορά $5{,}40$ €.
\sol{8} Μπουφάν: $0{,}85\cdot120=102$ €. Μπότα: $0{,}8\cdot85=68$ €. Μαζί $170$ €.
\sol{9} $250\cdot1{,}1=275$ και $275\cdot1{,}1=302{,}5$ €. Συνολική αύξηση $\frac{52{,}5}{250}=0{,}21=21\%$ (όχι $20\%$).
\sol{10} Έτος $1$: τόκος $100$ €, κεφάλαιο $5\,100$. Έτος $2$: τόκος $102$ €. Σύνολο τόκων $202$ € και ποσό $5\,202$ €.
\sol{11} $0{,}03x=60$, $x=2\,000$ €.
\sol{12} $12+0{,}05\cdot800=52$ €, με ΦΠΑ $1{,}24\cdot52=64{,}48$ €.
\sol{13} Μετρητά $200$ €, υπόλοιπο $600$ €. Τόκος $3\cdot0{,}01\cdot600=18$ €, δηλαδή $\frac{18}{800}=2{,}25\%$ της αρχικής τιμής.
\sol{14} Ύπνος $33{,}3\%$, σχολείο $25\%$, διάβασμα $12{,}5\%$, αθλητισμός $8{,}3\%$, υπόλοιπα $20{,}8\%$. Μήκη ράβδων: $3{,}3$, $2{,}5$, $1{,}25$, $0{,}83$ και $2{,}08$ εκ. (το άθροισμα βγαίνει $10$ εκ. με στρογγυλοποίηση).
\sol{15} Η νέα τιμή είναι $0{,}8$ της παλιάς, και $1:0{,}8=1{,}25$. Πρέπει να ανέβει κατά $25\%$.
\sol{16} $100\to70\to56$. Η τιμή έχει μειωθεί κατά $44$ €, δηλαδή $44\%$.
\sol{17} Όχι: $100\to50\to75$, δηλαδή $25\%$ χαμηλότερα.
\sol{18} Πληρώνεις $2$ από τα $3$, άρα η έκπτωση είναι $\frac{1}{3}\approx33{,}3\%$.
\sol{19} $100-28-35-12=25\%$.
\sol{20} Με αξία $100$ € ο ΦΠΑ είναι $24$ € και η τιμή $124$ €. Το $24$ είναι τα $\frac{24}{124}\approx19{,}4\%$ της τελικής τιμής.
\sol{21} Ελεύθερη απάντηση.
\sol{22} Ελεύθερη απάντηση. Αθροίζουν τα ποσοστά σε $100\%$.

\lyseisE{Επανάληψη κεφαλαίου 5}
\solr{5}{1} $60\%$, $87{,}5\%$, $45\%$.
\solr{5}{2} $\frac{35}{100}=\frac{7}{20}$ και $\frac{12{,}5}{100}=\frac{1}{8}$.
\solr{5}{3} $36$ και $21$.
\solr{5}{4} $\frac{18}{72}=\frac{1}{4}=25\%$.
\solr{5}{5} $0{,}7\cdot140=98$ €.
\solr{5}{6} $1{,}12\cdot250=280$ €.
\solr{5}{7} $\frac{20}{80}=25\%$ αύξηση.
\solr{5}{8} $0{,}8x=64$, $x=80$ €.
\solr{5}{9} $1{,}24x=186$, $x=150$ €.
\solr{5}{10} Ένα έτος $60$ €, εννέα μήνες $0{,}75\cdot60=45$ €.
\solr{5}{11} Ο τέταρτος πήρε $10\%$. Ψήφοι: $180$, $120$, $60$ και $40$.
\solr{5}{12} $0{,}9\cdot0{,}9\cdot200=162$ €. Συνολική έκπτωση $38$ €, δηλαδή $19\%$ (όχι $20\%$).
\solr{5}{13} α) Σωστό. β) Σωστό. γ) Σωστό. δ) Λάθος: $10\permil=1\%$. ε) Σωστό ($10$ και $10$). στ) Λάθος: $0{,}5\%=0{,}005$.
\solr{5}{14} α) Λάθος: $44\%$. β) Λάθος: $0{,}9375$ της αρχικής. γ) Λάθος: τα $\frac{24}{124}$. δ) Σωστό. ε) Σωστό.


\lyseiskef{Α6}{Αλγεβρικές σχέσεις}
\lyseis{Α6.1}{Ιδιότητες της ισότητας και η έννοια της εξίσωσης}
\sol{1} α) Σωστό. \quad β) Λάθος: $x=7+3=10$. \quad γ) Σωστό. \quad δ) Σωστό (διαιρώ με $3$). \quad ε) Σωστό (πολλαπλασιάζω με $3$). \quad στ) Λάθος: $2\cdot6=12\neq6$, η λύση είναι το $3$.
\sol{2} $x+200=500$. Αφαιρώ $200$ και από τα δύο μέλη: $x=300$ g. Έλεγχος: $300+200=500$.
\sol{3} α) $\alpha=7-2=5$ \quad β) $\kappa=10+3=13$ \quad γ) $\lambda=3$ \quad δ) $\rho=45:9=5$.
\sol{4} α) Ναι: $6+3=9$. \quad β) Ναι: $9-2=7$. \quad γ) Όχι: $2\cdot20=40\neq10$ (η λύση είναι το $5$). \quad δ) Όχι: $12-22=-10\neq10$ (η λύση είναι το $2$). \quad ε) Ναι: $12:4=3$.
\sol{5} α) $33$ \quad β) $8$ \quad γ) $53$ \quad δ) $48$ \quad ε) $3$ \quad στ) $200$ \quad ζ) $2$.
\sol{6} α) $x=-5$, έλεγχος $-5\cdot(-5)=25$ \quad β) $x=-12+8=-4$, έλεγχος $-4-8=-12$ \quad γ) $x=3$, έλεγχος $-6\cdot3=-18$ \quad δ) $-x=11+8=19$, $x=-19$, έλεγχος $-8+19=11$.
\sol{7} α) $2{,}25$ \quad β) $x=\frac{9}{4}\cdot2=\frac{9}{2}$ \quad γ) $x=\frac{5}{6}+\frac{1}{3}=\frac{7}{6}$ \quad δ) $x=10\cdot\frac{3}{5}=6$.
\sol{8} α) $x+17=52$, $x=35$ \quad β) $3x=84$, $x=28$ \quad γ) $50-x=18$, $x=32$ \quad δ) $x:2=7$, $x=14$.
\sol{9} $2x=1\,200$, άρα $x=600$ g. Έλεγχος: $2\cdot600=1\,200$.
\sol{10} α) $4x=54$, $x=13{,}5$ € \quad β) $36-x=11{,}5$, $x=24{,}5$ €.
\sol{11} α) $4x=98$, $x=24{,}5$ cm. \quad β) $3x=4{,}5$, $x=1{,}5$ m.
\sol{12} α) Ταυτότητα (αντιμεταθετική ιδιότητα). \quad β) Αδύνατη. \quad γ) Αδύνατη: το αριστερό μέλος είναι πάντα κατά $4$ μικρότερο. \quad δ) Ταυτότητα: $5x-5x=0$ και $0=0$.
\sol{13} Κάθε αριθμός επί $0$ κάνει $0$, άρα η $0\cdot x=0$ ισχύει για κάθε $x$. Κανένας αριθμός επί $0$ δεν κάνει $5$, άρα η $0\cdot x=5$ δεν έχει λύση.
\sol{14} Στην πρώτη πρόσθεσε αντί να αφαιρέσει: από την $\alpha-x=\beta$ προκύπτει $x=\alpha-\beta$, δηλαδή $x=12-5=7$ (έλεγχος $12-7=5$). Στη δεύτερη πρέπει $x=24:6=4$ (έλεγχος $24:4=6$).
\sol{15} Ελεύθερη απάντηση. Π.χ. $x+4=10$ (λύση $6$), $x+1=x+2$ (αδύνατη), $x+3=3+x$ (ταυτότητα).
\sol{16} α) Ισοδύναμες: και οι δύο έχουν μόνη λύση το $4$. \quad β) Όχι: η πρώτη έχει λύση το $3$ και η δεύτερη το $1$. \quad γ) Ισοδύναμες: και οι δύο είναι αδύνατες, δηλαδή έχουν τις ίδιες (καμία) λύσεις.
\sol{17} $3x=2{,}4$, άρα $x=0{,}8$.
\sol{18} Ελεύθερη απάντηση. Π.χ. «Η Ζωή είχε $x$ ευρώ, πήρε άλλα $12$ και έχει τώρα $30$»: $x=30-12=18$.

\lyseis{6.2}{Λόγος και αναλογία}
\sol{1} α) $\frac{3}{4}$ \quad β) $\frac{30}{120}=\frac{1}{4}$ \quad γ) $\frac{250}{50}=5$ \quad δ) $\frac{0{,}6}{1{,}5}=\frac{2}{5}$ \quad ε) $\frac{3}{4}:\frac{1}{2}=\frac{3}{2}$.
\sol{2} α) Ναι. β) Όχι ($4\cdot9=36\neq35=7\cdot5$). γ) Ναι ($12\cdot30=360=18\cdot20$). δ) Ναι ($0{,}5\cdot12=6=2\cdot3$). ε) Όχι ($7\cdot25=175\neq168=8\cdot21$).
\sol{3} α) $15x=60$, $x=4$ \quad β) $12x=72$, $x=6$ \quad γ) $5x=60$, $x=12$ \quad δ) $21x=84$, $x=4$ \quad ε) $5x=15$, $x=3$.
\sol{4} $4\cdot15=60=6\cdot10$. Οι δύο λόγοι είναι ίσοι ($\frac{2}{3}$ ο καθένας).
\sol{5} $\frac{18}{27}=\frac{2}{3}$.
\sol{6} α) $4\cdot50\,000$ cm $=200\,000$ cm $=2$ km. β) $15$ km $=1\,500\,000$ cm, διά $200\,000$ δίνει $7{,}5$ cm. γ) $8\cdot1\,000$ cm $=80$ μ.
\sol{7} α) Σωστό. β) Λάθος: $\frac{200}{50}=4$. γ) Σωστό. δ) Σωστό. ε) Λάθος: έχουν ίσους λόγους περιμέτρων, όχι ίσες περιμέτρους.
\sol{8} $3{,}6\cdot5=18$ km και $5{,}2\cdot5=26$ km.
\sol{9} $42$ km $=4\,200\,000$ cm, διά $200\,000$ είναι $21$ cm.
\sol{10} Λόγος μεγέθυνσης $\frac{40}{10}=4$, άρα νέο ύψος $4\cdot15=60$ cm. Περίμετροι: $2(10+15)=50$ και $2(40+60)=200$, λόγος $\frac{200}{50}=4$.
\sol{11} Το $3\times1{,}2$ μ.: $300:25=12=120:10$.
\sol{12} Λόγος $1:4$, άρα σιρόπι $750:4=187{,}5$ ml.
\sol{13} $36\cdot25=900$ cm $=9$ μ.
\sol{14} $\frac{15}{x}=\frac{3}{2}$, $3x=30$, $x=10$ λίτρα.
\sol{15} Το πλάτος γίνεται $2$ φορές μεγαλύτερο ($\frac{2}{1}$) αλλά το ύψος μένει ίδιο ($\frac{1}{1}$): οι λόγοι δεν είναι ίσοι.
\sol{16} $\frac{4}{10}=0{,}4$ και $\frac{6}{15}=0{,}4$: όμοια. Για τα $4\times6$ και $8\times10$: $\frac{4}{8}=0{,}5$ και $\frac{6}{10}=0{,}6$: όχι όμοια.
\sol{17} Πολλαπλασιάζω και τα δύο μέλη με $\beta\delta$: $\frac{\alpha}{\beta}\cdot\beta\delta=\alpha\delta$ και $\frac{\gamma}{\delta}\cdot\beta\delta=\beta\gamma$, άρα $\alpha\delta=\beta\gamma$.
\sol{18} $\frac{12}{16}=\frac{3}{4}$ και $\frac{6}{8}=\frac{3}{4}$.
\sol{19} $7{,}5\cdot10$ km $=75$ km. Ο χάρτης $1:200\,000$ δείχνει περισσότερες λεπτομέρειες, γιατί η ίδια απόσταση σχεδιάζεται μεγαλύτερη.
\sol{20} Ελεύθερη απάντηση.
\sol{21} Ελεύθερη απάντηση. Οι περισσότερες οθόνες έχουν λόγο μήκους προς πλάτος $\frac{16}{9}$ ή παραπλήσιο.

\lyseisE{Επανάληψη κεφαλαίου 6}
\solr{6}{1} Πάνω σε άξονα είναι το $(4,\,0)$ (στον $Ox$) και το $(0,\,2)$ (στον $Oy$).
\solr{6}{2} Τέταρτη κορυφή $(1,\,5)$. Πλευρά $4$, εμβαδόν $16$.
\solr{6}{3} $40$ λεπτά προς $120$ λεπτά: $\frac{1}{3}$.
\solr{6}{4} α) $9x=72$, $x=8$ \quad β) $20x=60$, $x=3$.
\solr{6}{5} $6\cdot250\,000=1\,500\,000$ cm $=15$ km.
\solr{6}{6} Ο Α: ναι, $\alpha=4$. Ο Β: όχι ($\frac{3}{1}\neq\frac{5}{2}$).
\solr{6}{7} Για $x=8$: $y=6$. Για $y=9$: $x=9:0{,}75=12$.
\solr{6}{8} $\alpha=\frac{9}{6}=1{,}5$. Για $x=10$: $y=15$.
\solr{6}{9} Ένα κιλό $4{,}5$ €, τα $10$ κιλά $45$ €.
\solr{6}{10} Δουλειά $4\cdot15=60$ εργάτες-ημέρες. Με $6$ εργάτες $10$ ημέρες.
\solr{6}{11} $4\cdot x=36$, $x=9$ μ.
\solr{6}{12} Σύνολο μερών $5$, κάθε μέρος $30$ €: $60$ € και $90$ €.
\solr{6}{13} α) Σωστό. β) Λάθος: για $x=0$ δίνει $y=3$. γ) Σωστό. δ) Σωστό. ε) Λάθος: $\frac{200}{50}=4$.
\solr{6}{14} α) Σωστό. β) Λάθος: μισός χρόνος. γ) Σωστό. δ) Σωστό. ε) Σωστό.


\lyseiskef{7}{Θετικοί και αρνητικοί αριθμοί}
\lyseis{7.1}{Θετικοί και αρνητικοί αριθμοί}
\sol{1} α) $-5$ \quad β) $+12$ \quad γ) $-8$ \quad δ) $+15$ \quad ε) $-2$ \quad στ) $-30$.
\sol{2} Θετικοί: $+3{,}5$, $8$, $+0{,}01$. Αρνητικοί: $-7$, $-\frac{1}{2}$. Το $0$ δεν είναι ούτε θετικός ούτε αρνητικός.
\sol{3} α) Ομόσημοι. β) Ετερόσημοι. γ) Ομόσημοι. δ) Ετερόσημοι.
\sol{4} Ακέραιοι: $-4$, $0$, $\frac{9}{3}=3$, $10$.
\sol{5} Ο $-4$ είναι $4$ cm αριστερά του $0$, ο $-1{,}5$ είναι $1{,}5$ cm αριστερά, ο $2$ $2$ cm δεξιά και ο $3{,}5$ $3{,}5$ cm δεξιά.
\sol{6} $A(-4)$, $B(-1{,}5)$, $\Gamma(1)$, $\Delta(3{,}5)$.
\sol{7} α) Λάθος. β) Σωστό. γ) Σωστό. δ) Λάθος: ομόσημοι. ε) Σωστό. στ) Σωστό.
\sol{8} Αθήνα $+7$, Θεσσαλονίκη $-3$, Ιωάννινα $-5$, Χανιά $+11$, Φλώρινα $-10$. Στον κατακόρυφο άξονα, από πάνω προς τα κάτω: $+11$, $+7$, $0$, $-3$, $-5$, $-10$.
\sol{9} Από το $3$ στο $-2$ κατεβαίνει $5$ ορόφους ($3,\,2,\,1,\,0,\,-1,\,-2$).
\sol{10} Το ψάρι ($-45$ μ.) είναι πιο κοντά στην επιφάνεια ($0$). Το αεροπλάνο είναι πάνω από την επιφάνεια, στα $300$ μ.
\sol{11} Μυρτώ $+20$, Στέλιος $-7$, Φανή $+5$, Βασίλης $-12$.
\sol{12} Από το $-3$ ανεβαίνω $5$ μονάδες δεξιά: $-2,\,-1,\,0,\,1,\,2$. Η θερμοκρασία γίνεται $+2$°.
\sol{13} Παράδειγμα: ένας χειμωνιάτικος θερμοκρασία $-4$°C, ή ο λογαριασμός τράπεζας με $-50$ €, ή το βάθος της θάλασσας.
\sol{14} Το $0$ είναι το σημείο που χωρίζει τις δύο πλευρές: δεν είναι ούτε στη θετική ούτε στην αρνητική. Δεν έχει κατεύθυνση.
\sol{15} Το μηδέν είναι σύμβαση: ανάλογα με την κλίμακα, το $0$°C αντιστοιχεί στα $32$°F. Δεν σημαίνει «καθόλου θερμοκρασία».
\sol{16} Με τα άκρα: $-5,\dots,5$, δηλαδή $11$. Χωρίς τα άκρα $9$.
\sol{17} Ελεύθερη απάντηση. Π.χ. θερμοκρασία, υψόμετρο και βάθος, όροφοι, λογαριασμός τράπεζας, έτη π.Χ.
\sol{18} Ελεύθερη απάντηση.

\lyseis{7.2}{Απόλυτη τιμή, αντίθετοι, σύγκριση}
\sol{1} α) $7$ \quad β) $7$ \quad γ) $0$ \quad δ) $2{,}5$ \quad ε) $\frac{3}{4}$ \quad στ) $0{,}01$.
\sol{2} $-5$, $+8$, $0$, $-2{,}5$, $+\frac{1}{3}$, $-0{,}6$.
\sol{3} α) $4$ \quad β) $-3$ \quad γ) $-2$ \quad δ) $-6$.
\sol{4} α) $5$ ή $-5$ \quad β) $0$ \quad γ) $\frac{1}{2}$ ή $-\frac{1}{2}$ \quad δ) αδύνατη.
\sol{5} α) $-3>-8$ \quad β) $2>-9$ \quad γ) $-0{,}5<-0{,}05$ \quad δ) $-\frac{7}{2}=-3{,}5<-3$ \quad ε) $0>-0{,}001$ \quad στ) $-4<4$.
\sol{6} $-7<-3<-0{,}5<0<2{,}5<4<5$.
\sol{7} Μεγαλύτερος ο $-1{,}02$ και μικρότερος ο $-1{,}22$.
\sol{8} α) Σωστό. β) Λάθος: $|-5|=5$. γ) Σωστό. δ) Λάθος: $-9<-2$. ε) Σωστό. στ) Σωστό.
\sol{9} $-10<-5<-3<-1<7$: Φλώρινα, Μέτσοβο, Ιωάννινα, Θεσσαλονίκη, Αθήνα.
\sol{10} Το βιβλιοπωλείο ($|{-2}|=2$ μονάδες $=200$ μ.) είναι πιο κοντά από τον φούρνο ($300$ μ.). Η απόσταση μεταξύ τους είναι από το $-2$ στο $+3$: $5$ μονάδες $=500$ μ.
\sol{11} $-11\,034<-430<0<2\,917$.
\sol{12} Χρωστάει περισσότερο η Ζωή ($|{-15}|=15$ €). Το μεγαλύτερο ποσό έχει η Ρήνη ($+12$).
\sol{13} Το $\Lambda$ έχει τετμημένη $+7$. Η απόσταση $K\Lambda$ είναι $14$ μονάδες ($7$ αριστερά και $7$ δεξιά από το $0$).
\sol{14} Στον άξονα το $-10$ είναι πολύ πιο αριστερά από το $-3$, άρα μικρότερο. Στο θερμόμετρο $10$ βαθμούς κάτω από το μηδέν κάνει περισσότερο κρύο από $3$ βαθμούς κάτω. Το $10>3$ ισχύει για τις απόλυτες τιμές, όχι για τους αρνητικούς.
\sol{15} Το $\alpha$ και το $-\alpha$ έχουν την ίδια απόσταση από το $0$.
\sol{16} Το $|\alpha|=\alpha$ ισχύει για $\alpha\ge0$ και το $|\alpha|=-\alpha$ για $\alpha\le0$.
\sol{17} Για παράδειγμα $x=3$, $y=-3$. Ακέραιοι με απόλυτη τιμή μικρότερη από $2$: $-1,\,0,\,1$, δηλαδή $3$.
\sol{18} Για παράδειγμα $-1{,}5$ και $-0{,}25$, $-0{,}22$. Υπάρχουν άπειροι.
\sol{19} Ελεύθερη απάντηση.

\lyseis{7.3}{Πρόσθεση ρητών}
\sol{1} α) $3$ \quad β) $-4$ \quad γ) $-7$ \quad δ) $0$ \quad ε) $0$.
\sol{2} α) $-13$ \quad β) $5$ \quad γ) $-5$ \quad δ) $12$ \quad ε) $-12$ \quad στ) $-1$.
\sol{3} α) $0$ \quad β) $-2{,}5$ \quad γ) $-3{,}2$ \quad δ) $-\frac{3}{4}$ \quad ε) $\frac{1}{6}$ \quad στ) $-\frac{5}{6}+\frac{2}{6}=-\frac{3}{6}=-\frac{1}{2}$.
\sol{4} α) $7$ \quad β) $-3{,}5$ \quad γ) $4$ \quad δ) $\frac{1}{3}$.
\sol{5} α) $21-15=6$ \quad β) $(4+1)+(-1{,}5-2{,}5)=5-4=1$ \quad γ) $10+(-12)=-2$.
\sol{6} α) $0$ \quad β) $\big(37+(-37)\big)+\big((-18)+8\big)=0+(-10)=-10$.
\sol{7} α) Σωστό. β) Λάθος: το πρόσημο είναι του μεγαλύτερου κατ' απόλυτη τιμή. γ) Σωστό. δ) Λάθος: $-7$. ε) Σωστό. στ) Σωστό. ζ) Σωστό ($=3$).
\sol{8} Μεσημέρι $-6+9=3$°. Βράδυ $3+(-5)=-2$°.
\sol{9} Υπόλοιπα: $30,\ 18,\ 26,\ -9,\ 11$. Τελικό υπόλοιπο $11$ €.
\sol{10} $-2+5=3$ και $3+(-4)=-1$. Φτάνει στον όροφο $-1$.
\sol{11} $-18+7=-11$ και $-11-4=-15$. Βρίσκεται στα $-15$ μ.
\sol{12} $8+(-5)+(-7)+12=8$ πόντους.
\sol{13} $-45+80=35$ και $35-20=15$ €.
\sol{14} $12+(-7)+(-9)+4+6=6$.
\sol{15} Όταν είναι αντίθετοι.
\sol{16} Το πρόσημο του προσθετέου με τη μεγαλύτερη απόλυτη τιμή (και $0$ αν είναι αντίθετοι).
\sol{17} Για παράδειγμα $(-2)+(-3)$ και $(-8)+(+3)$.
\sol{18} Όχι: $(-3)+(-2)=-5<-3$.
\sol{19} $(-3)+(-4)=-7$ και $-(3+4)=-7$. Ίσα.
\sol{20} Ελεύθερη απάντηση.
\sol{21} Ελεύθερη απάντηση.

\lyseis{7.4}{Αφαίρεση ρητών}
\sol{1} α) $2$ \quad β) $-2$ \quad γ) $-3$ \quad δ) $3$ \quad ε) $-6$ \quad στ) $6$.
\sol{2} α) $7+(-9)=-2$ \quad β) $-4+6=2$ \quad γ) $3{,}5+1{,}5=5$ \quad δ) $-\frac{1}{2}-\frac{1}{3}=-\frac{5}{6}$ \quad ε) $-2{,}4-0{,}6=-3$ \quad στ) $8+8=16$.
\sol{3} α) $x=3-7=-4$ \quad β) $x=-2+5=3$ \quad γ) $x=4-9=-5$ \quad δ) $x=-6-(-6)=0$ \quad ε) $x=-3-2=-5$.
\sol{4} α) $10-3+8-2=13$ \quad β) $5-3+7-2+9=16$ \quad γ) $4-5-7+10=2$ \quad δ) $20+5+3=28$.
\sol{5} α) $-12$ \quad β) $-2{,}5+7-4{,}5-1=-1$ \quad γ) $\frac{1}{2}-\frac{3}{4}+\frac{1}{4}-1=-1$.
\sol{6} α) Λάθος. β) Σωστό. γ) Σωστό. δ) Σωστό. ε) Λάθος: $4+2=6$. στ) Σωστό ($3-5=-2$).
\sol{7} α) $5-(-3)=8$ \quad β) $-2-(-11)=-2+11=9$.
\sol{8} $2\,917-(-430)=2\,917+430=3\,347$ μ.
\sol{9} $250-(-206)=456$ χρόνια (περίπου, χωρίς έτος $0$: $455$).
\sol{10} $8-(-15)=23$ € (ή $-15-8=-23$ αν αφαιρέσουμε ανάποδα).
\sol{11} $(-9)-(-22)=13$ μ.
\sol{12} $(-2)-3=-5$: κατεβαίνει $5$ ορόφους.
\sol{13} $-7-4=-11$: πέφτει $11$ βαθμούς.
\sol{14} Το σωστό είναι $5-(-3)=5+3=8$. Στον άξονα από το $-3$ ως το $5$ είναι $8$ μονάδες. Ο Στέλιος αφαίρεσε το $3$ αντί για το $-3$.
\sol{15} Όχι: $3-(-2)=5>3$.
\sol{16} Θετικό αν $\alpha>\beta$, αρνητικό αν $\alpha<\beta$, μηδέν αν ίσα.
\sol{17} $10-(3+4)=3$ και $10-3-4=3$.
\sol{18} $x-(-5)=2$, άρα $x=2+(-5)=-3$.
\sol{19} Ελεύθερη απάντηση.
\sol{20} Ελεύθερη απάντηση.

\lyseis{7.5}{Πολλαπλασιασμός ρητών}
\sol{1} $-12,\ -8,\ -4,\ 0,\ 4,\ 8,\ 12$.
\sol{2} α) $-24$ \quad β) $-21$ \quad γ) $45$ \quad δ) $0$ \quad ε) $40$ \quad στ) $1$.
\sol{3} α) $-6$ \quad β) $3{,}6$ \quad γ) $-1$ \quad δ) $\frac{2}{3}$ \quad ε) $-6$ \quad στ) $\frac{1}{3}\cdot(-3)=-1$.
\sol{4} α) Αρνητικό ($3$ αρνητικοί). β) Θετικό ($4$). γ) Αρνητικό ($3$). δ) Είναι $0$.
\sol{5} α) $-8$ \quad β) $8$ \quad γ) $0$ \quad δ) $7$.
\sol{6} α) $-40$ \quad β) $25$ \quad γ) $-10+18=8$ \quad δ) $-594$ \quad ε) $-4-3=-7$.
\sol{7} α) $-12+10+6=4$ \quad β) $(-2)(-3)+3(-3)=6-9=-3$ \quad γ) $10-30=-20$.
\sol{8} α) Λάθος: $+12$. β) Σωστό. γ) Σωστό. δ) Σωστό. ε) Λάθος: αρνητικό. στ) Σωστό.
\sol{9} Σε $3$ ώρες $3\cdot(-2)=-6$ βαθμούς. Πριν $3$ ώρες $(-2)\cdot(-3)=+6$, δηλαδή τότε είχε $6$ βαθμούς περισσότερο.
\sol{10} $5\cdot(-4)=-20$ €.
\sol{11} Σε $8$ λεπτά: $8\cdot(-3)=-24$ μ. Πριν $5$ λεπτά: $(-3)\cdot(-5)=+15$, δηλαδή $15$ μ. πιο ψηλά.
\sol{12} $4\cdot(-2)=-8$ όροφοι (κατεβαίνει $8$).
\sol{13} $9\cdot(-4)=-36$ €.
\sol{14} $6\cdot(-0{,}25)=-1{,}5$ €.
\sol{15} Αν η θερμοκρασία πέφτει $2$ βαθμούς την ώρα, πηγαίνοντας $3$ ώρες πίσω έχουμε $6$ περισσότερους. Στο μοτίβο: $0\cdot(-3)=0$, $(-1)\cdot(-3)=3$, $(-2)\cdot(-3)=6$.
\sol{16} Γιατί το $\alpha\cdot\alpha$ έχει ομόσημους παράγοντες (ίδιο πρόσημο), άρα θετικό (ή $0$).
\sol{17} Άρτιο πλήθος αρνητικών, π.χ. $2$, $4$, $6$, $8$, $10$ αρνητικοί.
\sol{18} Ένας τουλάχιστον από τους δύο είναι $0$.
\sol{19} Με $100$ παράγοντες $+1$ και με $101$ παράγοντες $-1$.
\sol{20} Ελεύθερη απάντηση: $3\cdot(-5)=-15$, $2\cdot(-5)=-10$, \dots, $(-3)\cdot(-5)=15$.

\lyseis{7.6}{Διαίρεση ρητών}
\sol{1} α) $-4$ \quad β) $8$ \quad γ) $-3$ \quad δ) $0$ \quad ε) $3$ \quad στ) $4$.
\sol{2} α) $x=12:(-3)=-4$ \quad β) $x=-4$ \quad γ) $x=6\cdot(-2)=-12$ \quad δ) $x=1$ \quad ε) $x=-15$ \quad στ) $x=12:(-4)=-3$.
\sol{3} α) $-\frac{3}{2}$ \quad β) $\frac{2}{3}\cdot\frac{9}{4}=\frac{3}{2}$ \quad γ) $5\cdot\left(-\frac{5}{2}\right)=-\frac{25}{2}$ \quad δ) $-3$ \quad ε) $-0{,}4$ \quad στ) $1$.
\sol{4} $-\frac{1}{4}$, $-\frac{3}{2}$, $\frac{1}{5}$, $-2$, $3$.
\sol{5} α) $-4$ \quad β) $\frac{2}{3}\cdot\left(-\frac{9}{4}\right)=-\frac{3}{2}$ \quad γ) $(-5)\cdot(-2)=10$.
\sol{6} α) $-4+10=6$ \quad β) $-5+6=1$ \quad γ) $3\cdot(-4)=-12$.
\sol{7} α) Αρνητικό. β) Θετικό. γ) Αρνητικό. δ) Θετικό.
\sol{8} α) Λάθος: $+3$. β) Σωστό. γ) Λάθος: δεν ορίζεται. δ) Σωστό. ε) Σωστό. στ) Σωστό.
\sol{9} $(-24):6=-4$ €. Και $(-24):(-3)=8$ παιδιά.
\sol{10} $(-18):6=-3$ βαθμούς ανά ώρα. Και $(-12):4=-3$.
\sol{11} $(-36):9=-4$ μ. το λεπτό.
\sol{12} $(-15):(-2{,}5)=6$ ώρες.
\sol{13} Άθροισμα $-10$, μέσος όρος $-10:7\approx-1{,}4$ βαθμοί.
\sol{14} $(-8+4-10+6):4=-8:4=-2$ πόντους.
\sol{15} $(-45):(-5)=9$ παιδιά.
\sol{16} Η $0\cdot x=5$ δεν έχει λύση. Η $0\cdot x=0$ έχει άπειρες. Σε καμία περίπτωση δεν έχουμε μοναδική λύση, άρα η $\alpha:0$ δεν ορίζεται.
\sol{17} Ετερόσημοι δίνουν αρνητικό: $(-12):4=-3$. Έλεγχος: $4\cdot(-3)=-12$.
\sol{18} $-2$, $-2$, $+2$.
\sol{19} $\beta=\alpha:2=-5$.
\sol{20} $\frac{-6}{3}=-2$, $\frac{6}{-3}=-2$ και $-\frac{6}{3}=-2$: ίσα.
\sol{21} Ελεύθερη απάντηση. Π.χ. $-30$ € σε $5$ παιδιά: $-6$ € το καθένα.

\lyseis{7.7}{Δεκαδική μορφή ρητών}
\sol{1} α) $0{,}375$ \quad β) $0{,}8\overline{3}$ \quad γ) $0{,}\overline{36}$ \quad δ) $0{,}\overline{7}$ \quad ε) $0{,}45$ \quad στ) $0{,}1\overline{3}$ \quad ζ) $0{,}08\overline{3}$.
\sol{2} Ακριβείς: $\frac{7}{20}$ ($20=2^2\cdot5$), $\frac{11}{40}$ ($40=2^3\cdot5$), $\frac{3}{16}$ ($16=2^4$), $\frac{13}{25}$ ($25=5^2$). Περιοδικοί: $\frac{5}{12}$ και $\frac{9}{14}$.
\sol{3} $0{,}\overline{36}=0{,}36363636$, $2{,}1\overline{45}=2{,}14545454$, $0{,}\overline{142857}=0{,}14285714$.
\sol{4} α) $\frac{7}{9}$ \quad β) $\frac{15}{99}=\frac{5}{33}$ \quad γ) $\frac{4}{3}$ \quad δ) $\frac{81}{99}=\frac{9}{11}$ \quad ε) $\frac{7}{30}$.
\sol{5} α) $0{,}\overline{3}>0{,}33$ \quad β) $0{,}\overline{6}=\frac{2}{3}$ \quad γ) $0{,}7\overline{2}=0{,}7222\ldots>0{,}72$.
\sol{6} α) $0{,}67$ \quad β) $0{,}83$ \quad γ) $0{,}14$ \quad δ) $0{,}44$.
\sol{7} α) Σωστό. β) Λάθος: $0{,}\overline{3}$. γ) Σωστό. δ) Σωστό. ε) Σωστό. στ) Λάθος: $0{,}1\overline{6}$.
\sol{8} Η περίοδος είναι $142857$. $2:7=0{,}\overline{285714}$ και $3:7=0{,}\overline{428571}$: ίδια ψηφία με το $1:7$, ξεκινούν από άλλο σημείο.
\sol{9} Το $10:3=3{,}\overline{3}$. Στην πράξη το καθένα παίρνει $3{,}33$ € και μένει $1$ λεπτό.
\sol{10} $\frac{100}{12}=\frac{25}{3}=8{,}\overline{3}$ m/s $\approx8{,}33$ m/s.
\sol{11} $\frac{1}{3}=33{,}\overline{3}\%\approx33{,}3\%$ και $\frac{5}{6}=83{,}\overline{3}\%\approx83{,}3\%$.
\sol{12} $\frac{23}{3}=7{,}\overline{6}\approx7{,}67$.
\sol{13} Τα υπόλοιπα της διαίρεσης με το $7$ είναι $1,\dots,6$. Αφού υπάρχουν μόνο έξι διαφορετικά, μετά από το πολύ έξι βήματα επαναλαμβάνεται κάποιο και μαζί του και τα ψηφία.
\sol{14} $x=0{,}\overline{9}$, $10x=9{,}\overline{9}$, $10x-x=9$, $9x=9$, $x=1$.
\sol{15} $\frac{1}{11}=0{,}\overline{09}$ και $\frac{1}{13}=0{,}\overline{076923}$.
\sol{16} Για παράδειγμα $0{,}335$, αφού $0{,}333\ldots<0{,}335<0{,}34$.
\sol{17} Το $0{,}333$ είναι προσέγγιση. Αν τα $3$ της περιόδου συνεχίζονται για πάντα, $3\cdot0{,}\overline{3}=0{,}\overline{9}=1$, όπως $3\cdot\frac{1}{3}=1$.
\sol{18} Ακριβή αναπτύγματα έχουν τα $\frac{1}{2},\frac{1}{4},\frac{1}{5},\frac{1}{8},\frac{1}{10}$ (και $\frac{1}{16}$ κ.λπ.). Οι παρονομαστές τους έχουν μόνο παράγοντες $2$ και $5$.

\lyseis{7.8}{Δυνάμεις με εκθέτη φυσικό}
\sol{1} α) $32$ \quad β) $81$ \quad γ) $-8$ \quad δ) $16$ \quad ε) $-1$ \quad στ) $0$.
\sol{2} α) $9$ \quad β) $-9$ \quad γ) $-\frac{1}{8}$ \quad δ) $0{,}001$ \quad ε) $0{,}04$ \quad στ) $\frac{8}{27}$.
\sol{3} α) Θετικό (άρτιος εκθέτης). β) Αρνητικό. γ) Αρνητικό ($(-3)^4=81$ και μετά το μείον). δ) Αρνητικό ($(-2)^3=-8$ και $(-3)^2=9$).
\sol{4} α) $2^7$ \quad β) $5^3$ \quad γ) $3^6$ \quad δ) $7^6$ \quad ε) $(-2)^7$ \quad στ) $(-4)^2=16$.
\sol{5} α) $10^3=1\,000$ \quad β) $(-6)^2=36$ \quad γ) $\frac{8}{27}$ \quad δ) $\frac{9}{16}$.
\sol{6} α) $18+8-2=24$ \quad β) $12-16:(-2)=12+8=20$ \quad γ) $5-12=-7$ \quad δ) $(-1+1)\cdot10=0$.
\sol{7} α) Λάθος: $8$. β) Σωστό. γ) Λάθος: $-16$. δ) Σωστό. ε) Λάθος: $\alpha^5$. στ) Σωστό. ζ) Λάθος.
\sol{8} Πλήθη: $1,\,2,\,4,\,8,\,16,\,32,\,64,\,128,\,256,\,512,\,1\,024$. Ξεπερνά τα $1\,000$ μετά από $10$ ώρες.
\sol{9} Όγκος $12^3=1\,728$ cm$^3$. Επιφάνεια $6\cdot144=864$ cm$^2$.
\sol{10} Στο $8$ο: $2^7=128$. Στο $11$ο: $2^{10}=1\,024$.
\sol{11} $3^5=243$ βακτήρια.
\sol{12} Ο $20$ός είναι $20^2=400$ και ο $\nu$-οστός $\nu^2$.
\sol{13} $3\cdot2^4=48$ €.
\sol{14} $10^3\cdot10^4=10^7$ και $10^8:10^5=10^3$.
\sol{15} Το $2\cdot3$ είναι $2+2+2=6$ (πρόσθεση του $2$ τρεις φορές). Το $2^3=2\cdot2\cdot2=8$ (πολλαπλασιασμός του $2$ τρεις φορές).
\sol{16} $(3+4)^2=49$ ενώ $3^2+4^2=9+16=25$.
\sol{17} Όταν ο εκθέτης $\nu$ είναι άρτιος $(-\alpha)^\nu=\alpha^\nu$, όταν είναι περιττός $(-\alpha)^\nu=-\alpha^\nu$.
\sol{18} $1^\nu=1$, $0^\nu=0$ και $(-1)^\nu=1$ για άρτιο και $-1$ για περιττό $\nu$.
\sol{19} $2^3\cdot2^4=(2\cdot2\cdot2)\cdot(2\cdot2\cdot2\cdot2)=2\cdot2\cdot2\cdot2\cdot2\cdot2\cdot2=2^7$.
\sol{20} $3^4=81$ και $4^3=64$, άρα το $3^4$.
\sol{21} $2,\,4,\,8,\,16,\,32$ στρώματα, δηλαδή $2^\nu$.

\lyseis{7.9}{Δυνάμεις με εκθέτη ακέραιο}
\sol{1} α) $\frac{1}{2}$ \quad β) $\frac{1}{8}$ \quad γ) $\frac{1}{100}$ \quad δ) $\frac{1}{25}$ \quad ε) $1$ \quad στ) $1$.
\sol{2} α) $0{,}1$ \quad β) $0{,}001$ \quad γ) $0{,}00001$.
\sol{3} α) $-\frac{1}{8}$ \quad β) $\frac{1}{4}$ \quad γ) $-1$ \quad δ) $8$ \quad ε) $\frac{9}{4}$ \quad στ) $-\frac{4}{3}$.
\sol{4} α) $3^{-2}$ \quad β) $2^{-5}$ \quad γ) $10^{-4}$ \quad δ) $(-5)^{-3}$.
\sol{5} α) $2^2=4$ \quad β) $10^6$ \quad γ) $\frac{1}{9}$ \quad δ) $\frac{1}{100}$ \quad ε) $7^0=1$ \quad στ) $(-2)^2=4$.
\sol{6} α) $4+1-\frac{1}{2}=\frac{9}{2}$ \quad β) $\frac{1}{4}+\frac{1}{2}=\frac{3}{4}$ \quad γ) $10^{2-5+3}=10^0=1$ \quad δ) $\frac{1}{9}$ \quad ε) $1\cdot\frac{1}{4}=\frac{1}{4}$.
\sol{7} α) $-4$ \quad β) $-3$ \quad γ) $-3$ \quad δ) $-3$ (αφού $\left(\frac{1}{2}\right)^{-3}=8$) \quad ε) $0$.
\sol{8} α) Λάθος: $\frac{1}{2}$. β) Σωστό. γ) Σωστό. δ) Λάθος: $0{,}001$. ε) Σωστό. στ) Σωστό. ζ) Λάθος: $\frac{1}{4}$.
\sol{9} $10^3$ μ., $10^{-2}$ μ., $10^{-3}$ μ., $10^{-6}$ μ., $10^{-9}$ μ.
\sol{10} Μένει $\left(\frac{1}{2}\right)^5=\frac{1}{32}=2^{-5}$.
\sol{11} Μετά από $3$ ημέρες $64\cdot2^{-3}=8$ g. Πριν από $2$ ημέρες $64\cdot2^2=256$ g.
\sol{12} $3^4=81$, $3^3=27$, $3^2=9$, $3^1=3$, $3^0=1$, $3^{-1}=\frac{1}{3}$, $3^{-2}=\frac{1}{9}$, $3^{-3}=\frac{1}{27}$.
\sol{13} $400\cdot2^{-3}=50$ €.
\sol{14} Στις δυνάμεις του $2$ όταν ο εκθέτης μειώνεται κατά $1$ η τιμή γίνεται το μισό: $2^1=2$, $2^0=1$. Το $1$ είναι το μισό του $2$, όχι το $0$. Το ίδιο ισχύει για κάθε βάση.
\sol{15} $\left(\frac{2}{3}\right)^{-2}=\frac{1}{(2/3)^2}=\frac{1}{4/9}=\frac{9}{4}$ και $\left(\frac{3}{2}\right)^2=\frac{9}{4}$.
\sol{16} $2^{-3}=\frac{1}{8}$ (θετικό), ενώ $-2^3=-8$.
\sol{17} $\alpha^{-3}\cdot\alpha^3=\alpha^0$ και, αφού $\alpha^{-3}=\frac{1}{\alpha^3}$, το γινόμενο είναι $1$. Άρα $\alpha^0=1$.
\sol{18} Θα ήταν $\frac{1}{0}$, και δεν διαιρούμε με το $0$.
\sol{19} $-1$, $1$, $-1$. Το $(-1)^{-\nu}=(-1)^\nu$: $1$ για άρτιο και $-1$ για περιττό.
\sol{20} $0{,}1$ mm $=10^{-1}\cdot10^{-3}$ μ. $=10^{-4}$ μ.

\lyseis{7.10}{Τυποποιημένη μορφή μεγάλων και μικρών αριθμών}
\sol{1} α) $5\cdot10^{-4}$ \quad β) $7{,}2\cdot10^{-5}$ \quad γ) $3\cdot10^{-1}$ \quad δ) $4{,}5\cdot10^{4}$ \quad ε) $8\cdot10^{-7}$ \quad στ) $6{,}3\cdot10^{6}$.
\sol{2} α) $0{,}000032$ \quad β) $0{,}0091$ \quad γ) $0{,}4$ \quad δ) $2\,500\,000$ \quad ε) $0{,}000101$ \quad στ) $7$.
\sol{3} α) $3\cdot10^{-3}$ \quad β) $1{,}2\cdot10^{-3}$ \quad γ) είναι ήδη τυποποιημένο \quad δ) $1\cdot10^{3}$.
\sol{4} α) $>$ ($-5>-6$) \quad β) $<$ ($-3<-2$) \quad γ) $>$ ($4>-4$) \quad δ) $>$ ($3{,}1>3{,}09$ με ίδιο εκθέτη).
\sol{5} $5\cdot10^{-5}<1{,}2\cdot10^{-3}<3\cdot10^{-3}<9\cdot10^{-1}<2\cdot10^{2}$.
\sol{6} α) $6\cdot10^{2}$ \quad β) $3\cdot10^{10}$ \quad γ) $20\cdot10^{2}=2\cdot10^{3}$ \quad δ) $2\cdot10^{-3}$.
\sol{7} $0{,}00045$, $25\,000\,000$ και $0{,}000\,000\,007$.
\sol{8} α) Σωστό. β) Λάθος: $1{,}2\cdot10^{-3}$. γ) Λάθος: το $10^{-5}$ είναι μεγαλύτερο. δ) Σωστό. ε) Σωστό. στ) Λάθος: $5{,}5\cdot10^{-3}<5{,}5\cdot10^{-2}$.
\sol{9} $10^{-7}<2\cdot10^{-6}<7\cdot10^{-6}<7\cdot10^{-5}$ (ιός, βακτήριο, αιμοσφαίριο, τρίχα). Η τρίχα είναι $10$ φορές χοντρότερη από το αιμοσφαίριο.
\sol{10} $9{,}1\cdot10^{-31}$ έχει $30$ μηδενικά μετά την υποδιαστολή πριν από το $9$ ($0{,}000\ldots09\ldots$ με $30$ μηδενικά).
\sol{11} $\frac{1{,}5\cdot10^{11}}{3\cdot10^{8}}=5\cdot10^{2}=500$ s, δηλαδή περίπου $8{,}3$ λεπτά.
\sol{12} $2\cdot10^{-9}$ μ.
\sol{13} $500\cdot10^{-4}=5\cdot10^{-2}$ μ. $=5$ εκατοστά.
\sol{14} $10^{6}$, δηλαδή ένα εκατομμύριο κόκκοι.
\sol{15} Γιατί για να πάμε στον αριθμό από το $\alpha\ge1$ διαιρούμε με δύναμη του $10$ (μετακινούμε την υποδιαστολή αριστερά), και αυτό είναι ο αρνητικός εκθέτης $10^{-\nu}=\frac{1}{10^\nu}$.
\sol{16} $1=1\cdot10^{0}$ και $10=1\cdot10^{1}$.
\sol{17} Δεκαδικά: $0{,}002>0{,}0005$. Εκθέτες: $-3>-4$. Άρα $2\cdot10^{-3}>5\cdot10^{-4}$.
\sol{18} $3\cdot10^{-4}\cdot10^{6}=3\cdot10^{2}=300$.
\sol{19} Γιατί ο $12$ δεν είναι μικρότερος από το $10$. Γράφεται $1{,}2\cdot10^{-3}$.
\sol{20} Ελεύθερη απάντηση. Π.χ. άτομο $\approx10^{-10}$ μ., ιός $\approx10^{-7}$ μ., βακτήριο $\approx2\cdot10^{-6}$ μ.

\lyseisE{Επανάληψη κεφαλαίου 7}
\solr{7}{1} α) $-8$ \quad β) $-25$ \quad γ) $-40$.
\solr{7}{2} Ακέραιοι: $-3$, $\frac{4}{2}=2$, $-\frac{7}{7}=-1$, $0$. Ο $0{,}5$ δεν είναι.
\solr{7}{3} α) $7{,}5$ \quad β) $\frac{3}{4}$ \quad γ) $+\frac{2}{3}$ \quad δ) $\alpha=4$ ή $\alpha=-4$.
\solr{7}{4} $-3<-2{,}5<-\frac{1}{2}<0<3$.
\solr{7}{5} α) $5$ \quad β) $-5$ \quad γ) $-\frac{1}{6}$ \quad δ) $0$.
\solr{7}{6} α) $8$ \quad β) $-10$ \quad γ) $0$ \quad δ) $-\frac{3}{4}$.
\solr{7}{7} $-12+7+5-3+9=6$.
\solr{7}{8} α) $42$ \quad β) $-10$ \quad γ) $-6$ \quad δ) $(-2)\cdot(-3)=6$.
\solr{7}{9} α) $8$ \quad β) $-3$ \quad γ) $\left(-\frac{3}{4}\right)\cdot\left(-\frac{2}{9}\right)=\frac{1}{6}$ \quad δ) $0$.
\solr{7}{10} $\frac{5}{6}=0{,}8\overline{3}$, $\frac{7}{20}=0{,}35$, $0{,}\overline{18}=\frac{18}{99}=\frac{2}{11}$.
\solr{7}{11} α) $-32$ \quad β) $-16$ \quad γ) $9$ \quad δ) $1+\frac{1}{8}=\frac{9}{8}$.
\solr{7}{12} α) $2^2=4$ \quad β) $3^{-2}=\frac{1}{9}$ \quad γ) $10^{-2}$.
\solr{7}{13} $0{,}00032=3{,}2\cdot10^{-4}$, $45\,000\,000=4{,}5\cdot10^{7}$, και $(4\cdot10^{3})(5\cdot10^{-5})=20\cdot10^{-2}=2\cdot10^{-1}$.
\solr{7}{14} $\frac{(-5)-4}{3}=-3$ βαθμούς ανά ώρα.
\solr{7}{15} α) Λάθος: $5$. β) Σωστό. γ) Λάθος: $5-(-2)=7$. δ) Σωστό. ε) Λάθος: $-3$. στ) Σωστό.
\solr{7}{16} α) Σωστό. β) Λάθος: $\frac{1}{2}$. γ) Σωστό. δ) Σωστό. ε) Λάθος. στ) Σωστό.


\lyseismeros{Β}{ΓΕΩΜΕΤΡΙΑ}
\lyseiskef{Β1}{Μέτρηση}
\lyseis{Β1.1}{Βασικές έννοιες: σημείο, ευθεία, ημιευθεία, ευθύγραμμο τμήμα, γωνία, κύκλος}
\sol{1} α) σημείου. β) ευθεία. γ) άπειρες, μία και μοναδική. δ) ημιευθεία. ε) αντικείμενες.
\sol{2} Ελεύθερη κατασκευή. Προκύπτουν $3$ ευθείες ($\mathrm{ΑΒ}$, $\mathrm{ΑΓ}$, $\mathrm{ΒΓ}$) και $3$ ευθύγραμμα τμήματα ($\mathrm{ΑΒ}$, $\mathrm{ΑΓ}$, $\mathrm{ΒΓ}$).
\sol{3} α) $\mathrm{ΚΛ}$ \quad β) $\mathrm{Κx}$ (με το $x$ προς το $\Lambda$) \quad γ) η ευθεία $\mathrm{ΚΛ}$.
\sol{4} Ελεύθερη κατασκευή. Κάθε ευθεία χωρίζεται από το $\mathrm{Ο}$ σε δύο αντικείμενες ημιευθείες, άρα βλέπεις $2\cdot2=4$ ημιευθείες συνολικά.
\sol{5} Ελεύθερη κατασκευή. Π.χ. $x\hat{\mathrm{Ο}}y$ και $y\hat{\mathrm{Ο}}x$ (ή $\hat{\mathrm{Ο}}$). Η κυρτή γωνία είναι η μικρότερη περιοχή και η μη κυρτή η μεγαλύτερη.
\sol{6} Ελεύθερη κατασκευή. Και οι τρεις ακτίνες βγαίνουν $2{,}5$ cm: όλες οι ακτίνες του ίδιου κύκλου είναι ίσες.
\sol{7} α) Λάθος: η ευθεία δεν έχει άκρα. β) Σωστό. γ) Λάθος: έχουν διαφορετική αρχή. δ) Σωστό. ε) Σωστό. στ) Λάθος: έχουν το ίδιο κέντρο, αλλά οι ακτίνες τους είναι διαφορετικές.
\sol{8} Αρκούν δύο παλούκια. Τεντώνει το σχοινί ανάμεσά τους, το προεκτείνει προς τις δύο μεριές και φυτεύει κατά μήκος του.
\sol{9} Η δέσμη είναι ημιευθεία με αρχή το $\mathrm{Α}$. Οι δύο φακοί που φωτίζουν ο ένας τον άλλον δίνουν δύο ημιευθείες πάνω στην ίδια ευθεία, με αντίθετες κατευθύνσεις. Αν οι φακοί ήταν στο ίδιο σημείο, θα ήταν αντικείμενες ημιευθείες.
\sol{10} Κάθε ζευγάρι ορίζει ένα σχοινί: $\dfrac{5\cdot4}{2}=10$ σχοινιά.
\sol{11} Τέμνουν το μονοπάτι τα $\mathrm{ΚΜ}$ και $\Lambda\mathrm{Μ}$, γιατί ενώνουν σημεία από διαφορετικά ημιεπίπεδα. Το $\mathrm{ΚΛ}$ δεν το τέμνει.
\sol{12} Ο κύκλος έχει ακτίνα $1{,}2$ μ. Το φυτό στα $1{,}2$ μ. είναι πάνω στη γραμμή του κύκλου. Το φυτό στα $1{,}6$ μ. είναι έξω από τον κύκλο (και αυτό στα $0{,}9$ μ. είναι μέσα).
\sol{13} Η ευθεία είναι άπειρη προς τις δύο μεριές και δεν έχει άκρα, γι' αυτό δεν έχει μήκος. Το τμήμα έχει άκρα και μήκος. Η ημιευθεία είναι άπειρη προς μία μεριά, άρα δεν μετριέται.
\sol{14} Κάθε σημείο ορίζει δύο ημιευθείες, άρα $3\cdot2=6$.
\sol{15} Όχι: από δύο διαφορετικά σημεία περνά μία και μοναδική ευθεία. Αν περνούσαν δύο, δεν θα ήταν διαφορετικές.
\sol{16} Όχι. Ένα κοινό σημείο θα απείχε από το $\mathrm{Ο}$ ταυτόχρονα $3$ cm και $5$ cm, που είναι αδύνατον.
\sol{17} Ελεύθερη απάντηση. Π.χ. σημείο: το ίχνος της μύτης του μολυβιού· τμήμα: η ακμή του τραπεζιού· επίπεδο: το δάπεδο· ημιευθεία: η δέσμη ενός φακού.
\sol{18} Ελεύθερη κατασκευή. Η δεύτερη καρφίτσα ορίζει τη θέση του νήματος και το νήμα δείχνει την ευθεία. Ο κύκλος έχει ως κέντρο το σημείο που διάλεξες και περνά από το άλλο σημείο.

\lyseis{Β1.2}{Γωνία, γραμμή, σχήματα, ίσα σχήματα}
\sol{1} Κορυφή το $\mathrm{Ο}$, πλευρές οι ημιευθείες $\mathrm{Οx}$ και $\mathrm{Οy}$, εσωτερικό τα σημεία ανάμεσά τους.
\sol{2} Δύο γωνίες: μία κυρτή και μία μη κυρτή.
\sol{3} α) ευθεία (ανοιχτή) \quad β) ανοιχτή τεθλασμένη \quad γ) κλειστή καμπύλη \quad δ) ανοιχτή καμπύλη \quad ε) κλειστή τεθλασμένη. Τεθλασμένες είναι οι β και ε.
\sol{4} Τρίγωνο: $3,\,3,\,0$. Τετράπλευρο: $4,\,4,\,2$. Πεντάγωνο: $5,\,5,\,5$. Εξάγωνο: $6,\,6,\,9$. Οκτάγωνο: $8,\,8,\,20$.
\sol{5} $\mathrm{ΑΓ}$, $\mathrm{ΑΔ}$, $\mathrm{ΒΔ}$, $\mathrm{ΒΕ}$, $\Gamma\mathrm{Ε}$.
\sol{6} Ίσα είναι τα ζεύγη α και γ.
\sol{7} α) Σωστό. β) Σωστό. γ) Λάθος. δ) Λάθος. ε) Σωστό.
\sol{8} $5$ σχοινιά, όσες και οι διαγώνιοι.
\sol{9} Ελεύθερη απάντηση: ίσα είναι τα πλακάκια που εφαρμόζουν ακριβώς.
\sol{10} Ανοιχτή τεθλασμένη. Αν ενώσουμε το $\mathrm{Δ}$ με το $\mathrm{Α}$, παίρνουμε τετράπλευρο.
\sol{11} Τέσσερις γωνίες.
\sol{12} Ελεύθερη απάντηση. Π.χ. οι γωνίες του πίνακα, του τοίχου, του θρανίου.
\sol{13} Οι δύο ημιευθείες χωρίζουν το επίπεδο σε δύο μέρη. Κυρτή είναι αυτή που είναι το μικρότερο μέρος.
\sol{14} Από κάθε κορυφή φεύγουν $\nu-3$ διαγώνιοι (όχι προς τον εαυτό της και τις δύο γειτονικές). Σε όλες τις κορυφές $\nu(\nu-3)$, αλλά κάθε διαγώνιος μετράται δύο φορές, άρα $\frac{\nu(\nu-3)}{2}$. Στο δωδεκάγωνο $\frac{12\cdot9}{2}=54$.
\sol{15} Όχι: ένα ορθογώνιο $2\times6$ και ένα $3\times4$ έχουν ίδιο εμβαδόν αλλά δεν είναι ίσα.
\sol{16} Όχι με την έννοια του πολυγώνου που ξέρουμε, γιατί η γραμμή τέμνει τον εαυτό της και δεν ορίζει ένα εσωτερικό.
\sol{17} Ελεύθερη απάντηση.
\sol{18} Ελεύθερη απάντηση.

\lyseis{Β1.3}{Μέτρο γωνίας}
\sol{1} Ελεύθερη κατασκευή. Η $\hat{\gamma}$ είναι ορθή και η $\hat{\varepsilon}$ έχει πλευρές αντικείμενες ημιευθείες. Έλεγξε κάθε γωνία μετρώντας ξανά.
\sol{2} Ελεύθερη κατασκευή. Και οι δύο έχουν μέτρο $40^\circ$, άρα είναι ίσες: το μέτρο εξαρτάται μόνο από το άνοιγμα και όχι από το μήκος των πλευρών.
\sol{3} α) $12{,}5^\circ$, γιατί $30'=0{,}5^\circ$. β) $40{,}25^\circ$, γιατί $15'=0{,}25^\circ$. γ) $18^\circ45'$, γιατί $0{,}75\cdot60=45$. δ) $1^\circ=60'=3600''$.
\sol{4} Σε μοίρες: $45$, $40{,}5$, $40{,}25$, $45{,}17$. Σειρά: $40{,}25^\circ<40^\circ30'<45^\circ<45^\circ10'$.
\sol{5} α) $x\hat{\mathrm{Ο}}z=35^\circ+48^\circ=83^\circ$. β) $y\hat{\mathrm{Ο}}z=120^\circ-75^\circ=45^\circ$.
\sol{6} Ελεύθερη κατασκευή. Σε κάθε τρίγωνο το άθροισμα των τριών γωνιών βγαίνει περίπου $180^\circ$ (με μικρές αποκλίσεις λόγω μέτρησης).
\sol{7} α) Σωστό. β) Λάθος: το μέτρο δεν εξαρτάται από το μήκος των πλευρών. γ) Σωστό. δ) Λάθος: $1^\circ=60'$. ε) Σωστό. στ) Σωστό.
\sol{8} Συνολική στροφή $45^\circ+30^\circ=75^\circ$. Μετά την αντίθετη στροφή απομένει $75^\circ-20^\circ=55^\circ$ από την αρχική πορεία.
\sol{9} Κάθε ώρα αντιστοιχεί σε $360^\circ:12=30^\circ$. Στις $2{:}00$: $60^\circ$. Στις $5{:}00$: $150^\circ$. Στις $6{:}00$: $180^\circ$.
\sol{10} $84^\circ-50^\circ=34^\circ$.
\sol{11} $35^\circ+50^\circ=85^\circ$ και μετά $85^\circ-20^\circ=65^\circ$ ανοιχτή.
\sol{12} Η γωνία φαίνεται μικρότερη από την ορθή, άρα το μέτρο της είναι μικρότερο από $90^\circ$: το σωστό είναι $50^\circ$. Οι δύο κλίμακες δίνουν άθροισμα $130^\circ+50^\circ=180^\circ$.
\sol{13} Οι πλευρές μιας γωνίας είναι ημιευθείες και δεν έχουν μήκος· το μέτρο δείχνει μόνο πόσο ανοίγουν. Αν μακρύνουμε ή κονταίνουμε τη γραμμή που σχεδιάσαμε, το άνοιγμα δεν αλλάζει.
\sol{14} $x\hat{\mathrm{Ο}}t=70^\circ+50^\circ+40^\circ=160^\circ$.
\sol{15} Αντιγράφουμε τη μία γωνία σε διαφανές χαρτί και την τοποθετούμε πάνω στην άλλη με κοινή κορυφή και μία κοινή πλευρά. Αν οι άλλες πλευρές συμπίπτουν, είναι ίσες. Αλλιώς, η γωνία που μένει στο εσωτερικό της άλλης είναι η μικρότερη.
\sol{16} Ελεύθερη κατασκευή. α) $x\hat{\mathrm{Ο}}t=60^\circ$. β) $x\hat{\mathrm{Ο}}t=15^\circ$. Έλεγξε με το μοιρογνωμόνιο.
\sol{17} Ελεύθερη απάντηση: π.χ. η ανοιχτή πόρτα $\approx$ μερικές δεκάδες μοίρες, η γωνία του τραπεζιού $90^\circ$.
\sol{18} Ελεύθερη κατασκευή. Όσο ανοίγει ο διαβήτης, τόσο μεγαλώνει το μέτρο της γωνίας των σκελών του.

\lyseisE{Επανάληψη κεφαλαίου Β1}
\solr{Β1}{1} Τρεις ευθείες και έξι ευθείες.
\solr{Β1}{2} Πεντάγωνο $5$, εξάγωνο $9$.
\solr{Β1}{3} Το μέσο απέχει $3{,}6$ cm από κάθε άκρο. Επειδή $7{,}2$ cm $=72$ mm $>70$ mm, το $\mathrm{ΑΒ}$ είναι μεγαλύτερο.
\solr{Β1}{4} $\mathrm{ΑΔ}=2{,}5+3+4=9{,}5$ και $\mathrm{ΒΔ}=3+4=7$.
\solr{Β1}{5} $2(6{,}5+4)=21$.
\solr{Β1}{6} Διχοτόμος: κάθε μισή $62^\circ$. Παραπλήρωμα $56^\circ$. Συμπλήρωμα της $36^\circ$: $54^\circ$.
\solr{Β1}{7} $65^\circ$, $115^\circ$ και $115^\circ$.
\solr{Β1}{8} $360^\circ-(90+80+110)^\circ=80^\circ$.
\solr{Β1}{9} $150^\circ$, αμβλεία.
\solr{Β1}{10} Το πολύ $\frac{4\cdot3}{2}=6$ σημεία. Μία παράλληλη.
\solr{Β1}{11} $1{,}8+1{,}8=3{,}6$ cm.
\solr{Β1}{12} Διάμετρος $12$. Το $\mathrm{Α}$ μέσα ($4<6$), το $\mathrm{Β}$ έξω ($7>6$).
\solr{Β1}{13} $360:9=40^\circ$. Το τόξο είναι $150^\circ$.
\solr{Β1}{14} α) εξωτερική ($0$ κοινά) \quad β) εφαπτομένη ($1$) \quad γ) τέμνουσα ($2$).
\solr{Β1}{15} α) Σωστό. β) Λάθος. γ) Σωστό. δ) Σωστό. ε) Σωστό.
\solr{Β1}{16} α) Σωστό. β) Λάθος: $90^\circ$. γ) Σωστό. δ) Σωστό. ε) Σωστό. στ) Λάθος: είναι το μικρότερο.


\lyseiskef{Β2}{Μετασχηματισμοί}
\lyseis{Β2.1}{Συμμετρία ως προς άξονα}
\sol{1} Ελεύθερη κατασκευή. Η κάθετη από το $\mathrm{Α}$ στην $\varepsilon$ περνά από το μέσο του $\mathrm{ΑΑ}'$.
\sol{2} Το ίδιο το σημείο.
\sol{3} $3{,}5$ cm από την $\varepsilon$ και $\mathrm{ΑΑ}'=7$ cm.
\sol{4} Ελεύθερη κατασκευή. Βρίσκω τα συμμετρικά των τριών κορυφών και τα ενώνω.
\sol{5} Το συμμετρικό τρίγωνο έχει $\mathrm{Α}'\mathrm{Β}'=5$ cm και $\widehat{\mathrm{Β}'}=40^\circ$.
\sol{6} Το συμμετρικό τμήμα έχει μήκος $4{,}2$ cm.
\sol{7} α) Σωστό. β) Λάθος: κάθετο. γ) Σωστό. δ) Σωστό. ε) Λάθος: διατηρεί τα μήκη.
\sol{8} Στο συμμετρικό σημείο: $1{,}2$ μ. από το μονοπάτι, από την άλλη μεριά. Τα δύο δέντρα απέχουν $2{,}4$ μ.
\sol{9} Η αντανάκλαση φαίνεται $4$ μ. βαθιά, όσο ύψος έχει το δέντρο.
\sol{10} Πραγματικά είναι $9{:}00$ (ο ωροδείκτης στον καθρέφτη είναι αριστερά αντί δεξιά).
\sol{11} Μένουν ίδια τα $\mathrm{Α}$, $\mathrm{Η}$, $\mathrm{Μ}$, $\mathrm{Ο}$, $\mathrm{Τ}$.
\sol{12} Ελεύθερη κατασκευή.
\sol{13} Η δίπλωση φέρνει το $\mathrm{Α}$ πάνω στο $\mathrm{Α}'$ μόνο αν το $\mathrm{ΑΑ}'$ είναι κάθετο στον άξονα και ο άξονας το περνά από το μέσο του, αλλιώς τα δύο σημεία δεν θα έπεφταν το ένα πάνω στο άλλο.
\sol{14} Τα σημεία της ίδιας της $\varepsilon$.
\sol{15} Επιστρέφει στο $\mathrm{Α}$.
\sol{16} Το $\mathrm{Μ}$ ανήκει στην $\varepsilon$, άρα είναι συμμετρικό του εαυτού του. Άρα το $\mathrm{Α}'\mathrm{Β}'$ περνά από το $\mathrm{Μ}$.
\sol{17} Ελεύθερη απάντηση: σχηματίζεται ένα συμμετρικό σχήμα με άξονα την πτυχή.
\sol{18} Ελεύθερη απάντηση.

\lyseisE{Επανάληψη κεφαλαίου Β2}
\solr{Β2}{1} Το $\mathrm{Α}'$ είναι από την άλλη μεριά του $\xi$, πάλι $3{,}4$ cm από τον άξονα. Τα $\mathrm{Α}$ και $\mathrm{Α}'$ απέχουν $6{,}8$ cm.
\solr{Β2}{2} Ισόπλευρο: $3$. Ορθογώνιο: $2$. Ρόμβος: $2$. Κύκλος: άπειροι.
\solr{Β2}{3} $5$ cm (ισαπέχει από τα άκρα).
\solr{Β2}{4} $5$ cm.
\solr{Β2}{5} Η ευθεία που περνά από το $\mathrm{Α}$ και το μέσο της $\mathrm{ΒΓ}$ (το ύψος από το $\mathrm{Α}$).
\solr{Β2}{6} Ως προς το $0$: $2$. Ως προς το $1$: $4$ (αφού $\frac{-2+4}{2}=1$).
\solr{Β2}{7} $14$ cm.
\solr{Β2}{8} Άξονα συμμετρίας έχουν τα Α, Β, Η, Τ, Φ. Κέντρο συμμετρίας έχουν τα Η, Ν, Σ, Φ.
\solr{Β2}{9} Τέσσερις γωνίες $80^\circ$ και τέσσερις $100^\circ$.
\solr{Β2}{10} $2x+(x+60^\circ)=180^\circ$, $3x=120^\circ$, $x=40^\circ$. Οι γωνίες είναι $80^\circ$ και $100^\circ$.
\solr{Β2}{11} Όχι, γιατί οι εντός εναλλάξ δεν είναι ίσες.
\solr{Β2}{12} Της διαγωνίου $\mathrm{ΑΓ}$.
\solr{Β2}{13} Στο μέσο της διαγωνίου, $6{,}5$ από κάθε κορυφή.
\solr{Β2}{14} $\mathrm{Α}'\mathrm{Β}'=8{,}5$ cm και είναι παράλληλο προς το $\mathrm{ΑΒ}$.
\solr{Β2}{15} α) Σωστό. β) Σωστό. γ) Λάθος. δ) Σωστό.
\solr{Β2}{16} α) Σωστό. β) Λάθος. γ) Σωστό. δ) Σωστό.


\lyseiskef{Β3}{Τρίγωνα – Παραλληλόγραμμα – Τραπέζια}
\lyseis{Β3.1}{Στοιχεία και είδη τριγώνων}
\sol{1} α) ισόπλευρο β) ισοσκελές γ) σκαληνό.
\sol{2} $6{,}5+7+8{,}5=22$ cm.
\sol{3} Ελεύθερη κατασκευή. Στο ορθογώνιο δύο ύψη είναι οι κάθετες πλευρές, στο αμβλυγώνιο δύο ύψη βρίσκονται έξω από το τρίγωνο.
\sol{4} $\mathrm{ΒΜ}=4{,}4:2=2{,}2$ cm.
\sol{5} Οι τρεις διάμεσοι περνούν από το ίδιο σημείο.
\sol{6} Ελεύθερη κατασκευή. Οι δύο γωνίες είναι ίσες.
\sol{7} α) Σωστό. β) Λάθος. γ) Σωστό. δ) Σωστό. ε) Λάθος.
\sol{8} Ισοσκελές, περίμετρος $30+30+40=100$ cm.
\sol{9} $3\cdot60=180$ cm.
\sol{10} $24:3=8$ μ.
\sol{11} $(50-20):2=15$ cm.
\sol{12} Με $3,4,9$ όχι (γιατί $3+4=7<9$). Με $3,4,5$ ναι ($3+4=7>5$ κ.λπ.).
\sol{13} Οξυγώνιο, ορθογώνιο ή αμβλυγώνιο.
\sol{14} Όχι: ένα τρίγωνο με ορθή γωνία έχει τις άλλες δύο οξείες, άρα πλευρές όχι όλες ίσες. Το ισόπλευρο έχει όλες τις γωνίες $60^\circ$.
\sol{15} Δύο (στο αμβλυγώνιο).
\sol{16} Δεν υπάρχει ισόπλευρο ορθογώνιο ή αμβλυγώνιο. Υπάρχουν: ορθογώνιο σκαληνό, ορθογώνιο ισοσκελές, αμβλυγώνιο σκαληνό, αμβλυγώνιο ισοσκελές, οξυγώνιο σκαληνό, οξυγώνιο ισοσκελές, οξυγώνιο ισόπλευρο.
\sol{17} Ελεύθερη απάντηση.
\sol{18} Ελεύθερη κατασκευή.

\lyseis{Β3.2}{Μεσοκάθετος ευθύγραμμου τμήματος}
\sol{1} Ελεύθερη κατασκευή. Και με τους δύο τρόπους η ευθεία περνά από το σημείο που απέχει $3$ cm από το $\mathrm{Α}$ και σχηματίζει ορθές γωνίες με το $\mathrm{ΑΒ}$. Στον τρόπο με τον διαβήτη η ακτίνα πρέπει να είναι μεγαλύτερη από $3$ cm.
\sol{2} Ελεύθερη κατασκευή. Για κάθε σημείο της μεσοκαθέτου οι δύο αποστάσεις από τα $\mathrm{Α}$ και $\mathrm{Β}$ βγαίνουν ίσες (μέχρι τα σφάλματα της μέτρησης).
\sol{3} Κάθε σημείο της μεσοκαθέτου ισαπέχει από τα άκρα, άρα $\mathrm{ΜΒ}=\mathrm{ΜΑ}=5$ cm.
\sol{4} Πάνω στη μεσοκάθετο του $\mathrm{ΑΒ}$, γιατί κάθε σημείο που ισαπέχει από τα άκρα ανήκει σε αυτήν.
\sol{5} α) Σωστό. β) Λάθος: πρέπει να είναι και κάθετη στο τμήμα. γ) Σωστό. δ) Λάθος: ισαπέχουν μόνο τα σημεία της μεσοκαθέτου.
\sol{6} Ελεύθερη κατασκευή. Το $\mathrm{ΑΒ}$ χωρίζεται σε $4$ ίσα κομμάτια: το $\mathrm{Μ}$ το κόβει στη μέση και τα μέσα των $\mathrm{ΑΜ}$ και $\mathrm{ΜΒ}$ κόβουν το καθένα από τα δύο μισά στη μέση.
\sol{7} Ελεύθερη κατασκευή. Οι τρεις αποστάσεις βγαίνουν ίσες: το σημείο τομής ανήκει και στις δύο μεσοκαθέτους, άρα ισαπέχει από $\mathrm{Α}$, $\mathrm{Β}$ και από $\mathrm{Β}$, $\Gamma$, δηλαδή και από τα τρία σημεία.
\sol{8} Ελεύθερη κατασκευή. Γράφουμε δύο κύκλους ακτίνας $5$ cm με κέντρα $\mathrm{Α}$ και $\mathrm{Β}$. Τέμνονται σε δύο σημεία, που είναι οι δύο θέσεις. Βρίσκονται και τα δύο πάνω στη μεσοκάθετο του $\mathrm{ΑΒ}$, από τις δύο μεριές του $\mathrm{ΑΒ}$, σε ίση απόσταση από αυτό.
\sol{9} Φέρνουμε τη μεσοκάθετο του $\mathrm{ΑΒ}$. Το σημείο όπου τέμνει τον δρόμο $\varepsilon$ ισαπέχει από τα δύο χωριά, άρα εκεί ανοίγεται το πηγάδι. (Αν ο δρόμος $\varepsilon$ είναι παράλληλος στη μεσοκάθετο δεν υπάρχει τέτοιο σημείο.)
\sol{10} Φέρνουν τις μεσοκαθέτους των δύο ζευγαριών σπιτιών, π.χ. των $\mathrm{ΑΒ}$ και $\mathrm{ΒΓ}$. Το σημείο τομής τους ισαπέχει και από τα τρία σπίτια.
\sol{11} Φέρνει τις μεσοκαθέτους των $\mathrm{ΑΒ}$ και $\mathrm{ΒΓ}$. Το σημείο τομής τους ισαπέχει από τα τρία σημεία του χείλους, άρα είναι το κέντρο του καπακιού.
\sol{12} Πάνω στη μεσοκάθετο του τμήματος $\mathrm{ΑΒ}$ (σε οποιοδήποτε σημείο της).
\sol{13} Αν η ακτίνα ήταν μικρότερη από το μισό του $\mathrm{ΑΒ}$ οι δύο κύκλοι δεν θα έφταναν ο ένας τον άλλο και δεν θα είχαν κοινό σημείο. Αν ήταν ακριβώς το μισό, θα εφάπτονταν μόνο στο μέσο.
\sol{14} Η πτυχή είναι η μεσοκάθετος του $\mathrm{ΑΒ}$: αφού το $\mathrm{Α}$ πέφτει πάνω στο $\mathrm{Β}$, η πτυχή είναι άξονας συμμετρίας του τμήματος, άρα περνά από το μέσο του και είναι κάθετη σε αυτό.
\sol{15} Αν $\mathrm{ΑΒ}=6$ cm: δύο σημεία (οι δύο κύκλοι ακτίνας $5$ τέμνονται σε δύο σημεία, γιατί $5>3$). Αν $\mathrm{ΑΒ}=10$ cm: ένα σημείο, το μέσο του $\mathrm{ΑΒ}$ (οι κύκλοι εφάπτονται). Αν $\mathrm{ΑΒ}=12$ cm: κανένα (οι κύκλοι δεν τέμνονται).
\sol{16} Η $\mathrm{ΚΛ}$ είναι η μεσοκάθετος του $\mathrm{ΑΒ}$: περνά από το μέσο του και είναι κάθετη σε αυτό.
\sol{17} Ελεύθερη κατασκευή. Η πτυχή που προκύπτει πρέπει να συμπίπτει με τη μεσοκάθετο που σχεδίασες.
\sol{18} Ελεύθερη κατασκευή. Σχεδιάζουμε δύο χορδές $\mathrm{ΑΒ}$ και $\mathrm{ΒΓ}$ και τις μεσοκαθέτους τους. Το σημείο τομής είναι το κέντρο.

\lyseis{Β3.3}{Παραλληλόγραμμα και τραπέζια}
\sol{1} α) ορθογώνιο β) ρόμβος γ) τετράγωνο.
\sol{2} α) Σωστό. β) Λάθος. γ) Σωστό. δ) Λάθος.
\sol{3} Ρόμβος: $4\cdot6{,}5=26$ cm. Τετράγωνο: $4\cdot4=16$ cm.
\sol{4} Ελεύθερη κατασκευή.
\sol{5} Ελεύθερη κατασκευή. Το ύψος είναι κάθετο στις βάσεις.
\sol{6} Όλες οι πλευρές είναι ίσες, άρα ρόμβο (και ειδικά τετράγωνο, αν οι γωνίες είναι ορθές).
\sol{7} Δύο ισοσκελή τρίγωνα, αυτά που έχουν βάσεις τις βάσεις του τραπεζίου και κορυφή το σημείο τομής των διαγωνίων.
\sol{8} Τετράγωνο. Περίμετρος: $4\cdot8=32$ μ.
\sol{9} $5\cdot(4\cdot2)=40$ μ.
\sol{10} $120+80+50+50=300$ cm $=3$ μ.
\sol{11} $120:10=12$ ρόμβοι.
\sol{12} Ελεύθερη κατασκευή.
\sol{13} Τα τετράγωνα.
\sol{14} Το ορθογώνιο όταν όλες οι πλευρές είναι ίσες, ο ρόμβος όταν όλες οι γωνίες είναι ορθές.
\sol{15} Έχει μόνο δύο παράλληλες πλευρές, ενώ το παραλληλόγραμμο έχει και τα δύο ζεύγη.
\sol{16} Τετράπλευρο $\to$ παραλληλόγραμμο $\to$ ορθογώνιο και ρόμβος $\to$ τετράγωνο. Το τραπέζιο είναι ξεχωριστά κάτω από το τετράπλευρο.
\sol{17} Ελεύθερη απάντηση.
\sol{18} Ελεύθερη κατασκευή.

\lyseis{Β3.4}{Είδη τριγώνων και στοιχεία τους}
\sol{1} α) Σωστό. β) Λάθος: για να είναι οξυγώνιο πρέπει και οι τρεις γωνίες να είναι οξείες. γ) Λάθος: δύο πλευρές ίσες, όχι όλες. δ) Σωστό: όλες οι πλευρές είναι ίσες, οποιαδήποτε μπορεί να είναι βάση. ε) Λάθος: έχει το πολύ μία αμβλεία γωνία.
\sol{2} α) ισόπλευρο \quad β) ισοσκελές \quad γ) σκαληνό.
\sol{3} α) οξυγώνιο \quad β) ορθογώνιο \quad γ) αμβλυγώνιο.
\sol{4} Ελεύθερη κατασκευή. Στο οξυγώνιο και τα τρία ύψη είναι εντός του τριγώνου. Στο ορθογώνιο οι δύο κάθετες πλευρές είναι ύψη. Στο αμβλυγώνιο δύο ύψη βρίσκονται εκτός του τριγώνου.
\sol{5} Ελεύθερη κατασκευή. Οι τρεις διάμεσοι περνούν από το ίδιο σημείο, που βρίσκεται μέσα στο τρίγωνο.
\sol{6} Ελεύθερη κατασκευή. Οι δύο γωνίες είναι ίσες, γιατί η διχοτόμος χωρίζει τη γωνία στα δύο.
\sol{7} Το $\mathrm{Μ}$ είναι το μέσο της $\mathrm{ΒΓ}$, άρα $\mathrm{ΒΓ}=2\cdot2{,}6=5{,}2$ cm.
\sol{8} Ισοσκελές (δύο ίσες πλευρές), με βάση $40$ cm. Περίμετρος $30+30+40=100$ cm.
\sol{9} Περίμετρος $3\cdot60=180$ cm.
\sol{10} Η πλευρά είναι $24:3=8$ μ.
\sol{11} Οι δύο ίσες πλευρές έχουν μαζί $50-20=30$ cm, άρα η καθεμία είναι $30:2=15$ cm.
\sol{12} Βάση είναι η πλευρά που δεν είναι ίση με τις άλλες, δηλαδή τα $8$ μ. Περίμετρος $5{,}2+5{,}2+8=18{,}4$ μ.
\sol{13} Υπάρχουν: σκαληνό οξυγώνιο, σκαληνό ορθογώνιο, σκαληνό αμβλυγώνιο, ισοσκελές οξυγώνιο, ισοσκελές ορθογώνιο, ισοσκελές αμβλυγώνιο και ισόπλευρο οξυγώνιο. Ισόπλευρο ορθογώνιο ή αμβλυγώνιο δεν υπάρχει, γιατί όλες οι γωνίες του ισόπλευρου είναι ίσες (και $60^\circ$), άρα όλες οξείες. Ελεύθερα σχέδια για τα παραδείγματα.
\sol{14} Σε οξυγώνιο, σε ορθογώνιο ή σε αμβλυγώνιο (π.χ. ισοσκελές με γωνία βάσης $45^\circ$ είναι ορθογώνιο).
\sol{15} Όχι. Ένα ισόπλευρο τρίγωνο έχει όλες τις γωνίες ίσες (και ίσες με $60^\circ$), άρα δεν μπορεί να έχει ορθή γωνία.
\sol{16} Δύο, στο αμβλυγώνιο τρίγωνο. Στο οξυγώνιο κανένα, στο ορθογώνιο κανένα (οι κάθετες πλευρές είναι ύψη πάνω στις πλευρές).
\sol{17} Ελεύθερη απάντηση. Π.χ. φέτα πίτσας: ισοσκελές, πινακίδα «Δώσε προτεραιότητα»: ισόπλευρο.
\sol{18} Ελεύθερη κατασκευή. Διπλώνουμε ώστε οι δύο άκρες της πλευράς να συμπέσουν: η πτυχή τέμνει την πλευρά στο μέσο της. Η διάμεσος ενώνει το μέσο με την απέναντι κορυφή.

\lyseisE{Επανάληψη κεφαλαίου Β3}
\solr{Β3}{1} $25^\circ$ οξεία, $90^\circ$ ορθή, $135^\circ$ αμβλεία, $180^\circ$ ευθεία, $300^\circ$ μη κυρτή, $360^\circ$ πλήρης.
\solr{Β3}{2} $\mathrm{ΜΒ}=7$ cm, γιατί κάθε σημείο της μεσοκαθέτου ισαπέχει από τα άκρα. Σχεδιάζουμε δύο κύκλους ίδιας ακτίνας (μεγαλύτερης από το μισό του $\mathrm{ΑΒ}$) με κέντρα $\mathrm{Α}$, $\mathrm{Β}$, και φέρνουμε την ευθεία των δύο σημείων τομής τους.
\solr{Β3}{3} $x\hat{\mathrm{Ο}}\delta=112^\circ:2=56^\circ$. Τα σημεία της διχοτόμου ισαπέχουν από τις πλευρές, άρα απέχει $3$ cm και από την $\mathrm{Ο}y$.
\solr{Β3}{4} Ισοσκελές (δύο ίσες πλευρές). Περίμετρος $9+9+12=30$ cm.
\solr{Β3}{5} $\hat{\Gamma}=\hat{\mathrm{Β}}=65^\circ$. Η διάμεσος προς τη βάση: $\mathrm{ΒΜ}=10:2=5$ cm, και είναι και ύψος: $\mathrm{ΑΜ}\hat{\mathrm{Β}}=90^\circ$.
\solr{Β3}{6} $64^\circ$ (κατακορυφήν), και δύο γωνίες $180^\circ-64^\circ=116^\circ$.
\solr{Β3}{7} Συμπληρωματική: $90^\circ-38^\circ=52^\circ$. Παραπληρωματική: $180^\circ-38^\circ=142^\circ$.
\solr{Β3}{8} Είναι παράλληλες (κάθετες στην ίδια ευθεία). Από σημείο εκτός της $\varepsilon_1$ περνά μία μόνο παράλληλη προς αυτήν.
\solr{Β3}{9} Οι οξείες γωνίες είναι όλες $70^\circ$ και οι αμβλείες όλες $180^\circ-70^\circ=110^\circ$ (τέσσερις γωνίες η καθεμία).
\solr{Β3}{10} $180^\circ-(58^\circ+64^\circ)=58^\circ$. Οξυγώνιο (και ισοσκελές).
\solr{Β3}{11} $(180^\circ-36^\circ):2=72^\circ$ η καθεμία.
\solr{Β3}{12} $x+(x+20^\circ)=90^\circ$, $2x=70^\circ$, $x=35^\circ$. Οι γωνίες είναι $35^\circ$ και $55^\circ$.
\solr{Β3}{13} Περίμετρος $2\cdot(9+6)=30$ cm. Οι γωνίες είναι $75^\circ$, $105^\circ$, $75^\circ$, $105^\circ$.
\solr{Β3}{14} Οι γωνίες της μικρής βάσης είναι $180^\circ-58^\circ=122^\circ$ η καθεμία. Άρα $58^\circ$, $58^\circ$, $122^\circ$, $122^\circ$.
\solr{Β3}{15} Περίμετρος $4\cdot7{,}5=30$ cm. Οι διαγώνιοι διχοτομούνται και τέμνονται κάθετα.
\solr{Β3}{16} $\mathrm{ΒΔ}=15$ cm (ίσες διαγώνιοι). Οι διαγώνιοι διχοτομούνται, άρα η απόσταση είναι $15:2=7{,}5$ cm.
\solr{Β3}{17} Διάμετρος $2\cdot4=8$ cm. Το $\mathrm{Α}$ απέχει $5>4$ από το κέντρο, άρα είναι έξω από τον κύκλο.
\solr{Β3}{18} Για $3$ cm: δύο κοινά σημεία (τέμνουσα). Για $5$ cm: ένα (εφαπτομένη). Για $7$ cm: κανένα (εξωτερική).
\solr{Β3}{19} α) Λάθος. β) Σωστό. γ) Σωστό. δ) Σωστό. ε) Λάθος: δύο ύψη βρίσκονται εκτός.
\solr{Β3}{20} α) Σωστό. β) Λάθος: οι διαγώνιοι του ρόμβου είναι κάθετες, όχι πάντα ίσες. γ) Σωστό. δ) Λάθος: δύο μόνο. ε) Σωστό.


\end{document}
